\pdfmapfile{+kpfonts.map}
\documentclass[10pt]{book}
\usepackage{makeidx}
\usepackage{framed}
\setlength{\FrameSep}{0mm}
 
\makeindex
\usepackage{kpfonts}

\newcommand{\eq}{=}



\usepackage{amsopn}
\usepackage{amsmath}
\usepackage{amsthm}
\usepackage{url}
\usepackage{graphicx}
\usepackage{datetime}
\usepackage{amsfonts}
\usepackage{graphicx}
\usepackage{threeparttable}
\usepackage{wasysym}
\usepackage{emptypage}
\usepackage{titling}



\newcommand{\cpponly}[1]{}\newcommand{\javaonly}[1]{#1}\newcommand{\pcodeonly}[1]{}\newcommand{\notpcode}[1]{#1}\newcommand{\lang}{Java}

\allowdisplaybreaks[1]  % for ams math align environments

\hyphenation{Array-Stack}
\hyphenation{Fast-Array-Stack}
\hyphenation{Array-Queue}
\hyphenation{Array-Deque}
\hyphenation{Dual-Array-Deque}
\hyphenation{Root-ish-Array-Stack}
\hyphenation{Skip-list-Set}
\hyphenation{Skip-list-List}
\hyphenation{Hash-Table}
\hyphenation{Chained-Hash-Table}
\hyphenation{Linear-Hash-Table}
\hyphenation{Red-Black-Tree}
\hyphenation{Binary-Tree}
\hyphenation{Binary-Search-Tree}
\hyphenation{Scape-goat-Tree}
\hyphenation{Count-down-Tree}
\hyphenation{Dy-na-mite-Tree}
\hyphenation{Binary-Heap}
\hyphenation{Meld-able-Heap}
\hyphenation{Java-Script}

\usepackage{everysel}
\EverySelectfont{%
\hyphenchar\font=`\-% to allow hyphenation
}

\usepackage[sf,small,raggedright]{titlesec} % formatting titles
\titlespacing*{\section}{0pt}{24pt}{14pt}
\titlespacing*{\subsection}{0pt}{14pt}{14pt}
\usepackage{relsize,fancyvrb}  % formatting pseudocode
\usepackage{ods} % Personalization and commands

\newcommand{\javaimport}[1]{}
\newcommand{\cppimport}[1]{}
\newcommand{\pcodeimport}[1]{}

\htmlonly{
  \newcommand{\ScaleIfNeeded}{\textwidth}
  \newcommand{\HalfScaleIfNeeded}{\textwidth}
  \newcommand{\HeightScaleIfNeeded}{\textheight}
  \newcommand{\QuarterHeightScaleIfNeeded}{.25\textheight}
  \newcommand{\FifthHeightScaleIfNeeded}{.2\textheight}
  \newcommand{\fancyhead}[2][zzz]{}
  \newcommand{\fancyfoot}[2][zzz]{}
}

\newcommand{\chaplabel}[1]{\label{chap:#1}}
\newcommand{\Chapref}[1]{Chapter~\ref{chap:#1}}
\newcommand{\chapref}[1]{Chapter~\ref{chap:#1}}
\newcommand{\seclabel}[1]{\label{sec:#1}}
\newcommand{\Secref}[1]{Section~\ref{sec:#1}}
\newcommand{\secref}[1]{Section~\ref{sec:#1}}
\newcommand{\sref}[1]{\textsection~\ref{sec:#1}}

\newcommand{\alglabel}[1]{\label{alg:#1}}
\newcommand{\Algref}[1]{Algorithm~\ref{alg:#1}}
\newcommand{\algref}[1]{Algorithm~\ref{alg:#1}}

\newcommand{\applabel}[1]{\label{app:#1}}
\newcommand{\Appref}[1]{Appendix~\ref{app:#1}}
\newcommand{\appref}[1]{Appendix~\ref{app:#1}}

\newcommand{\tablabel}[1]{\label{tab:#1}}
\newcommand{\Tabref}[1]{Table~\ref{tab:#1}}
\newcommand{\tabref}[1]{Table~\ref{tab:#1}}

\newcommand{\figlabel}[1]{\label{fig:#1}}
\newcommand{\Figref}[1]{Figure~\ref{fig:#1}}
\newcommand{\figref}[1]{Figure~\ref{fig:#1}}

\newcommand{\eqlabel}[1]{\label{eq:#1}}
\newcommand{\myeqref}[1]{(\ref{eq:#1})}
\newcommand{\Eqref}[1]{Equation~(\ref{eq:#1})}

\theoremstyle{plain}
\newtheorem{thm}{Theorem}[chapter]
\newcommand{\thmlabel}[1]{\label{thm:#1}}
\newcommand{\thmref}[1]{Theorem~\ref{thm:#1}}

\newtheorem{lem}{Lemma}[chapter]
\newcommand{\lemlabel}[1]{\label{lem:#1}}
\newcommand{\lemref}[1]{Lemma~\ref{lem:#1}}

\newtheorem{cor}{Corollary}[chapter]
\newcommand{\corlabel}[1]{\label{cor:#1}}
\newcommand{\corref}[1]{Corollary~\ref{cor:#1}}

\theoremstyle{definition}

\newtheorem{exc}{Exercise}[chapter]
\newcommand{\exclabel}[1]{\label{exc:#1}}
\newcommand{\excref}[1]{Exercise~\ref{exc:#1}}


\newtheorem{prp}{Property}[chapter]
\newcommand{\prplabel}[1]{\label{prp:#1}}
\newcommand{\prpref}[1]{Property~\ref{prp:#1}}

\newcommand{\etal}{\emph{et al.}}
\newcommand{\voronoi}{Vorono\u\i}
\newcommand{\ceil}[1]{{\lceil #1 \rceil}}
\newcommand{\Ceil}[1]{{\left\lceil #1 \right\rceil}}
\newcommand{\floor}[1]{{\lfloor #1 \rfloor}}
\newcommand{\Floor}[1]{{\left\lfloor #1 \right\rfloor}}
\newcommand{\R}{\mathbb{R}}
\newcommand{\N}{\mathbb{N}}
\newcommand{\Z}{\mathbb{Z}}
\newcommand{\Sp}{\mathbb{S}}
\newcommand{\E}{\mathrm{E}}
\DeclareMathOperator{\ddiv}{div}

\usepackage{ods-colors}

\usepackage{tikz,gnuplot-lua-tikz}

\def\gpsetdashtype#1{}  

\usepackage[bookmarks]{hyperref}
\hypersetup{colorlinks=true, linkcolor=linkblue,  anchorcolor=linkblue,%
	citecolor=linkblue, filecolor=linkblue, menucolor=linkblue,%
	urlcolor=linkblue,%
    pdfauthor={Pat Morin},%
    pdftitle={Open Data Structures},%
    pdfsubject={Computer Science, Data Structures},%
    pdfkeywords={Data structures, algorithms}} 

\DeclareMathOperator{\bdiv}{div}

\title{Open Data Structures (in \lang)}
\author{Pat Morin}
\date{%
Edition 0.1G\cpponly{$\beta$}\pcodeonly{$\beta$}
\htmlonly{\\ \includegraphics[scale=0.90909,scale=0.5]{images/cc-by}}}

\pagenumbering{roman}


\begin{document}

\htmlonly{\newcommand{\thetitlepage}{
  \begin{center}\includegraphics[scale=0.90909]{images/tree3-thick}\end{center}
  \maketitle
}}
\thetitlepage

\cleardoublepage

\setlength{\baselineskip}{14pt}




\addtocontents{toc}{\protect\thispagestyle{empty}} % get rid of page number
\tableofcontents
\cleardoublepage

\fancyhead[RO,LE]{} % disable section numbers, for now
\pagestyle{fancy}
\clearpage
\chapter*{Acknowledgments}
\addcontentsline{toc}{chapter}{Acknowledgments}

I am grateful to Nima~Hoda, who spent a summer tirelessly proofreading
many of the chapters in this book; to the students in the Fall 2011
offering of COMP2402/2002, who put up with the first draft of this book
and spotted many typographic, grammatical, and factual errors; and to
Morgan~Tunzelmann at Athabasca University Press, for patiently editing
several near-final drafts.

\clearpage

\thispagestyle{empty}
\cleardoublepage

\fancyhead[CE]{\small Why This Book?} % chapter title, left center
\clearpage
\chapter*{Why This Book?}
\addcontentsline{toc}{chapter}{Why This Book?}

There are plenty of books that teach introductory data structures.
Some of them are very good.  Most of them cost money, and the vast
majority of computer science undergraduate students will shell out at
least some cash on a data structures book.

Several free data structures books are available online.  Some are very
good, but most of them are getting old.  The majority of these books
became free when their authors and/or publishers decided to stop updating
them.  Updating these books is usually not possible, for two reasons:
(1)~The copyright belongs to the author and/or publisher, either of whom
may not allow it.  (2)~The \emph{source code} for these books is often
not available.  That is, the Word, WordPerfect, FrameMaker, or \LaTeX\
source for the book is not available, and even the version of the software
that handles this source may not be available.

The goal of this project is to free undergraduate computer science
students from having to pay for an introductory data structures book.
I have decided to implement this goal by treating this book like an
Open Source
\index{Open Source}
software project.  The \LaTeX\ source, \lang\ source, and
build scripts for the book are available to download from the
author's website\footnote{\url{http://opendatastructures.org}}
and also, more importantly, on a reliable source code management
site.\footnote{\url{https://github.com/patmorin/ods}}

The source code available there is released under a Creative Commons
Attribution license, meaning that anyone is free to \emph{share}:
\index{share}
to copy, distribute and
transmit the work; and to \emph{remix}:
\index{remix}
to adapt the work, including the
right to make commercial use of the work.  The only condition on these
rights is \emph{attribution}: you must acknowledge that the derived work
contains code and/or text from \url{opendatastructures.org}.

Anyone can contribute corrections/fixes using the \texttt{git}
\index{git@\texttt{git}}
source-code management system.  Anyone can also fork the book's sources to
develop a separate version (for example, in another programming language).
My hope is that, by doing things this way, this book will continue to
be a useful textbook long after my interest in the project, or my pulse,
(whichever comes first) has waned.



\clearpage

\cleardoublepage

\cpponly{
  \clearpage
\cpponly{
\chapter*{Preface to the C++ Edition}
\addcontentsline{toc}{chapter}{Preface to the C++ Edition}

This book is intended to teach the design and analysis of basic data
structures and their implementation in an object-oriented language.
In this edition, the language happens to be C++.

This book is not intended to act as an introduction to the C++ programming
language.  Readers of this book need only be familiar with the basic
syntax of C++ and similar languages.  Those wishing to work with the
accompanying source code should have some experience programming in C++.

This book is also not intended as an introduction to the C++ Standard
Template Library or the generic programming paradigm that the STL
embodies.  This book describes implementations of several different data
structures, many of which are used in implementations of the STL. The
contents of this book may help an STL programmer understand how some of
the STL data structures are implemented and why these implementations
are efficient.
}

\clearpage

  \cleardoublepage
}

\fancyhead[CE]{\small\nouppercase{\leftmark}} % chapter title, left center
\fancyhead[CO]{\small\rightmarktitle} % section title, right center
\fancyhead[RO,LE]{\small\rightmarksection}

\clearpage
\chapter{Introduction}
\pagenumbering{arabic}

Every computer science curriculum in the world includes a course on data
structures and algorithms.  Data structures are \emph{that} important;
they improve our quality of life and even save lives on a regular basis.
Many multi-million and several multi-billion dollar companies have been
built around data structures.

How can this be?  If we stop to think about it, we realize that we
interact with data structures constantly.
\begin{itemize}
  \item  Open a file: File system 
    \index{file system}%
    data structures are used to locate
    the parts of that file on disk so they can be retrieved.  This isn't
    easy; disks contain hundreds of millions of blocks.  The contents of your
    file could be stored on any one of them.
  \item Look up a contact on your phone:  A data
    structure is used to look up a phone number in your contact list
    \index{contact list}%
    based on partial
    information even before you finish dialing/typing.  This isn't easy;
    your phone may contain information about a lot of people---everyone
    you have ever contacted via phone or email---and your phone doesn't
    have a very fast processor or a lot of memory.
  \item Log in to your favourite social network:
    \index{social network}%
    The network servers
    use your login information to look up your account information.
    This isn't easy; the most popular social networks have hundreds of
    millions of active users.
  \item Do a web search:
    \index{web search}%
    The search engine uses data structures to find
    the web pages containing your search terms.  This isn't easy; there
    are over 8.5 billion web pages on the Internet and each page contains
    a lot of potential search terms.
  \item Phone emergency services (9-1-1):
    \index{emergency services}\index{9-1-1}%
    The emergency services network
    looks up your phone number in a data structure that maps phone numbers
    to addresses so that police cars, ambulances, or fire trucks can be
    sent there without delay. This is important;  the person making the
    call may not be able to provide the exact address they are calling
    from and a delay can mean the difference between life or death.
\end{itemize}

\section{The Need for Efficiency}

In the next section, we look at the operations supported by the most
commonly used data structures.  Anyone with a bit of programming
experience will see that these operations are not hard to implement
correctly.  We can store the data in an array or a linked list and each
operation can be implemented by iterating over all the elements of the
array or list and possibly adding or removing an element.

This kind of implementation is easy, but not very efficient.  Does this
really matter?  Computers are becoming faster and faster. Maybe the
obvious implementation is good enough. Let's do some rough calculations
to find out.

\paragraph{Number of operations:}  Imagine an application with a
moderately-sized data set, say of one million ($10^6$), items.  It is
reasonable, in most applications, to assume that the application will
want to look up each item at least once.  This means we can expect to do
at least one million ($10^6$) searches in this data.  If each of these
$10^6$ searches inspects  each of the $10^6$ items, this gives a total
of $10^6\times 10^6=10^{12}$ (one thousand billion) inspections.

\paragraph{Processor speeds:} At the time of writing, even a very fast
desktop computer can not do more than one billion ($10^9$) operations per
second.\footnote{Computer speeds are at most a few gigahertz (billions
of cycles per second), and each operation typically takes a few cycles.}
This means that this application will take at least $10^{12}/10^9 = 1000$
seconds, or roughly 16 minutes and 40 seconds.  Sixteen minutes is an
eon in computer time, but a person might be willing to put up with it
(if he or she were headed out for a coffee break).


\paragraph{Bigger data sets:} Now consider a company like Google,
\index{Google}%
that
indexes over 8.5 billion web pages.  By our calculations, doing any kind
of query over this data would take at least 8.5 seconds.  We already
know that this isn't the case; web searches complete in much less than
8.5 seconds, and they do much more complicated queries than just asking
if a particular page is in their list of indexed pages.  At the time
of writing, Google receives approximately $4,500$ queries per second,
meaning that they would require at least $4,500\times 8.5=38,250$ very
fast servers just to keep up.

\paragraph{The solution:} 
These examples tell us that the obvious implementations of data structures
do not scale well when the number of items, \ensuremath{\mathtt{{\color{var}n}}}, in the data structure
and the number of operations, $m$, performed on the data structure
are both large.  In these cases, the time (measured in, say, machine
instructions) is roughly $\ensuremath{\mathtt{{\color{var}n}}}\times m$.

The solution, of course, is to carefully organize data within the data
structure so that not every operation requires every data item to be
inspected.  Although it sounds impossible at first, we will see data
structures where a search requires looking at only two items on average,
independent of the number of items stored in the data structure.  In our
billion instruction per second computer it takes only $0.000000002$
seconds to search in a data structure containing a billion items (or a
trillion, or a quadrillion, or even a quintillion items).

We will also see implementations of data structures that keep the
items in sorted order, where the number of items inspected during an
operation grows very slowly as a function of the number of items in
the data structure.  For example, we can maintain a sorted set of one
billion items while inspecting at most 60 items during any operation.
In our billion instruction per second computer, these operations take
$0.00000006$ seconds each.

The remainder of this chapter briefly reviews some of the main concepts
used throughout the rest of the book.   \secref{interfaces} describes
the interfaces implemented by all of the data structures described in
this book and should be considered required reading.  The remaining
sections discuss:
\begin{itemize}
  \item some mathematical review including exponentials, logarithms,
  factorials, asymptotic (big-Oh) notation, probability, and randomization;
  \item the model of computation; 
  \item correctness, running time, and space;
  \item an overview of the rest of the chapters; and
  \item the sample code and typesetting conventions.
\end{itemize}
A reader with or without a background in these areas can easily skip
them now and come back to them later if necessary.


\section{Interfaces}
\seclabel{interfaces}
When discussing data structures, it is important to understand the
difference between a data structure's interface and its implementation.
An interface describes what a data structure does, while an implementation
describes how the data structure does it.

An \emph{interface},
\index{interface}%
\index{abstract data type|see{interface}}%
sometimes also called an \emph{abstract data
type}, defines the set of operations supported by a data structure and
the semantics, or meaning, of those operations.  An interface tells us
nothing about how the data structure implements these operations; it only
provides a list of supported operations along with specifications about
what types of arguments each operation accepts and the value returned
by each operation.

A data structure \emph{implementation}, on the other hand, includes the
internal representation of the data structure as well as the definitions
of the algorithms that implement the operations supported by the data
structure.  Thus, there can be many implementations of a single interface.
For example, in \chapref{arrays}, we will see implementations of the
\texttt{List} interface using arrays and in \chapref{linkedlists} we will
see implementations of the \texttt{List} interface using pointer-based data
structures.  Each implements the same interface, \texttt{List},
but in different ways.

\subsection{The \texttt{Queue}, \texttt{Stack}, and \texttt{Deque} Interfaces}

The \texttt{Queue} interface represents a collection of elements to which we
can add elements and remove the next element.  More precisely, the operations
supported by the \texttt{Queue} interface are
\begin{itemize}
  \item \ensuremath{\mathtt{add({\color{var}x})}}: add the value \ensuremath{\mathtt{{\color{var}x}}} to the \texttt{Queue}
  \item \ensuremath{\mathtt{remove()}}: remove the next (previously added) value, \ensuremath{\mathtt{{\color{var}y}}}, from the \texttt{Queue} and return \ensuremath{\mathtt{{\color{var}y}}}
\end{itemize}
Notice that the \ensuremath{\mathtt{remove()}} operation takes no argument.  The \texttt{Queue}'s
\emph{queueing discipline} decides which element should be removed.
There are many possible queueing disciplines, the most common of which
include FIFO, priority, and LIFO.

A \emph{FIFO (first-in-first-out) \texttt{Queue}},
\index{FIFO queue}%
\index{queue!FIFO}%
which is illustrated in
\figref{queue}, removes items in the same order they were added, much
in the same way a queue (or line-up) works when checking out at a cash
register in a grocery store.  This is the most common kind of \texttt{Queue}
so the qualifier FIFO is often omitted.  In other texts, the \ensuremath{\mathtt{add({\color{var}x})}}
and \ensuremath{\mathtt{remove()}} operations on a FIFO \texttt{Queue} are often called \ensuremath{\mathtt{enqueue({\color{var}x})}}
and \ensuremath{\mathtt{dequeue()}}, respectively.

\begin{figure}
  \centering{\includegraphics[width=\ScaleIfNeeded]{figs/queue}}
  \caption[A FIFO queue]{A FIFO \texttt{Queue}.}
  \figlabel{queue}
\end{figure}

A \emph{priority \texttt{Queue}},
\index{priority queue}%
\index{priority queue|seealso{heap}}%
\index{queue!priority}%
illustrated in \figref{prioqueue}, always
removes the smallest element from the \texttt{Queue}, breaking ties arbitrarily.
This is similar to the way in which patients are triaged in a hospital
emergency room.  As patients arrive they are evaluated and then placed in
a waiting room.  When a doctor becomes available he or she first treats
the patient with the most life-threatening condition.  The \ensuremath{\mathtt{remove()}}
operation on a priority \texttt{Queue} is usually called \ensuremath{\mathtt{deleteMin()}} in
other texts.

\begin{figure}
  \centering{\includegraphics[width=\ScaleIfNeeded]{figs/prioqueue}}
  \caption[A priority queue]{A priority \texttt{Queue}.}
  \figlabel{prioqueue}
\end{figure}


A very common queueing discipline is the LIFO (last-in-first-out)
\index{LIFO queue}%
\index{LIFO queue|seealso{stack}}%
\index{queue!LIFO}%
\index{stack}%
discipline, illustrated in \figref{stack}.  In a \emph{LIFO Queue},
the most recently added element is the next one removed.  This is best
visualized in terms of a stack of plates; plates are placed on the top of
the stack and also removed from the top of the stack. This structure is
so common that it gets its own name: \texttt{Stack}.  Often, when discussing a
\texttt{Stack}, the names of \ensuremath{\mathtt{add({\color{var}x})}} and \ensuremath{\mathtt{remove()}} are changed to \ensuremath{\mathtt{push({\color{var}x})}}
and \ensuremath{\mathtt{pop()}}; this is to avoid confusing the LIFO and FIFO queueing
disciplines.

\begin{figure}
  \centering{\includegraphics[width=\ScaleIfNeeded]{figs/stack}}
  \caption[A stack]{A stack.}
  \figlabel{stack}
\end{figure}


A \texttt{Deque}
\index{deque}%
is a generalization of both the FIFO \texttt{Queue} and LIFO \texttt{Queue} (\texttt{Stack}).
A \texttt{Deque} represents a sequence of elements, with a front and a back.
Elements can be added at the front of the sequence or the back of the
sequence.  The names of the \texttt{Deque} operations are self-explanatory:
\ensuremath{\mathtt{addFirst({\color{var}x})}}, \ensuremath{\mathtt{removeFirst()}}, \ensuremath{\mathtt{addLast({\color{var}x})}}, and \ensuremath{\mathtt{removeLast()}}.  It is
worth noting that a \texttt{Stack} can be implemented using only \ensuremath{\mathtt{addFirst({\color{var}x})}}
and \ensuremath{\mathtt{removeFirst()}} while a FIFO \texttt{Queue} can be implemented using
\ensuremath{\mathtt{addLast({\color{var}x})}} and \ensuremath{\mathtt{removeFirst()}}.


\subsection{The \texttt{List} Interface: Linear Sequences}

This book will talk very little about the FIFO \texttt{Queue}, \texttt{Stack}, or
\texttt{Deque} interfaces.  This is because these interfaces are subsumed by the
\texttt{List} interface.  A \texttt{List},
\index{List@\texttt{List}}%
illustrated in \figref{list}, represents a
sequence, $\ensuremath{\mathtt{{\color{var}x}}}_0,\ldots,\ensuremath{\mathtt{{\color{var}x}}}_{\ensuremath{\mathtt{{\color{var}n}}}-1}$, of values.  The \texttt{List} interface
includes the following operations:

\begin{enumerate}
  \item \ensuremath{\mathtt{size()}}: return \ensuremath{\mathtt{{\color{var}n}}}, the length of the list
  \item \ensuremath{\mathtt{get({\color{var}i})}}: return the value $\ensuremath{\mathtt{{\color{var}x}}}_{\ensuremath{\mathtt{{\color{var}i}}}}$
  \item \ensuremath{\mathtt{set({\color{var}i},{\color{var}x})}}: set the value of $\ensuremath{\mathtt{{\color{var}x}}}_{\ensuremath{\mathtt{{\color{var}i}}}}$ equal to \ensuremath{\mathtt{{\color{var}x}}}
  \item \ensuremath{\mathtt{add({\color{var}i},{\color{var}x})}}: add \ensuremath{\mathtt{{\color{var}x}}} at position \ensuremath{\mathtt{{\color{var}i}}}, displacing
    $\ensuremath{\mathtt{{\color{var}x}}}_{\ensuremath{\mathtt{{\color{var}i}}}},\ldots,\ensuremath{\mathtt{{\color{var}x}}}_{\ensuremath{\mathtt{{\color{var}n}}}-1}$; \\ 
    Set $\ensuremath{\mathtt{{\color{var}x}}}_{j+1}=\ensuremath{\mathtt{{\color{var}x}}}_j$, for all
    $j\in\{\ensuremath{\mathtt{{\color{var}n}}}-1,\ldots,\ensuremath{\mathtt{{\color{var}i}}}\}$, increment \ensuremath{\mathtt{{\color{var}n}}}, and set $\ensuremath{\mathtt{{\color{var}x}}}_i=\ensuremath{\mathtt{{\color{var}x}}}$
  \item \ensuremath{\mathtt{remove({\color{var}i})}} remove the value $\ensuremath{\mathtt{{\color{var}x}}}_{\ensuremath{\mathtt{{\color{var}i}}}}$, displacing
    $\ensuremath{\mathtt{{\color{var}x}}}_{\ensuremath{\mathtt{{\color{var}i}+1}}},\ldots,\ensuremath{\mathtt{{\color{var}x}}}_{\ensuremath{\mathtt{{\color{var}n}}}-1}$; \\ 
    Set $\ensuremath{\mathtt{{\color{var}x}}}_{j}=\ensuremath{\mathtt{{\color{var}x}}}_{j+1}$, for all
    $j\in\{\ensuremath{\mathtt{{\color{var}i}}},\ldots,\ensuremath{\mathtt{{\color{var}n}}}-2\}$ and decrement \ensuremath{\mathtt{{\color{var}n}}}
\end{enumerate}
Notice that these operations are easily sufficient to implement the
\texttt{Deque} interface:
\begin{eqnarray*}
  \ensuremath{\mathtt{addFirst({\color{var}x})}} &\Rightarrow& \ensuremath{\mathtt{add(0,{\color{var}x})}} \\
  \ensuremath{\mathtt{removeFirst()}} &\Rightarrow& \ensuremath{\mathtt{remove(0)}}  \\
  \ensuremath{\mathtt{addLast({\color{var}x})}} &\Rightarrow& \ensuremath{\mathtt{add(size(),{\color{var}x})}} \\
  \ensuremath{\mathtt{removeLast()}} &\Rightarrow& \ensuremath{\mathtt{remove(size()-1)}}
\end{eqnarray*}

\begin{figure}
  \centering{\includegraphics[width=\ScaleIfNeeded]{figs/list}}
  \caption[A List]{A \texttt{List} represents a sequence indexed by
   $0,1,2,\ldots,\ensuremath{\mathtt{{\color{var}n}}}-1$.  In this \texttt{List} a call to \ensuremath{\mathtt{get(2)}} would return
   the value $c$.}
  \figlabel{list}
\end{figure}

Although we will normally not discuss the \texttt{Stack}, \texttt{Deque} and FIFO
\texttt{Queue} interfaces in subsequent chapters, the terms \texttt{Stack} and \texttt{Deque}
are sometimes used in the names of data structures that implement the
\texttt{List} interface.  When this happens, it highlights the fact that
these data structures can be used to implement the \texttt{Stack} or \texttt{Deque}
interface very efficiently.  For example, the \texttt{ArrayDeque} class is an
implementation of the \texttt{List} interface that implements all the \texttt{Deque}
operations in constant time per operation.


\subsection{The \texttt{USet} Interface: Unordered Sets}

The \texttt{USet}
\index{USet@\texttt{USet}}%
interface represents an unordered set of unique elements, which
mimics a mathematical \emph{set}. A \texttt{USet} contains \ensuremath{\mathtt{{\color{var}n}}} \emph{distinct}
elements; no element appears more than once; the elements are in no
specific order.  A \texttt{USet} supports the following operations:

\begin{enumerate}
  \item \ensuremath{\mathtt{size()}}: return the number, \ensuremath{\mathtt{{\color{var}n}}}, of elements in the set
  \item \ensuremath{\mathtt{add({\color{var}x})}}: add the element \ensuremath{\mathtt{{\color{var}x}}} to the set if not already present; \\
    Add \ensuremath{\mathtt{{\color{var}x}}} to the set provided that there
    is no element \ensuremath{\mathtt{{\color{var}y}}} in the set such that \ensuremath{\mathtt{{\color{var}x}}} equals \ensuremath{\mathtt{{\color{var}y}}}.  Return \ensuremath{\mathtt{{\color{var}true}}}
    if \ensuremath{\mathtt{{\color{var}x}}} was added to the set and \ensuremath{\mathtt{{\color{var}false}}} otherwise.
  \item \ensuremath{\mathtt{remove({\color{var}x})}}: remove \ensuremath{\mathtt{{\color{var}x}}} from the set; \\
    Find an element \ensuremath{\mathtt{{\color{var}y}}} in the set such that \ensuremath{\mathtt{{\color{var}x}}} equals
    \ensuremath{\mathtt{{\color{var}y}}} and remove \ensuremath{\mathtt{{\color{var}y}}}.  Return \ensuremath{\mathtt{{\color{var}y}}}, or \ensuremath{\mathtt{{\color{keyword}null}}} if no such element exists.
  \item \ensuremath{\mathtt{find({\color{var}x})}}: find \ensuremath{\mathtt{{\color{var}x}}} in the set if it exists; \\
    Find an element \ensuremath{\mathtt{{\color{var}y}}} in the set such that \ensuremath{\mathtt{{\color{var}y}}} equals
    \ensuremath{\mathtt{{\color{var}x}}}.  Return \ensuremath{\mathtt{{\color{var}y}}}, or \ensuremath{\mathtt{{\color{keyword}null}}} if no such element exists.
\end{enumerate}

These definitions are a bit fussy about distinguishing \ensuremath{\mathtt{{\color{var}x}}}, the element
we are removing or finding, from \ensuremath{\mathtt{{\color{var}y}}}, the element we may remove or find.
This is because \ensuremath{\mathtt{{\color{var}x}}} and \ensuremath{\mathtt{{\color{var}y}}} might actually be distinct objects that
are nevertheless treated as equal.\javaonly{\footnote{In Java, this is
done by overriding the class's \ensuremath{\mathtt{equals({\color{var}y})}} and \ensuremath{\mathtt{hashCode()}} methods.}}
Such a distinction is  useful because it allows for the creation of
\emph{dictionaries} or \emph{maps} that map keys onto values.
\index{dictionary}%
\index{map}%

To create a dictionary/map, one forms compound objects called \texttt{Pair}s,
\index{pair}%
each of which contains a \emph{key} and a \emph{value}. Two \texttt{Pair}s
are treated as equal if their keys are equal.  If we store some pair
$(\ensuremath{\mathtt{{\color{var}k}}},\ensuremath{\mathtt{{\color{var}v}}})$ in a \texttt{USet} and then later call the \ensuremath{\mathtt{find({\color{var}x})}} method using
the pair $\ensuremath{\mathtt{{\color{var}x}}}=(\ensuremath{\mathtt{{\color{var}k}}},\ensuremath{\mathtt{{\color{keyword}null}}})$ the result will be $\ensuremath{\mathtt{{\color{var}y}}}=(\ensuremath{\mathtt{{\color{var}k}}},\ensuremath{\mathtt{{\color{var}v}}})$.  In other
words, it is possible to recover the value, \ensuremath{\mathtt{{\color{var}v}}}, given only the key, \ensuremath{\mathtt{{\color{var}k}}}.


\subsection{The \texttt{SSet} Interface: Sorted Sets}
\seclabel{sset}

\index{SSet@\texttt{SSet}}%
The \texttt{SSet} interface represents a sorted set of elements.  An \texttt{SSet}
stores elements from some total order, so that any two elements \ensuremath{\mathtt{{\color{var}x}}}
and \ensuremath{\mathtt{{\color{var}y}}} can be compared.  In code examples, this will be done with a
method called \ensuremath{\mathtt{compare({\color{var}x},{\color{var}y})}} in which
\[
    \ensuremath{\mathtt{compare({\color{var}x},{\color{var}y})}} 
      \begin{cases}
        {}<0 & \text{if $\ensuremath{\mathtt{{\color{var}x}}}<\ensuremath{\mathtt{{\color{var}y}}}$} \\
        {}>0 & \text{if $\ensuremath{\mathtt{{\color{var}x}}}>\ensuremath{\mathtt{{\color{var}y}}}$} \\
        {}=0 & \text{if $\ensuremath{\mathtt{{\color{var}x}}}=\ensuremath{\mathtt{{\color{var}y}}}$}
      \end{cases}
\]
\index{compare@\ensuremath{\mathtt{compare({\color{var}x},{\color{var}y})}}}%
An \texttt{SSet} supports the \ensuremath{\mathtt{size()}}, \ensuremath{\mathtt{add({\color{var}x})}}, and \ensuremath{\mathtt{remove({\color{var}x})}} methods with
exactly the same semantics as in the \texttt{USet} interface.  The difference
between a \texttt{USet} and an \texttt{SSet} is in the \ensuremath{\mathtt{find({\color{var}x})}} method:
\begin{enumerate}
\setcounter{enumi}{3}
\item \ensuremath{\mathtt{find({\color{var}x})}}: locate \ensuremath{\mathtt{{\color{var}x}}} in the sorted set; \\
   Find the smallest element \ensuremath{\mathtt{{\color{var}y}}} in the set such that $\ensuremath{\mathtt{{\color{var}y}}} \ge \ensuremath{\mathtt{{\color{var}x}}}$.
   Return \ensuremath{\mathtt{{\color{var}y}}} or \ensuremath{\mathtt{{\color{keyword}null}}} if no such element exists.
\end{enumerate}

This version of the \ensuremath{\mathtt{find({\color{var}x})}} operation is sometimes referred to
as a \emph{successor search}.
\index{successor search}%
It differs in a fundamental way from
\ensuremath{\mathtt{USet.find({\color{var}x})}} since it returns a meaningful result even when there is
no element equal to \ensuremath{\mathtt{{\color{var}x}}} in the set.

The distinction between the \texttt{USet} and \texttt{SSet} \ensuremath{\mathtt{find({\color{var}x})}} operations is
very important and often missed.  The extra functionality provided
by an \texttt{SSet} usually comes with a price that includes both a larger
running time and a higher implementation complexity.  For example, most
of the \texttt{SSet} implementations discussed in this book all have \ensuremath{\mathtt{find({\color{var}x})}}
operations with running times that are logarithmic in the size of the set.
On the other hand, the implementation of a \texttt{USet} as a \texttt{ChainedHashTable}
in \chapref{hashing} has a \ensuremath{\mathtt{find({\color{var}x})}} operation that runs in constant
expected time.  When choosing which of these structures to use, one should
always use a \texttt{USet} unless the extra functionality offered by an \texttt{SSet}
is truly needed.


\section{Mathematical Background}

In this section, we review some mathematical notations and tools
used throughout this book, including logarithms, big-Oh notation, and
probability theory.  This review will be brief and is not intended as
an introduction. Readers who feel they are missing this background are
encouraged to read, and do exercises from, the appropriate sections of
the very good (and free) textbook on mathematics for computer science
\cite{llm11}.

\subsection{Exponentials and Logarithms}

\index{exponential}%
The expression $b^x$ denotes the number $b$ raised to the power of $x$.
If $x$ is a positive integer, then this is just the value of $b$
multiplied by itself $x-1$ times:
\[
    b^x = \underbrace{b\times b\times \cdots \times b}_{x} \enspace .
\]
When $x$ is a negative integer, $b^{-x}=1/b^{x}$.  When $x=0$, $b^x=1$.
When $b$ is not an integer, we can still define exponentiation in terms
of the exponential function $e^x$ (see below), which is itself defined in
terms of the exponential series, but this is best left to a calculus text.

\index{logarithm}%
In this book, the expression $\log_b k$ denotes the \emph{base-$b$ logarithm}
of $k$.  That is, the unique value $x$ that satisfies
\[
    b^{x} = k  \enspace .
\]
Most of the logarithms in this book are base 2 (\emph{binary logarithms}).
\index{binary logarithm}%
\index{logarithm!binary}%
For these, we omit the base, so that $\log k$ is shorthand for
$\log_2 k$.

An informal, but useful, way to think about logarithms is to think
of $\log_b k$ as the number of times we have to divide $k$ by $b$
before the result is less than or equal to 1.  For example, when one
does binary search, each comparison reduces the number of possible
answers by a factor of 2.  This is repeated until there is at most one
possible answer.  Therefore, the number of comparison done by binary
search when there are initially at most $n+1$ possible answers is at
most $\lceil\log_2(n+1)\rceil$.

\index{natural logarithm}%
\index{logarithm!natural}%
Another logarithm that comes up several times in this book is the
\emph{natural logarithm}.  Here we use the notation $\ln k$ to denote
$\log_e k$, where $e$ --- \emph{Euler's constant} --- is given by
\index{Euler's constant}%
\index{e@$e$ (Euler's constant)}%
\[
   e = \lim_{n\rightarrow\infty} \left(1+\frac{1}{n}\right)^n
   \approx  2.71828 \enspace .
\]
The natural logarithm comes up frequently because it is the value
of a particularly common integral:
\[
    \int_{1}^{k} 1/x\,\mathrm{d}x  = \ln k \enspace .
\]
Two of the most common manipulations we do with logarithms are removing
them from an exponent:
\[
    b^{\log_b k} = k
\]
and changing the base of a logarithm:
\[
    \log_b k = \frac{\log_a k}{\log_a b} \enspace .
\]
For example, we can use these two manipulations to compare the natural and binary logarithms
\[
   \ln k = \frac{\log k}{\log e} = \frac{\log k}{(\ln e)/(\ln 2)} = 
    (\ln 2)(\log k) \approx 0.693147\log k \enspace .
\]

\subsection{Factorials}
\seclabel{factorials}

\index{factorial}%
In one or two places in this book, the \emph{factorial} function is used.
For a non-negative integer $n$, the notation $n!$ (pronounced ``$n$ factorial'') is defined to mean 
\[
   n! = 1\cdot2\cdot3\cdot\cdots\cdot n \enspace .
\]
Factorials appear because $n!$ counts the number of distinct
permutations, i.e., orderings, of $n$ distinct elements.
\index{permutation}%
For the special case $n=0$, $0!$ is defined as 1. 

\index{Stirling's Approximation}%
The quantity $n!$ can be approximated using \emph{Stirling's Approximation}:
\[
	n! 
   = \sqrt{2\pi n}\left(\frac{n}{e}\right)^{n}e^{\alpha(n)} \enspace ,
\]
where
\[  
   \frac{1}{12n+1} <  \alpha(n) < \frac{1}{12n}  \enspace .
\]
Stirling's Approximation also approximates $\ln(n!)$:
\[
   \ln(n!) = n\ln n - n + \frac{1}{2}\ln(2\pi n) + \alpha(n)
\]
(In fact, Stirling's Approximation is most easily proven by approximating
$\ln(n!)=\ln 1 + \ln 2  + \cdots + \ln n$ by the integral
$\int_1^n \ln n\,\mathrm{d}n = n\ln n - n +1$.)

\index{binomial coefficients}%
Related to the factorial function are the \emph{binomial coefficients}.
For a non-negative integer $n$ and an integer $k\in\{0,\ldots,n\}$,
the notation $\binom{n}{k}$ denotes:
\[
   \binom{n}{k} = \frac{n!}{k!(n-k)!} \enspace .
\]
The binomial coefficient $\binom{n}{k}$ (pronounced ``$n$ choose $k$'')
counts the number of subsets of an $n$ element set that have size $k$,
i.e., the number of ways of choosing $k$ distinct integers from the
set $\{1,\ldots,n\}$.

\subsection{Asymptotic Notation}

\index{asymptotic notation}%
\index{big-Oh notation}%
\index{O@$O$ notation}%
When analyzing data structures in this book, we want to talk about
the running times of various operations.  The exact running times will,
of course, vary from computer to computer and even from run to run on an
individual computer.  When we talk about the running time of an operation
we are referring to the number of computer instructions performed during
the operation.  Even for simple code, this quantity can be difficult to
compute exactly.  Therefore, instead of analyzing running times exactly,
we will use the so-called \emph{big-Oh notation}: For a function $f(n)$,
$O(f(n))$ denotes a set of functions,
\[
   O(f(n)) = \left\{
     \begin{array}{l}
       g(n):\mbox{there exists $c>0$, and $n_0$ such that} \\
             \quad\mbox{$g(n) \le c\cdot f(n)$ for all $n\ge n_0$}   
     \end{array} \right\} \enspace .
\]
Thinking graphically, this set consists of the functions $g(n)$ where
$c\cdot f(n)$ starts to dominate $g(n)$ when $n$ is sufficiently large.

We generally use asymptotic notation to simplify functions.  For example,
in place of $5n\log n + 8n - 200$ we can write $O(n\log n)$.
This is proven as follows:
\begin{align*} 
       5n\log n + 8n - 200
        & \le 5n\log n + 8n \\
        & \le 5n\log n + 8n\log n & \mbox{ for $n\ge 2$ (so that $\log n \ge 1$)}
            \\
        & \le 13n\log n  \enspace .
\end{align*}
This demonstrates that the function $f(n)=5n\log n + 8n - 200$ is in
the set $O(n\log n)$ using the constants $c=13$ and $n_0 = 2$.

A number of useful shortcuts can be applied when using asymptotic
notation.  First:
\[ O(n^{c_1}) \subset O(n^{c_2}) \enspace ,\]
for any $c_1 < c_2$.  Second: For any constants $a,b,c > 0$,
\[ O(a) \subset O(\log n) \subset O(n^{b}) \subset O({c}^n) \enspace . \]
These inclusion relations can be multiplied by any positive value,
and they still hold. For example, multiplying by $n$ yields:
\[ O(n) \subset O(n\log n) \subset O(n^{1+b}) \subset O(n{c}^n) \enspace . \]

Continuing in a long and distinguished tradition, we will abuse this
notation by writing things like $f_1(n) = O(f(n))$ when what we really
mean is $f_1(n) \in O(f(n))$.  We will also make statements like ``the
running time of this operation is $O(f(n))$'' when this statement should
be ``the running time of this operation is \emph{a member of} $O(f(n))$.''
These shortcuts are mainly to avoid awkward language and to make it
easier to use asymptotic notation within strings of equations.

A particularly strange example of this occurs when we write statements like
\[
  T(n) = 2\log n + O(1)  \enspace .
\]
Again, this would be more correctly written as
\[
  T(n) \le 2\log n + [\mbox{some member of $O(1)$]}  \enspace .
\]

The expression $O(1)$ also brings up another issue. Since there is
no variable in this expression, it may not be clear which variable is
getting arbitrarily large.  Without context, there is no way to tell.
In the example above, since the only variable in the rest of the equation
is $n$, we can assume that this should be read as $T(n) = 2\log n +
O(f(n))$, where $f(n) = 1$.

Big-Oh notation is not new or unique to computer science.  It was
used by the number theorist Paul Bachmann as early as 1894, and is
immensely useful for describing the running times of computer algorithms.
Consider the following piece of code:


 \ \vspace{.5ex}\begin{Verbatim}[gobble=1,tabsize=2,frame=single,framerule=1.2pt,commandchars=\\@\$,label=\texttt{Simple},labelposition=topline]
	@\color@keyword$void$ snippet() {
		@\color@keyword$for$ (@\color@keyword$int$ @\color@var$i$ = 0; @\color@var$i$ < @\color@var$n$; @\color@var$i$++) 
			@\color@var$a$[@\color@var$i$] = @\color@var$i$;
	}
\end{Verbatim}

One execution of this method involves
\begin{itemize}
      \item $1$ assignment (\ensuremath{\mathtt{{\color{keyword}int}\, {\color{var}i}\, =\, 0}}),
      \item $\ensuremath{\mathtt{{\color{var}n}}}+1$ comparisons (\ensuremath{\mathtt{{\color{var}i} < {\color{var}n}}}),
      \item \ensuremath{\mathtt{{\color{var}n}}} increments (\ensuremath{\mathtt{{\color{var}i}++}}),
      \item \ensuremath{\mathtt{{\color{var}n}}} array offset calculations (\ensuremath{\mathtt{{\color{var}a}[{\color{var}i}]}}), and
      \item \ensuremath{\mathtt{{\color{var}n}}} indirect assignments (\ensuremath{\mathtt{{\color{var}a}[{\color{var}i}] = {\color{var}i}}}).
\end{itemize}
So we could write this running time as
\[
    T(\ensuremath{\mathtt{{\color{var}n}}})=a + b(\ensuremath{\mathtt{{\color{var}n}}}+1) + c\ensuremath{\mathtt{{\color{var}n}}} + d\ensuremath{\mathtt{{\color{var}n}}} + e\ensuremath{\mathtt{{\color{var}n}}} \enspace , 
\]
where $a$, $b$, $c$, $d$, and $e$ are constants that depend on the
machine running the code and represent the time to perform assignments,
comparisons, increment operations, array offset calculations, and indirect
assignments, respectively.  However, if this expression represents the
running time of two lines of code, then clearly this kind of analysis
will not be tractable to complicated code or algorithms.  Using big-Oh
notation, the running time can be simplified to
\[
    T(\ensuremath{\mathtt{{\color{var}n}}})= O(\ensuremath{\mathtt{{\color{var}n}}}) \enspace .
\]
Not only is this more compact, but it also gives nearly as much
information.  The fact that the running time depends on the constants $a$,
$b$, $c$, $d$, and $e$ in the above example means that, in general, it
will not be possible to compare two running times to know which is faster
without knowing the values of these constants.  Even if we make the
effort to determine these constants (say, through timing tests), then
our conclusion will only be valid for the machine we run our tests on.

Big-Oh notation allows us to reason at a much higher level, making
it possible to analyze more complicated functions.  If two algorithms
have the same big-Oh running time, then we won't know which is faster,
and there may not be a clear winner.  One may be faster on one machine,
and the other may be faster on a different machine.  However, if the
two algorithms have demonstrably different big-Oh running times, then
we can be certain that the one with the smaller running time will be
faster \emph{for large enough values of \ensuremath{\mathtt{{\color{var}n}}}}.

An example of how big-Oh notation allows us to compare two different
functions is shown in \figref{intro-asymptotics}, which compares the rate
of growth of $f_1(\ensuremath{\mathtt{{\color{var}n}}})=15\ensuremath{\mathtt{{\color{var}n}}}$ versus $f_2(n)=2\ensuremath{\mathtt{{\color{var}n}}}\log\ensuremath{\mathtt{{\color{var}n}}}$.  It might be
that $f_1(n)$  is the running time of a complicated linear time algorithm
while $f_2(n)$ is the running time of a considerably simpler algorithm
based on the divide-and-conquer paradigm.  This illustrates that,
although $f_1(\ensuremath{\mathtt{{\color{var}n}}})$ is greater than $f_2(n)$ for small values of \ensuremath{\mathtt{{\color{var}n}}},
the opposite is true for large values of \ensuremath{\mathtt{{\color{var}n}}}.  Eventually $f_1(\ensuremath{\mathtt{{\color{var}n}}})$
wins out, by an increasingly wide margin.  Analysis using big-Oh notation
told us that this would happen, since $O(\ensuremath{\mathtt{{\color{var}n}}})\subset O(\ensuremath{\mathtt{{\color{var}n}}}\log \ensuremath{\mathtt{{\color{var}n}}})$.

\begin{figure}
  \begin{center}
    \newlength{\tmpa}\setlength{\tmpa}{.98\linewidth}
    \addtolength{\tmpa}{-4mm}
    \resizebox{\tmpa}{!}{\input{images/bigoh-1.tex}}\\[4ex]
    \resizebox{.98\linewidth}{!}{\input{images/bigoh-2.tex}}
  \end{center}
  \caption{Plots of $15\ensuremath{\mathtt{{\color{var}n}}}$ versus $2\ensuremath{\mathtt{{\color{var}n}}}\log\ensuremath{\mathtt{{\color{var}n}}}$.}
  \figlabel{intro-asymptotics}
\end{figure}

In a few cases, we will use asymptotic notation on functions with more
than one variable. There seems to be no standard for this, but for our
purposes, the following definition is sufficient:
\[
   O(f(n_1,\ldots,n_k)) = 
   \left\{\begin{array}{@{}l@{}}
             g(n_1,\ldots,n_k):\mbox{there exists $c>0$, and $z$ such that} \\
             \quad \mbox{$g(n_1,\ldots,n_k) \le c\cdot f(n_1,\ldots,n_k)$} \\
             \qquad \mbox{for all $n_1,\ldots,n_k$ such that $g(n_1,\ldots,n_k)\ge z$}   
   \end{array}\right\} \enspace .
\]
This definition captures the situation we really care about:  when the
arguments $n_1,\ldots,n_k$ make $g$ take on large values.  This definition
also agrees with the univariate definition of $O(f(n))$ when $f(n)$
is an increasing function of $n$.  The reader should be warned that,
although this works for our purposes, other texts may treat multivariate
functions and asymptotic notation differently.


\subsection{Randomization and Probability}
\seclabel{randomization}

\index{randomization}%
\index{probability}%
\index{randomized data structure}%
\index{randomized algorithm}%
Some of the data structures presented in this book are \emph{randomized};
they make random choices that are independent of the data being stored
in them or the operations being performed on them.  For this reason,
performing the same set of operations more than once using these
structures could result in different running times.  When analyzing these
data structures we are interested in their average or \emph{expected}
running times.
\index{expected running time}%
\index{running time!expected}%

Formally, the running time of an operation on a randomized data structure
is a random variable, and we want to study its \emph{expected value}.
\index{expected value}%
For
a discrete random variable $X$ taking on values in some countable
universe $U$, the expected value of $X$, denoted by $\E[X]$, is given
by the formula
\[
    \E[X] = \sum_{x\in U} x\cdot\Pr\{X=x\} \enspace .
\]
Here $\Pr\{\mathcal{E}\}$ denotes the probability that the event
$\mathcal{E}$ occurs.  In all of the examples in this book, these
probabilities are only with respect to the random choices made by the
randomized data structure;  there is no assumption that the data stored
in the structure, nor the sequence of operations performed on the
data structure, is random.

One of the most important properties of expected values is \emph{linearity
of expectation}.
\index{linearity of expectation}%
For any two random variables $X$ and $Y$,
\[
   \E[X+Y] = \E[X] + \E[Y] \enspace .
\]
More generally, for any random variables $X_1,\ldots,X_k$,
\[
   \E\left[\sum_{i=1}^k X_i\right] = \sum_{i=1}^k \E[X_i] \enspace .
\]
Linearity of expectation allows us to break down complicated random variables (like the left hand sides of the above equations) into sums of simpler random variables (the right hand sides).

A useful trick, that we will use repeatedly, is defining \emph{indicator
random variables}.
\index{indicator random variable}%
These binary variables are useful when we want to
count something and are best illustrated by an example.  Suppose we toss
a fair coin $k$ times and we want to know the expected number of times
the coin turns up as heads.
\index{coin toss}%
Intuitively, we know the answer is $k/2$,
but if we try to prove it using the definition of expected value, we get
\begin{align*}
   \E[X] & = \sum_{i=0}^k i\cdot\Pr\{X=i\} \\
         & = \sum_{i=0}^k i\cdot\binom{k}{i}/2^k \\
         & = k\cdot \sum_{i=0}^{k-1}\binom{k-1}{i}/2^k \\
         & = k/2 \enspace .
\end{align*}
This requires that we know enough to calculate that $\Pr\{X=i\}
= \binom{k}{i}/2^k$, and that we know the binomial identities
$i\binom{k}{i}=k\binom{k-1}{i-1}$ and $\sum_{i=0}^{k} \binom{k}{i} = 2^{k}$.

Using indicator variables and linearity of expectation makes things
much easier.  For each $i\in\{1,\ldots,k\}$, define the indicator
random variable
\[
    I_i = \begin{cases}
           1 & \text{if the $i$th coin toss is heads} \\
           0 & \text{otherwise.}
          \end{cases}
\]
Then 
\[ \E[I_i] = (1/2)1 + (1/2)0 = 1/2 \enspace . \]
Now, $X=\sum_{i=1}^k I_i$, so
\begin{align*}
   \E[X] & = \E\left[\sum_{i=1}^k I_i\right] \\
         & = \sum_{i=1}^k \E[I_i] \\
         & = \sum_{i=1}^k 1/2 \\
         & = k/2 \enspace .
\end{align*}
This is a bit more long-winded, but doesn't require that we know any
magical identities or compute any non-trivial probabilities. Even better,
it agrees with the intuition that we expect half the coins to turn up as
heads precisely because each individual coin turns up as heads with
a probability of $1/2$.

\section{The Model of Computation}
\seclabel{model}

In this book, we will analyze the theoretical running times of operations
on the data structures we study.  To do this precisely, we need a
mathematical model of computation.  For this, we use the \emph{\ensuremath{\mathtt{{\color{var}w}}}-bit
word-RAM}
\index{word-RAM}%
\index{RAM}%
model.  RAM stands for Random Access Machine. In this model, we
have access to a random access memory consisting of \emph{cells}, each of
which stores a \ensuremath{\mathtt{{\color{var}w}}}-bit \emph{word}.
\index{word}%
This implies that a memory cell can
represent, for example, any integer in the set $\{0,\ldots,2^{\ensuremath{\mathtt{{\color{var}w}}}}-1\}$.

In the word-RAM model, basic operations on words take constant time.
This includes arithmetic operations (\ensuremath{\mathtt{+}}, \ensuremath{\mathtt{-}}, \ensuremath{\mathtt{*}}, \ensuremath{\mathtt{/}}, \ensuremath{\mathtt{\text{\ttfamily \%}}}), comparisons
($<$, $>$, $=$, $\le$, $\ge$), and bitwise boolean operations (bitwise-AND,
OR, and exclusive-OR).

Any cell can be read or written in constant time.  A computer's memory
is managed by a memory management system from which we can allocate or
deallocate a block of memory of any size we would like. Allocating a
block of memory of size $k$ takes $O(k)$ time and returns a reference
(a pointer) to the newly-allocated memory block.  This reference is
small enough to be represented by a single word.

The word-size \ensuremath{\mathtt{{\color{var}w}}} is a very important parameter of this model.  The only
assumption we will make about \ensuremath{\mathtt{{\color{var}w}}} is the lower-bound $\ensuremath{\mathtt{{\color{var}w}}} \ge \log \ensuremath{\mathtt{{\color{var}n}}}$,
where \ensuremath{\mathtt{{\color{var}n}}} is the number of elements stored in any of our data structures.
This is a fairly modest assumption, since otherwise a word is not even
big enough to count the number of elements stored in the data structure.

Space is measured in words, so that when we talk about the amount of
space used by a data structure, we are referring to the number of words
of memory used by the structure.  All of our data structures store values of
a generic type \ensuremath{\mathtt{{\color{keyword}T}}}, and we assume an element of type \ensuremath{\mathtt{{\color{keyword}T}}} occupies one word
of memory.  \javaonly{(In reality, we are storing references to objects
of type \ensuremath{\mathtt{{\color{keyword}T}}}, and these references occupy only one word of memory.)}

\javaonly{The \ensuremath{\mathtt{{\color{var}w}}}-bit word-RAM model is a fairly close match for the
(32-bit) Java Virtual Machine (JVM) when $\ensuremath{\mathtt{{\color{var}w}}}=32$.}\cpponly{The \ensuremath{\mathtt{{\color{var}w}}}-bit
word-RAM model is a fairly close match for modern desktop computers when
$\ensuremath{\mathtt{{\color{var}w}}}=32$ or $\ensuremath{\mathtt{{\color{var}w}}}=64$.} The data structures presented in this book don't
use any special tricks that are not implementable \javaonly{on the JVM and most
other architectures.}\cpponly{in C++ on most architectures.}

\section{Correctness, Time Complexity, and Space Complexity}

When studying the performance of a data structure, there are three things
that matter most:

\begin{description}
  \item[Correctness:] The data structure should correctly implement
    its interface.
    \index{correctness}%
  \item[Time complexity:] The running times of operations on the data
    structure should be as small as possible.
    \index{time complexity}%
    \index{complexity!time}%
  \item[Space complexity:] The data structure should use as little memory
    as possible.
    \index{space complexity}%
    \index{complexity!space}%
\end{description}


In this introductory text, we will take correctness as a given;  we
won't consider data structures that give incorrect answers to queries or
don't perform updates properly.  We will, however, see data structures
that make an extra effort to keep space usage to a minimum.  This won't
usually affect the (asymptotic) running times of operations, but can
make the data structures a little slower in practice.

When studying running times in the context of data structures we tend to
come across three different kinds of running time guarantees:

\begin{description}
\item[Worst-case running times:] 
  \index{running time}%
  \index{running time!worst-case}%
  \index{worst-case running time}%
  These are the strongest kind of running
  time guarantees.  If a data structure operation has a worst-case
  running time of $f(\ensuremath{\mathtt{{\color{var}n}}})$, then one of these operations \emph{never}
  takes longer than $f(\ensuremath{\mathtt{{\color{var}n}}})$ time.
\item[Amortized running times:]
  \index{running time!amortized}%
  \index{amortized running time}%
  If we say that the amortized running
  time of an operation in a data structure is $f(\ensuremath{\mathtt{{\color{var}n}}})$, then this means that
  the cost of a typical operation is at most $f(\ensuremath{\mathtt{{\color{var}n}}})$.  More precisely,
  if a data structure has an amortized running time of $f(\ensuremath{\mathtt{{\color{var}n}}})$,
  then a sequence of $m$ operations takes at most $mf(\ensuremath{\mathtt{{\color{var}n}}})$ time.
  Some individual operations may take more than $f(\ensuremath{\mathtt{{\color{var}n}}})$ time but the
  average, over the entire sequence of operations, is at most $f(\ensuremath{\mathtt{{\color{var}n}}})$.
\item[Expected running times:] 
  \index{running time!expected}%
  \index{expected running time}%
  If we say that the expected running time
  of an operation on a data structure is $f(\ensuremath{\mathtt{{\color{var}n}}})$, this means that the
  actual running time is a random variable (see \secref{randomization})
  and the expected value of this random variable is at most $f(\ensuremath{\mathtt{{\color{var}n}}})$.
  The randomization here is with respect to random choices made by the
  data structure.
\end{description}

To understand the difference between worst-case, amortized, and expected
running times, it helps to consider a financial example.  Consider the
cost of buying a house:
\paragraph{Worst-case versus amortized cost:}
\index{amortized cost}%
Suppose that a home costs \$120\,000.  In order to buy this home,
we might get a 120 month (10 year) mortgage with monthly payments of
\$1\,200 per month.  In this case, the worst-case monthly cost of paying
this mortgage is \$1\,200 per month.

If we have enough cash on hand, we might choose to buy the house outright,
with one payment of \$120\,000.  In this case, over a period of 10 years,
the amortized monthly cost of buying this house is
\[
   \$120\,000 / 120\text{ months} = \$1\,000\text{ per month} \enspace .
\]
This is much less than the \$1\,200 per month we would have to pay if
we took out a mortgage.

\paragraph{Worst-case versus expected cost:}
\index{expected cost}%
Next, consider the issue of fire insurance on our \$120\,000 home.
By studying hundreds of thousands of cases, insurance companies have
determined that the expected amount of fire damage caused to a home
like ours is \$10 per month.  This is a very small number, since most
homes never have fires, a few homes may have some small fires that
cause a bit of smoke damage, and a tiny number of homes burn right to
their foundations.  Based on this information, the insurance company
charges \$15 per month for fire insurance.

Now it's decision time. Should we pay the \$15 worst-case monthly cost
for fire insurance, or should we gamble and self-insure at an expected
cost of \$10 per month?  Clearly, the \$10 per month costs less \emph{in
expectation}, but we have to be able to accept the possibility that
the \emph{actual cost} may be much higher.  In the unlikely event that the
entire house burns down, the actual cost will be \$120\,000.

These financial examples also offer insight into why we sometimes settle
for an amortized or expected running time over a worst-case running time.
It is often possible to get a lower expected or amortized running time
than a worst-case running time. At the very least, it is very often possible
to get a much simpler data structure if one is willing to settle for
amortized or expected running times.

\section{Code Samples}

\pcodeonly{

The code samples in this book are written in pseudocode.
\index{pseudocode}%
These should be easy enough to read for anyone who has any programming experience in any of the most common programming languages of the last 40 years.  To get an idea of what the pseudocode in this book looks like, here is a function that computes the average of an array, \ensuremath{\mathtt{{\color{var}a}}}:
\pcodeimport{ods/Algorithms.average(a)}
As this code illustrates, assigning to a variable is done using the $\gets$
notation.
We use the convention that the length of an
array, \ensuremath{\mathtt{{\color{var}a}}}, is denoted by \ensuremath{\mathtt{len({\color{var}a})}} and array indices start at zero,
so that \ensuremath{\mathtt{range(len({\color{var}a}))}} are the valid indices for \ensuremath{\mathtt{{\color{var}a}}}.  To shorten
code, and sometimes make it easier to read, our pseudocode allows for
(sub)array assignments.  The following two functions are equivalent:
\pcodeimport{ods/Algorithms.left_shift_a(a).left_shift_b(a)}
The following code sets all the values in an array to zero:
\pcodeimport{ods/Algorithms.zero(a)}
When analyzing the running time of code like this, we have to be careful;
statements like 
    \ensuremath{\mathtt{{\color{var}a}[0:len({\color{var}a})] = 1}}
or
   \ensuremath{\mathtt{{\color{var}a}[1:len({\color{var}a})] = {\color{var}a}[0:len({\color{var}a})-1]}}
do not run in constant time. They run in $O(\ensuremath{\mathtt{len({\color{var}a})}})$ time.

We take similar shortcuts with variable assignments, so that the code
\ensuremath{\mathtt{{\color{var}x},{\color{var}y}=0,1}} sets \ensuremath{\mathtt{{\color{var}x}}} to zero and \ensuremath{\mathtt{{\color{var}y}}} to 1 and the code \ensuremath{\mathtt{{\color{var}x},{\color{var}y} = {\color{var}y},{\color{var}x}}} swaps
the values of the variables \ensuremath{\mathtt{{\color{var}x}}} and \ensuremath{\mathtt{{\color{var}y}}}.
\index{swap}

Our pseudocode uses a few operators that may be unfamiliar.  As is
standard in mathematics, (normal) division is denoted by the $/$ operator.
In many cases, we want to do integer division instead, in which case we
use the $\ensuremath{\mathtt{{\color{comment}//}}}$ operator, so that $\ensuremath{\mathtt{{\color{var}a}{\color{comment}//b}}} = \lfloor a/b\rfloor$ is the
integer part of $a/b$.  So, for example, $3/2=1.5$ but $\ensuremath{\mathtt{3{\color{comment}//2}}} = 1$.
\index{integer division}%
\index{div operator}%
Occasionally, we also use the $\bmod$ operator to obtain the
remainder from integer division, but this will be defined when it is used.
\index{mod operator}%
\index{div operator}%
Later in the book, we may use some bitwise operators including left-shift (\ensuremath{\mathtt{\text{\ttfamily <<}}}), right shift (\ensuremath{\mathtt{\text{\ttfamily >>}}}), bitwise and (\ensuremath{\mathtt{\text{\ttfamily \&}}}), and bitwise exclusive-or (\ensuremath{\mathtt{^}}).
\index{left shift}%
\index{\ensuremath{\mathtt{\text{\ttfamily <<}}}|see {left shift}}%
\index{right shift}%
\index{\ensuremath{\mathtt{\text{\ttfamily >>}}}|see {right shift}}%
\index{bitwise and}%
\index{\ensuremath{\mathtt{\text{\ttfamily \&}}}|see {bitwise and}}%
\index{\ensuremath{\mathtt{^}}|see {bitwise exclusive-or}}%

The pseudocode samples in this book are machine translations
of Python code that can be downloaded from the book's website.\footnote{ \url{http://opendatastructures.org}}  If you ever encounter an ambiguity in the pseudocode that you can't resolve yourself, then you can always refer
to the corresponding Python code.  If you don't read Python, the code is
also available in Java and C++.  If you can't decipher the pseudocode,
or read Python, C++, or Java, then you may not be ready for this book.
}

\notpcode{
The code samples in this book are written in the \lang\ programming
language.  However, to make the book accessible to readers not familiar
with all of \lang's constructs and keywords, the code samples have
been simplified.  For example, a reader won't find any of the keywords
\ensuremath{\mathtt{{\color{keyword}public}}}, \ensuremath{\mathtt{{\color{keyword}protected}}}, \ensuremath{\mathtt{{\color{keyword}private}}}, or \ensuremath{\mathtt{{\color{keyword}static}}}.  A reader also won't find
much discussion about class hierarchies.  Which interfaces a particular
class implements or which class it extends, if relevant to the discussion,
should be clear from the accompanying text.

These conventions should make the code samples understandable by
anyone with a background in any of the languages from the ALGOL tradition,
including B, C, C++, C\#, Objective-C, D, Java, JavaScript, and so on.
Readers who want the full details of all implementations are encouraged
to look at the \lang\ source code that accompanies this book.

This book mixes mathematical analyses of running times with \lang\
source code for the algorithms being analyzed.  This means that
some equations contain variables also found in the source code.
These variables are typeset consistently, both within the source code
and within equations.  The most common such variable is the variable \ensuremath{\mathtt{{\color{var}n}}}
\index{n@\ensuremath{\mathtt{{\color{var}n}}}}%
that, without exception, always refers to the number of items currently
stored in the data structure.
}

\section{List of Data Structures}

Tables~\ref{tab:summary-i} and \ref{tab:summary-ii} summarize the
performance of data structures in this book that implement each of the
interfaces, \texttt{List}, \texttt{USet}, and \texttt{SSet}, described in \secref{interfaces}.
\Figref{dependencies} shows the dependencies between various chapters in
this book.
\index{dependencies}%
A dashed arrow indicates only a weak dependency, in which
only a small part of the chapter depends on a previous chapter or only
the main results of the previous chapter.

\begin{table}
\vspace{56pt}
\begin{center}
\resizebox{.98\textwidth}{!}{
\begin{threeparttable}
\begin{tabular}{|l|l|l|l|} \hline
\multicolumn{4}{|c|}{\texttt{List} implementations} \\ \hline
 & \ensuremath{\mathtt{get({\color{var}i})}}/\ensuremath{\mathtt{set({\color{var}i},{\color{var}x})}} & \ensuremath{\mathtt{add({\color{var}i},{\color{var}x})}}/\ensuremath{\mathtt{remove({\color{var}i})}} & \\ \hline
\texttt{ArrayStack} & $O(1)$ & $O(1+\ensuremath{\mathtt{{\color{var}n}}}-\ensuremath{\mathtt{{\color{var}i}}})$\tnote{A} & \sref{arraystack} \\
\texttt{ArrayDeque} & $O(1)$ & $O(1+\min\{\ensuremath{\mathtt{{\color{var}i}}},\ensuremath{\mathtt{{\color{var}n}}}-\ensuremath{\mathtt{{\color{var}i}}}\})$\tnote{A} & \sref{arraydeque} \\
\texttt{DualArrayDeque} & $O(1)$ & $O(1+\min\{\ensuremath{\mathtt{{\color{var}i}}},\ensuremath{\mathtt{{\color{var}n}}}-\ensuremath{\mathtt{{\color{var}i}}}\})$\tnote{A} & \sref{dualarraydeque}\\
\texttt{RootishArrayStack} & $O(1)$ & $O(1+\ensuremath{\mathtt{{\color{var}n}}}-\ensuremath{\mathtt{{\color{var}i}}})$\tnote{A}  & \sref{rootisharraystack} \\
\texttt{DLList} & $O(1+\min\{\ensuremath{\mathtt{{\color{var}i}}},\ensuremath{\mathtt{{\color{var}n}}}-\ensuremath{\mathtt{{\color{var}i}}}\})$ & $O(1+\min\{\ensuremath{\mathtt{{\color{var}i}}},\ensuremath{\mathtt{{\color{var}n}}}-\ensuremath{\mathtt{{\color{var}i}}}\})$  & \sref{dllist} \\
\texttt{SEList} & $O(1+\min\{\ensuremath{\mathtt{{\color{var}i}}},\ensuremath{\mathtt{{\color{var}n}}}-\ensuremath{\mathtt{{\color{var}i}}}\}/\ensuremath{\mathtt{{\color{var}b}}})$ & $O(\ensuremath{\mathtt{{\color{var}b}}}+\min\{\ensuremath{\mathtt{{\color{var}i}}},\ensuremath{\mathtt{{\color{var}n}}}-\ensuremath{\mathtt{{\color{var}i}}}\}/\ensuremath{\mathtt{{\color{var}b}}})$\tnote{A}  & \sref{selist} \\
\texttt{SkiplistList} & $O(\log \ensuremath{\mathtt{{\color{var}n}}})$\tnote{E} & $O(\log \ensuremath{\mathtt{{\color{var}n}}})$\tnote{E}  & \sref{skiplistlist} \\ \hline
\multicolumn{4}{c}{} \\[2ex] \hline
\multicolumn{4}{|c|}{\texttt{USet} implementations} \\ \hline
 & \ensuremath{\mathtt{find({\color{var}x})}} & \ensuremath{\mathtt{add({\color{var}x})}}/\ensuremath{\mathtt{remove({\color{var}x})}} & \\ \hline
\texttt{ChainedHashTable} & $O(1)$\tnote{E} & $O(1)$\tnote{A,E} & \sref{hashtable} \\ 
\texttt{LinearHashTable} & $O(1)$\tnote{E} & $O(1)$\tnote{A,E} & \sref{linearhashtable} \\ \hline
\end{tabular}
\begin{tablenotes}
\item[A]{Denotes an \emph{amortized} running time.}
\item[E]{Denotes an \emph{expected} running time.}
\end{tablenotes}
\end{threeparttable}}
\end{center}
\caption[Summary of List and USet implementations.]{Summary of \texttt{List} and \texttt{USet} implementations.}
\tablabel{summary-i}
\end{table}

\begin{table}
\begin{center}
\begin{threeparttable}
\begin{tabular}{|l|l|l|l|} \hline
\multicolumn{4}{|c|}{\texttt{SSet} implementations} \\ \hline
 & \ensuremath{\mathtt{find({\color{var}x})}} & \ensuremath{\mathtt{add({\color{var}x})}}/\ensuremath{\mathtt{remove({\color{var}x})}} & \\ \hline
\texttt{SkiplistSSet} & $O(\log \ensuremath{\mathtt{{\color{var}n}}})$\tnote{E} & $O(\log \ensuremath{\mathtt{{\color{var}n}}})$\tnote{E} & \sref{skiplistset} \\ 
\texttt{Treap} & $O(\log \ensuremath{\mathtt{{\color{var}n}}})$\tnote{E} & $O(\log \ensuremath{\mathtt{{\color{var}n}}})$\tnote{E} & \sref{treap} \\ 
\texttt{ScapegoatTree} & $O(\log \ensuremath{\mathtt{{\color{var}n}}})$ & $O(\log \ensuremath{\mathtt{{\color{var}n}}})$\tnote{A} & \sref{scapegoattree} \\
\texttt{RedBlackTree} & $O(\log \ensuremath{\mathtt{{\color{var}n}}})$ & $O(\log \ensuremath{\mathtt{{\color{var}n}}})$ & \sref{redblacktree} \\ 
\texttt{BinaryTrie}\tnote{I} & $O(\ensuremath{\mathtt{{\color{var}w}}})$ & $O(\ensuremath{\mathtt{{\color{var}w}}})$ & \sref{binarytrie} \\ 
\texttt{XFastTrie}\tnote{I} & $O(\log \ensuremath{\mathtt{{\color{var}w}}})$\tnote{A,E} & $O(\ensuremath{\mathtt{{\color{var}w}}})$\tnote{A,E} & \sref{xfast} \\ 
\texttt{YFastTrie}\tnote{I} & $O(\log \ensuremath{\mathtt{{\color{var}w}}})$\tnote{A,E} & $O(\log \ensuremath{\mathtt{{\color{var}w}}})$\tnote{A,E} & \sref{yfast} \\ 
\javaonly{\texttt{BTree} & $O(\log \ensuremath{\mathtt{{\color{var}n}}})$ & $O(B+\log \ensuremath{\mathtt{{\color{var}n}}})$\tnote{A} & \sref{btree} \\ 
\texttt{BTree}\tnote{X} & $O(\log_B \ensuremath{\mathtt{{\color{var}n}}})$ & $O(\log_B \ensuremath{\mathtt{{\color{var}n}}})$ & \sref{btree} \\ } \hline
\multicolumn{4}{c}{} \\[2ex] \hline
\multicolumn{4}{|c|}{(Priority) \texttt{Queue} implementations} \\ \hline
 & \ensuremath{\mathtt{findMin()}} & \ensuremath{\mathtt{add({\color{var}x})}}/\ensuremath{\mathtt{remove()}} & \\ \hline
\texttt{BinaryHeap} & $O(1)$ & $O(\log \ensuremath{\mathtt{{\color{var}n}}})$\tnote{A} & \sref{binaryheap} \\ 
\texttt{MeldableHeap} & $O(1)$ & $O(\log \ensuremath{\mathtt{{\color{var}n}}})$\tnote{E} & \sref{meldableheap} \\ \hline
\end{tabular}
\begin{tablenotes}
\item[I]{This structure can only store \ensuremath{\mathtt{{\color{var}w}}}-bit integer data.}
\javaonly{\item[X]{This denotes the running time in the external-memory model; see \chapref{btree}.}}
\end{tablenotes}
\end{threeparttable}
\end{center}
\caption[Summary of SSet and priority Queue implementations.]{Summary of \texttt{SSet} and priority \texttt{Queue} implementations.}
\tablabel{summary-ii}
\end{table}

\begin{figure}
  \begin{center}
    \includegraphics[width=\ScaleIfNeeded]{figs/dependencies}
  \end{center}
  \caption{The dependencies between chapters in this book.}
  \figlabel{dependencies}
\end{figure}

\section{Discussion and Exercises}

The \texttt{List}, \texttt{USet}, and \texttt{SSet} interfaces described in
\secref{interfaces} are influenced by the Java Collections Framework
\cite{oracle_collections}.
\index{Java Collections Framework}%
These are essentially simplified versions of
the \texttt{List}, \texttt{Set}, \texttt{Map}, \texttt{SortedSet}, and \texttt{SortedMap} interfaces found in
the Java Collections Framework.  \javaonly{The accompanying source
code includes wrapper classes for making \texttt{USet} and \texttt{SSet} implementations
into \texttt{Set}, \texttt{Map}, \texttt{SortedSet}, and \texttt{SortedMap} implementations.}

For a superb (and free) treatment of the mathematics discussed in this
chapter, including asymptotic notation, logarithms, factorials, Stirling's
approximation, basic probability, and lots more, see the textbook by
Leyman, Leighton, and Meyer \cite{llm11}.  For a gentle calculus text
that includes formal definitions of exponentials and logarithms, see the
(freely available) classic text by Thompson \cite{t14}.

For more information on basic probability, especially as it relates to
computer science, see the textbook by Ross \cite{r01}.  Another good
reference, which covers both asymptotic notation and probability, is
the textbook by Graham, Knuth, and Patashnik \cite{gkp94}.

\javaonly{Readers wanting to brush up on their Java programming can find
many Java tutorials online \cite{oracle_tutorials}.}

\begin{exc}
  This exercise is designed to help familiarize the reader with choosing
  the right data structure for the right problem.  If implemented, the
  parts of this exercise should be done by making use of an implementation
  of the relevant interface (\texttt{Stack}, \texttt{Queue}, \texttt{Deque}, \texttt{USet}, or \texttt{SSet})
  provided by the \javaonly{Java Collections Framework}\cpponly{C++
  Standard Template Library}.

  Solve the following problems by reading a text file one line at a
  time and performing operations on each line in the appropriate data
  structure(s).  Your implementations should be fast enough that even
  files containing a million lines can be processed in a few seconds.
  \begin{enumerate}
    \item Read the input one line at a time and then write the lines out
      in reverse order, so that the last input line is printed first,
      then the second last input line, and so on.

    \item  Read the first 50 lines of input and then write them out in
      reverse order. Read the next 50 lines and then write them out in
      reverse order. Do this until there are no more lines left to read,
      at which point any remaining lines should be output in reverse
      order.

      In other words, your output will start with the 50th line, then
      the 49th, then the 48th, and so on down to the first line. This
      will be followed by the 100th line, followed by the 99th, and so
      on down to the 51st line. And so on.
      
      Your code should never have to store more than 50 lines at any
      given time.

    \item Read the input one line at a time.
      At any point after reading the first 42 lines, if some line is blank
      (i.e., a string of length 0), then output the line that occured
      42 lines prior to that one. For example, if Line 242 is blank,
      then your program should output line 200. This program should
      be implemented so that it never stores more than 43 lines of the
      input at any given time.

    \item Read the input one line at a time and write each line to the
      output if it is not a duplicate of some previous input line. Take
      special care so that a file with a lot of duplicate lines does not
      use more memory than what is required for the number of unique lines.

    \item Read the input one line at a time and write each line to the
      output only if you have already read this line before. (The end
      result is that you remove the first occurrence of each line.)
      Take special care so that a file with a lot of duplicate lines
      does not use more memory than what is required for the number of
      unique lines.

    \item Read the entire input one line at a time. Then output all lines
      sorted by length, with the shortest lines first. In the case where
      two lines have the same length, resolve their order using the usual
      ``sorted order.''  Duplicate lines should be printed only once.

    \item Do the same as the previous question except that duplicate lines
      should be printed the same number of times that they appear in the input.

    \item Read the entire input one line at a time and then output the
      even numbered lines (starting with the first line, line 0) followed
      by the odd-numbered lines.

    \item Read the entire input one line at a time and randomly permute
      the lines before outputting them.  To be clear: You should not
      modify the contents of any line. Instead, the same collection of
      lines should be printed, but in a random order.
  \end{enumerate}
\end{exc}

\begin{exc}
  \index{Dyck word}%
  A \emph{Dyck word} is a sequence of +1's and -1's with the property that
  the sum of any prefix of the sequence is never negative.  For example,
  $+1,-1,+1,-1$ is a Dyck word, but $+1,-1,-1,+1$ is not a Dyck word
  since the prefix $+1-1-1<0$.  Describe any relationship
  between Dyck words and \texttt{Stack} \ensuremath{\mathtt{push({\color{var}x})}} and \ensuremath{\mathtt{pop()}} operations.
\end{exc}

\begin{exc}
  \index{matched string}%
  \index{string!matched}%
  A \emph{matched string} is a sequence of \{, \}, (, ), [, and ]
  characters that are properly matched.  For example, ``\{\{()[]\}\}''
  is a matched string, but this ``\{\{()]\}'' is not, since the second \{
  is matched with a ].  Show how to use a stack so that, given a string
  of length \ensuremath{\mathtt{{\color{var}n}}}, you can determine if it is a matched string in $O(\ensuremath{\mathtt{{\color{var}n}}})$
  time.
\end{exc}

\begin{exc}
  Suppose you have a \texttt{Stack}, \ensuremath{\mathtt{{\color{var}s}}}, that supports only the \ensuremath{\mathtt{push({\color{var}x})}}
  and \ensuremath{\mathtt{pop()}} operations. Show how, using only a FIFO \texttt{Queue}, \ensuremath{\mathtt{{\color{var}q}}},
  you can reverse the order of all elements in \ensuremath{\mathtt{{\color{var}s}}}.
\end{exc}

\begin{exc}
  \index{Bag@\texttt{Bag}}%
  Using a \texttt{USet}, implement a \texttt{Bag}.  A \texttt{Bag} is like a \texttt{USet}---it
  supports the \ensuremath{\mathtt{add({\color{var}x})}}, \ensuremath{\mathtt{remove({\color{var}x})}} and \ensuremath{\mathtt{find({\color{var}x})}} methods---but it allows
  duplicate elements to be stored.  The \ensuremath{\mathtt{find({\color{var}x})}} operation in a \texttt{Bag}
  returns some element (if any) that is equal to \ensuremath{\mathtt{{\color{var}x}}}.  In addition,
  a \texttt{Bag} supports the \ensuremath{\mathtt{findAll({\color{var}x})}} operation that returns a list of
  all elements in the \texttt{Bag} that are equal to \ensuremath{\mathtt{{\color{var}x}}}.
\end{exc}

\begin{exc}
  From scratch, write and test implementations of the \texttt{List}, \texttt{USet}
  and \texttt{SSet} interfaces.  These do not have to be efficient.  They can
  be used later to test the correctness and performance of more efficient
  implementations.  (The easiest way to do this is to store the elements
  in an array.)
\end{exc}

\begin{exc}
  Work to improve the performance of your implementations from the
  previous question using any tricks you can think of.  Experiment and
  think about how you could improve the performance of \ensuremath{\mathtt{add({\color{var}i},{\color{var}x})}} and
  \ensuremath{\mathtt{remove({\color{var}i})}} in your \texttt{List} implementation.  Think about how you could
  improve the performance of the \ensuremath{\mathtt{find({\color{var}x})}} operation in your \texttt{USet}
  and \texttt{SSet} implementations.  This exercise is designed to give you a
  feel for  how difficult it can be to obtain efficient implementations
  of these interfaces.
\end{exc}





\clearpage

\clearpage
\chapter{Array-Based Lists}
\chaplabel{arrays}

In this chapter, we will study implementations of the \texttt{List} and \texttt{Queue}
interfaces where the underlying data is stored in an array, called the
\emph{backing array}.
\index{backing array}%
The following table summarizes the running times
of operations for the data structures presented in this chapter:
\newlength{\tabsep}
\setlength{\tabsep}{\itemsep}
\addtolength{\tabsep}{\parsep}
\addtolength{\tabsep}{-2pt}
\begin{center}
\vspace{\tabsep}
\begin{tabular}{|l|l|l|} \hline
 & \ensuremath{\mathtt{get({\color{var}i})}}/\ensuremath{\mathtt{set({\color{var}i},{\color{var}x})}} & \ensuremath{\mathtt{add({\color{var}i},{\color{var}x})}}/\ensuremath{\mathtt{remove({\color{var}i})}} \\ \hline
\texttt{ArrayStack} & $O(1)$ & $O(\ensuremath{\mathtt{{\color{var}n}}}-\ensuremath{\mathtt{{\color{var}i}}})$ \\
\texttt{ArrayDeque} & $O(1)$ & $O(\min\{\ensuremath{\mathtt{{\color{var}i}}},\ensuremath{\mathtt{{\color{var}n}}}-\ensuremath{\mathtt{{\color{var}i}}}\})$ \\
\texttt{DualArrayDeque} & $O(1)$ & $O(\min\{\ensuremath{\mathtt{{\color{var}i}}},\ensuremath{\mathtt{{\color{var}n}}}-\ensuremath{\mathtt{{\color{var}i}}}\})$ \\
\texttt{RootishArrayStack} & $O(1)$ & $O(\ensuremath{\mathtt{{\color{var}n}}}-\ensuremath{\mathtt{{\color{var}i}}})$ \\ \hline
\end{tabular}
\vspace{\tabsep}
\end{center}
Data structures that work by storing data in a single array have many
advantages and limitations in common:
\index{arrays}%
\begin{itemize}
  \item Arrays offer constant time access to any value in the array.
  This is what allows \ensuremath{\mathtt{get({\color{var}i})}} and \ensuremath{\mathtt{set({\color{var}i},{\color{var}x})}} to run in constant time.

  \item Arrays are not very dynamic.  Adding or removing an element
  near the middle of a list means that a large number of elements in the
  array need to be shifted to make room for the newly added element or
  to fill in the gap created by the deleted element.  This is why the
  operations \ensuremath{\mathtt{add({\color{var}i},{\color{var}x})}} and \ensuremath{\mathtt{remove({\color{var}i})}} have running times that depend
  on \ensuremath{\mathtt{{\color{var}n}}} and \ensuremath{\mathtt{{\color{var}i}}}.

  \item Arrays cannot expand or shrink.  When the number of elements in
  the data structure exceeds the size of the backing array, a new array needs
  to be allocated and the data from the old array needs to be copied
  into the new array.  This is an expensive operation.
\end{itemize}
The third point is important.  The running times cited in the table
above do not include the cost associated with growing and shrinking
the backing array.  We will see that, if carefully managed, the cost of
growing and shrinking the backing array does not add much to the cost of
an \emph{average} operation.  More precisely, if we start with an empty
data structure, and perform any sequence of $m$ \ensuremath{\mathtt{add({\color{var}i},{\color{var}x})}} or \ensuremath{\mathtt{remove({\color{var}i})}}
operations, then the total cost of growing and shrinking the backing
array, over the entire sequence of $m$ operations is $O(m)$.  Although
some individual operations are more expensive, the amortized cost,
when amortized over all $m$ operations, is only $O(1)$ per operation.

\cpponly{
In this chapter, and throughout this book, it will be convenient to
have arrays that keep track of their size.  The usual C++ arrays do
not do this, so we have defined a class,  \ensuremath{\mathtt{{\color{var}array}}}, that keeps track
of its length.  The implementation of this class is straightforward.
It is implemented as a standard C++ array, \ensuremath{\mathtt{{\color{var}a}}}, and an integer, \ensuremath{\mathtt{{\color{var}length}}}:}
\cpponly{
The size of an \ensuremath{\mathtt{{\color{var}array}}} is specified at the time of creation:
}
\cpponly{The elements of an array can be indexed:}
\cpponly{Finally, when one array is assigned to another, this is just
a pointer manipulation that takes constant time:}

\section{\texttt{ArrayStack}: Fast Stack Operations Using an Array}
\seclabel{arraystack}

\index{ArrayStack@\texttt{ArrayStack}}%
An \texttt{ArrayStack} implements the list interface using an array \ensuremath{\mathtt{{\color{var}a}}}, called
the \emph{backing array}.  The list element with index \ensuremath{\mathtt{{\color{var}i}}} is stored
in \ensuremath{\mathtt{{\color{var}a}[{\color{var}i}]}}.  At most times, \ensuremath{\mathtt{{\color{var}a}}} is larger than strictly necessary,
so an integer \ensuremath{\mathtt{{\color{var}n}}} is used to keep track of the number of elements
actually stored in \ensuremath{\mathtt{{\color{var}a}}}.  In this way, the list elements are stored in
\ensuremath{\mathtt{{\color{var}a}[0]}},\ldots,\ensuremath{\mathtt{{\color{var}a}[{\color{var}n}-1]}} and, at all times, $\ensuremath{\mathtt{{\color{var}a}.{\color{var}length}}} \ge \ensuremath{\mathtt{{\color{var}n}}}$.



 \ \vspace{.5ex}\begin{Verbatim}[gobble=1,tabsize=2,frame=single,framerule=1.2pt,commandchars=\\@\$,label=\texttt{ArrayStack},labelposition=topline]
	@\color@keyword$T$[] @\color@var$a$;
	@\color@keyword$int$ @\color@var$n$;
	@\color@keyword$int$ size() {
		@\color@var$return$ @\color@var$n$;
	}
\end{Verbatim}


\subsection{The Basics}

Accessing and modifying the elements of an \texttt{ArrayStack} using \ensuremath{\mathtt{get({\color{var}i})}}
and \ensuremath{\mathtt{set({\color{var}i},{\color{var}x})}} is trivial. After performing any necessary bounds-checking
we simply return or set, respectively, \ensuremath{\mathtt{{\color{var}a}[{\color{var}i}]}}.



 \ \vspace{.5ex}\begin{Verbatim}[gobble=1,tabsize=2,frame=single,framerule=1.2pt,commandchars=\\@\$,label=\texttt{ArrayStack},labelposition=topline]
	@\color@keyword$T$ get(@\color@keyword$int$ @\color@var$i$) {
		@\color@var$return$ @\color@var$a$[@\color@var$i$];
	}
	@\color@keyword$T$ set(@\color@keyword$int$ @\color@var$i$, @\color@keyword$T$ @\color@var$x$) {
		@\color@keyword$T$ @\color@var$y$ = @\color@var$a$[@\color@var$i$];
		@\color@var$a$[@\color@var$i$] = @\color@var$x$;
		@\color@var$return$ @\color@var$y$;
	}
\end{Verbatim}


The operations of adding and removing elements from an \texttt{ArrayStack}
are illustrated in \figref{arraystack}.  To implement the \ensuremath{\mathtt{add({\color{var}i},{\color{var}x})}}
operation, we first check if \ensuremath{\mathtt{{\color{var}a}}} is already full.  If so, we call
the method \ensuremath{\mathtt{resize()}} to increase the size of \ensuremath{\mathtt{{\color{var}a}}}.  How \ensuremath{\mathtt{resize()}}
is implemented will be discussed later.  For now, it is sufficient to
know that, after a call to \ensuremath{\mathtt{resize()}}, we can be sure that $\ensuremath{\mathtt{{\color{var}a}.{\color{var}length}}}
> \ensuremath{\mathtt{{\color{var}n}}}$.  With this out of the way, we now shift the elements
$\ensuremath{\mathtt{{\color{var}a}[{\color{var}i}]}},\ldots,\ensuremath{\mathtt{{\color{var}a}[{\color{var}n}-1]}}$ right by one position to make room for \ensuremath{\mathtt{{\color{var}x}}},
set \ensuremath{\mathtt{{\color{var}a}[{\color{var}i}]}} equal to \ensuremath{\mathtt{{\color{var}x}}}, and increment \ensuremath{\mathtt{{\color{var}n}}}.

\begin{figure}
  \begin{center}
    \includegraphics[scale=0.90909]{figs/arraystack}
  \end{center}
  \caption[Adding to an ArrayStack]{A sequence of \ensuremath{\mathtt{add({\color{var}i},{\color{var}x})}} and \ensuremath{\mathtt{remove({\color{var}i})}} operations on an
  \texttt{ArrayStack}.  Arrows denote elements being copied.  Operations that
  result in a call to \ensuremath{\mathtt{resize()}} are marked with an asterisk.}
  \figlabel{arraystack}
\end{figure}



 \ \vspace{.5ex}\begin{Verbatim}[gobble=1,tabsize=2,frame=single,framerule=1.2pt,commandchars=\\@\$,label=\texttt{ArrayStack},labelposition=topline]
	@\color@keyword$void$ add(@\color@keyword$int$ @\color@var$i$, @\color@keyword$T$ @\color@var$x$) {
		@\color@keyword$if$ (@\color@var$n$ + 1 > @\color@var$a$.@\color@var$length$) resize();
		@\color@keyword$for$ (@\color@keyword$int$ @\color@var$j$ = @\color@var$n$; @\color@var$j$ > @\color@var$i$; @\color@var$j$--) 
			@\color@var$a$[@\color@var$j$] = @\color@var$a$[@\color@var$j$-1];
		@\color@var$a$[@\color@var$i$] = @\color@var$x$;
		@\color@var$n$++;
	}
\end{Verbatim}

If we ignore the cost of the potential call to \ensuremath{\mathtt{resize()}}, then the cost
of the \ensuremath{\mathtt{add({\color{var}i},{\color{var}x})}} operation is proportional to the number of elements we
have to shift to make room for \ensuremath{\mathtt{{\color{var}x}}}.  Therefore the cost of this operation
(ignoring the cost of resizing \ensuremath{\mathtt{{\color{var}a}}}) is $O(\ensuremath{\mathtt{{\color{var}n}}}-\ensuremath{\mathtt{{\color{var}i}}})$.

Implementing the \ensuremath{\mathtt{remove({\color{var}i})}} operation is similar.  We shift the elements
$\ensuremath{\mathtt{{\color{var}a}[{\color{var}i}+1]}},\ldots,\ensuremath{\mathtt{{\color{var}a}[{\color{var}n}-1]}}$ left by one position (overwriting \ensuremath{\mathtt{{\color{var}a}[{\color{var}i}]}}) and
decrease the value of \ensuremath{\mathtt{{\color{var}n}}}.  After doing this, we check if \ensuremath{\mathtt{{\color{var}n}}} is getting
much smaller than \ensuremath{\mathtt{{\color{var}a}.{\color{var}length}}} by checking if $\ensuremath{\mathtt{{\color{var}a}.{\color{var}length}}} \ge 3\ensuremath{\mathtt{{\color{var}n}}}$. If so,
then we call \ensuremath{\mathtt{resize()}} to reduce the size of \ensuremath{\mathtt{{\color{var}a}}}.



 \ \vspace{.5ex}\begin{Verbatim}[gobble=1,tabsize=2,frame=single,framerule=1.2pt,commandchars=\\@\$,label=\texttt{ArrayStack},labelposition=topline]
	@\color@keyword$T$ remove(@\color@keyword$int$ @\color@var$i$) {
		@\color@keyword$T$ @\color@var$x$ = @\color@var$a$[@\color@var$i$];
		@\color@keyword$for$ (@\color@keyword$int$ @\color@var$j$ = @\color@var$i$; @\color@var$j$ < @\color@var$n$-1; @\color@var$j$++) 
			@\color@var$a$[@\color@var$j$] = @\color@var$a$[@\color@var$j$+1];
		@\color@var$n$--;
		@\color@keyword$if$ (@\color@var$a$.@\color@var$length$ >= 3*@\color@var$n$) resize();
		@\color@var$return$ @\color@var$x$;
	}
\end{Verbatim}

If we ignore the cost of the \ensuremath{\mathtt{resize()}} method, the cost of a \ensuremath{\mathtt{remove({\color{var}i})}}
operation is proportional to the number of elements we shift, which
is $O(\ensuremath{\mathtt{{\color{var}n}}}-\ensuremath{\mathtt{{\color{var}i}}})$.

\subsection{Growing and Shrinking}

The \ensuremath{\mathtt{resize()}} method is fairly straightforward; it allocates a new
array \ensuremath{\mathtt{{\color{var}b}}} whose size is $2\ensuremath{\mathtt{{\color{var}n}}}$ and copies the \ensuremath{\mathtt{{\color{var}n}}} elements of \ensuremath{\mathtt{{\color{var}a}}} into
the first \ensuremath{\mathtt{{\color{var}n}}} positions in \ensuremath{\mathtt{{\color{var}b}}}, and then sets \ensuremath{\mathtt{{\color{var}a}}} to \ensuremath{\mathtt{{\color{var}b}}}. Thus, after a call to \ensuremath{\mathtt{resize()}}, $\ensuremath{\mathtt{{\color{var}a}.{\color{var}length}}} = 2\ensuremath{\mathtt{{\color{var}n}}}$.



 \ \vspace{.5ex}\begin{Verbatim}[gobble=1,tabsize=2,frame=single,framerule=1.2pt,commandchars=\\@\$,label=\texttt{ArrayStack},labelposition=topline]
	@\color@keyword$void$ resize() {
		@\color@keyword$T$[] @\color@var$b$ = newArray(max(@\color@var$n$*2,1));
		@\color@keyword$for$ (@\color@keyword$int$ @\color@var$i$ = 0; @\color@var$i$ < @\color@var$n$; @\color@var$i$++) {
			@\color@var$b$[@\color@var$i$] = @\color@var$a$[@\color@var$i$];
		}
		@\color@var$a$ = @\color@var$b$;
	}
\end{Verbatim}


Analyzing the actual cost of the \ensuremath{\mathtt{resize()}} operation is easy. It
allocates an array \ensuremath{\mathtt{{\color{var}b}}} of size $2\ensuremath{\mathtt{{\color{var}n}}}$ and copies the \ensuremath{\mathtt{{\color{var}n}}} elements of \ensuremath{\mathtt{{\color{var}a}}}
into \ensuremath{\mathtt{{\color{var}b}}}.  This takes $O(\ensuremath{\mathtt{{\color{var}n}}})$ time.

The running time analysis from the previous section ignored the cost
of calls to \ensuremath{\mathtt{resize()}}.  In this section we analyze this cost using a
technique known as \emph{amortized analysis}.  This technique does not
try to determine the cost of resizing during each individual \ensuremath{\mathtt{add({\color{var}i},{\color{var}x})}}
and \ensuremath{\mathtt{remove({\color{var}i})}} operation.  Instead, it considers the cost of all calls to
\ensuremath{\mathtt{resize()}} during a sequence of $m$ calls to \ensuremath{\mathtt{add({\color{var}i},{\color{var}x})}} or \ensuremath{\mathtt{remove({\color{var}i})}}.
In particular, we will show:

\begin{lem}\lemlabel{arraystack-amortized}
  If an empty \texttt{ArrayStack} is created and any sequence of $m\ge 1$ calls
  to \ensuremath{\mathtt{add({\color{var}i},{\color{var}x})}} and \ensuremath{\mathtt{remove({\color{var}i})}} are performed, then the total time spent
  during all calls to \ensuremath{\mathtt{resize()}} is $O(m)$.
\end{lem}

\begin{proof}
  We will show that any time \ensuremath{\mathtt{resize()}} is called, the number of calls
  to \ensuremath{\mathtt{{\color{var}add}}} or \ensuremath{\mathtt{{\color{var}remove}}} since the last call to \ensuremath{\mathtt{resize()}} is at least
  $\ensuremath{\mathtt{{\color{var}n}}}/2-1$.  Therefore, if $\ensuremath{\mathtt{{\color{var}n}}}_i$ denotes the value of \ensuremath{\mathtt{{\color{var}n}}} during the
  $i$th call to \ensuremath{\mathtt{resize()}} and $r$ denotes the number of calls to
  \ensuremath{\mathtt{resize()}}, then the total number of calls to \ensuremath{\mathtt{add({\color{var}i},{\color{var}x})}} or
  \ensuremath{\mathtt{remove({\color{var}i})}} is at least
  \[
     \sum_{i=1}^{r} (\ensuremath{\mathtt{{\color{var}n}}}_i/2-1) \le m  \enspace ,
  \]
  which is equivalent to
  \[
    \sum_{i=1}^{r} \ensuremath{\mathtt{{\color{var}n}}}_i \le 2m + 2r  \enspace .
  \]
  On the other hand, the total time spent during all calls to \ensuremath{\mathtt{resize()}} is 
  \[
     \sum_{i=1}^{r} O(\ensuremath{\mathtt{{\color{var}n}}}_i) \le O(m+r) = O(m)  \enspace ,
  \]
  since $r$ is not more than $m$.  All that remains is to show that the
  number of calls to \ensuremath{\mathtt{add({\color{var}i},{\color{var}x})}} or \ensuremath{\mathtt{remove({\color{var}i})}} between the $(i-1)$th
  and the $i$th call to \ensuremath{\mathtt{resize()}} is at least $\ensuremath{\mathtt{{\color{var}n}}}_i/2$.

  There are two cases to consider. In the first case, \ensuremath{\mathtt{resize()}} is
  being called by \ensuremath{\mathtt{add({\color{var}i},{\color{var}x})}} because the backing array \ensuremath{\mathtt{{\color{var}a}}} is full, i.e.,
  $\ensuremath{\mathtt{{\color{var}a}.{\color{var}length}}} = \ensuremath{\mathtt{{\color{var}n}}}=\ensuremath{\mathtt{{\color{var}n}}}_i$.  Consider the previous call to \ensuremath{\mathtt{resize()}}:
  after this previous call, the size of \ensuremath{\mathtt{{\color{var}a}}} was \ensuremath{\mathtt{{\color{var}a}.{\color{var}length}}}, but the
  number of elements stored in \ensuremath{\mathtt{{\color{var}a}}} was at most $\ensuremath{\mathtt{{\color{var}a}.{\color{var}length}}}/2=\ensuremath{\mathtt{{\color{var}n}}}_i/2$.
  But now the number of elements stored in \ensuremath{\mathtt{{\color{var}a}}} is $\ensuremath{\mathtt{{\color{var}n}}}_i=\ensuremath{\mathtt{{\color{var}a}.{\color{var}length}}}$,
  so there must have been at least $\ensuremath{\mathtt{{\color{var}n}}}_i/2$ calls to \ensuremath{\mathtt{add({\color{var}i},{\color{var}x})}} since
  the previous call to \ensuremath{\mathtt{resize()}}.
  
  The second case occurs when \ensuremath{\mathtt{resize()}} is being called by
  \ensuremath{\mathtt{remove({\color{var}i})}} because $\ensuremath{\mathtt{{\color{var}a}.{\color{var}length}}} \ge 3\ensuremath{\mathtt{{\color{var}n}}}=3\ensuremath{\mathtt{{\color{var}n}}}_i$.  Again, after the
  previous call to \ensuremath{\mathtt{resize()}} the number of elements stored in \ensuremath{\mathtt{{\color{var}a}}} was
  at least $\ensuremath{\mathtt{{\color{var}a}.{\color{var}length}/2}}-1$.\footnote{The ${}-1$ in this formula accounts for
  the special case that occurs when $\ensuremath{\mathtt{{\color{var}n}}}=0$ and $\ensuremath{\mathtt{{\color{var}a}.{\color{var}length}}} = 1$.} Now there
  are $\ensuremath{\mathtt{{\color{var}n}}}_i\le\ensuremath{\mathtt{{\color{var}a}.{\color{var}length}}}/3$ elements stored in \ensuremath{\mathtt{{\color{var}a}}}.  Therefore, the number
  of \ensuremath{\mathtt{remove({\color{var}i})}} operations since the last call to \ensuremath{\mathtt{resize()}} is at least
  \begin{align*}
      R & \ge \ensuremath{\mathtt{{\color{var}a}.{\color{var}length}}}/2 - 1 - \ensuremath{\mathtt{{\color{var}a}.{\color{var}length}}}/3 \\
        & = \ensuremath{\mathtt{{\color{var}a}.{\color{var}length}}}/6 - 1 \\
        & = (\ensuremath{\mathtt{{\color{var}a}.{\color{var}length}}}/3)/2 - 1 \\
        & \ge \ensuremath{\mathtt{{\color{var}n}}}_i/2 -1\enspace .
  \end{align*}
  In either case, the number of calls to \ensuremath{\mathtt{add({\color{var}i},{\color{var}x})}} or \ensuremath{\mathtt{remove({\color{var}i})}} that
  occur between the $(i-1)$th call to \ensuremath{\mathtt{resize()}} and the $i$th call to
  \ensuremath{\mathtt{resize()}} is at least $\ensuremath{\mathtt{{\color{var}n}}}_i/2-1$, as required to complete the proof.
\end{proof}

\subsection{Summary}

The following theorem summarizes the performance of an \texttt{ArrayStack}:

\begin{thm}\thmlabel{arraystack}
  An \texttt{ArrayStack} implements the \texttt{List} interface.  Ignoring the cost of
  calls to \ensuremath{\mathtt{resize()}}, an \texttt{ArrayStack} supports the operations
  \begin{itemize}
    \item \ensuremath{\mathtt{get({\color{var}i})}} and \ensuremath{\mathtt{set({\color{var}i},{\color{var}x})}} in $O(1)$ time per operation; and
    \item \ensuremath{\mathtt{add({\color{var}i},{\color{var}x})}} and \ensuremath{\mathtt{remove({\color{var}i})}} in $O(1+\ensuremath{\mathtt{{\color{var}n}}}-\ensuremath{\mathtt{{\color{var}i}}})$ time per operation.
  \end{itemize}
  Furthermore, beginning with an empty \texttt{ArrayStack} and performing any
  sequence of $m$ \ensuremath{\mathtt{add({\color{var}i},{\color{var}x})}} and \ensuremath{\mathtt{remove({\color{var}i})}} operations results in a
  total of $O(m)$ time spent during all calls to \ensuremath{\mathtt{resize()}}.
\end{thm}

The \texttt{ArrayStack} is an efficient way to implement a \texttt{Stack}.
In particular, we can implement \ensuremath{\mathtt{push({\color{var}x})}} as \ensuremath{\mathtt{add({\color{var}n},{\color{var}x})}} and \ensuremath{\mathtt{pop()}}
as \ensuremath{\mathtt{remove({\color{var}n}-1)}}, in which case these operations will run in $O(1)$
amortized time.

\section{\texttt{FastArrayStack}: An Optimized ArrayStack}
\seclabel{fastarraystack}

\index{FastArrayStack@\texttt{FastArrayStack}}%
Much of the work done by an \texttt{ArrayStack} involves shifting (by
\ensuremath{\mathtt{add({\color{var}i},{\color{var}x})}} and \ensuremath{\mathtt{remove({\color{var}i})}}) and copying (by \ensuremath{\mathtt{resize()}}) of data.
\notpcode{In the implementations shown above, this was done using \ensuremath{\mathtt{{\color{keyword}for}}} loops.}%
\pcodeonly{In a naive implementation, this would be done using \ensuremath{\mathtt{{\color{keyword}for}}} loops.}
It
turns out that many programming environments have specific functions
that are very efficient at copying and moving blocks of data.  In the C
programming language, there are the \ensuremath{\mathtt{memcpy({\color{var}d},{\color{var}s},{\color{var}n})}} and \ensuremath{\mathtt{memmove({\color{var}d},{\color{var}s},{\color{var}n})}}
functions. In the C++ language there is the \ensuremath{\mathtt{{\color{var}std}::copy({\color{var}a0},{\color{var}a1},{\color{var}b})}} algorithm.
In Java there is the \ensuremath{\mathtt{System.arraycopy({\color{var}s},{\color{var}i},{\color{var}d},{\color{var}j},{\color{var}n})}} method.
\index{memcpy@\ensuremath{\mathtt{memcpy({\color{var}d},{\color{var}s},{\color{var}n})}}}%
\index{std::copy@\ensuremath{\mathtt{{\color{var}std}::copy({\color{var}a0},{\color{var}a1},{\color{var}b})}}}%
\index{System.arraycopy@\ensuremath{\mathtt{System.arraycopy({\color{var}s},{\color{var}i},{\color{var}d},{\color{var}j},{\color{var}n})}}}%



 \ \vspace{.5ex}\begin{Verbatim}[gobble=1,tabsize=2,frame=single,framerule=1.2pt,commandchars=\\@\$,label=\texttt{FastArrayStack},labelposition=topline]
	@\color@keyword$void$ resize() {
		@\color@keyword$T$[] @\color@var$b$ = newArray(max(2*@\color@var$n$,1));
		System.arraycopy(@\color@var$a$, 0, @\color@var$b$, 0, @\color@var$n$);
		@\color@var$a$ = @\color@var$b$;
	}
	@\color@keyword$void$ add(@\color@keyword$int$ @\color@var$i$, @\color@keyword$T$ @\color@var$x$) {
		@\color@keyword$if$ (@\color@var$n$ + 1 > @\color@var$a$.@\color@var$length$) resize();
		System.arraycopy(@\color@var$a$, @\color@var$i$, @\color@var$a$, @\color@var$i$+1, @\color@var$n$-@\color@var$i$); 
		@\color@var$a$[@\color@var$i$] = @\color@var$x$;
		@\color@var$n$++;
	}
	@\color@keyword$T$ remove(@\color@keyword$int$ @\color@var$i$) {
		@\color@keyword$T$ @\color@var$x$ = @\color@var$a$[@\color@var$i$];
		System.arraycopy(@\color@var$a$, @\color@var$i$+1, @\color@var$a$, @\color@var$i$, @\color@var$n$-@\color@var$i$-1);
		@\color@var$n$--; 
		@\color@keyword$if$ (@\color@var$a$.@\color@var$length$ >= 3*@\color@var$n$) resize();
		@\color@var$return$ @\color@var$x$;
	}
\end{Verbatim}


These functions are usually highly optimized and may even use special
machine instructions that can do this copying much faster than we could by
using a \ensuremath{\mathtt{{\color{keyword}for}}} loop.  Although using these functions does not asymptotically
decrease the running times, it can still be a worthwhile optimization.

\pcodeonly{In our C++ and Java implementations, the use of fast array
copying functions}
\notpcode{In the \lang\ implementations here, the use of the native\linebreak\javaonly{\ensuremath{\mathtt{System.arraycopy({\color{var}s},{\color{var}i},{\color{var}d},{\color{var}j},{\color{var}n})}}}\cpponly{\ensuremath{\mathtt{{\color{var}std}::copy({\color{var}a0},{\color{var}a1},{\color{var}b})}}}}
resulted in speedups of a factor between 2 and 3, depending on the types of
operations performed.  Your mileage may vary.

\section{\texttt{ArrayQueue}: An Array-Based Queue}
\seclabel{arrayqueue}

\index{ArrayQueue@\texttt{ArrayQueue}}%
In this section, we present the \texttt{ArrayQueue} data structure, which
implements a FIFO (first-in-first-out) queue; elements are removed (using
the \ensuremath{\mathtt{remove()}} operation) from the queue in the same order they are added
(using the \ensuremath{\mathtt{add({\color{var}x})}} operation).

Notice that an \texttt{ArrayStack} is a poor choice for an implementation of a
FIFO queue.  It is not a good choice because we must choose one end of
the list upon which to add elements and then remove elements from the
other end.  One of the two operations must work on the head of the list,
which involves calling \ensuremath{\mathtt{add({\color{var}i},{\color{var}x})}} or \ensuremath{\mathtt{remove({\color{var}i})}} with a value of $\ensuremath{\mathtt{{\color{var}i}}}=0$.
This gives a running time proportional to \ensuremath{\mathtt{{\color{var}n}}}.

To obtain an efficient array-based implementation of a queue, we
first notice that the problem would be easy if we had an infinite
array \ensuremath{\mathtt{{\color{var}a}}}.  We could maintain one index \ensuremath{\mathtt{{\color{var}j}}} that keeps track of the
next element to remove and an integer \ensuremath{\mathtt{{\color{var}n}}} that counts the number of
elements in the queue.  The queue elements would always be stored in
\[ \ensuremath{\mathtt{{\color{var}a}[{\color{var}j}]}},\ensuremath{\mathtt{{\color{var}a}[{\color{var}j}+1]}},\ldots,\ensuremath{\mathtt{{\color{var}a}[{\color{var}j}+{\color{var}n}-1]}} \enspace . \]
Initially, both \ensuremath{\mathtt{{\color{var}j}}} and \ensuremath{\mathtt{{\color{var}n}}} would be 
set to 0.  To add an element, we would place it in \ensuremath{\mathtt{{\color{var}a}[{\color{var}j}+{\color{var}n}]}} and increment \ensuremath{\mathtt{{\color{var}n}}}.
To remove an element, we would remove it from \ensuremath{\mathtt{{\color{var}a}[{\color{var}j}]}}, increment \ensuremath{\mathtt{{\color{var}j}}}, and
decrement \ensuremath{\mathtt{{\color{var}n}}}.

Of course, the problem with this solution is that it requires an infinite
array.  An \texttt{ArrayQueue} simulates this by using a finite array \ensuremath{\mathtt{{\color{var}a}}}
and \emph{modular arithmetic}.
\index{modular arithmetic}%
This is the kind of arithmetic used when
we are talking about the time of day.  For example 10:00 plus five
hours gives 3:00.  Formally, we say that
\[
    10 + 5 = 15 \equiv 3 \pmod{12} \enspace .
\]
We read the latter part of this equation as ``15 is congruent to 3 modulo
12.'' We can also treat $\bmod$ as a binary operator, so that
\[
   15 \bmod 12 = 3 \enspace .
\]

More generally, for an integer $a$ and positive integer $m$, $a \bmod m$
is the unique integer $r\in\{0,\ldots,m-1\}$ such that $a = r + km$ for
some integer $k$.  Less formally, the value $r$ is the remainder we get
when we divide $a$ by $m$.
\pcodeonly{In many programming languages, including C, C++, and Java,
the mod operate is represented using the \% symbol.}
\notpcode{In many programming languages, including
\javaonly{Java}\cpponly{C++}, the $\bmod$ operator is represented
using the \ensuremath{\mathtt{\text{\ttfamily \%}}} symbol.\footnote{This is sometimes referred to as the
\emph{brain-dead} mod operator, since it does not correctly implement
the mathematical mod operator when the first argument is negative.}}

Modular arithmetic is useful for simulating an infinite array,
since $\ensuremath{\mathtt{{\color{var}i}}}\bmod \ensuremath{\mathtt{{\color{var}a}.{\color{var}length}}}$ always gives a value in the range
$0,\ldots,\ensuremath{\mathtt{{\color{var}a}.{\color{var}length}-1}}$.  Using modular arithmetic we can store the
queue elements at array locations
\[ \ensuremath{\mathtt{{\color{var}a}[{\color{var}j}\text{\ttfamily \%}{\color{var}a}.{\color{var}length}]}},\ensuremath{\mathtt{{\color{var}a}[({\color{var}j}+1)\text{\ttfamily \%}{\color{var}a}.{\color{var}length}]}},\ldots,\ensuremath{\mathtt{{\color{var}a}[({\color{var}j}+{\color{var}n}-1)\text{\ttfamily \%}{\color{var}a}.{\color{var}length}]}}
\enspace. \]
This treats the array \ensuremath{\mathtt{{\color{var}a}}} like a \emph{circular array}
\index{circular array}%
\index{array!circular}%
in which array indices
larger than $\ensuremath{\mathtt{{\color{var}a}.{\color{var}length}}}-1$ ``wrap around'' to the beginning of
the array.

The only remaining thing to worry about is taking care that the number
of elements in the \texttt{ArrayQueue} does not exceed the size of \ensuremath{\mathtt{{\color{var}a}}}.



 \ \vspace{.5ex}\begin{Verbatim}[gobble=1,tabsize=2,frame=single,framerule=1.2pt,commandchars=\\@\$,label=\texttt{ArrayQueue},labelposition=topline]
	@\color@keyword$T$[] @\color@var$a$;
	@\color@keyword$int$ @\color@var$j$;
	@\color@keyword$int$ @\color@var$n$;
\end{Verbatim}


A sequence of \ensuremath{\mathtt{add({\color{var}x})}} and \ensuremath{\mathtt{remove()}} operations on an \texttt{ArrayQueue} is
illustrated in \figref{arrayqueue}.  To implement \ensuremath{\mathtt{add({\color{var}x})}}, we first
check if \ensuremath{\mathtt{{\color{var}a}}} is full and, if necessary, call \ensuremath{\mathtt{resize()}} to increase
the size of \ensuremath{\mathtt{{\color{var}a}}}.  Next, we store \ensuremath{\mathtt{{\color{var}x}}} in
\ensuremath{\mathtt{{\color{var}a}[({\color{var}j}+{\color{var}n})\text{\ttfamily \%}{\color{var}a}.{\color{var}length}]}} and increment \ensuremath{\mathtt{{\color{var}n}}}.

\begin{figure}
  \begin{center}
    \includegraphics[scale=0.90909]{figs/arrayqueue}
  \end{center}
  \caption[Adding and removing from an ArrayQueue]{A sequence of \ensuremath{\mathtt{add({\color{var}x})}} and \ensuremath{\mathtt{remove({\color{var}i})}} operations on an
  \texttt{ArrayQueue}.  Arrows denote elements being copied.  Operations that
  result in a call to \ensuremath{\mathtt{resize()}} are marked with an asterisk.}
  \figlabel{arrayqueue}
\end{figure}





 \ \vspace{.5ex}\begin{Verbatim}[gobble=1,tabsize=2,frame=single,framerule=1.2pt,commandchars=\\@\$,label=\texttt{ArrayQueue},labelposition=topline]
	@\color@var$boolean$ add(@\color@keyword$T$ @\color@var$x$) {
		@\color@keyword$if$ (@\color@var$n$ + 1 > @\color@var$a$.@\color@var$length$) resize();
		@\color@var$a$[(@\color@var$j$+@\color@var$n$) % @\color@var$a$.@\color@var$length$] = @\color@var$x$;
		@\color@var$n$++;
		@\color@var$return$ @\color@var$true$;
	}
\end{Verbatim}


To implement \ensuremath{\mathtt{remove()}}, we first store \ensuremath{\mathtt{{\color{var}a}[{\color{var}j}]}} so that we can return
it later.  Next, we decrement \ensuremath{\mathtt{{\color{var}n}}} and increment \ensuremath{\mathtt{{\color{var}j}}} (modulo \ensuremath{\mathtt{{\color{var}a}.{\color{var}length}}})
by setting $\ensuremath{\mathtt{{\color{var}j}}}=(\ensuremath{\mathtt{{\color{var}j}}}+1)\bmod \ensuremath{\mathtt{{\color{var}a}.{\color{var}length}}}$.  Finally, we return the stored
value of \ensuremath{\mathtt{{\color{var}a}[{\color{var}j}]}}. If necessary, we may call \ensuremath{\mathtt{resize()}} to decrease the
size of \ensuremath{\mathtt{{\color{var}a}}}.



 \ \vspace{.5ex}\begin{Verbatim}[gobble=1,tabsize=2,frame=single,framerule=1.2pt,commandchars=\\@\$,label=\texttt{ArrayQueue},labelposition=topline]
	@\color@keyword$T$ remove() { 
		@\color@keyword$if$ (@\color@var$n$ == 0) @\color@keyword$throw$ @\color@keyword$new$ NoSuchElementException();
		@\color@keyword$T$ @\color@var$x$ = @\color@var$a$[@\color@var$j$];
		@\color@var$j$ = (@\color@var$j$ + 1) % @\color@var$a$.@\color@var$length$;
		@\color@var$n$--;
		@\color@keyword$if$ (@\color@var$a$.@\color@var$length$ >= 3*@\color@var$n$) resize();
		@\color@var$return$ @\color@var$x$;
	}
\end{Verbatim}


Finally, the \ensuremath{\mathtt{resize()}} operation is very similar to the \ensuremath{\mathtt{resize()}}
operation of \texttt{ArrayStack}.  It allocates a new array, \ensuremath{\mathtt{{\color{var}b}}}, of size $2\ensuremath{\mathtt{{\color{var}n}}}$
and copies
\[
   \ensuremath{\mathtt{{\color{var}a}[{\color{var}j}]}},\ensuremath{\mathtt{{\color{var}a}[({\color{var}j}+1)\text{\ttfamily \%}{\color{var}a}.{\color{var}length}]}},\ldots,\ensuremath{\mathtt{{\color{var}a}[({\color{var}j}+{\color{var}n}-1)\text{\ttfamily \%}{\color{var}a}.{\color{var}length}]}}
\]
onto
\[
   \ensuremath{\mathtt{{\color{var}b}[0]}},\ensuremath{\mathtt{{\color{var}b}[1]}},\ldots,\ensuremath{\mathtt{{\color{var}b}[{\color{var}n}-1]}}
\]
and sets $\ensuremath{\mathtt{{\color{var}j}}}=0$.



 \ \vspace{.5ex}\begin{Verbatim}[gobble=1,tabsize=2,frame=single,framerule=1.2pt,commandchars=\\@\$,label=\texttt{ArrayQueue},labelposition=topline]
	@\color@keyword$void$ resize() {
		@\color@keyword$T$[] @\color@var$b$ = newArray(max(1,@\color@var$n$*2));
		@\color@keyword$for$ (@\color@keyword$int$ @\color@var$k$ = 0; @\color@var$k$ < @\color@var$n$; @\color@var$k$++) 
			@\color@var$b$[@\color@var$k$] = @\color@var$a$[(@\color@var$j$+@\color@var$k$) % @\color@var$a$.@\color@var$length$];
		@\color@var$a$ = @\color@var$b$;
		@\color@var$j$ = 0;
	}
\end{Verbatim}


\subsection{Summary}

The following theorem summarizes the performance of the \texttt{ArrayQueue}
data structure:

\begin{thm}
An \texttt{ArrayQueue} implements the (FIFO) \texttt{Queue} interface.  Ignoring the cost of
calls to \ensuremath{\mathtt{resize()}}, an \texttt{ArrayQueue} supports the operations
\ensuremath{\mathtt{add({\color{var}x})}} and \ensuremath{\mathtt{remove()}} in $O(1)$ time per operation.
Furthermore, beginning with an empty \texttt{ArrayQueue}, any sequence of $m$
\ensuremath{\mathtt{add({\color{var}i},{\color{var}x})}} and \ensuremath{\mathtt{remove({\color{var}i})}} operations results in a total of $O(m)$
time spent during all calls to \ensuremath{\mathtt{resize()}}.
\end{thm}


\section{\texttt{ArrayDeque}: Fast Deque Operations Using an Array}
\seclabel{arraydeque}

\index{ArrayDeque@\texttt{ArrayDeque}}%
The \texttt{ArrayQueue} from the previous section is a data structure for
representing a sequence that allows us to efficiently add to one
end of the sequence and remove from the other end.  The \texttt{ArrayDeque}
data structure allows for efficient addition and removal at both ends.
This structure implements the \texttt{List} interface by using the same circular
array technique used to represent an \texttt{ArrayQueue}.



 \ \vspace{.5ex}\begin{Verbatim}[gobble=1,tabsize=2,frame=single,framerule=1.2pt,commandchars=\\@\$,label=\texttt{ArrayDeque},labelposition=topline]
	@\color@keyword$T$[] @\color@var$a$;
	@\color@keyword$int$ @\color@var$j$;
	@\color@keyword$int$ @\color@var$n$;
\end{Verbatim}


The \ensuremath{\mathtt{get({\color{var}i})}} and \ensuremath{\mathtt{set({\color{var}i},{\color{var}x})}} operations on an \texttt{ArrayDeque} are
straightforward.  They get or set the array element $\ensuremath{\mathtt{{\color{var}a}[}}{\ensuremath{\mathtt{({\color{var}j}+{\color{var}i})}}\bmod
\ensuremath{\mathtt{{\color{var}a}.{\color{var}length}}}}\ensuremath{\mathtt{]}}$.



 \ \vspace{.5ex}\begin{Verbatim}[gobble=1,tabsize=2,frame=single,framerule=1.2pt,commandchars=\\@\$,label=\texttt{ArrayDeque},labelposition=topline]
	@\color@keyword$T$ get(@\color@keyword$int$ @\color@var$i$) {
		@\color@var$return$ @\color@var$a$[(@\color@var$j$+@\color@var$i$)%@\color@var$a$.@\color@var$length$];
	}
	@\color@keyword$T$ set(@\color@keyword$int$ @\color@var$i$, @\color@keyword$T$ @\color@var$x$) {
		@\color@keyword$T$ @\color@var$y$ = @\color@var$a$[(@\color@var$j$+@\color@var$i$)%@\color@var$a$.@\color@var$length$];
		@\color@var$a$[(@\color@var$j$+@\color@var$i$)%@\color@var$a$.@\color@var$length$] = @\color@var$x$;
		@\color@var$return$ @\color@var$y$;
	}
\end{Verbatim}


The implementation of \ensuremath{\mathtt{add({\color{var}i},{\color{var}x})}} is a little more interesting.  As
usual, we first check if \ensuremath{\mathtt{{\color{var}a}}} is full and, if necessary, call
\ensuremath{\mathtt{resize()}} to resize \ensuremath{\mathtt{{\color{var}a}}}.  Remember that we want this operation to be
fast when \ensuremath{\mathtt{{\color{var}i}}} is small (close to 0) or when \ensuremath{\mathtt{{\color{var}i}}} is large (close to
\ensuremath{\mathtt{{\color{var}n}}}).  Therefore, we check if $\ensuremath{\mathtt{{\color{var}i}}}<\ensuremath{\mathtt{{\color{var}n}}}/2$.  If so, we shift the
elements $\ensuremath{\mathtt{{\color{var}a}[0]}},\ldots,\ensuremath{\mathtt{{\color{var}a}[{\color{var}i}-1]}}$ left by one position.  Otherwise
($\ensuremath{\mathtt{{\color{var}i}}}\ge\ensuremath{\mathtt{{\color{var}n}}}/2$), we shift the elements $\ensuremath{\mathtt{{\color{var}a}[{\color{var}i}]}},\ldots,\ensuremath{\mathtt{{\color{var}a}[{\color{var}n}-1]}}$ right
by one position.  See \figref{arraydeque} for an illustration of
\ensuremath{\mathtt{add({\color{var}i},{\color{var}x})}} and \ensuremath{\mathtt{remove({\color{var}x})}} operations on an \texttt{ArrayDeque}.

\begin{figure}
  \begin{center}
    \includegraphics[scale=0.90909]{figs/arraydeque}
  \end{center}
  \caption[Adding and removing from an ArrayDeque]{A sequence of \ensuremath{\mathtt{add({\color{var}i},{\color{var}x})}} and \ensuremath{\mathtt{remove({\color{var}i})}} operations on an
  \texttt{ArrayDeque}.  Arrows denote elements being copied.}
  \figlabel{arraydeque}
\end{figure}




 \ \vspace{.5ex}\begin{Verbatim}[gobble=1,tabsize=2,frame=single,framerule=1.2pt,commandchars=\\@\$,label=\texttt{ArrayDeque},labelposition=topline]
	@\color@keyword$void$ add(@\color@keyword$int$ @\color@var$i$, @\color@keyword$T$ @\color@var$x$) {
		@\color@keyword$if$ (@\color@var$n$+1 > @\color@var$a$.@\color@var$length$) resize();
		@\color@keyword$if$ (@\color@var$i$ < @\color@var$n$/2) { @\color@comment$// shift a[0],..,a[i-1] left one position$
			@\color@var$j$ = (@\color@var$j$ == 0) ? @\color@var$a$.@\color@var$length$ - 1 : @\color@var$j$ - 1; @\color@comment$//(j-1)mod a.length$
			@\color@keyword$for$ (@\color@keyword$int$ @\color@var$k$ = 0; @\color@var$k$ <= @\color@var$i$-1; @\color@var$k$++)
				@\color@var$a$[(@\color@var$j$+@\color@var$k$)%@\color@var$a$.@\color@var$length$] = @\color@var$a$[(@\color@var$j$+@\color@var$k$+1)%@\color@var$a$.@\color@var$length$];
		} @\color@keyword$else$ { @\color@comment$// shift a[i],..,a[n-1] right one position$
			@\color@keyword$for$ (@\color@keyword$int$ @\color@var$k$ = @\color@var$n$; @\color@var$k$ > @\color@var$i$; @\color@var$k$--)
				@\color@var$a$[(@\color@var$j$+@\color@var$k$)%@\color@var$a$.@\color@var$length$] = @\color@var$a$[(@\color@var$j$+@\color@var$k$-1)%@\color@var$a$.@\color@var$length$];
		}
		@\color@var$a$[(@\color@var$j$+@\color@var$i$)%@\color@var$a$.@\color@var$length$] = @\color@var$x$;
		@\color@var$n$++;
	}
\end{Verbatim}


By doing the shifting in this way, we guarantee that \ensuremath{\mathtt{add({\color{var}i},{\color{var}x})}} never
has to shift more than $\min\{ \ensuremath{\mathtt{{\color{var}i}}}, \ensuremath{\mathtt{{\color{var}n}}}-\ensuremath{\mathtt{{\color{var}i}}} \}$ elements.  Thus, the running
time of the \ensuremath{\mathtt{add({\color{var}i},{\color{var}x})}} operation (ignoring the cost of a \ensuremath{\mathtt{resize()}}
operation) is $O(1+\min\{\ensuremath{\mathtt{{\color{var}i}}},\ensuremath{\mathtt{{\color{var}n}}}-\ensuremath{\mathtt{{\color{var}i}}}\})$.

The implementation of the \ensuremath{\mathtt{remove({\color{var}i})}} operation is similar.  It either
shifts elements $\ensuremath{\mathtt{{\color{var}a}[0]}},\ldots,\ensuremath{\mathtt{{\color{var}a}[{\color{var}i}-1]}}$ right by one position or shifts
the elements $\ensuremath{\mathtt{{\color{var}a}[{\color{var}i}+1]}},\ldots,\ensuremath{\mathtt{{\color{var}a}[{\color{var}n}-1]}}$ left by one position depending
on whether $\ensuremath{\mathtt{{\color{var}i}}}<\ensuremath{\mathtt{{\color{var}n}}}/2$.  Again, this means that \ensuremath{\mathtt{remove({\color{var}i})}} never spends
more than $O(1+\min\{\ensuremath{\mathtt{{\color{var}i}}},\ensuremath{\mathtt{{\color{var}n}}}-\ensuremath{\mathtt{{\color{var}i}}}\})$ time to shift elements.



 \ \vspace{.5ex}\begin{Verbatim}[gobble=1,tabsize=2,frame=single,framerule=1.2pt,commandchars=\\@\$,label=\texttt{ArrayDeque},labelposition=topline]
	@\color@keyword$T$ remove(@\color@keyword$int$ @\color@var$i$) {
		@\color@keyword$T$ @\color@var$x$ = @\color@var$a$[(@\color@var$j$+@\color@var$i$)%@\color@var$a$.@\color@var$length$];
		@\color@keyword$if$ (@\color@var$i$ < @\color@var$n$/2) {  @\color@comment$// shift a[0],..,[i-1] right one position$
			@\color@keyword$for$ (@\color@keyword$int$ @\color@var$k$ = @\color@var$i$; @\color@var$k$ > 0; @\color@var$k$--)
				@\color@var$a$[(@\color@var$j$+@\color@var$k$)%@\color@var$a$.@\color@var$length$] = @\color@var$a$[(@\color@var$j$+@\color@var$k$-1)%@\color@var$a$.@\color@var$length$];
			@\color@var$j$ = (@\color@var$j$ + 1) % @\color@var$a$.@\color@var$length$;
		} @\color@keyword$else$ { @\color@comment$// shift a[i+1],..,a[n-1] left one position$
			@\color@keyword$for$ (@\color@keyword$int$ @\color@var$k$ = @\color@var$i$; @\color@var$k$ < @\color@var$n$-1; @\color@var$k$++)
				@\color@var$a$[(@\color@var$j$+@\color@var$k$)%@\color@var$a$.@\color@var$length$] = @\color@var$a$[(@\color@var$j$+@\color@var$k$+1)%@\color@var$a$.@\color@var$length$];
		}
		@\color@var$n$--;
		@\color@keyword$if$ (3*@\color@var$n$ < @\color@var$a$.@\color@var$length$) resize();
		@\color@var$return$ @\color@var$x$;
	}
\end{Verbatim}


\subsection{Summary}

The following theorem summarizes the performance of the \texttt{ArrayDeque}
data structure:
\begin{thm}\thmlabel{arraydeque}
  An \texttt{ArrayDeque} implements the \texttt{List} interface.  Ignoring the cost of
  calls to \ensuremath{\mathtt{resize()}}, an \texttt{ArrayDeque} supports the operations
  \begin{itemize}
    \item \ensuremath{\mathtt{get({\color{var}i})}} and \ensuremath{\mathtt{set({\color{var}i},{\color{var}x})}} in $O(1)$ time per operation; and
    \item \ensuremath{\mathtt{add({\color{var}i},{\color{var}x})}} and \ensuremath{\mathtt{remove({\color{var}i})}} in $O(1+\min\{\ensuremath{\mathtt{{\color{var}i}}},\ensuremath{\mathtt{{\color{var}n}}}-\ensuremath{\mathtt{{\color{var}i}}}\})$ time
          per operation.
  \end{itemize}
  Furthermore, beginning with an empty \texttt{ArrayDeque}, performing any
  sequence of $m$ \ensuremath{\mathtt{add({\color{var}i},{\color{var}x})}} and \ensuremath{\mathtt{remove({\color{var}i})}} operations results in a
  total of $O(m)$ time spent during all calls to \ensuremath{\mathtt{resize()}}.
\end{thm}

\section{\texttt{DualArrayDeque}: Building a Deque from Two Stacks}
\seclabel{dualarraydeque}

\index{DualArrayDeque@\texttt{DualArrayDeque}}%
Next, we present a data structure, the \texttt{DualArrayDeque} that
achieves the same performance bounds as an \texttt{ArrayDeque} by using
two \texttt{ArrayStack}s.  Although the asymptotic performance of the
\texttt{DualArrayDeque} is no better than that of the \texttt{ArrayDeque}, it is
still worth studying, since it offers a good example of how to make a
sophisticated data structure by combining two simpler data structures.

A \texttt{DualArrayDeque} represents a list using two \texttt{ArrayStack}s.  Recall that
an \texttt{ArrayStack} is fast when the operations on it modify elements near
the end.  A \texttt{DualArrayDeque} places two \texttt{ArrayStack}s, called \ensuremath{\mathtt{{\color{var}front}}}
and \ensuremath{\mathtt{{\color{var}back}}}, back-to-back so that operations are fast at either end.



 \ \vspace{.5ex}\begin{Verbatim}[gobble=1,tabsize=2,frame=single,framerule=1.2pt,commandchars=\\@\$,label=\texttt{DualArrayDeque},labelposition=topline]
	List<@\color@keyword$T$> @\color@var$front$;
	List<@\color@keyword$T$> @\color@var$back$;
\end{Verbatim}


A \texttt{DualArrayDeque} does not explicitly store the number, \ensuremath{\mathtt{{\color{var}n}}},
of elements it contains.  It doesn't need to, since it contains
$\ensuremath{\mathtt{{\color{var}n}}}=\ensuremath{\mathtt{{\color{var}front}.size()}} + \ensuremath{\mathtt{{\color{var}back}.size()}}$ elements.  Nevertheless, when
analyzing the \texttt{DualArrayDeque} we will still use \ensuremath{\mathtt{{\color{var}n}}} to denote the number
of elements it contains.



 \ \vspace{.5ex}\begin{Verbatim}[gobble=1,tabsize=2,frame=single,framerule=1.2pt,commandchars=\\@\$,label=\texttt{DualArrayDeque},labelposition=topline]
	@\color@keyword$int$ size() {
		@\color@var$return$ @\color@var$front$.size() + @\color@var$back$.size();		
	}
\end{Verbatim}


The \ensuremath{\mathtt{{\color{var}front}}} \texttt{ArrayStack} stores the list elements that whose indices
are $0,\ldots,\ensuremath{\mathtt{{\color{var}front}.size()}}-1$, but stores them in reverse order.
The \ensuremath{\mathtt{{\color{var}back}}} \texttt{ArrayStack} contains list elements with indices
in $\ensuremath{\mathtt{{\color{var}front}.size()}},\ldots,\ensuremath{\mathtt{size()}}-1$ in the normal order.  In this way,
\ensuremath{\mathtt{get({\color{var}i})}} and \ensuremath{\mathtt{set({\color{var}i},{\color{var}x})}} translate into appropriate calls to \ensuremath{\mathtt{get({\color{var}i})}}
or \ensuremath{\mathtt{set({\color{var}i},{\color{var}x})}} on either \ensuremath{\mathtt{{\color{var}front}}} or \ensuremath{\mathtt{{\color{var}back}}}, which take $O(1)$ time per operation.



 \ \vspace{.5ex}\begin{Verbatim}[gobble=1,tabsize=2,frame=single,framerule=1.2pt,commandchars=\\@\$,label=\texttt{DualArrayDeque},labelposition=topline]
	@\color@keyword$T$ get(@\color@keyword$int$ @\color@var$i$) {
		@\color@keyword$if$ (@\color@var$i$ < @\color@var$front$.size()) {
			@\color@var$return$ @\color@var$front$.get(@\color@var$front$.size()-@\color@var$i$-1);
		} @\color@keyword$else$ {
			@\color@var$return$ @\color@var$back$.get(@\color@var$i$-@\color@var$front$.size());
		}
	}
	@\color@keyword$T$ set(@\color@keyword$int$ @\color@var$i$, @\color@keyword$T$ @\color@var$x$) {
		@\color@keyword$if$ (@\color@var$i$ < @\color@var$front$.size()) {
			@\color@var$return$ @\color@var$front$.set(@\color@var$front$.size()-@\color@var$i$-1, @\color@var$x$);
			
		} @\color@keyword$else$ {
			@\color@var$return$ @\color@var$back$.set(@\color@var$i$-@\color@var$front$.size(), @\color@var$x$);
		}
	}
\end{Verbatim}


Note that if an index $\ensuremath{\mathtt{{\color{var}i}}}<\ensuremath{\mathtt{{\color{var}front}.size()}}$, then it corresponds to the
element of \ensuremath{\mathtt{{\color{var}front}}} at position $\ensuremath{\mathtt{{\color{var}front}.size()}}-\ensuremath{\mathtt{{\color{var}i}}}-1$, since the
elements of \ensuremath{\mathtt{{\color{var}front}}} are stored in reverse order.

Adding and removing elements from a \texttt{DualArrayDeque} is illustrated in
\figref{dualarraydeque}.  The \ensuremath{\mathtt{add({\color{var}i},{\color{var}x})}} operation manipulates either \ensuremath{\mathtt{{\color{var}front}}}
or \ensuremath{\mathtt{{\color{var}back}}}, as appropriate:

\begin{figure}
  \begin{center}
    \includegraphics[scale=0.90909]{figs/dualarraydeque}
  \end{center}
  \caption[Adding and removing in a DualArrayDeque]{A sequence of \ensuremath{\mathtt{add({\color{var}i},{\color{var}x})}} and \ensuremath{\mathtt{remove({\color{var}i})}} operations on a
  \texttt{DualArrayDeque}.  Arrows denote elements being copied.  Operations that
  result in a rebalancing by \ensuremath{\mathtt{balance()}} are marked with an asterisk.}
  \figlabel{dualarraydeque}
\end{figure}





 \ \vspace{.5ex}\begin{Verbatim}[gobble=1,tabsize=2,frame=single,framerule=1.2pt,commandchars=\\@\$,label=\texttt{DualArrayDeque},labelposition=topline]
	@\color@keyword$void$ add(@\color@keyword$int$ @\color@var$i$, @\color@keyword$T$ @\color@var$x$) {
		@\color@keyword$if$ (@\color@var$i$ < @\color@var$front$.size()) { 
			@\color@var$front$.add(@\color@var$front$.size()-@\color@var$i$, @\color@var$x$);
		} @\color@keyword$else$ {
			@\color@var$back$.add(@\color@var$i$-@\color@var$front$.size(), @\color@var$x$);
		}
		balance();
	}
\end{Verbatim}


The \ensuremath{\mathtt{add({\color{var}i},{\color{var}x})}} method performs rebalancing of the two \texttt{ArrayStack}s
\ensuremath{\mathtt{{\color{var}front}}} and \ensuremath{\mathtt{{\color{var}back}}}, by calling the \ensuremath{\mathtt{balance()}} method.  The
implementation of \ensuremath{\mathtt{balance()}} is described below, but for now it is
sufficient to know that \ensuremath{\mathtt{balance()}} ensures that, unless $\ensuremath{\mathtt{size()}}<2$,
\ensuremath{\mathtt{{\color{var}front}.size()}} and \ensuremath{\mathtt{{\color{var}back}.size()}} do not differ by more than a factor
of 3.  In particular, $3\cdot\ensuremath{\mathtt{{\color{var}front}.size()}} \ge \ensuremath{\mathtt{{\color{var}back}.size()}}$ and
$3\cdot\ensuremath{\mathtt{{\color{var}back}.size()}} \ge \ensuremath{\mathtt{{\color{var}front}.size()}}$.

Next we analyze the cost of \ensuremath{\mathtt{add({\color{var}i},{\color{var}x})}}, ignoring the cost of calls to
\ensuremath{\mathtt{balance()}}. If $\ensuremath{\mathtt{{\color{var}i}}}<\ensuremath{\mathtt{{\color{var}front}.size()}}$, then \ensuremath{\mathtt{add({\color{var}i},{\color{var}x})}} gets implemented
by the call to $\ensuremath{\mathtt{{\color{var}front}.add({\color{var}front}.size()-{\color{var}i}-1,{\color{var}x})}}$.  Since \ensuremath{\mathtt{{\color{var}front}}} is an
\texttt{ArrayStack}, the cost of this is
\begin{equation}
  O(\ensuremath{\mathtt{{\color{var}front}.size()}}-(\ensuremath{\mathtt{{\color{var}front}.size()}}-\ensuremath{\mathtt{{\color{var}i}}}-1)+1) = O(\ensuremath{\mathtt{{\color{var}i}}}+1) \enspace .
  \eqlabel{das-front}
\end{equation}
On the other hand, if $\ensuremath{\mathtt{{\color{var}i}}}\ge\ensuremath{\mathtt{{\color{var}front}.size()}}$, then \ensuremath{\mathtt{add({\color{var}i},{\color{var}x})}} gets
implemented as $\ensuremath{\mathtt{{\color{var}back}.add({\color{var}i}-{\color{var}front}.size(),{\color{var}x})}}$.  The cost of this is
\begin{equation}
  O(\ensuremath{\mathtt{{\color{var}back}.size()}}-(\ensuremath{\mathtt{{\color{var}i}}}-\ensuremath{\mathtt{{\color{var}front}.size()}})+1) = O(\ensuremath{\mathtt{{\color{var}n}}}-\ensuremath{\mathtt{{\color{var}i}}}+1) \enspace .
  \eqlabel{das-back}
\end{equation}

Notice that the first case \myeqref{das-front} occurs when $\ensuremath{\mathtt{{\color{var}i}}}<\ensuremath{\mathtt{{\color{var}n}}}/4$.
The second case \myeqref{das-back} occurs when $\ensuremath{\mathtt{{\color{var}i}}}\ge 3\ensuremath{\mathtt{{\color{var}n}}}/4$.  When
$\ensuremath{\mathtt{{\color{var}n}}}/4\le\ensuremath{\mathtt{{\color{var}i}}}<3\ensuremath{\mathtt{{\color{var}n}}}/4$, we cannot be sure whether the operation affects
\ensuremath{\mathtt{{\color{var}front}}} or \ensuremath{\mathtt{{\color{var}back}}}, but in either case, the operation takes
$O(\ensuremath{\mathtt{{\color{var}n}}})=O(\ensuremath{\mathtt{{\color{var}i}}})=O(\ensuremath{\mathtt{{\color{var}n}}}-\ensuremath{\mathtt{{\color{var}i}}})$ time, since $\ensuremath{\mathtt{{\color{var}i}}}\ge \ensuremath{\mathtt{{\color{var}n}}}/4$ and $\ensuremath{\mathtt{{\color{var}n}}}-\ensuremath{\mathtt{{\color{var}i}}}>
\ensuremath{\mathtt{{\color{var}n}}}/4$.  Summarizing the situation, we have
\[
     \mbox{Running time of } \ensuremath{\mathtt{add({\color{var}i},{\color{var}x})}} \le 
          \begin{cases}
            O(1+ \ensuremath{\mathtt{{\color{var}i}}}) & \mbox{if $\ensuremath{\mathtt{{\color{var}i}}}< \ensuremath{\mathtt{{\color{var}n}}}/4$} \\
            O(\ensuremath{\mathtt{{\color{var}n}}}) & \mbox{if $\ensuremath{\mathtt{{\color{var}n}}}/4 \le \ensuremath{\mathtt{{\color{var}i}}} < 3\ensuremath{\mathtt{{\color{var}n}}}/4$} \\
            O(1+\ensuremath{\mathtt{{\color{var}n}}}-\ensuremath{\mathtt{{\color{var}i}}}) & \mbox{if $\ensuremath{\mathtt{{\color{var}i}}} \ge 3\ensuremath{\mathtt{{\color{var}n}}}/4$}
          \end{cases}
\]
Thus, the running time of \ensuremath{\mathtt{add({\color{var}i},{\color{var}x})}}, if we ignore the cost of the call
to \ensuremath{\mathtt{balance()}}, is $O(1+\min\{\ensuremath{\mathtt{{\color{var}i}}}, \ensuremath{\mathtt{{\color{var}n}}}-\ensuremath{\mathtt{{\color{var}i}}}\})$.

The \ensuremath{\mathtt{remove({\color{var}i})}} operation and its analysis resemble the \ensuremath{\mathtt{add({\color{var}i},{\color{var}x})}}
operation and analysis.



 \ \vspace{.5ex}\begin{Verbatim}[gobble=1,tabsize=2,frame=single,framerule=1.2pt,commandchars=\\@\$,label=\texttt{DualArrayDeque},labelposition=topline]
	@\color@keyword$T$ remove(@\color@keyword$int$ @\color@var$i$) {
		@\color@keyword$T$ @\color@var$x$;
		@\color@keyword$if$ (@\color@var$i$ < @\color@var$front$.size()) {
			@\color@var$x$ = @\color@var$front$.remove(@\color@var$front$.size()-@\color@var$i$-1);
		} @\color@keyword$else$ {
			@\color@var$x$ = @\color@var$back$.remove(@\color@var$i$-@\color@var$front$.size());
		}
		balance();
		@\color@var$return$ @\color@var$x$;
	}
\end{Verbatim}


\subsection{Balancing}

Finally, we turn to the \ensuremath{\mathtt{balance()}} operation performed by \ensuremath{\mathtt{add({\color{var}i},{\color{var}x})}}
and \ensuremath{\mathtt{remove({\color{var}i})}}.  This operation ensures that neither \ensuremath{\mathtt{{\color{var}front}}} nor \ensuremath{\mathtt{{\color{var}back}}}
becomes too big (or too small).  It ensures that, unless there are fewer
than two elements, each of \ensuremath{\mathtt{{\color{var}front}}} and \ensuremath{\mathtt{{\color{var}back}}} contain at least $\ensuremath{\mathtt{{\color{var}n}}}/4$
elements. If this is not the case, then it moves elements between them
so that \ensuremath{\mathtt{{\color{var}front}}} and \ensuremath{\mathtt{{\color{var}back}}} contain exactly $\lfloor\ensuremath{\mathtt{{\color{var}n}}}/2\rfloor$ elements
and $\lceil\ensuremath{\mathtt{{\color{var}n}}}/2\rceil$ elements, respectively.



 \ \vspace{.5ex}\begin{Verbatim}[gobble=1,tabsize=2,frame=single,framerule=1.2pt,commandchars=\\@\$,label=\texttt{DualArrayDeque},labelposition=topline]
	@\color@keyword$void$ balance() {
		@\color@keyword$int$ @\color@var$n$ = size();
		@\color@keyword$if$ (3*@\color@var$front$.size() < @\color@var$back$.size()) {
			@\color@keyword$int$ @\color@var$s$ = @\color@var$n$/2 - @\color@var$front$.size();
			List<@\color@keyword$T$> @\color@var$l1$ = newStack();
			List<@\color@keyword$T$> @\color@var$l2$ = newStack();
			@\color@var$l1$.addAll(@\color@var$back$.subList(0,@\color@var$s$));
			Collections.reverse(@\color@var$l1$);
			@\color@var$l1$.addAll(@\color@var$front$);
			@\color@var$l2$.addAll(@\color@var$back$.subList(@\color@var$s$, @\color@var$back$.size()));
			@\color@var$front$ = @\color@var$l1$;
			@\color@var$back$ = @\color@var$l2$;
		} @\color@keyword$else$ @\color@keyword$if$ (3*@\color@var$back$.size() < @\color@var$front$.size()) {
			@\color@keyword$int$ @\color@var$s$ = @\color@var$front$.size() - @\color@var$n$/2;
			List<@\color@keyword$T$> @\color@var$l1$ = newStack();
			List<@\color@keyword$T$> @\color@var$l2$ = newStack();
			@\color@var$l1$.addAll(@\color@var$front$.subList(@\color@var$s$, @\color@var$front$.size()));
			@\color@var$l2$.addAll(@\color@var$front$.subList(0, @\color@var$s$));
			Collections.reverse(@\color@var$l2$);
			@\color@var$l2$.addAll(@\color@var$back$);
			@\color@var$front$ = @\color@var$l1$;
			@\color@var$back$ = @\color@var$l2$;
		}
	}
\end{Verbatim}


Here there is little to analyze.  If the \ensuremath{\mathtt{balance()}} operation does
rebalancing, then it moves $O(\ensuremath{\mathtt{{\color{var}n}}})$ elements and this takes $O(\ensuremath{\mathtt{{\color{var}n}}})$
time. This is bad, since \ensuremath{\mathtt{balance()}} is called with each call to
\ensuremath{\mathtt{add({\color{var}i},{\color{var}x})}} and \ensuremath{\mathtt{remove({\color{var}i})}}.  However, the following lemma shows that, on
average, \ensuremath{\mathtt{balance()}} only spends a constant amount of time per operation.

\begin{lem}\lemlabel{dualarraydeque-amortized}
  If an empty \texttt{DualArrayDeque} is created and any sequence of $m\ge 1$ calls
  to \ensuremath{\mathtt{add({\color{var}i},{\color{var}x})}} and \ensuremath{\mathtt{remove({\color{var}i})}} are performed, then the total time spent
  during all calls to \ensuremath{\mathtt{balance()}} is $O(m)$.
\end{lem}

\begin{proof}
  We will show that, if \ensuremath{\mathtt{balance()}} is forced to shift elements, then
  the number of \ensuremath{\mathtt{add({\color{var}i},{\color{var}x})}} and \ensuremath{\mathtt{remove({\color{var}i})}} operations since the last
  time any elements were shifted by \ensuremath{\mathtt{balance()}} is at least $\ensuremath{\mathtt{{\color{var}n}}}/2-1$.
  As in the proof of \lemref{arraystack-amortized}, this is sufficient
  to prove that the total time spent by \ensuremath{\mathtt{balance()}} is $O(m)$.

  We will perform our analysis using a technique knows as the
  \emph{potential method}.
  \index{potential}%
  \index{potential method}%
  Define the \emph{potential}, $\Phi$, of the
  \texttt{DualArrayDeque} as the difference in size between \ensuremath{\mathtt{{\color{var}front}}} and \ensuremath{\mathtt{{\color{var}back}}}:
  \[  \Phi = |\ensuremath{\mathtt{{\color{var}front}.size()}} - \ensuremath{\mathtt{{\color{var}back}.size()}}| \enspace . \]
  The interesting thing about this potential is that a call to \ensuremath{\mathtt{add({\color{var}i},{\color{var}x})}}
  or \ensuremath{\mathtt{remove({\color{var}i})}} that does not do any balancing can increase the potential
  by at most 1.

  Observe that, immediately after a call to \ensuremath{\mathtt{balance()}} that shifts
  elements, the potential, $\Phi_0$, is at most 1, since
  \[ \Phi_0 = \left|\lfloor\ensuremath{\mathtt{{\color{var}n}}}/2\rfloor-\lceil\ensuremath{\mathtt{{\color{var}n}}}/2\rceil\right|\le 1  \enspace .\]

  Consider the situation immediately before a call to \ensuremath{\mathtt{balance()}} that
  shifts elements and suppose, without loss of generality, that \ensuremath{\mathtt{balance()}}
  is shifting elements because $3\ensuremath{\mathtt{{\color{var}front}.size()}} < \ensuremath{\mathtt{{\color{var}back}.size()}}$.
  Notice that, in this case,
  \begin{eqnarray*}
   \ensuremath{\mathtt{{\color{var}n}}} & = & \ensuremath{\mathtt{{\color{var}front}.size()}}+\ensuremath{\mathtt{{\color{var}back}.size()}} \\
       & < & \ensuremath{\mathtt{{\color{var}back}.size()}}/3+\ensuremath{\mathtt{{\color{var}back}.size()}} \\
       & = & \frac{4}{3}\ensuremath{\mathtt{{\color{var}back}.size()}}
  \end{eqnarray*}
  Furthermore, the potential at this point in time is
  \begin{eqnarray*}
  \Phi_1 & = & \ensuremath{\mathtt{{\color{var}back}.size()}} - \ensuremath{\mathtt{{\color{var}front}.size()}} \\
      &>& \ensuremath{\mathtt{{\color{var}back}.size()}} - \ensuremath{\mathtt{{\color{var}back}.size()}}/3 \\
      &=& \frac{2}{3}\ensuremath{\mathtt{{\color{var}back}.size()}} \\
      &>& \frac{2}{3}\times\frac{3}{4}\ensuremath{\mathtt{{\color{var}n}}} \\
      &=& \ensuremath{\mathtt{{\color{var}n}}}/2
  \end{eqnarray*}
  Therefore, the number of calls to \ensuremath{\mathtt{add({\color{var}i},{\color{var}x})}} or \ensuremath{\mathtt{remove({\color{var}i})}} since
  the last time \ensuremath{\mathtt{balance()}} shifted elements is at least $\Phi_1-\Phi_0
  > \ensuremath{\mathtt{{\color{var}n}}}/2-1$. This completes the proof.
\end{proof}

\subsection{Summary}

The following theorem summarizes the properties of a \texttt{DualArrayDeque}:

\begin{thm}\thmlabel{dualarraydeque}
  A \texttt{DualArrayDeque} implements the \texttt{List} interface.  Ignoring the
  cost of calls to \ensuremath{\mathtt{resize()}} and \ensuremath{\mathtt{balance()}}, a \texttt{DualArrayDeque}
  supports the operations
  \begin{itemize}
    \item \ensuremath{\mathtt{get({\color{var}i})}} and \ensuremath{\mathtt{set({\color{var}i},{\color{var}x})}} in $O(1)$ time per operation; and
    \item \ensuremath{\mathtt{add({\color{var}i},{\color{var}x})}} and \ensuremath{\mathtt{remove({\color{var}i})}} in $O(1+\min\{\ensuremath{\mathtt{{\color{var}i}}},\ensuremath{\mathtt{{\color{var}n}}}-\ensuremath{\mathtt{{\color{var}i}}}\})$ time
          per operation.
  \end{itemize}
  Furthermore, beginning with an empty \texttt{DualArrayDeque}, any sequence of $m$
  \ensuremath{\mathtt{add({\color{var}i},{\color{var}x})}} and \ensuremath{\mathtt{remove({\color{var}i})}} operations results in a total of $O(m)$
  time spent during all calls to \ensuremath{\mathtt{resize()}} and \ensuremath{\mathtt{balance()}}.
\end{thm}


\section{\texttt{RootishArrayStack}: A Space-Efficient Array Stack}
\seclabel{rootisharraystack}

\index{RootishArrayStack@\texttt{RootishArrayStack}}%
One of the drawbacks of all previous data structures in this chapter is
that, because they store their data in one or two arrays and they avoid
resizing these arrays too often, the arrays frequently are not very full.
For example, immediately after a \ensuremath{\mathtt{resize()}} operation on an \texttt{ArrayStack},
the backing array \ensuremath{\mathtt{{\color{var}a}}} is only half full.  Even worse, there are times
when only one third of \ensuremath{\mathtt{{\color{var}a}}} contains data.

In this section, we discuss the \texttt{RootishArrayStack} data structure,
that addresses the problem of wasted space.  The \texttt{RootishArrayStack}
stores \ensuremath{\mathtt{{\color{var}n}}} elements using $O(\sqrt{\ensuremath{\mathtt{{\color{var}n}}}})$ arrays.  In these arrays, at
most $O(\sqrt{\ensuremath{\mathtt{{\color{var}n}}}})$ array locations are unused at any time.  All
remaining array locations are used to store data.  Therefore, these
data structures waste at most $O(\sqrt{\ensuremath{\mathtt{{\color{var}n}}}})$ space when storing \ensuremath{\mathtt{{\color{var}n}}}
elements.

A \texttt{RootishArrayStack} stores its elements in a list of \ensuremath{\mathtt{{\color{var}r}}}
arrays called \emph{blocks} that are numbered $0,1,\ldots,\ensuremath{\mathtt{{\color{var}r}}}-1$.
See \figref{rootisharraystack}.  Block $b$ contains $b+1$ elements.
Therefore, all \ensuremath{\mathtt{{\color{var}r}}} blocks contain a total of
\[
  1+ 2+ 3+\cdots +\ensuremath{\mathtt{{\color{var}r}}} = \ensuremath{\mathtt{{\color{var}r}}}(\ensuremath{\mathtt{{\color{var}r}}}+1)/2
\]
elements.  The above formula can be obtained as shown in \figref{gauss}.

\begin{figure}
  \begin{center}
    \includegraphics[width=\ScaleIfNeeded]{figs/rootisharraystack}
  \end{center}
  \caption[Adding and removing in a RootishArrayStack]{A sequence of \ensuremath{\mathtt{add({\color{var}i},{\color{var}x})}} and \ensuremath{\mathtt{remove({\color{var}i})}} operations on a
  \texttt{RootishArrayStack}.  Arrows denote elements being copied. }
  \figlabel{rootisharraystack}
\end{figure}



 \ \vspace{.5ex}\begin{Verbatim}[gobble=1,tabsize=2,frame=single,framerule=1.2pt,commandchars=\\@\$,label=\texttt{RootishArrayStack},labelposition=topline]
	List<@\color@keyword$T$[]> @\color@var$blocks$;
	@\color@keyword$int$ @\color@var$n$;
\end{Verbatim}


\begin{figure}
  \begin{center}
    \includegraphics[scale=0.90909]{figs/gauss}
  \end{center}
  \caption{The number of white squares is $1+2+3+\cdots+\ensuremath{\mathtt{{\color{var}r}}}$.  The number of
  shaded squares is the same.  Together the white and shaded squares make a
  rectangle consisting of $\ensuremath{\mathtt{{\color{var}r}}}(\ensuremath{\mathtt{{\color{var}r}}}+1)$ squares.}
  \figlabel{gauss}
\end{figure}

As we might expect, the elements of the list are laid out in order
within the blocks.  The list element with index 0 is stored in block 0,
elements with list indices 1 and 2 are stored in block 1, elements with
list indices 3, 4, and 5 are stored in block 2, and so on.  The main
problem we have to address is that of determining, given an index $\ensuremath{\mathtt{{\color{var}i}}}$,
which block contains \ensuremath{\mathtt{{\color{var}i}}} as well as the index corresponding to \ensuremath{\mathtt{{\color{var}i}}}
within that block.

Determining the index of \ensuremath{\mathtt{{\color{var}i}}} within its block turns out to be easy. If
index \ensuremath{\mathtt{{\color{var}i}}} is in block \ensuremath{\mathtt{{\color{var}b}}}, then the number of elements in blocks
$0,\ldots,\ensuremath{\mathtt{{\color{var}b}}}-1$ is $\ensuremath{\mathtt{{\color{var}b}}}(\ensuremath{\mathtt{{\color{var}b}}}+1)/2$.  Therefore, \ensuremath{\mathtt{{\color{var}i}}} is stored at location
\[
     \ensuremath{\mathtt{{\color{var}j}}} = \ensuremath{\mathtt{{\color{var}i}}} - \ensuremath{\mathtt{{\color{var}b}}}(\ensuremath{\mathtt{{\color{var}b}}}+1)/2
\]
within block \ensuremath{\mathtt{{\color{var}b}}}.  Somewhat more challenging is the problem of determining
the value of \ensuremath{\mathtt{{\color{var}b}}}.  The number of elements that have indices less than
or equal to \ensuremath{\mathtt{{\color{var}i}}} is $\ensuremath{\mathtt{{\color{var}i}}}+1$.  On the other hand, the number of elements
in blocks $0,\ldots,b$ is $(\ensuremath{\mathtt{{\color{var}b}}}+1)(\ensuremath{\mathtt{{\color{var}b}}}+2)/2$.  Therefore, \ensuremath{\mathtt{{\color{var}b}}} is the smallest
integer such that
\[
    (\ensuremath{\mathtt{{\color{var}b}}}+1)(\ensuremath{\mathtt{{\color{var}b}}}+2)/2 \ge \ensuremath{\mathtt{{\color{var}i}}}+1 \enspace .
\]
We can rewrite this equation as
\[
    \ensuremath{\mathtt{{\color{var}b}}}^2 + 3\ensuremath{\mathtt{{\color{var}b}}} - 2\ensuremath{\mathtt{{\color{var}i}}} \ge  0 \enspace .
\]
The corresponding quadratic equation $\ensuremath{\mathtt{{\color{var}b}}}^2 + 3\ensuremath{\mathtt{{\color{var}b}}} - 2\ensuremath{\mathtt{{\color{var}i}}} =  0$ has two
solutions: $\ensuremath{\mathtt{{\color{var}b}}}=(-3 + \sqrt{9+8\ensuremath{\mathtt{{\color{var}i}}}}) / 2$ and $\ensuremath{\mathtt{{\color{var}b}}}=(-3 - \sqrt{9+8\ensuremath{\mathtt{{\color{var}i}}}}) / 2$.
The second solution makes no sense in our application since it always
gives a negative value. Therefore, we obtain the solution $\ensuremath{\mathtt{{\color{var}b}}} = (-3 +
\sqrt{9+8i}) / 2$.  In general, this solution is not an integer, but
going back to our inequality, we want the smallest integer $\ensuremath{\mathtt{{\color{var}b}}}$ such that 
$\ensuremath{\mathtt{{\color{var}b}}} \ge (-3 + \sqrt{9+8i}) / 2$.  This is simply
\[
   \ensuremath{\mathtt{{\color{var}b}}} = \left\lceil(-3 + \sqrt{9+8i}) / 2\right\rceil \enspace .
\]



 \ \vspace{.5ex}\begin{Verbatim}[gobble=1,tabsize=2,frame=single,framerule=1.2pt,commandchars=\\@\$,label=\texttt{RootishArrayStack},labelposition=topline]
	 @\color@keyword$int$ i2b(@\color@keyword$int$ @\color@var$i$) {
		@\color@keyword$double$ @\color@var$db$ = (-3.0 + Math.sqrt(9 + 8*@\color@var$i$)) / 2.0;
		@\color@keyword$int$ @\color@var$b$ = (@\color@keyword$int$)Math.ceil(@\color@var$db$);
		@\color@var$return$ @\color@var$b$; 
	}
\end{Verbatim}


With this out of the way, the \ensuremath{\mathtt{get({\color{var}i})}} and \ensuremath{\mathtt{set({\color{var}i},{\color{var}x})}} methods are straightforward.  We first compute the appropriate block \ensuremath{\mathtt{{\color{var}b}}} and the appropriate index \ensuremath{\mathtt{{\color{var}j}}} within the block and then perform the appropriate operation:



 \ \vspace{.5ex}\begin{Verbatim}[gobble=1,tabsize=2,frame=single,framerule=1.2pt,commandchars=\\@\$,label=\texttt{RootishArrayStack},labelposition=topline]
	@\color@keyword$T$ get(@\color@keyword$int$ @\color@var$i$) {
		@\color@keyword$int$ @\color@var$b$ = i2b(@\color@var$i$);
		@\color@keyword$int$ @\color@var$j$ = @\color@var$i$ - @\color@var$b$*(@\color@var$b$+1)/2;
		@\color@var$return$ @\color@var$blocks$.get(@\color@var$b$)[@\color@var$j$];
	}
	@\color@keyword$T$ set(@\color@keyword$int$ @\color@var$i$, @\color@keyword$T$ @\color@var$x$) {
		@\color@keyword$int$ @\color@var$b$ = i2b(@\color@var$i$);
		@\color@keyword$int$ @\color@var$j$ = @\color@var$i$ - @\color@var$b$*(@\color@var$b$+1)/2;
		@\color@keyword$T$ @\color@var$y$ = @\color@var$blocks$.get(@\color@var$b$)[@\color@var$j$];
		@\color@var$blocks$.get(@\color@var$b$)[@\color@var$j$] = @\color@var$x$;
		@\color@var$return$ @\color@var$y$;
	}
\end{Verbatim}


If we use any of the data structures in this chapter for representing the \ensuremath{\mathtt{{\color{var}blocks}}} list, then \ensuremath{\mathtt{get({\color{var}i})}} and \ensuremath{\mathtt{set({\color{var}i},{\color{var}x})}} will each run in constant time.

The \ensuremath{\mathtt{add({\color{var}i},{\color{var}x})}} method will, by now, look familiar.  We first check
to see if our data structure is full, by checking if the number of
blocks, \ensuremath{\mathtt{{\color{var}r}}}, is such that $\ensuremath{\mathtt{{\color{var}r}}}(\ensuremath{\mathtt{{\color{var}r}}}+1)/2 = \ensuremath{\mathtt{{\color{var}n}}}$. If so, we call \ensuremath{\mathtt{grow()}}
to add another block.  With this done, we shift elements with indices
$\ensuremath{\mathtt{{\color{var}i}}},\ldots,\ensuremath{\mathtt{{\color{var}n}}}-1$ to the right by one position to make room for the
new element with index \ensuremath{\mathtt{{\color{var}i}}}:



 \ \vspace{.5ex}\begin{Verbatim}[gobble=1,tabsize=2,frame=single,framerule=1.2pt,commandchars=\\@\$,label=\texttt{RootishArrayStack},labelposition=topline]
	@\color@keyword$void$ add(@\color@keyword$int$ @\color@var$i$, @\color@keyword$T$ @\color@var$x$) {
		@\color@keyword$int$ @\color@var$r$ = @\color@var$blocks$.size();
		@\color@keyword$if$ (@\color@var$r$*(@\color@var$r$+1)/2 < @\color@var$n$ + 1) grow();
		@\color@var$n$++;
		@\color@keyword$for$ (@\color@keyword$int$ @\color@var$j$ = @\color@var$n$-1; @\color@var$j$ > @\color@var$i$; @\color@var$j$--)
			set(@\color@var$j$, get(@\color@var$j$-1));
		set(@\color@var$i$, @\color@var$x$);
	}
\end{Verbatim}


The \ensuremath{\mathtt{grow()}} method does what we expect. It adds a new block:



 \ \vspace{.5ex}\begin{Verbatim}[gobble=1,tabsize=2,frame=single,framerule=1.2pt,commandchars=\\@\$,label=\texttt{RootishArrayStack},labelposition=topline]
	@\color@keyword$void$ grow() {
		@\color@var$blocks$.add(newArray(@\color@var$blocks$.size()+1));
	}
\end{Verbatim}


Ignoring the cost of the \ensuremath{\mathtt{grow()}} operation, the cost of an \ensuremath{\mathtt{add({\color{var}i},{\color{var}x})}}
operation is dominated by the cost of shifting and is therefore
$O(1+\ensuremath{\mathtt{{\color{var}n}}}-\ensuremath{\mathtt{{\color{var}i}}})$, just like an \texttt{ArrayStack}.

The \ensuremath{\mathtt{remove({\color{var}i})}} operation is similar to \ensuremath{\mathtt{add({\color{var}i},{\color{var}x})}}.  It shifts the
elements with indices $\ensuremath{\mathtt{{\color{var}i}}}+1,\ldots,\ensuremath{\mathtt{{\color{var}n}}}$ left by one position and then,
if there is more than one empty block, it calls the \ensuremath{\mathtt{shrink()}} method
to remove all but one of the unused blocks:



 \ \vspace{.5ex}\begin{Verbatim}[gobble=1,tabsize=2,frame=single,framerule=1.2pt,commandchars=\\@\$,label=\texttt{RootishArrayStack},labelposition=topline]
	@\color@keyword$T$ remove(@\color@keyword$int$ @\color@var$i$) {
		@\color@keyword$T$ @\color@var$x$ = get(@\color@var$i$);
		@\color@keyword$for$ (@\color@keyword$int$ @\color@var$j$ = @\color@var$i$; @\color@var$j$ < @\color@var$n$-1; @\color@var$j$++)
			set(@\color@var$j$, get(@\color@var$j$+1));
		@\color@var$n$--;
		@\color@keyword$int$ @\color@var$r$ = @\color@var$blocks$.size();
		@\color@keyword$if$ ((@\color@var$r$-2)*(@\color@var$r$-1)/2 >= @\color@var$n$)	shrink();
		@\color@var$return$ @\color@var$x$;
	}
\end{Verbatim}



 \ \vspace{.5ex}\begin{Verbatim}[gobble=1,tabsize=2,frame=single,framerule=1.2pt,commandchars=\\@\$,label=\texttt{RootishArrayStack},labelposition=topline]
	@\color@keyword$void$ shrink() {
		@\color@keyword$int$ @\color@var$r$ = @\color@var$blocks$.size();
		@\color@keyword$while$ (@\color@var$r$ > 0 && (@\color@var$r$-2)*(@\color@var$r$-1)/2 >= @\color@var$n$) {
			@\color@var$blocks$.remove(@\color@var$blocks$.size()-1);
			@\color@var$r$--;
		}
	}
\end{Verbatim}


Once again, ignoring the cost of the \ensuremath{\mathtt{shrink()}} operation, the cost of
a \ensuremath{\mathtt{remove({\color{var}i})}} operation is dominated by the cost of shifting  and is
therefore $O(\ensuremath{\mathtt{{\color{var}n}}}-\ensuremath{\mathtt{{\color{var}i}}})$.

\subsection{Analysis of Growing and Shrinking}

The above analysis of \ensuremath{\mathtt{add({\color{var}i},{\color{var}x})}} and \ensuremath{\mathtt{remove({\color{var}i})}} does not account for
the cost of \ensuremath{\mathtt{grow()}} and \ensuremath{\mathtt{shrink()}}.  Note that, unlike the
\ensuremath{\mathtt{ArrayStack.resize()}} operation, \ensuremath{\mathtt{grow()}} and \ensuremath{\mathtt{shrink()}} do not copy
any data.  They only allocate or free an array of size \ensuremath{\mathtt{{\color{var}r}}}.  In
some environments, this takes only constant time, while in others, it
may require time proportional to \ensuremath{\mathtt{{\color{var}r}}}.

We note that, immediately after a call to \ensuremath{\mathtt{grow()}} or \ensuremath{\mathtt{shrink()}}, the
situation is clear. The final block is completely empty, and all other
blocks are completely full.  Another call to \ensuremath{\mathtt{grow()}} or \ensuremath{\mathtt{shrink()}} will
not happen until at least $\ensuremath{\mathtt{{\color{var}r}}}-1$ elements have been added or removed.
Therefore, even if \ensuremath{\mathtt{grow()}} and \ensuremath{\mathtt{shrink()}} take $O(\ensuremath{\mathtt{{\color{var}r}}})$ time, this
cost can be amortized over at least $\ensuremath{\mathtt{{\color{var}r}}}-1$ \ensuremath{\mathtt{add({\color{var}i},{\color{var}x})}} or \ensuremath{\mathtt{remove({\color{var}i})}}
operations, so that the amortized cost of \ensuremath{\mathtt{grow()}} and \ensuremath{\mathtt{shrink()}} is
$O(1)$ per operation.

\subsection{Space Usage}
\seclabel{rootishspaceusage}

Next, we analyze the amount of extra space used by a \texttt{RootishArrayStack}.
In particular, we want to count any space used by a \texttt{RootishArrayStack} that is not an array element currently used to hold a list element.  We call all such space \emph{wasted space}.
\index{wasted space}%

The \ensuremath{\mathtt{remove({\color{var}i})}} operation ensures that a \texttt{RootishArrayStack} never has
more than two blocks that are not completely full.  The number of blocks,
\ensuremath{\mathtt{{\color{var}r}}}, used by a \texttt{RootishArrayStack} that stores \ensuremath{\mathtt{{\color{var}n}}} elements therefore
satisfies
\[
    (\ensuremath{\mathtt{{\color{var}r}}}-2)(\ensuremath{\mathtt{{\color{var}r}}}-1)/2 \le \ensuremath{\mathtt{{\color{var}n}}} \enspace .
\]
Again, using the quadratic equation on this gives
\[
   \ensuremath{\mathtt{{\color{var}r}}} \le \frac{1}{2}\left(3+\sqrt{8\ensuremath{\mathtt{{\color{var}n}}}+1}\right) = O(\sqrt{\ensuremath{\mathtt{{\color{var}n}}}}) \enspace .
\]
The last two blocks have sizes \ensuremath{\mathtt{{\color{var}r}}} and \ensuremath{\mathtt{{\color{var}r}-1}}, so the space wasted by these
two blocks is at most $2\ensuremath{\mathtt{{\color{var}r}}}-1 = O(\sqrt{\ensuremath{\mathtt{{\color{var}n}}}})$.  If we store the blocks
in (for example) an \texttt{ArrayStack}, then the amount of space wasted by the
\texttt{List} that stores those \ensuremath{\mathtt{{\color{var}r}}} blocks is also $O(\ensuremath{\mathtt{{\color{var}r}}})=O(\sqrt{\ensuremath{\mathtt{{\color{var}n}}}})$.  The
other space needed for storing \ensuremath{\mathtt{{\color{var}n}}} and other accounting information is $O(1)$.
Therefore, the total amount of wasted space in a \texttt{RootishArrayStack}
is $O(\sqrt{\ensuremath{\mathtt{{\color{var}n}}}})$.

Next, we argue that this space usage is optimal for any data structure
that starts out empty and can support the addition of one item at
a time. More precisely, we will show that, at some point during the
addition of \ensuremath{\mathtt{{\color{var}n}}} items, the data structure is wasting at least in $\sqrt{\ensuremath{\mathtt{{\color{var}n}}}}$ space (though it may be only wasted for a moment).

Suppose we start with an empty data structure and we add \ensuremath{\mathtt{{\color{var}n}}} items one
at a time.  At the end of this process, all \ensuremath{\mathtt{{\color{var}n}}} items are stored in
the structure and distributed among a collection of \ensuremath{\mathtt{{\color{var}r}}} memory blocks.
If $\ensuremath{\mathtt{{\color{var}r}}}\ge \sqrt{\ensuremath{\mathtt{{\color{var}n}}}}$, then the data structure must be using \ensuremath{\mathtt{{\color{var}r}}}
pointers (or references) to keep track of these \ensuremath{\mathtt{{\color{var}r}}} blocks, and these
pointers are wasted space.  On the other hand, if $\ensuremath{\mathtt{{\color{var}r}}} < \sqrt{\ensuremath{\mathtt{{\color{var}n}}}}$
then, by the pigeonhole principle, some block must have size at
least $\ensuremath{\mathtt{{\color{var}n}}}/\ensuremath{\mathtt{{\color{var}r}}} > \sqrt{\ensuremath{\mathtt{{\color{var}n}}}}$.  Consider the moment at which this block
was first allocated.  Immediately after it was allocated, this block
was empty, and was therefore wasting $\sqrt{\ensuremath{\mathtt{{\color{var}n}}}}$ space.  Therefore,
at some point in time during the insertion of \ensuremath{\mathtt{{\color{var}n}}} elements, the data
structure was wasting $\sqrt{\ensuremath{\mathtt{{\color{var}n}}}}$ space.

\subsection{Summary}

The following theorem summarizes our discussion of the \texttt{RootishArrayStack}
data structure:

\begin{thm}\thmlabel{rootisharraystack}
  A \texttt{RootishArrayStack} implements the \texttt{List} interface.  Ignoring the cost of
  calls to \ensuremath{\mathtt{grow()}} and \ensuremath{\mathtt{shrink()}}, a \texttt{RootishArrayStack} supports the operations
  \begin{itemize}
    \item \ensuremath{\mathtt{get({\color{var}i})}} and \ensuremath{\mathtt{set({\color{var}i},{\color{var}x})}} in $O(1)$ time per operation; and
    \item \ensuremath{\mathtt{add({\color{var}i},{\color{var}x})}} and \ensuremath{\mathtt{remove({\color{var}i})}} in $O(1+\ensuremath{\mathtt{{\color{var}n}}}-\ensuremath{\mathtt{{\color{var}i}}})$ time per operation.
  \end{itemize}
  Furthermore, beginning with an empty \texttt{RootishArrayStack}, any sequence of $m$
  \ensuremath{\mathtt{add({\color{var}i},{\color{var}x})}} and \ensuremath{\mathtt{remove({\color{var}i})}} operations results in a total of $O(m)$
  time spent during all calls to \ensuremath{\mathtt{grow()}} and \ensuremath{\mathtt{shrink()}}.

  The space (measured in words)\footnote{Recall \secref{model} for a
  discussion of how memory is measured.} used by a \texttt{RootishArrayStack}
  that stores \ensuremath{\mathtt{{\color{var}n}}} elements is $\ensuremath{\mathtt{{\color{var}n}}} +O(\sqrt{\ensuremath{\mathtt{{\color{var}n}}}})$.
\end{thm}

\notpcode{

\subsection{Computing Square Roots}

\index{square roots}%
A reader who has had some exposure to models of computation may notice
that the \texttt{RootishArrayStack}, as described above, does not fit into
the usual word-RAM model of computation (\secref{model}) because it
requires taking square roots.  The square root operation is generally
not considered a basic operation and is therefore not usually part of
the word-RAM model.

In this section, we show that the square root operation can be
implemented efficiently.  In particular, we show that for any integer
$\ensuremath{\mathtt{{\color{var}x}}}\in\{0,\ldots,\ensuremath{\mathtt{{\color{var}n}}}\}$,  $\lfloor\sqrt{\ensuremath{\mathtt{{\color{var}x}}}}\rfloor$ can be computed
in constant-time, after $O(\sqrt{\ensuremath{\mathtt{{\color{var}n}}}})$ preprocessing that creates two
arrays of length $O(\sqrt{\ensuremath{\mathtt{{\color{var}n}}}})$.  The following lemma shows that we
can reduce the problem of computing the square root of \ensuremath{\mathtt{{\color{var}x}}} to the square
root of a related value \ensuremath{\mathtt{{\color{var}x}'}}.

\begin{lem}\lemlabel{root}
Let $\ensuremath{\mathtt{{\color{var}x}}}\ge 1$ and let $\ensuremath{\mathtt{{\color{var}x}'}}=\ensuremath{\mathtt{{\color{var}x}}}-a$, where $0\le a\le\sqrt{\ensuremath{\mathtt{{\color{var}x}}}}$.  Then
   $\sqrt{x'} \ge \sqrt{\ensuremath{\mathtt{{\color{var}x}}}}-1$.
\end{lem}

\begin{proof}
It suffices to show that
\[
\sqrt{\ensuremath{\mathtt{{\color{var}x}}}-\sqrt{\ensuremath{\mathtt{{\color{var}x}}}}} \ge \sqrt{\ensuremath{\mathtt{{\color{var}x}}}}-1 \enspace .
\]
Square both sides of this inequality to get
\[
 \ensuremath{\mathtt{{\color{var}x}}}-\sqrt{\ensuremath{\mathtt{{\color{var}x}}}} \ge \ensuremath{\mathtt{{\color{var}x}}}-2\sqrt{\ensuremath{\mathtt{{\color{var}x}}}}+1
\]
and gather terms to get 
\[
 \sqrt{\ensuremath{\mathtt{{\color{var}x}}}} \ge 1
\]
which is clearly true for any $\ensuremath{\mathtt{{\color{var}x}}}\ge 1$.
\end{proof}

Start by restricting the problem a little, and assume that $2^{\ensuremath{\mathtt{{\color{var}r}}}} \le
\ensuremath{\mathtt{{\color{var}x}}} < 2^{\ensuremath{\mathtt{{\color{var}r}}}+1}$, so that $\lfloor\log \ensuremath{\mathtt{{\color{var}x}}}\rfloor=\ensuremath{\mathtt{{\color{var}r}}}$, i.e., \ensuremath{\mathtt{{\color{var}x}}} is an
integer having $\ensuremath{\mathtt{{\color{var}r}}}+1$ bits in its binary representation.  We can take
$\ensuremath{\mathtt{{\color{var}x}'}}=\ensuremath{\mathtt{{\color{var}x}}} - (\ensuremath{\mathtt{{\color{var}x}}}\bmod 2^{\lfloor r/2\rfloor})$.  Now, \ensuremath{\mathtt{{\color{var}x}'}} satisfies
the conditions of \lemref{root}, so $\sqrt{\ensuremath{\mathtt{{\color{var}x}}}}-\sqrt{\ensuremath{\mathtt{{\color{var}x}'}}} \le 1$.
Furthermore, \ensuremath{\mathtt{{\color{var}x}'}} has all of its lower-order $\lfloor \ensuremath{\mathtt{{\color{var}r}}}/2\rfloor$ bits
equal to 0, so there are only
\[
  2^{\ensuremath{\mathtt{{\color{var}r}}}+1-\lfloor \ensuremath{\mathtt{{\color{var}r}}}/2\rfloor} \le 4\cdot2^{\ensuremath{\mathtt{{\color{var}r}}}/2} \le 4\sqrt{\ensuremath{\mathtt{{\color{var}x}}}}
\]
possible values of \ensuremath{\mathtt{{\color{var}x}'}}.  This means that we can use an array, \ensuremath{\mathtt{{\color{var}sqrttab}}},
that stores the value of $\lfloor\sqrt{\ensuremath{\mathtt{{\color{var}x}'}}}\rfloor$ for each possible
value of \ensuremath{\mathtt{{\color{var}x}'}}.  A little more precisely, we have
\[
   \ensuremath{\mathtt{{\color{var}sqrttab}}}[i] 
    = \left\lfloor
       \sqrt{i 2^{\lfloor \ensuremath{\mathtt{{\color{var}r}}}/2\rfloor}}
      \right\rfloor \enspace .
\]
In this way, $\ensuremath{\mathtt{{\color{var}sqrttab}}}[i]$ is within 2 of $\sqrt{\ensuremath{\mathtt{{\color{var}x}}}}$ for all
$\ensuremath{\mathtt{{\color{var}x}}}\in\{i2^{\lfloor \ensuremath{\mathtt{{\color{var}r}}}/2\rfloor},\ldots,(i+1)2^{\lfloor \ensuremath{\mathtt{{\color{var}r}}}/2\rfloor}-1\}$.
Stated another way, the array entry 
$\ensuremath{\mathtt{{\color{var}s}}}=\ensuremath{\mathtt{{\color{var}sqrttab}}}[\ensuremath{\mathtt{{\color{var}x}}}\ensuremath{\mathtt{\text{\ttfamily >>}}}\lfloor \ensuremath{\mathtt{{\color{var}r}}}/2\rfloor]$ is either equal to
$\lfloor\sqrt{\ensuremath{\mathtt{{\color{var}x}}}}\rfloor$,
$\lfloor\sqrt{\ensuremath{\mathtt{{\color{var}x}}}}\rfloor-1$, or
$\lfloor\sqrt{\ensuremath{\mathtt{{\color{var}x}}}}\rfloor-2$.  From \ensuremath{\mathtt{{\color{var}s}}} we can determine the value
of $\lfloor\sqrt{\ensuremath{\mathtt{{\color{var}x}}}}\rfloor$ by
incrementing \ensuremath{\mathtt{{\color{var}s}}} until 
$(\ensuremath{\mathtt{{\color{var}s}}}+1)^2 > \ensuremath{\mathtt{{\color{var}x}}}$.
} % notpcode


 \ \vspace{.5ex}\begin{Verbatim}[gobble=1,tabsize=2,frame=single,framerule=1.2pt,commandchars=\\@\$,label=\texttt{FastSqrt},labelposition=topline]
	@\color@keyword$int$ sqrt(@\color@keyword$int$ @\color@var$x$, @\color@keyword$int$ @\color@var$r$) {
		@\color@keyword$int$ @\color@var$s$ = @\color@var$sqrtab$[@\color@var$x$>>@\color@var$r$/2];
		@\color@keyword$while$ ((@\color@var$s$+1)*(@\color@var$s$+1) <= @\color@var$x$) @\color@var$s$++; @\color@comment$// executes at most twice$
		@\color@var$return$ @\color@var$s$;
	}
\end{Verbatim}

\notpcode{
Now, this only works for $\ensuremath{\mathtt{{\color{var}x}}}\in\{2^{\ensuremath{\mathtt{{\color{var}r}}}},\ldots,2^{\ensuremath{\mathtt{{\color{var}r}}}+1}-1\}$ and
\ensuremath{\mathtt{{\color{var}sqrttab}}} is a special table that only works for a particular value
of $\ensuremath{\mathtt{{\color{var}r}}}=\lfloor\log \ensuremath{\mathtt{{\color{var}x}}}\rfloor$.  To overcome this, we could compute
$\lfloor\log \ensuremath{\mathtt{{\color{var}n}}}\rfloor$ different \ensuremath{\mathtt{{\color{var}sqrttab}}} arrays, one for each possible
value of $\lfloor\log \ensuremath{\mathtt{{\color{var}x}}}\rfloor$. The sizes of these tables form an exponential sequence whose largest value is at most $4\sqrt{\ensuremath{\mathtt{{\color{var}n}}}}$, so the total size of all tables is $O(\sqrt{\ensuremath{\mathtt{{\color{var}n}}}})$.

However, it turns out that more than one \ensuremath{\mathtt{{\color{var}sqrttab}}} array is unnecessary;
we only need one \ensuremath{\mathtt{{\color{var}sqrttab}}} array for the value $\ensuremath{\mathtt{{\color{var}r}}}=\lfloor\log
\ensuremath{\mathtt{{\color{var}n}}}\rfloor$.  Any value \ensuremath{\mathtt{{\color{var}x}}} with $\log\ensuremath{\mathtt{{\color{var}x}}}=\ensuremath{\mathtt{{\color{var}r}'}}<\ensuremath{\mathtt{{\color{var}r}}}$ can be \emph{upgraded}
by multiplying \ensuremath{\mathtt{{\color{var}x}}} by $2^{\ensuremath{\mathtt{{\color{var}r}}}-\ensuremath{\mathtt{{\color{var}r}'}}}$ and using the equation
\[
    \sqrt{2^{\ensuremath{\mathtt{{\color{var}r}}}-\ensuremath{\mathtt{{\color{var}r}'}}}x} = 2^{(\ensuremath{\mathtt{{\color{var}r}}}-\ensuremath{\mathtt{{\color{var}r}}}')/2}\sqrt{\ensuremath{\mathtt{{\color{var}x}}}} \enspace .
\]
The quantity $2^{\ensuremath{\mathtt{{\color{var}r}}}-\ensuremath{\mathtt{{\color{var}r}}}'}x$ is in the range
$\{2^{\ensuremath{\mathtt{{\color{var}r}}}},\ldots,2^{\ensuremath{\mathtt{{\color{var}r}}}+1}-1\}$ so we can look up its square root
in \ensuremath{\mathtt{{\color{var}sqrttab}}}.  The following code implements this idea to compute
$\lfloor\sqrt{\ensuremath{\mathtt{{\color{var}x}}}}\rfloor$ for all non-negative integers \ensuremath{\mathtt{{\color{var}x}}} in the
range $\{0,\ldots,2^{30}-1\}$ using an array, \ensuremath{\mathtt{{\color{var}sqrttab}}}, of size $2^{16}$.
} % notpcode


 \ \vspace{.5ex}\begin{Verbatim}[gobble=1,tabsize=2,frame=single,framerule=1.2pt,commandchars=\\@\$,label=\texttt{FastSqrt},labelposition=topline]
	@\color@keyword$int$ sqrt(@\color@keyword$int$ @\color@var$x$) {
		@\color@keyword$int$ @\color@var$rp$ = log(@\color@var$x$);
		@\color@keyword$int$ @\color@var$upgrade$ = ((@\color@var$r$-@\color@var$rp$)/2) * 2;
		@\color@keyword$int$ @\color@var$xp$ = @\color@var$x$ << @\color@var$upgrade$;  @\color@comment$// xp has r or r-1 bits$
		@\color@keyword$int$ @\color@var$s$ = @\color@var$sqrtab$[@\color@var$xp$>>(@\color@var$r$/2)] >> (@\color@var$upgrade$/2);
		@\color@keyword$while$ ((@\color@var$s$+1)*(@\color@var$s$+1) <= @\color@var$x$) @\color@var$s$++;  @\color@comment$// executes at most twice$
		@\color@var$return$ @\color@var$s$;
	}
\end{Verbatim}

\notpcode{
Something we have taken for granted thus far is the question of how
to compute
$\ensuremath{\mathtt{{\color{var}r}}}'=\lfloor\log\ensuremath{\mathtt{{\color{var}x}}}\rfloor$.  Again, this is a problem that can be solved
with an array, \ensuremath{\mathtt{{\color{var}logtab}}}, of size $2^{\ensuremath{\mathtt{{\color{var}r}}}/2}$.  In this case, the
code is particularly simple, since $\lfloor\log \ensuremath{\mathtt{{\color{var}x}}}\rfloor$ is just the
index of the most significant 1 bit in the binary representation of \ensuremath{\mathtt{{\color{var}x}}}.
This means that, for $\ensuremath{\mathtt{{\color{var}x}}}>2^{\ensuremath{\mathtt{{\color{var}r}}}/2}$, we can right-shift the bits of
\ensuremath{\mathtt{{\color{var}x}}} by $\ensuremath{\mathtt{{\color{var}r}}}/2$ positions before using it as an index into \ensuremath{\mathtt{{\color{var}logtab}}}.
The following code does this using an array \ensuremath{\mathtt{{\color{var}logtab}}} of size $2^{16}$ to compute
$\lfloor\log \ensuremath{\mathtt{{\color{var}x}}}\rfloor$ for all \ensuremath{\mathtt{{\color{var}x}}} in the range $\{1,\ldots,2^{32}-1\}$.
} % notpcode


 \ \vspace{.5ex}\begin{Verbatim}[gobble=1,tabsize=2,frame=single,framerule=1.2pt,commandchars=\\@\$,label=\texttt{FastSqrt},labelposition=topline]
	@\color@keyword$int$ log(@\color@keyword$int$ @\color@var$x$) {
		@\color@keyword$if$ (@\color@var$x$ >= @\color@var$halfint$)
			@\color@var$return$ 16 + @\color@var$logtab$[@\color@var$x$>>>16];
		@\color@var$return$ @\color@var$logtab$[@\color@var$x$];
	}
\end{Verbatim}

\notpcode{
Finally, for completeness, we include the following code that initializes \ensuremath{\mathtt{{\color{var}logtab}}} and \ensuremath{\mathtt{{\color{var}sqrttab}}}:
} % notpcode


 \ \vspace{.5ex}\begin{Verbatim}[gobble=1,tabsize=2,frame=single,framerule=1.2pt,commandchars=\\@\$,label=\texttt{FastSqrt},labelposition=topline]
	@\color@keyword$void$ inittabs() {
		@\color@var$sqrtab$ = @\color@keyword$new$ @\color@keyword$int$[1<<(@\color@var$r$/2)];
		@\color@var$logtab$ = @\color@keyword$new$ @\color@keyword$int$[1<<(@\color@var$r$/2)];
		@\color@keyword$for$ (@\color@keyword$int$ @\color@var$d$ = 0; @\color@var$d$ < @\color@var$r$/2; @\color@var$d$++) 
			Arrays.fill(@\color@var$logtab$, 1<<@\color@var$d$, 2<<@\color@var$d$, @\color@var$d$);
		@\color@keyword$int$ @\color@var$s$ = 1<<(@\color@var$r$/4);                    @\color@comment$// sqrt(2^(r/2))$
		@\color@keyword$for$ (@\color@keyword$int$ @\color@var$i$ = 0; @\color@var$i$ < 1<<(@\color@var$r$/2); @\color@var$i$++) {
			@\color@keyword$if$ ((@\color@var$s$+1)*(@\color@var$s$+1) <= @\color@var$i$ << (@\color@var$r$/2)) @\color@var$s$++; @\color@comment$// sqrt increases$
			@\color@var$sqrtab$[@\color@var$i$] = @\color@var$s$;
		}
	}
\end{Verbatim}

\notpcode{
To summarize, the computations done by the \ensuremath{\mathtt{i2b({\color{var}i})}} method can be
implemented in constant time on the word-RAM using $O(\sqrt{n})$ extra
memory to store the \ensuremath{\mathtt{{\color{var}sqrttab}}} and \ensuremath{\mathtt{{\color{var}logtab}}} arrays.  These arrays can be
rebuilt when \ensuremath{\mathtt{{\color{var}n}}} increases or decreases by a factor of two, and the cost
of this rebuilding can be amortized over the number of \ensuremath{\mathtt{add({\color{var}i},{\color{var}x})}} and
\ensuremath{\mathtt{remove({\color{var}i})}} operations that caused the change in \ensuremath{\mathtt{{\color{var}n}}} in the same way that
the cost of \ensuremath{\mathtt{resize()}} is analyzed in the \texttt{ArrayStack} implementation.
} % notpcode

\section{Discussion and Exercises}

Most of the data structures described in this chapter are folklore. They
can be found in implementations dating back over 30 years.  For example,
implementations of stacks, queues, and deques, which generalize easily
to the \texttt{ArrayStack}, \texttt{ArrayQueue} and \texttt{ArrayDeque} structures described
here, are discussed by Knuth \cite[Section~2.2.2]{k97v1}.

Brodnik \etal\ \cite{bcdms99} seem to have been the first to describe
the \texttt{RootishArrayStack} and prove a $\sqrt{n}$ lower-bound like that
in \secref{rootishspaceusage}.  They also present a different structure
that uses a more sophisticated choice of block sizes in order to avoid
computing square roots in the \ensuremath{\mathtt{i2b({\color{var}i})}} method.  Within their scheme,
the block containing \ensuremath{\mathtt{{\color{var}i}}} is block $\lfloor\log (\ensuremath{\mathtt{{\color{var}i}}}+1)\rfloor$, which
is simply the index of the leading 1 bit in the binary representation
of $\ensuremath{\mathtt{{\color{var}i}}}+1$.  Some computer architectures provide an instruction for
computing the index of the leading 1-bit in an integer. \javaonly{In
Java, the \texttt{Integer} class provides a method \ensuremath{\mathtt{numberOfLeadingZeros({\color{var}i})}}
from which one can easily compute $\lfloor\log (\ensuremath{\mathtt{{\color{var}i}}}+1)\rfloor$.}

A structure related to the \texttt{RootishArrayStack} is the two-level
\emph{tiered-vector} of Goodrich and Kloss \cite{gk99}.
\index{tiered-vector}%
This structure
supports the \ensuremath{\mathtt{get({\color{var}i},{\color{var}x})}} and \ensuremath{\mathtt{set({\color{var}i},{\color{var}x})}} operations in constant time and
\ensuremath{\mathtt{add({\color{var}i},{\color{var}x})}} and \ensuremath{\mathtt{remove({\color{var}i})}} in $O(\sqrt{\ensuremath{\mathtt{{\color{var}n}}}})$ time.  These running times
are similar to what can be achieved with the more careful implementation
of a \texttt{RootishArrayStack} discussed in \excref{rootisharraystack-fast}.

\javaonly{
\begin{exc}
  In the \texttt{ArrayStack} implementation, after the first call to \ensuremath{\mathtt{remove({\color{var}i})}},
  the backing array, \ensuremath{\mathtt{{\color{var}a}}}, contains $\ensuremath{\mathtt{{\color{var}n}}}+1$ non-\ensuremath{\mathtt{{\color{keyword}null}}} values despite
  the fact that the \texttt{ArrayStack} only contains \ensuremath{\mathtt{{\color{var}n}}} elements.  Where is
  the extra non-\ensuremath{\mathtt{{\color{keyword}null}}} value?  Discuss any consequences this non-\ensuremath{\mathtt{{\color{keyword}null}}}
  value might have on the Java Runtime Environment's memory manager.
  \index{Java Runtime Environment}%
  \index{memory manager}%
\end{exc}
}

\begin{exc}
  The \texttt{List} method \ensuremath{\mathtt{addAll({\color{var}i},{\color{var}c})}} inserts all elements of the \texttt{Collection}
  \ensuremath{\mathtt{{\color{var}c}}} into the list at position \ensuremath{\mathtt{{\color{var}i}}}.  (The \ensuremath{\mathtt{add({\color{var}i},{\color{var}x})}} method is a special
  case where $\ensuremath{\mathtt{{\color{var}c}}}=\{\ensuremath{\mathtt{{\color{var}x}}}\}$.)  Explain why, for the data structures
  in this chapter, it is not efficient to implement \ensuremath{\mathtt{addAll({\color{var}i},{\color{var}c})}} by
  repeated calls to \ensuremath{\mathtt{add({\color{var}i},{\color{var}x})}}.  Design and implement a more efficient
  implementation.
\end{exc}

\begin{exc}
  Design and implement a \emph{\texttt{RandomQueue}}.
  \index{RandomQueue@\texttt{RandomQueue}}%
  This is an implementation
  of the \texttt{Queue} interface in which the \ensuremath{\mathtt{remove()}} operation removes
  an element that is chosen uniformly at random among all the elements
  currently in the queue.  (Think of a \texttt{RandomQueue} as a bag in which
  we can add elements or reach in and blindly remove some random element.)
  The \ensuremath{\mathtt{add({\color{var}x})}} and \ensuremath{\mathtt{remove()}} operations in a \texttt{RandomQueue} should run
  in amortized constant time per operation.
\end{exc}

\begin{exc}
  Design and implement a \texttt{Treque} (triple-ended queue). 
  \index{Treque@\texttt{Treque}}%
  This is a \texttt{List}
  implementation in which \ensuremath{\mathtt{get({\color{var}i})}} and \ensuremath{\mathtt{set({\color{var}i},{\color{var}x})}} run in constant time
  and \ensuremath{\mathtt{add({\color{var}i},{\color{var}x})}} and \ensuremath{\mathtt{remove({\color{var}i})}} run in time
  \[
     O(1+\min\{\ensuremath{\mathtt{{\color{var}i}}}, \ensuremath{\mathtt{{\color{var}n}}}-\ensuremath{\mathtt{{\color{var}i}}}, |\ensuremath{\mathtt{{\color{var}n}}}/2-\ensuremath{\mathtt{{\color{var}i}}}|\}) \enspace .
  \]
  In other words, modifications are fast if they are near either
  end or near the middle of the list.
\end{exc}

\begin{exc}
  Implement a method \ensuremath{\mathtt{rotate({\color{var}a},{\color{var}r})}} that ``rotates'' the array \ensuremath{\mathtt{{\color{var}a}}}
  so that \ensuremath{\mathtt{{\color{var}a}[{\color{var}i}]}} moves to $\ensuremath{\mathtt{{\color{var}a}}}[(\ensuremath{\mathtt{{\color{var}i}}}+\ensuremath{\mathtt{{\color{var}r}}})\bmod \ensuremath{\mathtt{{\color{var}a}.{\color{var}length}}}]$, for all
  $\ensuremath{\mathtt{{\color{var}i}}}\in\{0,\ldots,\ensuremath{\mathtt{{\color{var}a}.{\color{var}length}}}\}$.
\end{exc}

\begin{exc}
  Implement a method \ensuremath{\mathtt{rotate({\color{var}r})}} that ``rotates'' a \texttt{List} so that
  list item \ensuremath{\mathtt{{\color{var}i}}} becomes list item $(\ensuremath{\mathtt{{\color{var}i}}}+\ensuremath{\mathtt{{\color{var}r}}})\bmod \ensuremath{\mathtt{{\color{var}n}}}$.  When run on
  an \texttt{ArrayDeque}, or a \texttt{DualArrayDeque}, \ensuremath{\mathtt{rotate({\color{var}r})}} should run in
  $O(1+\min\{\ensuremath{\mathtt{{\color{var}r}}},\ensuremath{\mathtt{{\color{var}n}}}-\ensuremath{\mathtt{{\color{var}r}}}\})$ time.
\end{exc}

\begin{exc}
  \pcodeonly{This exercise is left out of the \lang\ edition.}
  \notpcode{
  Modify the \texttt{ArrayDeque} implementation so that the shifting
  done by \ensuremath{\mathtt{add({\color{var}i},{\color{var}x})}}, \ensuremath{\mathtt{remove({\color{var}i})}}, and \ensuremath{\mathtt{resize()}} is done using
  the faster \ensuremath{\mathtt{System.arraycopy({\color{var}s},{\color{var}i},{\color{var}d},{\color{var}j},{\color{var}n})}} method.}
\end{exc}

\begin{exc}
  Modify the \texttt{ArrayDeque} implementation so that it does not use the
  \ensuremath{\mathtt{\text{\ttfamily \%}}} operator (which is expensive on some systems).  Instead, it
  should make use of the fact that, if \ensuremath{\mathtt{{\color{var}a}.{\color{var}length}}} is a power of 2,
  then 
  \[  \ensuremath{\mathtt{{\color{var}k}\text{\ttfamily \%}{\color{var}a}.{\color{var}length}}}=\ensuremath{\mathtt{{\color{var}k}\text{\ttfamily \&}({\color{var}a}.{\color{var}length}-1)}} \enspace .
  \]
  (Here, \ensuremath{\mathtt{\text{\ttfamily \&}}} is the bitwise-and operator.)
\end{exc}

\begin{exc}
  Design and implement a variant of \texttt{ArrayDeque} that does not do any
  modular arithmetic at all.  Instead, all the data sits in a consecutive
  block, in order, inside an array.  When the data overruns the beginning
  or the end of this array, a modified \ensuremath{\mathtt{rebuild()}} operation is performed.
  The amortized cost of all operations should be the same as in an
  \texttt{ArrayDeque}.

  \noindent Hint: Getting this to work is really all about how you implement
  the \ensuremath{\mathtt{rebuild()}} operation.  You would like \ensuremath{\mathtt{rebuild()}} to put the data
  structure into a state where the data cannot run off either end until
  at least $\ensuremath{\mathtt{{\color{var}n}}}/2$ operations have been performed.

  Test the performance of your implementation against the \texttt{ArrayDeque}.
  Optimize your implementation (by using \ensuremath{\mathtt{System.arraycopy({\color{var}a},{\color{var}i},{\color{var}b},{\color{var}i},{\color{var}n})}})
  and see if you can get it to outperform the \texttt{ArrayDeque} implementation.
\end{exc}

\begin{exc}
  Design and implement a version of a \texttt{RootishArrayStack} that has
  only $O(\sqrt{\ensuremath{\mathtt{{\color{var}n}}}})$ wasted space, but that can perform \ensuremath{\mathtt{add({\color{var}i},{\color{var}x})}}
  and \ensuremath{\mathtt{remove({\color{var}i},{\color{var}x})}} operations in $O(1+\min\{\ensuremath{\mathtt{{\color{var}i}}},\ensuremath{\mathtt{{\color{var}n}}}-\ensuremath{\mathtt{{\color{var}i}}}\})$ time.
\end{exc}

\begin{exc}\exclabel{rootisharraystack-fast}
  Design and implement a version of a \texttt{RootishArrayStack} that has
  only $O(\sqrt{\ensuremath{\mathtt{{\color{var}n}}}})$ wasted space, but that can perform \ensuremath{\mathtt{add({\color{var}i},{\color{var}x})}}
  and \ensuremath{\mathtt{remove({\color{var}i},{\color{var}x})}} operations in $O(1+\min\{\sqrt{\ensuremath{\mathtt{{\color{var}n}}}},\ensuremath{\mathtt{{\color{var}n}}}-\ensuremath{\mathtt{{\color{var}i}}}\})$
  time. (For an idea on how to do this, see \secref{selist}.)
\end{exc}

\begin{exc}
  Design and implement a version of a \texttt{RootishArrayStack} that has
  only $O(\sqrt{\ensuremath{\mathtt{{\color{var}n}}}})$ wasted space, but that can perform \ensuremath{\mathtt{add({\color{var}i},{\color{var}x})}} and
  \ensuremath{\mathtt{remove({\color{var}i},{\color{var}x})}} operations in $O(1+\min\{\ensuremath{\mathtt{{\color{var}i}}},\sqrt {\ensuremath{\mathtt{{\color{var}n}}}},\ensuremath{\mathtt{{\color{var}n}}}-\ensuremath{\mathtt{{\color{var}i}}}\})$ time.
  (See \secref{selist} for ideas on how to achieve this.)
\end{exc}

\begin{exc}
  Design and implement a \texttt{CubishArrayStack}.
  \index{CubishArrayStack@\texttt{CubishArrayStack}}%
  This three level structure
  implements the \texttt{List} interface using $O(\ensuremath{\mathtt{{\color{var}n}}}^{2/3})$ wasted space.
  In this structure, \ensuremath{\mathtt{get({\color{var}i})}} and \ensuremath{\mathtt{set({\color{var}i},{\color{var}x})}} take constant time; while
  \ensuremath{\mathtt{add({\color{var}i},{\color{var}x})}} and \ensuremath{\mathtt{remove({\color{var}i})}} take $O(\ensuremath{\mathtt{{\color{var}n}}}^{1/3})$ amortized time.
\end{exc}



\clearpage

\clearpage
\chapter{Linked Lists}
\chaplabel{linkedlists}

\index{linked list}%
In this chapter, we continue to study implementations of the \texttt{List}
interface, this time using pointer-based data structures rather than
arrays.  The structures in this chapter are made up of nodes that
contain the list items.  Using references (pointers), the nodes are
linked together into a sequence.  We first study singly-linked lists,
which can implement \texttt{Stack} and (FIFO) \texttt{Queue} operations in constant
time per operation and then move on to doubly-linked lists, which can
implement \texttt{Deque} operations in constant time.

Linked lists have advantages and disadvantages when compared to array-based
implementations of the \texttt{List} interface.  The primary disadvantage is that
we lose the ability to access any element using \ensuremath{\mathtt{get({\color{var}i})}} or \ensuremath{\mathtt{set({\color{var}i},{\color{var}x})}}
in constant time.  Instead, we have to walk through the list, one element
at a time, until we reach the \ensuremath{\mathtt{{\color{var}i}}}th element.  The primary advantage is
that they are more dynamic:  Given a reference to any list node \ensuremath{\mathtt{{\color{var}u}}}, we
can delete \ensuremath{\mathtt{{\color{var}u}}} or insert a node adjacent to \ensuremath{\mathtt{{\color{var}u}}} in constant time. This
is true no matter where \ensuremath{\mathtt{{\color{var}u}}} is in the list.


\section{\texttt{SLList}: A Singly-Linked List}
\seclabel{sllist}

\index{SLList@\texttt{SLList}}%
\index{linked list!singly-}%
\index{singly-linked list}%
An \texttt{SLList} (singly-linked list) is a sequence of \texttt{Node}s.  Each node
\ensuremath{\mathtt{{\color{var}u}}} stores a data value \ensuremath{\mathtt{{\color{var}u}.{\color{var}x}}} and a reference \ensuremath{\mathtt{{\color{var}u}.{\color{var}next}}} to the next node in
the sequence.  For the last node \ensuremath{\mathtt{{\color{var}w}}} in the sequence, $\ensuremath{\mathtt{{\color{var}w}.{\color{var}next}}} = \ensuremath{\mathtt{{\color{keyword}null}}}$



 \ \vspace{.5ex}\begin{Verbatim}[gobble=1,tabsize=2,frame=single,framerule=1.2pt,commandchars=\\@\$,label=\texttt{SLList},labelposition=topline]
	@\color@keyword$class$ Node {
		@\color@keyword$T$ @\color@var$x$;
		Node @\color@var$next$;
	}
\end{Verbatim}


For efficiency, an \texttt{SLList} uses variables \ensuremath{\mathtt{{\color{var}head}}} and \ensuremath{\mathtt{{\color{var}tail}}} to keep
track of the first and last node in the sequence, as well as an integer
\ensuremath{\mathtt{{\color{var}n}}} to keep track of the length of the sequence:


 \ \vspace{.5ex}\begin{Verbatim}[gobble=1,tabsize=2,frame=single,framerule=1.2pt,commandchars=\\@\$,label=\texttt{SLList},labelposition=topline]
	Node @\color@var$head$;
	Node @\color@var$tail$;
	@\color@keyword$int$ @\color@var$n$;
\end{Verbatim}

A sequence of \texttt{Stack} and \texttt{Queue} operations on an \texttt{SLList} is
illustrated in \figref{sllist}.

\begin{figure}
  \begin{center}
    \includegraphics[width=\ScaleIfNeeded]{figs/sllist}
  \end{center}
  \caption[A sequence of Queue and Stack operations on an SLList]{A sequence of \texttt{Queue} (\ensuremath{\mathtt{add({\color{var}x})}} and \ensuremath{\mathtt{remove()}}) and \texttt{Stack} (\ensuremath{\mathtt{push({\color{var}x})}} and \ensuremath{\mathtt{pop()}}) operations on an \texttt{SLList}.}
  \figlabel{sllist}
\end{figure}


An \texttt{SLList} can efficiently implement the \texttt{Stack} operations \ensuremath{\mathtt{push()}}
and \ensuremath{\mathtt{pop()}} by adding and removing elements at the head of the sequence.
The \ensuremath{\mathtt{push()}} operation simply creates a new node \ensuremath{\mathtt{{\color{var}u}}} with data value \ensuremath{\mathtt{{\color{var}x}}},
sets \ensuremath{\mathtt{{\color{var}u}.{\color{var}next}}} to the old head of the list and makes \ensuremath{\mathtt{{\color{var}u}}} the new head
of the list. Finally, it increments \ensuremath{\mathtt{{\color{var}n}}} since the size of the \texttt{SLList}
has increased by one:



 \ \vspace{.5ex}\begin{Verbatim}[gobble=1,tabsize=2,frame=single,framerule=1.2pt,commandchars=\\@\$,label=\texttt{SLList},labelposition=topline]
	@\color@keyword$T$ push(@\color@keyword$T$ @\color@var$x$) {
		Node @\color@var$u$ = @\color@keyword$new$ Node();
		@\color@var$u$.@\color@var$x$ = @\color@var$x$;
		@\color@var$u$.@\color@var$next$ = @\color@var$head$;
		@\color@var$head$ = @\color@var$u$;
		@\color@keyword$if$ (@\color@var$n$ == 0)
			@\color@var$tail$ = @\color@var$u$;
		@\color@var$n$++;
		@\color@var$return$ @\color@var$x$;
	}
\end{Verbatim}


The \ensuremath{\mathtt{pop()}} operation, after checking that the \texttt{SLList} is not empty,
removes the head by setting $\ensuremath{\mathtt{{\color{var}head}={\color{var}head}.{\color{var}next}}}$ and decrementing \ensuremath{\mathtt{{\color{var}n}}}.
A special case occurs when the last element is being removed, in which case \ensuremath{\mathtt{{\color{var}tail}}} is set to \ensuremath{\mathtt{{\color{keyword}null}}}:



 \ \vspace{.5ex}\begin{Verbatim}[gobble=1,tabsize=2,frame=single,framerule=1.2pt,commandchars=\\@\$,label=\texttt{SLList},labelposition=topline]
	@\color@keyword$T$ pop() {
		@\color@keyword$if$ (@\color@var$n$ == 0)	@\color@var$return$ @\color@keyword$null$;
		@\color@keyword$T$ @\color@var$x$ = @\color@var$head$.@\color@var$x$;
		@\color@var$head$ = @\color@var$head$.@\color@var$next$;
		@\color@keyword$if$ (--@\color@var$n$ == 0) @\color@var$tail$ = @\color@keyword$null$;
		@\color@var$return$ @\color@var$x$;
	}	
\end{Verbatim}


Clearly, both the \ensuremath{\mathtt{push({\color{var}x})}} and \ensuremath{\mathtt{pop()}} operations run in $O(1)$ time.

\subsection{Queue Operations}

An \texttt{SLList} can also implement the FIFO queue operations \ensuremath{\mathtt{add({\color{var}x})}} and
\ensuremath{\mathtt{remove()}} in constant time.  Removals are done from the head of the list,
and are identical to the \ensuremath{\mathtt{pop()}} operation:



 \ \vspace{.5ex}\begin{Verbatim}[gobble=1,tabsize=2,frame=single,framerule=1.2pt,commandchars=\\@\$,label=\texttt{SLList},labelposition=topline]
	@\color@keyword$T$ remove() {
		@\color@keyword$if$ (@\color@var$n$ == 0)	@\color@var$return$ @\color@keyword$null$;
		@\color@keyword$T$ @\color@var$x$ = @\color@var$head$.@\color@var$x$;
		@\color@var$head$ = @\color@var$head$.@\color@var$next$;
		@\color@keyword$if$ (--@\color@var$n$ == 0) @\color@var$tail$ = @\color@keyword$null$;
		@\color@var$return$ @\color@var$x$;
	}	
\end{Verbatim}


Additions, on the other hand, are done at the tail of the list.  In most
cases, this is done by setting $\ensuremath{\mathtt{{\color{var}tail}.{\color{var}next}}}=\ensuremath{\mathtt{{\color{var}u}}}$, where \ensuremath{\mathtt{{\color{var}u}}} is the newly
created node that contains \ensuremath{\mathtt{{\color{var}x}}}.  However, a special case occurs when
$\ensuremath{\mathtt{{\color{var}n}}}=0$, in which case $\ensuremath{\mathtt{{\color{var}tail}}}=\ensuremath{\mathtt{{\color{var}head}}}=\ensuremath{\mathtt{{\color{keyword}null}}}$.  In this case, both \ensuremath{\mathtt{{\color{var}tail}}}
and \ensuremath{\mathtt{{\color{var}head}}} are set to \ensuremath{\mathtt{{\color{var}u}}}.



 \ \vspace{.5ex}\begin{Verbatim}[gobble=1,tabsize=2,frame=single,framerule=1.2pt,commandchars=\\@\$,label=\texttt{SLList},labelposition=topline]
	@\color@var$boolean$ add(@\color@keyword$T$ @\color@var$x$) {
		Node @\color@var$u$ = @\color@keyword$new$ Node();
		@\color@var$u$.@\color@var$x$ = @\color@var$x$;
		@\color@keyword$if$ (@\color@var$n$ == 0) {
			@\color@var$head$ = @\color@var$u$;
		} @\color@keyword$else$ {
			@\color@var$tail$.@\color@var$next$ = @\color@var$u$;
		}
		@\color@var$tail$ = @\color@var$u$;
		@\color@var$n$++;
		@\color@var$return$ @\color@var$true$;
	}
\end{Verbatim}


Clearly, both \ensuremath{\mathtt{add({\color{var}x})}} and \ensuremath{\mathtt{remove()}} take constant time.

\subsection{Summary}

The following theorem summarizes the performance of an \texttt{SLList}:

\begin{thm}\thmlabel{sllist}
  An \texttt{SLList} implements the \texttt{Stack} and (FIFO) \texttt{Queue} interfaces.  
  The \ensuremath{\mathtt{push({\color{var}x})}}, \ensuremath{\mathtt{pop()}}, \ensuremath{\mathtt{add({\color{var}x})}} and \ensuremath{\mathtt{remove()}} operations run
  in $O(1)$ time per operation.
\end{thm}

An \texttt{SLList} nearly implements the full set of \texttt{Deque} operations.
The only missing operation is removing from the tail of an \texttt{SLList}.
Removing from the tail of an \texttt{SLList} is difficult because it requires
updating the value of \ensuremath{\mathtt{{\color{var}tail}}} so that it points to the node \ensuremath{\mathtt{{\color{var}w}}}
that precedes \ensuremath{\mathtt{{\color{var}tail}}} in the \texttt{SLList}; this is the node \ensuremath{\mathtt{{\color{var}w}}} such that
$\ensuremath{\mathtt{{\color{var}w}.{\color{var}next}}}=\ensuremath{\mathtt{{\color{var}tail}}}$.  Unfortunately, the only way to get to \ensuremath{\mathtt{{\color{var}w}}} is by
traversing the \texttt{SLList} starting at \ensuremath{\mathtt{{\color{var}head}}} and taking $\ensuremath{\mathtt{{\color{var}n}}}-2$ steps.

\section{\texttt{DLList}: A Doubly-Linked List}
\seclabel{dllist}

\index{DLList@\texttt{DLList}}%
\index{doubly-linked list}%
\index{linked list!doubly-}%
A \texttt{DLList} (doubly-linked list) is very similar to an \texttt{SLList} except
that each node \ensuremath{\mathtt{{\color{var}u}}} in a \texttt{DLList} has references to both the node \ensuremath{\mathtt{{\color{var}u}.{\color{var}next}}}
that follows it and the node \ensuremath{\mathtt{{\color{var}u}.{\color{var}prev}}} that precedes it.



 \ \vspace{.5ex}\begin{Verbatim}[gobble=1,tabsize=2,frame=single,framerule=1.2pt,commandchars=\\@\$,label=\texttt{DLList},labelposition=topline]
	@\color@keyword$class$ Node {
		@\color@keyword$T$ @\color@var$x$;
		Node @\color@var$prev$, @\color@var$next$;
	}
\end{Verbatim}


When implementing an \texttt{SLList}, we saw that there were always several
special cases to worry about. For example, removing the last element
from an \texttt{SLList} or adding an element to an empty \texttt{SLList} requires care
to ensure that \ensuremath{\mathtt{{\color{var}head}}} and \ensuremath{\mathtt{{\color{var}tail}}} are correctly updated.  In a \texttt{DLList},
the number of these special cases increases considerably.  Perhaps the
cleanest way to take care of all these special cases in a \texttt{DLList} is to
introduce a \ensuremath{\mathtt{{\color{var}dummy}}} node.
\index{dummy node}%
This is a node that does not contain any data,
but acts as a placeholder so that there are no special nodes; every node
has both a \ensuremath{\mathtt{{\color{var}next}}} and a \ensuremath{\mathtt{{\color{var}prev}}}, with \ensuremath{\mathtt{{\color{var}dummy}}} acting as the node that
follows the last node in the list and that precedes the first node in
the list.  In this way, the nodes of the list are (doubly-)linked into
a cycle, as illustrated in \figref{dllist}.

\begin{figure}
  \begin{center}
    \includegraphics[width=\ScaleIfNeeded]{figs/dllist2}
  \end{center}
  \caption[A DLList]{A \texttt{DLList} containing a,b,c,d,e.}
  \figlabel{dllist}
\end{figure}





 \ \vspace{.5ex}\begin{Verbatim}[gobble=1,tabsize=2,frame=single,framerule=1.2pt,commandchars=\\@\$,label=\texttt{DLList},labelposition=topline]
	@\color@keyword$int$ @\color@var$n$;
	Node @\color@var$dummy$;
	DLList() {
		@\color@var$dummy$ = @\color@keyword$new$ Node();
		@\color@var$dummy$.@\color@var$next$ = @\color@var$dummy$;
		@\color@var$dummy$.@\color@var$prev$ = @\color@var$dummy$;
		@\color@var$n$ = 0;
	}
\end{Verbatim}


Finding the node with a particular index in a \texttt{DLList} is easy;  we can
either start at the head of the list (\ensuremath{\mathtt{{\color{var}dummy}.{\color{var}next}}}) and work forward,
or start at the tail of the list (\ensuremath{\mathtt{{\color{var}dummy}.{\color{var}prev}}}) and work backward.
This allows us to reach the \ensuremath{\mathtt{{\color{var}i}}}th node in $O(1+\min\{\ensuremath{\mathtt{{\color{var}i}}},\ensuremath{\mathtt{{\color{var}n}}}-\ensuremath{\mathtt{{\color{var}i}}}\})$ time:



 \ \vspace{.5ex}\begin{Verbatim}[gobble=1,tabsize=2,frame=single,framerule=1.2pt,commandchars=\\@\$,label=\texttt{DLList},labelposition=topline]
	Node getNode(@\color@keyword$int$ @\color@var$i$) {
		Node @\color@var$p$ = @\color@keyword$null$;
		@\color@keyword$if$ (@\color@var$i$ < @\color@var$n$ / 2) {
			@\color@var$p$ = @\color@var$dummy$.@\color@var$next$;
			@\color@keyword$for$ (@\color@keyword$int$ @\color@var$j$ = 0; @\color@var$j$ < @\color@var$i$; @\color@var$j$++)
				@\color@var$p$ = @\color@var$p$.@\color@var$next$;
		} @\color@keyword$else$ {
			@\color@var$p$ = @\color@var$dummy$;
			@\color@keyword$for$ (@\color@keyword$int$ @\color@var$j$ = @\color@var$n$; @\color@var$j$ > @\color@var$i$; @\color@var$j$--)
				@\color@var$p$ = @\color@var$p$.@\color@var$prev$;
		}
		@\color@var$return$ @\color@var$p$;
	}
\end{Verbatim}


The \ensuremath{\mathtt{get({\color{var}i})}} and \ensuremath{\mathtt{set({\color{var}i},{\color{var}x})}} operations are now also easy.  We first find the \ensuremath{\mathtt{{\color{var}i}}}th node and then get or set its \ensuremath{\mathtt{{\color{var}x}}} value:



 \ \vspace{.5ex}\begin{Verbatim}[gobble=1,tabsize=2,frame=single,framerule=1.2pt,commandchars=\\@\$,label=\texttt{DLList},labelposition=topline]
	@\color@keyword$T$ get(@\color@keyword$int$ @\color@var$i$) {
		@\color@var$return$ getNode(@\color@var$i$).@\color@var$x$;
	}
	@\color@keyword$T$ set(@\color@keyword$int$ @\color@var$i$, @\color@keyword$T$ @\color@var$x$) {
		Node @\color@var$u$ = getNode(@\color@var$i$);
		@\color@keyword$T$ @\color@var$y$ = @\color@var$u$.@\color@var$x$;
		@\color@var$u$.@\color@var$x$ = @\color@var$x$;
		@\color@var$return$ @\color@var$y$;
	}
\end{Verbatim}


The running time of these operations is dominated by the time it takes
to find the \ensuremath{\mathtt{{\color{var}i}}}th node, and is therefore $O(1+\min\{\ensuremath{\mathtt{{\color{var}i}}},\ensuremath{\mathtt{{\color{var}n}}}-\ensuremath{\mathtt{{\color{var}i}}}\})$.

\subsection{Adding and Removing}

If we have a reference to a node \ensuremath{\mathtt{{\color{var}w}}} in a \texttt{DLList} and we want to insert a
node \ensuremath{\mathtt{{\color{var}u}}} before \ensuremath{\mathtt{{\color{var}w}}}, then this is just a matter of setting $\ensuremath{\mathtt{{\color{var}u}.{\color{var}next}}}=\ensuremath{\mathtt{{\color{var}w}}}$,
$\ensuremath{\mathtt{{\color{var}u}.{\color{var}prev}}}=\ensuremath{\mathtt{{\color{var}w}.{\color{var}prev}}}$, and then adjusting \ensuremath{\mathtt{{\color{var}u}.{\color{var}prev}.{\color{var}next}}} and \ensuremath{\mathtt{{\color{var}u}.{\color{var}next}.{\color{var}prev}}}.  (See \figref{dllist-addbefore}.)
Thanks to the dummy node, there is no need to worry about \ensuremath{\mathtt{{\color{var}w}.{\color{var}prev}}}
or \ensuremath{\mathtt{{\color{var}w}.{\color{var}next}}} not existing.



 \ \vspace{.5ex}\begin{Verbatim}[gobble=1,tabsize=2,frame=single,framerule=1.2pt,commandchars=\\@\$,label=\texttt{DLList},labelposition=topline]
	Node addBefore(Node @\color@var$w$, @\color@keyword$T$ @\color@var$x$) {
		Node @\color@var$u$ = @\color@keyword$new$ Node();
		@\color@var$u$.@\color@var$x$ = @\color@var$x$;
		@\color@var$u$.@\color@var$prev$ = @\color@var$w$.@\color@var$prev$;
		@\color@var$u$.@\color@var$next$ = @\color@var$w$;
		@\color@var$u$.@\color@var$next$.@\color@var$prev$ = @\color@var$u$;
		@\color@var$u$.@\color@var$prev$.@\color@var$next$ = @\color@var$u$;
		@\color@var$n$++;
		@\color@var$return$ @\color@var$u$;
	}
\end{Verbatim}


\begin{figure}
   \begin{center}
      \includegraphics[scale=0.90909]{figs/dllist-addbefore}
   \end{center}
   \caption[Adding to a DLList]{Adding the node \ensuremath{\mathtt{{\color{var}u}}} before the node \ensuremath{\mathtt{{\color{var}w}}}
      in a \texttt{DLList}.}
   \figlabel{dllist-addbefore}
\end{figure}

Now, the list operation \ensuremath{\mathtt{add({\color{var}i},{\color{var}x})}} is trivial to implement.  We find the
\ensuremath{\mathtt{{\color{var}i}}}th node in the \texttt{DLList} and insert a new node \ensuremath{\mathtt{{\color{var}u}}} that contains \ensuremath{\mathtt{{\color{var}x}}}
just before it.



 \ \vspace{.5ex}\begin{Verbatim}[gobble=1,tabsize=2,frame=single,framerule=1.2pt,commandchars=\\@\$,label=\texttt{DLList},labelposition=topline]
	@\color@keyword$void$ add(@\color@keyword$int$ @\color@var$i$, @\color@keyword$T$ @\color@var$x$) {
		addBefore(getNode(@\color@var$i$), @\color@var$x$);
	}
\end{Verbatim}


The only non-constant part of the running time of \ensuremath{\mathtt{add({\color{var}i},{\color{var}x})}} is the time
it takes to find the \ensuremath{\mathtt{{\color{var}i}}}th node (using \ensuremath{\mathtt{getNode({\color{var}i})}}).  Thus, \ensuremath{\mathtt{add({\color{var}i},{\color{var}x})}}
runs in $O(1+\min\{\ensuremath{\mathtt{{\color{var}i}}}, \ensuremath{\mathtt{{\color{var}n}}}-\ensuremath{\mathtt{{\color{var}i}}}\})$ time.

Removing a node \ensuremath{\mathtt{{\color{var}w}}} from a \texttt{DLList} is easy.  We only need to adjust
pointers at \ensuremath{\mathtt{{\color{var}w}.{\color{var}next}}} and \ensuremath{\mathtt{{\color{var}w}.{\color{var}prev}}} so that they skip over \ensuremath{\mathtt{{\color{var}w}}}.  Again, the
use of the dummy node eliminates the need to consider any special cases:



 \ \vspace{.5ex}\begin{Verbatim}[gobble=1,tabsize=2,frame=single,framerule=1.2pt,commandchars=\\@\$,label=\texttt{DLList},labelposition=topline]
	@\color@keyword$void$ remove(Node @\color@var$w$) {
		@\color@var$w$.@\color@var$prev$.@\color@var$next$ = @\color@var$w$.@\color@var$next$;
		@\color@var$w$.@\color@var$next$.@\color@var$prev$ = @\color@var$w$.@\color@var$prev$;
		@\color@var$n$--;
	}
\end{Verbatim}


Now the \ensuremath{\mathtt{remove({\color{var}i})}} operation is trivial. We find the node with index \ensuremath{\mathtt{{\color{var}i}}} and remove it:



 \ \vspace{.5ex}\begin{Verbatim}[gobble=1,tabsize=2,frame=single,framerule=1.2pt,commandchars=\\@\$,label=\texttt{DLList},labelposition=topline]
	@\color@keyword$T$ remove(@\color@keyword$int$ @\color@var$i$) {
		Node @\color@var$w$ = getNode(@\color@var$i$);
		remove(@\color@var$w$);
		@\color@var$return$ @\color@var$w$.@\color@var$x$;
	}
\end{Verbatim}


Again, the only expensive part of this operation is finding the \ensuremath{\mathtt{{\color{var}i}}}th node
using \ensuremath{\mathtt{getNode({\color{var}i})}}, so \ensuremath{\mathtt{remove({\color{var}i})}} runs in $O(1+\min\{\ensuremath{\mathtt{{\color{var}i}}}, \ensuremath{\mathtt{{\color{var}n}}}-\ensuremath{\mathtt{{\color{var}i}}}\})$
time.

\subsection{Summary}

The following theorem summarizes the performance of a \texttt{DLList}:

\begin{thm}\thmlabel{dllist}
  A \texttt{DLList} implements the \texttt{List} interface.  In this implementation,
  the \ensuremath{\mathtt{get({\color{var}i})}}, \ensuremath{\mathtt{set({\color{var}i},{\color{var}x})}}, \ensuremath{\mathtt{add({\color{var}i},{\color{var}x})}} and \ensuremath{\mathtt{remove({\color{var}i})}} operations run
  in $O(1+\min\{\ensuremath{\mathtt{{\color{var}i}}},\ensuremath{\mathtt{{\color{var}n}}}-\ensuremath{\mathtt{{\color{var}i}}}\})$ time per operation.
\end{thm}

It is worth noting that, if we ignore the cost of the \ensuremath{\mathtt{getNode({\color{var}i})}}
operation, then all operations on a \texttt{DLList} take constant time.
Thus, the only expensive part of operations on a \texttt{DLList} is finding
the relevant node.  Once we have the relevant node, adding, removing,
or accessing the data at that node takes only constant time.

This is in sharp contrast to the array-based \texttt{List} implementations of
\chapref{arrays}; in those implementations, the relevant array
item can be found in constant time. However, addition or removal requires
shifting elements in the array and, in general, takes non-constant time.

For this reason, linked list structures are well-suited to applications
where references to list nodes can be obtained through external means.
\javaonly{An example of this is the \texttt{LinkedHashSet} data structure found in the
Java Collections Framework, in which a set of items is stored in a
doubly-linked list and the nodes of the doubly-linked list are stored
in a hash table (discussed in \chapref{hashing}).  When elements
are removed from a \texttt{LinkedHashSet}, the hash table is used to find the
relevant list node in constant time and then the list node is deleted
(also in constant time).}
\cpponly{For example, pointers to the nodes of a linked list could be
stored in a \texttt{USet}.  Then, to remove an item \ensuremath{\mathtt{{\color{var}x}}} from the linked list,
the node that contains \ensuremath{\mathtt{{\color{var}x}}} can be found quickly using the \texttt{Uset} and
the node can be removed from the list in constant time.}

\section{\texttt{SEList}: A Space-Efficient Linked List}
\seclabel{selist}

\index{SEList@\texttt{SEList}}%
\index{linked list!space-efficient}%
One of the drawbacks of linked lists (besides the time it takes to access
elements that are deep within the list) is their space usage.  Each node
in a \texttt{DLList} requires an additional two references to the next and
previous nodes in the list.  Two of the fields in a \texttt{Node} are dedicated
to maintaining the list, and only one of the fields is for storing data!

An \texttt{SEList} (space-efficient list) reduces this wasted space using
a simple idea: Rather than store individual elements in a \texttt{DLList},
we store a block (array) containing several items. More precisely, an
\texttt{SEList} is parameterized by a \emph{block size} \ensuremath{\mathtt{{\color{var}b}}}. Each individual
node in an \texttt{SEList} stores a block that can hold up to \ensuremath{\mathtt{{\color{var}b}+1}} elements.

For reasons that will become clear later, it will be helpful if we can
do \texttt{Deque} operations on each block.  The data structure that we choose
for this is a \texttt{BDeque} (bounded deque),
\index{BDeque@\texttt{BDeque}}%
\index{bounded deque}%
\index{deque!bounded}%
derived from the \texttt{ArrayDeque}
structure described in \secref{arraydeque}.  The \texttt{BDeque} differs from the
\texttt{ArrayDeque} in one small way: When a new \texttt{BDeque} is created, the size
of the backing array \ensuremath{\mathtt{{\color{var}a}}} is fixed at \ensuremath{\mathtt{{\color{var}b}+1}} and never grows or shrinks.
The important property of a \texttt{BDeque} is that it allows for the addition
or removal of elements at either the front or back in constant time. This
will be useful as elements are shifted from one block to another.



 \ \vspace{.5ex}\begin{Verbatim}[gobble=1,tabsize=2,frame=single,framerule=1.2pt,commandchars=\\@\$,label=\texttt{SEList},labelposition=topline]
	@\color@keyword$class$ BDeque @\color@keyword$extends$ ArrayDeque<@\color@keyword$T$> {
		BDeque() {
			super(SEList.@\color@var$this$.type());
			@\color@var$a$ = newArray(@\color@var$b$+1);
		}
		@\color@keyword$void$ resize() { }
	}
\end{Verbatim}


\notpcode{
An \texttt{SEList} is then a doubly-linked list of blocks:
} % notpcode


 \ \vspace{.5ex}\begin{Verbatim}[gobble=1,tabsize=2,frame=single,framerule=1.2pt,commandchars=\\@\$,label=\texttt{SEList},labelposition=topline]
	@\color@keyword$class$ Node {
		BDeque @\color@var$d$;
		Node @\color@var$prev$, @\color@var$next$;
	}
\end{Verbatim}



 \ \vspace{.5ex}\begin{Verbatim}[gobble=1,tabsize=2,frame=single,framerule=1.2pt,commandchars=\\@\$,label=\texttt{SEList},labelposition=topline]
	@\color@keyword$int$ @\color@var$n$;
	Node @\color@var$dummy$;
\end{Verbatim}


\pcodeonly{
An \texttt{SEList} is just a doubly-linked list of blocks. In addition to \ensuremath{\mathtt{{\color{var}next}}} and \ensuremath{\mathtt{{\color{var}prev}}} pointers, each node \ensuremath{\mathtt{{\color{var}u}}} in an \texttt{SEList} contains a \texttt{BDeque}, \ensuremath{\mathtt{{\color{var}u}.{\color{var}d}}}.
}

\subsection{Space Requirements}

An \texttt{SEList} places very tight restrictions on the number of elements
in a block: Unless a block is the last block, then that block contains
at least $\ensuremath{\mathtt{{\color{var}b}}}-1$ and at most $\ensuremath{\mathtt{{\color{var}b}}}+1$ elements.  This means that, if an
\texttt{SEList} contains \ensuremath{\mathtt{{\color{var}n}}} elements, then it has at most
\[
    \ensuremath{\mathtt{{\color{var}n}}}/(\ensuremath{\mathtt{{\color{var}b}}}-1) + 1 = O(\ensuremath{\mathtt{{\color{var}n}}}/\ensuremath{\mathtt{{\color{var}b}}})
\]
blocks.  The \texttt{BDeque} for each block contains an array of length $\ensuremath{\mathtt{{\color{var}b}}}+1$
but, for every block except the last, at most a constant amount of
space is wasted in this array.  The remaining memory used by a block is
also constant.  This means that the wasted space in an \texttt{SEList} is only
$O(\ensuremath{\mathtt{{\color{var}b}}}+\ensuremath{\mathtt{{\color{var}n}}}/\ensuremath{\mathtt{{\color{var}b}}})$.  By choosing a value of \ensuremath{\mathtt{{\color{var}b}}} within a constant factor
of $\sqrt{\ensuremath{\mathtt{{\color{var}n}}}}$, we can make the space-overhead of an SEList approach
the $\sqrt{\ensuremath{\mathtt{{\color{var}n}}}}$ lower bound given in \secref{rootishspaceusage}.

\subsection{Finding Elements}

The first challenge we face with an \texttt{SEList} is finding the list item
with a given index \ensuremath{\mathtt{{\color{var}i}}}.  Note that the location of an element consists
of two parts: 
\begin{enumerate}
  \item The node \ensuremath{\mathtt{{\color{var}u}}} that contains the block that contains the element
  with index \ensuremath{\mathtt{{\color{var}i}}}; and
  \item the index \ensuremath{\mathtt{{\color{var}j}}} of the element within its block.
\end{enumerate}



 \ \vspace{.5ex}\begin{Verbatim}[gobble=1,tabsize=2,frame=single,framerule=1.2pt,commandchars=\\@\$,label=\texttt{SEList},labelposition=topline]
	@\color@keyword$class$ Location {
		Node @\color@var$u$;
		@\color@keyword$int$ @\color@var$j$;
		Location(Node @\color@var$u$, @\color@keyword$int$ @\color@var$j$) {
			@\color@var$this$.@\color@var$u$ = @\color@var$u$;
			@\color@var$this$.@\color@var$j$ = @\color@var$j$;
		}
	}
\end{Verbatim}


To find the block that contains a particular element, we proceed the same
way as we do in a \texttt{DLList}.  We either start at the front of the list and
traverse in the forward direction, or at the back of the list and traverse
backwards until we reach the node we want.  The only difference is that,
each time we move from one node to the next, we skip over a whole block
of elements.

\pcodeimport{ods/SEList.get_location(i)}


 \ \vspace{.5ex}\begin{Verbatim}[gobble=1,tabsize=2,frame=single,framerule=1.2pt,commandchars=\\@\$,label=\texttt{SEList},labelposition=topline]
	Location getLocation(@\color@keyword$int$ @\color@var$i$) {
		@\color@keyword$if$ (@\color@var$i$ < @\color@var$n$/2) {
			Node @\color@var$u$ = @\color@var$dummy$.@\color@var$next$;
			@\color@keyword$while$ (@\color@var$i$ >= @\color@var$u$.@\color@var$d$.size()) {
				@\color@var$i$ -= @\color@var$u$.@\color@var$d$.size();
				@\color@var$u$ = @\color@var$u$.@\color@var$next$;
			}
			@\color@var$return$ @\color@keyword$new$ Location(@\color@var$u$, @\color@var$i$);
		} @\color@keyword$else$ {
			Node @\color@var$u$ = @\color@var$dummy$;
			@\color@keyword$int$ @\color@var$idx$ = @\color@var$n$;
			@\color@keyword$while$ (@\color@var$i$ < @\color@var$idx$) {
				@\color@var$u$ = @\color@var$u$.@\color@var$prev$;
				@\color@var$idx$ -= @\color@var$u$.@\color@var$d$.size();
			}
			@\color@var$return$ @\color@keyword$new$ Location(@\color@var$u$, @\color@var$i$-@\color@var$idx$);
		}
	}
\end{Verbatim}


Remember that, with the exception of at most one block, each block
contains at least $\ensuremath{\mathtt{{\color{var}b}}}-1$ elements, so each step in our search gets
us $\ensuremath{\mathtt{{\color{var}b}}}-1$ elements closer to the element we are looking for.  If we
are searching forward, this means that we reach the node we want after
$O(1+\ensuremath{\mathtt{{\color{var}i}}}/\ensuremath{\mathtt{{\color{var}b}}})$ steps.  If we search backwards, then we reach the node we
want after $O(1+(\ensuremath{\mathtt{{\color{var}n}}}-\ensuremath{\mathtt{{\color{var}i}}})/\ensuremath{\mathtt{{\color{var}b}}})$ steps.  The algorithm takes the smaller
of these two quantities depending on the value of \ensuremath{\mathtt{{\color{var}i}}}, so the time to
locate the item with index \ensuremath{\mathtt{{\color{var}i}}} is $O(1+\min\{\ensuremath{\mathtt{{\color{var}i}}},\ensuremath{\mathtt{{\color{var}n}}}-\ensuremath{\mathtt{{\color{var}i}}}\}/\ensuremath{\mathtt{{\color{var}b}}})$.

Once we know how to locate the item with index \ensuremath{\mathtt{{\color{var}i}}}, the \ensuremath{\mathtt{get({\color{var}i})}} and
\ensuremath{\mathtt{set({\color{var}i},{\color{var}x})}} operations translate into getting or setting a particular
index in the correct block:



 \ \vspace{.5ex}\begin{Verbatim}[gobble=1,tabsize=2,frame=single,framerule=1.2pt,commandchars=\\@\$,label=\texttt{SEList},labelposition=topline]
	@\color@keyword$T$ get(@\color@keyword$int$ @\color@var$i$) {
		Location @\color@var$l$ = getLocation(@\color@var$i$);
		@\color@var$return$ @\color@var$l$.@\color@var$u$.@\color@var$d$.get(@\color@var$l$.@\color@var$j$);
	}
	@\color@keyword$T$ set(@\color@keyword$int$ @\color@var$i$, @\color@keyword$T$ @\color@var$x$) {
		Location @\color@var$l$ = getLocation(@\color@var$i$);
		@\color@keyword$T$ @\color@var$y$ = @\color@var$l$.@\color@var$u$.@\color@var$d$.get(@\color@var$l$.@\color@var$j$);
		@\color@var$l$.@\color@var$u$.@\color@var$d$.set(@\color@var$l$.@\color@var$j$,@\color@var$x$);
		@\color@var$return$ @\color@var$y$;
	}
\end{Verbatim}


The running times of these operations are dominated by the time it takes
to locate the item, so they also run in $O(1+\min\{\ensuremath{\mathtt{{\color{var}i}}},\ensuremath{\mathtt{{\color{var}n}}}-\ensuremath{\mathtt{{\color{var}i}}}\}/\ensuremath{\mathtt{{\color{var}b}}})$
time.

\subsection{Adding an Element}

Adding elements to an \texttt{SEList} is a little more complicated.  Before
considering the general case, we consider the easier operation, \ensuremath{\mathtt{add({\color{var}x})}},
in which \ensuremath{\mathtt{{\color{var}x}}} is added to the end of the list.  If the last block is full
(or does not exist because there are no blocks yet), then we first
allocate a new block and append it to the list of blocks.  Now that
we are sure that the last block exists and is not full, we append \ensuremath{\mathtt{{\color{var}x}}}
to the last block.



 \ \vspace{.5ex}\begin{Verbatim}[gobble=1,tabsize=2,frame=single,framerule=1.2pt,commandchars=\\@\$,label=\texttt{SEList},labelposition=topline]
	@\color@var$boolean$ add(@\color@keyword$T$ @\color@var$x$) {
		Node @\color@var$last$ = @\color@var$dummy$.@\color@var$prev$;
		@\color@keyword$if$ (@\color@var$last$ == @\color@var$dummy$ || @\color@var$last$.@\color@var$d$.size() == @\color@var$b$+1) {
			@\color@var$last$ = addBefore(@\color@var$dummy$);
		}
		@\color@var$last$.@\color@var$d$.add(@\color@var$x$);
		@\color@var$n$++;
		@\color@var$return$ @\color@var$true$;
	}
\end{Verbatim}

\pcodeimport{ods/SEList.append(x)}

Things get more complicated when we add to the interior of the list
using \ensuremath{\mathtt{add({\color{var}i},{\color{var}x})}}.  We first locate \ensuremath{\mathtt{{\color{var}i}}} to get the node \ensuremath{\mathtt{{\color{var}u}}} whose block
contains the \ensuremath{\mathtt{{\color{var}i}}}th list item.  The problem is that we want to insert
\ensuremath{\mathtt{{\color{var}x}}} into \ensuremath{\mathtt{{\color{var}u}}}'s block, but we have to be prepared for the case where
\ensuremath{\mathtt{{\color{var}u}}}'s block already contains $\ensuremath{\mathtt{{\color{var}b}}}+1$ elements, so that it is full and
there is no room for \ensuremath{\mathtt{{\color{var}x}}}.

Let $\ensuremath{\mathtt{{\color{var}u}}}_0,\ensuremath{\mathtt{{\color{var}u}}}_1,\ensuremath{\mathtt{{\color{var}u}}}_2,\ldots$ denote \ensuremath{\mathtt{{\color{var}u}}}, \ensuremath{\mathtt{{\color{var}u}.{\color{var}next}}}, \ensuremath{\mathtt{{\color{var}u}.{\color{var}next}.{\color{var}next}}},
and so on.  We explore $\ensuremath{\mathtt{{\color{var}u}}}_0,\ensuremath{\mathtt{{\color{var}u}}}_1,\ensuremath{\mathtt{{\color{var}u}}}_2,\ldots$ looking for a node
that can provide space for \ensuremath{\mathtt{{\color{var}x}}}.  Three cases can occur during our
space exploration (see \figref{selist-add}):

\begin{figure}
  \noindent
  \begin{center}
    \begin{tabular}{@{}l@{}}
      \includegraphics[width=\ScaleIfNeeded]{figs/selist-add-a}\\[4ex]
      \includegraphics[width=\ScaleIfNeeded]{figs/selist-add-b}\\[4ex]
      \includegraphics[width=\ScaleIfNeeded]{figs/selist-add-c}\\
    \end{tabular}
  \end{center}
  \caption[SEList add]{The three cases that occur during the addition of an item \ensuremath{\mathtt{{\color{var}x}}} in the interior of an \texttt{SEList}.  (This \texttt{SEList} has block size $\ensuremath{\mathtt{{\color{var}b}}}=3$.)}
  \figlabel{selist-add}
\end{figure}


\begin{enumerate}
\item We quickly (in $r+1\le \ensuremath{\mathtt{{\color{var}b}}}$ steps) find a node $\ensuremath{\mathtt{{\color{var}u}}}_r$ whose block
is not full.  In this case, we perform $r$ shifts of an element from
one block into the next, so that the free space in $\ensuremath{\mathtt{{\color{var}u}}}_r$ becomes a
free space in $\ensuremath{\mathtt{{\color{var}u}}}_0$.  We can then insert \ensuremath{\mathtt{{\color{var}x}}} into $\ensuremath{\mathtt{{\color{var}u}}}_0$'s block.

\item We quickly (in $r+1\le \ensuremath{\mathtt{{\color{var}b}}}$ steps) run off the end of the list
of blocks.  In this case, we add a new empty block to the end of the
list of blocks and proceed as in the first case.

\item After \ensuremath{\mathtt{{\color{var}b}}} steps we do not find any block that is not full.
In this case, $\ensuremath{\mathtt{{\color{var}u}}}_0,\ldots,\ensuremath{\mathtt{{\color{var}u}}}_{\ensuremath{\mathtt{{\color{var}b}}}-1}$ is a sequence of \ensuremath{\mathtt{{\color{var}b}}} blocks
that each contain $\ensuremath{\mathtt{{\color{var}b}}}+1$ elements.  We insert a new block $\ensuremath{\mathtt{{\color{var}u}}}_{\ensuremath{\mathtt{{\color{var}b}}}}$
at the end of this sequence and \emph{spread} the original $\ensuremath{\mathtt{{\color{var}b}}}(\ensuremath{\mathtt{{\color{var}b}}}+1)$
elements so that each block of $\ensuremath{\mathtt{{\color{var}u}}}_0,\ldots,\ensuremath{\mathtt{{\color{var}u}}}_{\ensuremath{\mathtt{{\color{var}b}}}}$ contains exactly
\ensuremath{\mathtt{{\color{var}b}}} elements.  Now $\ensuremath{\mathtt{{\color{var}u}}}_0$'s block contains only \ensuremath{\mathtt{{\color{var}b}}} elements so it has
room for us to insert \ensuremath{\mathtt{{\color{var}x}}}.
\end{enumerate}



 \ \vspace{.5ex}\begin{Verbatim}[gobble=1,tabsize=2,frame=single,framerule=1.2pt,commandchars=\\@\$,label=\texttt{SEList},labelposition=topline]
	@\color@keyword$void$ add(@\color@keyword$int$ @\color@var$i$, @\color@keyword$T$ @\color@var$x$) {
		@\color@keyword$if$ (@\color@var$i$ == @\color@var$n$) {
			add(@\color@var$x$);
			@\color@var$return$;
		}
		Location @\color@var$l$ = getLocation(@\color@var$i$);
		Node @\color@var$u$ = @\color@var$l$.@\color@var$u$;
		@\color@keyword$int$ @\color@var$r$ = 0;
		@\color@keyword$while$ (@\color@var$r$ < @\color@var$b$ && @\color@var$u$ != @\color@var$dummy$ && @\color@var$u$.@\color@var$d$.size() == @\color@var$b$+1) {
			@\color@var$u$ = @\color@var$u$.@\color@var$next$;
			@\color@var$r$++;
		}
		@\color@keyword$if$ (@\color@var$r$ == @\color@var$b$) {      @\color@comment$// b blocks each with b+1 elements$
			spread(@\color@var$l$.@\color@var$u$);
			@\color@var$u$ = @\color@var$l$.@\color@var$u$;
		} 
		@\color@keyword$if$ (@\color@var$u$ == @\color@var$dummy$) {  @\color@comment$// ran off the end - add new node$
			@\color@var$u$ = addBefore(@\color@var$u$);
		}
		@\color@keyword$while$ (@\color@var$u$ != @\color@var$l$.@\color@var$u$) { @\color@comment$// work backwards, shifting elements$
			@\color@var$u$.@\color@var$d$.add(0, @\color@var$u$.@\color@var$prev$.@\color@var$d$.remove(@\color@var$u$.@\color@var$prev$.@\color@var$d$.size()-1));
			@\color@var$u$ = @\color@var$u$.@\color@var$prev$;
		}
		@\color@var$u$.@\color@var$d$.add(@\color@var$l$.@\color@var$j$, @\color@var$x$);
		@\color@var$n$++;
	}
\end{Verbatim}


The running time of the \ensuremath{\mathtt{add({\color{var}i},{\color{var}x})}} operation depends on which of
the three cases above occurs.  Cases~1 and 2 involve examining and
shifting elements through at most \ensuremath{\mathtt{{\color{var}b}}} blocks and take $O(\ensuremath{\mathtt{{\color{var}b}}})$ time.
Case~3 involves calling the \ensuremath{\mathtt{spread({\color{var}u})}} method, which  moves $\ensuremath{\mathtt{{\color{var}b}}}(\ensuremath{\mathtt{{\color{var}b}}}+1)$
elements and takes $O(\ensuremath{\mathtt{{\color{var}b}}}^2)$ time.  If we ignore the cost of Case~3
(which we will account for later with amortization) this means that
the total running time to locate \ensuremath{\mathtt{{\color{var}i}}} and perform the insertion of \ensuremath{\mathtt{{\color{var}x}}}
is $O(\ensuremath{\mathtt{{\color{var}b}}}+\min\{\ensuremath{\mathtt{{\color{var}i}}},\ensuremath{\mathtt{{\color{var}n}}}-\ensuremath{\mathtt{{\color{var}i}}}\}/\ensuremath{\mathtt{{\color{var}b}}})$.

\subsection{Removing an Element}

Removing an element from an \texttt{SEList} is similar to adding an element.
We first locate the node \ensuremath{\mathtt{{\color{var}u}}} that contains the element with index \ensuremath{\mathtt{{\color{var}i}}}.
Now, we have to be prepared for the case where we cannot remove an element
from \ensuremath{\mathtt{{\color{var}u}}} without causing \ensuremath{\mathtt{{\color{var}u}}}'s block to become smaller than $\ensuremath{\mathtt{{\color{var}b}}}-1$.

Again, let $\ensuremath{\mathtt{{\color{var}u}}}_0,\ensuremath{\mathtt{{\color{var}u}}}_1,\ensuremath{\mathtt{{\color{var}u}}}_2,\ldots$ denote \ensuremath{\mathtt{{\color{var}u}}}, \ensuremath{\mathtt{{\color{var}u}.{\color{var}next}}}, \ensuremath{\mathtt{{\color{var}u}.{\color{var}next}.{\color{var}next}}},
and so on.  We examine $\ensuremath{\mathtt{{\color{var}u}}}_0,\ensuremath{\mathtt{{\color{var}u}}}_1,\ensuremath{\mathtt{{\color{var}u}}}_2,\ldots$ in order to look
for a node from which we can borrow an element to make the size of
$\ensuremath{\mathtt{{\color{var}u}}}_0$'s block at least $\ensuremath{\mathtt{{\color{var}b}}}-1$.  There are three cases to consider
(see \figref{selist-remove}):

\begin{figure}
  \noindent
  \begin{center}
    \begin{tabular}{l}
      \includegraphics[scale=0.90909]{figs/selist-remove-a}\\[4ex]
      \includegraphics[scale=0.90909]{figs/selist-remove-b}\\[4ex]
      \includegraphics[scale=0.90909]{figs/selist-remove-c}\\
    \end{tabular}
  \end{center}
  \caption[SEList remove]{The three cases that occur during the removal of an item \ensuremath{\mathtt{{\color{var}x}}} in the interior of an \texttt{SEList}.  (This \texttt{SEList} has block size $\ensuremath{\mathtt{{\color{var}b}}}=3$.)}
  \figlabel{selist-remove}
\end{figure}


\begin{enumerate}
\item We quickly (in $r+1\le \ensuremath{\mathtt{{\color{var}b}}}$ steps) find a node whose block contains
more than $\ensuremath{\mathtt{{\color{var}b}}}-1$ elements. In this case, we perform $r$ shifts of an
element from one block into the previous one, so that the extra element
in $\ensuremath{\mathtt{{\color{var}u}}}_r$ becomes an extra element in $\ensuremath{\mathtt{{\color{var}u}}}_0$.  We can then remove the
appropriate element from $\ensuremath{\mathtt{{\color{var}u}}}_0$'s block.

\item We quickly (in $r+1\le \ensuremath{\mathtt{{\color{var}b}}}$ steps) run off the end of the list of
blocks.  In this case, $\ensuremath{\mathtt{{\color{var}u}}}_r$ is the last block, and there is no need
for $\ensuremath{\mathtt{{\color{var}u}}}_r$'s block to contain at least $\ensuremath{\mathtt{{\color{var}b}}}-1$ elements.  Therefore,
we proceed as above, borrowing an element from $\ensuremath{\mathtt{{\color{var}u}}}_r$ to make an extra
element in $\ensuremath{\mathtt{{\color{var}u}}}_0$.  If this causes $\ensuremath{\mathtt{{\color{var}u}}}_r$'s block to become empty,
then we remove it.

\item After \ensuremath{\mathtt{{\color{var}b}}} steps, we do not find any block containing more than
$\ensuremath{\mathtt{{\color{var}b}}}-1$ elements.  In this case, $\ensuremath{\mathtt{{\color{var}u}}}_0,\ldots,\ensuremath{\mathtt{{\color{var}u}}}_{\ensuremath{\mathtt{{\color{var}b}}}-1}$ is a sequence
of \ensuremath{\mathtt{{\color{var}b}}} blocks that each contain $\ensuremath{\mathtt{{\color{var}b}}}-1$ elements.  We \emph{gather}
these $\ensuremath{\mathtt{{\color{var}b}}}(\ensuremath{\mathtt{{\color{var}b}}}-1)$ elements into $\ensuremath{\mathtt{{\color{var}u}}}_0,\ldots,\ensuremath{\mathtt{{\color{var}u}}}_{\ensuremath{\mathtt{{\color{var}b}}}-2}$ so that each
of these $\ensuremath{\mathtt{{\color{var}b}}}-1$ blocks contains exactly \ensuremath{\mathtt{{\color{var}b}}} elements and we remove
$\ensuremath{\mathtt{{\color{var}u}}}_{\ensuremath{\mathtt{{\color{var}b}}}-1}$, which is now empty.  Now $\ensuremath{\mathtt{{\color{var}u}}}_0$'s block contains \ensuremath{\mathtt{{\color{var}b}}}
elements and we can then remove the appropriate element from it.
\end{enumerate}



 \ \vspace{.5ex}\begin{Verbatim}[gobble=1,tabsize=2,frame=single,framerule=1.2pt,commandchars=\\@\$,label=\texttt{SEList},labelposition=topline]
	@\color@keyword$T$ remove(@\color@keyword$int$ @\color@var$i$) {
		Location @\color@var$l$ = getLocation(@\color@var$i$);
		@\color@keyword$T$ @\color@var$y$ = @\color@var$l$.@\color@var$u$.@\color@var$d$.get(@\color@var$l$.@\color@var$j$);
		Node @\color@var$u$ = @\color@var$l$.@\color@var$u$;
		@\color@keyword$int$ @\color@var$r$ = 0;
		@\color@keyword$while$ (@\color@var$r$ < @\color@var$b$ && @\color@var$u$ != @\color@var$dummy$ && @\color@var$u$.@\color@var$d$.size() == @\color@var$b$-1) {
			@\color@var$u$ = @\color@var$u$.@\color@var$next$;
			@\color@var$r$++;
		}
		@\color@keyword$if$ (@\color@var$r$ == @\color@var$b$) {  @\color@comment$// b blocks each with b-1 elements$
			gather(@\color@var$l$.@\color@var$u$);
		}
		@\color@var$u$ = @\color@var$l$.@\color@var$u$;
		@\color@var$u$.@\color@var$d$.remove(@\color@var$l$.@\color@var$j$);
		@\color@keyword$while$ (@\color@var$u$.@\color@var$d$.size() < @\color@var$b$-1 && @\color@var$u$.@\color@var$next$ != @\color@var$dummy$) {
			@\color@var$u$.@\color@var$d$.add(@\color@var$u$.@\color@var$next$.@\color@var$d$.remove(0));
			@\color@var$u$ = @\color@var$u$.@\color@var$next$;
		}
		@\color@keyword$if$ (@\color@var$u$.@\color@var$d$.isEmpty()) remove(@\color@var$u$);
		@\color@var$n$--;
		@\color@var$return$ @\color@var$y$;
	}
\end{Verbatim}


Like the \ensuremath{\mathtt{add({\color{var}i},{\color{var}x})}} operation, the running time of the \ensuremath{\mathtt{remove({\color{var}i})}}
operation is $O(\ensuremath{\mathtt{{\color{var}b}}}+\min\{\ensuremath{\mathtt{{\color{var}i}}},\ensuremath{\mathtt{{\color{var}n}}}-\ensuremath{\mathtt{{\color{var}i}}}\}/\ensuremath{\mathtt{{\color{var}b}}})$ if we ignore the cost of
the \ensuremath{\mathtt{gather({\color{var}u})}} method that occurs in Case~3.

\subsection{Amortized Analysis of Spreading and Gathering}

Next, we consider the cost of the \ensuremath{\mathtt{gather({\color{var}u})}} and \ensuremath{\mathtt{spread({\color{var}u})}} methods
that may be executed by the \ensuremath{\mathtt{add({\color{var}i},{\color{var}x})}} and \ensuremath{\mathtt{remove({\color{var}i})}} methods.  For the
sake of completeness, here they are:



 \ \vspace{.5ex}\begin{Verbatim}[gobble=1,tabsize=2,frame=single,framerule=1.2pt,commandchars=\\@\$,label=\texttt{SEList},labelposition=topline]
	@\color@keyword$void$ spread(Node @\color@var$u$) {
		Node @\color@var$w$ = @\color@var$u$;
		@\color@keyword$for$ (@\color@keyword$int$ @\color@var$j$ = 0; @\color@var$j$ < @\color@var$b$; @\color@var$j$++) {
			@\color@var$w$ = @\color@var$w$.@\color@var$next$;
		}
		@\color@var$w$ = addBefore(@\color@var$w$);
		@\color@keyword$while$ (@\color@var$w$ != @\color@var$u$) {
			@\color@keyword$while$ (@\color@var$w$.@\color@var$d$.size() < @\color@var$b$)
				@\color@var$w$.@\color@var$d$.add(0,@\color@var$w$.@\color@var$prev$.@\color@var$d$.remove(@\color@var$w$.@\color@var$prev$.@\color@var$d$.size()-1));
			@\color@var$w$ = @\color@var$w$.@\color@var$prev$;
		}
	}
\end{Verbatim}



 \ \vspace{.5ex}\begin{Verbatim}[gobble=1,tabsize=2,frame=single,framerule=1.2pt,commandchars=\\@\$,label=\texttt{SEList},labelposition=topline]
	@\color@keyword$void$ gather(Node @\color@var$u$) {
		Node @\color@var$w$ = @\color@var$u$;
		@\color@keyword$for$ (@\color@keyword$int$ @\color@var$j$ = 0; @\color@var$j$ < @\color@var$b$-1; @\color@var$j$++) {
			@\color@keyword$while$ (@\color@var$w$.@\color@var$d$.size() < @\color@var$b$)
				@\color@var$w$.@\color@var$d$.add(@\color@var$w$.@\color@var$next$.@\color@var$d$.remove(0));
			@\color@var$w$ = @\color@var$w$.@\color@var$next$;
		}
		remove(@\color@var$w$);
	}
\end{Verbatim}


The running time of each of these methods is dominated by the two
nested loops.  Both the inner and outer loops execute at most
$\ensuremath{\mathtt{{\color{var}b}}}+1$ times, so the total running time of each of these methods
is $O((\ensuremath{\mathtt{{\color{var}b}}}+1)^2)=O(\ensuremath{\mathtt{{\color{var}b}}}^2)$. However, the following lemma shows that
these methods execute on at most one out of every \ensuremath{\mathtt{{\color{var}b}}} calls to \ensuremath{\mathtt{add({\color{var}i},{\color{var}x})}}
or \ensuremath{\mathtt{remove({\color{var}i})}}.

\begin{lem}\lemlabel{selist-amortized}
  If an empty \texttt{SEList} is created and any sequence of $m\ge 1$ calls
  to \ensuremath{\mathtt{add({\color{var}i},{\color{var}x})}} and \ensuremath{\mathtt{remove({\color{var}i})}} is performed, then the total time
  spent during all calls to \ensuremath{\mathtt{spread()}} and \ensuremath{\mathtt{gather()}} is $O(\ensuremath{\mathtt{{\color{var}b}}}m)$.
\end{lem}

\begin{proof}
  We will use the potential method of amortized analysis.
  \index{potential method}%
  We say that
  a node \ensuremath{\mathtt{{\color{var}u}}} is \emph{fragile} if \ensuremath{\mathtt{{\color{var}u}}}'s block does not contain \ensuremath{\mathtt{{\color{var}b}}}
  elements (so that \ensuremath{\mathtt{{\color{var}u}}} is either the last node, or contains $\ensuremath{\mathtt{{\color{var}b}}}-1$
  or $\ensuremath{\mathtt{{\color{var}b}}}+1$ elements).  Any node whose block contains \ensuremath{\mathtt{{\color{var}b}}} elements is
  \emph{rugged}. Define the \emph{potential} of an \texttt{SEList} as the number
  of fragile nodes it contains.  We will consider only the \ensuremath{\mathtt{add({\color{var}i},{\color{var}x})}}
  operation and its relation to the number of calls to \ensuremath{\mathtt{spread({\color{var}u})}}.
  The analysis of \ensuremath{\mathtt{remove({\color{var}i})}} and \ensuremath{\mathtt{gather({\color{var}u})}} is identical.

  Notice that, if Case~1 occurs during the \ensuremath{\mathtt{add({\color{var}i},{\color{var}x})}} method, then
  only one node, $\ensuremath{\mathtt{{\color{var}u}}}_r$ has the size of its block changed. Therefore,
  at most one node, namely $\ensuremath{\mathtt{{\color{var}u}}}_r$, goes from being rugged to being
  fragile.  If Case~2 occurs, then a new node is created, and this node
  is fragile, but no other node changes size, so the number of fragile
  nodes increases by one.  Thus, in either Case~1 or Case~2 the potential
  of the SEList increases by at most one.

  Finally, if Case~3 occurs, it is because $\ensuremath{\mathtt{{\color{var}u}}}_0,\ldots,\ensuremath{\mathtt{{\color{var}u}}}_{\ensuremath{\mathtt{{\color{var}b}}}-1}$
  are all fragile nodes.  Then $\ensuremath{\mathtt{spread(}}u_0\ensuremath{\mathtt{)}}$ is called and these \ensuremath{\mathtt{{\color{var}b}}}
  fragile nodes are replaced with $\ensuremath{\mathtt{{\color{var}b}}}+1$ rugged nodes.  Finally, \ensuremath{\mathtt{{\color{var}x}}}
  is added to $\ensuremath{\mathtt{{\color{var}u}}}_0$'s block, making $\ensuremath{\mathtt{{\color{var}u}}}_0$ fragile.  In total the
  potential decreases by $\ensuremath{\mathtt{{\color{var}b}}}-1$.

  In summary, the potential starts at 0 (there are no nodes in the list).
  Each time Case~1 or Case~2 occurs, the potential increases by at
  most 1.  Each time Case~3 occurs, the potential decreases by $\ensuremath{\mathtt{{\color{var}b}}}-1$.
  The potential (which counts the number of fragile nodes) is never
  less than 0.  We conclude that, for every occurrence of Case~3, there
  are at least $\ensuremath{\mathtt{{\color{var}b}}}-1$ occurrences of Case~1 or Case~2.  Thus, for every
  call to \ensuremath{\mathtt{spread({\color{var}u})}} there are at least \ensuremath{\mathtt{{\color{var}b}}} calls to \ensuremath{\mathtt{add({\color{var}i},{\color{var}x})}}.  This
  completes the proof.
\end{proof}

\subsection{Summary}

The following theorem summarizes the performance of the \texttt{SEList} data
structure:

\begin{thm}\thmlabel{selist}
  An \texttt{SEList} implements the \texttt{List} interface.  Ignoring the cost of
  calls to \ensuremath{\mathtt{spread({\color{var}u})}} and \ensuremath{\mathtt{gather({\color{var}u})}}, an \texttt{SEList} with block size \ensuremath{\mathtt{{\color{var}b}}}
  supports the operations
  \begin{itemize}
    \item \ensuremath{\mathtt{get({\color{var}i})}} and \ensuremath{\mathtt{set({\color{var}i},{\color{var}x})}} in $O(1+\min\{\ensuremath{\mathtt{{\color{var}i}}},\ensuremath{\mathtt{{\color{var}n}}}-\ensuremath{\mathtt{{\color{var}i}}}\}/\ensuremath{\mathtt{{\color{var}b}}})$ time per operation; and
    \item \ensuremath{\mathtt{add({\color{var}i},{\color{var}x})}} and \ensuremath{\mathtt{remove({\color{var}i})}} in $O(\ensuremath{\mathtt{{\color{var}b}}}+\min\{\ensuremath{\mathtt{{\color{var}i}}},\ensuremath{\mathtt{{\color{var}n}}}-\ensuremath{\mathtt{{\color{var}i}}}\}/\ensuremath{\mathtt{{\color{var}b}}})$ time per operation.
  \end{itemize}
  Furthermore, beginning with an empty \texttt{SEList}, any sequence of $m$
  \ensuremath{\mathtt{add({\color{var}i},{\color{var}x})}} and \ensuremath{\mathtt{remove({\color{var}i})}} operations results in a total of $O(\ensuremath{\mathtt{{\color{var}b}}}m)$
  time spent during all calls to \ensuremath{\mathtt{spread({\color{var}u})}} and \ensuremath{\mathtt{gather({\color{var}u})}}.

  The space (measured in words)\footnote{Recall \secref{model} for a
  discussion of how memory is measured.} used by an \texttt{SEList}
  that stores \ensuremath{\mathtt{{\color{var}n}}} elements is $\ensuremath{\mathtt{{\color{var}n}}} +O(\ensuremath{\mathtt{{\color{var}b}}} + \ensuremath{\mathtt{{\color{var}n}}}/\ensuremath{\mathtt{{\color{var}b}}})$.
\end{thm}

The \texttt{SEList} is a trade-off between an \texttt{ArrayList} and a \texttt{DLList} where
the relative mix of these two structures depends on the block size \ensuremath{\mathtt{{\color{var}b}}}.
At the extreme $\ensuremath{\mathtt{{\color{var}b}}}=2$, each \texttt{SEList} node stores at most three values,
which is not much different than a \texttt{DLList}. At the other extreme,
$\ensuremath{\mathtt{{\color{var}b}}}>\ensuremath{\mathtt{{\color{var}n}}}$, all the elements are stored in a single array, just like in
an \texttt{ArrayList}.  In between these two extremes lies a trade-off between
the time it takes to add or remove a list item and the time it takes to
locate a particular list item.

\section{Discussion and Exercises}

Both singly-linked and doubly-linked lists are established techniques,
having been used in programs for over 40 years.  They are discussed,
for example, by Knuth \cite[Sections~2.2.3--2.2.5]{k97v1}.  Even the
\texttt{SEList} data structure seems to be a well-known data structures exercise.
The \texttt{SEList} is sometimes referred to as an \emph{unrolled linked list}
\cite{sra94}.
\index{unrolled linked list|seealso{\texttt{SEList}}}%
\index{linked list!unrolled|seealso{\texttt{SEList}}}%

Another way to save space in a doubly-linked list is to use 
so-called XOR-lists.
\index{XOR-list}%
In an XOR-list, each node, \ensuremath{\mathtt{{\color{var}u}}}, contains only one
pointer, called \ensuremath{\mathtt{{\color{var}u}.{\color{var}nextprev}}}, that holds the bitwise exclusive-or of \ensuremath{\mathtt{{\color{var}u}.{\color{var}prev}}}
and \ensuremath{\mathtt{{\color{var}u}.{\color{var}next}}}.  The list itself needs to store two pointers, one to the \ensuremath{\mathtt{{\color{var}dummy}}}
node and one to \ensuremath{\mathtt{{\color{var}dummy}.{\color{var}next}}} (the first node, or \ensuremath{\mathtt{{\color{var}dummy}}} if the list is
empty). This technique uses the fact that, if we have pointers to \ensuremath{\mathtt{{\color{var}u}}}
and \ensuremath{\mathtt{{\color{var}u}.{\color{var}prev}}}, then we can extract \ensuremath{\mathtt{{\color{var}u}.{\color{var}next}}} using the formula
\[
   \ensuremath{\mathtt{{\color{var}u}.{\color{var}next}}} = \ensuremath{\mathtt{{\color{var}u}.{\color{var}prev}}} \verb+^+ \ensuremath{\mathtt{{\color{var}u}.{\color{var}nextprev}}} \enspace .
\]
(Here \verb+^+ computes the bitwise exclusive-or of its two arguments.)
This technique complicates the code a little and is not possible in
some languages, like Java and Python, that have garbage collection but gives a
doubly-linked list implementation that requires only one pointer per node.
See Sinha's magazine article \cite{s04} for a detailed discussion of
XOR-lists.

\begin{exc}
  Why is it not possible to use a dummy node in an \texttt{SLList} to avoid
  all the special cases that occur in the operations \ensuremath{\mathtt{push({\color{var}x})}}, \ensuremath{\mathtt{pop()}},
  \ensuremath{\mathtt{add({\color{var}x})}}, and \ensuremath{\mathtt{remove()}}?
\end{exc}

\begin{exc}
  Design and implement an \texttt{SLList} method, \ensuremath{\mathtt{secondLast()}}, that returns
  the second-last element of an \texttt{SLList}.  Do this without using the
  member variable, \ensuremath{\mathtt{{\color{var}n}}}, that keeps track of the size of the list.
\end{exc}

\begin{exc}
  Implement the \texttt{List} operations \ensuremath{\mathtt{get({\color{var}i})}}, \ensuremath{\mathtt{set({\color{var}i},{\color{var}x})}},
  \ensuremath{\mathtt{add({\color{var}i},{\color{var}x})}} and \ensuremath{\mathtt{remove({\color{var}i})}} on an \texttt{SLList}.  Each of these operations
  should run in $O(1+\ensuremath{\mathtt{{\color{var}i}}})$ time.
\end{exc}

\begin{exc}
  Design and implement an \texttt{SLList} method, \ensuremath{\mathtt{reverse()}} that reverses the
  order of elements in an \texttt{SLList}.  This method should run in $O(\ensuremath{\mathtt{{\color{var}n}}})$
  time, should not use recursion, should not use any secondary data
  structures, and should not create any new nodes.
\end{exc}

\begin{exc}
  Design and implement \texttt{SLList} and \texttt{DLList} methods called \ensuremath{\mathtt{checkSize()}}.
  These methods walk through the list and count the number of nodes to
  see if this matches the value, \ensuremath{\mathtt{{\color{var}n}}}, stored in the list.  These methods
  return nothing, but throw an exception if the size they compute does
  not match the value of \ensuremath{\mathtt{{\color{var}n}}}.
\end{exc}

\begin{exc}
  Try to recreate the code for the \ensuremath{\mathtt{addBefore({\color{var}w})}} operation that creates a
  node, \ensuremath{\mathtt{{\color{var}u}}}, and adds it in a \texttt{DLList} just before the node \ensuremath{\mathtt{{\color{var}w}}}.  Do not
  refer to this chapter.  Even if your code does not exactly match the
  code given in this book it may still be correct.  Test it and see if
  it works.
\end{exc}

The next few exercises involve performing manipulations on \texttt{DLList}s.
You should complete them without allocating any new nodes or temporary
arrays.  They can all be done only by changing the \ensuremath{\mathtt{{\color{var}prev}}} and \ensuremath{\mathtt{{\color{var}next}}}
values of existing nodes.

\begin{exc}
  Write a \texttt{DLList} method \ensuremath{\mathtt{isPalindrome()}} that returns \ensuremath{\mathtt{{\color{var}true}}} if the
  list is a \emph{palindrome},
  \index{palindrome}%
  i.e., the element at position \ensuremath{\mathtt{{\color{var}i}}} is equal to
  the element at position $\ensuremath{\mathtt{{\color{var}n}}}-i-1$ for all $i\in\{0,\ldots,\ensuremath{\mathtt{{\color{var}n}}}-1\}$.
  Your code should run in $O(\ensuremath{\mathtt{{\color{var}n}}})$ time.
\end{exc}

\begin{exc}
  Implement a method \ensuremath{\mathtt{rotate({\color{var}r})}} that ``rotates'' a \texttt{DLList} so that list
  item \ensuremath{\mathtt{{\color{var}i}}} becomes list item $(\ensuremath{\mathtt{{\color{var}i}}}+\ensuremath{\mathtt{{\color{var}r}}})\bmod \ensuremath{\mathtt{{\color{var}n}}}$.  This method should
  run in $O(1+\min\{\ensuremath{\mathtt{{\color{var}r}}},\ensuremath{\mathtt{{\color{var}n}}}-\ensuremath{\mathtt{{\color{var}r}}}\})$ time and should not modify any nodes in
  the list.
\end{exc}


\begin{exc}\exclabel{linkedlist-truncate}
  Write a method, \ensuremath{\mathtt{truncate({\color{var}i})}}, that truncates a \texttt{DLList} at position
  \ensuremath{\mathtt{{\color{var}i}}}.  After executing this method, the size of the list will be \ensuremath{\mathtt{{\color{var}i}}} and
  it should contain only the elements at indices $0,\ldots,\ensuremath{\mathtt{{\color{var}i}}}-1$.  The
  return value is another \texttt{DLList} that contains the elements at indices
  $\ensuremath{\mathtt{{\color{var}i}}},\ldots,\ensuremath{\mathtt{{\color{var}n}}}-1$.  This method should run in $O(\min\{\ensuremath{\mathtt{{\color{var}i}}},\ensuremath{\mathtt{{\color{var}n}}}-\ensuremath{\mathtt{{\color{var}i}}}\})$
  time.
\end{exc}

\begin{exc}
  Write a \texttt{DLList} method, \ensuremath{\mathtt{absorb({\color{var}l2})}}, that takes as an argument
  a \texttt{DLList}, \ensuremath{\mathtt{{\color{var}l2}}}, empties it and appends its contents, in order,
  to the receiver.  For example, if \ensuremath{\mathtt{{\color{var}l1}}} contains $a,b,c$ and \ensuremath{\mathtt{{\color{var}l2}}}
  contains $d,e,f$, then after calling \ensuremath{\mathtt{{\color{var}l1}.absorb({\color{var}l2})}}, \ensuremath{\mathtt{{\color{var}l1}}} will contain
  $a,b,c,d,e,f$ and \ensuremath{\mathtt{{\color{var}l2}}} will be empty.
\end{exc}

\begin{exc}
  Write a method \ensuremath{\mathtt{deal()}} that removes all the elements with odd-numbered
  indices from a \texttt{DLList} and return a \texttt{DLList} containing these elements.
  For example, if \ensuremath{\mathtt{{\color{var}l1}}}, contains the elements $a,b,c,d,e,f$, then after
  calling \ensuremath{\mathtt{{\color{var}l1}.deal()}}, \ensuremath{\mathtt{{\color{var}l1}}} should contain $a,c,e$ and a list containing
  $b,d,f$ should be returned.
\end{exc}

\begin{exc}
  Write a method, \ensuremath{\mathtt{reverse()}}, that reverses the order of elements in
  a \texttt{DLList}.  
\end{exc}

\begin{exc}\exclabel{dllist-sort}
  \index{merge-sort}%
  This exercise walks you through an implementation of the merge-sort
  algorithm for sorting a \texttt{DLList}, as discussed in \secref{merge-sort}.
  \javaonly{In your implementation, perform comparisons between elements
  using the \ensuremath{\mathtt{compareTo({\color{var}x})}} method so that the resulting implementation can
  sort any \texttt{DLList} containing elements that implement the \texttt{Comparable}
  interface.}
  \begin{enumerate}
    \item Write a \texttt{DLList} method called \ensuremath{\mathtt{takeFirst({\color{var}l2})}}.
       This method takes the first node from \ensuremath{\mathtt{{\color{var}l2}}} and appends it to the the
       receiving list.  This is equivalent to \ensuremath{\mathtt{add(size(),{\color{var}l2}.remove(0))}},
       except that it should not create a new node.
    \item Write a \texttt{DLList} static method, \ensuremath{\mathtt{merge({\color{var}l1},{\color{var}l2})}}, that takes two
       sorted lists \ensuremath{\mathtt{{\color{var}l1}}} and \ensuremath{\mathtt{{\color{var}l2}}}, merges them, and returns a new sorted
       list containing the result.  This causes \ensuremath{\mathtt{{\color{var}l1}}} and \ensuremath{\mathtt{{\color{var}l2}}} to be emptied
       in the proces.  For example, if \ensuremath{\mathtt{{\color{var}l1}}} contains $a,c,d$ and \ensuremath{\mathtt{{\color{var}l2}}} contains
       $b,e,f$, then this method returns a new list containing $a,b,c,d,e,f$.
    \item Write a \texttt{DLList} method \ensuremath{\mathtt{sort()}} that sorts the elements
       contained in the list using the merge sort algorithm.
       This recursive algorithm works in the following way:
       \begin{enumerate}
          \item If the list contains 0 or 1 elements then there is
            nothing to do.  Otherwise,
          \item Using the \ensuremath{\mathtt{truncate(size()/2)}} method, split the list
            into two lists of approximately equal length, \ensuremath{\mathtt{{\color{var}l1}}} and \ensuremath{\mathtt{{\color{var}l2}}};
          \item Recursively sort \ensuremath{\mathtt{{\color{var}l1}}};
          \item Recursively sort \ensuremath{\mathtt{{\color{var}l2}}}; and, finally,
          \item Merge \ensuremath{\mathtt{{\color{var}l1}}} and \ensuremath{\mathtt{{\color{var}l2}}} into a single sorted list.
       \end{enumerate}
  \end{enumerate}
\end{exc}


The next few exercises are more advanced and require a clear
understanding of what happens to the minimum value stored in a \texttt{Stack}
or \texttt{Queue} as items are added and removed.

\begin{exc}
  \index{MinStack@\texttt{MinStack}}%
  Design and implement a \texttt{MinStack} data structure that can store
  comparable elements and supports the stack operations \ensuremath{\mathtt{push({\color{var}x})}},
  \ensuremath{\mathtt{pop()}}, and \ensuremath{\mathtt{size()}}, as well as the \ensuremath{\mathtt{min()}} operation, which
  returns the minimum value currently stored in the data structure.
  All operations should run in constant time.
\end{exc}

\begin{exc}
  \index{MinQueue@\texttt{MinQueue}}%
  Design and implement a \texttt{MinQueue} data structure that can store
  comparable elements and supports the queue operations \ensuremath{\mathtt{add({\color{var}x})}},
  \ensuremath{\mathtt{remove()}}, and \ensuremath{\mathtt{size()}}, as well as the \ensuremath{\mathtt{min()}} operation, which
  returns the minimum value currently stored in the data structure.
  All operations should run in constant amortized time.
\end{exc}

\begin{exc}
  \index{MinDeque@\texttt{MinDeque}}%
  Design and implement a \texttt{MinDeque} data structure that can store
  comparable elements and supports all the deque operations \ensuremath{\mathtt{addFirst({\color{var}x})}},
  \ensuremath{\mathtt{addLast({\color{var}x})}} \ensuremath{\mathtt{removeFirst()}}, \ensuremath{\mathtt{removeLast()}} and \ensuremath{\mathtt{size()}}, and the
  \ensuremath{\mathtt{min()}} operation, which returns the minimum value currently stored in
  the data structure.  All operations should run in constant amortized
  time.
\end{exc}

The next exercises are designed to test the reader's understanding of
the implementation and analysis of the space-efficient \texttt{SEList}:

\begin{exc}
  Prove that, if an \texttt{SEList} is used like a \texttt{Stack} (so that the
  only modifications to the \texttt{SEList} are done using $\ensuremath{\mathtt{push({\color{var}x})}}\equiv
  \ensuremath{\mathtt{add(size(),{\color{var}x})}}$ and $\ensuremath{\mathtt{pop()}}\equiv \ensuremath{\mathtt{remove(size()-1)}}$), then these
  operations run in constant amortized time, independent of the value
  of \ensuremath{\mathtt{{\color{var}b}}}.
\end{exc}

\begin{exc}
  Design and implement of a version of an \texttt{SEList} that supports all
  the \texttt{Deque} operations in constant amortized time per operation,
  independent of the value of \ensuremath{\mathtt{{\color{var}b}}}.
\end{exc}

\begin{exc}
  Explain how to use the bitwise exclusive-or operator, \verb+^+, to
  swap the values of two \ensuremath{\mathtt{{\color{keyword}int}}} variables without using a third variable.
\end{exc}






\clearpage

\clearpage
\chapter{Skiplists}

\chaplabel{skiplists}

In this chapter, we discuss a beautiful data structure: the skiplist,
which has a variety of applications.  Using a skiplist we can implement a
\texttt{List} that has $O(\log n)$ time implementations of \ensuremath{\mathtt{get({\color{var}i})}}, \ensuremath{\mathtt{set({\color{var}i},{\color{var}x})}},
\ensuremath{\mathtt{add({\color{var}i},{\color{var}x})}}, and \ensuremath{\mathtt{remove({\color{var}i})}}. We can also implement an \texttt{SSet} in which
all operations run in $O(\log \ensuremath{\mathtt{{\color{var}n}}})$ expected time.

The efficiency of skiplists relies on their use of randomization.
When a new element is added to a skiplist, the skiplist uses random coin
tosses to determine the height of the new element.  The performance of
skiplists is expressed in terms of expected running times and path
lengths. This expectation is taken over the random coin tosses used by
the skiplist.  In the implementation, the random coin tosses used by a
skiplist are simulated using a pseudo-random number (or bit) generator.

\section{The Basic Structure}

\index{skiplist}%
Conceptually, a skiplist is a sequence of singly-linked lists
$L_0,\ldots,L_h$. Each list $L_r$ contains a subset of the items
in $L_{r-1}$.  We start with the input list $L_0$ that contains \ensuremath{\mathtt{{\color{var}n}}}
items and construct $L_1$ from $L_0$, $L_2$ from $L_1$, and so on.
The items in $L_r$ are obtained by tossing a coin for each element, \ensuremath{\mathtt{{\color{var}x}}},
in $L_{r-1}$ and including \ensuremath{\mathtt{{\color{var}x}}} in $L_r$ if the coin turns up as heads.
This process ends when we create a list $L_r$ that is empty.  An example
of a skiplist is shown in \figref{skiplist}.

\begin{figure}
  \begin{center}
    \includegraphics[width=\ScaleIfNeeded]{figs/skiplist}
  \end{center}
  \caption{A skiplist containing seven elements.}
  \figlabel{skiplist}
\end{figure}

For an element, \ensuremath{\mathtt{{\color{var}x}}}, in a skiplist, we call the \emph{height}
\index{height!of a skiplist}%
of \ensuremath{\mathtt{{\color{var}x}}} the
largest value $r$ such that \ensuremath{\mathtt{{\color{var}x}}} appears in $L_r$.  Thus, for example,
elements that only appear in $L_0$ have height $0$.  If we spend a few
moments thinking about it, we notice that the height of \ensuremath{\mathtt{{\color{var}x}}} corresponds
to the following experiment:  Toss a coin repeatedly until it comes
up as tails.  How many times did it come up as heads?  The answer, not
surprisingly, is that the expected height of a node is 1. (We expect to
toss the coin twice before getting tails, but we don't count the last
toss.) The \emph{height} of a skiplist is the height of its tallest node.

At the head of every list is a special node, called the \emph{sentinel},
\index{sentinel node}%
that acts as a dummy node for the list. The key property of skiplists
is that there is a short path, called the \emph{search path}, 
\index{search path!in a skiplist}%
from the
sentinel in $L_h$ to every node in $L_0$.  Remembering how to construct
a search path for a node, \ensuremath{\mathtt{{\color{var}u}}}, is easy (see \figref{skiplist-searchpath})
:  Start at the top left corner of your skiplist (the sentinel in $L_h$)
and always go right unless that would overshoot \ensuremath{\mathtt{{\color{var}u}}}, in which case you
should take a step down into the list below.

More precisely, to construct the search path for the node \ensuremath{\mathtt{{\color{var}u}}} in $L_0$,
we start at the sentinel, \ensuremath{\mathtt{{\color{var}w}}}, in $L_h$.  Next, we examine \ensuremath{\mathtt{{\color{var}w}.{\color{var}next}}}.
If \ensuremath{\mathtt{{\color{var}w}.{\color{var}next}}} contains an item that appears before \ensuremath{\mathtt{{\color{var}u}}} in $L_0$, then
we set $\ensuremath{\mathtt{{\color{var}w}}}=\ensuremath{\mathtt{{\color{var}w}.{\color{var}next}}}$.  Otherwise, we move down and continue the search
at the occurrence of \ensuremath{\mathtt{{\color{var}w}}} in the list $L_{h-1}$.  We continue this way
until we reach the predecessor of \ensuremath{\mathtt{{\color{var}u}}} in $L_0$. 
\begin{figure}
  \begin{center}
    \includegraphics[width=\ScaleIfNeeded]{figs/skiplist-searchpath}
  \end{center}
  \caption{The search path for the node containing $4$ in a skiplist.}
  \figlabel{skiplist-searchpath}
\end{figure}

The following result, which we will prove in \secref{skiplist-analysis},
shows that the search path is quite short:

\begin{lem}\lemlabel{skiplist-searchpath}
The expected length of the search path for any node, \ensuremath{\mathtt{{\color{var}u}}}, in $L_0$ is at
most $2\log \ensuremath{\mathtt{{\color{var}n}}} + O(1) = O(\log \ensuremath{\mathtt{{\color{var}n}}})$.
\end{lem}

A space-efficient way to implement a skiplist is to define a \texttt{Node},
\ensuremath{\mathtt{{\color{var}u}}}, as consisting of a data value, \ensuremath{\mathtt{{\color{var}x}}}, and an array, \ensuremath{\mathtt{{\color{var}next}}}, of
pointers, where \ensuremath{\mathtt{{\color{var}u}.{\color{var}next}[{\color{var}i}]}} points to \ensuremath{\mathtt{{\color{var}u}}}'s successor in the list
$L_{\ensuremath{\mathtt{{\color{var}i}}}}$.  In this way, the data, \ensuremath{\mathtt{{\color{var}x}}}, in a node is
\javaonly{referenced}\cpponly{stored}\pcodeonly{referenced} 
only once, even though \ensuremath{\mathtt{{\color{var}x}}} may appear in several lists.



 \ \vspace{.5ex}\begin{Verbatim}[gobble=1,tabsize=2,frame=single,framerule=1.2pt,commandchars=\\@\$,label=\texttt{SkiplistSSet},labelposition=topline]
	@\color@keyword$class$ Node<@\color@keyword$T$> {
		@\color@keyword$T$ @\color@var$x$;
		Node<@\color@keyword$T$>[] @\color@var$next$;
		Node(@\color@keyword$T$ @\color@var$ix$, @\color@keyword$int$ @\color@var$h$) {
			@\color@var$x$ = @\color@var$ix$;
			@\color@var$next$ = Array.newInstance(Node.@\color@keyword$class$, @\color@var$h$+1);
		}
		@\color@keyword$int$ height() {
			@\color@var$return$ @\color@var$next$.@\color@var$length$ - 1;
		}
	}
\end{Verbatim}


The next two sections of this chapter discuss two different applications
of skiplists.  In each of these applications, $L_0$ stores the main
structure (a list of elements or a sorted set of elements).
The primary difference between these structures is in how
a search path is navigated; in particular, they differ in how
they decide if a search path should go down into $L_{r-1}$ or go right
within $L_r$.

\section{\texttt{SkiplistSSet}: An Efficient \texttt{SSet}}
\seclabel{skiplistset}

\index{SkiplistSSet@\texttt{SkiplistSSet}}%
A \texttt{SkiplistSSet} uses a skiplist structure to implement the \texttt{SSet}
interface.   When used in this way, the list $L_0$ stores the elements of
the \texttt{SSet} in sorted order.  The \ensuremath{\mathtt{find({\color{var}x})}} method works by following
the search path for the smallest value \ensuremath{\mathtt{{\color{var}y}}} such that $\ensuremath{\mathtt{{\color{var}y}}}\ge\ensuremath{\mathtt{{\color{var}x}}}$:



 \ \vspace{.5ex}\begin{Verbatim}[gobble=1,tabsize=2,frame=single,framerule=1.2pt,commandchars=\\@\$,label=\texttt{SkiplistSSet},labelposition=topline]
	Node<@\color@keyword$T$> findPredNode(@\color@keyword$T$ @\color@var$x$) {
		Node<@\color@keyword$T$> @\color@var$u$ = @\color@var$sentinel$;
		@\color@keyword$int$ @\color@var$r$ = @\color@var$h$;
		@\color@keyword$while$ (@\color@var$r$ >= 0) {
			@\color@keyword$while$ (@\color@var$u$.@\color@var$next$[@\color@var$r$] != @\color@keyword$null$ && compare(@\color@var$u$.@\color@var$next$[@\color@var$r$].@\color@var$x$,@\color@var$x$) < 0)
				@\color@var$u$ = @\color@var$u$.@\color@var$next$[@\color@var$r$];   @\color@comment$// go right in list r$
			@\color@var$r$--;               @\color@comment$// go down into list r-1$
		}
		@\color@var$return$ @\color@var$u$;
	}
	@\color@keyword$T$ find(@\color@keyword$T$ @\color@var$x$) {
		Node<@\color@keyword$T$> @\color@var$u$ = findPredNode(@\color@var$x$);
		@\color@var$return$ @\color@var$u$.@\color@var$next$[0] == @\color@keyword$null$ ? @\color@keyword$null$ : @\color@var$u$.@\color@var$next$[0].@\color@var$x$;
	}
\end{Verbatim}


Following the search path for \ensuremath{\mathtt{{\color{var}y}}} is easy:  when situated at
some node, \ensuremath{\mathtt{{\color{var}u}}}, in  $L_{\ensuremath{\mathtt{{\color{var}r}}}}$, we look right to \ensuremath{\mathtt{{\color{var}u}.{\color{var}next}[{\color{var}r}].{\color{var}x}}}.
If $\ensuremath{\mathtt{{\color{var}x}}}>\ensuremath{\mathtt{{\color{var}u}.{\color{var}next}[{\color{var}r}].{\color{var}x}}}$, then we take a step to the right in
$L_{\ensuremath{\mathtt{{\color{var}r}}}}$; otherwise, we move down into $L_{\ensuremath{\mathtt{{\color{var}r}}}-1}$.  Each step
(right or down) in this search takes only constant time; thus, by
\lemref{skiplist-searchpath}, the expected running time of \ensuremath{\mathtt{find({\color{var}x})}}
is $O(\log \ensuremath{\mathtt{{\color{var}n}}})$.

Before we can add an element to a \texttt{SkipListSSet}, we need a method to
simulate tossing coins to determine the height, \ensuremath{\mathtt{{\color{var}k}}}, of a new node.
We do so by picking a random integer, \ensuremath{\mathtt{{\color{var}z}}}, and counting the number of
trailing $1$s in the binary representation of \ensuremath{\mathtt{{\color{var}z}}}:\footnote{This method
does not exactly replicate the coin-tossing experiment since the value of
\ensuremath{\mathtt{{\color{var}k}}} will always be less than the number of bits in an \ensuremath{\mathtt{{\color{keyword}int}}}.  However,
this will have negligible impact unless the number of elements in the
structure is much greater than $2^{32}=4294967296$.}



 \ \vspace{.5ex}\begin{Verbatim}[gobble=1,tabsize=2,frame=single,framerule=1.2pt,commandchars=\\@\$,label=\texttt{SkiplistSSet},labelposition=topline]
	@\color@keyword$int$ pickHeight() {
		@\color@keyword$int$ @\color@var$z$ = @\color@var$rand$.nextInt();
		@\color@keyword$int$ @\color@var$k$ = 0;
		@\color@keyword$int$ @\color@var$m$ = 1;
		@\color@keyword$while$ ((@\color@var$z$ & @\color@var$m$) != 0) {
			@\color@var$k$++;
			@\color@var$m$ <<= 1;
		}
		@\color@var$return$ @\color@var$k$;
	}
\end{Verbatim}


To implement the \ensuremath{\mathtt{add({\color{var}x})}} method in a \texttt{SkiplistSSet} we search for \ensuremath{\mathtt{{\color{var}x}}}
and then splice \ensuremath{\mathtt{{\color{var}x}}} into a few lists $L_0$,\ldots,$L_{\ensuremath{\mathtt{{\color{var}k}}}}$, where \ensuremath{\mathtt{{\color{var}k}}}
is selected using the \ensuremath{\mathtt{pickHeight()}} method. The easiest way to do this
is to use an array, \ensuremath{\mathtt{{\color{var}stack}}}, that keeps track of the nodes at which
the search path goes down from some list $L_{\ensuremath{\mathtt{{\color{var}r}}}}$ into $L_{\ensuremath{\mathtt{{\color{var}r}}}-1}$.
More precisely, \ensuremath{\mathtt{{\color{var}stack}[{\color{var}r}]}} is the node in $L_{\ensuremath{\mathtt{{\color{var}r}}}}$ where the search path
proceeded down into $L_{\ensuremath{\mathtt{{\color{var}r}}}-1}$.  The nodes that we modify to insert \ensuremath{\mathtt{{\color{var}x}}}
are precisely the nodes $\ensuremath{\mathtt{{\color{var}stack}[0]}},\ldots,\ensuremath{\mathtt{{\color{var}stack}[{\color{var}k}]}}$.  The following
code implements this algorithm for \ensuremath{\mathtt{add({\color{var}x})}}:
\label{pg:skiplist-add}


 \ \vspace{.5ex}\begin{Verbatim}[gobble=1,tabsize=2,frame=single,framerule=1.2pt,commandchars=\\@\$,label=\texttt{SkiplistSSet},labelposition=topline]
	@\color@var$boolean$ add(@\color@keyword$T$ @\color@var$x$) {
		Node<@\color@keyword$T$> @\color@var$u$ = @\color@var$sentinel$;
		@\color@keyword$int$ @\color@var$r$ = @\color@var$h$;
		@\color@keyword$int$ @\color@var$comp$ = 0;
		@\color@keyword$while$ (@\color@var$r$ >= 0) {
			@\color@keyword$while$ (@\color@var$u$.@\color@var$next$[@\color@var$r$] != @\color@keyword$null$ 
			       && (@\color@var$comp$ = compare(@\color@var$u$.@\color@var$next$[@\color@var$r$].@\color@var$x$,@\color@var$x$)) < 0)
				@\color@var$u$ = @\color@var$u$.@\color@var$next$[@\color@var$r$];
			@\color@keyword$if$ (@\color@var$u$.@\color@var$next$[@\color@var$r$] != @\color@keyword$null$ && @\color@var$comp$ == 0) @\color@var$return$ @\color@var$false$;
			@\color@var$stack$[@\color@var$r$--] = @\color@var$u$;          @\color@comment$// going down, store u$
		}
		Node<@\color@keyword$T$> @\color@var$w$ = @\color@keyword$new$ Node<@\color@keyword$T$>(@\color@var$x$, pickHeight());
		@\color@keyword$while$ (@\color@var$h$ < @\color@var$w$.height())
			@\color@var$stack$[++@\color@var$h$] = @\color@var$sentinel$;   @\color@comment$// height increased$
		@\color@keyword$for$ (@\color@keyword$int$ @\color@var$i$ = 0; @\color@var$i$ < @\color@var$w$.@\color@var$next$.@\color@var$length$; @\color@var$i$++) {
			@\color@var$w$.@\color@var$next$[@\color@var$i$] = @\color@var$stack$[@\color@var$i$].@\color@var$next$[@\color@var$i$];
			@\color@var$stack$[@\color@var$i$].@\color@var$next$[@\color@var$i$] = @\color@var$w$;
		}
		@\color@var$n$++;
		@\color@var$return$ @\color@var$true$;
	}
\end{Verbatim}


\begin{figure}
  \begin{center}
    \includegraphics[width=\ScaleIfNeeded]{figs/skiplist-add}
  \end{center}
  \caption[Adding to a skiplist]{Adding the node containing $3.5$ to a skiplist.  The nodes stored in \ensuremath{\mathtt{{\color{var}stack}}}
  are highlighted.}
  \figlabel{skiplist-add}
\end{figure}

Removing an element, \ensuremath{\mathtt{{\color{var}x}}}, is done in a similar way, except that there
is no need for \ensuremath{\mathtt{{\color{var}stack}}} to keep track of the search path.  The removal
can be done as we are following the search path.  We search for \ensuremath{\mathtt{{\color{var}x}}}
and each time the search moves downward from a node \ensuremath{\mathtt{{\color{var}u}}}, we check if
$\ensuremath{\mathtt{{\color{var}u}.{\color{var}next}.{\color{var}x}}}=\ensuremath{\mathtt{{\color{var}x}}}$ and if so, we splice \ensuremath{\mathtt{{\color{var}u}}} out of the list:


 \ \vspace{.5ex}\begin{Verbatim}[gobble=1,tabsize=2,frame=single,framerule=1.2pt,commandchars=\\@\$,label=\texttt{SkiplistSSet},labelposition=topline]
	@\color@var$boolean$ remove(@\color@keyword$T$ @\color@var$x$) {
		@\color@var$boolean$ @\color@var$removed$ = @\color@var$false$;
		Node<@\color@keyword$T$> @\color@var$u$ = @\color@var$sentinel$;
		@\color@keyword$int$ @\color@var$r$ = @\color@var$h$;
		@\color@keyword$int$ @\color@var$comp$ = 0;
		@\color@keyword$while$ (@\color@var$r$ >= 0) {
			@\color@keyword$while$ (@\color@var$u$.@\color@var$next$[@\color@var$r$] != @\color@keyword$null$ 
			       && (@\color@var$comp$ = compare(@\color@var$u$.@\color@var$next$[@\color@var$r$].@\color@var$x$, @\color@var$x$)) < 0) {
				@\color@var$u$ = @\color@var$u$.@\color@var$next$[@\color@var$r$];
			}
			@\color@keyword$if$ (@\color@var$u$.@\color@var$next$[@\color@var$r$] != @\color@keyword$null$ && @\color@var$comp$ == 0) {
				@\color@var$removed$ = @\color@var$true$;
				@\color@var$u$.@\color@var$next$[@\color@var$r$] = @\color@var$u$.@\color@var$next$[@\color@var$r$].@\color@var$next$[@\color@var$r$];
				@\color@keyword$if$ (@\color@var$u$ == @\color@var$sentinel$ && @\color@var$u$.@\color@var$next$[@\color@var$r$] == @\color@keyword$null$)
					@\color@var$h$--;  @\color@comment$// height has gone down$
			}
			@\color@var$r$--;
		}
		@\color@keyword$if$ (@\color@var$removed$) @\color@var$n$--;
		@\color@var$return$ @\color@var$removed$;
	}
\end{Verbatim}


\begin{figure}
  \begin{center}
    \includegraphics[width=\ScaleIfNeeded]{figs/skiplist-remove}
  \end{center}
  \caption{Removing the node containing $3$ from a skiplist.}
  \figlabel{skiplist-remove}
\end{figure}

\subsection{Summary}

The following theorem summarizes the performance of skiplists when used to
implement sorted sets:

\begin{thm}\thmlabel{skiplist}
\texttt{SkiplistSSet} implements the \texttt{SSet} interface. A \texttt{SkiplistSSet} supports
the operations \ensuremath{\mathtt{add({\color{var}x})}}, \ensuremath{\mathtt{remove({\color{var}x})}}, and \ensuremath{\mathtt{find({\color{var}x})}} in $O(\log \ensuremath{\mathtt{{\color{var}n}}})$
expected time per operation.
\end{thm}

\section{\texttt{SkiplistList}: An Efficient Random-Access \texttt{List}}
\seclabel{skiplistlist}

\index{SkiplistList@\texttt{SkiplistList}}%
A \texttt{SkiplistList} implements the \texttt{List} interface using a skiplist
structure.  In a \texttt{SkiplistList}, $L_0$ contains the elements of the
list in the order in which they appear in the list.   As in a
\texttt{SkiplistSSet}, elements can be added, removed, and accessed in $O(\log
\ensuremath{\mathtt{{\color{var}n}}})$ time.

For this to be possible, we need a way to follow the search path for
the \ensuremath{\mathtt{{\color{var}i}}}th element in $L_0$.  The easiest way to do this is to define
the notion of the \emph{length} of an edge in some list, $L_{\ensuremath{\mathtt{{\color{var}r}}}}$.
We define the length of every edge in $L_{0}$ as 1.  The length of an edge, \ensuremath{\mathtt{{\color{var}e}}},
in $L_{\ensuremath{\mathtt{{\color{var}r}}}}$, $\ensuremath{\mathtt{{\color{var}r}}}>0$, is defined as the sum of the lengths of the edges below \ensuremath{\mathtt{{\color{var}e}}}
in $L_{\ensuremath{\mathtt{{\color{var}r}}}-1}$.  Equivalently, the length of \ensuremath{\mathtt{{\color{var}e}}} is
the number of edges in $L_0$ below \ensuremath{\mathtt{{\color{var}e}}}.  See \figref{skiplist-lengths} for
an example of a skiplist with the lengths of its edges shown.  Since the
edges of skiplists are stored in arrays, the lengths can be stored the same
way:

\begin{figure}
  \begin{center}
    \includegraphics[width=\ScaleIfNeeded]{figs/skiplist-lengths}
  \end{center}
  \caption{The lengths of the edges in a skiplist.}
  \figlabel{skiplist-lengths}
\end{figure}



 \ \vspace{.5ex}\begin{Verbatim}[gobble=1,tabsize=2,frame=single,framerule=1.2pt,commandchars=\\@\$,label=\texttt{SkiplistList},labelposition=topline]
	@\color@keyword$class$ Node {
		@\color@keyword$T$ @\color@var$x$;
		Node[] @\color@var$next$;
		@\color@keyword$int$[] @\color@var$length$;
		Node(@\color@keyword$T$ @\color@var$ix$, @\color@keyword$int$ @\color@var$h$) {
			@\color@var$x$ = @\color@var$ix$;
			@\color@var$next$ = Array.newInstance(Node.@\color@keyword$class$, @\color@var$h$+1);
			@\color@var$length$ = @\color@keyword$new$ @\color@keyword$int$[@\color@var$h$+1];
		}
		@\color@keyword$int$ height() {
			@\color@var$return$ @\color@var$next$.@\color@var$length$ - 1;
		}
	}
\end{Verbatim}


The useful property of this definition of length is that, if we are
currently at a node that is at position \ensuremath{\mathtt{{\color{var}j}}} in $L_0$ and we follow an
edge of length $\ell$, then we move to a node whose position, in $L_0$,
is $\ensuremath{\mathtt{{\color{var}j}}}+\ell$.  In this way, while following a search path, we can keep
track of the position, \ensuremath{\mathtt{{\color{var}j}}}, of the current node in $L_0$.  When at a
node, \ensuremath{\mathtt{{\color{var}u}}}, in $L_{\ensuremath{\mathtt{{\color{var}r}}}}$, we go right if \ensuremath{\mathtt{{\color{var}j}}} plus the length of the edge
\ensuremath{\mathtt{{\color{var}u}.{\color{var}next}[{\color{var}r}]}} is less than \ensuremath{\mathtt{{\color{var}i}}}. Otherwise, we go down into $L_{\ensuremath{\mathtt{{\color{var}r}}}-1}$.



 \ \vspace{.5ex}\begin{Verbatim}[gobble=1,tabsize=2,frame=single,framerule=1.2pt,commandchars=\\@\$,label=\texttt{SkiplistList},labelposition=topline]
	Node findPred(@\color@keyword$int$ @\color@var$i$) {
		Node @\color@var$u$ = @\color@var$sentinel$;
		@\color@keyword$int$ @\color@var$r$ = @\color@var$h$;
		@\color@keyword$int$ @\color@var$j$ = -1;   @\color@comment$// index of the current node in list 0$
		@\color@keyword$while$ (@\color@var$r$ >= 0) {
			@\color@keyword$while$ (@\color@var$u$.@\color@var$next$[@\color@var$r$] != @\color@keyword$null$ && @\color@var$j$ + @\color@var$u$.@\color@var$length$[@\color@var$r$] < @\color@var$i$) {
				@\color@var$j$ += @\color@var$u$.@\color@var$length$[@\color@var$r$];
				@\color@var$u$ = @\color@var$u$.@\color@var$next$[@\color@var$r$];
			}
			@\color@var$r$--;
		}
		@\color@var$return$ @\color@var$u$;
	}
\end{Verbatim}



 \ \vspace{.5ex}\begin{Verbatim}[gobble=1,tabsize=2,frame=single,framerule=1.2pt,commandchars=\\@\$,label=\texttt{SkiplistList},labelposition=topline]
	@\color@keyword$T$ get(@\color@keyword$int$ @\color@var$i$) {
		@\color@var$return$ findPred(@\color@var$i$).@\color@var$next$[0].@\color@var$x$;
	}
	@\color@keyword$T$ set(@\color@keyword$int$ @\color@var$i$, @\color@keyword$T$ @\color@var$x$) {
		Node @\color@var$u$ = findPred(@\color@var$i$).@\color@var$next$[0];
		@\color@keyword$T$ @\color@var$y$ = @\color@var$u$.@\color@var$x$;
		@\color@var$u$.@\color@var$x$ = @\color@var$x$;
		@\color@var$return$ @\color@var$y$;
	}
\end{Verbatim}


Since the hardest part of the operations \ensuremath{\mathtt{get({\color{var}i})}} and \ensuremath{\mathtt{set({\color{var}i},{\color{var}x})}} is
finding the \ensuremath{\mathtt{{\color{var}i}}}th node in $L_0$, these operations run in
$O(\log \ensuremath{\mathtt{{\color{var}n}}})$ time.

Adding an element to a \texttt{SkiplistList} at a position, \ensuremath{\mathtt{{\color{var}i}}}, is fairly
simple.  Unlike in a \texttt{SkiplistSSet}, we are sure that a new
node will actually be added, so we can do the addition at the same time
as we search for the new node's location. We first pick the height, \ensuremath{\mathtt{{\color{var}k}}},
of the newly inserted node, \ensuremath{\mathtt{{\color{var}w}}}, and then follow the search path for \ensuremath{\mathtt{{\color{var}i}}}.
Any time the search path moves down from $L_{\ensuremath{\mathtt{{\color{var}r}}}}$ with $\ensuremath{\mathtt{{\color{var}r}}}\le \ensuremath{\mathtt{{\color{var}k}}}$, we
splice \ensuremath{\mathtt{{\color{var}w}}} into $L_{\ensuremath{\mathtt{{\color{var}r}}}}$.  The only extra care needed is to ensure that
the lengths of edges are updated properly.  See \figref{skiplist-addix}.

\begin{figure}
  \begin{center}
    \includegraphics[width=\ScaleIfNeeded]{figs/skiplist-addix}
  \end{center}
  \caption[Adding to a SkiplistList]{Adding an element to a \texttt{SkiplistList}.}
  \figlabel{skiplist-addix}
\end{figure}

Note that, each time the search path goes down at a node, \ensuremath{\mathtt{{\color{var}u}}}, in $L_{\ensuremath{\mathtt{{\color{var}r}}}}$,
the length of the edge \ensuremath{\mathtt{{\color{var}u}.{\color{var}next}[{\color{var}r}]}} increases by one, since we are adding
an element below that edge at position \ensuremath{\mathtt{{\color{var}i}}}.  Splicing  the node \ensuremath{\mathtt{{\color{var}w}}} between two nodes,
\ensuremath{\mathtt{{\color{var}u}}} and \ensuremath{\mathtt{{\color{var}z}}}, works as shown in \figref{skiplist-lengths-splice}.  While
following the search path we are already keeping track of the position,
\ensuremath{\mathtt{{\color{var}j}}}, of \ensuremath{\mathtt{{\color{var}u}}} in $L_0$.  Therefore, we know that the length of the edge from
\ensuremath{\mathtt{{\color{var}u}}} to \ensuremath{\mathtt{{\color{var}w}}} is $\ensuremath{\mathtt{{\color{var}i}}}-\ensuremath{\mathtt{{\color{var}j}}}$.  We can also deduce the length of the edge
from \ensuremath{\mathtt{{\color{var}w}}}  to \ensuremath{\mathtt{{\color{var}z}}} from the length, $\ell$, of the edge from \ensuremath{\mathtt{{\color{var}u}}} to \ensuremath{\mathtt{{\color{var}z}}}.
Therefore, we can splice in \ensuremath{\mathtt{{\color{var}w}}} and update the lengths of the edges in
constant time.

\begin{figure}
  \begin{center}
    \includegraphics[scale=0.90909]{figs/skiplist-lengths-splice}
  \end{center}
  \caption[Adding to a SkiplistList]{Updating the lengths of edges while splicing a node 
   \ensuremath{\mathtt{{\color{var}w}}} into a skiplist.}
  \figlabel{skiplist-lengths-splice}
\end{figure}

This sounds more complicated than it is, for the code is actually
quite simple:



 \ \vspace{.5ex}\begin{Verbatim}[gobble=1,tabsize=2,frame=single,framerule=1.2pt,commandchars=\\@\$,label=\texttt{SkiplistList},labelposition=topline]
	@\color@keyword$void$ add(@\color@keyword$int$ @\color@var$i$, @\color@keyword$T$ @\color@var$x$) {
		Node @\color@var$w$ = @\color@keyword$new$ Node(@\color@var$x$, pickHeight());
		@\color@keyword$if$ (@\color@var$w$.height() > @\color@var$h$) 
			@\color@var$h$ = @\color@var$w$.height();
		add(@\color@var$i$, @\color@var$w$);
	}
\end{Verbatim}



 \ \vspace{.5ex}\begin{Verbatim}[gobble=1,tabsize=2,frame=single,framerule=1.2pt,commandchars=\\@\$,label=\texttt{SkiplistList},labelposition=topline]
	Node add(@\color@keyword$int$ @\color@var$i$, Node @\color@var$w$) {
		Node @\color@var$u$ = @\color@var$sentinel$;
		@\color@keyword$int$ @\color@var$k$ = @\color@var$w$.height();
		@\color@keyword$int$ @\color@var$r$ = @\color@var$h$;
		@\color@keyword$int$ @\color@var$j$ = -1; @\color@comment$// index of u$
		@\color@keyword$while$ (@\color@var$r$ >= 0) {
			@\color@keyword$while$ (@\color@var$u$.@\color@var$next$[@\color@var$r$] != @\color@keyword$null$ && @\color@var$j$+@\color@var$u$.@\color@var$length$[@\color@var$r$] < @\color@var$i$) {
				@\color@var$j$ += @\color@var$u$.@\color@var$length$[@\color@var$r$];
				@\color@var$u$ = @\color@var$u$.@\color@var$next$[@\color@var$r$];
			}
			@\color@var$u$.@\color@var$length$[@\color@var$r$]++; @\color@comment$// accounts for new node in list 0$
			@\color@keyword$if$ (@\color@var$r$ <= @\color@var$k$) {
				@\color@var$w$.@\color@var$next$[@\color@var$r$] = @\color@var$u$.@\color@var$next$[@\color@var$r$];
				@\color@var$u$.@\color@var$next$[@\color@var$r$] = @\color@var$w$;
				@\color@var$w$.@\color@var$length$[@\color@var$r$] = @\color@var$u$.@\color@var$length$[@\color@var$r$] - (@\color@var$i$ - @\color@var$j$);
				@\color@var$u$.@\color@var$length$[@\color@var$r$] = @\color@var$i$ - @\color@var$j$;
			}
			@\color@var$r$--;
		}
		@\color@var$n$++;
		@\color@var$return$ @\color@var$u$;
	}
\end{Verbatim}



By now, the implementation of 
the \ensuremath{\mathtt{remove({\color{var}i})}} operation in a \texttt{SkiplistList} should be obvious.  We follow the search path for the node at position \ensuremath{\mathtt{{\color{var}i}}}.  Each time the search path takes a step down from a node, \ensuremath{\mathtt{{\color{var}u}}}, at level \ensuremath{\mathtt{{\color{var}r}}} we decrement the length of the edge leaving \ensuremath{\mathtt{{\color{var}u}}} at that level.  We also check if \ensuremath{\mathtt{{\color{var}u}.{\color{var}next}[{\color{var}r}]}} is the element of rank \ensuremath{\mathtt{{\color{var}i}}} and, if so, splice it out of the list at that level.   An example is shown in \figref{skiplist-removei}.
\begin{figure}
  \begin{center}
    \includegraphics[width=\ScaleIfNeeded]{figs/skiplist-removei}
  \end{center}
  \caption[Removing an element from a SkiplistList]{Removing an element from a \texttt{SkiplistList}.}
  \figlabel{skiplist-removei}
\end{figure}


 \ \vspace{.5ex}\begin{Verbatim}[gobble=1,tabsize=2,frame=single,framerule=1.2pt,commandchars=\\@\$,label=\texttt{SkiplistList},labelposition=topline]
	@\color@keyword$T$ remove(@\color@keyword$int$ @\color@var$i$) {
		@\color@keyword$T$ @\color@var$x$ = @\color@keyword$null$;
		Node @\color@var$u$ = @\color@var$sentinel$;
		@\color@keyword$int$ @\color@var$r$ = @\color@var$h$;
		@\color@keyword$int$ @\color@var$j$ = -1; @\color@comment$// index of node u$
		@\color@keyword$while$ (@\color@var$r$ >= 0) {
			@\color@keyword$while$ (@\color@var$u$.@\color@var$next$[@\color@var$r$] != @\color@keyword$null$ && @\color@var$j$+@\color@var$u$.@\color@var$length$[@\color@var$r$] < @\color@var$i$) {
				@\color@var$j$ += @\color@var$u$.@\color@var$length$[@\color@var$r$];
				@\color@var$u$ = @\color@var$u$.@\color@var$next$[@\color@var$r$];
			}
			@\color@var$u$.@\color@var$length$[@\color@var$r$]--;  @\color@comment$// for the node we are removing$
			@\color@keyword$if$ (@\color@var$j$ + @\color@var$u$.@\color@var$length$[@\color@var$r$] + 1 == @\color@var$i$ && @\color@var$u$.@\color@var$next$[@\color@var$r$] != @\color@keyword$null$) {
				@\color@var$x$ = @\color@var$u$.@\color@var$next$[@\color@var$r$].@\color@var$x$;
				@\color@var$u$.@\color@var$length$[@\color@var$r$] += @\color@var$u$.@\color@var$next$[@\color@var$r$].@\color@var$length$[@\color@var$r$];
				@\color@var$u$.@\color@var$next$[@\color@var$r$] = @\color@var$u$.@\color@var$next$[@\color@var$r$].@\color@var$next$[@\color@var$r$];
				@\color@keyword$if$ (@\color@var$u$ == @\color@var$sentinel$ && @\color@var$u$.@\color@var$next$[@\color@var$r$] == @\color@keyword$null$)
					@\color@var$h$--;
			}
			@\color@var$r$--;
		}
		@\color@var$n$--;
		@\color@var$return$ @\color@var$x$;
	}
\end{Verbatim}


\subsection{Summary}

The following theorem summarizes the performance of the \texttt{SkiplistList}
data structure:

\begin{thm}\thmlabel{skiplistlist}
  A \texttt{SkiplistList} implements the \texttt{List} interface.  A \texttt{SkiplistList}
  supports the operations \ensuremath{\mathtt{get({\color{var}i})}}, \ensuremath{\mathtt{set({\color{var}i},{\color{var}x})}}, \ensuremath{\mathtt{add({\color{var}i},{\color{var}x})}}, and
  \ensuremath{\mathtt{remove({\color{var}i})}} in $O(\log \ensuremath{\mathtt{{\color{var}n}}})$ expected time per operation.
\end{thm}


\section{Analysis of Skiplists}
\seclabel{skiplist-analysis}

In this section, we analyze the expected height, size, and length of
the search path in a skiplist.  This section requires a background in
basic probability.  Several proofs are based on the following basic
observation about coin tosses.


\begin{lem}\lemlabel{coin-tosses}
  \index{coin toss}%
  Let $T$ be the number of times a fair coin is tossed up to and including
  the first time the coin comes up heads.  Then $\E[T]=2$.
\end{lem}

\begin{proof}
  Suppose we stop tossing the coin the first time it comes up
  heads. Define the indicator variable
  \[ I_{i} = \left\{\begin{array}{ll}
     0 & \mbox{if the coin is tossed less than $i$ times} \\
     1 & \mbox{if the coin is tossed $i$ or more times} 
     \end{array}\right.
  \]
  Note that $I_i=1$ if and only if the first $i-1$ coin tosses are tails,
  so $\E[I_i]=\Pr\{I_i=1\}=1/2^{i-1}$.  Observe that $T$, the total
  number of coin tosses, can be written as $T=\sum_{i=1}^{\infty} I_i$.
  Therefore,
  \begin{align*}
    \E[T] & =  \E\left[\sum_{i=1}^\infty I_i\right] \\
     & =  \sum_{i=1}^\infty \E\left[I_i\right] \\
     & =  \sum_{i=1}^\infty 1/2^{i-1} \\
     & =  1 + 1/2 + 1/4 + 1/8 + \cdots \\
     & =  2 \enspace .   \qedhere
  \end{align*} 
\end{proof}

The next two lemmata tell us that skiplists have linear size:

\begin{lem}\lemlabel{skiplist-size1}
  The expected number of nodes in a skiplist containing $\ensuremath{\mathtt{{\color{var}n}}}$ elements,
  not including occurrences of the sentinel, is $2\ensuremath{\mathtt{{\color{var}n}}}$.
\end{lem}

\begin{proof}
  The probability that any particular element, \ensuremath{\mathtt{{\color{var}x}}}, is included in list
  $L_{\ensuremath{\mathtt{{\color{var}r}}}}$ is $1/2^{\ensuremath{\mathtt{{\color{var}r}}}}$, so the expected number of nodes in $L_{\ensuremath{\mathtt{{\color{var}r}}}}$
  is $\ensuremath{\mathtt{{\color{var}n}}}/2^{\ensuremath{\mathtt{{\color{var}r}}}}$.\footnote{See \secref{randomization} to see how this
  is derived using indicator variables and linearity of expectation.}
  Therefore, the total expected number of nodes in all lists is
  \[ \sum_{\ensuremath{\mathtt{{\color{var}r}}}=0}^\infty \ensuremath{\mathtt{{\color{var}n}}}/2^{\ensuremath{\mathtt{{\color{var}r}}}} = \ensuremath{\mathtt{{\color{var}n}}}(1+1/2+1/4+1/8+\cdots) = 2\ensuremath{\mathtt{{\color{var}n}}} \enspace . \qedhere \]
\end{proof}

\begin{lem}\lemlabel{skiplist-height}
  The expected height of a skiplist containing \ensuremath{\mathtt{{\color{var}n}}} elements is at most
  $\log \ensuremath{\mathtt{{\color{var}n}}} + 2$.
\end{lem}

\begin{proof}
  For each $\ensuremath{\mathtt{{\color{var}r}}}\in\{1,2,3,\ldots,\infty\}$, 
  define the indicator random variable
  \[ I_{\ensuremath{\mathtt{{\color{var}r}}}} = \left\{\begin{array}{ll}
     0 & \mbox{if $L_{\ensuremath{\mathtt{{\color{var}r}}}}$ is empty} \\
     1 & \mbox{if $L_{\ensuremath{\mathtt{{\color{var}r}}}}$ is non-empty}
     \end{array}\right.
  \]
  The height, \ensuremath{\mathtt{{\color{var}h}}}, of the skiplist is then given by
  \[
       \ensuremath{\mathtt{{\color{var}h}}} = \sum_{r=1}^\infty I_{\ensuremath{\mathtt{{\color{var}r}}}} \enspace .
  \]
  Note that $I_{\ensuremath{\mathtt{{\color{var}r}}}}$ is never more than the length, $|L_{\ensuremath{\mathtt{{\color{var}r}}}}|$, of $L_{\ensuremath{\mathtt{{\color{var}r}}}}$, so 
  \[
     \E[I_{\ensuremath{\mathtt{{\color{var}r}}}}] \le \E[|L_{\ensuremath{\mathtt{{\color{var}r}}}}|] = \ensuremath{\mathtt{{\color{var}n}}}/2^{\ensuremath{\mathtt{{\color{var}r}}}} \enspace .
  \]
  Therefore, we have
  \begin{align*}
       \E[\ensuremath{\mathtt{{\color{var}h}}}] &= \E\left[\sum_{r=1}^\infty I_{\ensuremath{\mathtt{{\color{var}r}}}}\right] \\
        &= \sum_{\ensuremath{\mathtt{{\color{var}r}}}=1}^{\infty} E[I_{\ensuremath{\mathtt{{\color{var}r}}}}] \\
        &= \sum_{\ensuremath{\mathtt{{\color{var}r}}}=1}^{\lfloor\log \ensuremath{\mathtt{{\color{var}n}}}\rfloor} E[I_{\ensuremath{\mathtt{{\color{var}r}}}}]
                 + \sum_{r=\lfloor\log \ensuremath{\mathtt{{\color{var}n}}}\rfloor+1}^{\infty} E[I_{\ensuremath{\mathtt{{\color{var}r}}}}]  \\
        &\le \sum_{\ensuremath{\mathtt{{\color{var}r}}}=1}^{\lfloor\log \ensuremath{\mathtt{{\color{var}n}}}\rfloor} 1
                 + \sum_{r=\lfloor\log \ensuremath{\mathtt{{\color{var}n}}}\rfloor+1}^{\infty} \ensuremath{\mathtt{{\color{var}n}}}/2^{\ensuremath{\mathtt{{\color{var}r}}}} \\
        &\le \log \ensuremath{\mathtt{{\color{var}n}}}
                 + \sum_{\ensuremath{\mathtt{{\color{var}r}}}=0}^\infty 1/2^{\ensuremath{\mathtt{{\color{var}r}}}} \\
        &= \log \ensuremath{\mathtt{{\color{var}n}}} + 2 \enspace . \qedhere
  \end{align*}
\end{proof}

\begin{lem}\lemlabel{skiplist-size2}
  The expected number of nodes in a skiplist containing $\ensuremath{\mathtt{{\color{var}n}}}$ elements,
  including all occurrences of the sentinel, is $2\ensuremath{\mathtt{{\color{var}n}}}+O(\log \ensuremath{\mathtt{{\color{var}n}}})$.
\end{lem}

\begin{proof}
  By \lemref{skiplist-size1}, the expected number of nodes, not
  including the sentinel, is $2\ensuremath{\mathtt{{\color{var}n}}}$.  The number of occurrences of
  the sentinel is equal to the height, $\ensuremath{\mathtt{{\color{var}h}}}$, of the skiplist so, by
  \lemref{skiplist-height} the expected number of occurrences of the
  sentinel is at most $\log \ensuremath{\mathtt{{\color{var}n}}}+2 = O(\log \ensuremath{\mathtt{{\color{var}n}}})$.
\end{proof}



\begin{lem}
The expected length of a search path in a skiplist is at most $2\log \ensuremath{\mathtt{{\color{var}n}}} + O(1)$.
\end{lem}

\begin{proof}
  The easiest way to see this is to consider the \emph{reverse search
  path} for a node, \ensuremath{\mathtt{{\color{var}x}}}.  This path starts at the predecessor of \ensuremath{\mathtt{{\color{var}x}}}
  in $L_0$.  At any point in time, if the path can go up a level, then
  it does.  If it cannot go up a level then it goes left.  Thinking about
  this for a few moments will convince us that the reverse search path for
  \ensuremath{\mathtt{{\color{var}x}}} is identical to the search path for \ensuremath{\mathtt{{\color{var}x}}}, except that it is reversed.

  The number of nodes that the reverse search path visits at a particular
  level, \ensuremath{\mathtt{{\color{var}r}}}, is related to the following experiment:  Toss a coin.
  If the coin comes up as heads, then move up and stop. Otherwise, move
  left and repeat the experiment.  The number of coin tosses before
  the heads represents the number of steps to the left that a reverse
  search path takes at a particular level.\footnote{Note that this
  might overcount the number of steps to the left, since the experiment
  should end either at the first heads or when the search path reaches
  the sentinel, whichever comes first. This is not a problem since the
  lemma is only stating an upper bound.} \lemref{coin-tosses} tells us
  that the expected number of coin tosses before the first heads is 1.

  Let $S_{\ensuremath{\mathtt{{\color{var}r}}}}$ denote the number of steps the forward search path takes at level
  $\ensuremath{\mathtt{{\color{var}r}}}$ that go to the right.   We have just argued that $\E[S_{\ensuremath{\mathtt{{\color{var}r}}}}]\le
  1$.  Furthermore, $S_{\ensuremath{\mathtt{{\color{var}r}}}}\le |L_{\ensuremath{\mathtt{{\color{var}r}}}}|$, since we can't take more steps
  in $L_{\ensuremath{\mathtt{{\color{var}r}}}}$ than the length of $L_{\ensuremath{\mathtt{{\color{var}r}}}}$, so
  \[
    \E[S_{\ensuremath{\mathtt{{\color{var}r}}}}] \le \E[|L_{\ensuremath{\mathtt{{\color{var}r}}}}|] = \ensuremath{\mathtt{{\color{var}n}}}/2^{\ensuremath{\mathtt{{\color{var}r}}}} \enspace .
  \]
  We can now finish as in the proof of \lemref{skiplist-height}.
  Let $S$ be  the length of the search path for some node, \ensuremath{\mathtt{{\color{var}u}}}, in a
  skiplist, and let $\ensuremath{\mathtt{{\color{var}h}}}$ be the height of the skiplist.  Then
  \begin{align*}
      \E[S] 
         &= \E\left[ \ensuremath{\mathtt{{\color{var}h}}} + \sum_{\ensuremath{\mathtt{{\color{var}r}}}=0}^\infty S_{\ensuremath{\mathtt{{\color{var}r}}}} \right] \\
         &= \E[\ensuremath{\mathtt{{\color{var}h}}}] + \sum_{\ensuremath{\mathtt{{\color{var}r}}}=0}^\infty \E[S_{\ensuremath{\mathtt{{\color{var}r}}}}]  \\
         &= \E[\ensuremath{\mathtt{{\color{var}h}}}] + \sum_{\ensuremath{\mathtt{{\color{var}r}}}=0}^{\lfloor\log \ensuremath{\mathtt{{\color{var}n}}}\rfloor} \E[S_{\ensuremath{\mathtt{{\color{var}r}}}}] 
              + \sum_{\ensuremath{\mathtt{{\color{var}r}}}=\lfloor\log \ensuremath{\mathtt{{\color{var}n}}}\rfloor+1}^\infty \E[S_{\ensuremath{\mathtt{{\color{var}r}}}}] \\
         &\le \E[\ensuremath{\mathtt{{\color{var}h}}}] + \sum_{\ensuremath{\mathtt{{\color{var}r}}}=0}^{\lfloor\log \ensuremath{\mathtt{{\color{var}n}}}\rfloor} 1
              + \sum_{r=\lfloor\log \ensuremath{\mathtt{{\color{var}n}}}\rfloor+1}^\infty \ensuremath{\mathtt{{\color{var}n}}}/2^{\ensuremath{\mathtt{{\color{var}r}}}} \\
         &\le \E[\ensuremath{\mathtt{{\color{var}h}}}] + \sum_{\ensuremath{\mathtt{{\color{var}r}}}=0}^{\lfloor\log \ensuremath{\mathtt{{\color{var}n}}}\rfloor} 1
              + \sum_{\ensuremath{\mathtt{{\color{var}r}}}=0}^{\infty} 1/2^{\ensuremath{\mathtt{{\color{var}r}}}} \\
         &\le \E[\ensuremath{\mathtt{{\color{var}h}}}] + \sum_{\ensuremath{\mathtt{{\color{var}r}}}=0}^{\lfloor\log \ensuremath{\mathtt{{\color{var}n}}}\rfloor} 1
              + \sum_{\ensuremath{\mathtt{{\color{var}r}}}=0}^{\infty} 1/2^{\ensuremath{\mathtt{{\color{var}r}}}} \\
         &\le \E[\ensuremath{\mathtt{{\color{var}h}}}] + \log \ensuremath{\mathtt{{\color{var}n}}} + 3 \\
         &\le 2\log \ensuremath{\mathtt{{\color{var}n}}} + 5  \enspace . \qedhere
  \end{align*}
\end{proof}


The following theorem summarizes the results in this section:
\begin{thm}
A skiplist containing $\ensuremath{\mathtt{{\color{var}n}}}$ elements has expected size $O(\ensuremath{\mathtt{{\color{var}n}}})$ and the
expected length of the search path for any particular element is at most
$2\log \ensuremath{\mathtt{{\color{var}n}}} + O(1)$.
\end{thm}





\section{Discussion and Exercises}

Skiplists were introduced by Pugh \cite{p91} who also presented
a number of applications and extensions of skiplists \cite{p89}.  Since then they
have been studied extensively.  Several researchers have done very
precise analyses of the expected length and variance of the length of the
search path for the \ensuremath{\mathtt{{\color{var}i}}}th element in a skiplist \cite{kp94,kmp95,pmp92}.
Deterministic versions \cite{mps92}, biased versions \cite{bbg02,esss01},
and self-adjusting versions \cite{bdl08} of skiplists have all been
developed.  Skiplist implementations have been written for various
languages and frameworks and have been used in open-source database
systems \cite{skipdb,redis}. A variant of skiplists is used in the HP-UX
operating system kernel's process management structures \cite{hpux}.
\javaonly{Skiplists are even part of the Java 1.6 API \cite{oracle_jdk6}.}


\begin{exc}
  Illustrate the search paths for 2.5 and 5.5 on the skiplist in 
  \figref{skiplist}.
\end{exc}

\begin{exc}
  Illustrate the addition of the values 0.5 (with a height of 1) and
  then 3.5 (with a height of 2) to the skiplist in \figref{skiplist}.
\end{exc}

\begin{exc}
  Illustrate the removal of the values 1 and then 3  from the skiplist
  in \figref{skiplist}.
\end{exc}

\begin{exc}
  Illustrate the execution of \ensuremath{\mathtt{remove(2)}} on the \texttt{SkiplistList}
  in \figref{skiplist-lengths}.
\end{exc}

\begin{exc}
  Illustrate the execution of \ensuremath{\mathtt{add(3,{\color{var}x})}} on the \texttt{SkiplistList}
  in \figref{skiplist-lengths}.  Assume that \ensuremath{\mathtt{pickHeight()}} selects a height
  of 4 for the newly created node.
\end{exc}

\begin{exc}\exclabel{skiplist-changes}
  Show that, during an \ensuremath{\mathtt{add({\color{var}x})}} or a \ensuremath{\mathtt{remove({\color{var}x})}} operation, the expected
  number of pointers in a \texttt{SkiplistSet} that get changed is constant.
\end{exc}

\begin{exc}\exclabel{skiplist-opt}
  Suppose that, instead of promoting an element from $L_{i-1}$ into $L_i$
  based on a coin toss, we promote it with some probability $p$, $0 <
  p < 1$.
  \begin{enumerate}
   \item Show that, with this modification, the expected length of a
     search path is at most $(1/p)\log_{1/p} \ensuremath{\mathtt{{\color{var}n}}} + O(1)$.
   \item What is the value of $p$ that minimizes the preceding expression?
   \item What is the expected height of the skiplist? 
   \item What is the expected number of nodes in the skiplist?
  \end{enumerate}
\end{exc}


\begin{exc}\exclabel{skiplist-opt-2}
  The \ensuremath{\mathtt{find({\color{var}x})}} method in a \texttt{SkiplistSet} sometimes performs
  \emph{redundant comparisons}; these occur when \ensuremath{\mathtt{{\color{var}x}}} is compared to the
  same value more than once.  They can occur when, for some node, \ensuremath{\mathtt{{\color{var}u}}},
  $\ensuremath{\mathtt{{\color{var}u}.{\color{var}next}[{\color{var}r}]}} = \ensuremath{\mathtt{{\color{var}u}.{\color{var}next}[{\color{var}r}-1]}}$.  Show how these redundant comparisons
  happen and modify \ensuremath{\mathtt{find({\color{var}x})}} so that they are avoided.  Analyze the
  expected number of comparisons done by your modified \ensuremath{\mathtt{find({\color{var}x})}} method.
\end{exc}

\begin{exc}
  Design and implement a version of a skiplist that implements the
  \texttt{SSet} interface, but also allows fast access to elements by rank.
  That is, it also supports the function \ensuremath{\mathtt{get({\color{var}i})}}, which returns the
  element whose rank is \ensuremath{\mathtt{{\color{var}i}}} in $O(\log \ensuremath{\mathtt{{\color{var}n}}})$ expected time. (The rank
  of an element \ensuremath{\mathtt{{\color{var}x}}} in an \texttt{SSet} is the number of elements in the \texttt{SSet}
  that are less than \ensuremath{\mathtt{{\color{var}x}}}.)
\end{exc}

\begin{exc}
  \index{finger}%
  \index{finger search!in a skiplist}%
  A \emph{finger} in a skiplist is an array that stores the
  sequence of nodes on a search path at which the search path
  goes down. (The variable \ensuremath{\mathtt{{\color{var}stack}}} in the \ensuremath{\mathtt{add({\color{var}x})}} code on
  page~\pageref{pg:skiplist-add} is a finger;  the shaded nodes in
  \figref{skiplist-add} show the contents of the finger.)  One can
  think of a finger as pointing out the path to a node in the lowest
  list, $L_0$.

  A \emph{finger search} implements the \ensuremath{\mathtt{find({\color{var}x})}} operation using a
  finger, by walking up the list using the finger until reaching a node
  \ensuremath{\mathtt{{\color{var}u}}} such that $\ensuremath{\mathtt{{\color{var}u}.{\color{var}x}}} < \ensuremath{\mathtt{{\color{var}x}}}$ and $\ensuremath{\mathtt{{\color{var}u}.{\color{var}next}}}=\ensuremath{\mathtt{{\color{keyword}null}}}$ or $\ensuremath{\mathtt{{\color{var}u}.{\color{var}next}.{\color{var}x}}} >
  \ensuremath{\mathtt{{\color{var}x}}}$ and then performing a normal search for \ensuremath{\mathtt{{\color{var}x}}} starting from \ensuremath{\mathtt{{\color{var}u}}}.
  It is possible to prove that the expected number of steps required
  for a finger search is $O(1+\log r)$, where $r$ is the number values
  in $L_0$ between \ensuremath{\mathtt{{\color{var}x}}} and the value pointed to by the finger.

  Implement a subclass of \texttt{Skiplist} called \texttt{SkiplistWithFinger} that
  implements \ensuremath{\mathtt{find({\color{var}x})}} operations using an internal finger.  This subclass
  stores a finger, which is then used so that every \ensuremath{\mathtt{find({\color{var}x})}} operation
  is implemented as a finger search.  During each \ensuremath{\mathtt{find({\color{var}x})}} operation
  the finger is updated so that each \ensuremath{\mathtt{find({\color{var}x})}} operation uses, as a
  starting point, a finger that points to the result of the previous
  \ensuremath{\mathtt{find({\color{var}x})}} operation.
\end{exc}

\begin{exc}\exclabel{skiplist-truncate}
  Write a method, \ensuremath{\mathtt{truncate({\color{var}i})}}, that truncates a \texttt{SkiplistList}
  at position \ensuremath{\mathtt{{\color{var}i}}}.  After the execution of this method, the size
  of the list is \ensuremath{\mathtt{{\color{var}i}}} and it contains only the elements at indices
  $0,\ldots,\ensuremath{\mathtt{{\color{var}i}}}-1$.  The return value is another \texttt{SkiplistList} that
  contains the elements at indices $\ensuremath{\mathtt{{\color{var}i}}},\ldots,\ensuremath{\mathtt{{\color{var}n}}}-1$.  This method
  should run in $O(\log \ensuremath{\mathtt{{\color{var}n}}})$ time.
\end{exc}

\begin{exc}
  Write a \texttt{SkiplistList} method, \ensuremath{\mathtt{absorb({\color{var}l2})}}, that takes as an
  argument a \texttt{SkiplistList}, \ensuremath{\mathtt{{\color{var}l2}}}, empties it and appends its contents,
  in order, to the receiver.  For example, if \ensuremath{\mathtt{{\color{var}l1}}} contains $a,b,c$
  and \ensuremath{\mathtt{{\color{var}l2}}} contains $d,e,f$, then after calling \ensuremath{\mathtt{{\color{var}l1}.absorb({\color{var}l2})}}, \ensuremath{\mathtt{{\color{var}l1}}}
  will contain $a,b,c,d,e,f$ and \ensuremath{\mathtt{{\color{var}l2}}} will be empty. This method should
  run in $O(\log \ensuremath{\mathtt{{\color{var}n}}})$ time.
\end{exc}

\begin{exc}
  Using the ideas from the space-efficient list, \texttt{SEList}, design
  and implement a space-efficient \texttt{SSet}, \texttt{SESSet}.  To do this, store the
  data, in order, in an \texttt{SEList}, and store the blocks of this \texttt{SEList}
  in an \texttt{SSet}. If the original \texttt{SSet} implementation uses $O(\ensuremath{\mathtt{{\color{var}n}}})$
  space to store \ensuremath{\mathtt{{\color{var}n}}} elements, then the \texttt{SESSet} will use enough space
  for \ensuremath{\mathtt{{\color{var}n}}} elements plus $O(\ensuremath{\mathtt{{\color{var}n}}}/\ensuremath{\mathtt{{\color{var}b}}}+\ensuremath{\mathtt{{\color{var}b}}})$ wasted space.
\end{exc}

\begin{exc}
  Using an \texttt{SSet} as your underlying structure, design and implement an
  application that reads a (large) text file and allows you to search,
  interactively, for any substring contained in the text.  As the user
  types their query, a matching part of the text (if any) should appear 
  as a result.

  \noindent  Hint 1: Every substring is a prefix of some suffix, so it
  suffices to store all suffixes of the text file.

  \noindent Hint 2:  Any suffix can be represented compactly as a single
  integer indicating where the suffix begins in the text.

  \noindent Test your application on some large texts, such as some of
  the books available at Project Gutenberg \cite{gutenberg}.  If done
  correctly, your applications will be very responsive; there should be
  no noticeable lag between typing keystrokes and seeing the results.
\end{exc}

\begin{exc}
  \index{skiplist!versus binary search tree}%
  \index{binary search tree!versus skiplist}%
  (This exercise should be done after reading about binary search trees,
  in \secref{binarysearchtree}.)  Compare skiplists with binary search
  trees in the following ways:  
  \begin{enumerate}
     \item Explain how removing some edges of a skiplist leads to
       a structure that looks like a binary tree and is similar to a
       binary search tree.
     \item Skiplists and binary search trees each use about the same
       number of pointers (2 per node).  Skiplists make better use of
       those pointers, though. Explain why.
  \end{enumerate}
\end{exc}

\clearpage

\clearpage
\chapter{Hash Tables}
\chaplabel{hashtables}
\chaplabel{hashing}

Hash tables are an efficient method of storing a small number,
\ensuremath{\mathtt{{\color{var}n}}}, of integers from a large range $U=\{0,\ldots,2^{\ensuremath{\mathtt{{\color{var}w}}}}-1\}$.
The term \emph{hash table}
\index{hash table}%
includes a broad range of data structures.  The first part of this
chapter focuses on two of the most common implementations of hash tables:
hashing with chaining and linear probing.

Very often hash tables store types of data that are not integers.
In this case, an integer \emph{hash code}
\index{hash code}%
is associated with each data
item and is used in the hash table.  The second part of this chapter
discusses how such hash codes are generated.

Some of the methods used in this chapter require random choices of
integers in some specific range.  In the code samples, some of these
``random'' integers are hard-coded constants.  These constants were
obtained using random bits generated from atmospheric noise.


\section{\texttt{ChainedHashTable}: Hashing with Chaining}
\seclabel{hashtable}

\index{ChainedHashTable@\texttt{ChainedHashTable}}%
\index{chaining}%
\index{hashing with chaining}%
A \texttt{ChainedHashTable} data structure uses \emph{hashing with chaining} to store
data as an array, \ensuremath{\mathtt{{\color{var}t}}}, of lists.  An integer, \ensuremath{\mathtt{{\color{var}n}}}, keeps track of the
total number of items in all lists (see \figref{chainedhashtable}):


 \ \vspace{.5ex}\begin{Verbatim}[gobble=1,tabsize=2,frame=single,framerule=1.2pt,commandchars=\\@\$,label=\texttt{ChainedHashTable},labelposition=topline]
	List<@\color@keyword$T$>[] @\color@var$t$;
	@\color@keyword$int$ @\color@var$n$;
\end{Verbatim}

\begin{figure}
   \begin{center}
     \includegraphics[width=\ScaleIfNeeded]{figs/chainedhashtable}
   \end{center}
   \caption[A ChainedHashTable]{An example of a \texttt{ChainedHashTable} with $\ensuremath{\mathtt{{\color{var}n}}}=14$ and $\ensuremath{\mathtt{{\color{var}t}.{\color{var}length}}}=16$.  In this example $\ensuremath{\mathtt{hash({\color{var}x})}}=6$}
   \figlabel{chainedhashtable}
\end{figure}
\index{hash value}%
\index{hash(x)@\ensuremath{\mathtt{hash({\color{var}x})}}}%
The \emph{hash value} of a data item \ensuremath{\mathtt{{\color{var}x}}}, denoted \ensuremath{\mathtt{hash({\color{var}x})}} is a value
in the range $\{0,\ldots,\ensuremath{\mathtt{{\color{var}t}.{\color{var}length}}}-1\}$.  All items with hash value \ensuremath{\mathtt{{\color{var}i}}}
are stored in the list at \ensuremath{\mathtt{{\color{var}t}[{\color{var}i}]}}.  To ensure that lists don't get too
long, we maintain the invariant
\[
    \ensuremath{\mathtt{{\color{var}n}}} \le \ensuremath{\mathtt{{\color{var}t}.{\color{var}length}}}
\]
so that the average number of elements stored in one of these lists is 
$\ensuremath{\mathtt{{\color{var}n}}}/\ensuremath{\mathtt{{\color{var}t}.{\color{var}length}}} \le 1$.

To add an element, \ensuremath{\mathtt{{\color{var}x}}}, to the hash table, we first check if
the length of \ensuremath{\mathtt{{\color{var}t}}} needs to be increased and, if so, we grow \ensuremath{\mathtt{{\color{var}t}}}.
With this out of the way we hash \ensuremath{\mathtt{{\color{var}x}}} to get an integer, \ensuremath{\mathtt{{\color{var}i}}}, in the
range $\{0,\ldots,\ensuremath{\mathtt{{\color{var}t}.{\color{var}length}}}-1\}$, and we append \ensuremath{\mathtt{{\color{var}x}}} to the list
\ensuremath{\mathtt{{\color{var}t}[{\color{var}i}]}}:


 \ \vspace{.5ex}\begin{Verbatim}[gobble=1,tabsize=2,frame=single,framerule=1.2pt,commandchars=\\@\$,label=\texttt{ChainedHashTable},labelposition=topline]
	@\color@var$boolean$ add(@\color@keyword$T$ @\color@var$x$) {
		@\color@keyword$if$ (find(@\color@var$x$) != @\color@keyword$null$) @\color@var$return$ @\color@var$false$;
		@\color@keyword$if$ (@\color@var$n$+1 > @\color@var$t$.@\color@var$length$) resize();
		@\color@var$t$[hash(@\color@var$x$)].add(@\color@var$x$);
		@\color@var$n$++;
		@\color@var$return$ @\color@var$true$;
	}
\end{Verbatim}

Growing the table,
if necessary, involves doubling the length of \ensuremath{\mathtt{{\color{var}t}}} and reinserting
all elements into the new table.  This strategy is exactly the same
as the one used in the implementation of \texttt{ArrayStack} and the same
result applies: The cost of growing is only constant when amortized
over a sequence of insertions (see~\lemref{arraystack-amortized} on
page~\pageref{lem:arraystack-amortized}).

Besides growing, the only other work done when adding a new value \ensuremath{\mathtt{{\color{var}x}}} to a
\texttt{ChainedHashTable} involves appending \ensuremath{\mathtt{{\color{var}x}}} to the list \ensuremath{\mathtt{{\color{var}t}[hash({\color{var}x})]}}.  For
any of the list implementations described in Chapters~\ref{chap:arrays}
or \ref{chap:linkedlists}, this takes only constant time.

To remove an element, \ensuremath{\mathtt{{\color{var}x}}}, from the hash table, we iterate over the list
\ensuremath{\mathtt{{\color{var}t}[hash({\color{var}x})]}} until we find \ensuremath{\mathtt{{\color{var}x}}} so that we can remove it:


 \ \vspace{.5ex}\begin{Verbatim}[gobble=1,tabsize=2,frame=single,framerule=1.2pt,commandchars=\\@\$,label=\texttt{ChainedHashTable},labelposition=topline]
	@\color@keyword$T$ remove(@\color@keyword$T$ @\color@var$x$) {
		Iterator<@\color@keyword$T$> @\color@var$it$ = @\color@var$t$[hash(@\color@var$x$)].iterator();
		@\color@keyword$while$ (@\color@var$it$.hasNext()) {
			@\color@keyword$T$ @\color@var$y$ = @\color@var$it$.next();
			@\color@keyword$if$ (@\color@var$y$.equals(@\color@var$x$)) {
				@\color@var$it$.remove();
				@\color@var$n$--;
				@\color@var$return$ @\color@var$y$;
			}
		}
		@\color@var$return$ @\color@keyword$null$;
	}
\end{Verbatim}

This takes $O(\ensuremath{\mathtt{{\color{var}n}}}_{\ensuremath{\mathtt{hash({\color{var}x})}}})$ time, where $\ensuremath{\mathtt{{\color{var}n}}}_{\ensuremath{\mathtt{{\color{var}i}}}}$ denotes the length
of the list stored at \ensuremath{\mathtt{{\color{var}t}[{\color{var}i}]}}.

Searching for the element \ensuremath{\mathtt{{\color{var}x}}} in a hash table is similar.  We perform
a linear search on the list \ensuremath{\mathtt{{\color{var}t}[hash({\color{var}x})]}}:


 \ \vspace{.5ex}\begin{Verbatim}[gobble=1,tabsize=2,frame=single,framerule=1.2pt,commandchars=\\@\$,label=\texttt{ChainedHashTable},labelposition=topline]
	@\color@keyword$T$ find(Object @\color@var$x$) {
		@\color@keyword$for$ (@\color@keyword$T$ @\color@var$y$ : @\color@var$t$[hash(@\color@var$x$)])
			@\color@keyword$if$ (@\color@var$y$.equals(@\color@var$x$))
				@\color@var$return$ @\color@var$y$;
		@\color@var$return$ @\color@keyword$null$;
	}
\end{Verbatim}

Again, this takes time proportional to the length of the list \ensuremath{\mathtt{{\color{var}t}[hash({\color{var}x})]}}.

The performance of a hash table depends critically on the choice of the
hash function.  A good hash function will spread the elements evenly
among the \ensuremath{\mathtt{{\color{var}t}.{\color{var}length}}} lists, so that the expected size of the list
\ensuremath{\mathtt{{\color{var}t}[hash({\color{var}x})]}} is $O(\ensuremath{\mathtt{{\color{var}n}}}/\ensuremath{\mathtt{{\color{var}t}.{\color{var}length})}} = O(1)$.  On the other hand, a bad
hash function will hash all values (including \ensuremath{\mathtt{{\color{var}x}}}) to the same table
location, in which case the size of the list \ensuremath{\mathtt{{\color{var}t}[hash({\color{var}x})]}} will be \ensuremath{\mathtt{{\color{var}n}}}.
In the next section we describe a good hash function.

\subsection{Multiplicative Hashing}
\seclabel{multihash}

\index{hashing!multiplicative}%
\index{multiplicative hashing}%
Multiplicative hashing is an efficient method of generating hash
values based on modular arithmetic (discussed in \secref{arrayqueue})
and integer division.  It uses the $\ddiv$ operator, which calculates
the integral part of a quotient, while discarding the remainder.
Formally, for any integers $a\ge 0$ and $b\ge 1$, $a\ddiv b = \lfloor
a/b\rfloor$.

In multiplicative hashing, we use a hash table of size $2^{\ensuremath{\mathtt{{\color{var}d}}}}$ for some
integer \ensuremath{\mathtt{{\color{var}d}}} (called the \emph{dimension}).  The formula for hashing an
integer $\ensuremath{\mathtt{{\color{var}x}}}\in\{0,\ldots,2^{\ensuremath{\mathtt{{\color{var}w}}}}-1\}$ is
\[
    \ensuremath{\mathtt{hash({\color{var}x})}} = ((\ensuremath{\mathtt{{\color{var}z}}}\cdot\ensuremath{\mathtt{{\color{var}x}}}) \bmod 2^{\ensuremath{\mathtt{{\color{var}w}}}}) \ddiv 2^{\ensuremath{\mathtt{{\color{var}w}}}-\ensuremath{\mathtt{{\color{var}d}}}} \enspace .
\]
Here, \ensuremath{\mathtt{{\color{var}z}}} is a randomly chosen \emph{odd} integer in
$\{1,\ldots,2^{\ensuremath{\mathtt{{\color{var}w}}}}-1\}$.  This hash function can be realized very
efficiently by observing that, by default, operations on integers
are already done modulo $2^{\ensuremath{\mathtt{{\color{var}w}}}}$ where $\ensuremath{\mathtt{{\color{var}w}}}$ is the number of bits
in an integer.\footnote{This is true for most programming languages
including C, C\#, C++, and Java. Notable exceptions are Python and
Ruby, in which the result of a fixed-length \ensuremath{\mathtt{{\color{var}w}}}-bit integer operation
that overflows is upgraded to a variable-length representation.}  (See
\figref{multihashing}.) Furthermore, integer division by $2^{\ensuremath{\mathtt{{\color{var}w}}}-\ensuremath{\mathtt{{\color{var}d}}}}$
is equivalent to dropping the rightmost $\ensuremath{\mathtt{{\color{var}w}}}-\ensuremath{\mathtt{{\color{var}d}}}$ bits in a binary
representation (which is implemented by shifting the bits right by
$\ensuremath{\mathtt{{\color{var}w}}}-\ensuremath{\mathtt{{\color{var}d}}}$ using the \javaonly{\ensuremath{\mathtt{\text{\ttfamily >>>}}}}\cpponly{\ensuremath{\mathtt{\text{\ttfamily >>}}}}\pcodeonly{\ensuremath{\mathtt{\text{\ttfamily >>}}}}
operator).  \notpcode{In this way, the code that implements the above
formula is simpler than the formula itself:}


 \ \vspace{.5ex}\begin{Verbatim}[gobble=1,tabsize=2,frame=single,framerule=1.2pt,commandchars=\\@\$,label=\texttt{ChainedHashTable},labelposition=topline]
	@\color@keyword$int$ hash(Object @\color@var$x$) {
		return (@\color@var$z$ * @\color@var$x$.hashCode()) >>> (@\color@var$w$-@\color@var$d$);
	}
\end{Verbatim}


\begin{figure}
  \begin{center}
    \resizebox{.98\textwidth}{!}{
    \setlength{\arrayrulewidth}{1pt}
    \begin{tabular}{|lr@{}r|}\hline
    $2^\ensuremath{\mathtt{{\color{var}w}}}$ (4294967296)&            \ensuremath{\mathtt{1}}&\ensuremath{\mathtt{00000000000000000000000000000000}} \\
    \ensuremath{\mathtt{{\color{var}z}}} (4102541685)&                   &\ensuremath{\mathtt{11110100100001111101000101110101}} \\
    \ensuremath{\mathtt{{\color{var}x}}} (42) &                          &\ensuremath{\mathtt{00000000000000000000000000101010}} \\
    $\ensuremath{\mathtt{{\color{var}z}}}\cdot\ensuremath{\mathtt{{\color{var}x}}}$ &             \ensuremath{\mathtt{101000}}&\ensuremath{\mathtt{00011110010010000101110100110010}} \\
    $(\ensuremath{\mathtt{{\color{var}z}}}\cdot\ensuremath{\mathtt{{\color{var}x}}})\bmod 2^{\ensuremath{\mathtt{{\color{var}w}}}}$ &      &\ensuremath{\mathtt{00011110010010000101110100110010}} \\
    $((\ensuremath{\mathtt{{\color{var}z}}}\cdot\ensuremath{\mathtt{{\color{var}x}}})\bmod 2^{\ensuremath{\mathtt{{\color{var}w}}}})\ddiv 2^{\ensuremath{\mathtt{{\color{var}w}}}-\ensuremath{\mathtt{{\color{var}d}}}}$ &&
                      \multicolumn{1}{@{}l|}{\ensuremath{\mathtt{00011110}}} \\\hline
    \end{tabular}}
    \setlength{\arrayrulewidth}{.4pt}
  \end{center}
  \caption{The operation of the multiplicative hash function with $\ensuremath{\mathtt{{\color{var}w}}}=32$
    and $\ensuremath{\mathtt{{\color{var}d}}}=8$.}
  \figlabel{multihashing}
\end{figure}

The following lemma, whose proof is deferred until later in this section,
shows that multiplicative hashing does a good job of avoiding collisions:

\begin{lem}\lemlabel{universal-hashing}
  Let \ensuremath{\mathtt{{\color{var}x}}} and \ensuremath{\mathtt{{\color{var}y}}} be any two values in $\{0,\ldots,2^{\ensuremath{\mathtt{{\color{var}w}}}}-1\}$ with
  $\ensuremath{\mathtt{{\color{var}x}}}\neq \ensuremath{\mathtt{{\color{var}y}}}$. Then $\Pr\{\ensuremath{\mathtt{hash({\color{var}x})}}=\ensuremath{\mathtt{hash({\color{var}y})}}\} \le 2/2^{\ensuremath{\mathtt{{\color{var}d}}}}$.
\end{lem}

With \lemref{universal-hashing}, the performance of \ensuremath{\mathtt{remove({\color{var}x})}}, and
\ensuremath{\mathtt{find({\color{var}x})}} are easy to analyze:

\begin{lem}
  For any data value \ensuremath{\mathtt{{\color{var}x}}}, the expected length of the list \ensuremath{\mathtt{{\color{var}t}[hash({\color{var}x})]}}
  is at most $\ensuremath{\mathtt{{\color{var}n}}}_{\ensuremath{\mathtt{{\color{var}x}}}} + 2$, where $\ensuremath{\mathtt{{\color{var}n}}}_{\ensuremath{\mathtt{{\color{var}x}}}}$ is the number of
  occurrences of \ensuremath{\mathtt{{\color{var}x}}} in the hash table.
\end{lem}

\begin{proof}
  Let $S$ be the (multi-)set of elements stored in the hash table that
  are not equal to \ensuremath{\mathtt{{\color{var}x}}}.  For an element $\ensuremath{\mathtt{{\color{var}y}}}\in S$, define the indicator
  variable
    \[ I_{\ensuremath{\mathtt{{\color{var}y}}}} = \left\{\begin{array}{ll}
       1 & \mbox{if $\ensuremath{\mathtt{hash({\color{var}x})}}=\ensuremath{\mathtt{hash({\color{var}y})}}$} \\
       0 & \mbox{otherwise}
       \end{array}\right.
    \]
  and notice that, by \lemref{universal-hashing}, $\E[I_{\ensuremath{\mathtt{{\color{var}y}}}}] \le
  2/2^{\ensuremath{\mathtt{{\color{var}d}}}}=2/\ensuremath{\mathtt{{\color{var}t}.{\color{var}length}}}$.  The expected length of the list \ensuremath{\mathtt{{\color{var}t}[hash({\color{var}x})]}}
  is given by
  \begin{eqnarray*}
   \E\left[\ensuremath{\mathtt{{\color{var}t}[hash({\color{var}x})].size()}}\right] &=& \E\left[\ensuremath{\mathtt{{\color{var}n}}}_{\ensuremath{\mathtt{{\color{var}x}}}} + \sum_{\ensuremath{\mathtt{{\color{var}y}}}\in S} I_{\ensuremath{\mathtt{{\color{var}y}}}}\right] \\
    &=& \ensuremath{\mathtt{{\color{var}n}}}_{\ensuremath{\mathtt{{\color{var}x}}}} + \sum_{\ensuremath{\mathtt{{\color{var}y}}}\in S} \E [I_{\ensuremath{\mathtt{{\color{var}y}}}} ] \\
    &\le& \ensuremath{\mathtt{{\color{var}n}}}_{\ensuremath{\mathtt{{\color{var}x}}}} + \sum_{\ensuremath{\mathtt{{\color{var}y}}}\in S} 2/\ensuremath{\mathtt{{\color{var}t}.{\color{var}length}}} \\
    &\le& \ensuremath{\mathtt{{\color{var}n}}}_{\ensuremath{\mathtt{{\color{var}x}}}} + \sum_{\ensuremath{\mathtt{{\color{var}y}}}\in S} 2/\ensuremath{\mathtt{{\color{var}n}}} \\
    &\le& \ensuremath{\mathtt{{\color{var}n}}}_{\ensuremath{\mathtt{{\color{var}x}}}} + (\ensuremath{\mathtt{{\color{var}n}}}-\ensuremath{\mathtt{{\color{var}n}}}_{\ensuremath{\mathtt{{\color{var}x}}}})2/\ensuremath{\mathtt{{\color{var}n}}} \\
    &\le& \ensuremath{\mathtt{{\color{var}n}}}_{\ensuremath{\mathtt{{\color{var}x}}}} + 2 \enspace ,
  \end{eqnarray*}
  as required.
\end{proof}

Now, we want to prove \lemref{universal-hashing}, but first we need a
result from number theory.  In the following proof, we use the notation
$(b_r,\ldots,b_0)_2$ to denote $\sum_{i=0}^r b_i2^i$, where each $b_i$
is a bit, either 0 or 1.  In other words, $(b_r,\ldots,b_0)_2$ is
the integer whose binary representation is given by $b_r,\ldots,b_0$.
We use $\star$ to denote a bit of unknown value.

\begin{lem}\lemlabel{hashing-mapping}
  Let $S$ be the set of odd integers in $\{1,\ldots,2^{\ensuremath{\mathtt{{\color{var}w}}}}-1\}$; let $q$
  and $i$ be any two elements in $S$.  Then there is exactly one value
  $\ensuremath{\mathtt{{\color{var}z}}}\in S$ such that $\ensuremath{\mathtt{{\color{var}z}}}q\bmod 2^{\ensuremath{\mathtt{{\color{var}w}}}} = i$.
\end{lem}

\begin{proof}
  Since the number of choices for $\ensuremath{\mathtt{{\color{var}z}}}$ and $i$ is the same, it is
  sufficient to prove that there is \emph{at most} one value $\ensuremath{\mathtt{{\color{var}z}}}\in S$ that
  satisfies $\ensuremath{\mathtt{{\color{var}z}}}q\bmod 2^{\ensuremath{\mathtt{{\color{var}w}}}} = i$.

  Suppose, for the sake of contradiction, that there are two such values
  \ensuremath{\mathtt{{\color{var}z}}} and \ensuremath{\mathtt{{\color{var}z}'}}, with $\ensuremath{\mathtt{{\color{var}z}}}>\ensuremath{\mathtt{{\color{var}z}}}'$.  Then
  \[
     \ensuremath{\mathtt{{\color{var}z}}}q\bmod 2^{\ensuremath{\mathtt{{\color{var}w}}}} = \ensuremath{\mathtt{{\color{var}z}}}'q \bmod 2^{\ensuremath{\mathtt{{\color{var}w}}}} = i
  \]
  So
  \[ 
     (\ensuremath{\mathtt{{\color{var}z}}}-\ensuremath{\mathtt{{\color{var}z}}}')q\bmod 2^{\ensuremath{\mathtt{{\color{var}w}}}} = 0 
  \]
  But this means that 
  \begin{equation}
    (\ensuremath{\mathtt{{\color{var}z}}}-\ensuremath{\mathtt{{\color{var}z}}}')q = k 2^{\ensuremath{\mathtt{{\color{var}w}}}} \eqlabel{factors} 
  \end{equation}
  for some integer $k$.  Thinking in terms of binary numbers, we have 
  \[
    (\ensuremath{\mathtt{{\color{var}z}}}-\ensuremath{\mathtt{{\color{var}z}}}')q = k\cdot(1,\underbrace{0,\ldots,0}_{\ensuremath{\mathtt{{\color{var}w}}}})_2 \enspace ,
  \]
  so that the \ensuremath{\mathtt{{\color{var}w}}} trailing bits in the binary representation of
  $(\ensuremath{\mathtt{{\color{var}z}}}-\ensuremath{\mathtt{{\color{var}z}}}')q$ are all 0's.

  Furthermore $k\neq 0$, since $q\neq 0$ and $\ensuremath{\mathtt{{\color{var}z}}}-\ensuremath{\mathtt{{\color{var}z}}}'\neq 0$.  Since $q$
  is odd, it has no trailing 0's in its binary representation:
  \[
    q = (\star,\ldots,\star,1)_2 \enspace .
  \]
  Since $|\ensuremath{\mathtt{{\color{var}z}}}-\ensuremath{\mathtt{{\color{var}z}}}'| < 2^{\ensuremath{\mathtt{{\color{var}w}}}}$, $\ensuremath{\mathtt{{\color{var}z}}}-\ensuremath{\mathtt{{\color{var}z}}}'$ has fewer than \ensuremath{\mathtt{{\color{var}w}}} trailing
  0's in its binary representation:
  \[
    \ensuremath{\mathtt{{\color{var}z}}}-\ensuremath{\mathtt{{\color{var}z}}}' = (\star,\ldots,\star,1,\underbrace{0,\ldots,0}_{<\ensuremath{\mathtt{{\color{var}w}}}})_2
      \enspace .
  \]
  Therefore, the product $(\ensuremath{\mathtt{{\color{var}z}}}-\ensuremath{\mathtt{{\color{var}z}}}')q$ has fewer than \ensuremath{\mathtt{{\color{var}w}}} trailing 0's in
  its binary representation:
  \[
   (\ensuremath{\mathtt{{\color{var}z}}}-\ensuremath{\mathtt{{\color{var}z}}}')q = (\star,\cdots,\star,1,\underbrace{0,\ldots,0}_{<\ensuremath{\mathtt{{\color{var}w}}}})_2 
    \enspace .
  \]
  Therefore $(\ensuremath{\mathtt{{\color{var}z}}}-\ensuremath{\mathtt{{\color{var}z}}}')q$ cannot satisfy \myeqref{factors}, yielding a
  contradiction and completing the proof.
\end{proof}

The utility of \lemref{hashing-mapping} comes from the following
observation:  If \ensuremath{\mathtt{{\color{var}z}}} is chosen uniformly at random from $S$, then \ensuremath{\mathtt{{\color{var}zt}}}
is uniformly distributed over $S$.  In the following proof, it helps
to think of the binary representation of \ensuremath{\mathtt{{\color{var}z}}}, which consists of $\ensuremath{\mathtt{{\color{var}w}}}-1$
random bits followed by a 1.

\begin{proof}[Proof of \lemref{universal-hashing}]
  First we note that the condition $\ensuremath{\mathtt{hash({\color{var}x})}}=\ensuremath{\mathtt{hash({\color{var}y})}}$ is equivalent to
  the statement ``the highest-order \ensuremath{\mathtt{{\color{var}d}}} bits of $\ensuremath{\mathtt{{\color{var}z}}} \ensuremath{\mathtt{{\color{var}x}}}\bmod2^{\ensuremath{\mathtt{{\color{var}w}}}}$
  and the highest-order \ensuremath{\mathtt{{\color{var}d}}} bits of $\ensuremath{\mathtt{{\color{var}z}}} \ensuremath{\mathtt{{\color{var}y}}}\bmod 2^{\ensuremath{\mathtt{{\color{var}w}}}}$ are the same.''
  A necessary condition of that statement is that the highest-order
  \ensuremath{\mathtt{{\color{var}d}}} bits in the binary representation of $\ensuremath{\mathtt{{\color{var}z}}}(\ensuremath{\mathtt{{\color{var}x}}}-\ensuremath{\mathtt{{\color{var}y}}})\bmod 2^{\ensuremath{\mathtt{{\color{var}w}}}}$
  are either all 0's or all 1's.  That is,
  \begin{equation}
      \ensuremath{\mathtt{{\color{var}z}}}(\ensuremath{\mathtt{{\color{var}x}}}-\ensuremath{\mathtt{{\color{var}y}}})\bmod 2^{\ensuremath{\mathtt{{\color{var}w}}}} = 
      (\underbrace{0,\ldots,0}_{\ensuremath{\mathtt{{\color{var}d}}}},\underbrace{\star,\ldots,\star}_{\ensuremath{\mathtt{{\color{var}w}}}-\ensuremath{\mathtt{{\color{var}d}}}})_2 
      \eqlabel{all-zeros}
  \end{equation}
  when $\ensuremath{\mathtt{{\color{var}zx}}}\bmod 2^{\ensuremath{\mathtt{{\color{var}w}}}} > \ensuremath{\mathtt{{\color{var}zy}}}\bmod 2^{\ensuremath{\mathtt{{\color{var}w}}}}$ or
  \begin{equation}
      \ensuremath{\mathtt{{\color{var}z}}}(\ensuremath{\mathtt{{\color{var}x}}}-\ensuremath{\mathtt{{\color{var}y}}})\bmod 2^{\ensuremath{\mathtt{{\color{var}w}}}} = 
      (\underbrace{1,\ldots,1}_{\ensuremath{\mathtt{{\color{var}d}}}},\underbrace{\star,\ldots,\star}_{\ensuremath{\mathtt{{\color{var}w}}}-\ensuremath{\mathtt{{\color{var}d}}}})_2 
       \enspace .
      \eqlabel{all-ones}
  \end{equation}
  when $\ensuremath{\mathtt{{\color{var}zx}}}\bmod 2^{\ensuremath{\mathtt{{\color{var}w}}}} < \ensuremath{\mathtt{{\color{var}zy}}}\bmod 2^{\ensuremath{\mathtt{{\color{var}w}}}}$.
  Therefore, we only have to bound the probability that 
  $\ensuremath{\mathtt{{\color{var}z}}}(\ensuremath{\mathtt{{\color{var}x}}}-\ensuremath{\mathtt{{\color{var}y}}})\bmod 2^{\ensuremath{\mathtt{{\color{var}w}}}}$ looks like \myeqref{all-zeros} or \myeqref{all-ones}.
  
  Let $q$ be the unique odd integer such that $(\ensuremath{\mathtt{{\color{var}x}}}-\ensuremath{\mathtt{{\color{var}y}}})\bmod
  2^{\ensuremath{\mathtt{{\color{var}w}}}}=q2^r$ for some integer $r\ge 0$. By
  \lemref{hashing-mapping}, the binary representation of $\ensuremath{\mathtt{{\color{var}z}}}q\bmod
  2^{\ensuremath{\mathtt{{\color{var}w}}}}$ has $\ensuremath{\mathtt{{\color{var}w}}}-1$ random bits, followed by a 1:
  \[
   \ensuremath{\mathtt{{\color{var}z}}}q\bmod 2^{\ensuremath{\mathtt{{\color{var}w}}}}  = (\underbrace{b_{\ensuremath{\mathtt{{\color{var}w}}}-1},\ldots,b_{1}}_{\ensuremath{\mathtt{{\color{var}w}}}-1},1)_2
  \]
  Therefore, the binary representation of $\ensuremath{\mathtt{{\color{var}z}}}(\ensuremath{\mathtt{{\color{var}x}}}-\ensuremath{\mathtt{{\color{var}y}}})\bmod 2^{\ensuremath{\mathtt{{\color{var}w}}}}=\ensuremath{\mathtt{{\color{var}z}}}q2^r\bmod 2^{\ensuremath{\mathtt{{\color{var}w}}}}$ has
  $\ensuremath{\mathtt{{\color{var}w}}}-r-1$ random bits, followed by a 1, followed by $r$ 0's:
  \[
  \ensuremath{\mathtt{{\color{var}z}}}(\ensuremath{\mathtt{{\color{var}x}}}-\ensuremath{\mathtt{{\color{var}y}}})\bmod 2^{\ensuremath{\mathtt{{\color{var}w}}}}  =
  \ensuremath{\mathtt{{\color{var}z}}}q2^{r}\bmod 2^{\ensuremath{\mathtt{{\color{var}w}}}} =
      (\underbrace{b_{\ensuremath{\mathtt{{\color{var}w}}}-r-1},\ldots,b_{1}}_{\ensuremath{\mathtt{{\color{var}w}}}-r-1},1,\underbrace{0,0,\ldots,0}_{r})_2
  \]
  We can now finish the proof:  If $r > \ensuremath{\mathtt{{\color{var}w}}}-\ensuremath{\mathtt{{\color{var}d}}}$, then the \ensuremath{\mathtt{{\color{var}d}}}
  higher order bits of $\ensuremath{\mathtt{{\color{var}z}}}(\ensuremath{\mathtt{{\color{var}x}}}-\ensuremath{\mathtt{{\color{var}y}}})\bmod 2^{\ensuremath{\mathtt{{\color{var}w}}}}$  contain both 0's
  and 1's, so the probability that $\ensuremath{\mathtt{{\color{var}z}}}(\ensuremath{\mathtt{{\color{var}x}}}-\ensuremath{\mathtt{{\color{var}y}}})\bmod 2^{\ensuremath{\mathtt{{\color{var}w}}}}$ looks
  like \myeqref{all-zeros} or \myeqref{all-ones} is 0.  If $\ensuremath{\mathtt{{\color{var}r}}}=\ensuremath{\mathtt{{\color{var}w}}}-\ensuremath{\mathtt{{\color{var}d}}}$,
  then the probability of looking like \myeqref{all-zeros} is 0, but the
  probability of looking like \myeqref{all-ones} is $1/2^{\ensuremath{\mathtt{{\color{var}d}}}-1}=2/2^{\ensuremath{\mathtt{{\color{var}d}}}}$
  (since we must have $b_1,\ldots,b_{d-1}=1,\ldots,1$).  If $r < \ensuremath{\mathtt{{\color{var}w}}}-\ensuremath{\mathtt{{\color{var}d}}}$,
  then we must have $b_{\ensuremath{\mathtt{{\color{var}w}}}-r-1},\ldots,b_{\ensuremath{\mathtt{{\color{var}w}}}-r-\ensuremath{\mathtt{{\color{var}d}}}}=0,\ldots,0$ or
  $b_{\ensuremath{\mathtt{{\color{var}w}}}-r-1},\ldots,b_{\ensuremath{\mathtt{{\color{var}w}}}-r-\ensuremath{\mathtt{{\color{var}d}}}}=1,\ldots,1$.  The probability of each
  of these cases is $1/2^{\ensuremath{\mathtt{{\color{var}d}}}}$ and they are mutually exclusive, so the
  probability of either of these cases is $2/2^{\ensuremath{\mathtt{{\color{var}d}}}}$.  This completes
  the proof.
\end{proof}

\subsection{Summary}

The following theorem summarizes the performance of a \texttt{ChainedHashTable}
data structure:

\begin{thm}\thmlabel{hashtable}
  A \texttt{ChainedHashTable} implements the \texttt{USet} interface.  Ignoring the cost of
  calls to \ensuremath{\mathtt{grow()}}, a \texttt{ChainedHashTable} supports the operations \ensuremath{\mathtt{add({\color{var}x})}},
  \ensuremath{\mathtt{remove({\color{var}x})}}, and \ensuremath{\mathtt{find({\color{var}x})}} in $O(1)$ expected time per operation.

  Furthermore, beginning with an empty \texttt{ChainedHashTable}, any sequence of
  $m$ \ensuremath{\mathtt{add({\color{var}x})}} and \ensuremath{\mathtt{remove({\color{var}x})}} operations results in a total of $O(m)$
  time spent during all calls to \ensuremath{\mathtt{grow()}}.
\end{thm}

\section{\texttt{LinearHashTable}: Linear Probing}
\seclabel{linearhashtable}

\index{LinearHashTable@\texttt{LinearHashTable}}%
The \texttt{ChainedHashTable} data structure uses an array of lists, where the \ensuremath{\mathtt{{\color{var}i}}}th
list stores all elements \ensuremath{\mathtt{{\color{var}x}}} such that $\ensuremath{\mathtt{hash({\color{var}x})}}=\ensuremath{\mathtt{{\color{var}i}}}$.  An alternative,
called \emph{open addressing}
\index{open addressing}%
is to store the elements directly in an
array, \ensuremath{\mathtt{{\color{var}t}}}, with each array location in \ensuremath{\mathtt{{\color{var}t}}} storing at most one value.
This approach is taken by the \texttt{LinearHashTable} described in this
section.  In some places, this data structure is described as \emph{open
addressing with linear probing}.
\index{linear probing}%

The main idea behind a \texttt{LinearHashTable} is that we would, ideally, like
to store the element \ensuremath{\mathtt{{\color{var}x}}} with hash value \ensuremath{\mathtt{{\color{var}i}=hash({\color{var}x})}} in the table location
\ensuremath{\mathtt{{\color{var}t}[{\color{var}i}]}}.  If we cannot do this (because some element is already stored
there) then we try to store it at location $\ensuremath{\mathtt{{\color{var}t}}}[(\ensuremath{\mathtt{{\color{var}i}}}+1)\bmod\ensuremath{\mathtt{{\color{var}t}.{\color{var}length}}}]$;
if that's not possible, then we try $\ensuremath{\mathtt{{\color{var}t}}}[(\ensuremath{\mathtt{{\color{var}i}}}+2)\bmod\ensuremath{\mathtt{{\color{var}t}.{\color{var}length}}}]$,
and so on, until we find a place for \ensuremath{\mathtt{{\color{var}x}}}.

There are three types of entries stored in \ensuremath{\mathtt{{\color{var}t}}}: 
\begin{enumerate}
  \item data values: actual values in the \texttt{USet} that we are representing;
  \item \ensuremath{\mathtt{{\color{keyword}null}}} values: at array locations where no data has ever been
  stored; and
  \item \ensuremath{\mathtt{{\color{var}del}}} values: at array locations where data was once stored but that
   has since been deleted.
\end{enumerate}
In addition to the counter, \ensuremath{\mathtt{{\color{var}n}}}, that keeps track of the number of elements
in the \texttt{LinearHashTable}, a counter, \ensuremath{\mathtt{{\color{var}q}}}, keeps track of the number of
elements of Types~1 and 3.  That is, \ensuremath{\mathtt{{\color{var}q}}} is equal to \ensuremath{\mathtt{{\color{var}n}}} plus the number
of \ensuremath{\mathtt{{\color{var}del}}} values in \ensuremath{\mathtt{{\color{var}t}}}.  To make this work efficiently, we need
\ensuremath{\mathtt{{\color{var}t}}} to be considerably larger than \ensuremath{\mathtt{{\color{var}q}}}, so that there are lots of \ensuremath{\mathtt{{\color{keyword}null}}}
values in \ensuremath{\mathtt{{\color{var}t}}}.  The operations on a \texttt{LinearHashTable} therefore maintain
the invariant that $\ensuremath{\mathtt{{\color{var}t}.{\color{var}length}}}\ge 2\ensuremath{\mathtt{{\color{var}q}}}$.

To summarize, a \texttt{LinearHashTable} contains an array, \ensuremath{\mathtt{{\color{var}t}}}, that stores
data elements, and integers \ensuremath{\mathtt{{\color{var}n}}} and \ensuremath{\mathtt{{\color{var}q}}} that keep track of the number of
data elements and non-\ensuremath{\mathtt{{\color{keyword}null}}} values of \ensuremath{\mathtt{{\color{var}t}}}, respectively.  Because many
hash functions only work for table sizes that are a power of 2, we also
keep an integer \ensuremath{\mathtt{{\color{var}d}}} and maintain the invariant that $\ensuremath{\mathtt{{\color{var}t}.{\color{var}length}}}=2^\ensuremath{\mathtt{{\color{var}d}}}$.


 \ \vspace{.5ex}\begin{Verbatim}[gobble=1,tabsize=2,frame=single,framerule=1.2pt,commandchars=\\@\$,label=\texttt{LinearHashTable},labelposition=topline]
	@\color@keyword$T$[] @\color@var$t$;   @\color@comment$// the table$
	@\color@keyword$int$ @\color@var$n$;   @\color@comment$// the size$
	@\color@keyword$int$ @\color@var$d$;   @\color@comment$// t.length = 2^d$
	@\color@keyword$int$ @\color@var$q$;   @\color@comment$// number of non-null entries in t$
\end{Verbatim}


The \ensuremath{\mathtt{find({\color{var}x})}} operation in a \texttt{LinearHashTable} is simple.  We start
at array entry \ensuremath{\mathtt{{\color{var}t}[{\color{var}i}]}} where $\ensuremath{\mathtt{{\color{var}i}}}=\ensuremath{\mathtt{hash({\color{var}x})}}$ and search entries \ensuremath{\mathtt{{\color{var}t}[{\color{var}i}]}},
$\ensuremath{\mathtt{{\color{var}t}}}[(\ensuremath{\mathtt{{\color{var}i}}}+1)\bmod \ensuremath{\mathtt{{\color{var}t}.{\color{var}length}}}]$, $\ensuremath{\mathtt{{\color{var}t}}}[(\ensuremath{\mathtt{{\color{var}i}}}+2)\bmod \ensuremath{\mathtt{{\color{var}t}.{\color{var}length}}}]$, and so on,
until we find an index \ensuremath{\mathtt{{\color{var}i}'}} such that, either, \ensuremath{\mathtt{{\color{var}t}[{\color{var}i}']={\color{var}x}}}, or \ensuremath{\mathtt{{\color{var}t}[{\color{var}i}']={\color{keyword}null}}}.
In the former case we return \ensuremath{\mathtt{{\color{var}t}[{\color{var}i}']}}. In the latter case, we conclude
that \ensuremath{\mathtt{{\color{var}x}}} is not contained in the hash table and return \ensuremath{\mathtt{{\color{keyword}null}}}.


 \ \vspace{.5ex}\begin{Verbatim}[gobble=1,tabsize=2,frame=single,framerule=1.2pt,commandchars=\\@\$,label=\texttt{LinearHashTable},labelposition=topline]
	@\color@keyword$T$ find(@\color@keyword$T$ @\color@var$x$) {
		@\color@keyword$int$ @\color@var$i$ = hash(@\color@var$x$);
		@\color@keyword$while$ (@\color@var$t$[@\color@var$i$] != @\color@keyword$null$) {
			@\color@keyword$if$ (@\color@var$t$[@\color@var$i$] != @\color@var$del$ && @\color@var$x$.equals(@\color@var$t$[@\color@var$i$])) @\color@var$return$ @\color@var$t$[@\color@var$i$];
			@\color@var$i$ = (@\color@var$i$ == @\color@var$t$.@\color@var$length$-1) ? 0 : @\color@var$i$ + 1; @\color@comment$// increment i$
		}
		@\color@var$return$ @\color@keyword$null$;
	}
\end{Verbatim}


The \ensuremath{\mathtt{add({\color{var}x})}} operation is also fairly easy to implement. After checking
that \ensuremath{\mathtt{{\color{var}x}}} is not already stored in the table (using \ensuremath{\mathtt{find({\color{var}x})}}), we search
\ensuremath{\mathtt{{\color{var}t}[{\color{var}i}]}}, $\ensuremath{\mathtt{{\color{var}t}}}[(\ensuremath{\mathtt{{\color{var}i}}}+1)\bmod \ensuremath{\mathtt{{\color{var}t}.{\color{var}length}}}]$, $\ensuremath{\mathtt{{\color{var}t}}}[(\ensuremath{\mathtt{{\color{var}i}}}+2)\bmod \ensuremath{\mathtt{{\color{var}t}.{\color{var}length}}}]$,
and so on, until we find a \ensuremath{\mathtt{{\color{keyword}null}}} or \ensuremath{\mathtt{{\color{var}del}}} and store \ensuremath{\mathtt{{\color{var}x}}} at that location,
increment \ensuremath{\mathtt{{\color{var}n}}}, and \ensuremath{\mathtt{{\color{var}q}}}, if appropriate.


 \ \vspace{.5ex}\begin{Verbatim}[gobble=1,tabsize=2,frame=single,framerule=1.2pt,commandchars=\\@\$,label=\texttt{LinearHashTable},labelposition=topline]
	@\color@var$boolean$ add(@\color@keyword$T$ @\color@var$x$) {
		@\color@keyword$if$ (find(@\color@var$x$) != @\color@keyword$null$) @\color@var$return$ @\color@var$false$;
		@\color@keyword$if$ (2*(@\color@var$q$+1) > @\color@var$t$.@\color@var$length$) resize(); @\color@comment$// max 50% occupancy$
		@\color@keyword$int$ @\color@var$i$ = hash(@\color@var$x$);
		@\color@keyword$while$ (@\color@var$t$[@\color@var$i$] != @\color@keyword$null$ && @\color@var$t$[@\color@var$i$] != @\color@var$del$)
			@\color@var$i$ = (@\color@var$i$ == @\color@var$t$.@\color@var$length$-1) ? 0 : @\color@var$i$ + 1; @\color@comment$// increment i$
		@\color@keyword$if$ (@\color@var$t$[@\color@var$i$] == @\color@keyword$null$) @\color@var$q$++;
		@\color@var$n$++;
		@\color@var$t$[@\color@var$i$] = @\color@var$x$;
		@\color@var$return$ @\color@var$true$;
	}
\end{Verbatim}


By now, the implementation of the \ensuremath{\mathtt{remove({\color{var}x})}} operation should be obvious.
We search \ensuremath{\mathtt{{\color{var}t}[{\color{var}i}]}}, $\ensuremath{\mathtt{{\color{var}t}}}[(\ensuremath{\mathtt{{\color{var}i}}}+1)\bmod \ensuremath{\mathtt{{\color{var}t}.{\color{var}length}}}]$, $\ensuremath{\mathtt{{\color{var}t}}}[(\ensuremath{\mathtt{{\color{var}i}}}+2)\bmod
\ensuremath{\mathtt{{\color{var}t}.{\color{var}length}}}]$, and so on until we find an index \ensuremath{\mathtt{{\color{var}i}'}} such that \ensuremath{\mathtt{{\color{var}t}[{\color{var}i}']={\color{var}x}}}
or \ensuremath{\mathtt{{\color{var}t}[{\color{var}i}']={\color{keyword}null}}}.  In the former case, we set \ensuremath{\mathtt{{\color{var}t}[{\color{var}i}']={\color{var}del}}} and return
\ensuremath{\mathtt{{\color{var}true}}}.  In the latter case we conclude that \ensuremath{\mathtt{{\color{var}x}}} was not stored in the
table (and therefore cannot be deleted) and return \ensuremath{\mathtt{{\color{var}false}}}.


 \ \vspace{.5ex}\begin{Verbatim}[gobble=1,tabsize=2,frame=single,framerule=1.2pt,commandchars=\\@\$,label=\texttt{LinearHashTable},labelposition=topline]
	@\color@keyword$T$ remove(@\color@keyword$T$ @\color@var$x$) {
		@\color@keyword$int$ @\color@var$i$ = hash(@\color@var$x$);
		@\color@keyword$while$ (@\color@var$t$[@\color@var$i$] != @\color@keyword$null$) {
			@\color@keyword$T$ @\color@var$y$ = @\color@var$t$[@\color@var$i$];
			@\color@keyword$if$ (@\color@var$y$ != @\color@var$del$ && @\color@var$x$.equals(@\color@var$y$)) { 
				@\color@var$t$[@\color@var$i$] = @\color@var$del$;
				@\color@var$n$--;
				@\color@keyword$if$ (8*@\color@var$n$ < @\color@var$t$.@\color@var$length$) resize(); @\color@comment$// min 12.5% occupancy$
				@\color@var$return$ @\color@var$y$;
			}
			@\color@var$i$ = (@\color@var$i$ == @\color@var$t$.@\color@var$length$-1) ? 0 : @\color@var$i$ + 1;  @\color@comment$// increment i$
		}
		@\color@var$return$ @\color@keyword$null$;
	}
\end{Verbatim}


The correctness of the \ensuremath{\mathtt{find({\color{var}x})}}, \ensuremath{\mathtt{add({\color{var}x})}}, and \ensuremath{\mathtt{remove({\color{var}x})}} methods is
easy to verify, though it relies on the use of \ensuremath{\mathtt{{\color{var}del}}} values.  Notice
that none of these operations ever sets a non-\ensuremath{\mathtt{{\color{keyword}null}}} entry to \ensuremath{\mathtt{{\color{keyword}null}}}.
Therefore, when we reach an index \ensuremath{\mathtt{{\color{var}i}'}} such that \ensuremath{\mathtt{{\color{var}t}[{\color{var}i}']={\color{keyword}null}}}, this is
a proof that the element, \ensuremath{\mathtt{{\color{var}x}}}, that we are searching for is not stored
in the table; \ensuremath{\mathtt{{\color{var}t}[{\color{var}i}']}} has always been \ensuremath{\mathtt{{\color{keyword}null}}}, so there is no reason that
a previous \ensuremath{\mathtt{add({\color{var}x})}} operation would have proceeded beyond index \ensuremath{\mathtt{{\color{var}i}'}}.

The \ensuremath{\mathtt{resize()}} method is called by \ensuremath{\mathtt{add({\color{var}x})}} when the number of non-\ensuremath{\mathtt{{\color{keyword}null}}}
entries exceeds $\ensuremath{\mathtt{{\color{var}t}.{\color{var}length}}}/2$ or by \ensuremath{\mathtt{remove({\color{var}x})}} when the number of
data entries is less than \ensuremath{\mathtt{{\color{var}t}.{\color{var}length}/8}}.  The \ensuremath{\mathtt{resize()}} method works
like the \ensuremath{\mathtt{resize()}} methods in other array-based data structures.
We find the smallest non-negative integer \ensuremath{\mathtt{{\color{var}d}}} such that $2^{\ensuremath{\mathtt{{\color{var}d}}}}
\ge 3\ensuremath{\mathtt{{\color{var}n}}}$.  We reallocate the array \ensuremath{\mathtt{{\color{var}t}}} so that it has size $2^{\ensuremath{\mathtt{{\color{var}d}}}}$,
and then we insert all the elements in the old version of \ensuremath{\mathtt{{\color{var}t}}} into the
newly-resized copy of \ensuremath{\mathtt{{\color{var}t}}}.  While doing this, we reset \ensuremath{\mathtt{{\color{var}q}}} equal to \ensuremath{\mathtt{{\color{var}n}}}
since the newly-allocated \ensuremath{\mathtt{{\color{var}t}}} contains no \ensuremath{\mathtt{{\color{var}del}}} values.


 \ \vspace{.5ex}\begin{Verbatim}[gobble=1,tabsize=2,frame=single,framerule=1.2pt,commandchars=\\@\$,label=\texttt{LinearHashTable},labelposition=topline]
	@\color@keyword$void$ resize() {
		@\color@var$d$ = 1;
		@\color@keyword$while$ ((1<<@\color@var$d$) < 3*@\color@var$n$) @\color@var$d$++;
		@\color@keyword$T$[] @\color@var$told$ = @\color@var$t$;
		@\color@var$t$ = newArray(1<<@\color@var$d$);
		@\color@var$q$ = @\color@var$n$;
		@\color@comment$// insert everything from told$
		@\color@keyword$for$ (@\color@keyword$int$ @\color@var$k$ = 0; @\color@var$k$ < @\color@var$told$.@\color@var$length$; @\color@var$k$++) {
			@\color@keyword$if$ (@\color@var$told$[@\color@var$k$] != @\color@keyword$null$ && @\color@var$told$[@\color@var$k$] != @\color@var$del$) {
				@\color@keyword$int$ @\color@var$i$ = hash(@\color@var$told$[@\color@var$k$]);
				@\color@keyword$while$ (@\color@var$t$[@\color@var$i$] != @\color@keyword$null$) 
					@\color@var$i$ = (@\color@var$i$ == @\color@var$t$.@\color@var$length$-1) ? 0 : @\color@var$i$ + 1;
				@\color@var$t$[@\color@var$i$] = @\color@var$told$[@\color@var$k$];
			}
		}
	}
\end{Verbatim}


\subsection{Analysis of Linear Probing}

Notice that each operation, \ensuremath{\mathtt{add({\color{var}x})}}, \ensuremath{\mathtt{remove({\color{var}x})}}, or \ensuremath{\mathtt{find({\color{var}x})}}, finishes
as soon as (or before) it discovers the first \ensuremath{\mathtt{{\color{keyword}null}}} entry in \ensuremath{\mathtt{{\color{var}t}}}.
The intuition behind the analysis of linear probing is that, since at
least half the elements in \ensuremath{\mathtt{{\color{var}t}}} are equal to \ensuremath{\mathtt{{\color{keyword}null}}}, an operation should
not take long to complete because it will very quickly come across a
\ensuremath{\mathtt{{\color{keyword}null}}} entry.  We shouldn't rely too heavily on this intuition, though,
because it would lead us to (the incorrect) conclusion that the expected
number of locations in \ensuremath{\mathtt{{\color{var}t}}} examined by an operation is at most 2.

For the rest of this section, we will assume that all hash values are
independently and uniformly distributed in $\{0,\ldots,\ensuremath{\mathtt{{\color{var}t}.{\color{var}length}}}-1\}$.
This is not a realistic assumption, but it will make it possible for
us to analyze linear probing.  Later in this section we will describe a
method, called tabulation hashing, that produces a hash function that is
``good enough'' for linear probing.  We will also assume that all indices
into the positions of \ensuremath{\mathtt{{\color{var}t}}} are taken modulo \ensuremath{\mathtt{{\color{var}t}.{\color{var}length}}}, so that \ensuremath{\mathtt{{\color{var}t}[{\color{var}i}]}}
is really a shorthand for $\ensuremath{\mathtt{{\color{var}t}}}[\ensuremath{\mathtt{{\color{var}i}}}\bmod\ensuremath{\mathtt{{\color{var}t}.{\color{var}length}}}]$.

\index{run}%
We say that a \emph{run of length $k$ that starts at \ensuremath{\mathtt{{\color{var}i}}}} occurs when all
the table entries $\ensuremath{\mathtt{{\color{var}t}[{\color{var}i}]}}, \ensuremath{\mathtt{{\color{var}t}[{\color{var}i}+1]}},\ldots,\ensuremath{\mathtt{{\color{var}t}}}[\ensuremath{\mathtt{{\color{var}i}}}+k-1]$ are non-\ensuremath{\mathtt{{\color{keyword}null}}}
and $\ensuremath{\mathtt{{\color{var}t}}}[\ensuremath{\mathtt{{\color{var}i}}}-1]=\ensuremath{\mathtt{{\color{var}t}}}[\ensuremath{\mathtt{{\color{var}i}}}+k]=\ensuremath{\mathtt{{\color{keyword}null}}}$.  The number of non-\ensuremath{\mathtt{{\color{keyword}null}}} elements of
\ensuremath{\mathtt{{\color{var}t}}} is exactly \ensuremath{\mathtt{{\color{var}q}}} and the \ensuremath{\mathtt{add({\color{var}x})}} method ensures that, at all times,
$\ensuremath{\mathtt{{\color{var}q}}}\le\ensuremath{\mathtt{{\color{var}t}.{\color{var}length}}}/2$.  There are \ensuremath{\mathtt{{\color{var}q}}} elements $\ensuremath{\mathtt{{\color{var}x}}}_1,\ldots,\ensuremath{\mathtt{{\color{var}x}}}_{\ensuremath{\mathtt{{\color{var}q}}}}$
that have been inserted into \ensuremath{\mathtt{{\color{var}t}}} since the last \ensuremath{\mathtt{resize()}} operation.
By our assumption, each of these has a hash value, $\ensuremath{\mathtt{{\color{var}hash}}}(\ensuremath{\mathtt{{\color{var}x}}}_j)$,
that is uniform and independent of the rest.  With this setup, we can
prove the main lemma required to analyze linear probing.

\begin{lem}\lemlabel{linear-probing}
Fix a value $\ensuremath{\mathtt{{\color{var}i}}}\in\{0,\ldots,\ensuremath{\mathtt{{\color{var}t}.{\color{var}length}}}-1\}$.  Then the probability that
a run of length $k$ starts at \ensuremath{\mathtt{{\color{var}i}}} is $O(c^k)$ for some constant $0<c<1$.
\end{lem}

\begin{proof}
If a run of length $k$ starts at \ensuremath{\mathtt{{\color{var}i}}}, then there are exactly $k$
elements $\ensuremath{\mathtt{{\color{var}x}}}_j$ such that $\ensuremath{\mathtt{{\color{var}hash}}}(\ensuremath{\mathtt{{\color{var}x}}}_j)\in\{\ensuremath{\mathtt{{\color{var}i}}},\ldots,\ensuremath{\mathtt{{\color{var}i}}}+k-1\}$.
The probability that this occurs is exactly
\[
  p_k  = \binom{\ensuremath{\mathtt{{\color{var}q}}}}{k}\left(\frac{k}{\ensuremath{\mathtt{{\color{var}t}.{\color{var}length}}}}\right)^k\left(\frac{\ensuremath{\mathtt{{\color{var}t}.{\color{var}length}}}-k}{\ensuremath{\mathtt{{\color{var}t}.{\color{var}length}}}}\right)^{\ensuremath{\mathtt{{\color{var}q}}}-k} \enspace ,
\]
since, for each choice of $k$ elements, these $k$ elements must hash to
one of the $k$ locations and the remaining $\ensuremath{\mathtt{{\color{var}q}}}-k$ elements must hash to
the other $\ensuremath{\mathtt{{\color{var}t}.{\color{var}length}}}-k$ table locations.\footnote{Note that $p_k$ is greater
than the probability that a run of length $k$ starts at \ensuremath{\mathtt{{\color{var}i}}}, since the definition of $p_k$ does not include the requirement $\ensuremath{\mathtt{{\color{var}t}}}[\ensuremath{\mathtt{{\color{var}i}}}-1]=\ensuremath{\mathtt{{\color{var}t}}}[\ensuremath{\mathtt{{\color{var}i}}}+k]=\ensuremath{\mathtt{{\color{keyword}null}}}$.}

In the following derivation we will cheat a little and replace $r!$ with
$(r/e)^r$. Stirling's Approximation (\secref{factorials}) shows that this
is only a factor of $O(\sqrt{r})$ from the truth.  This is just done to
make the derivation simpler; \excref{linear-probing} asks the reader to
redo the calculation more rigorously using Stirling's Approximation in
its entirety.

The value of $p_k$ is maximized when \ensuremath{\mathtt{{\color{var}t}.{\color{var}length}}} is minimum, and the data
structure maintains the invariant that $\ensuremath{\mathtt{{\color{var}t}.{\color{var}length}}} \ge 2\ensuremath{\mathtt{{\color{var}q}}}$, so
\begin{align*}
   p_k & \le \binom{\ensuremath{\mathtt{{\color{var}q}}}}{k}\left(\frac{k}{2\ensuremath{\mathtt{{\color{var}q}}}}\right)^k\left(\frac{2\ensuremath{\mathtt{{\color{var}q}}}-k}{2\ensuremath{\mathtt{{\color{var}q}}}}\right)^{\ensuremath{\mathtt{{\color{var}q}}}-k} \\
  & = \left(\frac{\ensuremath{\mathtt{{\color{var}q}}}!}{(\ensuremath{\mathtt{{\color{var}q}}}-k)!k!}\right)\left(\frac{k}{2\ensuremath{\mathtt{{\color{var}q}}}}\right)^k\left(\frac{2\ensuremath{\mathtt{{\color{var}q}}}-k}{2\ensuremath{\mathtt{{\color{var}q}}}}\right)^{\ensuremath{\mathtt{{\color{var}q}}}-k} \\
  & \approx \left(\frac{\ensuremath{\mathtt{{\color{var}q}}}^{\ensuremath{\mathtt{{\color{var}q}}}}}{(\ensuremath{\mathtt{{\color{var}q}}}-k)^{\ensuremath{\mathtt{{\color{var}q}}}-k}k^k}\right)\left(\frac{k}{2\ensuremath{\mathtt{{\color{var}q}}}}\right)^k\left(\frac{2\ensuremath{\mathtt{{\color{var}q}}}-k}{2\ensuremath{\mathtt{{\color{var}q}}}}\right)^{\ensuremath{\mathtt{{\color{var}q}}}-k} && \text{[Stirling's approximation]} \\
  & = \left(\frac{\ensuremath{\mathtt{{\color{var}q}}}^{k}\ensuremath{\mathtt{{\color{var}q}}}^{\ensuremath{\mathtt{{\color{var}q}}}-k}}{(\ensuremath{\mathtt{{\color{var}q}}}-k)^{\ensuremath{\mathtt{{\color{var}q}}}-k}k^k}\right)\left(\frac{k}{2\ensuremath{\mathtt{{\color{var}q}}}}\right)^k\left(\frac{2\ensuremath{\mathtt{{\color{var}q}}}-k}{2\ensuremath{\mathtt{{\color{var}q}}}}\right)^{\ensuremath{\mathtt{{\color{var}q}}}-k} \\
 & = \left(\frac{\ensuremath{\mathtt{{\color{var}q}}}k}{2\ensuremath{\mathtt{{\color{var}q}}}k}\right)^k
     \left(\frac{\ensuremath{\mathtt{{\color{var}q}}}(2\ensuremath{\mathtt{{\color{var}q}}}-k)}{2\ensuremath{\mathtt{{\color{var}q}}}(\ensuremath{\mathtt{{\color{var}q}}}-k)}\right)^{\ensuremath{\mathtt{{\color{var}q}}}-k} \\
 & = \left(\frac{1}{2}\right)^k
     \left(\frac{(2\ensuremath{\mathtt{{\color{var}q}}}-k)}{2(\ensuremath{\mathtt{{\color{var}q}}}-k)}\right)^{\ensuremath{\mathtt{{\color{var}q}}}-k} \\
 & = \left(\frac{1}{2}\right)^k
     \left(1+\frac{k}{2(\ensuremath{\mathtt{{\color{var}q}}}-k)}\right)^{\ensuremath{\mathtt{{\color{var}q}}}-k} \\
 & \le \left(\frac{\sqrt{e}}{2}\right)^k \enspace .
\end{align*}
(In the last step, we use the inequality $(1+1/x)^x \le e$, which holds
for all $x>0$.)  Since $\sqrt{e}/{2}< 0.824360636 < 1$, this completes
the proof.
\end{proof}

Using \lemref{linear-probing} to prove upper-bounds on the expected
running time of \ensuremath{\mathtt{find({\color{var}x})}}, \ensuremath{\mathtt{add({\color{var}x})}}, and \ensuremath{\mathtt{remove({\color{var}x})}} is now fairly
straightforward.  Consider the simplest case, where we execute
\ensuremath{\mathtt{find({\color{var}x})}} for some value \ensuremath{\mathtt{{\color{var}x}}} that has never been stored in the
\texttt{LinearHashTable}.  In this case, $\ensuremath{\mathtt{{\color{var}i}}}=\ensuremath{\mathtt{hash({\color{var}x})}}$ is a random value in
$\{0,\ldots,\ensuremath{\mathtt{{\color{var}t}.{\color{var}length}}}-1\}$ independent of the contents of \ensuremath{\mathtt{{\color{var}t}}}.  If \ensuremath{\mathtt{{\color{var}i}}}
is part of a run of length $k$, then the time it takes to execute the
\ensuremath{\mathtt{find({\color{var}x})}} operation is at most $O(1+k)$.  Thus, the expected running time
can be upper-bounded by
\[
  O\left(1 + \left(\frac{1}{\ensuremath{\mathtt{{\color{var}t}.{\color{var}length}}}}\right)\sum_{i=1}^{\ensuremath{\mathtt{{\color{var}t}.{\color{var}length}}}}\sum_{k=0}^{\infty} k\Pr\{\text{\ensuremath{\mathtt{{\color{var}i}}} is part of a run of length $k$}\}\right) \enspace .
\]
Note that each run of length $k$ contributes to the inner sum $k$ times for a total contribution of $k^2$,  so the above sum can be rewritten as
\begin{align*}
  & { } O\left(1 + \left(\frac{1}{\ensuremath{\mathtt{{\color{var}t}.{\color{var}length}}}}\right)\sum_{i=1}^{\ensuremath{\mathtt{{\color{var}t}.{\color{var}length}}}}\sum_{k=0}^{\infty} k^2\Pr\{\mbox{\ensuremath{\mathtt{{\color{var}i}}} starts a run of length $k$}\}\right) \\
  & \le O\left(1 + \left(\frac{1}{\ensuremath{\mathtt{{\color{var}t}.{\color{var}length}}}}\right)\sum_{i=1}^{\ensuremath{\mathtt{{\color{var}t}.{\color{var}length}}}}\sum_{k=0}^{\infty} k^2p_k\right) \\
  & = O\left(1 + \sum_{k=0}^{\infty} k^2p_k\right) \\
  & = O\left(1 + \sum_{k=0}^{\infty} k^2\cdot O(c^k)\right) \\
  & = O(1) \enspace .
\end{align*}
The last step in this derivation comes from the fact that
$\sum_{k=0}^{\infty} k^2\cdot O(c^k)$ is an exponentially decreasing
series.\footnote{In the terminology of many calculus texts, this sum
passes the ratio test:  There exists a positive integer $k_0$ such
that, for all $k\ge k_0$, $\frac{(k+1)^2c^{k+1}}{k^2c^k} < 1$.}
Therefore, we conclude that the expected running time of the \ensuremath{\mathtt{find({\color{var}x})}} operation
for a value \ensuremath{\mathtt{{\color{var}x}}} that is not contained in a \texttt{LinearHashTable} is $O(1)$.

If we ignore the cost of the \ensuremath{\mathtt{resize()}} operation, then the above
analysis gives us all we need to analyze the cost of operations on a
\texttt{LinearHashTable}.

First of all, the analysis of \ensuremath{\mathtt{find({\color{var}x})}} given above applies to the
\ensuremath{\mathtt{add({\color{var}x})}} operation when \ensuremath{\mathtt{{\color{var}x}}} is not contained in the table.  To analyze
the \ensuremath{\mathtt{find({\color{var}x})}} operation when \ensuremath{\mathtt{{\color{var}x}}} is contained in the table, we need only
note that this is the same as the cost of the \ensuremath{\mathtt{add({\color{var}x})}} operation that
previously added \ensuremath{\mathtt{{\color{var}x}}} to the table.  Finally, the cost of a \ensuremath{\mathtt{remove({\color{var}x})}}
operation is the same as the cost of a \ensuremath{\mathtt{find({\color{var}x})}} operation.

In summary, if we ignore the cost of calls to \ensuremath{\mathtt{resize()}}, all operations
on a \texttt{LinearHashTable} run in $O(1)$ expected time.  Accounting for the
cost of resize can be done using the same type of amortized analysis
performed for the \texttt{ArrayStack} data structure in \secref{arraystack}.

\subsection{Summary}

The following theorem summarizes the performance of the \texttt{LinearHashTable}
data structure:

\begin{thm}\thmlabel{linear-probing}
  A \texttt{LinearHashTable} implements the \texttt{USet} interface.  Ignoring the
  cost of calls to \ensuremath{\mathtt{resize()}}, a \texttt{LinearHashTable} supports the operations
  \ensuremath{\mathtt{add({\color{var}x})}}, \ensuremath{\mathtt{remove({\color{var}x})}}, and \ensuremath{\mathtt{find({\color{var}x})}} in $O(1)$ expected time per
  operation.  

  Furthermore, beginning with an empty \texttt{LinearHashTable}, any
  sequence of $m$ \ensuremath{\mathtt{add({\color{var}x})}} and \ensuremath{\mathtt{remove({\color{var}x})}} operations results in a total
  of $O(m)$ time spent during all calls to \ensuremath{\mathtt{resize()}}.
\end{thm}

\subsection{Tabulation Hashing}
\seclabel{tabulation}

\index{tabulation hashing}%
While analyzing the \texttt{LinearHashTable} structure, we made a very strong
assumption:  That for any set of elements, $\{\ensuremath{\mathtt{{\color{var}x}}}_1,\ldots,\ensuremath{\mathtt{{\color{var}x}}}_\ensuremath{\mathtt{{\color{var}n}}}\}$,
the hash values $\ensuremath{\mathtt{{\color{var}hash}}}($x$_1),\ldots,\ensuremath{\mathtt{{\color{var}hash}}}(\ensuremath{\mathtt{{\color{var}x}}}_\ensuremath{\mathtt{{\color{var}n}}})$ are independently and
uniformly distributed over the set $\{0,\ldots,\ensuremath{\mathtt{{\color{var}t}.{\color{var}length}}}-1\}$.  One way to
achieve this is to store a giant array, \ensuremath{\mathtt{{\color{var}tab}}}, of length $2^{\ensuremath{\mathtt{{\color{var}w}}}}$,
where each entry is a random \ensuremath{\mathtt{{\color{var}w}}}-bit integer, independent of all the
other entries.  In this way, we could implement \ensuremath{\mathtt{hash({\color{var}x})}} by extracting
a \ensuremath{\mathtt{{\color{var}d}}}-bit integer from \ensuremath{\mathtt{{\color{var}tab}[{\color{var}x}.hashCode()]}}:


 \ \vspace{.5ex}\begin{Verbatim}[gobble=1,tabsize=2,frame=single,framerule=1.2pt,commandchars=\\@\$,label=\texttt{LinearHashTable},labelposition=topline]
	@\color@keyword$int$ idealHash(@\color@keyword$T$ @\color@var$x$) {
		@\color@var$return$ @\color@var$tab$[@\color@var$x$.hashCode() >>> @\color@var$w$-@\color@var$d$];
	}
\end{Verbatim}

\pcodeonly{Here, \ensuremath{\mathtt{\text{\ttfamily >>}}}, is the \emph{bitwise right shift} operator, so
\ensuremath{\mathtt{{\color{var}x}.hashCode() \text{\ttfamily >>} {\color{var}w}-{\color{var}d}}} extracts the \ensuremath{\mathtt{{\color{var}d}}} most significant bits of \ensuremath{\mathtt{{\color{var}x}}}'s \ensuremath{\mathtt{{\color{var}w}}}-bit hash code.}

Unfortunately, storing an array of size $2^{\ensuremath{\mathtt{{\color{var}w}}}}$ is prohibitive in terms
of memory usage.  The approach used by \emph{tabulation hashing} is to,
instead, treat \ensuremath{\mathtt{{\color{var}w}}}-bit integers as being comprised of $\ensuremath{\mathtt{{\color{var}w}}}/\ensuremath{\mathtt{{\color{var}r}}}$
integers, each having only $\ensuremath{\mathtt{{\color{var}r}}}$ bits.  In this way, tabulation hashing
only needs $\ensuremath{\mathtt{{\color{var}w}}}/\ensuremath{\mathtt{{\color{var}r}}}$ arrays each of length $2^{\ensuremath{\mathtt{{\color{var}r}}}}$.  All the entries
in these arrays are independent random \ensuremath{\mathtt{{\color{var}w}}}-bit integers.  To obtain the value
of \ensuremath{\mathtt{hash({\color{var}x})}} we split \ensuremath{\mathtt{{\color{var}x}.hashCode()}} up into $\ensuremath{\mathtt{{\color{var}w}}}/\ensuremath{\mathtt{{\color{var}r}}}$ \ensuremath{\mathtt{{\color{var}r}}}-bit integers
and use these as indices into these arrays.  We then combine all these
values with the bitwise exclusive-or operator to obtain \ensuremath{\mathtt{hash({\color{var}x})}}.
The following code shows how this works when $\ensuremath{\mathtt{{\color{var}w}}}=32$ and $\ensuremath{\mathtt{{\color{var}r}}}=4$:


 \ \vspace{.5ex}\begin{Verbatim}[gobble=1,tabsize=2,frame=single,framerule=1.2pt,commandchars=\\@\$,label=\texttt{LinearHashTable},labelposition=topline]
	@\color@keyword$int$ hash(@\color@keyword$T$ @\color@var$x$) {
		@\color@keyword$int$ @\color@var$h$ = @\color@var$x$.hashCode();
		return (@\color@var$tab$[0][@\color@var$h$&0xff] 
		         ^ @\color@var$tab$[1][(@\color@var$h$>>>8)&0xff]
		         ^ @\color@var$tab$[2][(@\color@var$h$>>>16)&0xff] 
		         ^ @\color@var$tab$[3][(@\color@var$h$>>>24)&0xff])
		  	    >>> (@\color@var$w$-@\color@var$d$);
	}
\end{Verbatim}

In this case, \ensuremath{\mathtt{{\color{var}tab}}} is a two-dimensional array with four columns and
$2^{32/4}=256$ rows.
\pcodeonly{
Quantities like \ensuremath{\mathtt{0xff}}, used above, are \emph{hexadecimal numbers}
\index{hexadecimal numbers}%
whose digits have 16 possible values 0--9, which have their usual
meaning and a--f, which denote 10--15.  The number $\ensuremath{\mathtt{0xff}}=15\cdot 16 + 15 = 255$.
The \ensuremath{\mathtt{\text{\ttfamily \&}}} symbol is the \emph{bitwise and} operator, so
code like \ensuremath{\mathtt{{\color{var}h} \text{\ttfamily >>} 8 \text{\ttfamily \&} 0xff}} extracts bits with index 8 through 15 of \ensuremath{\mathtt{{\color{var}h}}}.}

One can easily verify that, for any \ensuremath{\mathtt{{\color{var}x}}}, \ensuremath{\mathtt{hash({\color{var}x})}} is uniformly
distributed over $\{0,\ldots,2^{\ensuremath{\mathtt{{\color{var}d}}}}-1\}$.  With a little work, one
can even verify that any pair of values have independent hash values.
This implies tabulation hashing could be used in place of multiplicative
hashing for the \texttt{ChainedHashTable} implementation.

However, it is not true that any set of \ensuremath{\mathtt{{\color{var}n}}} distinct values gives a set
of \ensuremath{\mathtt{{\color{var}n}}} independent hash values.  Nevertheless, when tabulation hashing is
used, the bound of \thmref{linear-probing} still holds.  References for
this are provided at the end of this chapter.

\section{Hash Codes}

\index{hash code}%
The hash tables discussed in the previous section are used to associate
data with integer keys consisting of \ensuremath{\mathtt{{\color{var}w}}} bits.  In many cases, we
have keys that are not integers.  They may be strings, objects, arrays,
or other compound structures.  To use hash tables for these types of data,
we must map these data types to \ensuremath{\mathtt{{\color{var}w}}}-bit hash codes.  Hash code mappings should
have the following properties:

\begin{enumerate}
  \item If \ensuremath{\mathtt{{\color{var}x}}} and \ensuremath{\mathtt{{\color{var}y}}} are equal, then \ensuremath{\mathtt{{\color{var}x}.hashCode()}} and \ensuremath{\mathtt{{\color{var}y}.hashCode()}}
  are equal.

  \item If \ensuremath{\mathtt{{\color{var}x}}} and \ensuremath{\mathtt{{\color{var}y}}} are not equal, then the probability that
  $\ensuremath{\mathtt{{\color{var}x}.hashCode()}}=\ensuremath{\mathtt{{\color{var}y}.hashCode()}}$ should be small (close to
  $1/2^{\ensuremath{\mathtt{{\color{var}w}}}}$).
\end{enumerate}

The first property ensures that if we store \ensuremath{\mathtt{{\color{var}x}}} in a hash table and later
look up a value \ensuremath{\mathtt{{\color{var}y}}} equal to \ensuremath{\mathtt{{\color{var}x}}}, then we will find \ensuremath{\mathtt{{\color{var}x}}}---as we should.
The second property minimizes the loss from converting our objects
to integers.  It ensures that unequal objects usually have different
hash codes and so are likely to be stored at different locations in
our hash table.

\subsection{Hash Codes for Primitive Data Types}

\index{hash code!for primitive data}%
Small primitive data types like \ensuremath{\mathtt{{\color{keyword}char}}}, \ensuremath{\mathtt{{\color{keyword}byte}}}, \ensuremath{\mathtt{{\color{keyword}int}}}, and \ensuremath{\mathtt{{\color{keyword}float}}} are
usually easy to find hash codes for.  These data types always have a
binary representation and this binary representation usually consists of
\ensuremath{\mathtt{{\color{var}w}}} or fewer bits. \javaonly{(For example, in Java, \ensuremath{\mathtt{{\color{keyword}byte}}} is an 8-bit
type and \ensuremath{\mathtt{{\color{keyword}float}}} is a 32-bit type.)}\cpponly{(For example, in C++ \ensuremath{\mathtt{{\color{keyword}char}}}
is typically an 8-bit type and \ensuremath{\mathtt{{\color{keyword}float}}} is a 32-bit type.)} In these
cases, we just treat these bits as the representation of an integer in
the range $\{0,\ldots,2^\ensuremath{\mathtt{{\color{var}w}}}-1\}$.  If two values are different, they get
different hash codes.  If they are the same, they get the same hash code.

A few primitive data types are made up of more than \ensuremath{\mathtt{{\color{var}w}}} bits, usually
$c\ensuremath{\mathtt{{\color{var}w}}}$ bits for some constant integer $c$. (Java's \ensuremath{\mathtt{{\color{keyword}long}}} and \ensuremath{\mathtt{{\color{keyword}double}}}
types are examples of this with $c=2$.)  These data types can be treated
as compound objects made of $c$ parts, as described in the next section.

\subsection{Hash Codes for Compound Objects}
\seclabel{stringhash}

\index{hash code!for compound objects}%
For a compound object, we want to create a hash code by combining the
individual hash codes of the object's constituent parts.  This is not
as easy as it sounds.  Although one can find many hacks for this (for
example, combining the hash codes with bitwise exclusive-or operations),
many of these hacks turn out to be easy to foil (see Exercises~\ref{exc:hash-hack-first}--\ref{exc:hash-hack-last}).
However, if one is willing to do arithmetic with $2\ensuremath{\mathtt{{\color{var}w}}}$ bits of
precision, then there are simple and robust methods available.
Suppose we have an object made up of several parts
$P_0,\ldots,P_{r-1}$ whose hash codes are $\ensuremath{\mathtt{{\color{var}x}}}_0,\ldots,\ensuremath{\mathtt{{\color{var}x}}}_{r-1}$.
Then we can choose mutually independent random \ensuremath{\mathtt{{\color{var}w}}}-bit integers
$\ensuremath{\mathtt{{\color{var}z}}}_0,\ldots,\ensuremath{\mathtt{{\color{var}z}}}_{r-1}$ and a random $2\ensuremath{\mathtt{{\color{var}w}}}$-bit odd integer \ensuremath{\mathtt{{\color{var}z}}} and
compute a hash code for our object with
\[
   h(\ensuremath{\mathtt{{\color{var}x}}}_0,\ldots,\ensuremath{\mathtt{{\color{var}x}}}_{r-1}) =  
   \left(\left(\ensuremath{\mathtt{{\color{var}z}}}\sum_{i=0}^{r-1} \ensuremath{\mathtt{{\color{var}z}}}_i \ensuremath{\mathtt{{\color{var}x}}}_i\right)\bmod 2^{2\ensuremath{\mathtt{{\color{var}w}}}}\right)
   \ddiv 2^{\ensuremath{\mathtt{{\color{var}w}}}} \enspace .
\]
Note that this hash code has a final step (multiplying by \ensuremath{\mathtt{{\color{var}z}}} and
dividing by $2^{\ensuremath{\mathtt{{\color{var}w}}}}$) that uses the multiplicative hash function
from \secref{multihash} to take the $2\ensuremath{\mathtt{{\color{var}w}}}$-bit intermediate result and
reduce it to a \ensuremath{\mathtt{{\color{var}w}}}-bit final result.  Here is an example of this method applied to a simple compound object with three parts \ensuremath{\mathtt{{\color{var}x0}}}, \ensuremath{\mathtt{{\color{var}x1}}}, and \ensuremath{\mathtt{{\color{var}x2}}}:


 \ \vspace{.5ex}\begin{Verbatim}[gobble=1,tabsize=2,frame=single,framerule=1.2pt,commandchars=\\@\$,label=\texttt{Point3D},labelposition=topline]
	@\color@keyword$int$ hashCode() {
		@\color@comment$// random numbers from rand.org$
		@\color@keyword$long$[] @\color@var$z$ = {0x2058cc50L, 0xcb19137eL, 0x2cb6b6fdL}; 
		@\color@keyword$long$ @\color@var$zz$ = 0xbea0107e5067d19dL;

		@\color@comment$// convert (unsigned) hashcodes to long$
		@\color@keyword$long$ @\color@var$h0$ = @\color@var$x0$.hashCode() & ((1L<<32)-1);
		@\color@keyword$long$ @\color@var$h1$ = @\color@var$x1$.hashCode() & ((1L<<32)-1);
		@\color@keyword$long$ @\color@var$h2$ = @\color@var$x2$.hashCode() & ((1L<<32)-1);
		
		return (@\color@keyword$int$)(((@\color@var$z$[0]*@\color@var$h0$ + @\color@var$z$[1]*@\color@var$h1$ + @\color@var$z$[2]*@\color@var$h2$)*@\color@var$zz$)
		             >>> 32);
	}
\end{Verbatim}

\pcodeimport{ods/Point3d.hashCode()}
The following theorem shows that, in addition to being straightforward to implement, this method is provably good:

\begin{thm}\thmlabel{multihash}
Let $\ensuremath{\mathtt{{\color{var}x}}}_0,\ldots,\ensuremath{\mathtt{{\color{var}x}}}_{r-1}$ and $\ensuremath{\mathtt{{\color{var}y}}}_0,\ldots,\ensuremath{\mathtt{{\color{var}y}}}_{r-1}$ each be sequences of \ensuremath{\mathtt{{\color{var}w}}} bit integers in $\{0,\ldots,2^{\ensuremath{\mathtt{{\color{var}w}}}}-1\}$ and assume $\ensuremath{\mathtt{{\color{var}x}}}_i \neq \ensuremath{\mathtt{{\color{var}y}}}_i$ for at least one index $i\in\{0,\ldots,r-1\}$. Then 
\[
   \Pr\{ h(\ensuremath{\mathtt{{\color{var}x}}}_0,\ldots,\ensuremath{\mathtt{{\color{var}x}}}_{r-1}) =  h(\ensuremath{\mathtt{{\color{var}y}}}_0,\ldots,\ensuremath{\mathtt{{\color{var}y}}}_{r-1}) \} 
        \le 3/2^{\ensuremath{\mathtt{{\color{var}w}}}} \enspace .  
\] 
\end{thm}

\begin{proof}
  We will first ignore the final multiplicative hashing step and see how
  that step contributes later.  Define:
  \[
    h'(\ensuremath{\mathtt{{\color{var}x}}}_0,\ldots,\ensuremath{\mathtt{{\color{var}x}}}_{r-1}) =  
       \left(\sum_{j=0}^{r-1} \ensuremath{\mathtt{{\color{var}z}}}_j \ensuremath{\mathtt{{\color{var}x}}}_j\right)\bmod 2^{2\ensuremath{\mathtt{{\color{var}w}}}} \enspace .
  \]
  Suppose that $h'(\ensuremath{\mathtt{{\color{var}x}}}_0,\ldots,\ensuremath{\mathtt{{\color{var}x}}}_{r-1}) =  h'(\ensuremath{\mathtt{{\color{var}y}}}_0,\ldots,\ensuremath{\mathtt{{\color{var}y}}}_{r-1})$.
  We can rewrite this as:
  \begin{equation}  \eqlabel{bighash-x}
      \ensuremath{\mathtt{{\color{var}z}}}_i(\ensuremath{\mathtt{{\color{var}x}}}_i-\ensuremath{\mathtt{{\color{var}y}}}_i) \bmod 2^{2\ensuremath{\mathtt{{\color{var}w}}}} = t
  \end{equation}
  where 
  \[
     t = \left(\sum_{j=0}^{i-1} \ensuremath{\mathtt{{\color{var}z}}}_j(\ensuremath{\mathtt{{\color{var}y}}}_j-\ensuremath{\mathtt{{\color{var}x}}}_j) + \sum_{j=i+1}^{r-1} \ensuremath{\mathtt{{\color{var}z}}}_j(\ensuremath{\mathtt{{\color{var}y}}}_j-\ensuremath{\mathtt{{\color{var}x}}}_j)\right) \bmod 2^{2\ensuremath{\mathtt{{\color{var}w}}}}
  \]
  If we assume, without loss of generality that $\ensuremath{\mathtt{{\color{var}x}}}_i> \ensuremath{\mathtt{{\color{var}y}}}_i$, then
  \myeqref{bighash-x} becomes
  \begin{equation}
      \ensuremath{\mathtt{{\color{var}z}}}_i(\ensuremath{\mathtt{{\color{var}x}}}_i-\ensuremath{\mathtt{{\color{var}y}}}_i) = t \eqlabel{bighash-xx} \enspace ,
  \end{equation}
  since each of $\ensuremath{\mathtt{{\color{var}z}}}_i$ and $(\ensuremath{\mathtt{{\color{var}x}}}_i-\ensuremath{\mathtt{{\color{var}y}}}_i)$ is at most $2^{\ensuremath{\mathtt{{\color{var}w}}}}-1$, so
  their product is at most $2^{2\ensuremath{\mathtt{{\color{var}w}}}}-2^{\ensuremath{\mathtt{{\color{var}w}}}+1}+1 < 2^{2\ensuremath{\mathtt{{\color{var}w}}}}-1$.
  By assumption, $\ensuremath{\mathtt{{\color{var}x}}}_i-\ensuremath{\mathtt{{\color{var}y}}}_i\neq 0$, so \myeqref{bighash-xx} has at most one
  solution in $\ensuremath{\mathtt{{\color{var}z}}}_i$.  Therefore, since $\ensuremath{\mathtt{{\color{var}z}}}_i$ and $t$ are
  independent ($\ensuremath{\mathtt{{\color{var}z}}}_0,\ldots,\ensuremath{\mathtt{{\color{var}z}}}_{r-1}$ are mutually independent), the
  probability that we select $\ensuremath{\mathtt{{\color{var}z}}}_i$
  so that $h'(\ensuremath{\mathtt{{\color{var}x}}}_0,\ldots,\ensuremath{\mathtt{{\color{var}x}}}_{r-1})=h'(\ensuremath{\mathtt{{\color{var}y}}}_0,\ldots,\ensuremath{\mathtt{{\color{var}y}}}_{r-1})$ is at most
  $1/2^{\ensuremath{\mathtt{{\color{var}w}}}}$.

  The final step of the hash function is to apply multiplicative hashing
  to reduce our $2\ensuremath{\mathtt{{\color{var}w}}}$-bit intermediate result $h'(\ensuremath{\mathtt{{\color{var}x}}}_0,\ldots,\ensuremath{\mathtt{{\color{var}x}}}_{r-1})$ to
  a \ensuremath{\mathtt{{\color{var}w}}}-bit final result $h(\ensuremath{\mathtt{{\color{var}x}}}_0,\ldots,\ensuremath{\mathtt{{\color{var}x}}}_{r-1})$.  By \thmref{multihash},
  if $h'(\ensuremath{\mathtt{{\color{var}x}}}_0,\ldots,\ensuremath{\mathtt{{\color{var}x}}}_{r-1})\neq h'(\ensuremath{\mathtt{{\color{var}y}}}_0,\ldots,\ensuremath{\mathtt{{\color{var}y}}}_{r-1})$, then
  $\Pr\{h(\ensuremath{\mathtt{{\color{var}x}}}_0,\ldots,\ensuremath{\mathtt{{\color{var}x}}}_{r-1}) = h(\ensuremath{\mathtt{{\color{var}y}}}_0,\ldots,\ensuremath{\mathtt{{\color{var}y}}}_{r-1})\} \le 2/2^{\ensuremath{\mathtt{{\color{var}w}}}}$.

  To summarize, 
  \begin{align*}
    & \Pr\left\{\begin{array}{l}
          h(\ensuremath{\mathtt{{\color{var}x}}}_0,\ldots,\ensuremath{\mathtt{{\color{var}x}}}_{r-1}) \\
          \quad = h(\ensuremath{\mathtt{{\color{var}y}}}_0,\ldots,\ensuremath{\mathtt{{\color{var}y}}}_{r-1})\end{array}\right\} \\
      &= \Pr\left\{\begin{array}{ll}
            \mbox{$h'(\ensuremath{\mathtt{{\color{var}x}}}_0,\ldots,\ensuremath{\mathtt{{\color{var}x}}}_{r-1}) = h'(\ensuremath{\mathtt{{\color{var}y}}}_0,\ldots,\ensuremath{\mathtt{{\color{var}y}}}_{r-1})$ or} \\
            \mbox{$h'(\ensuremath{\mathtt{{\color{var}x}}}_0,\ldots,\ensuremath{\mathtt{{\color{var}x}}}_{r-1}) \neq h'(\ensuremath{\mathtt{{\color{var}y}}}_0,\ldots,\ensuremath{\mathtt{{\color{var}y}}}_{r-1})$} \\
                 \quad  \mbox{and
$\ensuremath{\mathtt{{\color{var}z}}}h'(\ensuremath{\mathtt{{\color{var}x}}}_0,\ldots,\ensuremath{\mathtt{{\color{var}x}}}_{r-1})\ddiv2^{\ensuremath{\mathtt{{\color{var}w}}}} = \ensuremath{\mathtt{{\color{var}z}}} h'(\ensuremath{\mathtt{{\color{var}y}}}_0,\ldots,\ensuremath{\mathtt{{\color{var}y}}}_{r-1})\ddiv 2^{\ensuremath{\mathtt{{\color{var}w}}}}$}
          \end{array}\right\} \\
      &\le 1/2^{\ensuremath{\mathtt{{\color{var}w}}}} + 2/2^{\ensuremath{\mathtt{{\color{var}w}}}} = 3/2^{\ensuremath{\mathtt{{\color{var}w}}}} \enspace . \qedhere
  \end{align*}
\end{proof}


\index{hash code!for strings}%
\index{hash code!for arrays}%
\subsection{Hash Codes for Arrays and Strings}
\seclabel{polyhash}

The method from the previous section works well for objects that have a
fixed, constant, number of components.  However, it breaks down when we
want to use it with objects that have a variable number of components,
since it requires a random \ensuremath{\mathtt{{\color{var}w}}}-bit integer $\ensuremath{\mathtt{{\color{var}z}}}_i$ for each component.
We could use a pseudorandom sequence to generate as many $\ensuremath{\mathtt{{\color{var}z}}}_i$'s as we
need, but then the $\ensuremath{\mathtt{{\color{var}z}}}_i$'s are not mutually independent, and it becomes
difficult to prove that the pseudorandom numbers don't interact badly
with the hash function we are using.  In particular, the values of $t$
and $\ensuremath{\mathtt{{\color{var}z}}}_i$ in the proof of \thmref{multihash} are no longer independent.

\index{prime field}%
A more rigorous approach is to base our hash codes on polynomials over
prime fields;  these are just regular polynomials that are evaluated
modulo some prime number, \ensuremath{\mathtt{{\color{var}p}}}. This method is based on the following
theorem, which says that polynomials over prime fields behave pretty-much
like usual polynomials:

\begin{thm}\thmlabel{prime-polynomial}
 Let $\ensuremath{\mathtt{{\color{var}p}}}$ be a prime number, and let $f(\ensuremath{\mathtt{{\color{var}z}}}) = \ensuremath{\mathtt{{\color{var}x}}}_0\ensuremath{\mathtt{{\color{var}z}}}^0 + \ensuremath{\mathtt{{\color{var}x}}}_1\ensuremath{\mathtt{{\color{var}z}}}^1 +
 \cdots + \ensuremath{\mathtt{{\color{var}x}}}_{r-1}\ensuremath{\mathtt{{\color{var}z}}}^{r-1}$ be a non-trivial polynomial with coefficients
 $\ensuremath{\mathtt{{\color{var}x}}}_i\in\{0,\ldots,\ensuremath{\mathtt{{\color{var}p}}}-1\}$. Then the equation $f(\ensuremath{\mathtt{{\color{var}z}}})\bmod \ensuremath{\mathtt{{\color{var}p}}} = 0$
 has at most $r-1$ solutions for $\ensuremath{\mathtt{{\color{var}z}}}\in\{0,\ldots,p-1\}$.
\end{thm}

To use \thmref{prime-polynomial}, we hash a sequence of integers
$\ensuremath{\mathtt{{\color{var}x}}}_0,\ldots,\ensuremath{\mathtt{{\color{var}x}}}_{r-1}$ with each $\ensuremath{\mathtt{{\color{var}x}}}_i\in \{0,\ldots,\ensuremath{\mathtt{{\color{var}p}}}-2\}$ using
a random integer $\ensuremath{\mathtt{{\color{var}z}}}\in\{0,\ldots,\ensuremath{\mathtt{{\color{var}p}}}-1\}$ via the formula
\[
   h(\ensuremath{\mathtt{{\color{var}x}}}_0,\ldots,\ensuremath{\mathtt{{\color{var}x}}}_{r-1}) 
    = \left(\ensuremath{\mathtt{{\color{var}x}}}_0\ensuremath{\mathtt{{\color{var}z}}}^0+\cdots+\ensuremath{\mathtt{{\color{var}x}}}_{r-1}\ensuremath{\mathtt{{\color{var}z}}}^{r-1}+(\ensuremath{\mathtt{{\color{var}p}}}-1)\ensuremath{\mathtt{{\color{var}z}}}^r \right)\bmod \ensuremath{\mathtt{{\color{var}p}}} \enspace .
\]

Note the extra $(\ensuremath{\mathtt{{\color{var}p}}}-1)\ensuremath{\mathtt{{\color{var}z}}}^r$ term at the end of the formula.  It helps
to think of $(\ensuremath{\mathtt{{\color{var}p}}}-1)$ as the last element, $\ensuremath{\mathtt{{\color{var}x}}}_r$, in the sequence
$\ensuremath{\mathtt{{\color{var}x}}}_0,\ldots,\ensuremath{\mathtt{{\color{var}x}}}_{r}$.  Note that this element differs from every other
element in the sequence (each of which is in the set $\{0,\ldots,\ensuremath{\mathtt{{\color{var}p}}}-2\}$).
We can think of $\ensuremath{\mathtt{{\color{var}p}}}-1$ as an end-of-sequence marker.

The following theorem, which considers the case of two sequences of
the same length, shows that this hash function gives a good return for
the small amount of randomization needed to choose \ensuremath{\mathtt{{\color{var}z}}}:

\begin{thm}\thmlabel{stringhash-eqlen}
  Let $\ensuremath{\mathtt{{\color{var}p}}}>2^{\ensuremath{\mathtt{{\color{var}w}}}}+1$ be a prime, let $\ensuremath{\mathtt{{\color{var}x}}}_0,\ldots,\ensuremath{\mathtt{{\color{var}x}}}_{r-1}$ and
  $\ensuremath{\mathtt{{\color{var}y}}}_0,\ldots,\ensuremath{\mathtt{{\color{var}y}}}_{r-1}$ each be sequences of \ensuremath{\mathtt{{\color{var}w}}}-bit integers in
  $\{0,\ldots,2^{\ensuremath{\mathtt{{\color{var}w}}}}-1\}$, and assume $\ensuremath{\mathtt{{\color{var}x}}}_i \neq \ensuremath{\mathtt{{\color{var}y}}}_i$ for at least
  one index $i\in\{0,\ldots,r-1\}$. Then
  \[
     \Pr\{ h(\ensuremath{\mathtt{{\color{var}x}}}_0,\ldots,\ensuremath{\mathtt{{\color{var}x}}}_{r-1}) =  h(\ensuremath{\mathtt{{\color{var}y}}}_0,\ldots,\ensuremath{\mathtt{{\color{var}y}}}_{r-1}) \} 
          \le (r-1)/\ensuremath{\mathtt{{\color{var}p}}} \enspace .
  \] 
\end{thm}

\begin{proof}
  The equation $h(\ensuremath{\mathtt{{\color{var}x}}}_0,\ldots,\ensuremath{\mathtt{{\color{var}x}}}_{r-1}) =  h(\ensuremath{\mathtt{{\color{var}y}}}_0,\ldots,\ensuremath{\mathtt{{\color{var}y}}}_{r-1})$
  can be rewritten as
  \begin{equation}  \eqlabel{strhash-eqlen}
    \left(
       (\ensuremath{\mathtt{{\color{var}x}}}_0-\ensuremath{\mathtt{{\color{var}y}}}_0)\ensuremath{\mathtt{{\color{var}z}}}^0+\cdots+(\ensuremath{\mathtt{{\color{var}x}}}_{r-1}-\ensuremath{\mathtt{{\color{var}y}}}_{r-1})\ensuremath{\mathtt{{\color{var}z}}}^{r-1} 
    \right)\bmod \ensuremath{\mathtt{{\color{var}p}}} = 0.
  \end{equation}
  Since $\ensuremath{\mathtt{{\color{var}x}}}_\ensuremath{\mathtt{{\color{var}i}}}\neq \ensuremath{\mathtt{{\color{var}y}}}_\ensuremath{\mathtt{{\color{var}i}}}$, this polynomial is non-trivial.  Therefore,
  by \thmref{prime-polynomial}, it has at most $r-1$ solutions in \ensuremath{\mathtt{{\color{var}z}}}.
  The probability that we pick \ensuremath{\mathtt{{\color{var}z}}} to be one of these solutions is therefore
  at most $(r-1)/\ensuremath{\mathtt{{\color{var}p}}}$.
\end{proof}

Note that this hash function also deals with the case in which two
sequences have different lengths, even when one of the sequences is a
prefix of the other. This is because this function effectively hashes
the infinite sequence
\[
  \ensuremath{\mathtt{{\color{var}x}}}_0,\ldots,\ensuremath{\mathtt{{\color{var}x}}}_{r-1}, \ensuremath{\mathtt{{\color{var}p}}}-1,0,0,\ldots \enspace .
\]
This guarantees that if we have two sequences of length $r$ and $r'$
with $r > r'$, then these two sequences differ at index $i=r$.  In
this case, \myeqref{strhash-eqlen} becomes
\[
  \left(
     \sum_{i=0}^{i=r'-1}(\ensuremath{\mathtt{{\color{var}x}}}_i-\ensuremath{\mathtt{{\color{var}y}}}_i)\ensuremath{\mathtt{{\color{var}z}}}^i + (\ensuremath{\mathtt{{\color{var}x}}}_{r'} - \ensuremath{\mathtt{{\color{var}p}}} + 1)\ensuremath{\mathtt{{\color{var}z}}}^{r'}
     +\sum_{i=r'+1}^{i=r-1}\ensuremath{\mathtt{{\color{var}x}}}_i\ensuremath{\mathtt{{\color{var}z}}}^i + (\ensuremath{\mathtt{{\color{var}p}}}-1)\ensuremath{\mathtt{{\color{var}z}}}^{r}
  \right)\bmod \ensuremath{\mathtt{{\color{var}p}}} = 0 \enspace ,
\]
which, by \thmref{prime-polynomial}, has at most $r$ solutions in $\ensuremath{\mathtt{{\color{var}z}}}$.
This combined with \thmref{stringhash-eqlen} suffice to prove the
following more general theorem:

\begin{thm}\thmlabel{stringhash}
  Let $\ensuremath{\mathtt{{\color{var}p}}}>2^{\ensuremath{\mathtt{{\color{var}w}}}}+1$ be a prime, let $\ensuremath{\mathtt{{\color{var}x}}}_0,\ldots,\ensuremath{\mathtt{{\color{var}x}}}_{r-1}$ and
  $\ensuremath{\mathtt{{\color{var}y}}}_0,\ldots,\ensuremath{\mathtt{{\color{var}y}}}_{r'-1}$ be distinct sequences of \ensuremath{\mathtt{{\color{var}w}}}-bit integers in
  $\{0,\ldots,2^{\ensuremath{\mathtt{{\color{var}w}}}}-1\}$. Then
  \[
     \Pr\{ h(\ensuremath{\mathtt{{\color{var}x}}}_0,\ldots,\ensuremath{\mathtt{{\color{var}x}}}_{r-1}) =  h(\ensuremath{\mathtt{{\color{var}y}}}_0,\ldots,\ensuremath{\mathtt{{\color{var}y}}}_{r-1}) \} 
          \le \max\{r,r'\}/\ensuremath{\mathtt{{\color{var}p}}}  \enspace .  
  \] 
\end{thm}

The following example code shows how this hash function is applied to
an object that contains an array, \ensuremath{\mathtt{{\color{var}x}}}, of values:


 \ \vspace{.5ex}\begin{Verbatim}[gobble=1,tabsize=2,frame=single,framerule=1.2pt,commandchars=\\@\$,label=\texttt{GeomVector},labelposition=topline]
	@\color@keyword$int$ hashCode() {
		@\color@keyword$long$ @\color@var$p$ = (1L<<32)-5;   @\color@comment$// prime: 2^32 - 5$
		@\color@keyword$long$ @\color@var$z$ = 0x64b6055aL;  @\color@comment$// 32 bits from random.org$
		@\color@keyword$int$ @\color@var$z2$ = 0x5067d19d;   @\color@comment$// random odd 32 bit number$
		@\color@keyword$long$ @\color@var$s$ = 0;
		@\color@keyword$long$ @\color@var$zi$ = 1;
		@\color@keyword$for$ (@\color@keyword$int$ @\color@var$i$ = 0; @\color@var$i$ < @\color@var$x$.@\color@var$length$; @\color@var$i$++) {
			@\color@comment$// reduce to 31 bits$
			@\color@keyword$long$ @\color@var$xi$ = (@\color@var$x$[@\color@var$i$].hashCode() * @\color@var$z2$) >>> 1; 
			@\color@var$s$ = (@\color@var$s$ + @\color@var$zi$ * @\color@var$xi$) % @\color@var$p$;
			@\color@var$zi$ = (@\color@var$zi$ * @\color@var$z$) % @\color@var$p$;	
		}
		@\color@var$s$ = (@\color@var$s$ + @\color@var$zi$ * (@\color@var$p$-1)) % @\color@var$p$;
		return (@\color@keyword$int$)@\color@var$s$;
	}
\end{Verbatim}

\pcodeimport{ods/GeomVector.hashCode()}

The preceding code sacrifices some collision probability for implementation
convenience.  In particular, it applies the multiplicative hash function
from \secref{multihash}, with $\ensuremath{\mathtt{{\color{var}d}}}=31$ to reduce \ensuremath{\mathtt{{\color{var}x}[{\color{var}i}].hashCode()}} to a
31-bit value.  This is so that the additions and multiplications that are
done modulo the prime $\ensuremath{\mathtt{{\color{var}p}}}=2^{32}-5$ can be carried out using unsigned
63-bit arithmetic.  Thus the probability of two different sequences,
the longer of which has length $r$, having the same hash code is at most
\[
    2/2^{31} + r/(2^{32}-5)
\]
rather than the $r/(2^{32}-5)$ specified in \thmref{stringhash}.

\section{Discussion and Exercises}

Hash tables and hash codes represent an enormous and active field
of research that is just touched upon in this chapter.  The online
Bibliography on Hashing \cite{hashing}
\index{Bibliography on Hashing}%
contains nearly 2000 entries.

A variety of different hash table implementations exist. The one
described in \secref{hashtable} is known as \emph{hashing with chaining}
\index{hashing with chaining}%
(each array entry contains a chain (\texttt{List}) of elements).  Hashing with
chaining dates back to an internal IBM memorandum authored by H. P. Luhn
and dated January 1953. This memorandum also seems to be one of the
earliest references to linked lists.

\index{open addressing}%
An alternative to hashing with chaining is that used by \emph{open
addressing} schemes, where all data is stored directly in an
array.  These schemes include the \texttt{LinearHashTable} structure of
\secref{linearhashtable}. This idea was also proposed, independently, by
a group at IBM in the 1950s.  Open addressing schemes must deal with the
problem of \emph{collision resolution}: 
\index{collision resolution}%
the case where two values hash
to the same array location.  Different strategies exist for collision
resolution; these provide different performance guarantees and often
require more sophisticated hash functions than the ones described here.

Yet another category of hash table implementations are the so-called
\emph{perfect hashing} methods.
\index{perfect hashing}%
These are methods in which \ensuremath{\mathtt{find({\color{var}x})}}
operations take $O(1)$ time in the worst-case.  For static data sets,
this can be accomplished by finding \emph{perfect hash functions}
\index{perfect hash function}%
\index{hash function!perfect}%
for
the data; these are functions that map each piece of data to a unique
array location.  For data that changes over time, perfect hashing
methods include \emph{FKS two-level hash tables}
\index{two-level hash table}%
\index{hash table!two-level}%
\cite{fks84,dkkmrt94}
and \emph{cuckoo hashing} \cite{pr04}.
\index{cuckoo hashing}%
\index{hash table!cuckoo}%

The hash functions presented in this chapter are probably among the most
practical methods currently known that can be proven to work well for any
set of data.  Other provably good methods date back to the pioneering work
of Carter and Wegman who introduced the notion of \emph{universal hashing}
\index{universal hashing}%
\index{hashing!universal}%
and described several hash functions for different scenarios \cite{cw79}.
Tabulation hashing, described in \secref{tabulation}, is due to Carter
and Wegman \cite{cw79}, but its analysis, when applied to linear probing
(and several other hash table schemes) is due to P\v{a}tra\c{s}cu and
Thorup \cite{pt12}.

The idea of \emph{multiplicative hashing}
\index{multiplicative hashing}%
\index{hashing!multiplicative}%
is very old and seems to be
part of the hashing folklore \cite[Section~6.4]{k97v3}.  However, the
idea of choosing the multiplier \ensuremath{\mathtt{{\color{var}z}}} to be a random \emph{odd} number,
and the analysis in \secref{multihash} is due to Dietzfelbinger \etal\
\cite{dhkp97}.  This version of multiplicative hashing is one of the
simplest, but its collision probability of $2/2^{\ensuremath{\mathtt{{\color{var}d}}}}$ is a factor of two
larger than what one could expect with a random function from $2^{\ensuremath{\mathtt{{\color{var}w}}}}\to
2^{\ensuremath{\mathtt{{\color{var}d}}}}$.  The \emph{multiply-add hashing}
\index{hashing!multiply-add}%
\index{multiply-add hashing}%
method uses the function
\[
   h(\ensuremath{\mathtt{{\color{var}x}}}) = ((\ensuremath{\mathtt{{\color{var}z}}}\ensuremath{\mathtt{{\color{var}x}}} + b) \bmod 2^{\ensuremath{\mathtt{2w}}}) \ddiv 2^{\ensuremath{\mathtt{2w}}-\ensuremath{\mathtt{{\color{var}d}}}}
\]
where \ensuremath{\mathtt{{\color{var}z}}} and \ensuremath{\mathtt{{\color{var}b}}} are each randomly chosen from $\{0,\ldots,2^{\ensuremath{\mathtt{2w}}}-1\}$.
Multiply-add hashing has a collision probability of only $1/2^{\ensuremath{\mathtt{{\color{var}d}}}}$
\cite{d96}, but requires $2\ensuremath{\mathtt{{\color{var}w}}}$-bit precision arithmetic.

There are a number of methods of obtaining hash codes from fixed-length
sequences of \ensuremath{\mathtt{{\color{var}w}}}-bit integers.  One particularly fast method
\cite{bhkkr99} is the function
\[\begin{array}{l}
  h(\ensuremath{\mathtt{{\color{var}x}}}_0,\ldots,\ensuremath{\mathtt{{\color{var}x}}}_{r-1}) \\
   \quad = \left(\sum_{i=0}^{r/2-1} ((\ensuremath{\mathtt{{\color{var}x}}}_{2i}+\ensuremath{\mathtt{{\color{var}a}}}_{2i})\bmod 2^{\ensuremath{\mathtt{{\color{var}w}}}})((\ensuremath{\mathtt{{\color{var}x}}}_{2i+1}+\ensuremath{\mathtt{{\color{var}a}}}_{2i+1})\bmod 2^{\ensuremath{\mathtt{{\color{var}w}}}})\right) \bmod 2^{2\ensuremath{\mathtt{{\color{var}w}}}}
\end{array}
\]
where $r$ is even and $\ensuremath{\mathtt{{\color{var}a}}}_0,\ldots,\ensuremath{\mathtt{{\color{var}a}}}_{r-1}$ are randomly chosen from
$\{0,\ldots,2^{\ensuremath{\mathtt{{\color{var}w}}}}\}$. This yields a $2\ensuremath{\mathtt{{\color{var}w}}}$-bit hash code that has
collision probability $1/2^{\ensuremath{\mathtt{{\color{var}w}}}}$.  This can be reduced to a \ensuremath{\mathtt{{\color{var}w}}}-bit hash
code using multiplicative (or multiply-add) hashing. This method is fast
because it requires only $r/2$ $2\ensuremath{\mathtt{{\color{var}w}}}$-bit multiplications whereas the
method described in \secref{stringhash} requires $r$ multiplications.
(The $\bmod$ operations occur implicitly by using \ensuremath{\mathtt{{\color{var}w}}} and $2\ensuremath{\mathtt{{\color{var}w}}}$-bit
arithmetic for the additions and multiplications, respectively.)

The method from \secref{polyhash} of using polynomials over prime fields
to hash variable-length arrays and strings is due to Dietzfelbinger \etal\
\cite{dgmp92}.  Due to its use of the $\bmod$ operator which relies
on a costly machine instruction, it is, unfortunately, not very fast.
Some variants of this method choose the prime \ensuremath{\mathtt{{\color{var}p}}} to be one of the form
$2^{\ensuremath{\mathtt{{\color{var}w}}}}-1$, in which case the $\bmod$ operator can be replaced with
addition (\ensuremath{\mathtt{+}}) and bitwise-and (\ensuremath{\mathtt{\text{\ttfamily \&}}}) operations \cite[Section~3.6]{k97v2}.
Another option is to apply one of the fast methods for fixed-length
strings to blocks of length $c$ for some constant $c>1$ and then apply
the prime field method to the resulting sequence of $\lceil r/c\rceil$
hash codes.


\begin{exc}
  A certain university assigns each of its students student numbers the
  first time they register for any course.  These numbers are sequential
  integers that started at 0 many years ago and are now in the millions.
  Suppose we have a class of one hundred first year students and we want
  to assign them hash codes based on their student numbers.  Does it
  make more sense to use the first two digits or the last two digits of
  their student number?  Justify your answer.
\end{exc}

\begin{exc}
  Consider the hashing scheme in \secref{multihash}, and suppose
  $\ensuremath{\mathtt{{\color{var}n}}}=2^{\ensuremath{\mathtt{{\color{var}d}}}}$ and $\ensuremath{\mathtt{{\color{var}d}}}\le \ensuremath{\mathtt{{\color{var}w}}}/2$.
  \begin{enumerate}
    \item Show that, for any choice
      of the muliplier, \ensuremath{\mathtt{{\color{var}z}}}, there exists \ensuremath{\mathtt{{\color{var}n}}} values that all have
      the same hash code. (Hint: This is easy, and doesn't require any
      number theory.)
    \item Given the multiplier, \ensuremath{\mathtt{{\color{var}z}}}, describe \ensuremath{\mathtt{{\color{var}n}}} values that all
      have the same hash code.  (Hint: This is harder, and requires
      some basic number theory.)
  \end{enumerate}
\end{exc}

\begin{exc}
  Prove that the bound $2/2^{\ensuremath{\mathtt{{\color{var}d}}}}$ in \lemref{universal-hashing} is
  the best possible bound by showing that, if $x=2^{\ensuremath{\mathtt{{\color{var}w}}}-\ensuremath{\mathtt{{\color{var}d}}}-2}$ and
  $\ensuremath{\mathtt{{\color{var}y}}}=3\ensuremath{\mathtt{{\color{var}x}}}$, then $\Pr\{\ensuremath{\mathtt{hash({\color{var}x})}}=\ensuremath{\mathtt{hash({\color{var}y})}}\}=2/2^{\ensuremath{\mathtt{{\color{var}d}}}}$.  (Hint look
  at the binary representations of $\ensuremath{\mathtt{{\color{var}zx}}}$ and $\ensuremath{\mathtt{{\color{var}z}}}3\ensuremath{\mathtt{{\color{var}x}}}$ and use the fact
  that $\ensuremath{\mathtt{{\color{var}z}}}3\ensuremath{\mathtt{{\color{var}x}}} = \ensuremath{\mathtt{{\color{var}z}}}x\ensuremath{\mathtt{+2}}z\ensuremath{\mathtt{{\color{var}x}}}$.)
\end{exc}

\begin{exc}\exclabel{linear-probing}
  Reprove \lemref{linear-probing} using the full version of Stirling's
  Approximation given in \secref{factorials}.
\end{exc}

\begin{exc}
  Consider the following simplified version of the code for adding
  an element \ensuremath{\mathtt{{\color{var}x}}} to a \texttt{LinearHashTable}, which simply stores \ensuremath{\mathtt{{\color{var}x}}} in the
  first \ensuremath{\mathtt{{\color{keyword}null}}} array entry it finds.  Explain why this could be very slow
  by giving an example of a sequence of $O(\ensuremath{\mathtt{{\color{var}n}}})$ \ensuremath{\mathtt{add({\color{var}x})}}, \ensuremath{\mathtt{remove({\color{var}x})}},
  and \ensuremath{\mathtt{find({\color{var}x})}} operations that would take on the order of $\ensuremath{\mathtt{{\color{var}n}}}^2$
  time to execute.


 \ \vspace{.5ex}\begin{Verbatim}[gobble=1,tabsize=2,frame=single,framerule=1.2pt,commandchars=\\@\$,label=\texttt{LinearHashTable},labelposition=topline]
	@\color@var$boolean$ addSlow(@\color@keyword$T$ @\color@var$x$) {
		@\color@keyword$if$ (2*(@\color@var$q$+1) > @\color@var$t$.@\color@var$length$) resize(); @\color@comment$// max 50% occupancy$
		@\color@keyword$int$ @\color@var$i$ = hash(@\color@var$x$);
		@\color@keyword$while$ (@\color@var$t$[@\color@var$i$] != @\color@keyword$null$) {
			@\color@keyword$if$ (@\color@var$t$[@\color@var$i$] != @\color@var$del$ && @\color@var$x$.equals(@\color@var$t$[@\color@var$i$])) @\color@var$return$ @\color@var$false$;
			@\color@var$i$ = (@\color@var$i$ == @\color@var$t$.@\color@var$length$-1) ? 0 : @\color@var$i$ + 1; @\color@comment$// increment i$
		}
		@\color@var$t$[@\color@var$i$] = @\color@var$x$;
		@\color@var$n$++; @\color@var$q$++;
		@\color@var$return$ @\color@var$true$;
	}
\end{Verbatim}

\end{exc}

\begin{exc}
  Early versions of the Java \ensuremath{\mathtt{hashCode()}} method for the \texttt{String} class
  worked by not using all of the characters found in long strings. For
  example, for a sixteen character string, the hash code was computed
  using only the eight even-indexed characters. Explain why this was a
  very bad idea by giving an example of large set of strings that all
  have the same hash code.
\end{exc}

\begin{exc}\exclabel{hash-hack-first}
  Suppose you have an object made up of two \ensuremath{\mathtt{{\color{var}w}}}-bit integers, \ensuremath{\mathtt{{\color{var}x}}} and \ensuremath{\mathtt{{\color{var}y}}}.
  Show why $\ensuremath{\mathtt{{\color{var}x}}}\oplus\ensuremath{\mathtt{{\color{var}y}}}$ does not make a good hash code for your object.
  Give an example of a large set of objects that would all have hash
  code 0.
\end{exc}

\begin{exc}
  Suppose you have an object made up of two \ensuremath{\mathtt{{\color{var}w}}}-bit integers, \ensuremath{\mathtt{{\color{var}x}}} and \ensuremath{\mathtt{{\color{var}y}}}.
  Show why $\ensuremath{\mathtt{{\color{var}x}}}+\ensuremath{\mathtt{{\color{var}y}}}$ does not make a good hash code for your object.
  Give an example of a large set of objects that would all have the same
  hash code.
\end{exc}

\begin{exc}\exclabel{hash-hack-last}
  Suppose you have an object made up of two \ensuremath{\mathtt{{\color{var}w}}}-bit integers, \ensuremath{\mathtt{{\color{var}x}}}
  and \ensuremath{\mathtt{{\color{var}y}}}.  Suppose that the hash code for your object is defined
  by some deterministic function $h(\ensuremath{\mathtt{{\color{var}x}}},\ensuremath{\mathtt{{\color{var}y}}})$ that produces a single
  \ensuremath{\mathtt{{\color{var}w}}}-bit integer.  Prove that there exists a large set of objects that
  have the same hash code.
\end{exc}

\begin{exc}
  Let $p=2^{\ensuremath{\mathtt{{\color{var}w}}}}-1$ for some positive integer \ensuremath{\mathtt{{\color{var}w}}}.  Explain why, for a
  positive integer $x$
  \[
      (x\bmod 2^{\ensuremath{\mathtt{{\color{var}w}}}}) + (x\ddiv 2^{\ensuremath{\mathtt{{\color{var}w}}}}) \equiv x \bmod (2^{\ensuremath{\mathtt{{\color{var}w}}}}-1) \enspace .
  \]
  \javaonly{
  (This gives an algorithm for computing $x \bmod (2^{\ensuremath{\mathtt{{\color{var}w}}}}-1)$ by
  repeatedly setting
  \javaonly{\[
    \ensuremath{\mathtt{{\color{var}x} = {\color{var}x}\text{\ttfamily \&}((1\text{\ttfamily <<}{\color{var}w})-1) + {\color{var}x}\text{\ttfamily >>>}{\color{var}w}}}
  \]}
  \cpponly{\[
    \ensuremath{\mathtt{{\color{var}x} = {\color{var}x}\text{\ttfamily \&}((1\text{\ttfamily <<}{\color{var}w})-1) + {\color{var}x}\text{\ttfamily >>}{\color{var}w}}}
  \]}
  until $\ensuremath{\mathtt{{\color{var}x}}} \le 2^{\ensuremath{\mathtt{{\color{var}w}}}}-1$.)}
\end{exc}

\begin{exc}
  Find some commonly used hash table implementation such as the \javaonly{Java
  Collection Framework \texttt{HashMap},}\cpponly{C++ STL \ensuremath{\mathtt{{\color{var}unordered}\_map}},}
  the \texttt{HashTable} or \texttt{LinearHashTable} implementations in this book,
  and design a program that stores integers in this data
  structure so that there are integers, \ensuremath{\mathtt{{\color{var}x}}}, such that \ensuremath{\mathtt{find({\color{var}x})}} takes
  linear time.  That is, find a set of \ensuremath{\mathtt{{\color{var}n}}} integers for which there are
  $c\ensuremath{\mathtt{{\color{var}n}}}$ elements that hash to the same table location.

  Depending on how good the implementation is, you may be able to do this
  just by inspecting the code for the implementation, or you may have to
  write some code that does trial insertions and searches, timing how long
  it takes to add and find particular values. (This can be, and has been,
  used to launch denial of service attacks on web servers \cite{cw03}.)
  \index{algorithmic complexity attack}%
\end{exc}

\clearpage

\clearpage
\chapter{Binary Trees}
\chaplabel{binarytrees}

This chapter introduces one of the most fundamental structures in computer
science: binary trees.  The use of the word \emph{tree}
\index{tree}%
\index{tree!binary}%
\index{binary tree}%
here comes from
the fact that, when we draw them, the resultant drawing often resembles
the trees found in a forest.  There are many ways of ways of defining
binary trees.  Mathematically, a \emph{binary tree} is a connected,
undirected, finite graph with no cycles, and no vertex of degree greater
than three.

For most computer science applications, binary trees are \emph{rooted:}
\index{tree!rooted}%
\index{rooted tree}%
A special node, \ensuremath{\mathtt{{\color{var}r}}}, of degree at most two is called the \emph{root}
of the tree.  For every node, $\ensuremath{\mathtt{{\color{var}u}}}\neq \ensuremath{\mathtt{{\color{var}r}}}$, the second node on the
path from \ensuremath{\mathtt{{\color{var}u}}} to \ensuremath{\mathtt{{\color{var}r}}} is called the \emph{parent} of \ensuremath{\mathtt{{\color{var}u}}}.
\index{parent}%
Each of the
other nodes adjacent to \ensuremath{\mathtt{{\color{var}u}}} is called a \emph{child}
\index{child} of \ensuremath{\mathtt{{\color{var}u}}}. Most of the
binary trees we are interested in are \emph{ordered},
\index{ordered tree}%
\index{tree!ordered}%
so we distinguish
between the \emph{left child} and \emph{right child} of \ensuremath{\mathtt{{\color{var}u}}}.
\index{left child}%
\index{child!left}%
\index{right child}%
\index{child!right}%

In illustrations, binary trees are usually drawn from the root
downward, with the root at the top of the drawing and the left and right
children respectively given by left and right positions in the drawing
(\figref{bintree-orientation}).  For example, \figref{binary-tree}.a shows
a binary tree with nine nodes.

\begin{figure}
  \begin{center}
    \includegraphics[scale=0.90909]{figs/bintree-traverse-1} 
  \end{center}
  \caption[Parent, left child, and right child]{The parent, left child, and right child of the node \ensuremath{\mathtt{{\color{var}u}}}
    in a \texttt{BinaryTree}.}
  \figlabel{bintree-orientation}
\end{figure}


\begin{figure}
  \begin{center}
    \begin{tabular}{cc}
      \includegraphics[width=\HalfScaleIfNeeded]{figs/bintree-1} &
      \includegraphics[width=\HalfScaleIfNeeded]{figs/bintree-2} \\
      (a) & (b)
    \end{tabular}
  \end{center}
  \caption{A binary tree with (a)~nine real nodes and (b)~ten external nodes.}
  \figlabel{binary-tree}
\end{figure}

Because binary trees are so important, a certain terminology has developed
for them: The \emph{depth}
\index{depth}%
of a node, \ensuremath{\mathtt{{\color{var}u}}}, in a binary tree is the
length of the path from \ensuremath{\mathtt{{\color{var}u}}} to the root of the tree.   If a node, \ensuremath{\mathtt{{\color{var}w}}},
is on the path from \ensuremath{\mathtt{{\color{var}u}}} to \ensuremath{\mathtt{{\color{var}r}}}, then \ensuremath{\mathtt{{\color{var}w}}} is called an \emph{ancestor}
\index{ancestor}%
of \ensuremath{\mathtt{{\color{var}u}}} and \ensuremath{\mathtt{{\color{var}u}}} a \emph{descendant}
\index{descendant}%
of \ensuremath{\mathtt{{\color{var}w}}}.  The \emph{subtree} of a
node, \ensuremath{\mathtt{{\color{var}u}}}, is the binary tree that is rooted at \ensuremath{\mathtt{{\color{var}u}}} and contains all
of \ensuremath{\mathtt{{\color{var}u}}}'s descendants.  The \emph{height}
\index{height!in a tree} of a node, \ensuremath{\mathtt{{\color{var}u}}}, is the length
of the longest path from \ensuremath{\mathtt{{\color{var}u}}} to one of its descendants.  The \emph{height} of
\index{height!of a tree}%
a tree is the height of its root.
A node, \ensuremath{\mathtt{{\color{var}u}}}, is a \emph{leaf}
\index{leaf}%
if it has no children.

We sometimes think of the tree as being augmented with \emph{external
nodes}. Any node that does not have a left child has an external
node as its left child, and, correspondingly, any node that does
not have a right child has an external node as its right child (see
\figref{binary-tree}.b).  It is easy to verify, by induction, that a
binary tree with $\ensuremath{\mathtt{{\color{var}n}}}\ge 1$ real nodes has $\ensuremath{\mathtt{{\color{var}n}}}+1$ external nodes.


\section{\texttt{BinaryTree}: A Basic Binary Tree}

\index{BinaryTree@\texttt{BinaryTree}}%
The simplest way to represent a node, \ensuremath{\mathtt{{\color{var}u}}}, in a binary tree is to
explicitly store the (at most three) neighbours of \ensuremath{\mathtt{{\color{var}u}}}\notpcode{:}\pcodeonly{.}


 \ \vspace{.5ex}\begin{Verbatim}[gobble=1,tabsize=2,frame=single,framerule=1.2pt,commandchars=\\@\$,label=\texttt{BinaryTree},labelposition=topline]
	@\color@keyword$class$ BTNode<Node @\color@keyword$extends$ BTNode<Node>> {
		Node @\color@var$left$;
		Node @\color@var$right$;
		Node @\color@var$parent$;	
	}
\end{Verbatim}

When one of these three neighbours is not present, we set it to \ensuremath{\mathtt{{\color{var}nil}}}.
In this way, both external nodes of the tree and the parent of the root
correspond to the value \ensuremath{\mathtt{{\color{var}nil}}}.

The binary tree itself can then be represented by a
\javaonly{reference}\cpponly{pointer}\pcodeonly{reference} to its root node, \ensuremath{\mathtt{{\color{var}r}}}:


 \ \vspace{.5ex}\begin{Verbatim}[gobble=1,tabsize=2,frame=single,framerule=1.2pt,commandchars=\\@\$,label=\texttt{BinaryTree},labelposition=topline]
	Node @\color@var$r$;
\end{Verbatim}


We can compute the depth of a node, \ensuremath{\mathtt{{\color{var}u}}}, in a binary tree by counting
the number of steps on the path from \ensuremath{\mathtt{{\color{var}u}}} to the root:


 \ \vspace{.5ex}\begin{Verbatim}[gobble=1,tabsize=2,frame=single,framerule=1.2pt,commandchars=\\@\$,label=\texttt{BinaryTree},labelposition=topline]
	@\color@keyword$int$ depth(Node @\color@var$u$) {
		@\color@keyword$int$ @\color@var$d$ = 0;
		@\color@keyword$while$ (@\color@var$u$ != @\color@var$r$) {
			@\color@var$u$ = @\color@var$u$.@\color@var$parent$;
			@\color@var$d$++;
		}
		@\color@var$return$ @\color@var$d$;
	}
\end{Verbatim}



\subsection{Recursive Algorithms}

\index{recursive algorithm}%
Using recursive algorithms makes it very easy to compute facts about
binary trees. For example, to compute the size of (number of nodes in)
a binary tree rooted at node \ensuremath{\mathtt{{\color{var}u}}}, we recursively compute the sizes of the
two subtrees rooted at the children of \ensuremath{\mathtt{{\color{var}u}}}, sum up these sizes, and add one:



 \ \vspace{.5ex}\begin{Verbatim}[gobble=1,tabsize=2,frame=single,framerule=1.2pt,commandchars=\\@\$,label=\texttt{BinaryTree},labelposition=topline]
	@\color@keyword$int$ size(Node @\color@var$u$) {
		@\color@keyword$if$ (@\color@var$u$ == @\color@var$nil$) @\color@var$return$ 0;
		@\color@var$return$ 1 + size(@\color@var$u$.@\color@var$left$) + size(@\color@var$u$.@\color@var$right$);
	}
\end{Verbatim}


To compute the height of a node \ensuremath{\mathtt{{\color{var}u}}}, we can compute the height of \ensuremath{\mathtt{{\color{var}u}}}'s
two subtrees, take the maximum, and add one:



 \ \vspace{.5ex}\begin{Verbatim}[gobble=1,tabsize=2,frame=single,framerule=1.2pt,commandchars=\\@\$,label=\texttt{BinaryTree},labelposition=topline]
	@\color@keyword$int$ height(Node @\color@var$u$) {
		@\color@keyword$if$ (@\color@var$u$ == @\color@var$nil$) @\color@var$return$ -1;
		@\color@var$return$ 1 + max(height(@\color@var$u$.@\color@var$left$), height(@\color@var$u$.@\color@var$right$));
	}
\end{Verbatim}


\subsection{Traversing Binary Trees}
\seclabel{bintree:traversal}

\index{traversal!of a binary tree}%
\index{tree traversal}%
\index{binary-tree traversal}%
The two algorithms from the previous section both use recursion to visit
all the nodes in a binary tree.  Each of them visits the nodes of the
binary tree in the same order as the following code:


 \ \vspace{.5ex}\begin{Verbatim}[gobble=1,tabsize=2,frame=single,framerule=1.2pt,commandchars=\\@\$,label=\texttt{BinaryTree},labelposition=topline]
	@\color@keyword$void$ traverse(Node @\color@var$u$) {
		@\color@keyword$if$ (@\color@var$u$ == @\color@var$nil$) @\color@var$return$;
		traverse(@\color@var$u$.@\color@var$left$);
		traverse(@\color@var$u$.@\color@var$right$);
	}
\end{Verbatim}


Using recursion this way produces very short, simple code, but it can
also be problematic.  The maximum depth of the recursion is given by
the maximum depth of a node in the binary tree, i.e., the tree's height.
If the height of the tree is very large, then this recursion could very
well use more stack space than is available, causing a crash.

To traverse a binary tree without recursion, you can use an algorithm that
relies on where it came from to determine where it will go next.  See
\figref{bintree-traverse}.  If we arrive at a node \ensuremath{\mathtt{{\color{var}u}}} from \ensuremath{\mathtt{{\color{var}u}.{\color{var}parent}}},
then the next thing to do is to visit \ensuremath{\mathtt{{\color{var}u}.{\color{var}left}}}.  If we arrive at \ensuremath{\mathtt{{\color{var}u}}}
from \ensuremath{\mathtt{{\color{var}u}.{\color{var}left}}}, then the next thing to do is to visit \ensuremath{\mathtt{{\color{var}u}.{\color{var}right}}}.  If we
arrive at \ensuremath{\mathtt{{\color{var}u}}} from \ensuremath{\mathtt{{\color{var}u}.{\color{var}right}}}, then we are done visiting \ensuremath{\mathtt{{\color{var}u}}}'s subtree,
and so we return to \ensuremath{\mathtt{{\color{var}u}.{\color{var}parent}}}.  The following code implements this
idea, with code included for handling the cases where any of \ensuremath{\mathtt{{\color{var}u}.{\color{var}left}}},
\ensuremath{\mathtt{{\color{var}u}.{\color{var}right}}}, or \ensuremath{\mathtt{{\color{var}u}.{\color{var}parent}}} is \ensuremath{\mathtt{{\color{var}nil}}}:


 \ \vspace{.5ex}\begin{Verbatim}[gobble=1,tabsize=2,frame=single,framerule=1.2pt,commandchars=\\@\$,label=\texttt{BinaryTree},labelposition=topline]
	@\color@keyword$void$ traverse2() {
		Node @\color@var$u$ = @\color@var$r$, @\color@var$prev$ = @\color@var$nil$, @\color@var$next$;
		@\color@keyword$while$ (@\color@var$u$ != @\color@var$nil$) {
			@\color@keyword$if$ (@\color@var$prev$ == @\color@var$u$.@\color@var$parent$) {
				@\color@keyword$if$ (@\color@var$u$.@\color@var$left$ != @\color@var$nil$) @\color@var$next$ = @\color@var$u$.@\color@var$left$;
				@\color@keyword$else$ @\color@keyword$if$ (@\color@var$u$.@\color@var$right$ != @\color@var$nil$) @\color@var$next$ = @\color@var$u$.@\color@var$right$;
				@\color@keyword$else$ @\color@var$next$ = @\color@var$u$.@\color@var$parent$;
			} @\color@keyword$else$ @\color@keyword$if$ (@\color@var$prev$ == @\color@var$u$.@\color@var$left$) {
				@\color@keyword$if$ (@\color@var$u$.@\color@var$right$ != @\color@var$nil$) @\color@var$next$ = @\color@var$u$.@\color@var$right$;
				@\color@keyword$else$ @\color@var$next$ = @\color@var$u$.@\color@var$parent$;
			} @\color@keyword$else$ {
				@\color@var$next$ = @\color@var$u$.@\color@var$parent$;
			}
			@\color@var$prev$ = @\color@var$u$;
			@\color@var$u$ = @\color@var$next$;
		}
	}
\end{Verbatim}


\begin{figure}
  \begin{center}
    \begin{tabular}{cc}
      \includegraphics[scale=0.90909]{figs/bintree-traverse-2}
      \includegraphics[scale=0.90909]{figs/bintree-3}
    \end{tabular}
  \end{center}
  \caption[Traversing a BinaryTree]{The three cases that occur at node
    \ensuremath{\mathtt{{\color{var}u}}} when traversing a binary tree non-recursively, and the resultant
    traversal of the tree.}
  \figlabel{bintree-traverse}
\end{figure}

The same facts that can be computed with recursive algorithms can also be
computed in this way, without recursion. For example, to compute the size
of the tree we keep a counter, \ensuremath{\mathtt{{\color{var}n}}}, and increment \ensuremath{\mathtt{{\color{var}n}}} whenever visiting
a node for the first time:


 \ \vspace{.5ex}\begin{Verbatim}[gobble=1,tabsize=2,frame=single,framerule=1.2pt,commandchars=\\@\$,label=\texttt{BinaryTree},labelposition=topline]
	@\color@keyword$int$ size2() {
		Node @\color@var$u$ = @\color@var$r$, @\color@var$prev$ = @\color@var$nil$, @\color@var$next$;
		@\color@keyword$int$ @\color@var$n$ = 0;
		@\color@keyword$while$ (@\color@var$u$ != @\color@var$nil$) {
			@\color@keyword$if$ (@\color@var$prev$ == @\color@var$u$.@\color@var$parent$) {
				@\color@var$n$++;
				@\color@keyword$if$ (@\color@var$u$.@\color@var$left$ != @\color@var$nil$) @\color@var$next$ = @\color@var$u$.@\color@var$left$;
				@\color@keyword$else$ @\color@keyword$if$ (@\color@var$u$.@\color@var$right$ != @\color@var$nil$) @\color@var$next$ = @\color@var$u$.@\color@var$right$;
				@\color@keyword$else$ @\color@var$next$ = @\color@var$u$.@\color@var$parent$;
			} @\color@keyword$else$ @\color@keyword$if$ (@\color@var$prev$ == @\color@var$u$.@\color@var$left$) {
				@\color@keyword$if$ (@\color@var$u$.@\color@var$right$ != @\color@var$nil$) @\color@var$next$ = @\color@var$u$.@\color@var$right$;
				@\color@keyword$else$ @\color@var$next$ = @\color@var$u$.@\color@var$parent$;
			} @\color@keyword$else$ {
				@\color@var$next$ = @\color@var$u$.@\color@var$parent$;
			}
			@\color@var$prev$ = @\color@var$u$;
			@\color@var$u$ = @\color@var$next$;
		}
		@\color@var$return$ @\color@var$n$;
	}
\end{Verbatim}


In some implementations of binary trees, the \ensuremath{\mathtt{{\color{var}parent}}} field is not used.
When this is the case, a non-recursive implementation is still possible,
but the implementation has to use a \texttt{List} (or \texttt{Stack}) to keep track
of the path from the current node to the root.

A special kind of traversal that does not fit the pattern of the above
functions is the \emph{breadth-first traversal}.
\index{breadth-first traversal}%
\index{traversal!breadth-first}%
In a breadth-first
traversal, the nodes are visited level-by-level starting at the root and
moving down, visiting the nodes at each level from left to right (see
\figref{bintree-bfs}). This is similar to the way that we would read a
page of English text.   Breadth-first traversal is implemented using a
queue, \ensuremath{\mathtt{{\color{var}q}}}, that initially contains only the root, \ensuremath{\mathtt{{\color{var}r}}}.  At each step,
we extract the next node, \ensuremath{\mathtt{{\color{var}u}}}, from \ensuremath{\mathtt{{\color{var}q}}}, process \ensuremath{\mathtt{{\color{var}u}}} and add \ensuremath{\mathtt{{\color{var}u}.{\color{var}left}}}
and \ensuremath{\mathtt{{\color{var}u}.{\color{var}right}}} (if they are non-\ensuremath{\mathtt{{\color{var}nil}}}) to \ensuremath{\mathtt{{\color{var}q}}}:


 \ \vspace{.5ex}\begin{Verbatim}[gobble=1,tabsize=2,frame=single,framerule=1.2pt,commandchars=\\@\$,label=\texttt{BinaryTree},labelposition=topline]
	@\color@keyword$void$ bfTraverse() {
		Queue<Node> @\color@var$q$ = @\color@keyword$new$ LinkedList<Node>();
		@\color@keyword$if$ (@\color@var$r$ != @\color@var$nil$) @\color@var$q$.add(@\color@var$r$);
		@\color@keyword$while$ (!@\color@var$q$.isEmpty()) {
			Node @\color@var$u$ = @\color@var$q$.remove();
			@\color@keyword$if$ (@\color@var$u$.@\color@var$left$ != @\color@var$nil$) @\color@var$q$.add(@\color@var$u$.@\color@var$left$);
			@\color@keyword$if$ (@\color@var$u$.@\color@var$right$ != @\color@var$nil$) @\color@var$q$.add(@\color@var$u$.@\color@var$right$);
		}
	}
\end{Verbatim}


\begin{figure}
  \begin{center}
    \includegraphics[scale=0.90909]{figs/bintree-4}
  \end{center}
  \caption{During a breadth-first traversal, the nodes of a binary tree
  are visited
level-by-level, and left-to-right within each level.}
  \figlabel{bintree-bfs}
\end{figure}





\section{\texttt{BinarySearchTree}: An Unbalanced Binary Search Tree}
\seclabel{binarysearchtree}

\index{BinarySearchTree@\texttt{BinarySearchTree}}%
\index{binary search tree}%
\index{binary tree!search}%
A \texttt{BinarySearchTree} is a special kind of binary tree in which each node, \ensuremath{\mathtt{{\color{var}u}}},
also stores a data value, \ensuremath{\mathtt{{\color{var}u}.{\color{var}x}}}, from some total order.  The data values in
a binary search tree obey the \emph{binary search tree property}:
\index{binary search tree property}%
For
a node, \ensuremath{\mathtt{{\color{var}u}}}, every data value stored in the subtree rooted at \ensuremath{\mathtt{{\color{var}u}.{\color{var}left}}}
is less than \ensuremath{\mathtt{{\color{var}u}.{\color{var}x}}} and every data value stored in the subtree rooted at
\ensuremath{\mathtt{{\color{var}u}.{\color{var}right}}} is greater than \ensuremath{\mathtt{{\color{var}u}.{\color{var}x}}}.  An example of a \texttt{BinarySearchTree} is shown in \figref{bst}.

\begin{figure}
  \begin{center}
    \includegraphics[scale=0.90909]{figs/bst-example}
  \end{center}
  \caption{A binary search tree.}
  \figlabel{bst}
\end{figure}


\subsection{Searching}

\index{search path!in a binary search tree}%
The binary search tree property is extremely useful because it allows
us to quickly locate a value, \ensuremath{\mathtt{{\color{var}x}}}, in a binary search tree.  To do this we start
searching for \ensuremath{\mathtt{{\color{var}x}}} at the root, \ensuremath{\mathtt{{\color{var}r}}}.  When examining a node, \ensuremath{\mathtt{{\color{var}u}}}, there
are three cases:
\begin{enumerate}
\item If $\ensuremath{\mathtt{{\color{var}x}}}< \ensuremath{\mathtt{{\color{var}u}.{\color{var}x}}}$, then the search proceeds to \ensuremath{\mathtt{{\color{var}u}.{\color{var}left}}};
\item If $\ensuremath{\mathtt{{\color{var}x}}}> \ensuremath{\mathtt{{\color{var}u}.{\color{var}x}}}$, then the search proceeds to \ensuremath{\mathtt{{\color{var}u}.{\color{var}right}}};
\item If $\ensuremath{\mathtt{{\color{var}x}}}= \ensuremath{\mathtt{{\color{var}u}.{\color{var}x}}}$, then we have found the node \ensuremath{\mathtt{{\color{var}u}}} containing \ensuremath{\mathtt{{\color{var}x}}}.
\end{enumerate}
The search terminates when Case~3 occurs or when \ensuremath{\mathtt{{\color{var}u}={\color{var}nil}}}.  In the
former case, we found \ensuremath{\mathtt{{\color{var}x}}}.  In the latter case, we conclude that \ensuremath{\mathtt{{\color{var}x}}}
is not in the binary search tree.


 \ \vspace{.5ex}\begin{Verbatim}[gobble=1,tabsize=2,frame=single,framerule=1.2pt,commandchars=\\@\$,label=\texttt{BinarySearchTree},labelposition=topline]
	@\color@keyword$T$ findEQ(@\color@keyword$T$ @\color@var$x$) {
		Node @\color@var$u$ = @\color@var$r$;
		@\color@keyword$while$ (@\color@var$u$ != @\color@var$nil$) {
			@\color@keyword$int$ @\color@var$comp$ = compare(@\color@var$x$, @\color@var$u$.@\color@var$x$);
			@\color@keyword$if$ (@\color@var$comp$ < 0) 
				@\color@var$u$ = @\color@var$u$.@\color@var$left$;
			@\color@keyword$else$ @\color@keyword$if$ (@\color@var$comp$ > 0)
				@\color@var$u$ = @\color@var$u$.@\color@var$right$;
			@\color@keyword$else$
				@\color@var$return$ @\color@var$u$.@\color@var$x$;
		}
		@\color@var$return$ @\color@keyword$null$;
	}
\end{Verbatim}


Two examples of searches in a binary search tree are shown in
\figref{bst-search}.  As the second example shows, even if we don't
find \ensuremath{\mathtt{{\color{var}x}}} in the tree, we still gain some valuable information.  If we
look at the last node, \ensuremath{\mathtt{{\color{var}u}}}, at which Case~1 occurred, we see that \ensuremath{\mathtt{{\color{var}u}.{\color{var}x}}}
is the smallest value in the tree that is greater than \ensuremath{\mathtt{{\color{var}x}}}.  Similarly,
the last node at which Case~2 occurred contains the largest value in the
tree that is less than \ensuremath{\mathtt{{\color{var}x}}}.  Therefore, by keeping track of the last
node, \ensuremath{\mathtt{{\color{var}z}}}, at which Case~1 occurs, a \texttt{BinarySearchTree} can implement
the \ensuremath{\mathtt{find({\color{var}x})}} operation that returns the smallest value stored in the
tree that is greater than or equal to \ensuremath{\mathtt{{\color{var}x}}}:


 \ \vspace{.5ex}\begin{Verbatim}[gobble=1,tabsize=2,frame=single,framerule=1.2pt,commandchars=\\@\$,label=\texttt{BinarySearchTree},labelposition=topline]
	@\color@keyword$T$ find(@\color@keyword$T$ @\color@var$x$) {
		Node @\color@var$w$ = @\color@var$r$, @\color@var$z$ = @\color@var$nil$;
		@\color@keyword$while$ (@\color@var$w$ != @\color@var$nil$) {
			@\color@keyword$int$ @\color@var$comp$ = compare(@\color@var$x$, @\color@var$w$.@\color@var$x$);
			@\color@keyword$if$ (@\color@var$comp$ < 0) {
				@\color@var$z$ = @\color@var$w$;
				@\color@var$w$ = @\color@var$w$.@\color@var$left$;
			} @\color@keyword$else$ @\color@keyword$if$ (@\color@var$comp$ > 0) {
				@\color@var$w$ = @\color@var$w$.@\color@var$right$;
			} @\color@keyword$else$ {
				@\color@var$return$ @\color@var$w$.@\color@var$x$;
			}
		}
		@\color@var$return$ @\color@var$z$ == @\color@var$nil$ ? @\color@keyword$null$ : @\color@var$z$.@\color@var$x$;
	}
\end{Verbatim}


\begin{figure}
  \begin{center}
    \begin{tabular}{cc}
    \includegraphics[width=\HalfScaleIfNeeded]{figs/bst-example-2} &
    \includegraphics[width=\HalfScaleIfNeeded]{figs/bst-example-3} \\
    (a) & (b)
    \end{tabular}
  \end{center}
  \caption{An example of (a)~a successful search (for $6$) and (b)~an unsuccessful search (for $10$) in a binary search tree.}
  \figlabel{bst-search}
\end{figure}


\subsection{Addition}

To add a new value, \ensuremath{\mathtt{{\color{var}x}}}, to a \texttt{BinarySearchTree}, we first search for
\ensuremath{\mathtt{{\color{var}x}}}. If we find it, then there is no need to insert it.  Otherwise,
we store \ensuremath{\mathtt{{\color{var}x}}} at a leaf child of the last node, \ensuremath{\mathtt{{\color{var}p}}}, encountered during the
search for \ensuremath{\mathtt{{\color{var}x}}}. Whether the new node is the left or right child of \ensuremath{\mathtt{{\color{var}p}}} depends on the result of comparing \ensuremath{\mathtt{{\color{var}x}}} and \ensuremath{\mathtt{{\color{var}p}.{\color{var}x}}}.


 \ \vspace{.5ex}\begin{Verbatim}[gobble=1,tabsize=2,frame=single,framerule=1.2pt,commandchars=\\@\$,label=\texttt{BinarySearchTree},labelposition=topline]
	@\color@var$boolean$ add(@\color@keyword$T$ @\color@var$x$) {
		Node @\color@var$p$ = findLast(@\color@var$x$);
		@\color@var$return$ addChild(@\color@var$p$, newNode(@\color@var$x$));		
	}
\end{Verbatim}



 \ \vspace{.5ex}\begin{Verbatim}[gobble=1,tabsize=2,frame=single,framerule=1.2pt,commandchars=\\@\$,label=\texttt{BinarySearchTree},labelposition=topline]
	Node findLast(@\color@keyword$T$ @\color@var$x$) {
		Node @\color@var$w$ = @\color@var$r$, @\color@var$prev$ = @\color@var$nil$;
		@\color@keyword$while$ (@\color@var$w$ != @\color@var$nil$) {
			@\color@var$prev$ = @\color@var$w$;
			@\color@keyword$int$ @\color@var$comp$ = compare(@\color@var$x$, @\color@var$w$.@\color@var$x$);
			@\color@keyword$if$ (@\color@var$comp$ < 0) {
				@\color@var$w$ = @\color@var$w$.@\color@var$left$;
			} @\color@keyword$else$ @\color@keyword$if$ (@\color@var$comp$ > 0) {
				@\color@var$w$ = @\color@var$w$.@\color@var$right$;
			} @\color@keyword$else$ {
				@\color@var$return$ @\color@var$w$;
			}
		}
		@\color@var$return$ @\color@var$prev$;
	}
\end{Verbatim}



 \ \vspace{.5ex}\begin{Verbatim}[gobble=1,tabsize=2,frame=single,framerule=1.2pt,commandchars=\\@\$,label=\texttt{BinarySearchTree},labelposition=topline]
	@\color@var$boolean$ addChild(Node @\color@var$p$, Node @\color@var$u$) {
		@\color@keyword$if$ (@\color@var$p$ == @\color@var$nil$) {
			@\color@var$r$ = @\color@var$u$;              @\color@comment$// inserting into empty tree$
		} @\color@keyword$else$ {
			@\color@keyword$int$ @\color@var$comp$ = compare(@\color@var$u$.@\color@var$x$, @\color@var$p$.@\color@var$x$);
			@\color@keyword$if$ (@\color@var$comp$ < 0) {
				@\color@var$p$.@\color@var$left$ = @\color@var$u$;
			} @\color@keyword$else$ @\color@keyword$if$ (@\color@var$comp$ > 0) {
				@\color@var$p$.@\color@var$right$ = @\color@var$u$;
			} @\color@keyword$else$ {
				@\color@var$return$ @\color@var$false$;   @\color@comment$// u.x is already in the tree$
			}
			@\color@var$u$.@\color@var$parent$ = @\color@var$p$;
		}
		@\color@var$n$++;
		@\color@var$return$ @\color@var$true$;		
	}
\end{Verbatim}

An example is shown in \figref{bst-insert}. The most time-consuming
part of this process is the initial search for \ensuremath{\mathtt{{\color{var}x}}}, which takes an
amount of time proportional to the height of the newly added node \ensuremath{\mathtt{{\color{var}u}}}.
In the worst case, this is equal to the height of the \texttt{BinarySearchTree}.


\begin{figure}
  \begin{center}
    \begin{tabular}{cc}
    \includegraphics[width=\HalfScaleIfNeeded]{figs/bst-example-4} &
    \includegraphics[width=\HalfScaleIfNeeded]{figs/bst-example-5} 
    \end{tabular}
  \end{center}
  \caption{Inserting the value $8.5$ into a binary search tree.}
  \figlabel{bst-insert}
\end{figure}


\subsection{Removal}

Deleting a value stored in a node, \ensuremath{\mathtt{{\color{var}u}}}, of a \texttt{BinarySearchTree} is a
little more difficult.  If \ensuremath{\mathtt{{\color{var}u}}} is a leaf, then we can just detach \ensuremath{\mathtt{{\color{var}u}}}
from its parent.  Even better: If \ensuremath{\mathtt{{\color{var}u}}} has only one child, then we can
splice \ensuremath{\mathtt{{\color{var}u}}} from the tree by having \ensuremath{\mathtt{{\color{var}u}.{\color{var}parent}}} adopt \ensuremath{\mathtt{{\color{var}u}}}'s child (see
\figref{bst-splice}):


 \ \vspace{.5ex}\begin{Verbatim}[gobble=1,tabsize=2,frame=single,framerule=1.2pt,commandchars=\\@\$,label=\texttt{BinarySearchTree},labelposition=topline]
	@\color@keyword$void$ splice(Node @\color@var$u$) {
		Node @\color@var$s$, @\color@var$p$;
		@\color@keyword$if$ (@\color@var$u$.@\color@var$left$ != @\color@var$nil$) {
			@\color@var$s$ = @\color@var$u$.@\color@var$left$;
		} @\color@keyword$else$ {
			@\color@var$s$ = @\color@var$u$.@\color@var$right$;
		}
		@\color@keyword$if$ (@\color@var$u$ == @\color@var$r$) {
			@\color@var$r$ = @\color@var$s$;
			@\color@var$p$ = @\color@var$nil$;
		} @\color@keyword$else$ {
			@\color@var$p$ = @\color@var$u$.@\color@var$parent$;
			@\color@keyword$if$ (@\color@var$p$.@\color@var$left$ == @\color@var$u$) {
				@\color@var$p$.@\color@var$left$ = @\color@var$s$;
			} @\color@keyword$else$ {
				@\color@var$p$.@\color@var$right$ = @\color@var$s$; 
			}
		}
		@\color@keyword$if$ (@\color@var$s$ != @\color@var$nil$) {
			@\color@var$s$.@\color@var$parent$ = @\color@var$p$;
		}
		@\color@var$n$--;
	}
\end{Verbatim}


\begin{figure}
  \begin{center}
    \includegraphics[scale=0.90909]{figs/bst-splice}
  \end{center}
  \caption{Removing a leaf ($6$) or a node with only one child ($9$) is easy.}
  \figlabel{bst-splice}
\end{figure}

Things get tricky, though, when \ensuremath{\mathtt{{\color{var}u}}} has two children.  In this case,
the simplest thing to do is to find a node, \ensuremath{\mathtt{{\color{var}w}}}, that has less than
two children and such that \ensuremath{\mathtt{{\color{var}w}.{\color{var}x}}} can replace \ensuremath{\mathtt{{\color{var}u}.{\color{var}x}}}.  To maintain
the binary search tree property, the value \ensuremath{\mathtt{{\color{var}w}.{\color{var}x}}} should be close to the
value of \ensuremath{\mathtt{{\color{var}u}.{\color{var}x}}}.  For example, choosing \ensuremath{\mathtt{{\color{var}w}}} such that \ensuremath{\mathtt{{\color{var}w}.{\color{var}x}}} is the smallest
value greater than \ensuremath{\mathtt{{\color{var}u}.{\color{var}x}}} will work.  Finding the node \ensuremath{\mathtt{{\color{var}w}}} is easy; it is
the smallest value in the subtree rooted at \ensuremath{\mathtt{{\color{var}u}.{\color{var}right}}}.  This node can
be easily removed because it has no left child (see \figref{bst-remove}).


 \ \vspace{.5ex}\begin{Verbatim}[gobble=1,tabsize=2,frame=single,framerule=1.2pt,commandchars=\\@\$,label=\texttt{BinarySearchTree},labelposition=topline]
	@\color@keyword$void$ remove(Node @\color@var$u$) {
		@\color@keyword$if$ (@\color@var$u$.@\color@var$left$ == @\color@var$nil$ || @\color@var$u$.@\color@var$right$ == @\color@var$nil$) {
			splice(@\color@var$u$);
		} @\color@keyword$else$ {
			Node @\color@var$w$ = @\color@var$u$.@\color@var$right$;
			@\color@keyword$while$ (@\color@var$w$.@\color@var$left$ != @\color@var$nil$) 
				@\color@var$w$ = @\color@var$w$.@\color@var$left$;
			@\color@var$u$.@\color@var$x$ = @\color@var$w$.@\color@var$x$;
			splice(@\color@var$w$);
		}
	}
\end{Verbatim}

\pcodeimport{ods/BinarySearchTree.remove_node(u)}

\begin{figure}
  \begin{center}
    \begin{tabular}{cc}
    \includegraphics[width=\HalfScaleIfNeeded]{figs/bst-delete-1}
    \includegraphics[width=\HalfScaleIfNeeded]{figs/bst-delete-2}
    \end{tabular}
  \end{center}
  \caption[Deleting from a BinarySearchTree]{Deleting a value ($11$) from a node, \ensuremath{\mathtt{{\color{var}u}}}, with two children is done by replacing \ensuremath{\mathtt{{\color{var}u}}}'s value with the smallest value in the right subtree of \ensuremath{\mathtt{{\color{var}u}}}.}
  \figlabel{bst-remove}
\end{figure}

\subsection{Summary}

The \ensuremath{\mathtt{find({\color{var}x})}}, \ensuremath{\mathtt{add({\color{var}x})}}, and \ensuremath{\mathtt{remove({\color{var}x})}} operations in a
\texttt{BinarySearchTree} each involve following a path from the root of the
tree to some node in the tree. Without knowing more about the shape of
the tree it is difficult to say much about the length of this path,
except that it is less than \ensuremath{\mathtt{{\color{var}n}}}, the number of nodes in the tree.
The following (unimpressive) theorem summarizes the performance of the
\texttt{BinarySearchTree} data structure:

\begin{thm}\thmlabel{bst}
  \texttt{BinarySearchTree} implements the \texttt{SSet} interface and 
  supports the operations \ensuremath{\mathtt{add({\color{var}x})}}, \ensuremath{\mathtt{remove({\color{var}x})}},
  and \ensuremath{\mathtt{find({\color{var}x})}} in $O(\ensuremath{\mathtt{{\color{var}n}}})$ time per operation.
\end{thm}

\thmref{bst} compares poorly with \thmref{skiplist}, which shows
that the \texttt{SkiplistSSet} structure can implement the \texttt{SSet} interface
with $O(\log \ensuremath{\mathtt{{\color{var}n}}})$ expected time per operation.  The problem with the
\texttt{BinarySearchTree} structure is that it can become \emph{unbalanced}.
Instead of looking like the tree in \figref{bst} it can look like a long
chain of \ensuremath{\mathtt{{\color{var}n}}} nodes, all but the last having exactly one child.

There are a number of ways of avoiding unbalanced binary search
trees, all of which lead to data structures that have $O(\log
\ensuremath{\mathtt{{\color{var}n}}})$ time operations. In \chapref{rbs} we show how $O(\log \ensuremath{\mathtt{{\color{var}n}}})$
\emph{expected} time operations can be achieved with randomization.
In \chapref{scapegoat} we show how $O(\log \ensuremath{\mathtt{{\color{var}n}}})$ \emph{amortized}
time operations can be achieved with partial rebuilding operations.
In \chapref{redblack} we show how $O(\log \ensuremath{\mathtt{{\color{var}n}}})$ \emph{worst-case}
time operations can be achieved by simulating a tree that is not binary:
one in which nodes can have up to four children.

\section{Discussion and Exercises}

Binary trees have been used to model relationships for thousands
of years.  One reason for this is that binary trees naturally model
(pedigree) family trees.
\index{family tree}%
\index{pedigree family tree}%
These are the family trees in which the root is
a person, the left and right children are the person's parents, and so
on, recursively.  In more recent centuries binary trees have also been
used to model species trees
\index{species tree} in biology, where the leaves of the tree
represent extant species and the internal nodes of the tree represent
\emph{speciation events}
\index{speciation event} in which two populations of a single species
evolve into two separate species.

Binary search trees appear to have been discovered independently by
several groups in the 1950s \cite[Section~6.2.2]{k97v3}.  Further
references to specific kinds of binary search trees are provided in
subsequent chapters.

When implementing a binary tree from scratch, there are several design
decisions to be made.  One of these is the question of whether or not
each node stores a pointer to its parent.  If most of the operations
simply follow a root-to-leaf path, then parent pointers are unnecessary,
waste space, and are a potential source of coding errors.  On the other
hand, the lack of parent pointers means that tree traversals must be done
recursively or with the use of an explicit stack.  Some other methods
(like inserting or deleting into some kinds of balanced binary search
trees) are also complicated by the lack of parent pointers.

Another design decision is concerned with how to store the parent, left
child, and right child pointers at a node.  In the implementation given
here, these pointers are stored as separate variables.   Another option
is to store them in an array, \ensuremath{\mathtt{{\color{var}p}}}, of length 3, so that \ensuremath{\mathtt{{\color{var}u}.{\color{var}p}[0]}} is the
left child of \ensuremath{\mathtt{{\color{var}u}}}, \ensuremath{\mathtt{{\color{var}u}.{\color{var}p}[1]}} is the right child of \ensuremath{\mathtt{{\color{var}u}}}, and \ensuremath{\mathtt{{\color{var}u}.{\color{var}p}[2]}} is
the parent of \ensuremath{\mathtt{{\color{var}u}}}.  Using an array this way means that some sequences
of \ensuremath{\mathtt{{\color{keyword}if}}} statements can be simplified into algebraic expressions.

An example of such a simplification occurs during tree traversal. If
a traversal arrives at a node \ensuremath{\mathtt{{\color{var}u}}} from \ensuremath{\mathtt{{\color{var}u}.{\color{var}p}[{\color{var}i}]}}, then the next node in
the traversal is $\ensuremath{\mathtt{{\color{var}u}.{\color{var}p}}}[(\ensuremath{\mathtt{{\color{var}i}}}+1)\bmod 3]$.  Similar examples occur when
there is left-right symmetry.  For example, the sibling of \ensuremath{\mathtt{{\color{var}u}.{\color{var}p}[{\color{var}i}]}} is
$\ensuremath{\mathtt{{\color{var}u}.{\color{var}p}}}[(\ensuremath{\mathtt{{\color{var}i}}}+1)\bmod 2]$.  This trick works whether \ensuremath{\mathtt{{\color{var}u}.{\color{var}p}[{\color{var}i}]}} is a left
child ($\ensuremath{\mathtt{{\color{var}i}}}=0$) or a right child ($\ensuremath{\mathtt{{\color{var}i}}}=1$) of \ensuremath{\mathtt{{\color{var}u}}}.  In some cases this
means that some complicated code that would otherwise need to have both a
left version and right version can be written only once. See the methods
\ensuremath{\mathtt{rotateLeft({\color{var}u})}} and \ensuremath{\mathtt{rotateRight({\color{var}u})}} on page~\pageref{page:rotations}
for an example.

\begin{exc}
  Prove that a binary tree having $\ensuremath{\mathtt{{\color{var}n}}}\ge 1$ nodes has $\ensuremath{\mathtt{{\color{var}n}}}-1$ edges.
\end{exc}

\begin{exc}
  Prove that a binary tree having $\ensuremath{\mathtt{{\color{var}n}}}\ge 1$ real nodes has $\ensuremath{\mathtt{{\color{var}n}}}+1$
  external nodes.
\end{exc}

\begin{exc}
  Prove that, if a binary tree, $T$, has at least one leaf, then either
  (a)~$T$'s root has at most one child or (b)~$T$ has more than
  one leaf.
\end{exc}

\begin{exc}
  Implement a non-recursive method, \ensuremath{\mathtt{size2({\color{var}u})}}, that computes the size
  of the subtree rooted at node \ensuremath{\mathtt{{\color{var}u}}}.
\end{exc}

\begin{exc}
  Write a non-recursive method, \ensuremath{\mathtt{height2({\color{var}u})}}, that computes the height
  of node \ensuremath{\mathtt{{\color{var}u}}} in a \texttt{BinaryTree}.
\end{exc}

\begin{exc}
  A binary tree is \emph{size-balanced}
  \index{size-balanced}%
  \index{binary search tree!size-balanced}%
  if, for every node \ensuremath{\mathtt{{\color{var}u}}}, the size of
  the subtrees rooted at \ensuremath{\mathtt{{\color{var}u}.{\color{var}left}}} and \ensuremath{\mathtt{{\color{var}u}.{\color{var}right}}} differ by at most one.
  Write a recursive method, \ensuremath{\mathtt{isBalanced()}}, that tests if a binary tree
  is balanced.  Your method should run in $O(\ensuremath{\mathtt{{\color{var}n}}})$ time.  (Be sure to
  test your code on some large trees with different shapes; it is easy
  to write a method that takes much longer than $O(\ensuremath{\mathtt{{\color{var}n}}})$ time.)
\end{exc}

\index{traversal!pre-order}%
\index{traversal!post-order}%
\index{traversal!in-order}%
\index{pre-order traversal}%
\index{post-order traversal}%
\index{in-order traversal}%
A \emph{pre-order} traversal of a binary tree is a traversal that visits
each node, \ensuremath{\mathtt{{\color{var}u}}}, before any of its children.  An \emph{in-order} traversal
visits \ensuremath{\mathtt{{\color{var}u}}} after visiting all the nodes in \ensuremath{\mathtt{{\color{var}u}}}'s left subtree but before
visiting any of the nodes in \ensuremath{\mathtt{{\color{var}u}}}'s right subtree.  A \emph{post-order}
traversal visits \ensuremath{\mathtt{{\color{var}u}}} only after visiting all other nodes in \ensuremath{\mathtt{{\color{var}u}}}'s subtree.
The pre/in/post-order numbering of a tree labels the nodes of a tree with
the integers $0,\ldots,\ensuremath{\mathtt{{\color{var}n}}}-1$ in the order that they are encountered
by a pre/in/post-order traversal.  See \figref{binarytree-numbering}
for an example.

\begin{figure}
  \begin{center}
    \includegraphics[scale=0.90909]{figs/binarytree-numbering-1}
    \includegraphics[scale=0.90909]{figs/binarytree-numbering-2} \\[2ex]
    \includegraphics[scale=0.90909]{figs/binarytree-numbering-3}
  \end{center}
  \caption{Pre-order, post-order, and in-order numberings of a binary tree.}
  \figlabel{binarytree-numbering}
\end{figure}

\begin{exc}
  \index{number!pre-order}%
  \index{number!post-order}%
  \index{number!in-order}%
  \index{pre-order number}%
  \index{post-order number}%
  \index{in-order number}%
  Create a subclass of \texttt{BinaryTree} whose nodes have fields for storing
  pre-order, post-order, and in-order numbers.  Write recursive methods
  \ensuremath{\mathtt{preOrderNumber()}}, \ensuremath{\mathtt{inOrderNumber()}}, and \ensuremath{\mathtt{postOrderNumbers()}} that
  assign these numbers correctly. These methods should each run in
  $O(\ensuremath{\mathtt{{\color{var}n}}})$ time.
\end{exc}

\begin{exc}\exclabel{tree-traversal}
  Implement the non-recursive functions \ensuremath{\mathtt{nextPreOrder({\color{var}u})}}, \ensuremath{\mathtt{nextInOrder({\color{var}u})}}, and
  \ensuremath{\mathtt{nextPostOrder({\color{var}u})}} that return the node that follows \ensuremath{\mathtt{{\color{var}u}}} in a pre-order,
  in-order, or post-order traversal, respectively.   These functions
  should take amortized constant time; if we start at any node
  \ensuremath{\mathtt{{\color{var}u}}} and repeatedly call one of these functions and assign the return
  value to \ensuremath{\mathtt{{\color{var}u}}} until $\ensuremath{\mathtt{{\color{var}u}}}=\ensuremath{\mathtt{{\color{keyword}null}}}$, then the cost of all these calls should
  be $O(\ensuremath{\mathtt{{\color{var}n}}})$.
\end{exc}

\begin{exc}
  Suppose we are given a binary tree with pre-, post-, and in-order numbers
  assigned to the nodes.  Show how these numbers can be used to answer
  each of the following questions in constant time:
  \begin{enumerate}
    \item Given a node \ensuremath{\mathtt{{\color{var}u}}}, determine the size of the subtree rooted at \ensuremath{\mathtt{{\color{var}u}}}.
    \item Given a node \ensuremath{\mathtt{{\color{var}u}}}, determine the depth of \ensuremath{\mathtt{{\color{var}u}}}.
    \item Given two nodes \ensuremath{\mathtt{{\color{var}u}}} and \ensuremath{\mathtt{{\color{var}w}}}, determine if \ensuremath{\mathtt{{\color{var}u}}} is an ancestor of \ensuremath{\mathtt{{\color{var}w}}}
  \end{enumerate}
\end{exc}

\begin{exc}
  Suppose you are given a list of nodes with pre-order and in-order
  numbers assigned.  Prove that there is at most one possible tree with
  this pre-order/in-order numbering and show how to construct it.
\end{exc}

\begin{exc}
  Show that the shape of any binary tree on \ensuremath{\mathtt{{\color{var}n}}} nodes can be represented
  using at most $2(\ensuremath{\mathtt{{\color{var}n}}}-1)$ bits.  (Hint: think about recording what
  happens during a traversal and then playing back that recording to
  reconstruct the tree.)
\end{exc}

\begin{exc}
  Illustrate what happens when we add the values $3.5$ and then 4.5 to
  the binary search tree in \figref{bst}.
\end{exc}

\begin{exc}
  Illustrate what happens when we remove the values $3$ and then 5 from
  the binary search tree in \figref{bst}.
\end{exc}

\begin{exc}
  Implement a \texttt{BinarySearchTree} method, \ensuremath{\mathtt{getLE({\color{var}x})}},
  that returns a list of all items in the tree that are less than or
  equal to \ensuremath{\mathtt{{\color{var}x}}}.  The running time of your method should be $O(\ensuremath{\mathtt{{\color{var}n}}}'+\ensuremath{\mathtt{{\color{var}h}}})$
  where $\ensuremath{\mathtt{{\color{var}n}}}'$ is the number of items less than or equal to \ensuremath{\mathtt{{\color{var}x}}} and \ensuremath{\mathtt{{\color{var}h}}}
  is the height of the tree.
\end{exc}

\begin{exc}
  Describe how to add the elements $\{1,\ldots,\ensuremath{\mathtt{{\color{var}n}}}\}$ to an initially
  empty \texttt{BinarySearchTree} in such a way that the resulting tree has
  height $\ensuremath{\mathtt{{\color{var}n}}}-1$.  How many ways are there to do this?
\end{exc}

\begin{exc}
  If we have some \texttt{BinarySearchTree} and perform the operations \ensuremath{\mathtt{add({\color{var}x})}}
  followed by \ensuremath{\mathtt{remove({\color{var}x})}} (with the same value of \ensuremath{\mathtt{{\color{var}x}}}) do we necessarily
  return to the original tree?
\end{exc}

\begin{exc}
  Can a \ensuremath{\mathtt{remove({\color{var}x})}} operation increase the height of any node in a
  \texttt{BinarySearchTree}?  If so, by how much?
\end{exc}

\begin{exc}
  Can an \ensuremath{\mathtt{add({\color{var}x})}} operation increase the height of any node in a
  \texttt{BinarySearchTree}?  Can it increase the height of the tree?  If so,
  by how much?
\end{exc}

\begin{exc}
  Design and implement a version of \texttt{BinarySearchTree} in which each node,
  \ensuremath{\mathtt{{\color{var}u}}}, maintains values \ensuremath{\mathtt{{\color{var}u}.{\color{var}size}}} (the size of the subtree rooted at \ensuremath{\mathtt{{\color{var}u}}}),
  \ensuremath{\mathtt{{\color{var}u}.{\color{var}depth}}} (the depth of \ensuremath{\mathtt{{\color{var}u}}}), and \ensuremath{\mathtt{{\color{var}u}.{\color{var}height}}} (the height of the subtree
  rooted at \ensuremath{\mathtt{{\color{var}u}}}).  

  These values should be maintained, even during calls to the \ensuremath{\mathtt{add({\color{var}x})}}
  and \ensuremath{\mathtt{remove({\color{var}x})}} operations, but this should not increase the cost of
  these operations by more than a constant factor.
\end{exc}

\clearpage

\clearpage
\chapter{Random Binary Search Trees}
\chaplabel{rbs}

In this chapter, we present a binary search tree structure that uses
randomization to achieve $O(\log \ensuremath{\mathtt{{\color{var}n}}})$ expected time for all operations.

\section{Random Binary Search Trees}
\seclabel{rbst}

Consider the two binary search trees shown in \figref{rbs-lvc}, each of
which has $\ensuremath{\mathtt{{\color{var}n}}}=15$ nodes.  The one on the left is a list and the other
is a perfectly balanced binary search tree. The one on the left has a
height of $\ensuremath{\mathtt{{\color{var}n}}}-1=14$ and the one on the right has a height of three.

\begin{figure}
  \begin{center}
    \begin{tabular}{cc}
      \includegraphics[scale=0.90909,scale=0.95]{figs/bst-path} &
      \includegraphics[scale=0.90909,scale=0.95]{figs/bst-balanced}
    \end{tabular}
  \end{center}
  \caption{Two binary search trees containing the integers $0,\ldots,14$.}
  \figlabel{rbs-lvc}
\end{figure}

Imagine how these two trees could have been constructed.  The one on
the left occurs if we start with an empty \texttt{BinarySearchTree} and add
the sequence
\[
    \langle 0,1,2,3,4,5,6,7,8,9,10,11,12,13,14 \rangle \enspace .
\]
No other sequence of additions will create this tree (as you can prove
by induction on \ensuremath{\mathtt{{\color{var}n}}}). On the other hand, the tree on the right can be
created by the sequence
\[
    \langle 7,3,11,1,5,9,13,0,2,4,6,8,10,12,14 \rangle  \enspace .
\]
Other sequences work as well, including
\[
    \langle 7,3,1,5,0,2,4,6,11,9,13,8,10,12,14 \rangle  \enspace ,
\]
and
\[
    \langle 7,3,1,11,5,0,2,4,6,9,13,8,10,12,14 \rangle \enspace .
\]
In fact, there are $21,964,800$ addition sequences that generate the
tree on the right and only one that generates the tree on the left.

The above example gives some anecdotal evidence that, if we choose a
random permutation of $0,\ldots,14$, and add it into a binary search
tree, then we are more likely to get a very balanced tree (the right
side of \figref{rbs-lvc}) than we are to get a very unbalanced tree
(the left side of \figref{rbs-lvc}).

We can formalize this notion by studying random binary search trees.
A \emph{random binary search tree}
\index{random binary search tree}%
\index{binary search tree!random}%
of size \ensuremath{\mathtt{{\color{var}n}}} is obtained in the
following way:  Take a random permutation, $\ensuremath{\mathtt{{\color{var}x}}}_0,\ldots,\ensuremath{\mathtt{{\color{var}x}}}_{\ensuremath{\mathtt{{\color{var}n}}}-1}$,
of the integers $0,\ldots,\ensuremath{\mathtt{{\color{var}n}}}-1$ and add its elements, one by one,
into a \texttt{BinarySearchTree}.  By \emph{random permutation}
\index{permutation!random}%
\index{random permutation}%
we mean that
each of the possible $\ensuremath{\mathtt{{\color{var}n}}}!$ permutations (orderings) of $0,\ldots,\ensuremath{\mathtt{{\color{var}n}}}-1$
is equally likely, so that the probability of obtaining any particular
permutation is $1/\ensuremath{\mathtt{{\color{var}n}}}!$.

Note that the values $0,\ldots,\ensuremath{\mathtt{{\color{var}n}}}-1$ could be replaced by any ordered
set of \ensuremath{\mathtt{{\color{var}n}}} elements without changing any of the properties of the
random binary search tree.  The element $\ensuremath{\mathtt{{\color{var}x}}}\in\{0,\ldots,\ensuremath{\mathtt{{\color{var}n}}}-1\}$ is
simply standing in for the element of rank \ensuremath{\mathtt{{\color{var}x}}} in an ordered set of
size \ensuremath{\mathtt{{\color{var}n}}}.

Before we can present our main result about random binary search trees,
we must take some time for a short digression to discuss a type of number
that comes up frequently when studying randomized structures. For a
non-negative integer, $k$, the $k$-th \emph{harmonic number},
\index{harmonic number}%
\index{H@$H_k$ (harmonic number)}%
denoted
$H_k$, is defined as
\[
  H_k = 1 + 1/2 + 1/3 + \cdots + 1/k \enspace .
\] 
The harmonic number $H_k$ has no simple closed form, but it is very
closely related to the natural logarithm of $k$.  In particular,
\[
  \ln k < H_k \le \ln k + 1  \enspace .
\]
\newcommand{\hint}{\int_1^k\! (1/x)\, \mathrm{d}x}%
Readers who have studied calculus might notice that this is because
the integral $\hint = \ln k$.  Keeping in mind that an integral can be
interpreted as the area between a curve and the $x$-axis, the value of
$H_k$ can be lower-bounded by the integral $\hint$ and upper-bounded by
$1+ \hint$.  (See \figref{harmonic-integral} for a graphical explanation.)

\begin{figure}
  \begin{center}
    \begin{tabular}{cc}
      \includegraphics[width=\HalfScaleIfNeeded]{figs/harmonic-2} 
        & \includegraphics[width=\HalfScaleIfNeeded]{figs/harmonic-3}
    \end{tabular}
  \end{center}
  \caption{The $k$th harmonic number $H_k=\sum_{i=1}^k 1/i$ is upper- and lower-bounded by two integrals. The value of these integrals is given by the 
  area of the shaded region, while the value of $H_k$ is given by the area of
  the rectangles.}
  \figlabel{harmonic-integral}
\end{figure}


\begin{lem}\lemlabel{rbs}
  In a random binary search tree of size \ensuremath{\mathtt{{\color{var}n}}}, the following statements hold:
  \begin{enumerate}
    \item For any $\ensuremath{\mathtt{{\color{var}x}}}\in\{0,\ldots,\ensuremath{\mathtt{{\color{var}n}}}-1\}$, the expected length of the
    search path for \ensuremath{\mathtt{{\color{var}x}}} is $H_{\ensuremath{\mathtt{{\color{var}x}}}+1} + H_{\ensuremath{\mathtt{{\color{var}n}}}-\ensuremath{\mathtt{{\color{var}x}}}} - O(1)$.\footnote{The
    expressions $\ensuremath{\mathtt{{\color{var}x}}}+1$ and $\ensuremath{\mathtt{{\color{var}n}}}-\ensuremath{\mathtt{{\color{var}x}}}$ can be interpreted respectively
    as the number of elements in the tree less than or equal to \ensuremath{\mathtt{{\color{var}x}}}
    and the number of elements in the tree greater than or equal to \ensuremath{\mathtt{{\color{var}x}}}.}
    \item For any $\ensuremath{\mathtt{{\color{var}x}}}\in(-1,n)\setminus\{0,\ldots,\ensuremath{\mathtt{{\color{var}n}}}-1\}$, the
    expected length of the search path for \ensuremath{\mathtt{{\color{var}x}}} is $H_{\lceil\ensuremath{\mathtt{{\color{var}x}}}\rceil}
    + H_{\ensuremath{\mathtt{{\color{var}n}}}-\lceil\ensuremath{\mathtt{{\color{var}x}}}\rceil}$.
  \end{enumerate}
\end{lem}

We will prove \lemref{rbs} in the next section.  For now, consider what
the two parts of \lemref{rbs} tell us.  The first part tells us that if
we search for an element in a tree of size \ensuremath{\mathtt{{\color{var}n}}}, then the expected length
of the search path is at most $2\ln n + O(1)$.  The second part tells
us the same thing about searching for a value not stored in the tree.
When we compare the two parts of the lemma, we see that it is only
slightly faster to search for something that is in the tree compared to
something that is not.


\subsection{Proof of \lemref{rbs}}

The key observation needed to prove \lemref{rbs} is the following:
The search path for a value \ensuremath{\mathtt{{\color{var}x}}} in the open interval $(-1,\ensuremath{\mathtt{{\color{var}n}}})$ in a
random binary search tree, $T$, contains the node with key $i < \ensuremath{\mathtt{{\color{var}x}}}$
if, and only if, in the random permutation used to create $T$, $i$
appears before any of $\{i+1,i+2,\ldots,\lfloor\ensuremath{\mathtt{{\color{var}x}}}\rfloor\}$.

To see this, refer to \figref{rbst-records} and notice that until
some value in $\{i,i+1,\ldots,\lfloor\ensuremath{\mathtt{{\color{var}x}}}\rfloor\}$ is added, the search
paths for each value in the open interval $(i-1,\lfloor\ensuremath{\mathtt{{\color{var}x}}}\rfloor+1)$
are identical.  (Remember that for two values to have
different search paths, there must be some element in the tree that
compares differently with them.)  Let $j$ be the first element in
$\{i,i+1,\ldots,\lfloor\ensuremath{\mathtt{{\color{var}x}}}\rfloor\}$ to appear in the random permutation.
Notice that $j$ is now and will always be on the search path for \ensuremath{\mathtt{{\color{var}x}}}.
If $j\neq i$ then the node $\ensuremath{\mathtt{{\color{var}u}}}_j$ containing $j$ is created before the
node $\ensuremath{\mathtt{{\color{var}u}}}_i$ that contains $i$.  Later, when $i$ is added, it will be
added to the subtree rooted at $\ensuremath{\mathtt{{\color{var}u}}}_j\ensuremath{\mathtt{.{\color{var}left}}}$, since $i<j$.  On the other
hand, the search path for \ensuremath{\mathtt{{\color{var}x}}} will never visit this subtree because it
will proceed to $\ensuremath{\mathtt{{\color{var}u}}}_j\ensuremath{\mathtt{.{\color{var}right}}}$ after visiting $\ensuremath{\mathtt{{\color{var}u}}}_j$.

\begin{figure}
  \begin{center}
    \includegraphics[width=\ScaleIfNeeded]{figs/rbst-records}
  \end{center}
  \caption[The search path in a random binary search tree]{The value $i<\ensuremath{\mathtt{{\color{var}x}}}$ is on the search path for \ensuremath{\mathtt{{\color{var}x}}} if and only
   if $i$ is the first element among $\{i,i+1,\ldots,\lfloor\ensuremath{\mathtt{{\color{var}x}}}\rfloor\}$ added to the tree.}
  \figlabel{rbst-records}
\end{figure}

Similarly, for $i>\ensuremath{\mathtt{{\color{var}x}}}$, $i$ appears in the search path for \ensuremath{\mathtt{{\color{var}x}}}
if and only if $i$ appears before any of $\{\lceil\ensuremath{\mathtt{{\color{var}x}}}\rceil,
\lceil\ensuremath{\mathtt{{\color{var}x}}}\rceil+1,\ldots,i-1\}$ in the random permutation used to
create $T$.

Notice that, if we start with a random permutation of $\{0,\ldots,\ensuremath{\mathtt{{\color{var}n}}}\}$,
then the subsequences containing only $\{i,i+1,\ldots,\lfloor\ensuremath{\mathtt{{\color{var}x}}}\rfloor\}$
and $\{\lceil\ensuremath{\mathtt{{\color{var}x}}}\rceil, \lceil\ensuremath{\mathtt{{\color{var}x}}}\rceil+1,\ldots,i-1\}$ are also random
permutations of their respective elements.  Each element, then, in the
subsets $\{i,i+1,\ldots,\lfloor\ensuremath{\mathtt{{\color{var}x}}}\rfloor\}$ and $\{\lceil\ensuremath{\mathtt{{\color{var}x}}}\rceil,
\lceil\ensuremath{\mathtt{{\color{var}x}}}\rceil+1,\ldots,i-1\}$ is equally likely to appear before
any other in its subset in the random permutation used to create $T$.
So we have
\[
  \Pr\{\mbox{$i$ is on the search path for \ensuremath{\mathtt{{\color{var}x}}}}\}
  = \left\{ \begin{array}{ll}
     1/(\lfloor\ensuremath{\mathtt{{\color{var}x}}}\rfloor-i+1) & \mbox{if $i < \ensuremath{\mathtt{{\color{var}x}}}$} \\
     1/(i-\lceil\ensuremath{\mathtt{{\color{var}x}}}\rceil+1) & \mbox{if $i > \ensuremath{\mathtt{{\color{var}x}}}$} 
     \end{array}\right . \enspace .
\]

With this observation, the proof of \lemref{rbs}
involves some simple calculations with harmonic numbers:

\begin{proof}[Proof of \lemref{rbs}]
Let $I_i$ be the indicator random variable that is equal to one when $i$
appears on the search path for \ensuremath{\mathtt{{\color{var}x}}} and zero otherwise.  Then the length
of the search path is given by
\[
  \sum_{i\in\{0,\ldots,\ensuremath{\mathtt{{\color{var}n}}}-1\}\setminus\{\ensuremath{\mathtt{{\color{var}x}}}\}} I_i
\]
so, if $\ensuremath{\mathtt{{\color{var}x}}}\in\{0,\ldots,\ensuremath{\mathtt{{\color{var}n}}}-1\}$, the expected length of the search
path is given by (see \figref{rbst-probs}.a)
\begin{align*}
  \E\left[\sum_{i=0}^{\ensuremath{\mathtt{{\color{var}x}}}-1} I_i + \sum_{i=\ensuremath{\mathtt{{\color{var}x}}}+1}^{\ensuremath{\mathtt{{\color{var}n}}}-1} I_i\right]
   & =  \sum_{i=0}^{\ensuremath{\mathtt{{\color{var}x}}}-1} \E\left[I_i\right]
         + \sum_{i=\ensuremath{\mathtt{{\color{var}x}}}+1}^{\ensuremath{\mathtt{{\color{var}n}}}-1} \E\left[I_i\right] \\
   & = \sum_{i=0}^{\ensuremath{\mathtt{{\color{var}x}}}-1} 1/(\lfloor\ensuremath{\mathtt{{\color{var}x}}}\rfloor-i+1)
         + \sum_{i=\ensuremath{\mathtt{{\color{var}x}}}+1}^{\ensuremath{\mathtt{{\color{var}n}}}-1} 1/(i-\lceil\ensuremath{\mathtt{{\color{var}x}}}\rceil+1) \\
   & = \sum_{i=0}^{\ensuremath{\mathtt{{\color{var}x}}}-1} 1/(\ensuremath{\mathtt{{\color{var}x}}}-i+1)
         + \sum_{i=\ensuremath{\mathtt{{\color{var}x}}}+1}^{\ensuremath{\mathtt{{\color{var}n}}}-1} 1/(i-\ensuremath{\mathtt{{\color{var}x}}}+1) \\
   & = \frac{1}{2}+\frac{1}{3}+\cdots+\frac{1}{\ensuremath{\mathtt{{\color{var}x}}}+1} \\
   & \quad {} + \frac{1}{2}+\frac{1}{3}+\cdots+\frac{1}{\ensuremath{\mathtt{{\color{var}n}}}-\ensuremath{\mathtt{{\color{var}x}}}} \\
   & = H_{\ensuremath{\mathtt{{\color{var}x}}}+1} + H_{\ensuremath{\mathtt{{\color{var}n}}}-\ensuremath{\mathtt{{\color{var}x}}}} - 2  \enspace .
\end{align*}
The corresponding calculations for a search value
$\ensuremath{\mathtt{{\color{var}x}}}\in(-1,n)\setminus\{0,\ldots,\ensuremath{\mathtt{{\color{var}n}}}-1\}$ are almost identical (see
\figref{rbst-probs}.b).
\end{proof}

\begin{figure}
  \begin{center}
    \begin{tabular}{@{}c@{}}
      \includegraphics[width=\ScaleIfNeeded]{figs/rbst-probs-a} \\ (a) \\[2ex]
      \includegraphics[width=\ScaleIfNeeded]{figs/rbst-probs-b} \\ (b) \\[2ex]
    \end{tabular}
  \end{center}
  \caption[The probabilities of an element being on a search path]{The probabilities of an element being on the search path for \ensuremath{\mathtt{{\color{var}x}}}
   when (a)~\ensuremath{\mathtt{{\color{var}x}}} is an integer and (b)~when \ensuremath{\mathtt{{\color{var}x}}} is not an integer.}
  \figlabel{rbst-probs}
\end{figure}

\subsection{Summary}

The following theorem summarizes the performance of a random binary
search tree:

\begin{thm}\thmlabel{rbs}
A random binary search tree can be constructed in $O(\ensuremath{\mathtt{{\color{var}n}}}\log \ensuremath{\mathtt{{\color{var}n}}})$ time.
In a random binary search tree, the \ensuremath{\mathtt{find({\color{var}x})}} operation takes $O(\log
\ensuremath{\mathtt{{\color{var}n}}})$ expected time.
\end{thm}

We should emphasize again that the expectation in \thmref{rbs} is with
respect to the random permutation used to create the random binary
search tree.  In particular, it does not depend on a random choice of
\ensuremath{\mathtt{{\color{var}x}}}; it is true for every value of \ensuremath{\mathtt{{\color{var}x}}}.


\section{\texttt{Treap}: A Randomized Binary Search Tree}
\seclabel{treap}

\index{Treap@\texttt{Treap}}%
The problem with random binary search trees is, of course, that they are
not dynamic.  They don't support the \ensuremath{\mathtt{add({\color{var}x})}} or \ensuremath{\mathtt{remove({\color{var}x})}} operations
needed to implement the \texttt{SSet} interface.  In this section we describe
a data structure called a \texttt{Treap} that uses \lemref{rbs} to implement
the \texttt{SSet} interface.\footnote{The name \texttt{Treap} comes from the fact
that this data structure is simultaneously a binary search \textbf{tr}ee
(\secref{binarysearchtree}) and a h\textbf{eap} (\chapref{heaps}).}

A node in a \texttt{Treap} is like a node in a \texttt{BinarySearchTree} in that it has
a data value, \ensuremath{\mathtt{{\color{var}x}}}, but it also contains a unique numerical \emph{priority},
\ensuremath{\mathtt{{\color{var}p}}}, that is assigned at random:


 \ \vspace{.5ex}\begin{Verbatim}[gobble=1,tabsize=2,frame=single,framerule=1.2pt,commandchars=\\@\$,label=\texttt{Treap},labelposition=topline]
	@\color@keyword$class$ Node<@\color@keyword$T$> @\color@keyword$extends$ BSTNode<Node<@\color@keyword$T$>,@\color@keyword$T$> {
		@\color@keyword$int$ @\color@var$p$;
	}
\end{Verbatim}

In addition to being a binary search tree, the nodes in a \texttt{Treap}
also obey the \emph{heap property}:
\begin{itemize}
\item (Heap Property)  At every node \ensuremath{\mathtt{{\color{var}u}}}, except the root, 
      $\ensuremath{\mathtt{{\color{var}u}.{\color{var}parent}.{\color{var}p}}} < \ensuremath{\mathtt{{\color{var}u}.{\color{var}p}}}$.
      \index{heap property}%
\end{itemize}
In other words, each node has a priority smaller than that of its two
children.  An example is shown in \figref{treap}.

\begin{figure}
  \begin{center}
    \includegraphics[width=\ScaleIfNeeded]{figs/treap}
  \end{center}
  \caption[A Treap]{An example of a \texttt{Treap} containing the integers
  $0,\ldots,9$. Each node, \ensuremath{\mathtt{{\color{var}u}}}, is illustrated as a box containing $\ensuremath{\mathtt{{\color{var}u}.{\color{var}x}}},\ensuremath{\mathtt{{\color{var}u}.{\color{var}p}}}$.}
  \figlabel{treap}
\end{figure}

The heap and binary search tree conditions together ensure that, once
the key (\ensuremath{\mathtt{{\color{var}x}}}) and priority (\ensuremath{\mathtt{{\color{var}p}}}) for each node are defined, the
shape of the \texttt{Treap} is completely determined. The heap property tells us that
the node with minimum priority has to be the root, \ensuremath{\mathtt{{\color{var}r}}}, of the \texttt{Treap}.
The binary search tree property tells us that all nodes with keys smaller
than \ensuremath{\mathtt{{\color{var}r}.{\color{var}x}}} are stored in the subtree rooted at \ensuremath{\mathtt{{\color{var}r}.{\color{var}left}}} and all nodes
with keys larger than \ensuremath{\mathtt{{\color{var}r}.{\color{var}x}}} are stored in the subtree rooted at \ensuremath{\mathtt{{\color{var}r}.{\color{var}right}}}.

The important point about the priority values in a \texttt{Treap} is that they
are unique and assigned at random.  Because of this, there are
two equivalent ways we can think about a \texttt{Treap}.  As defined above, a
\texttt{Treap} obeys the heap and binary search tree properties.  Alternatively,
we can think of a \texttt{Treap} as a \texttt{BinarySearchTree} whose nodes
were added in increasing order of priority.  For example, the \texttt{Treap}
in \figref{treap} can be obtained by adding the sequence of $(\ensuremath{\mathtt{{\color{var}x}}},\ensuremath{\mathtt{{\color{var}p}}})$
values 
\[
  \langle
   (3,1), (1,6), (0,9), (5,11), (4,14), (9,17), (7,22), (6,42), (8,49), (2,99)
  \rangle
\]
into a \texttt{BinarySearchTree}.

Since the priorities are chosen randomly, this is equivalent to taking a
random permutation of the keys---in this case the permutation is
\[
  \langle 3, 1, 0, 5, 9, 4, 7, 6, 8, 2 \rangle
\]
---and adding these to a \texttt{BinarySearchTree}.  But this means that the
shape of a treap is identical to that of a random binary search tree.
In particular, if we replace each key \ensuremath{\mathtt{{\color{var}x}}} by its rank,\footnote{The
rank of an element \ensuremath{\mathtt{{\color{var}x}}} in a set $S$ of elements is the number of
elements in $S$ that are less than \ensuremath{\mathtt{{\color{var}x}}}.} then \lemref{rbs} applies.
Restating \lemref{rbs} in terms of \texttt{Treap}s, we have:
\begin{lem}\lemlabel{rbs-treap}
  In a \texttt{Treap} that stores a set $S$ of \ensuremath{\mathtt{{\color{var}n}}} keys, the following statements hold:
  \begin{enumerate}
    \item For any $\ensuremath{\mathtt{{\color{var}x}}}\in S$, the expected length of
    the search path for \ensuremath{\mathtt{{\color{var}x}}} is $H_{r(\ensuremath{\mathtt{{\color{var}x}}})+1} + H_{\ensuremath{\mathtt{{\color{var}n}}}-r(\ensuremath{\mathtt{{\color{var}x}}})} - O(1)$.
    \item For any $\ensuremath{\mathtt{{\color{var}x}}}\not\in S$, the expected length of the
    search path for \ensuremath{\mathtt{{\color{var}x}}} is $H_{r(\ensuremath{\mathtt{{\color{var}x}}})} + H_{\ensuremath{\mathtt{{\color{var}n}}}-r(\ensuremath{\mathtt{{\color{var}x}}})}$.
  \end{enumerate}
  Here, $r(\ensuremath{\mathtt{{\color{var}x}}})$ denotes the rank of \ensuremath{\mathtt{{\color{var}x}}} in the set $S\cup\{\ensuremath{\mathtt{{\color{var}x}}}\}$.
\end{lem}
Again, we emphasize that the expectation in \lemref{rbs-treap} is taken
over the random choices of the priorities for each node.  It does not
require any assumptions about the randomness in the keys.

\lemref{rbs-treap} tells us that \texttt{Treap}s can implement the \ensuremath{\mathtt{find({\color{var}x})}}
operation efficiently. However, the real benefit of a \texttt{Treap} is that
it can support the \ensuremath{\mathtt{add({\color{var}x})}} and \ensuremath{\mathtt{delete({\color{var}x})}} operations.  To
do this, it needs to perform rotations in order to maintain the heap property.  Refer to \figref{rotations}.
A \emph{rotation}
\index{rotation}%
in a binary
search tree is a local modification that takes a parent \ensuremath{\mathtt{{\color{var}u}}} of a node \ensuremath{\mathtt{{\color{var}w}}}
and makes \ensuremath{\mathtt{{\color{var}w}}} the parent of \ensuremath{\mathtt{{\color{var}u}}}, while preserving the binary search tree
property. Rotations come in two flavours: \emph{left} or \emph{right}
depending on whether \ensuremath{\mathtt{{\color{var}w}}} is a right or left child of \ensuremath{\mathtt{{\color{var}u}}}, respectively.
\index{left rotation}%
\index{right rotation}%

\begin{figure}
  \begin{center}
     \includegraphics[width=\ScaleIfNeeded]{figs/rotation}
  \end{center}
  \caption{Left and right rotations in a binary search tree.}
  \figlabel{rotations}
\end{figure}

The code that implements this has to handle these two possibilities and
be careful of a boundary case (when \ensuremath{\mathtt{{\color{var}u}}} is the root), so the actual code
is a little longer than \figref{rotations} would lead a reader to believe:


 \ \vspace{.5ex}\begin{Verbatim}[gobble=1,tabsize=2,frame=single,framerule=1.2pt,commandchars=\\@\$,label=\texttt{BinarySearchTree},labelposition=topline]
	@\color@keyword$void$ rotateLeft(Node @\color@var$u$) {
		Node @\color@var$w$ = @\color@var$u$.@\color@var$right$;
		@\color@var$w$.@\color@var$parent$ = @\color@var$u$.@\color@var$parent$;
		@\color@keyword$if$ (@\color@var$w$.@\color@var$parent$ != @\color@var$nil$) {
			@\color@keyword$if$ (@\color@var$w$.@\color@var$parent$.@\color@var$left$ == @\color@var$u$) {
				@\color@var$w$.@\color@var$parent$.@\color@var$left$ = @\color@var$w$;
			} @\color@keyword$else$ {
				@\color@var$w$.@\color@var$parent$.@\color@var$right$ = @\color@var$w$;
			}
		}
		@\color@var$u$.@\color@var$right$ = @\color@var$w$.@\color@var$left$;
		@\color@keyword$if$ (@\color@var$u$.@\color@var$right$ != @\color@var$nil$) {
			@\color@var$u$.@\color@var$right$.@\color@var$parent$ = @\color@var$u$;
		}
		@\color@var$u$.@\color@var$parent$ = @\color@var$w$;
		@\color@var$w$.@\color@var$left$ = @\color@var$u$;
		@\color@keyword$if$ (@\color@var$u$ == @\color@var$r$) { @\color@var$r$ = @\color@var$w$; @\color@var$r$.@\color@var$parent$ = @\color@var$nil$; }
	}	
	@\color@keyword$void$ rotateRight(Node @\color@var$u$) {
		Node @\color@var$w$ = @\color@var$u$.@\color@var$left$;
		@\color@var$w$.@\color@var$parent$ = @\color@var$u$.@\color@var$parent$;
		@\color@keyword$if$ (@\color@var$w$.@\color@var$parent$ != @\color@var$nil$) {
			@\color@keyword$if$ (@\color@var$w$.@\color@var$parent$.@\color@var$left$ == @\color@var$u$) {
				@\color@var$w$.@\color@var$parent$.@\color@var$left$ = @\color@var$w$;
			} @\color@keyword$else$ {
				@\color@var$w$.@\color@var$parent$.@\color@var$right$ = @\color@var$w$;
			}
		}
		@\color@var$u$.@\color@var$left$ = @\color@var$w$.@\color@var$right$;
		@\color@keyword$if$ (@\color@var$u$.@\color@var$left$ != @\color@var$nil$) {
			@\color@var$u$.@\color@var$left$.@\color@var$parent$ = @\color@var$u$;
		}
		@\color@var$u$.@\color@var$parent$ = @\color@var$w$;
		@\color@var$w$.@\color@var$right$ = @\color@var$u$;
		@\color@keyword$if$ (@\color@var$u$ == @\color@var$r$) { @\color@var$r$ = @\color@var$w$; @\color@var$r$.@\color@var$parent$ = @\color@var$nil$; }
	}
\end{Verbatim}

\label{page:rotations}
In terms of the \texttt{Treap} data structure, the most important property of a
rotation is that the depth of \ensuremath{\mathtt{{\color{var}w}}} decreases by one while the depth of \ensuremath{\mathtt{{\color{var}u}}}
increases by one.

Using rotations, we can implement the \ensuremath{\mathtt{add({\color{var}x})}} operation as follows:
We create a new node, \ensuremath{\mathtt{{\color{var}u}}}, assign \ensuremath{\mathtt{{\color{var}u}.{\color{var}x}={\color{var}x}}}, and pick a random value
for \ensuremath{\mathtt{{\color{var}u}.{\color{var}p}}}.  Next we add \ensuremath{\mathtt{{\color{var}u}}} using the usual \ensuremath{\mathtt{add({\color{var}x})}} algorithm
for a \texttt{BinarySearchTree}, so that \ensuremath{\mathtt{{\color{var}u}}} is now a leaf of the \texttt{Treap}.
At this point, our \texttt{Treap} satisfies the binary search tree property,
but not necessarily the heap property.  In particular, it may be the
case that \ensuremath{\mathtt{{\color{var}u}.{\color{var}parent}.{\color{var}p} > {\color{var}u}.{\color{var}p}}}.  If this is the case, then we perform a
rotation at node \ensuremath{\mathtt{{\color{var}w}}}=\ensuremath{\mathtt{{\color{var}u}.{\color{var}parent}}} so that \ensuremath{\mathtt{{\color{var}u}}} becomes the parent of \ensuremath{\mathtt{{\color{var}w}}}.
If \ensuremath{\mathtt{{\color{var}u}}} continues to violate the heap property, we will have to repeat this, decreasing \ensuremath{\mathtt{{\color{var}u}}}'s depth by one every time, until
\ensuremath{\mathtt{{\color{var}u}}} either becomes the root or $\ensuremath{\mathtt{{\color{var}u}.{\color{var}parent}.{\color{var}p}}} < \ensuremath{\mathtt{{\color{var}u}.{\color{var}p}}}$.


 \ \vspace{.5ex}\begin{Verbatim}[gobble=1,tabsize=2,frame=single,framerule=1.2pt,commandchars=\\@\$,label=\texttt{Treap},labelposition=topline]
	@\color@var$boolean$ add(@\color@keyword$T$ @\color@var$x$) {
		Node<@\color@keyword$T$> @\color@var$u$ = newNode();
		@\color@var$u$.@\color@var$x$ = @\color@var$x$;
		@\color@var$u$.@\color@var$p$ = @\color@var$rand$.nextInt();
		@\color@keyword$if$ (@\color@var$super$.add(@\color@var$u$)) {
			bubbleUp(@\color@var$u$);
			@\color@var$return$ @\color@var$true$;
		}
		@\color@var$return$ @\color@var$false$;
	}
	@\color@keyword$void$ bubbleUp(Node<@\color@keyword$T$> @\color@var$u$) {
		@\color@keyword$while$ (@\color@var$u$.@\color@var$parent$ != @\color@var$nil$ && @\color@var$u$.@\color@var$parent$.@\color@var$p$ > @\color@var$u$.@\color@var$p$) {
			@\color@keyword$if$ (@\color@var$u$.@\color@var$parent$.@\color@var$right$ == @\color@var$u$) {
				rotateLeft(@\color@var$u$.@\color@var$parent$);
			} @\color@keyword$else$ {
				rotateRight(@\color@var$u$.@\color@var$parent$);
			}
		}
		@\color@keyword$if$ (@\color@var$u$.@\color@var$parent$ == @\color@var$nil$) {
			@\color@var$r$ = @\color@var$u$;
		}
	}
\end{Verbatim}

An example of an \ensuremath{\mathtt{add({\color{var}x})}} operation is shown in \figref{treap-add}.

\begin{figure}
  \begin{center}
  \includegraphics[width=\ScaleIfNeeded]{figs/treap-insert-a} \\
  \includegraphics[width=\ScaleIfNeeded]{figs/treap-insert-b} \\
  \includegraphics[width=\ScaleIfNeeded]{figs/treap-insert-c} \\
  \end{center}
  \caption[Adding to a Treap]{Adding the value 1.5 into the \texttt{Treap} from \figref{treap}.}
  \figlabel{treap-add}
\end{figure}

The running time of the \ensuremath{\mathtt{add({\color{var}x})}} operation is given by the time it
takes to follow the search path for \ensuremath{\mathtt{{\color{var}x}}} plus the number of rotations
performed to move the newly-added node, \ensuremath{\mathtt{{\color{var}u}}}, up to its correct location
in the \texttt{Treap}.  By \lemref{rbs-treap}, the expected length of the
search path is at most $2\ln \ensuremath{\mathtt{{\color{var}n}}}+O(1)$.  Furthermore, each rotation
decreases the depth of \ensuremath{\mathtt{{\color{var}u}}}.   This stops if \ensuremath{\mathtt{{\color{var}u}}} becomes the root, so
the expected number of rotations cannot exceed the expected length of
the search path.  Therefore, the expected running time of the \ensuremath{\mathtt{add({\color{var}x})}}
operation in a \texttt{Treap} is $O(\log \ensuremath{\mathtt{{\color{var}n}}})$.  (\excref{treap-rotates}
asks you to show that the expected number of rotations performed during
an addition is actually only $O(1)$.)

The \ensuremath{\mathtt{remove({\color{var}x})}} operation in a \texttt{Treap} is the opposite of the \ensuremath{\mathtt{add({\color{var}x})}}
operation.  We search for the node, \ensuremath{\mathtt{{\color{var}u}}}, containing \ensuremath{\mathtt{{\color{var}x}}}, then perform
rotations to move \ensuremath{\mathtt{{\color{var}u}}} downwards until it becomes a leaf, and then we
splice \ensuremath{\mathtt{{\color{var}u}}} from the \texttt{Treap}.  Notice that, to move \ensuremath{\mathtt{{\color{var}u}}} downwards, we can
perform either a left or right rotation at \ensuremath{\mathtt{{\color{var}u}}}, which will replace \ensuremath{\mathtt{{\color{var}u}}}
with \ensuremath{\mathtt{{\color{var}u}.{\color{var}right}}} or \ensuremath{\mathtt{{\color{var}u}.{\color{var}left}}}, respectively.
The choice is made by the first of the following that apply:
\begin{enumerate}
\item If \ensuremath{\mathtt{{\color{var}u}.{\color{var}left}}} and \ensuremath{\mathtt{{\color{var}u}.{\color{var}right}}} are both \ensuremath{\mathtt{{\color{keyword}null}}}, then \ensuremath{\mathtt{{\color{var}u}}} is a leaf and no rotation is performed.
\item If \ensuremath{\mathtt{{\color{var}u}.{\color{var}left}}} (or \ensuremath{\mathtt{{\color{var}u}.{\color{var}right}}}) is \ensuremath{\mathtt{{\color{keyword}null}}}, then perform a right (or left, respectively) rotation at \ensuremath{\mathtt{{\color{var}u}}}.
\item If $\ensuremath{\mathtt{{\color{var}u}.{\color{var}left}.{\color{var}p}}} < \ensuremath{\mathtt{{\color{var}u}.{\color{var}right}.{\color{var}p}}}$ (or $\ensuremath{\mathtt{{\color{var}u}.{\color{var}left}.{\color{var}p}}} > \ensuremath{\mathtt{{\color{var}u}.{\color{var}right}.{\color{var}p}}})$, then perform a right rotation (or left rotation, respectively) at \ensuremath{\mathtt{{\color{var}u}}}.
\end{enumerate}
These three rules ensure that the \texttt{Treap} doesn't become disconnected and that the heap property is restored once \ensuremath{\mathtt{{\color{var}u}}} is removed.


 \ \vspace{.5ex}\begin{Verbatim}[gobble=1,tabsize=2,frame=single,framerule=1.2pt,commandchars=\\@\$,label=\texttt{Treap},labelposition=topline]
	@\color@var$boolean$ remove(@\color@keyword$T$ @\color@var$x$) {
		Node<@\color@keyword$T$> @\color@var$u$ = findLast(@\color@var$x$);
		@\color@keyword$if$ (@\color@var$u$ != @\color@var$nil$ && compare(@\color@var$u$.@\color@var$x$, @\color@var$x$) == 0) {
			trickleDown(@\color@var$u$);
			splice(@\color@var$u$);
			@\color@var$return$ @\color@var$true$;
		}
		@\color@var$return$ @\color@var$false$;
	}
	@\color@keyword$void$ trickleDown(Node<@\color@keyword$T$> @\color@var$u$) {
		@\color@keyword$while$ (@\color@var$u$.@\color@var$left$ != @\color@var$nil$ || @\color@var$u$.@\color@var$right$ != @\color@var$nil$) {
			@\color@keyword$if$ (@\color@var$u$.@\color@var$left$ == @\color@var$nil$) {
				rotateLeft(@\color@var$u$);
			} @\color@keyword$else$ @\color@keyword$if$ (@\color@var$u$.@\color@var$right$ == @\color@var$nil$) {
				rotateRight(@\color@var$u$);
			} @\color@keyword$else$ @\color@keyword$if$ (@\color@var$u$.@\color@var$left$.@\color@var$p$ < @\color@var$u$.@\color@var$right$.@\color@var$p$) {
				rotateRight(@\color@var$u$);
			} @\color@keyword$else$ {
				rotateLeft(@\color@var$u$);
			}
			@\color@keyword$if$ (@\color@var$r$ == @\color@var$u$) {
				@\color@var$r$ = @\color@var$u$.@\color@var$parent$;
			}
		}
	}
\end{Verbatim}

An example of the \ensuremath{\mathtt{remove({\color{var}x})}} operation is shown in \figref{treap-remove}.
\begin{figure}
  \begin{center}
  \includegraphics[height=\QuarterHeightScaleIfNeeded]{figs/treap-delete-a} \\
  \includegraphics[height=\QuarterHeightScaleIfNeeded]{figs/treap-delete-b} \\
  \includegraphics[height=\QuarterHeightScaleIfNeeded]{figs/treap-delete-c} \\
  \includegraphics[height=\QuarterHeightScaleIfNeeded]{figs/treap-delete-d} 
  \end{center}
  \caption[Removing from a treap]{Removing the value 9 from the \texttt{Treap} in \figref{treap}.}
  \figlabel{treap-remove}
\end{figure}

The trick to analyze the running time of the \ensuremath{\mathtt{remove({\color{var}x})}} operation is
to notice that this operation reverses the \ensuremath{\mathtt{add({\color{var}x})}} operation.
In particular, if we were to reinsert \ensuremath{\mathtt{{\color{var}x}}}, using the same priority \ensuremath{\mathtt{{\color{var}u}.{\color{var}p}}},
then the \ensuremath{\mathtt{add({\color{var}x})}} operation would do exactly the same number of rotations
and would restore the \texttt{Treap} to exactly the same state it was in before
the \ensuremath{\mathtt{remove({\color{var}x})}} operation took place.  (Reading from bottom-to-top,
\figref{treap-remove} illustrates the addition of the value 9 into a
\texttt{Treap}.) This means that the expected running time of the \ensuremath{\mathtt{remove({\color{var}x})}}
on a \texttt{Treap} of size \ensuremath{\mathtt{{\color{var}n}}} is proportional to the expected running time
of the \ensuremath{\mathtt{add({\color{var}x})}} operation on a \texttt{Treap} of size $\ensuremath{\mathtt{{\color{var}n}}}-1$.  We conclude
that the expected running time of \ensuremath{\mathtt{remove({\color{var}x})}} is $O(\log \ensuremath{\mathtt{{\color{var}n}}})$.

\subsection{Summary}

The following theorem summarizes the performance of the \texttt{Treap} data
structure:

\begin{thm}
A \texttt{Treap} implements the \texttt{SSet} interface. A \texttt{Treap} supports
the operations \ensuremath{\mathtt{add({\color{var}x})}}, \ensuremath{\mathtt{remove({\color{var}x})}}, and \ensuremath{\mathtt{find({\color{var}x})}} in $O(\log \ensuremath{\mathtt{{\color{var}n}}})$
expected time per operation.
\end{thm}

It is worth comparing the \texttt{Treap} data structure to the \texttt{SkiplistSSet}
data structure.  Both implement the \texttt{SSet} operations in $O(\log \ensuremath{\mathtt{{\color{var}n}}})$
expected time per operation.  In both data structures, \ensuremath{\mathtt{add({\color{var}x})}} and
\ensuremath{\mathtt{remove({\color{var}x})}} involve a search and then a constant number of pointer changes
(see \excref{treap-rotates} below).  Thus, for both these structures,
the expected length of the search path is the critical value in assessing
their performance.  In a \texttt{SkiplistSSet}, the expected length of a search
path is
\[
     2\log \ensuremath{\mathtt{{\color{var}n}}} + O(1) \enspace ,
\]
In a \texttt{Treap}, the expected length of a search path is 
\[
    2\ln \ensuremath{\mathtt{{\color{var}n}}} +O(1) \approx 1.386\log \ensuremath{\mathtt{{\color{var}n}}}  + O(1) \enspace .
\]
Thus, the search paths in a \texttt{Treap} are considerably shorter and this
translates into noticeably faster operations on \texttt{Treap}s than \texttt{Skiplist}s.
\excref{skiplist-opt} in \chapref{skiplists} shows how the
expected length of the search path in a \texttt{Skiplist} can be reduced to
\[
     e\ln \ensuremath{\mathtt{{\color{var}n}}} + O(1) \approx 1.884\log \ensuremath{\mathtt{{\color{var}n}}} + O(1) 
\]
by using biased coin tosses.  Even with this optimization, the expected
length of search paths in a \texttt{SkiplistSSet} is noticeably longer than in
a \texttt{Treap}.

\section{Discussion and Exercises}

Random binary search trees have been studied extensively.  Devroye
\cite{d88} gives a proof of \lemref{rbs} and related results. There are
much stronger results in the literature as well, the most impressive
of which is due to Reed \cite{r03}, who shows that the expected height
of a random binary search tree is
\[
  \alpha\ln n - \beta\ln\ln n + O(1)
\]
where $\alpha\approx4.31107$ is the unique solution on the
interval $[2,\infty)$ of the equation $\alpha\ln((2e/\alpha))=1$ and
$\beta=\frac{3}{2\ln(\alpha/2)}$ .  Furthermore, the variance of the
height is constant.

The name \texttt{Treap} was coined by Seidel and Aragon \cite{as96} who discussed
\texttt{Treap}s and some of their variants.  However, their basic structure was
studied much earlier by Vuillemin \cite{v80} who called them Cartesian
trees.

One possible space-optimization of the \texttt{Treap} data structure 
is the elimination of the explicit storage of the priority \ensuremath{\mathtt{{\color{var}p}}}
in each node. Instead, the priority of a node, \ensuremath{\mathtt{{\color{var}u}}}, is computed by
hashing \ensuremath{\mathtt{{\color{var}u}}}'s address in memory\javaonly{ (in 32-bit Java, this is equivalent
to hashing \ensuremath{\mathtt{{\color{var}u}.hashCode()}})}.  Although a number of hash functions will
probably work well for this in practice, for the important parts of the
proof of \lemref{rbs} to remain valid, the hash function should be randomized
and have the \emph{min-wise independent property}:
\index{min-wise independence}%
For any distinct
values $x_1,\ldots,x_k$, each of the hash values $h(x_1),\ldots,h(x_k)$
should be distinct with high probability and, for each $i\in\{1,\ldots,k\}$,
\[
   \Pr\{h(x_i) = \min\{h(x_1),\ldots,h(x_k)\}\} \le c/k
\]
for some constant $c$.
One such class of hash functions that is easy to implement and fairly
fast is \emph{tabulation hashing} (\secref{tabulation}).
\index{tabulation hashing}%
\index{hashing!tabulation}%

Another \texttt{Treap} variant that doesn't store priorities at each node is
the randomized binary search tree
\index{randomized binary search tree}%
\index{binary search tree!randomized}%
of Mart\'\i nez and Roura \cite{mr98}.
In this variant, every node, \ensuremath{\mathtt{{\color{var}u}}}, stores the size, \ensuremath{\mathtt{{\color{var}u}.{\color{var}size}}}, of the
subtree rooted at \ensuremath{\mathtt{{\color{var}u}}}.  Both the \ensuremath{\mathtt{add({\color{var}x})}} and \ensuremath{\mathtt{remove({\color{var}x})}} algorithms are
randomized. The algorithm for adding \ensuremath{\mathtt{{\color{var}x}}} to the subtree rooted at \ensuremath{\mathtt{{\color{var}u}}}
does the following:
\begin{enumerate}
   \item With probability $1/(\ensuremath{\mathtt{size({\color{var}u})}}+1)$, the value \ensuremath{\mathtt{{\color{var}x}}} is added
   the usual way, as a leaf, and rotations are then done to bring \ensuremath{\mathtt{{\color{var}x}}}
   up to the root of this subtree.
   \item Otherwise (with probability $1-1/(\ensuremath{\mathtt{size({\color{var}u})}}+1)$), the value \ensuremath{\mathtt{{\color{var}x}}}
   is recursively added into one of the two subtrees rooted at \ensuremath{\mathtt{{\color{var}u}.{\color{var}left}}}
   or \ensuremath{\mathtt{{\color{var}u}.{\color{var}right}}}, as appropriate.
\end{enumerate}
The first case corresponds to an \ensuremath{\mathtt{add({\color{var}x})}} operation in a \texttt{Treap} where
\ensuremath{\mathtt{{\color{var}x}}}'s node receives a random priority that is smaller than any of the
\ensuremath{\mathtt{size({\color{var}u})}} priorities in \ensuremath{\mathtt{{\color{var}u}}}'s subtree, and this case occurs with exactly
the same probability.

Removing a value \ensuremath{\mathtt{{\color{var}x}}} from a randomized binary search tree is similar
to the process of removing from a \texttt{Treap}.  We find the node, \ensuremath{\mathtt{{\color{var}u}}},
that contains \ensuremath{\mathtt{{\color{var}x}}} and then perform rotations that repeatedly increase
the depth of \ensuremath{\mathtt{{\color{var}u}}} until it becomes a leaf, at which point we can splice
it from the tree.  The choice of whether to perform a left or right
rotation at each step is randomized.
\begin{enumerate}
  \item With probability \ensuremath{\mathtt{{\color{var}u}.{\color{var}left}.{\color{var}size}/({\color{var}u}.{\color{var}size}-1)}}, we perform a right
  rotation at \ensuremath{\mathtt{{\color{var}u}}}, making \ensuremath{\mathtt{{\color{var}u}.{\color{var}left}}} the root of the subtree that was
  formerly rooted at \ensuremath{\mathtt{{\color{var}u}}}.
  \item  With probability \ensuremath{\mathtt{{\color{var}u}.{\color{var}right}.{\color{var}size}/({\color{var}u}.{\color{var}size}-1)}}, we perform a left
  rotation at \ensuremath{\mathtt{{\color{var}u}}}, making \ensuremath{\mathtt{{\color{var}u}.{\color{var}right}}} the root of the subtree that was
  formerly rooted at \ensuremath{\mathtt{{\color{var}u}}}.
\end{enumerate}
Again, we can easily verify that these are exactly the same probabilities
that the removal algorithm in a \texttt{Treap} will perform a left or right
rotation of \ensuremath{\mathtt{{\color{var}u}}}.

Randomized binary search trees have the disadvantage, compared to treaps,
that when adding and removing elements they make many random choices, and
they must maintain the sizes of subtrees.  One advantage of randomized
binary search trees over treaps is that subtree sizes can serve another
useful purpose, namely to provide access by rank in $O(\log \ensuremath{\mathtt{{\color{var}n}}})$ expected
time (see \excref{treap-get}).  In comparison, the random priorities
stored in treap nodes have no use other than keeping the treap balanced.

\begin{exc}
  Illustrate the addition of 4.5 (with priority 7) and then 7.5 (with
  priority 20) on the \texttt{Treap} in \figref{treap}.
\end{exc}

\begin{exc}
  Illustrate the removal of 5 and then 7 on the \texttt{Treap} in \figref{treap}.
\end{exc}

\begin{exc}
  Prove the assertion that there are $21,964,800$ sequences that generate
  the tree on the right hand side of \figref{rbs-lvc}.  (Hint: Give a
  recursive formula for the number of sequences that generate a complete
  binary tree of height $h$ and evaluate this formula for $h=3$.)
\end{exc}

\begin{exc}
  Design and implement the \ensuremath{\mathtt{permute({\color{var}a})}} method that takes as input an
  array, \ensuremath{\mathtt{{\color{var}a}}}, that contains \ensuremath{\mathtt{{\color{var}n}}} distinct values and randomly permutes \ensuremath{\mathtt{{\color{var}a}}}.
  The method should run in $O(\ensuremath{\mathtt{{\color{var}n}}})$ time and you should prove that each
  of the $\ensuremath{\mathtt{{\color{var}n}}}!$ possible permutations of \ensuremath{\mathtt{{\color{var}a}}} is equally probable. 
\end{exc}

\begin{exc}\exclabel{treap-rotates}
  Use both parts of \lemref{rbs-treap} to prove that the expected number
  of rotations performed by an \ensuremath{\mathtt{add({\color{var}x})}} operation (and hence also a
  \ensuremath{\mathtt{remove({\color{var}x})}} operation) is $O(1)$.
\end{exc}

\begin{exc}
  Modify the \texttt{Treap} implementation given here so that it does not
  explicitly store priorities.  Instead, it should simulate them by
  hashing the \ensuremath{\mathtt{hashCode()}} of each node.
\end{exc}

\begin{exc}
  Suppose that a binary search tree stores, at each node, \ensuremath{\mathtt{{\color{var}u}}}, the height,
  \ensuremath{\mathtt{{\color{var}u}.{\color{var}height}}}, of the subtree rooted at \ensuremath{\mathtt{{\color{var}u}}}, and the size, \ensuremath{\mathtt{{\color{var}u}.{\color{var}size}}} of
  the subtree rooted at \ensuremath{\mathtt{{\color{var}u}}}. 
  \begin{enumerate}
    \item Show how, if we perform a left or right
      rotation at \ensuremath{\mathtt{{\color{var}u}}}, then these two quantities can be updated, in
      constant time, for all nodes affected by the rotation.
    \item Explain why the same result is not possible if we try to
      also store the depth, \ensuremath{\mathtt{{\color{var}u}.{\color{var}depth}}}, of each node \ensuremath{\mathtt{{\color{var}u}}}.
  \end{enumerate}
\end{exc}

\begin{exc}
  Design and implement an algorithm that constructs a \texttt{Treap} from a
  sorted array, \ensuremath{\mathtt{{\color{var}a}}}, of \ensuremath{\mathtt{{\color{var}n}}} elements.  This method should run in $O(\ensuremath{\mathtt{{\color{var}n}}})$
  worst-case time and should construct a \texttt{Treap} that is indistinguishable
  from one in which the elements of \ensuremath{\mathtt{{\color{var}a}}} were added one at a time using
  the \ensuremath{\mathtt{add({\color{var}x})}} method.
\end{exc}


\begin{exc}
  \index{finger}%
  \index{finger search!in a treap}%
  This exercise works out the details of how one can efficiently search
  a \texttt{Treap} given a pointer that is close to the node we are searching for.
  \begin{enumerate}
    \item Design and implement a \texttt{Treap} implementation in which each
      node keeps track of the minimum and maximum values in its subtree.
    \item Using this extra information, add a \ensuremath{\mathtt{fingerFind({\color{var}x},{\color{var}u})}} method
      that executes the \ensuremath{\mathtt{find({\color{var}x})}} operation with the help of a pointer
      to the node \ensuremath{\mathtt{{\color{var}u}}} (which is hopefully not far from the node that
      contains \ensuremath{\mathtt{{\color{var}x}}}).  This operation should start at \ensuremath{\mathtt{{\color{var}u}}} and walk upwards
      until it reaches a node \ensuremath{\mathtt{{\color{var}w}}} such that $\ensuremath{\mathtt{{\color{var}w}.{\color{var}min}}}\le \ensuremath{\mathtt{{\color{var}x}}}\le \ensuremath{\mathtt{{\color{var}w}.{\color{var}max}}}$.
      From that point onwards, it should perform a standard search
      for \ensuremath{\mathtt{{\color{var}x}}} starting from \ensuremath{\mathtt{{\color{var}w}}}.  (One can show that \ensuremath{\mathtt{fingerFind({\color{var}x},{\color{var}u})}}
      takes $O(1+\log r)$ time, where $r$ is the number of elements in
      the treap whose value is between \ensuremath{\mathtt{{\color{var}x}}} and \ensuremath{\mathtt{{\color{var}u}.{\color{var}x}}}.)
    \item Extend your implementation into a version of a treap that
      starts all its \ensuremath{\mathtt{find({\color{var}x})}} operations from the node most recently
      found by \ensuremath{\mathtt{find({\color{var}x})}}.
  \end{enumerate}
\end{exc}

\begin{exc}\exclabel{treap-get}
  Design and implement a version of a \texttt{Treap} that includes a \ensuremath{\mathtt{get({\color{var}i})}}
  operation that returns the key with rank \ensuremath{\mathtt{{\color{var}i}}} in the \texttt{Treap}.  (Hint:
  Have each node, \ensuremath{\mathtt{{\color{var}u}}}, keep track of the size of the subtree rooted
  at \ensuremath{\mathtt{{\color{var}u}}}.)
\end{exc}

\begin{exc}
  \index{TreapList@\texttt{TreapList}}%
  Implement a \texttt{TreapList}, an implementation of the \texttt{List} interface
  as a treap.  Each node in the treap should store a list item, and an
  in-order traversal of the treap finds the items in the same order that
  they occur in the list.  All the \texttt{List} operations \ensuremath{\mathtt{get({\color{var}i})}}, \ensuremath{\mathtt{set({\color{var}i},{\color{var}x})}},
  \ensuremath{\mathtt{add({\color{var}i},{\color{var}x})}} and \ensuremath{\mathtt{remove({\color{var}i})}} should run in $O(\log \ensuremath{\mathtt{{\color{var}n}}})$ expected time.
\end{exc}



\begin{exc}\exclabel{treap-split}
  Design and implement a version of a \texttt{Treap} that supports the \ensuremath{\mathtt{split({\color{var}x})}}
  operation.  This operation removes all values from the \texttt{Treap} that
  are greater than \ensuremath{\mathtt{{\color{var}x}}} and returns a second \texttt{Treap} that contains all
  the removed values.

  \noindent Example: the code \ensuremath{\mathtt{{\color{var}t2} = {\color{var}t}.split({\color{var}x})}} removes from \ensuremath{\mathtt{{\color{var}t}}} all values
  greater than \ensuremath{\mathtt{{\color{var}x}}} and returns a new \texttt{Treap} \ensuremath{\mathtt{{\color{var}t2}}} containing all
  these values. The \ensuremath{\mathtt{split({\color{var}x})}} operation should run in $O(\log \ensuremath{\mathtt{{\color{var}n}}})$
  expected time.

  \noindent Warning: For this modification to work properly and still allow the
  \ensuremath{\mathtt{size()}} method to run in constant time, it is necessary to implement
  the modifications in \excref{treap-get}.
\end{exc}

\begin{exc}\exclabel{treap-join}
  Design and implement a version of a \texttt{Treap} that supports the
  \ensuremath{\mathtt{absorb({\color{var}t2})}} operation, which can be thought of as the inverse of
  the \ensuremath{\mathtt{split({\color{var}x})}} operation.  This operation removes all values from the
  \texttt{Treap} \ensuremath{\mathtt{{\color{var}t2}}} and adds them to the receiver.  This operation presupposes
  that the smallest value in \ensuremath{\mathtt{{\color{var}t2}}} is greater than the largest value in
  the receiver.  The \ensuremath{\mathtt{absorb({\color{var}t2})}} operation should run in $O(\log \ensuremath{\mathtt{{\color{var}n}}})$
  expected time.
\end{exc}

\begin{exc}
  Implement Martinez's randomized binary search trees, as discussed in
  this section.  Compare the performance of your implementation with
  that of the \texttt{Treap} implementation.
\end{exc}


\clearpage

\clearpage
\chapter{Scapegoat Trees}
\chaplabel{scapegoat}

In this chapter, we study a binary search tree data structure, the
\texttt{ScapegoatTree}.  This structure is based on the common wisdom that,
when something goes wrong, the first thing people tend to do is find
someone to blame (the \emph{scapegoat}).
\index{scapegoat}%
Once blame is firmly
established, we can leave the scapegoat to fix the problem.

A \texttt{ScapegoatTree} keeps itself balanced by \emph{partial rebuilding
operations}.
\index{partial rebuilding}%
\index{binary search tree!partial rebuilding}%
During a partial rebuilding operation, an entire subtree is
deconstructed and rebuilt into a perfectly balanced subtree.  There are
many ways of rebuilding a subtree rooted at node \ensuremath{\mathtt{{\color{var}u}}} into a perfectly
balanced tree.  One of the simplest is to traverse \ensuremath{\mathtt{{\color{var}u}}}'s subtree,
gathering all its nodes into an array, \ensuremath{\mathtt{{\color{var}a}}}, and then to recursively
build a balanced subtree using \ensuremath{\mathtt{{\color{var}a}}}.  If we let $\ensuremath{\mathtt{{\color{var}m}}}=\ensuremath{\mathtt{{\color{var}a}.{\color{var}length}}}/2$,
then the element \ensuremath{\mathtt{{\color{var}a}[{\color{var}m}]}} becomes the root of the new subtree,
$\ensuremath{\mathtt{{\color{var}a}}}[0],\ldots,\ensuremath{\mathtt{{\color{var}a}}}[\ensuremath{\mathtt{{\color{var}m}}}-1]$ get stored recursively in the left subtree
and $\ensuremath{\mathtt{{\color{var}a}}}[\ensuremath{\mathtt{{\color{var}m}}}+1],\ldots,\ensuremath{\mathtt{{\color{var}a}}}[\ensuremath{\mathtt{{\color{var}a}.{\color{var}length}}}-1]$ get stored recursively in the
right subtree.


 \ \vspace{.5ex}\begin{Verbatim}[gobble=1,tabsize=2,frame=single,framerule=1.2pt,commandchars=\\@\$,label=\texttt{ScapegoatTree},labelposition=topline]
	@\color@keyword$void$ rebuild(Node<@\color@keyword$T$> @\color@var$u$) {
		@\color@keyword$int$ @\color@var$ns$ = size(@\color@var$u$);
		Node<@\color@keyword$T$> @\color@var$p$ = @\color@var$u$.@\color@var$parent$;
		Node<@\color@keyword$T$>[] @\color@var$a$ =  Array.newInstance(Node.@\color@keyword$class$, @\color@var$ns$);
		packIntoArray(@\color@var$u$, @\color@var$a$, 0);
		@\color@keyword$if$ (@\color@var$p$ == @\color@var$nil$) {
			@\color@var$r$ = buildBalanced(@\color@var$a$, 0, @\color@var$ns$);
			@\color@var$r$.@\color@var$parent$ = @\color@var$nil$;
		} @\color@keyword$else$ @\color@keyword$if$ (@\color@var$p$.@\color@var$right$ == @\color@var$u$) {
			@\color@var$p$.@\color@var$right$ = buildBalanced(@\color@var$a$, 0, @\color@var$ns$);
			@\color@var$p$.@\color@var$right$.@\color@var$parent$ = @\color@var$p$;
		} @\color@keyword$else$ {
			@\color@var$p$.@\color@var$left$ = buildBalanced(@\color@var$a$, 0, @\color@var$ns$);
			@\color@var$p$.@\color@var$left$.@\color@var$parent$ = @\color@var$p$;
		}
	}
	@\color@keyword$int$ packIntoArray(Node<@\color@keyword$T$> @\color@var$u$, Node<@\color@keyword$T$>[] @\color@var$a$, @\color@keyword$int$ @\color@var$i$) {
		@\color@keyword$if$ (@\color@var$u$ == @\color@var$nil$) {
			@\color@var$return$ @\color@var$i$;
		}
		@\color@var$i$ = packIntoArray(@\color@var$u$.@\color@var$left$, @\color@var$a$, @\color@var$i$);
		@\color@var$a$[@\color@var$i$++] = @\color@var$u$;
		@\color@var$return$ packIntoArray(@\color@var$u$.@\color@var$right$, @\color@var$a$, @\color@var$i$);
	}
	Node<@\color@keyword$T$> buildBalanced(Node<@\color@keyword$T$>[] @\color@var$a$, @\color@keyword$int$ @\color@var$i$, @\color@keyword$int$ @\color@var$ns$) {
		@\color@keyword$if$ (@\color@var$ns$ == 0)
			@\color@var$return$ @\color@var$nil$;
		@\color@keyword$int$ @\color@var$m$ = @\color@var$ns$ / 2;
		@\color@var$a$[@\color@var$i$ + @\color@var$m$].@\color@var$left$ = buildBalanced(@\color@var$a$, @\color@var$i$, @\color@var$m$);
		@\color@keyword$if$ (@\color@var$a$[@\color@var$i$ + @\color@var$m$].@\color@var$left$ != @\color@var$nil$)
			@\color@var$a$[@\color@var$i$ + @\color@var$m$].@\color@var$left$.@\color@var$parent$ = @\color@var$a$[@\color@var$i$ + @\color@var$m$];
		@\color@var$a$[@\color@var$i$ + @\color@var$m$].@\color@var$right$ = buildBalanced(@\color@var$a$, @\color@var$i$ + @\color@var$m$ + 1, @\color@var$ns$ - @\color@var$m$ - 1);
		@\color@keyword$if$ (@\color@var$a$[@\color@var$i$ + @\color@var$m$].@\color@var$right$ != @\color@var$nil$)
			@\color@var$a$[@\color@var$i$ + @\color@var$m$].@\color@var$right$.@\color@var$parent$ = @\color@var$a$[@\color@var$i$ + @\color@var$m$];
		@\color@var$return$ @\color@var$a$[@\color@var$i$ + @\color@var$m$];
	}
\end{Verbatim}

A call to \ensuremath{\mathtt{rebuild({\color{var}u})}} takes $O(\ensuremath{\mathtt{size({\color{var}u})}})$ time.  The resulting subtree
has minimum height; there is no tree of smaller height that
has \ensuremath{\mathtt{size({\color{var}u})}} nodes.


\section{\texttt{ScapegoatTree}: A Binary Search Tree with Partial Rebuilding}
\seclabel{scapegoattree}


\index{ScapegoatTree@\texttt{ScapegoatTree}}%
A \texttt{ScapegoatTree} is a \texttt{BinarySearchTree} that, in addition to keeping
track of the number, \ensuremath{\mathtt{{\color{var}n}}}, of  nodes in the tree also keeps a counter, \ensuremath{\mathtt{{\color{var}q}}},
that maintains an upper-bound on the number of nodes. 


 \ \vspace{.5ex}\begin{Verbatim}[gobble=1,tabsize=2,frame=single,framerule=1.2pt,commandchars=\\@\$,label=\texttt{ScapegoatTree},labelposition=topline]
	@\color@keyword$int$ @\color@var$q$;
\end{Verbatim}

At all times, \ensuremath{\mathtt{{\color{var}n}}} and \ensuremath{\mathtt{{\color{var}q}}} obey the following inequalities:
\[
      \ensuremath{\mathtt{{\color{var}q}}}/2 \le  \ensuremath{\mathtt{{\color{var}n}}} \le \ensuremath{\mathtt{{\color{var}q}}}  \enspace .
\]
In addition, a \texttt{ScapegoatTree} has logarithmic height; at all times, the height of the scapegoat tree does not exceed
\begin{equation}
     \log_{3/2} \ensuremath{\mathtt{{\color{var}q}}} \le \log_{3/2} 2\ensuremath{\mathtt{{\color{var}n}}} < \log_{3/2} \ensuremath{\mathtt{{\color{var}n}}} + 2\enspace .
     \eqlabel{scapegoat-height}
\end{equation}
Even with this constraint, a \texttt{ScapegoatTree} can look surprisingly unbalanced.  The tree in \figref{scapegoat-example} has $\ensuremath{\mathtt{{\color{var}q}}}=\ensuremath{\mathtt{{\color{var}n}}}=10$ and height $5<\log_{3/2}10 \approx 5.679$.

\begin{figure}
  \begin{center}
    \includegraphics[scale=0.90909]{figs/scapegoat-insert-1}
  \end{center}
  \caption[A ScapegoatTree]{A \texttt{ScapegoatTree} with 10 nodes and height 5.}
  \figlabel{scapegoat-example}
\end{figure}

Implementing the \ensuremath{\mathtt{find({\color{var}x})}} operation in a \texttt{ScapegoatTree} is done
using the standard algorithm for searching in a \texttt{BinarySearchTree}
(see \secref{binarysearchtree}).  This takes time proportional to the
height of the tree which, by \myeqref{scapegoat-height} is $O(\log \ensuremath{\mathtt{{\color{var}n}}})$.

To implement the \ensuremath{\mathtt{add({\color{var}x})}} operation, we first increment \ensuremath{\mathtt{{\color{var}n}}} and \ensuremath{\mathtt{{\color{var}q}}}
and then use the usual algorithm for adding \ensuremath{\mathtt{{\color{var}x}}} to a binary search
tree; we search for \ensuremath{\mathtt{{\color{var}x}}} and then add a new leaf \ensuremath{\mathtt{{\color{var}u}}} with $\ensuremath{\mathtt{{\color{var}u}.{\color{var}x}}}=\ensuremath{\mathtt{{\color{var}x}}}$.
At this point, we may get lucky and the depth of \ensuremath{\mathtt{{\color{var}u}}} might not exceed
$\log_{3/2}\ensuremath{\mathtt{{\color{var}q}}}$. If so, then we leave well enough alone and don't do
anything else.

Unfortunately, it will sometimes happen that $\ensuremath{\mathtt{depth({\color{var}u})}} > \log_{3/2}
\ensuremath{\mathtt{{\color{var}q}}}$.  In this case, we need to reduce the height.  This isn't a big
job; there is only one node, namely \ensuremath{\mathtt{{\color{var}u}}}, whose depth exceeds $\log_{3/2}
\ensuremath{\mathtt{{\color{var}q}}}$.  To fix \ensuremath{\mathtt{{\color{var}u}}}, we walk from \ensuremath{\mathtt{{\color{var}u}}} back up to the root looking for a
\emph{scapegoat}, \ensuremath{\mathtt{{\color{var}w}}}.  The scapegoat, \ensuremath{\mathtt{{\color{var}w}}}, is a very unbalanced node.
It has the property that
\begin{equation}
   \frac{\ensuremath{\mathtt{size({\color{var}w}.{\color{var}child})}}}{\ensuremath{\mathtt{size({\color{var}w})}}} > \frac{2}{3} \enspace ,
   \eqlabel{scapegoat}
\end{equation}
where \ensuremath{\mathtt{{\color{var}w}.{\color{var}child}}} is the child of \ensuremath{\mathtt{{\color{var}w}}} on the path from the root to \ensuremath{\mathtt{{\color{var}u}}}.
We'll very shortly prove that a scapegoat exists.  For now, we can
take it for granted.  Once we've found the scapegoat \ensuremath{\mathtt{{\color{var}w}}}, we completely
destroy the subtree rooted at \ensuremath{\mathtt{{\color{var}w}}} and rebuild it into a perfectly balanced
binary search tree.  We know, from \myeqref{scapegoat}, that, even before
the addition of \ensuremath{\mathtt{{\color{var}u}}}, \ensuremath{\mathtt{{\color{var}w}}}'s subtree was not a complete binary tree.
Therefore, when we rebuild \ensuremath{\mathtt{{\color{var}w}}}, the height decreases by at least 1 so that the height of the \texttt{ScapegoatTree} is once again at most $\log_{3/2}\ensuremath{\mathtt{{\color{var}q}}}$.



 \ \vspace{.5ex}\begin{Verbatim}[gobble=1,tabsize=2,frame=single,framerule=1.2pt,commandchars=\\@\$,label=\texttt{ScapegoatTree},labelposition=topline]
	@\color@var$boolean$ add(@\color@keyword$T$ @\color@var$x$) {
		@\color@comment$// first do basic insertion keeping track of depth$
		Node<@\color@keyword$T$> @\color@var$u$ = newNode(@\color@var$x$);
		@\color@keyword$int$ @\color@var$d$ = addWithDepth(@\color@var$u$);
		@\color@keyword$if$ (@\color@var$d$ > log32(@\color@var$q$)) {
			@\color@comment$// depth exceeded, find scapegoat$
			Node<@\color@keyword$T$> @\color@var$w$ = @\color@var$u$.@\color@var$parent$;
			@\color@keyword$while$ (3*size(@\color@var$w$) <= 2*size(@\color@var$w$.@\color@var$parent$))
				@\color@var$w$ = @\color@var$w$.@\color@var$parent$;
			rebuild(@\color@var$w$.@\color@var$parent$);
		}
		@\color@var$return$ @\color@var$d$ >= 0;
	}
\end{Verbatim}


\begin{figure}
  \begin{center}
    \begin{tabular}{cc}
      \includegraphics[scale=0.90909]{figs/scapegoat-insert-3} &
      \includegraphics[scale=0.90909]{figs/scapegoat-insert-4} 
    \end{tabular}
  \end{center}
  \caption[Adding to a scapegoat tree]{Inserting 3.5 into a \texttt{ScapegoatTree} increases its height to 6, which violates \myeqref{scapegoat-height} since $6 > \log_{3/2} 11 \approx 5.914$.  A scapegoat is found at the node containing 5.}
\end{figure}
If we ignore the cost of finding the scapegoat \ensuremath{\mathtt{{\color{var}w}}} and rebuilding the
subtree rooted at \ensuremath{\mathtt{{\color{var}w}}}, then the running time of \ensuremath{\mathtt{add({\color{var}x})}} is dominated
by the initial search, which takes $O(\log \ensuremath{\mathtt{{\color{var}q}}}) = O(\log \ensuremath{\mathtt{{\color{var}n}}})$ time.
We will account for the cost of finding the scapegoat and rebuilding
using amortized analysis in the next section.

The implementation of \ensuremath{\mathtt{remove({\color{var}x})}} in a \texttt{ScapegoatTree} is very simple.
We search for \ensuremath{\mathtt{{\color{var}x}}} and remove it using the usual algorithm for removing a
node from a \texttt{BinarySearchTree}.  (Note that this can never increase the
height of the tree.)  Next, we decrement \ensuremath{\mathtt{{\color{var}n}}}, but leave \ensuremath{\mathtt{{\color{var}q}}} unchanged.
Finally, we check if $\ensuremath{\mathtt{{\color{var}q}}} > 2\ensuremath{\mathtt{{\color{var}n}}}$ and, if so, then we \emph{rebuild the entire
tree} into a perfectly balanced binary search tree and set $\ensuremath{\mathtt{{\color{var}q}}}=\ensuremath{\mathtt{{\color{var}n}}}$.


 \ \vspace{.5ex}\begin{Verbatim}[gobble=1,tabsize=2,frame=single,framerule=1.2pt,commandchars=\\@\$,label=\texttt{ScapegoatTree},labelposition=topline]
	@\color@var$boolean$ remove(@\color@keyword$T$ @\color@var$x$) {
		@\color@keyword$if$ (@\color@var$super$.remove(@\color@var$x$)) {
			@\color@keyword$if$ (2*@\color@var$n$ < @\color@var$q$) {
				rebuild(@\color@var$r$);
				@\color@var$q$ = @\color@var$n$;
			}
			@\color@var$return$ @\color@var$true$;
		}
		@\color@var$return$ @\color@var$false$;
	}
\end{Verbatim}

Again, if we ignore the cost of rebuilding, the running time of the
\ensuremath{\mathtt{remove({\color{var}x})}} operation is proportional to the height of the tree, and is
therefore $O(\log \ensuremath{\mathtt{{\color{var}n}}})$.

\subsection{Analysis of Correctness and Running-Time}

In this section, we analyze the correctness and amortized running time
of operations on a \texttt{ScapegoatTree}.  We first prove the correctness by
showing that, when the \ensuremath{\mathtt{add({\color{var}x})}} operation results in a node that violates
Condition \myeqref{scapegoat-height}, then we can always find a scapegoat:

\begin{lem}
  Let \ensuremath{\mathtt{{\color{var}u}}} be a node of depth $h>\log_{3/2} \ensuremath{\mathtt{{\color{var}q}}}$ in a \texttt{ScapegoatTree}.
  Then there exists a node $\ensuremath{\mathtt{{\color{var}w}}}$ on the path from \ensuremath{\mathtt{{\color{var}u}}} to the root
  such that
  \[
     \frac{\ensuremath{\mathtt{size({\color{var}w})}}}{\ensuremath{\mathtt{size(parent({\color{var}w}))}}} > 2/3 \enspace .
  \]
\end{lem}

\begin{proof}
  Suppose, for the sake of contradiction, that this is not the case, and
  \[
     \frac{\ensuremath{\mathtt{size({\color{var}w})}}}{\ensuremath{\mathtt{size(parent({\color{var}w}))}}} \le 2/3 \enspace .
  \]
  for all nodes \ensuremath{\mathtt{{\color{var}w}}} on the path from \ensuremath{\mathtt{{\color{var}u}}} to the root.  Denote the path
  from the root to \ensuremath{\mathtt{{\color{var}u}}} as $\ensuremath{\mathtt{{\color{var}r}}}=\ensuremath{\mathtt{{\color{var}u}}}_0,\ldots,\ensuremath{\mathtt{{\color{var}u}}}_h=\ensuremath{\mathtt{{\color{var}u}}}$.  Then, we have
  $\ensuremath{\mathtt{size({\color{var}u}}}_0\ensuremath{\mathtt{)}}=\ensuremath{\mathtt{{\color{var}n}}}$,
  $\ensuremath{\mathtt{size({\color{var}u}}}_1\ensuremath{\mathtt{)}}\le\frac{2}{3}\ensuremath{\mathtt{{\color{var}n}}}$, 
  $\ensuremath{\mathtt{size({\color{var}u}}}_2\ensuremath{\mathtt{)}}\le\frac{4}{9}\ensuremath{\mathtt{{\color{var}n}}}$ and, more generally,
  \[
  \ensuremath{\mathtt{size({\color{var}u}}}_i\ensuremath{\mathtt{)}}\le\left(\frac{2}{3}\right)^i\ensuremath{\mathtt{{\color{var}n}}} \enspace .
  \]
  But this gives a contradiction, since $\ensuremath{\mathtt{size({\color{var}u})}}\ge 1$, hence
  \[
    1 \le \ensuremath{\mathtt{size({\color{var}u})}} \le \left(\frac{2}{3}\right)^h\ensuremath{\mathtt{{\color{var}n}}}
   < \left(\frac{2}{3}\right)^{\log_{3/2} \ensuremath{\mathtt{{\color{var}q}}}}\ensuremath{\mathtt{{\color{var}n}}}
   \le \left(\frac{2}{3}\right)^{\log_{3/2} \ensuremath{\mathtt{{\color{var}n}}}}\ensuremath{\mathtt{{\color{var}n}}}
   = \left(\frac{1}{\ensuremath{\mathtt{{\color{var}n}}}}\right) \ensuremath{\mathtt{{\color{var}n}}}
   = 1 \enspace . \qedhere
  \]
\end{proof}

Next, we analyze the parts of the running time that are not yet
accounted for.  There are two parts:  The cost of calls to \ensuremath{\mathtt{size({\color{var}u})}}
when searching for scapegoat nodes, and the cost of calls to \ensuremath{\mathtt{rebuild({\color{var}w})}}
when we find a scapegoat \ensuremath{\mathtt{{\color{var}w}}}.  The cost of calls to \ensuremath{\mathtt{size({\color{var}u})}} can be
related to the cost of calls to \ensuremath{\mathtt{rebuild({\color{var}w})}}, as follows:
\begin{lem}
During a call to \ensuremath{\mathtt{add({\color{var}x})}} in a \texttt{ScapegoatTree}, the cost of finding the scapegoat \ensuremath{\mathtt{{\color{var}w}}} and rebuilding the subtree rooted at \ensuremath{\mathtt{{\color{var}w}}} is $O(\ensuremath{\mathtt{size({\color{var}w})}})$.
\end{lem}

\begin{proof}
The cost of rebuilding the scapegoat node \ensuremath{\mathtt{{\color{var}w}}}, once we find it, is
$O(\ensuremath{\mathtt{size({\color{var}w})}})$.  When searching for the scapegoat node, we call \ensuremath{\mathtt{size({\color{var}u})}}
on a sequence of nodes $\ensuremath{\mathtt{{\color{var}u}}}_0,\ldots,\ensuremath{\mathtt{{\color{var}u}}}_k$ until we find the scapegoat
$\ensuremath{\mathtt{{\color{var}u}}}_k=\ensuremath{\mathtt{{\color{var}w}}}$.  However, since $\ensuremath{\mathtt{{\color{var}u}}}_k$ is the first node in this sequence
that is a scapegoat, we know that
\[
  \ensuremath{\mathtt{size({\color{var}u}}}_{i}\ensuremath{\mathtt{)}} < \frac{2}{3}\ensuremath{\mathtt{size({\color{var}u}}}_{i+1}\ensuremath{\mathtt{)}}
\]
for all $i\in\{0,\ldots,k-2\}$.  Therefore, the cost of all calls to \ensuremath{\mathtt{size({\color{var}u})}} is
\begin{eqnarray*}
 O\left( \sum_{i=0}^k \ensuremath{\mathtt{size({\color{var}u}}}_{k-i}\ensuremath{\mathtt{)}} \right)
 &=& O\left(
  \ensuremath{\mathtt{size({\color{var}u}}}_k\ensuremath{\mathtt{)}} 
  + \sum_{i=0}^{k-1} \ensuremath{\mathtt{size({\color{var}u}}}_{k-i-1}\ensuremath{\mathtt{)}}
  \right) \\
 &=& O\left(
  \ensuremath{\mathtt{size({\color{var}u}}}_k\ensuremath{\mathtt{)}} 
  + \sum_{i=0}^{k-1} \left(\frac{2}{3}\right)^i\ensuremath{\mathtt{size({\color{var}u}}}_{k}\ensuremath{\mathtt{)}}
  \right) \\
&=& O\left(
  \ensuremath{\mathtt{size({\color{var}u}}}_k\ensuremath{\mathtt{)}}\left(1+ 
   \sum_{i=0}^{k-1} \left(\frac{2}{3}\right)^i
  \right)\right) \\
&=& O(\ensuremath{\mathtt{size({\color{var}u}}}_k\ensuremath{\mathtt{)}}) = O(\ensuremath{\mathtt{size({\color{var}w})}}) \enspace ,
\end{eqnarray*}
where the last line follows from the fact that the sum is a geometrically decreasing series.
\end{proof}

All that remains is to prove an upper-bound on the cost of all calls to
\ensuremath{\mathtt{rebuild({\color{var}u})}} during a sequence of $m$ operations:

\begin{lem}\lemlabel{scapegoat-amortized}
  Starting with an empty \texttt{ScapegoatTree} any sequence of $m$ \ensuremath{\mathtt{add({\color{var}x})}}
  and \ensuremath{\mathtt{remove({\color{var}x})}} operations causes at most $O(m\log m)$ time to be used
  by \ensuremath{\mathtt{rebuild({\color{var}u})}} operations.
\end{lem}

\begin{proof}
  To prove this, we will use a \emph{credit scheme}.
  \index{credit scheme}%
  We imagine that each node
  stores a number of credits.  Each credit can pay for some constant,
  $c$, units of time spent rebuilding.  The scheme gives out a total of
  $O(m\log m)$ credits and every call to \ensuremath{\mathtt{rebuild({\color{var}u})}} is paid for with
  credits stored at \ensuremath{\mathtt{{\color{var}u}}}.

  During an insertion or deletion, we give one credit to each node on the
  path to the inserted node, or deleted node, \ensuremath{\mathtt{{\color{var}u}}}.  In this way we hand
  out at most $\log_{3/2}\ensuremath{\mathtt{{\color{var}q}}}\le \log_{3/2}m$ credits per operation.
  During a deletion we also store an additional credit ``on the side.''
  Thus, in total we give out at most $O(m\log m)$ credits.  All that
  remains is to show that these credits are sufficient to pay for all
  calls to \ensuremath{\mathtt{rebuild({\color{var}u})}}.

  If we call \ensuremath{\mathtt{rebuild({\color{var}u})}} during an insertion, it is because \ensuremath{\mathtt{{\color{var}u}}} is
  a scapegoat.  Suppose, without loss of generality, that
  \[
    \frac{\ensuremath{\mathtt{size({\color{var}u}.{\color{var}left})}}}{\ensuremath{\mathtt{size({\color{var}u})}}} > \frac{2}{3} \enspace .
  \]
  Using the fact that
  \[
    \ensuremath{\mathtt{size({\color{var}u})}} = 1 + \ensuremath{\mathtt{size({\color{var}u}.{\color{var}left})}} + \ensuremath{\mathtt{size({\color{var}u}.{\color{var}right})}} 
  \]
  we deduce that
  \[
    \frac{1}{2}\ensuremath{\mathtt{size({\color{var}u}.{\color{var}left})}} > \ensuremath{\mathtt{size({\color{var}u}.{\color{var}right})}}  \enspace 
  \]
  and therefore
  \[
    \ensuremath{\mathtt{size({\color{var}u}.{\color{var}left})}} - \ensuremath{\mathtt{size({\color{var}u}.{\color{var}right})}} > \frac{1}{2}\ensuremath{\mathtt{size({\color{var}u}.{\color{var}left})}} >
    \frac{1}{3}\ensuremath{\mathtt{size({\color{var}u})}}  \enspace .
  \]
  Now, the last time a subtree containing \ensuremath{\mathtt{{\color{var}u}}} was rebuilt (or when \ensuremath{\mathtt{{\color{var}u}}}
  was inserted, if a subtree containing \ensuremath{\mathtt{{\color{var}u}}} was never rebuilt), we had
  \[
    \ensuremath{\mathtt{size({\color{var}u}.{\color{var}left})}} - \ensuremath{\mathtt{size({\color{var}u}.{\color{var}right})}} \le 1 \enspace .
  \]
  Therefore, the number of \ensuremath{\mathtt{add({\color{var}x})}} or \ensuremath{\mathtt{remove({\color{var}x})}} operations that have
  affected \ensuremath{\mathtt{{\color{var}u}.{\color{var}left}}} or \ensuremath{\mathtt{{\color{var}u}.{\color{var}right}}} since then is at least
  \[
    \frac{1}{3}\ensuremath{\mathtt{size({\color{var}u})}} - 1 \enspace . 
  \]
  and there are therefore at least this many credits stored at \ensuremath{\mathtt{{\color{var}u}}}
  that are available to pay for the $O(\ensuremath{\mathtt{size({\color{var}u})}})$ time it takes to
  call \ensuremath{\mathtt{rebuild({\color{var}u})}}.

  If we call \ensuremath{\mathtt{rebuild({\color{var}u})}} during a deletion, it is because $\ensuremath{\mathtt{{\color{var}q}}} > 2\ensuremath{\mathtt{{\color{var}n}}}$.
  In this case, we have $\ensuremath{\mathtt{{\color{var}q}}}-\ensuremath{\mathtt{{\color{var}n}}}> \ensuremath{\mathtt{{\color{var}n}}}$ credits stored ``on the side,'' and
  we use these to pay for the $O(\ensuremath{\mathtt{{\color{var}n}}})$ time it takes to rebuild the root.
  This completes the proof.
\end{proof}

\subsection{Summary}
The following theorem summarizes the performance of the \texttt{ScapegoatTree} data structure:

\begin{thm}\thmlabel{scapegoat}
  A \texttt{ScapegoatTree} implements the \texttt{SSet} interface. Ignoring the cost
  of \ensuremath{\mathtt{rebuild({\color{var}u})}} operations, a \texttt{ScapegoatTree} supports the operations
  \ensuremath{\mathtt{add({\color{var}x})}}, \ensuremath{\mathtt{remove({\color{var}x})}}, and \ensuremath{\mathtt{find({\color{var}x})}} in $O(\log \ensuremath{\mathtt{{\color{var}n}}})$ time per operation.
  
  Furthermore, beginning with an empty \texttt{ScapegoatTree}, any sequence of $m$
  \ensuremath{\mathtt{add({\color{var}x})}} and \ensuremath{\mathtt{remove({\color{var}x})}} operations results in a total of $O(m\log m)$
  time spent during all calls to \ensuremath{\mathtt{rebuild({\color{var}u})}}.
\end{thm}

\section{Discussion and Exercises}

The term \emph{scapegoat tree} is due to Galperin and Rivest \cite{gr93},
who define and analyze these trees.  However, the same structure
was discovered earlier by Andersson \cite{a89,a99}, who called them
\emph{general balanced trees}
\index{general balanced tree}%
since they can have any shape as long as
their height is small.

Experimenting with the \texttt{ScapegoatTree} implementation will reveal that
it is often considerably slower than the other \texttt{SSet} implementations
in this book. This may be somewhat surprising, since height bound of
\[
   \log_{3/2}\ensuremath{\mathtt{{\color{var}q}}} \approx 1.709\log \ensuremath{\mathtt{{\color{var}n}}} + O(1)
\] 
is better than the expected length of a search path in a \texttt{Skiplist} and
not too far from that of a \texttt{Treap}.  The implementation could be optimized
by storing the sizes of subtrees explicitly at each node or by reusing
already computed subtree sizes (Exercises~\ref{exc:scapegoat-quicksize}
and \ref{exc:scapegoat-explicitsize}).  Even with these optimizations,
there will always be sequences of \ensuremath{\mathtt{add({\color{var}x})}} and \ensuremath{\mathtt{delete({\color{var}x})}} operation for
which a \texttt{ScapegoatTree} takes longer than other \texttt{SSet} implementations.

This gap in performance is due to the fact that, unlike the other \texttt{SSet}
implementations discussed in this book, a \texttt{ScapegoatTree} can spend a lot
of time restructuring itself.  \excref{scapegoat-nlogn} asks you to prove
that there are sequences of \ensuremath{\mathtt{{\color{var}n}}} operations in which a \texttt{ScapegoatTree}
will spend on the order of $\ensuremath{\mathtt{{\color{var}n}}}\log \ensuremath{\mathtt{{\color{var}n}}}$ time in calls to \ensuremath{\mathtt{rebuild({\color{var}u})}}.
This is in contrast to other \texttt{SSet} implementations discussed in this
book, which only make $O(\ensuremath{\mathtt{{\color{var}n}}})$ structural changes during a sequence
of \ensuremath{\mathtt{{\color{var}n}}} operations.  This is, unfortunately, a necessary consequence of
the fact that a \texttt{ScapegoatTree} does all its restructuring by calls to
\ensuremath{\mathtt{rebuild({\color{var}u})}} \cite{d90}.

Despite their lack of performance, there are applications in which a
\texttt{ScapegoatTree} could be the right choice.  This would occur any time
there is additional data associated with nodes that cannot be updated
in constant time when a rotation is performed, but that can be updated
during a \ensuremath{\mathtt{rebuild({\color{var}u})}} operation.  In such cases, the \texttt{ScapegoatTree}
and related structures based on partial rebuilding may work.  An example of such an application is outlined in \excref{list-order-maintenance}.

\begin{exc}
  Illustrate the addition of the values 1.5 and then 1.6 on the
  \texttt{ScapegoatTree} in \figref{scapegoat-example}.
\end{exc}

\begin{exc}
  Illustrate what happens when the sequence $1,5,2,4,3$ is added to an
  empty \texttt{ScapegoatTree}, and show where the credits described in the
  proof of \lemref{scapegoat-amortized} go, and how they are used during
  this sequence of additions.
\end{exc}

\begin{exc}\exclabel{scapegoat-nlogn}
  Show that, if we start with an empty \texttt{ScapegoatTree} and call \ensuremath{\mathtt{add({\color{var}x})}}
  for $\ensuremath{\mathtt{{\color{var}x}}}=1,2,3,\ldots,\ensuremath{\mathtt{{\color{var}n}}}$, then the total time spent during calls to
  \ensuremath{\mathtt{rebuild({\color{var}u})}} is at least $c\ensuremath{\mathtt{{\color{var}n}}}\log \ensuremath{\mathtt{{\color{var}n}}}$ for some constant $c>0$.
\end{exc}

\begin{exc}
  The \texttt{ScapegoatTree}, as described in this chapter, guarantees that the
  length of the search path does not exceed $\log_{3/2}\ensuremath{\mathtt{{\color{var}q}}}$.
  \begin{enumerate}
    \item  Design, analyze, and implement a modified version of
      \texttt{ScapegoatTree} where the length of the search path does not exceed
      $\log_{\ensuremath{\mathtt{{\color{var}b}}}} \ensuremath{\mathtt{{\color{var}q}}}$, where \ensuremath{\mathtt{{\color{var}b}}} is a parameter with $1<\ensuremath{\mathtt{{\color{var}b}}}<2$.
    \item What does your analysis and/or your experiments say about the
      amortized cost of \ensuremath{\mathtt{find({\color{var}x})}}, \ensuremath{\mathtt{add({\color{var}x})}} and \ensuremath{\mathtt{remove({\color{var}x})}} as a function
      of \ensuremath{\mathtt{{\color{var}n}}} and \ensuremath{\mathtt{{\color{var}b}}}?
  \end{enumerate}
\end{exc}

\begin{exc}\exclabel{scapegoat-quicksize}
  Modify the \ensuremath{\mathtt{add({\color{var}x})}} method of the \texttt{ScapegoatTree} so that it does not
  waste any time recomputing the sizes of subtrees that have already
  been computed.  This is possible because, by the time the method
  wants to compute \ensuremath{\mathtt{size({\color{var}w})}}, it has already computed one of \ensuremath{\mathtt{size({\color{var}w}.{\color{var}left})}}
  or \ensuremath{\mathtt{size({\color{var}w}.{\color{var}right})}}.  Compare the performance of your modified
  implementation with the implementation given here.
\end{exc}

\begin{exc}\exclabel{scapegoat-explicitsize}
  Implement a second version of the \texttt{ScapegoatTree} data structure that
  explicitly stores and maintains the sizes of the subtree rooted at
  each node.  Compare the performance of the resulting implementation
  with that of the original \texttt{ScapegoatTree} implementation as well as
  the implementation from \excref{scapegoat-quicksize}.
\end{exc}

\begin{exc}
  Reimplement the \ensuremath{\mathtt{rebuild({\color{var}u})}} method discussed at the beginning of this
  chapter so that it does not require the use of an array to store the
  nodes of the subtree being rebuilt.  Instead, it should use recursion
  to first connect the nodes into a linked list and then convert this
  linked list into a perfectly balanced binary tree.  (There are
  very elegant recursive implementations of both steps.)
\end{exc}

\begin{exc}
  \index{WeightBalancedTree@\texttt{WeightBalancedTree}}%
  Analyze and implement a \texttt{WeightBalancedTree}. This is a tree in
  which each node \ensuremath{\mathtt{{\color{var}u}}}, except the root, maintains the \emph{balance
  invariant} that $\ensuremath{\mathtt{size({\color{var}u})}} \le (2/3)\ensuremath{\mathtt{size({\color{var}u}.{\color{var}parent})}}$.  The \ensuremath{\mathtt{add({\color{var}x})}} and
  \ensuremath{\mathtt{remove({\color{var}x})}} operations are identical to the standard \texttt{BinarySearchTree}
  operations, except that any time the balance invariant is violated at
  a node \ensuremath{\mathtt{{\color{var}u}}}, the subtree rooted at \ensuremath{\mathtt{{\color{var}u}.{\color{var}parent}}} is rebuilt.
  Your analysis should show that operations on a \texttt{WeightBalancedTree}
  run in $O(\log\ensuremath{\mathtt{{\color{var}n}}})$ amortized time.  
\end{exc}

\begin{exc}
  \index{CountdownTree@\texttt{CountdownTree}}%
  Analyze and implement a \texttt{CountdownTree}.  In a \texttt{CountdownTree} each
  node \ensuremath{\mathtt{{\color{var}u}}} keeps a \emph{timer} \ensuremath{\mathtt{{\color{var}u}.{\color{var}t}}}.  The \ensuremath{\mathtt{add({\color{var}x})}} and \ensuremath{\mathtt{remove({\color{var}x})}}
  operations are exactly the same as in a standard \texttt{BinarySearchTree}
  except that, whenever one of these operations affects \ensuremath{\mathtt{{\color{var}u}}}'s subtree,
  \ensuremath{\mathtt{{\color{var}u}.{\color{var}t}}} is decremented.  When $\ensuremath{\mathtt{{\color{var}u}.{\color{var}t}}}=0$ the entire subtree rooted at \ensuremath{\mathtt{{\color{var}u}}}
  is rebuilt into a perfectly balanced binary search tree.  When a node
  \ensuremath{\mathtt{{\color{var}u}}} is involved in a rebuilding operation (either because \ensuremath{\mathtt{{\color{var}u}}} is rebuilt
  or one of \ensuremath{\mathtt{{\color{var}u}}}'s ancestors is rebuilt) \ensuremath{\mathtt{{\color{var}u}.{\color{var}t}}} is reset to $\ensuremath{\mathtt{size({\color{var}u})}}/3$.

  Your analysis should show that operations on a \texttt{CountdownTree} run
  in $O(\log \ensuremath{\mathtt{{\color{var}n}}})$ amortized time.  (Hint: First show that each node \ensuremath{\mathtt{{\color{var}u}}}
  satisfies some version of a balance invariant.)
\end{exc}

\begin{exc}
  \index{DynamiteTree@\texttt{DynamiteTree}}%
  Analyze and implement a \texttt{DynamiteTree}.  In a \texttt{DynamiteTree} each
  node \ensuremath{\mathtt{{\color{var}u}}} keeps tracks of the size of the subtree rooted at \ensuremath{\mathtt{{\color{var}u}}} in a
  variable \ensuremath{\mathtt{{\color{var}u}.{\color{var}size}}}.  The \ensuremath{\mathtt{add({\color{var}x})}} and \ensuremath{\mathtt{remove({\color{var}x})}} operations are exactly
  the same as in a standard \texttt{BinarySearchTree} except that, whenever one
  of these operations affects a node \ensuremath{\mathtt{{\color{var}u}}}'s subtree, \ensuremath{\mathtt{{\color{var}u}}} \emph{explodes}
  with probability $1/\ensuremath{\mathtt{{\color{var}u}.{\color{var}size}}}$.  When \ensuremath{\mathtt{{\color{var}u}}} explodes, its entire subtree
  is rebuilt into a perfectly balanced binary search tree.

  Your analysis should show that operations on a \texttt{DynamiteTree} run
  in $O(\log \ensuremath{\mathtt{{\color{var}n}}})$ expected time. 
\end{exc}
 

\begin{exc}\exclabel{list-order-maintenance}
  \index{Sequence@\texttt{Sequence}}%
  Design and implement a \texttt{Sequence} data structure that maintains a
  sequence (list) of elements.  It supports these operations:
  \begin{itemize}
    \item \ensuremath{\mathtt{addAfter({\color{var}e})}}: Add a new element after the element \ensuremath{\mathtt{{\color{var}e}}} in the
    sequence.  Return the newly added element.  (If \ensuremath{\mathtt{{\color{var}e}}} is null,
    the new element is added at the beginning of the sequence.)
    \item \ensuremath{\mathtt{remove({\color{var}e})}}: Remove \ensuremath{\mathtt{{\color{var}e}}} from the sequence.
    \item \ensuremath{\mathtt{testBefore({\color{var}e1},{\color{var}e2})}}: return \ensuremath{\mathtt{{\color{var}true}}} if and only if \ensuremath{\mathtt{{\color{var}e1}}} comes
    before \ensuremath{\mathtt{{\color{var}e2}}} in the sequence.
  \end{itemize}
  The first two operations should run in $O(\log \ensuremath{\mathtt{{\color{var}n}}})$ amortized time.
  The third operation should run in constant time.

  The \texttt{Sequence} data structure can be implemented by storing the elements
  in something like a \texttt{ScapegoatTree}, in the same order that they occur
  in the sequence.  To implement \ensuremath{\mathtt{testBefore({\color{var}e1},{\color{var}e2})}} in constant time,
  each element \ensuremath{\mathtt{{\color{var}e}}} is labelled with an integer that encodes the path from
  the root to \ensuremath{\mathtt{{\color{var}e}}}.  In this way, \ensuremath{\mathtt{testBefore({\color{var}e1},{\color{var}e2})}} can be implemented
  by comparing the labels of \ensuremath{\mathtt{{\color{var}e1}}} and \ensuremath{\mathtt{{\color{var}e2}}}.
\end{exc}


\clearpage

\clearpage
\chapter{Red-Black Trees}
\chaplabel{redblack}

\index{binary search tree!red-black}%
\index{red-black tree}%
In this chapter, we present red-black trees, a version of binary search
trees with logarithmic height.  Red-black trees are one of the most
widely used data structures.  They appear as the primary search structure
in many library implementations, including the Java Collections Framework
and several implementations of the C++ Standard Template Library. They
are also used within the Linux operating system kernel.  There are
several reasons for the popularity of red-black trees:
\begin{enumerate}
\item A red-black tree storing \ensuremath{\mathtt{{\color{var}n}}} values has height at most $2\log \ensuremath{\mathtt{{\color{var}n}}}$.
\item The \ensuremath{\mathtt{add({\color{var}x})}} and \ensuremath{\mathtt{remove({\color{var}x})}} operations on a red-black tree run
   in $O(\log \ensuremath{\mathtt{{\color{var}n}}})$ \emph{worst-case} time.
\item The amortized number of rotations performed during an \ensuremath{\mathtt{add({\color{var}x})}}
   or \ensuremath{\mathtt{remove({\color{var}x})}} operation is constant.
\end{enumerate}
The first two of these properties already put red-black trees 
ahead of skiplists, treaps, and scapegoat trees.
Skiplists and treaps rely on randomization and their $O(\log \ensuremath{\mathtt{{\color{var}n}}})$
running times are only expected. Scapegoat trees have a guaranteed
bound on their height, but \ensuremath{\mathtt{add({\color{var}x})}} and \ensuremath{\mathtt{remove({\color{var}x})}} only run in $O(\log
\ensuremath{\mathtt{{\color{var}n}}})$ amortized time.  The third property is just icing on the cake. It
tells us that  that the time needed to add or remove an element \ensuremath{\mathtt{{\color{var}x}}} is
dwarfed by the time it takes to find \ensuremath{\mathtt{{\color{var}x}}}.\footnote{Note that skiplists and
treaps also have this property in the expected sense. See
Exercises~\ref{exc:skiplist-changes} and \ref{exc:treap-rotates}.}

However, the nice properties of red-black trees come with a price:
implementation complexity. Maintaining a bound of $2\log \ensuremath{\mathtt{{\color{var}n}}}$ on the
height is not easy. It requires a careful analysis of a number of cases.
We must ensure that the implementation does exactly the right
thing in each case.  One misplaced rotation or change of colour produces
a bug that can be very difficult to understand and track down.

Rather than jumping directly into the implementation of red-black trees,
we will first provide some background on a related data structure:
2-4 trees.  This will give some insight into how red-black trees were
discovered and why efficiently maintaining them is even possible.

\section{2-4 Trees}
\seclabel{twofour}

A 2-4 tree is a rooted tree with the following properties:
\begin{prp}[height]
  All leaves have the same depth.
\end{prp}
\begin{prp}[degree]
  Every internal node has 2, 3, or 4 children.
\end{prp}
An example of a 2-4 tree is shown in \figref{twofour-example}.
\begin{figure}
  \begin{center}
    \includegraphics[scale=0.90909]{figs/24rb-2}
  \end{center}
  \caption{A 2-4 tree of height 3.}
  \figlabel{twofour-example}
\end{figure}
The properties of 2-4 trees imply that their height is logarithmic in
the number of leaves:
\begin{lem}\lemlabel{twofour-height}
  A 2-4 tree with \ensuremath{\mathtt{{\color{var}n}}} leaves has height at most $\log \ensuremath{\mathtt{{\color{var}n}}}$.
\end{lem}

\begin{proof}
  The lower-bound of 2 on the number of children of an internal node
  implies that, if the height of a 2-4 tree is $h$, then it has at least
  $2^h$ leaves.  In other words,
  \[
     \ensuremath{\mathtt{{\color{var}n}}} \ge 2^h \enspace .
  \]
  Taking logarithms on both sides of this inequality gives $h \le \log \ensuremath{\mathtt{{\color{var}n}}}$.
\end{proof}

\subsection{Adding a Leaf}

Adding a leaf to a 2-4 tree is easy (see \figref{twofour-add}).  If we
want to add a leaf \ensuremath{\mathtt{{\color{var}u}}} as the child of some node \ensuremath{\mathtt{{\color{var}w}}} on the second-last
level, then we simply make \ensuremath{\mathtt{{\color{var}u}}} a child of \ensuremath{\mathtt{{\color{var}w}}}.  This certainly maintains
the height property, but could violate the degree property;  if \ensuremath{\mathtt{{\color{var}w}}}
had four children prior to adding \ensuremath{\mathtt{{\color{var}u}}}, then \ensuremath{\mathtt{{\color{var}w}}} now has five children.
In this case, we \emph{split}
\index{split}%
\ensuremath{\mathtt{{\color{var}w}}} into two nodes, \ensuremath{\mathtt{{\color{var}w}}} and \ensuremath{\mathtt{{\color{var}w}}}', having
two and three children, respectively.
But now \ensuremath{\mathtt{{\color{var}w}}}' has no parent,
so we recursively make \ensuremath{\mathtt{{\color{var}w}}}' a child of \ensuremath{\mathtt{{\color{var}w}}}'s parent.  Again, this may
cause \ensuremath{\mathtt{{\color{var}w}}}'s parent to have too many children in which case we split it.
This process goes on until we reach a node that has fewer than four children,
or until we split the root, \ensuremath{\mathtt{{\color{var}r}}}, into two nodes \ensuremath{\mathtt{{\color{var}r}}} and \ensuremath{\mathtt{{\color{var}r}'}}.  In the
latter case, we make a new root that has \ensuremath{\mathtt{{\color{var}r}}} and \ensuremath{\mathtt{{\color{var}r}'}} as children.
This simultaneously increases the depth of all leaves and so maintains
the height property.

\begin{figure}
  \begin{center}
   \begin{tabular}{c}
     \includegraphics[scale=0.90909]{figs/24tree-add-1} \\
     \includegraphics[scale=0.90909]{figs/24tree-add-2} \\
     \includegraphics[scale=0.90909]{figs/24tree-add-3}
   \end{tabular}
  \end{center}
  \caption[Adding a leaf to a 2-4 Tree]{Adding a leaf to a 2-4 Tree.
  This process stops after one split because \ensuremath{\mathtt{{\color{var}w}.{\color{var}parent}}} has a degree of less
  than 4 before the addition.}
  \figlabel{twofour-add}
\end{figure}

Since the height of the 2-4 tree is never more than $\log \ensuremath{\mathtt{{\color{var}n}}}$, the
process of adding a leaf finishes after at most $\log \ensuremath{\mathtt{{\color{var}n}}}$ steps.

\subsection{Removing a Leaf}

Removing a leaf from a 2-4 tree is a little more tricky (see
\figref{twofour-remove}).  To remove a leaf \ensuremath{\mathtt{{\color{var}u}}} from its parent \ensuremath{\mathtt{{\color{var}w}}}, we
just remove it.  If \ensuremath{\mathtt{{\color{var}w}}} had only two children prior to the removal of \ensuremath{\mathtt{{\color{var}u}}},
then \ensuremath{\mathtt{{\color{var}w}}} is left with only one child and violates the degree property.

\begin{figure}
  \begin{center}
   \begin{tabular}{c}
     \includegraphics[height=\FifthHeightScaleIfNeeded]{figs/24tree-remove-1} \\
     \includegraphics[height=\FifthHeightScaleIfNeeded]{figs/24tree-remove-2} \\
     \includegraphics[height=\FifthHeightScaleIfNeeded]{figs/24tree-remove-3} \\
     \includegraphics[height=\FifthHeightScaleIfNeeded]{figs/24tree-remove-4} \\
     \includegraphics[height=\FifthHeightScaleIfNeeded]{figs/24tree-remove-5} \\
   \end{tabular}
  \end{center}
  \caption[Removing a leaf from a 2-4 Tree]{Removing a leaf from a
    2-4 Tree.  This process goes all the way to the root because each of
    \ensuremath{\mathtt{{\color{var}u}}}'s ancestors and their siblings have only two children.}
  \figlabel{twofour-remove}
\end{figure}

To correct this, we look at \ensuremath{\mathtt{{\color{var}w}}}'s sibling, \ensuremath{\mathtt{{\color{var}w}'}}.  The node \ensuremath{\mathtt{{\color{var}w}'}} is
sure to exist since \ensuremath{\mathtt{{\color{var}w}}}'s parent had at least two children.  If \ensuremath{\mathtt{{\color{var}w}'}}
has three or four children, then we take one of these children from \ensuremath{\mathtt{{\color{var}w}'}}
and give it to \ensuremath{\mathtt{{\color{var}w}}}. Now \ensuremath{\mathtt{{\color{var}w}}} has two children and \ensuremath{\mathtt{{\color{var}w}'}} has two or three
children and we are done.

On the other hand, if \ensuremath{\mathtt{{\color{var}w}'}} has only two children, then we \emph{merge}
\index{merge}%
\ensuremath{\mathtt{{\color{var}w}}} and \ensuremath{\mathtt{{\color{var}w}'}} into a single node, \ensuremath{\mathtt{{\color{var}w}}}, that has three children.  Next we
recursively remove \ensuremath{\mathtt{{\color{var}w}'}} from the parent of \ensuremath{\mathtt{{\color{var}w}'}}.  This process ends
when we reach a node, \ensuremath{\mathtt{{\color{var}u}}}, where \ensuremath{\mathtt{{\color{var}u}}} or its sibling has more than two
children, or when we reach the root.  In the latter case, if the root
is left with only one child, then we delete the root and make its child
the new root.  Again, this simultaneously decreases the height of every
leaf and therefore maintains the height property.

Again, since the height of the tree is never more than $\log \ensuremath{\mathtt{{\color{var}n}}}$,
the process of removing a leaf finishes after at most $\log \ensuremath{\mathtt{{\color{var}n}}}$ steps.

\section{\texttt{RedBlackTree}: A Simulated 2-4 Tree}
\seclabel{redblacktree}

A red-black tree is a binary search tree in which each node, \ensuremath{\mathtt{{\color{var}u}}},
has a \emph{colour}
\index{colour}%
which is either \emph{red} or \emph{black}.  Red is
represented by the value $0$ and black by the value $1$.
\index{red node}%
\index{black node}%


 \ \vspace{.5ex}\begin{Verbatim}[gobble=1,tabsize=2,frame=single,framerule=1.2pt,commandchars=\\@\$,label=\texttt{RedBlackTree},labelposition=topline]
	@\color@keyword$class$ Node<@\color@keyword$T$> @\color@keyword$extends$ BSTNode<Node<@\color@keyword$T$>,@\color@keyword$T$> {
		@\color@keyword$byte$ @\color@var$colour$;
	}
\end{Verbatim}


Before and after any operation on a red-black tree, the following two
properties are satisfied. Each property is defined both in terms of the
colours red and black, and in terms of the numeric values 0 and 1.
\begin{prp}[black-height]
  \index{black-height property}%
  There are the same number of black nodes on every root to leaf
  path. (The sum of the colours on any root to leaf path is the same.)
\end{prp}

\begin{prp}[no-red-edge]
  \index{no-red-edge property}%
  No two red nodes are adjacent.  (For any node \ensuremath{\mathtt{{\color{var}u}}}, except the root,
  $\ensuremath{\mathtt{{\color{var}u}.{\color{var}colour}}} + \ensuremath{\mathtt{{\color{var}u}.{\color{var}parent}.{\color{var}colour}}} \ge 1$.)
\end{prp}
Notice that we can always colour the root, \ensuremath{\mathtt{{\color{var}r}}}, of a red-black tree black
without violating either of these two properties, so we will assume
that the root is black, and the algorithms for updating a red-black
tree will maintain this.  Another trick that simplifies red-black trees
is to treat the external nodes (represented by \ensuremath{\mathtt{{\color{var}nil}}}) as black nodes.
This way, every real node, \ensuremath{\mathtt{{\color{var}u}}}, of a red-black tree has exactly two
children, each with a well-defined colour.  An example of a red-black
tree is shown in \figref{redblack-example}.

\begin{figure}
  \begin{center}
    \includegraphics[scale=0.90909]{figs/24rb-1}
  \end{center}
  \caption[A red-black tree]{An example of a red-black tree with black-height 3.  External (\ensuremath{\mathtt{{\color{var}nil}}}) nodes are drawn as squares.} 
  \figlabel{redblack-example}
\end{figure}


\subsection{Red-Black Trees and 2-4 Trees}

At first it might seem surprising that a red-black tree can be efficiently
updated to maintain the black-height and no-red-edge properties, and
it seems unusual to even consider these as useful properties.  However,
red-black trees were designed to be an efficient simulation of 2-4 trees
as binary trees.

Refer to \figref{twofour-redblack}.
Consider any red-black tree, $T$, having \ensuremath{\mathtt{{\color{var}n}}} nodes and perform the
following transformation: Remove each red node \ensuremath{\mathtt{{\color{var}u}}} and connect \ensuremath{\mathtt{{\color{var}u}}}'s two
children directly to the (black) parent of \ensuremath{\mathtt{{\color{var}u}}}.  After this transformation
we are left with a tree $T'$ having only black nodes.
\begin{figure}
  \begin{center}
    \begin{tabular}{cc}
      \includegraphics[scale=0.90909]{figs/24rb-3} \\
      \includegraphics[scale=0.90909]{figs/24rb-2}
    \end{tabular}
  \end{center}
  \caption{Every red-black tree has a corresponding 2-4 tree.}
  \figlabel{twofour-redblack}
\end{figure}

Every internal node in $T'$ has two, three, or four children: A black
node that started out with two black children will still have two
black children after this transformation.  A black node that started
out with one red and one black child will have three children after this
transformation.  A black node that started out with two red children will
have four children after this transformation.  Furthermore, the black-height
property now guarantees that every root-to-leaf path in $T'$ has the
same length.  In other words, $T'$ is a 2-4 tree!

The 2-4 tree $T'$ has $\ensuremath{\mathtt{{\color{var}n}}}+1$ leaves that correspond to the $\ensuremath{\mathtt{{\color{var}n}}}+1$
external nodes of the red-black tree.  Therefore, this tree has height
at most $\log (\ensuremath{\mathtt{{\color{var}n}}}+1)$. Now, every root to leaf path in the 2-4 tree corresponds
to a path from the root of the red-black tree $T$ to an external node.
The first and last node in this path are black and at most one out of
every two internal nodes is red, so this path has at most $\log(\ensuremath{\mathtt{{\color{var}n}}}+1)$
black nodes and at most $\log(\ensuremath{\mathtt{{\color{var}n}}}+1)-1$ red nodes.  Therefore, the longest path from the root to any \emph{internal} node in $T$ is at most
\[
   2\log(\ensuremath{\mathtt{{\color{var}n}}}+1) -2 \le 2\log \ensuremath{\mathtt{{\color{var}n}}} \enspace ,
\]
for any $\ensuremath{\mathtt{{\color{var}n}}}\ge 1$.  This proves the most important property of
red-black trees:
\begin{lem}
The height of red-black tree with \ensuremath{\mathtt{{\color{var}n}}} nodes is at most $2\log \ensuremath{\mathtt{{\color{var}n}}}$.
\end{lem}

Now that we have seen the relationship between 2-4 trees and
red-black trees, it is not hard to believe that we can efficiently
maintain a red-black tree while adding and removing elements.  

We have already seen that adding an element in a \texttt{BinarySearchTree}
can be done by adding a new leaf.  Therefore, to implement \ensuremath{\mathtt{add({\color{var}x})}} in a
red-black tree we need a method of simulating splitting a node with five
children in a 2-4 tree.  A 2-4 tree node with five children is represented
by a black node that has two red children, one of which also has a red
child. We can ``split'' this node by colouring it red and colouring its
two children black.  An example of this is shown in \figref{rb-split}.

\begin{figure}
  \begin{center}
   \begin{tabular}{c}
     \includegraphics[scale=0.90909]{figs/rb-split-1} \\
     \includegraphics[scale=0.90909]{figs/rb-split-2} \\
     \includegraphics[scale=0.90909]{figs/rb-split-3} \\
   \end{tabular}
  \end{center}
  \caption[Simulating a 2-4 tree]{Simulating a 2-4 tree split operation
    during an addition in a red-black tree.  (This simulates the 2-4
    tree addition shown in \figref{twofour-add}.)}
  \figlabel{rb-split}
\end{figure}

Similarly, implementing \ensuremath{\mathtt{remove({\color{var}x})}} requires a method of merging two nodes
and borrowing a child from a sibling.  Merging two nodes is the inverse of
a split (shown in \figref{rb-split}), and involves colouring two (black)
siblings red and colouring their (red) parent black.  Borrowing from
a sibling is the most complicated of the procedures and involves both
rotations and recolouring nodes.

Of course, during all of this we must still maintain the no-red-edge
property and the black-height property.  While it is no longer surprising
that this can be done, there are a large number of cases that have to
be considered if we try to do a direct simulation of a 2-4 tree by a
red-black tree.  At some point, it just becomes simpler to disregard the
underlying 2-4 tree and work directly towards maintaining the properties
of the red-black tree.

\subsection{Left-Leaning Red-Black Trees}

\index{red-black tree}%
\index{left-leaning red-black tree}%
No single definition of red-black trees exists.  Rather, there is
a family of structures that manage to maintain the black-height
and no-red-edge properties during \ensuremath{\mathtt{add({\color{var}x})}} and \ensuremath{\mathtt{remove({\color{var}x})}}
operations. Different structures do this in different ways.
Here, we implement a data structure that we call a \texttt{RedBlackTree}.
\index{RedBlackTree@\texttt{RedBlackTree}}%
This structure implements a particular variant of red-black trees that
satisfies an additional property:
\begin{prp}[left-leaning]\prplabel{left-leaning}\prplabel{redblack-last}
  \index{left-leaning property}%
  At any node \ensuremath{\mathtt{{\color{var}u}}}, if \ensuremath{\mathtt{{\color{var}u}.{\color{var}left}}} is black, then \ensuremath{\mathtt{{\color{var}u}.{\color{var}right}}} is black.
\end{prp}
Note that the red-black tree shown in \figref{redblack-example}  does
not satisfy the left-leaning property;  it is violated by the parent of
the red node in the rightmost path.

The reason for maintaining the left-leaning property is that it reduces
the number of cases encountered when updating the tree during \ensuremath{\mathtt{add({\color{var}x})}}
and \ensuremath{\mathtt{remove({\color{var}x})}} operations.  In terms of 2-4 trees, it implies that every
2-4 tree has a unique representation:  A node of degree two becomes
a black node with two black children.  A node of degree three becomes
a black node whose left child is red and whose right child is black.
A node of degree four becomes a black node with two red children.

Before we describe the implementation of \ensuremath{\mathtt{add({\color{var}x})}} and \ensuremath{\mathtt{remove({\color{var}x})}} in
detail, we first present some simple subroutines used by these methods
that are illustrated in \figref{redblack-flippullpush}.  The first two
subroutines are for manipulating colours while preserving the black-height
property. The \ensuremath{\mathtt{pushBlack({\color{var}u})}} method takes as input a black node \ensuremath{\mathtt{{\color{var}u}}}
that has two red children and colours \ensuremath{\mathtt{{\color{var}u}}} red and its two children black.
The \ensuremath{\mathtt{pullBlack({\color{var}u})}} method reverses this operation:


 \ \vspace{.5ex}\begin{Verbatim}[gobble=1,tabsize=2,frame=single,framerule=1.2pt,commandchars=\\@\$,label=\texttt{RedBlackTree},labelposition=topline]
	@\color@keyword$void$ pushBlack(Node<@\color@keyword$T$> @\color@var$u$) {
		@\color@var$u$.@\color@var$colour$--;
		@\color@var$u$.@\color@var$left$.@\color@var$colour$++;
		@\color@var$u$.@\color@var$right$.@\color@var$colour$++;
	}
	@\color@keyword$void$ pullBlack(Node<@\color@keyword$T$> @\color@var$u$) {
		@\color@var$u$.@\color@var$colour$++;
		@\color@var$u$.@\color@var$left$.@\color@var$colour$--;
		@\color@var$u$.@\color@var$right$.@\color@var$colour$--;
	}
\end{Verbatim}


\begin{figure}
  \begin{center}
    \includegraphics[width=\ScaleIfNeeded]{figs/flippullpush}
  \end{center}
  \caption{Flips, pulls and pushes}
  \figlabel{redblack-flippullpush}
\end{figure}

The \ensuremath{\mathtt{flipLeft({\color{var}u})}} method swaps the colours of \ensuremath{\mathtt{{\color{var}u}}} and \ensuremath{\mathtt{{\color{var}u}.{\color{var}right}}}
and then performs a left rotation at \ensuremath{\mathtt{{\color{var}u}}}.  This method reverses the
colours of these two nodes as well as their parent-child relationship:


 \ \vspace{.5ex}\begin{Verbatim}[gobble=1,tabsize=2,frame=single,framerule=1.2pt,commandchars=\\@\$,label=\texttt{RedBlackTree},labelposition=topline]
	@\color@keyword$void$ flipLeft(Node<@\color@keyword$T$> @\color@var$u$) {
		swapColors(@\color@var$u$, @\color@var$u$.@\color@var$right$);
		rotateLeft(@\color@var$u$);
	}
\end{Verbatim}

The \ensuremath{\mathtt{flipLeft({\color{var}u})}} operation
is especially useful in restoring the left-leaning property at a node
\ensuremath{\mathtt{{\color{var}u}}} that violates it (because \ensuremath{\mathtt{{\color{var}u}.{\color{var}left}}} is black and \ensuremath{\mathtt{{\color{var}u}.{\color{var}right}}} is red).
In this special case, we can be assured that this operation preserves both
the black-height and no-red-edge properties.  The \ensuremath{\mathtt{flipRight({\color{var}u})}} operation
is symmetric with \ensuremath{\mathtt{flipLeft({\color{var}u})}}, when the roles of left and right are reversed.


 \ \vspace{.5ex}\begin{Verbatim}[gobble=1,tabsize=2,frame=single,framerule=1.2pt,commandchars=\\@\$,label=\texttt{RedBlackTree},labelposition=topline]
	@\color@keyword$void$ flipRight(Node<@\color@keyword$T$> @\color@var$u$) {
		swapColors(@\color@var$u$, @\color@var$u$.@\color@var$left$);
		rotateRight(@\color@var$u$);
	}
\end{Verbatim}


\subsection{Addition}

To implement \ensuremath{\mathtt{add({\color{var}x})}} in a \texttt{RedBlackTree}, we perform a standard
\texttt{BinarySearchTree} insertion to add a new leaf, \ensuremath{\mathtt{{\color{var}u}}}, with $\ensuremath{\mathtt{{\color{var}u}.{\color{var}x}}}=\ensuremath{\mathtt{{\color{var}x}}}$ and
set $\ensuremath{\mathtt{{\color{var}u}.{\color{var}colour}}}=\ensuremath{\mathtt{{\color{var}red}}}$.  Note that this does not change the black height
of any node, so it does not violate the black-height property.  It may,
however, violate the left-leaning property (if \ensuremath{\mathtt{{\color{var}u}}} is the right child of
its parent), and it may violate the no-red-edge property (if \ensuremath{\mathtt{{\color{var}u}}}'s parent
is \ensuremath{\mathtt{{\color{var}red}}}).  To restore these properties, we call the method \ensuremath{\mathtt{addFixup({\color{var}u})}}.


 \ \vspace{.5ex}\begin{Verbatim}[gobble=1,tabsize=2,frame=single,framerule=1.2pt,commandchars=\\@\$,label=\texttt{RedBlackTree},labelposition=topline]
	@\color@var$boolean$ add(@\color@keyword$T$ @\color@var$x$) {
		Node<@\color@keyword$T$> @\color@var$u$ = newNode(@\color@var$x$);
		@\color@var$u$.@\color@var$colour$ = @\color@var$red$;
		@\color@var$boolean$ @\color@var$added$ = add(@\color@var$u$);
		@\color@keyword$if$ (@\color@var$added$)
			addFixup(@\color@var$u$);
		@\color@var$return$ @\color@var$added$;
	}
\end{Verbatim}


Illustrated in \figref{rb-addfix}, the \ensuremath{\mathtt{addFixup({\color{var}u})}} method takes
as input a node \ensuremath{\mathtt{{\color{var}u}}} whose colour is red and which may violate the
no-red-edge property and/or the left-leaning property.  The following
discussion is probably impossible to follow without referring to
\figref{rb-addfix} or recreating it on a piece of paper.  Indeed, the
reader may wish to study this figure before continuing.

\begin{figure}
  \begin{center}
    \includegraphics[width=\ScaleIfNeeded]{figs/rb-addfix}
  \end{center}
  \caption{A single round in the process of fixing Property~2 after
  an insertion.}
  \figlabel{rb-addfix}
\end{figure}

If \ensuremath{\mathtt{{\color{var}u}}} is the root of the tree, then we can colour \ensuremath{\mathtt{{\color{var}u}}} black to restore
both properties.  If \ensuremath{\mathtt{{\color{var}u}}}'s sibling is also red, then \ensuremath{\mathtt{{\color{var}u}}}'s parent must be
black, so both the left-leaning and no-red-edge properties already hold.

Otherwise, we first determine if \ensuremath{\mathtt{{\color{var}u}}}'s parent, \ensuremath{\mathtt{{\color{var}w}}}, violates the
left-leaning property and, if so, perform a \ensuremath{\mathtt{flipLeft({\color{var}w})}} operation and
set $\ensuremath{\mathtt{{\color{var}u}}}=\ensuremath{\mathtt{{\color{var}w}}}$.  This leaves us in a well-defined state:  \ensuremath{\mathtt{{\color{var}u}}} is the left
child of its parent, \ensuremath{\mathtt{{\color{var}w}}}, so \ensuremath{\mathtt{{\color{var}w}}} now satisfies the left-leaning property.
All that remains is to ensure the no-red-edge property at \ensuremath{\mathtt{{\color{var}u}}}.  We only
have to worry about the case in which \ensuremath{\mathtt{{\color{var}w}}} is red, since otherwise \ensuremath{\mathtt{{\color{var}u}}}
already satisfies the no-red-edge property.

Since we are not done yet, \ensuremath{\mathtt{{\color{var}u}}} is red and \ensuremath{\mathtt{{\color{var}w}}} is red.  The no-red-edge
property (which is only violated by \ensuremath{\mathtt{{\color{var}u}}} and not by \ensuremath{\mathtt{{\color{var}w}}}) implies that
\ensuremath{\mathtt{{\color{var}u}}}'s grandparent \ensuremath{\mathtt{{\color{var}g}}} exists and is black.  If \ensuremath{\mathtt{{\color{var}g}}}'s right child is red,
then the left-leaning property ensures that both \ensuremath{\mathtt{{\color{var}g}}}'s children are red,
and a call to \ensuremath{\mathtt{pushBlack({\color{var}g})}} makes \ensuremath{\mathtt{{\color{var}g}}} red and \ensuremath{\mathtt{{\color{var}w}}} black.  This restores
the no-red-edge property at \ensuremath{\mathtt{{\color{var}u}}}, but may cause it to be violated at \ensuremath{\mathtt{{\color{var}g}}},
so the whole process starts over with $\ensuremath{\mathtt{{\color{var}u}}}=\ensuremath{\mathtt{{\color{var}g}}}$.

If \ensuremath{\mathtt{{\color{var}g}}}'s right child is black, then a call to \ensuremath{\mathtt{flipRight({\color{var}g})}} makes
\ensuremath{\mathtt{{\color{var}w}}} the (black) parent of \ensuremath{\mathtt{{\color{var}g}}} and gives \ensuremath{\mathtt{{\color{var}w}}} two red children, \ensuremath{\mathtt{{\color{var}u}}} and
\ensuremath{\mathtt{{\color{var}g}}}. This ensures that \ensuremath{\mathtt{{\color{var}u}}} satisfies the no-red-edge property and \ensuremath{\mathtt{{\color{var}g}}}
satisfies the left-leaning property.  In this case we can stop.


 \ \vspace{.5ex}\begin{Verbatim}[gobble=1,tabsize=2,frame=single,framerule=1.2pt,commandchars=\\@\$,label=\texttt{RedBlackTree},labelposition=topline]
	@\color@keyword$void$ addFixup(Node<@\color@keyword$T$> @\color@var$u$) {
		@\color@keyword$while$ (@\color@var$u$.@\color@var$colour$ == @\color@var$red$) {
			@\color@keyword$if$ (@\color@var$u$ == @\color@var$r$) { @\color@comment$// u is the root - done$
				@\color@var$u$.@\color@var$colour$ = @\color@var$black$;
				@\color@var$return$;
			}
			Node<@\color@keyword$T$> @\color@var$w$ = @\color@var$u$.@\color@var$parent$;
			@\color@keyword$if$ (@\color@var$w$.@\color@var$left$.@\color@var$colour$ == @\color@var$black$) { @\color@comment$// ensure left-leaning$
				flipLeft(@\color@var$w$);
				@\color@var$u$ = @\color@var$w$;
				@\color@var$w$ = @\color@var$u$.@\color@var$parent$;
			}
			@\color@keyword$if$ (@\color@var$w$.@\color@var$colour$ == @\color@var$black$)
				@\color@var$return$; @\color@comment$// no red-red edge = done$
			Node<@\color@keyword$T$> @\color@var$g$ = @\color@var$w$.@\color@var$parent$; @\color@comment$// grandparent of u$
			@\color@keyword$if$ (@\color@var$g$.@\color@var$right$.@\color@var$colour$ == @\color@var$black$) {
				flipRight(@\color@var$g$);
				@\color@var$return$;
			} @\color@keyword$else$ {
				pushBlack(@\color@var$g$);
				@\color@var$u$ = @\color@var$g$;
			}
		}
	}
\end{Verbatim}


The \ensuremath{\mathtt{insertFixup({\color{var}u})}} method takes constant time per iteration and each
iteration either finishes or moves \ensuremath{\mathtt{{\color{var}u}}} closer to the root.  Therefore,
the \ensuremath{\mathtt{insertFixup({\color{var}u})}} method finishes after $O(\log \ensuremath{\mathtt{{\color{var}n}}})$ iterations in
$O(\log \ensuremath{\mathtt{{\color{var}n}}})$ time.

\subsection{Removal}

The \ensuremath{\mathtt{remove({\color{var}x})}} operation in a \texttt{RedBlackTree} is the most complicated
to implement, and this is true of all known red-black tree variants.
Just like the \ensuremath{\mathtt{remove({\color{var}x})}} operation in a \texttt{BinarySearchTree},
this operation boils down to finding a node \ensuremath{\mathtt{{\color{var}w}}} with only one child,
\ensuremath{\mathtt{{\color{var}u}}}, and splicing \ensuremath{\mathtt{{\color{var}w}}} out of the tree by having \ensuremath{\mathtt{{\color{var}w}.{\color{var}parent}}} adopt \ensuremath{\mathtt{{\color{var}u}}}.

The problem with this is that, if \ensuremath{\mathtt{{\color{var}w}}} is black, then the black-height
property will now be violated at \ensuremath{\mathtt{{\color{var}w}.{\color{var}parent}}}.  We may avoid this problem,
temporarily, by adding \ensuremath{\mathtt{{\color{var}w}.{\color{var}colour}}} to \ensuremath{\mathtt{{\color{var}u}.{\color{var}colour}}}.  Of course, this introduces
two other problems:  (1)~if \ensuremath{\mathtt{{\color{var}u}}} and \ensuremath{\mathtt{{\color{var}w}}} both started out black, then
$\ensuremath{\mathtt{{\color{var}u}.{\color{var}colour}}}+\ensuremath{\mathtt{{\color{var}w}.{\color{var}colour}}}=2$ (double black), which is an invalid colour.
If \ensuremath{\mathtt{{\color{var}w}}} was red, then it is replaced by a black node \ensuremath{\mathtt{{\color{var}u}}}, which may
violate the left-leaning property at $\ensuremath{\mathtt{{\color{var}u}.{\color{var}parent}}}$.  Both of these
problems can be resolved with a call to the \ensuremath{\mathtt{removeFixup({\color{var}u})}} method.


 \ \vspace{.5ex}\begin{Verbatim}[gobble=1,tabsize=2,frame=single,framerule=1.2pt,commandchars=\\@\$,label=\texttt{RedBlackTree},labelposition=topline]
	@\color@var$boolean$ remove(@\color@keyword$T$ @\color@var$x$) {
		Node<@\color@keyword$T$> @\color@var$u$ = findLast(@\color@var$x$);
		@\color@keyword$if$ (@\color@var$u$ == @\color@var$nil$ || compare(@\color@var$u$.@\color@var$x$, @\color@var$x$) != 0)
			@\color@var$return$ @\color@var$false$;
		Node<@\color@keyword$T$> @\color@var$w$ = @\color@var$u$.@\color@var$right$;
		@\color@keyword$if$ (@\color@var$w$ == @\color@var$nil$) {
			@\color@var$w$ = @\color@var$u$;
			@\color@var$u$ = @\color@var$w$.@\color@var$left$;
		} @\color@keyword$else$ {
			@\color@keyword$while$ (@\color@var$w$.@\color@var$left$ != @\color@var$nil$)
				@\color@var$w$ = @\color@var$w$.@\color@var$left$;
			@\color@var$u$.@\color@var$x$ = @\color@var$w$.@\color@var$x$;
			@\color@var$u$ = @\color@var$w$.@\color@var$right$;
		}
		splice(@\color@var$w$);
		@\color@var$u$.@\color@var$colour$ += @\color@var$w$.@\color@var$colour$;
		@\color@var$u$.@\color@var$parent$ = @\color@var$w$.@\color@var$parent$;
		removeFixup(@\color@var$u$);
		@\color@var$return$ @\color@var$true$;
	}
\end{Verbatim}


The \ensuremath{\mathtt{removeFixup({\color{var}u})}} method takes as its input a node \ensuremath{\mathtt{{\color{var}u}}} whose colour is black
(1) or double-black (2).  If \ensuremath{\mathtt{{\color{var}u}}} is double-black, then \ensuremath{\mathtt{removeFixup({\color{var}u})}}
performs a series of rotations and recolouring operations that move the
double-black node up the tree until it can be eliminated.  During this
process, the node \ensuremath{\mathtt{{\color{var}u}}} changes until, at the end of this process, \ensuremath{\mathtt{{\color{var}u}}}
refers to the root of the subtree that has been changed.  The root of
this subtree may have changed colour.  In particular, it may have gone
from red to black, so the \ensuremath{\mathtt{removeFixup({\color{var}u})}} method finishes by checking
if \ensuremath{\mathtt{{\color{var}u}}}'s parent violates the left-leaning property and, if so, fixing it.


 \ \vspace{.5ex}\begin{Verbatim}[gobble=1,tabsize=2,frame=single,framerule=1.2pt,commandchars=\\@\$,label=\texttt{RedBlackTree},labelposition=topline]
	@\color@keyword$void$ removeFixup(Node<@\color@keyword$T$> @\color@var$u$) {
		@\color@keyword$while$ (@\color@var$u$.@\color@var$colour$ > @\color@var$black$) {
			@\color@keyword$if$ (@\color@var$u$ == @\color@var$r$) {
				@\color@var$u$.@\color@var$colour$ = @\color@var$black$;
			} @\color@keyword$else$ @\color@keyword$if$ (@\color@var$u$.@\color@var$parent$.@\color@var$left$.@\color@var$colour$ == @\color@var$red$) {
				@\color@var$u$ = removeFixupCase1(@\color@var$u$);
			} @\color@keyword$else$ @\color@keyword$if$ (@\color@var$u$ == @\color@var$u$.@\color@var$parent$.@\color@var$left$) {
				@\color@var$u$ = removeFixupCase2(@\color@var$u$);
			} @\color@keyword$else$ { 
				@\color@var$u$ = removeFixupCase3(@\color@var$u$);
			}
		}
		@\color@keyword$if$ (@\color@var$u$ != @\color@var$r$) { @\color@comment$// restore left-leaning property if needed$
			Node<@\color@keyword$T$> @\color@var$w$ = @\color@var$u$.@\color@var$parent$;
			@\color@keyword$if$ (@\color@var$w$.@\color@var$right$.@\color@var$colour$ == @\color@var$red$ && @\color@var$w$.@\color@var$left$.@\color@var$colour$ == @\color@var$black$) {
				flipLeft(@\color@var$w$);
			}
		}
	}
\end{Verbatim}


The \ensuremath{\mathtt{removeFixup({\color{var}u})}} method is illustrated in \figref{rb-removefix}.
Again, the following text will be difficult, if not impossible, to follow
without referring to \figref{rb-removefix}.  Each iteration of the loop
in \ensuremath{\mathtt{removeFixup({\color{var}u})}} processes the double-black node \ensuremath{\mathtt{{\color{var}u}}}, based on one
of four cases:

\begin{figure}
  \begin{center}
    \includegraphics[height=\HeightScaleIfNeeded]{figs/rb-removefix}
  \end{center}
  \caption{A single round in the process of eliminating a double-black node
   after a removal.}
  \figlabel{rb-removefix}
\end{figure}

\noindent
Case 0: \ensuremath{\mathtt{{\color{var}u}}} is the root.  This is the easiest case to treat.  We recolour
\ensuremath{\mathtt{{\color{var}u}}} to be black (this does not violate any of the red-black tree
properties).

\noindent 
Case 1: \ensuremath{\mathtt{{\color{var}u}}}'s sibling, \ensuremath{\mathtt{{\color{var}v}}}, is red.  In this case, \ensuremath{\mathtt{{\color{var}u}}}'s sibling is the
left child of its parent, \ensuremath{\mathtt{{\color{var}w}}} (by the left-leaning property).  We perform
a right-flip at \ensuremath{\mathtt{{\color{var}w}}} and then proceed to the next iteration.  Note that
this action causes \ensuremath{\mathtt{{\color{var}w}}}'s parent to violate the left-leaning property and
the depth of \ensuremath{\mathtt{{\color{var}u}}} to increase.  However, it also implies that the next
iteration will be in Case~3 with \ensuremath{\mathtt{{\color{var}w}}} coloured red.  When examining Case~3
below, we will see that the process will stop during the next iteration.


 \ \vspace{.5ex}\begin{Verbatim}[gobble=1,tabsize=2,frame=single,framerule=1.2pt,commandchars=\\@\$,label=\texttt{RedBlackTree},labelposition=topline]
	Node<@\color@keyword$T$> removeFixupCase1(Node<@\color@keyword$T$> @\color@var$u$) {
		flipRight(@\color@var$u$.@\color@var$parent$);
		@\color@var$return$ @\color@var$u$;
	}
\end{Verbatim}


\noindent
Case 2: \ensuremath{\mathtt{{\color{var}u}}}'s sibling, \ensuremath{\mathtt{{\color{var}v}}}, is black, and \ensuremath{\mathtt{{\color{var}u}}} is the left child of its
parent, \ensuremath{\mathtt{{\color{var}w}}}.  In this case, we call \ensuremath{\mathtt{pullBlack({\color{var}w})}}, making \ensuremath{\mathtt{{\color{var}u}}} black,
\ensuremath{\mathtt{{\color{var}v}}} red, and darkening the colour of \ensuremath{\mathtt{{\color{var}w}}} to black or double-black.
At this point, \ensuremath{\mathtt{{\color{var}w}}} does not satisfy the left-leaning property, so we
call \ensuremath{\mathtt{flipLeft({\color{var}w})}} to fix this.

At this point, \ensuremath{\mathtt{{\color{var}w}}} is red and \ensuremath{\mathtt{{\color{var}v}}} is the root of the subtree with which
we started.  We need to check if \ensuremath{\mathtt{{\color{var}w}}} causes the no-red-edge property to
be violated.  We do this by inspecting \ensuremath{\mathtt{{\color{var}w}}}'s right child, \ensuremath{\mathtt{{\color{var}q}}}.  If \ensuremath{\mathtt{{\color{var}q}}}
is black, then \ensuremath{\mathtt{{\color{var}w}}} satisfies the no-red-edge property and we can continue
the next iteration with $\ensuremath{\mathtt{{\color{var}u}}}=\ensuremath{\mathtt{{\color{var}v}}}$.

Otherwise (\ensuremath{\mathtt{{\color{var}q}}} is red), so both the no-red-edge property and the left-leaning
properties are violated at \ensuremath{\mathtt{{\color{var}q}}} and \ensuremath{\mathtt{{\color{var}w}}}, respectively.  The left-leaning
property is restored with a call to  
\ensuremath{\mathtt{rotateLeft({\color{var}w})}}, but the no-red-edge
property is still violated.  At this point, \ensuremath{\mathtt{{\color{var}q}}} is the left child of
\ensuremath{\mathtt{{\color{var}v}}}, \ensuremath{\mathtt{{\color{var}w}}} is the left child of \ensuremath{\mathtt{{\color{var}q}}}, \ensuremath{\mathtt{{\color{var}q}}} and \ensuremath{\mathtt{{\color{var}w}}} are both red, and \ensuremath{\mathtt{{\color{var}v}}}
is black or double-black.  A \ensuremath{\mathtt{flipRight({\color{var}v})}}  makes \ensuremath{\mathtt{{\color{var}q}}} the parent of
both \ensuremath{\mathtt{{\color{var}v}}} and \ensuremath{\mathtt{{\color{var}w}}}.  Following this up by a \ensuremath{\mathtt{pushBlack({\color{var}q})}} makes both \ensuremath{\mathtt{{\color{var}v}}}
and \ensuremath{\mathtt{{\color{var}w}}} black and sets the colour of \ensuremath{\mathtt{{\color{var}q}}} back to the original colour of \ensuremath{\mathtt{{\color{var}w}}}.

At this point, the double-black node is has been eliminated and the
no-red-edge and black-height properties are reestablished.  Only one possible problem remains: the right child of \ensuremath{\mathtt{{\color{var}v}}} may be red, in which
case the left-leaning property would be violated.  We check this and
perform a \ensuremath{\mathtt{flipLeft({\color{var}v})}} to correct it if necessary.


 \ \vspace{.5ex}\begin{Verbatim}[gobble=1,tabsize=2,frame=single,framerule=1.2pt,commandchars=\\@\$,label=\texttt{RedBlackTree},labelposition=topline]
	Node<@\color@keyword$T$> removeFixupCase2(Node<@\color@keyword$T$> @\color@var$u$) {
		Node<@\color@keyword$T$> @\color@var$w$ = @\color@var$u$.@\color@var$parent$;
		Node<@\color@keyword$T$> @\color@var$v$ = @\color@var$w$.@\color@var$right$;
		pullBlack(@\color@var$w$); @\color@comment$// w.left$
		flipLeft(@\color@var$w$); @\color@comment$// w is now red$
		Node<@\color@keyword$T$> @\color@var$q$ = @\color@var$w$.@\color@var$right$;
		@\color@keyword$if$ (@\color@var$q$.@\color@var$colour$ == @\color@var$red$) { @\color@comment$// q-w is red-red$
			rotateLeft(@\color@var$w$);
			flipRight(@\color@var$v$);
			pushBlack(@\color@var$q$);
			@\color@keyword$if$ (@\color@var$v$.@\color@var$right$.@\color@var$colour$ == @\color@var$red$)
				flipLeft(@\color@var$v$);
			@\color@var$return$ @\color@var$q$;
		} @\color@keyword$else$ {
			@\color@var$return$ @\color@var$v$;
		}
	}
\end{Verbatim}


\noindent
Case 3: \ensuremath{\mathtt{{\color{var}u}}}'s sibling is black and \ensuremath{\mathtt{{\color{var}u}}} is the right child of its parent,
\ensuremath{\mathtt{{\color{var}w}}}.  This case is symmetric to Case~2 and is handled mostly the same way.
The only differences come from the fact that the left-leaning property
is asymmetric, so it requires different handling.

As before, we begin with a call to \ensuremath{\mathtt{pullBlack({\color{var}w})}}, which makes \ensuremath{\mathtt{{\color{var}v}}} red
and \ensuremath{\mathtt{{\color{var}u}}} black.  A call to \ensuremath{\mathtt{flipRight({\color{var}w})}} promotes \ensuremath{\mathtt{{\color{var}v}}} to the root of
the subtree.  At this point \ensuremath{\mathtt{{\color{var}w}}} is red, and the code branches two ways
depending on the colour of \ensuremath{\mathtt{{\color{var}w}}}'s left child, \ensuremath{\mathtt{{\color{var}q}}}.

If \ensuremath{\mathtt{{\color{var}q}}} is red, then the code finishes up exactly the same way as Case~2
does, but is even simpler since there is no danger of \ensuremath{\mathtt{{\color{var}v}}} not
satisfying the left-leaning property.

The more complicated case occurs when \ensuremath{\mathtt{{\color{var}q}}} is black.  In this case,
we examine the colour of \ensuremath{\mathtt{{\color{var}v}}}'s left child.  If it is red, then \ensuremath{\mathtt{{\color{var}v}}} has
two red children and its extra black can be pushed down with a call to
\ensuremath{\mathtt{pushBlack({\color{var}v})}}.  At this point, \ensuremath{\mathtt{{\color{var}v}}} now has \ensuremath{\mathtt{{\color{var}w}}}'s original colour, and we
are done.

If \ensuremath{\mathtt{{\color{var}v}}}'s left child is black, then \ensuremath{\mathtt{{\color{var}v}}} violates the left-leaning property,
and we restore this with a call to \ensuremath{\mathtt{flipLeft({\color{var}v})}}.  We then return the
node \ensuremath{\mathtt{{\color{var}v}}} so that the next iteration of \ensuremath{\mathtt{removeFixup({\color{var}u})}} then continues
with $\ensuremath{\mathtt{{\color{var}u}}}=\ensuremath{\mathtt{{\color{var}v}}}$.


 \ \vspace{.5ex}\begin{Verbatim}[gobble=1,tabsize=2,frame=single,framerule=1.2pt,commandchars=\\@\$,label=\texttt{RedBlackTree},labelposition=topline]
	Node<@\color@keyword$T$> removeFixupCase3(Node<@\color@keyword$T$> @\color@var$u$) {
		Node<@\color@keyword$T$> @\color@var$w$ = @\color@var$u$.@\color@var$parent$;
		Node<@\color@keyword$T$> @\color@var$v$ = @\color@var$w$.@\color@var$left$;
		pullBlack(@\color@var$w$);       
		flipRight(@\color@var$w$); @\color@comment$// w is now red$
		Node<@\color@keyword$T$> @\color@var$q$ = @\color@var$w$.@\color@var$left$;
		@\color@keyword$if$ (@\color@var$q$.@\color@var$colour$ == @\color@var$red$) { @\color@comment$// q-w is red-red$
			rotateRight(@\color@var$w$);
			flipLeft(@\color@var$v$);
			pushBlack(@\color@var$q$);
			@\color@var$return$ @\color@var$q$;
		} @\color@keyword$else$ {
			@\color@keyword$if$ (@\color@var$v$.@\color@var$left$.@\color@var$colour$ == @\color@var$red$) {
				pushBlack(@\color@var$v$); @\color@comment$// both v's children are red$
				@\color@var$return$ @\color@var$v$;
			} @\color@keyword$else$ { @\color@comment$// ensure left-leaning$
				flipLeft(@\color@var$v$);
				@\color@var$return$ @\color@var$w$;
			}
		}				
	}
\end{Verbatim}


Each iteration of \ensuremath{\mathtt{removeFixup({\color{var}u})}} takes constant time.  Cases~2 and 3
either finish or move \ensuremath{\mathtt{{\color{var}u}}} closer to the root of the tree.  Case~0 (where
\ensuremath{\mathtt{{\color{var}u}}} is the root) always terminates and Case~1 leads immediately to Case~3,
which also terminates.  Since the height of the tree is at most $2\log
\ensuremath{\mathtt{{\color{var}n}}}$, we conclude that there are at most $O(\log \ensuremath{\mathtt{{\color{var}n}}})$ iterations of
\ensuremath{\mathtt{removeFixup({\color{var}u})}}, so \ensuremath{\mathtt{removeFixup({\color{var}u})}} runs in $O(\log \ensuremath{\mathtt{{\color{var}n}}})$ time.


\section{Summary}
\seclabel{redblack-summary}

The following theorem summarizes the performance of the \texttt{RedBlackTree} data structure:

\begin{thm}
  A \texttt{RedBlackTree} implements the \texttt{SSet} interface and
  supports the operations \ensuremath{\mathtt{add({\color{var}x})}}, \ensuremath{\mathtt{remove({\color{var}x})}}, and \ensuremath{\mathtt{find({\color{var}x})}} in $O(\log
  \ensuremath{\mathtt{{\color{var}n}}})$ worst-case time per operation.
\end{thm}

Not included in the above theorem is the following extra bonus:

\begin{thm}\thmlabel{redblack-amortized}
  Beginning with an empty \texttt{RedBlackTree}, any sequence of $m$
  \ensuremath{\mathtt{add({\color{var}x})}} and \ensuremath{\mathtt{remove({\color{var}x})}} operations results in a total of $O(m)$
  time spent during all calls \ensuremath{\mathtt{addFixup({\color{var}u})}} and \ensuremath{\mathtt{removeFixup({\color{var}u})}}. 
\end{thm}

We only sketch a proof of \thmref{redblack-amortized}. By comparing
\ensuremath{\mathtt{addFixup({\color{var}u})}} and \ensuremath{\mathtt{removeFixup({\color{var}u})}} with the algorithms for adding or
removing a leaf in a 2-4 tree, we can convince ourselves that this
property is inherited from a 2-4 tree.  In particular, if we can show
that the total time spent splitting, merging, and borrowing in a 2-4
tree is $O(m)$, then this implies \thmref{redblack-amortized}.

The proof of this theorem for 2-4 trees uses the potential
method
\index{potential method}%
of amortized analysis.\footnote{See the proofs of
\lemref{dualarraydeque-amortized} and \lemref{selist-amortized} for
other applications of the potential method.} Define the potential of an
internal node \ensuremath{\mathtt{{\color{var}u}}} in a 2-4 tree as
\[
  \Phi(\ensuremath{\mathtt{{\color{var}u}}}) = 
    \begin{cases} 
      1 & \text{if \ensuremath{\mathtt{{\color{var}u}}} has 2 children} \\ 
      0 & \text{if \ensuremath{\mathtt{{\color{var}u}}} has 3 children} \\ 
      3 & \text{if \ensuremath{\mathtt{{\color{var}u}}} has 4 children}  
    \end{cases}
\]
and the potential of a 2-4 tree as the sum of the potentials of its nodes.
When a split occurs, it is because a node with four children becomes
two nodes, with two and three children.  This means that the overall
potential drops by $3-1-0 = 2$. When a merge occurs, two nodes that used
to have two children are replaced by one node with three children. The
result is a drop in potential of $2-0=2$.  Therefore, for every split
or merge, the potential decreases by two.

Next notice that, if we ignore splitting and merging of nodes, there are
only a constant number of nodes whose number of children is changed by
the addition or removal of a leaf.  When adding a node, one node has
its number of children increase by one, increasing the potential by
at most three.  During the removal of a leaf, one node has its number
of children decrease by one, increasing the potential by at most one,
and two nodes may be involved in a borrowing operation, increasing their
total potential by at most one.

To summarize, each merge and split causes the potential to drop by
at least two.  Ignoring merging and splitting, each addition or removal
causes the potential to rise by at most three, and the potential is always
non-negative.  Therefore, the number of splits and merges caused by $m$
additions or removals on an initially empty tree is at most $3m/2$.
\thmref{redblack-amortized} is a consequence of this analysis and the
correspondence between 2-4 trees and red-black trees.

\section{Discussion and Exercises}

Red-black trees were first introduced by Guibas and Sedgewick \cite{gs78}.
Despite their high implementation complexity they are found in some of
the most commonly used libraries and applications.  Most algorithms and
data structures textbooks discuss some variant of red-black trees.

Andersson \cite{a93} describes a left-leaning version of balanced trees
that is similar to red-black trees but has the additional constraint
that any node has at most one red child.  This implies that these trees
simulate 2-3 trees rather than 2-4 trees.  They are significantly simpler,
though, than the \texttt{RedBlackTree} structure presented in this chapter.

Sedgewick \cite{s08} describes two versions of left-leaning red-black
trees.  These use recursion along with a simulation of top-down splitting
and merging in 2-4 trees. The combination of these two techniques makes
for particularly short and elegant code.

A related, and older, data structure is the \emph{AVL tree} \cite{avl62}.
\index{AVL tree}%
AVL trees are \emph{height-balanced}:
\index{height-balanced}%
\index{binary search tree!height balanced}%
At each node $u$, the height
of the subtree rooted at \ensuremath{\mathtt{{\color{var}u}.{\color{var}left}}} and the subtree rooted at \ensuremath{\mathtt{{\color{var}u}.{\color{var}right}}}
differ by at most one.  It follows immediately that, if $F(h)$ is the
minimum number of leaves in a tree of height $h$, then $F(h)$ obeys the
Fibonacci recurrence
\[
   F(h) = F(h-1) + F(h-2)
\]
with base cases $F(0)=1$ and $F(1)=1$.  This means $F(h)$ is approximately
$\varphi^h/\sqrt{5}$, where $\varphi=(1+\sqrt{5})/2\approx1.61803399$ is the
\emph{golden ratio}.  (More precisely, $|\varphi^h/\sqrt{5} - F(h)|\le 1/2$.)
Arguing as in the proof of \lemref{twofour-height}, this implies
\[
   h \le \log_\varphi \ensuremath{\mathtt{{\color{var}n}}} \approx 1.440420088\log \ensuremath{\mathtt{{\color{var}n}}} \enspace ,
\]
so AVL trees have smaller height than red-black trees.  The height
balancing can be maintained during \ensuremath{\mathtt{add({\color{var}x})}} and \ensuremath{\mathtt{remove({\color{var}x})}} operations
by walking back up the path to the root and performing a rebalancing
operation at each node \ensuremath{\mathtt{{\color{var}u}}} where the height of \ensuremath{\mathtt{{\color{var}u}}}'s left and right
subtrees differ by two.  See \figref{avl-rebalance}.

\begin{figure}
  \begin{center}
    \includegraphics[scale=0.90909]{figs/avl-rebalance}
  \end{center}
  \caption{Rebalancing in an AVL tree.  At most two rotations are required
  to convert a node whose subtrees have a height of $h$ and $h+2$ into a node
  whose subtrees each have a height of at most $h+1$.}
  \figlabel{avl-rebalance}
\end{figure}

Andersson's variant of red-black trees, Sedgewick's variant of red-black
trees, and AVL trees are all simpler to implement than the \texttt{RedBlackTree}
structure defined here.  Unfortunately, none of them can guarantee that
the amortized time spent rebalancing is $O(1)$ per update.  In particular,
there is no analogue of \thmref{redblack-amortized} for those structures.

\begin{figure}
  \centering{\includegraphics[scale=0.90909]{figs/redblack-example}}
  \caption{A red-black tree on which to practice.}
  \figlabel{redblack-example2}
\end{figure}

\begin{exc}
  Illustrate the 2-4 tree that corresponds to the \texttt{RedBlackTree} in
  \figref{redblack-example2}.
\end{exc}

\begin{exc}
  Illustrate the addition of 13, then 3.5, then 3.3 on the \texttt{RedBlackTree}
  in \figref{redblack-example2}.
\end{exc}

\begin{exc}
  Illustrate the removal of 11, then 9, then 5 on the \texttt{RedBlackTree} in 
  \figref{redblack-example2}.
\end{exc}

\begin{exc}
  Show that, for arbitrarily large values of \ensuremath{\mathtt{{\color{var}n}}}, there are red-black
  trees with \ensuremath{\mathtt{{\color{var}n}}} nodes that have height $2\log \ensuremath{\mathtt{{\color{var}n}}}-O(1)$.
\end{exc}

\begin{exc}
  Consider the operations \ensuremath{\mathtt{pushBlack({\color{var}u})}} and \ensuremath{\mathtt{pullBlack({\color{var}u})}}.  What do
  these operations do to the underlying 2-4 tree that is being simulated
  by the red-black tree?
\end{exc}

\begin{exc}
  Show that, for arbitrarily large values of \ensuremath{\mathtt{{\color{var}n}}}, there exist sequences
  of \ensuremath{\mathtt{add({\color{var}x})}} and \ensuremath{\mathtt{remove({\color{var}x})}} operations that lead to red-black trees
  with \ensuremath{\mathtt{{\color{var}n}}} nodes that have height $2\log \ensuremath{\mathtt{{\color{var}n}}}-O(1)$.
\end{exc}



\begin{exc}
  Why does the method \ensuremath{\mathtt{remove({\color{var}x})}} in the \texttt{RedBlackTree} implementation
  perform the assignment \ensuremath{\mathtt{{\color{var}u}.{\color{var}parent}={\color{var}w}.{\color{var}parent}}}?  Shouldn't this already
  be done by the call to \ensuremath{\mathtt{splice({\color{var}w})}}?
\end{exc}

\begin{exc}
  Suppose a 2-4 tree, $T$, has $\ensuremath{\mathtt{{\color{var}n}}}_\ell$ leaves and $\ensuremath{\mathtt{{\color{var}n}}}_i$ internal nodes.
  \begin{enumerate}
    \item What is the minimum value of $\ensuremath{\mathtt{{\color{var}n}}}_i$, as a function of $\ensuremath{\mathtt{{\color{var}n}}}_\ell$?
    \item What is the maximum value of $\ensuremath{\mathtt{{\color{var}n}}}_i$, as a function of $\ensuremath{\mathtt{{\color{var}n}}}_\ell$?
    \item If $T'$ is a red-black tree that represents $T$, then how many red
     nodes does $T'$ have?
  \end{enumerate}
\end{exc}

\begin{exc}
  Suppose you are given a binary search tree with \ensuremath{\mathtt{{\color{var}n}}} nodes and a
  height of at most $2\log \ensuremath{\mathtt{{\color{var}n}}}-2$.  Is it always possible to colour the
  nodes red and black so that the tree satisfies the black-height and
  no-red-edge properties?  If so, can it also be made to satisfy the
  left-leaning property?
\end{exc}

\begin{exc}\exclabel{redblack-merge}
  Suppose you have two red-black trees $T_1$ and $T_2$ that have the
  same black height, $h$, and such that the largest key in $T_1$ is smaller
  than the smallest key in $T_2$.  Show how to merge $T_1$ and $T_2$
  into a single red-black tree in $O(h)$ time.
\end{exc}

\begin{exc}
  Extend your solution to \excref{redblack-merge} to the case where the
  two trees $T_1$ and $T_2$ have different black heights, $h_1\neq h_2$.
  The running-time should be $O(\max\{h_1,h_2\})$.
\end{exc}



\begin{exc}
  Prove that, during an \ensuremath{\mathtt{add({\color{var}x})}} operation, an AVL tree must perform
  at most one rebalancing operation (that involves at most two rotations;
  see \figref{avl-rebalance}).  Give an example of an AVL tree and a
  \ensuremath{\mathtt{remove({\color{var}x})}} operation on that tree that requires on the order of $\log
  \ensuremath{\mathtt{{\color{var}n}}}$ rebalancing operations.
\end{exc}

\begin{exc}
  Implement an \texttt{AVLTree} class that implements AVL trees as described
  above.  Compare its performance to that of the \texttt{RedBlackTree}
  implementation.   Which implementation has a faster \ensuremath{\mathtt{find({\color{var}x})}} operation?
\end{exc}

\begin{exc}
  Design and implement a series of experiments that compare the relative
  performance of \ensuremath{\mathtt{find({\color{var}x})}}, \ensuremath{\mathtt{add({\color{var}x})}}, and \ensuremath{\mathtt{remove({\color{var}x})}} for the \texttt{SSet} implemeentations \texttt{SkiplistSSet},
  \texttt{ScapegoatTree}, \texttt{Treap}, and \texttt{RedBlackTree}.  Be sure to include
  multiple test scenarios, including cases where the data is random,
  already sorted, is removed in random order, is removed in sorted order,
  and so on.
\end{exc}

\clearpage

\clearpage
\chapter{Heaps}
\chaplabel{heaps}

In this chapter, we discuss two implementations of the extremely useful
priority \texttt{Queue} data structure.  Both of these structures are a special
kind of binary tree called a \emph{heap}, 
\index{heap}%
\index{binary heap}%
\index{heap!binary}%
which means ``a disorganized
pile.''  This is in contrast to binary search trees that can be thought
of as a highly organized pile.

The first heap implementation uses an array to simulate a complete binary
tree.  This very fast implementation is the basis of one of the fastest
known sorting algorithms, namely heapsort (see \secref{heapsort}).
The second implementation is based on more flexible binary trees.
It supports a \ensuremath{\mathtt{meld({\color{var}h})}} operation that allows the priority queue to
absorb the elements of a second priority queue \ensuremath{\mathtt{{\color{var}h}}}.

\section{\texttt{BinaryHeap}: An Implicit Binary Tree}
\seclabel{binaryheap}

\index{BinaryHeap@\texttt{BinaryHeap}}%
Our first implementation of a (priority) \texttt{Queue} is based on a technique
that is over four hundred years old.  \emph{Eytzinger's method}
\index{Eytzinger's method}%
allows us
to represent a complete binary tree as an array by laying out the nodes
of the tree in breadth-first order (see \secref{bintree:traversal}).
In this way, the root is stored at position 0, the root's left child is
stored at position 1, the root's right child at position 2, the left
child of the left child of the root is stored at position 3, and so
on. See \figref{eytzinger}.

\begin{figure}
  \begin{center}
    \includegraphics[scale=0.90909]{figs/eytzinger}
  \end{center}
  \caption{Eytzinger's method represents a complete binary tree as an array.}
  \figlabel{eytzinger}
\end{figure}

If we apply Eytzinger's method to a sufficiently large tree, some
patterns emerge.  The left child of the node at index \ensuremath{\mathtt{{\color{var}i}}} is at index
$\ensuremath{\mathtt{left({\color{var}i})}}=2\ensuremath{\mathtt{{\color{var}i}}}+1$ and the right child of the node at index \ensuremath{\mathtt{{\color{var}i}}} is at
index $\ensuremath{\mathtt{right({\color{var}i})}}=2\ensuremath{\mathtt{{\color{var}i}}}+2$.  The parent of the node at index \ensuremath{\mathtt{{\color{var}i}}} is at
index $\ensuremath{\mathtt{parent({\color{var}i})}}=(\ensuremath{\mathtt{{\color{var}i}}}-1)/2$.


 \ \vspace{.5ex}\begin{Verbatim}[gobble=1,tabsize=2,frame=single,framerule=1.2pt,commandchars=\\@\$,label=\texttt{BinaryHeap},labelposition=topline]
	@\color@keyword$int$ left(@\color@keyword$int$ @\color@var$i$) {
		@\color@var$return$ 2*@\color@var$i$ + 1;
	}
	@\color@keyword$int$ right(@\color@keyword$int$ @\color@var$i$) {
		@\color@var$return$ 2*@\color@var$i$ + 2;
	}
	@\color@keyword$int$ parent(@\color@keyword$int$ @\color@var$i$) {
		return (@\color@var$i$-1)/2;
	}
\end{Verbatim}


A \texttt{BinaryHeap} uses this technique to implicitly represent a complete
binary tree in which the elements are \emph{heap-ordered}:
\index{heap-ordered binary tree}%
\index{binary tree!heap-ordered}%
\index{heap order}%
The value
stored at any index \ensuremath{\mathtt{{\color{var}i}}} is not smaller than the value stored at index
\ensuremath{\mathtt{parent({\color{var}i})}}, with the exception of the root value, $\ensuremath{\mathtt{{\color{var}i}}}=0$.  It follows
that the smallest value in the priority \texttt{Queue} is therefore stored at
position 0 (the root).

In a \texttt{BinaryHeap}, the \ensuremath{\mathtt{{\color{var}n}}} elements are stored in an array \ensuremath{\mathtt{{\color{var}a}}}:


 \ \vspace{.5ex}\begin{Verbatim}[gobble=1,tabsize=2,frame=single,framerule=1.2pt,commandchars=\\@\$,label=\texttt{BinaryHeap},labelposition=topline]
	@\color@keyword$T$[] @\color@var$a$;
	@\color@keyword$int$ @\color@var$n$;
\end{Verbatim}


Implementing the \ensuremath{\mathtt{add({\color{var}x})}} operation is fairly straightforward.  As with
all array-based structures, we first check to see if \ensuremath{\mathtt{{\color{var}a}}} is full (by checking if $\ensuremath{\mathtt{{\color{var}a}.{\color{var}length}}}=\ensuremath{\mathtt{{\color{var}n}}}$) and, if so, we grow \ensuremath{\mathtt{{\color{var}a}}}.  Next, we place \ensuremath{\mathtt{{\color{var}x}}} at location
\ensuremath{\mathtt{{\color{var}a}[{\color{var}n}]}} and increment \ensuremath{\mathtt{{\color{var}n}}}.  At this point, all that remains is to ensure
that we maintain the heap property.  We do this by repeatedly swapping
\ensuremath{\mathtt{{\color{var}x}}} with its parent until \ensuremath{\mathtt{{\color{var}x}}} is no longer smaller than its parent.
See \figref{heap-insert}.


 \ \vspace{.5ex}\begin{Verbatim}[gobble=1,tabsize=2,frame=single,framerule=1.2pt,commandchars=\\@\$,label=\texttt{BinaryHeap},labelposition=topline]
	@\color@var$boolean$ add(@\color@keyword$T$ @\color@var$x$) {
		@\color@keyword$if$ (@\color@var$n$ + 1 > @\color@var$a$.@\color@var$length$) resize();
		@\color@var$a$[@\color@var$n$++] = @\color@var$x$;
		bubbleUp(@\color@var$n$-1);
		@\color@var$return$ @\color@var$true$;
	}
	@\color@keyword$void$ bubbleUp(@\color@keyword$int$ @\color@var$i$) {
		@\color@keyword$int$ @\color@var$p$ = parent(@\color@var$i$);
		@\color@keyword$while$ (@\color@var$i$ > 0 && compare(@\color@var$a$[@\color@var$i$], @\color@var$a$[@\color@var$p$]) < 0) {
			swap(@\color@var$i$,@\color@var$p$);
			@\color@var$i$ = @\color@var$p$;
			@\color@var$p$ = parent(@\color@var$i$);
		}
	}
\end{Verbatim}


\begin{figure}
  \begin{center}
    \includegraphics[height=\QuarterHeightScaleIfNeeded]{figs/heap-insert-1} \\
    \includegraphics[height=\QuarterHeightScaleIfNeeded]{figs/heap-insert-2} \\
    \includegraphics[height=\QuarterHeightScaleIfNeeded]{figs/heap-insert-3} \\
    \includegraphics[height=\QuarterHeightScaleIfNeeded]{figs/heap-insert-4} \\
  \end{center}
  \caption[Adding to a BinaryHeap]{Adding the value 6 to a \texttt{BinaryHeap}.}
  \figlabel{heap-insert}
\end{figure}

Implementing the \ensuremath{\mathtt{remove()}} operation, which removes the smallest value
from the heap, is a little trickier.  We know where the smallest value is
(at the root), but we need to replace it after we remove it and ensure
that we maintain the heap property.

The easiest way to do this is to replace the root with the value \ensuremath{\mathtt{{\color{var}a}[{\color{var}n}-1]}}, delete
that value, and decrement \ensuremath{\mathtt{{\color{var}n}}}.  Unfortunately, the new root element is now
probably not the smallest element, so it needs to be moved downwards.
We do this by repeatedly comparing this element to its two children.
If it is the smallest of the three then we are done.  Otherwise, we swap
this element with the smallest of its two children and continue.


 \ \vspace{.5ex}\begin{Verbatim}[gobble=1,tabsize=2,frame=single,framerule=1.2pt,commandchars=\\@\$,label=\texttt{BinaryHeap},labelposition=topline]
	@\color@keyword$T$ remove() {
		@\color@keyword$T$ @\color@var$x$ = @\color@var$a$[0];
		@\color@var$a$[0] = @\color@var$a$[--@\color@var$n$];
		trickleDown(0);
		@\color@keyword$if$ (3*@\color@var$n$ < @\color@var$a$.@\color@var$length$) resize();
		@\color@var$return$ @\color@var$x$;
	}
	@\color@keyword$void$ trickleDown(@\color@keyword$int$ @\color@var$i$) {
		@\color@keyword$do$ {
			@\color@keyword$int$ @\color@var$j$ = -1;
			@\color@keyword$int$ @\color@var$r$ = right(@\color@var$i$);
			@\color@keyword$if$ (@\color@var$r$ < @\color@var$n$ && compare(@\color@var$a$[@\color@var$r$], @\color@var$a$[@\color@var$i$]) < 0) {
				@\color@keyword$int$ @\color@var$l$ = left(@\color@var$i$);
				@\color@keyword$if$ (compare(@\color@var$a$[@\color@var$l$], @\color@var$a$[@\color@var$r$]) < 0) {
					@\color@var$j$ = @\color@var$l$;
				} @\color@keyword$else$ {
					@\color@var$j$ = @\color@var$r$;
				}
			} @\color@keyword$else$ {
				@\color@keyword$int$ @\color@var$l$ = left(@\color@var$i$);
				@\color@keyword$if$ (@\color@var$l$ < @\color@var$n$ && compare(@\color@var$a$[@\color@var$l$], @\color@var$a$[@\color@var$i$]) < 0) {
					@\color@var$j$ = @\color@var$l$;
				}
			}
			@\color@keyword$if$ (@\color@var$j$ >= 0)	swap(@\color@var$i$, @\color@var$j$);
			@\color@var$i$ = @\color@var$j$;
		} @\color@keyword$while$ (@\color@var$i$ >= 0);
	}
\end{Verbatim}


\begin{figure}
  \begin{center}
    \includegraphics[height=\QuarterHeightScaleIfNeeded]{figs/heap-remove-1} \\
    \includegraphics[height=\QuarterHeightScaleIfNeeded]{figs/heap-remove-2} \\
    \includegraphics[height=\QuarterHeightScaleIfNeeded]{figs/heap-remove-3} \\
    \includegraphics[height=\QuarterHeightScaleIfNeeded]{figs/heap-remove-4} \\
  \end{center}
  \caption[Removing from a BinaryHeap]{Removing the minimum value, 4, from a \texttt{BinaryHeap}.}
  \figlabel{heap-remove}
\end{figure}



As with other array-based structures, we will ignore the time spent in
calls to \ensuremath{\mathtt{resize()}}, since these can be accounted for using the amortization
argument from \lemref{arraystack-amortized}.  The running times of
both \ensuremath{\mathtt{add({\color{var}x})}} and \ensuremath{\mathtt{remove()}} then depend on the height of the (implicit)
binary tree.  Luckily, this is a \emph{complete}
\index{binary tree!complete}%
\index{complete binary tree}%
binary tree;  every level
except the last has the maximum possible number of nodes.  Therefore,
if the height of this tree is $h$, then it has at least $2^h$ nodes.
Stated another way
\[
  \ensuremath{\mathtt{{\color{var}n}}} \ge 2^h \enspace .
\]  
Taking logarithms on both sides of this equation gives
\[
   h \le \log \ensuremath{\mathtt{{\color{var}n}}} \enspace .
\]
Therefore, both the \ensuremath{\mathtt{add({\color{var}x})}} and \ensuremath{\mathtt{remove()}} operation run in $O(\log \ensuremath{\mathtt{{\color{var}n}}})$ time.

\subsection{Summary}

The following theorem summarizes the performance of a \texttt{BinaryHeap}:

\begin{thm}\thmlabel{binaryheap}
  A \texttt{BinaryHeap} implements the (priority) \texttt{Queue} interface.  Ignoring
  the cost of calls to \ensuremath{\mathtt{resize()}}, a \texttt{BinaryHeap} supports the operations
  \ensuremath{\mathtt{add({\color{var}x})}} and \ensuremath{\mathtt{remove()}} in $O(\log \ensuremath{\mathtt{{\color{var}n}}})$ time per operation.

  Furthermore, beginning with an empty \texttt{BinaryHeap}, any sequence of $m$
  \ensuremath{\mathtt{add({\color{var}x})}} and \ensuremath{\mathtt{remove()}} operations results in a total of $O(m)$
  time spent during all calls to \ensuremath{\mathtt{resize()}}.
\end{thm}

\section{\texttt{MeldableHeap}: A Randomized Meldable Heap}
\seclabel{meldableheap}

\index{MeldableHeap@\texttt{MeldableHeap}}%
In this section, we describe the \texttt{MeldableHeap}, a priority \texttt{Queue}
implementation in which the underlying structure is also a heap-ordered
binary tree.  However, unlike a \texttt{BinaryHeap} in which the underlying
binary tree is completely defined by the number of elements, there
are no restrictions on the shape of the binary tree that underlies
a \texttt{MeldableHeap}; anything goes.

The \ensuremath{\mathtt{add({\color{var}x})}} and \ensuremath{\mathtt{remove()}} operations in a \texttt{MeldableHeap} are
implemented in terms of the \ensuremath{\mathtt{merge({\color{var}h1},{\color{var}h2})}} operation.  This operation
takes two heap nodes \ensuremath{\mathtt{{\color{var}h1}}} and \ensuremath{\mathtt{{\color{var}h2}}} and merges them, returning a heap
node that is the root of a heap that contains all elements in the subtree
rooted at \ensuremath{\mathtt{{\color{var}h1}}} and all elements in the subtree rooted at \ensuremath{\mathtt{{\color{var}h2}}}.

The nice thing about a \ensuremath{\mathtt{merge({\color{var}h1},{\color{var}h2})}} operation is that it can be
defined recursively. See \figref{meldable-merge}.  If either \ensuremath{\mathtt{{\color{var}h1}}} or
\ensuremath{\mathtt{{\color{var}h2}}} is \ensuremath{\mathtt{{\color{var}nil}}}, then we are merging with an empty set, so we return \ensuremath{\mathtt{{\color{var}h2}}}
or \ensuremath{\mathtt{{\color{var}h1}}}, respectively.  Otherwise, assume $\ensuremath{\mathtt{{\color{var}h1}.{\color{var}x}}} \le \ensuremath{\mathtt{{\color{var}h2}.{\color{var}x}}}$ since,
if $\ensuremath{\mathtt{{\color{var}h1}.{\color{var}x}}} > \ensuremath{\mathtt{{\color{var}h2}.{\color{var}x}}}$, then we can reverse the roles of \ensuremath{\mathtt{{\color{var}h1}}} and \ensuremath{\mathtt{{\color{var}h2}}}.
Then we know that the root of the merged heap will contain \ensuremath{\mathtt{{\color{var}h1}.{\color{var}x}}}, and
we can recursively merge \ensuremath{\mathtt{{\color{var}h2}}} with \ensuremath{\mathtt{{\color{var}h1}.{\color{var}left}}} or \ensuremath{\mathtt{{\color{var}h1}.{\color{var}right}}}, as we wish.
This is where randomization comes in, and we toss a coin to decide
whether to merge \ensuremath{\mathtt{{\color{var}h2}}} with \ensuremath{\mathtt{{\color{var}h1}.{\color{var}left}}} or \ensuremath{\mathtt{{\color{var}h1}.{\color{var}right}}}:


 \ \vspace{.5ex}\begin{Verbatim}[gobble=1,tabsize=2,frame=single,framerule=1.2pt,commandchars=\\@\$,label=\texttt{MeldableHeap},labelposition=topline]
	Node<@\color@keyword$T$> merge(Node<@\color@keyword$T$> @\color@var$h1$, Node<@\color@keyword$T$> @\color@var$h2$) {
		@\color@keyword$if$ (@\color@var$h1$ == @\color@var$nil$) @\color@var$return$ @\color@var$h2$;
		@\color@keyword$if$ (@\color@var$h2$ == @\color@var$nil$) @\color@var$return$ @\color@var$h1$;
		@\color@keyword$if$ (compare(@\color@var$h2$.@\color@var$x$, @\color@var$h1$.@\color@var$x$) < 0) @\color@var$return$ merge(@\color@var$h2$, @\color@var$h1$);
		@\color@comment$// now we know h1.x <= h2.x$
		@\color@keyword$if$ (@\color@var$rand$.nextBoolean()) {
			@\color@var$h1$.@\color@var$left$ = merge(@\color@var$h1$.@\color@var$left$, @\color@var$h2$);
			@\color@var$h1$.@\color@var$left$.@\color@var$parent$ = @\color@var$h1$;
		} @\color@keyword$else$ {
			@\color@var$h1$.@\color@var$right$ = merge(@\color@var$h1$.@\color@var$right$, @\color@var$h2$);
			@\color@var$h1$.@\color@var$right$.@\color@var$parent$ = @\color@var$h1$;
		}
		@\color@var$return$ @\color@var$h1$;
	}
\end{Verbatim}


\begin{figure}
  \centering{\includegraphics[width=\ScaleIfNeeded]{figs/meldable-merge}}
  \caption[Merging in a MeldableHeap]{Merging \ensuremath{\mathtt{{\color{var}h1}}} and \ensuremath{\mathtt{{\color{var}h2}}} is done by merging \ensuremath{\mathtt{{\color{var}h2}}} with one of
  \ensuremath{\mathtt{{\color{var}h1}.{\color{var}left}}} or \ensuremath{\mathtt{{\color{var}h1}.{\color{var}right}}}.}
  \figlabel{meldable-merge}
\end{figure}

In the next section, we show that \ensuremath{\mathtt{merge({\color{var}h1},{\color{var}h2})}} runs in $O(\log \ensuremath{\mathtt{{\color{var}n}}})$
expected time, where \ensuremath{\mathtt{{\color{var}n}}} is the total number of elements in \ensuremath{\mathtt{{\color{var}h1}}} and \ensuremath{\mathtt{{\color{var}h2}}}.

With access to a \ensuremath{\mathtt{merge({\color{var}h1},{\color{var}h2})}} operation, the \ensuremath{\mathtt{add({\color{var}x})}} operation is easy.  We create a new node \ensuremath{\mathtt{{\color{var}u}}} containing \ensuremath{\mathtt{{\color{var}x}}} and then merge \ensuremath{\mathtt{{\color{var}u}}} with the root of our heap:


 \ \vspace{.5ex}\begin{Verbatim}[gobble=1,tabsize=2,frame=single,framerule=1.2pt,commandchars=\\@\$,label=\texttt{MeldableHeap},labelposition=topline]
	@\color@var$boolean$ add(@\color@keyword$T$ @\color@var$x$) {
		Node<@\color@keyword$T$> @\color@var$u$ = newNode();
		@\color@var$u$.@\color@var$x$ = @\color@var$x$;
		@\color@var$r$ = merge(@\color@var$u$, @\color@var$r$);
		@\color@var$r$.@\color@var$parent$ = @\color@var$nil$;
		@\color@var$n$++;
		@\color@var$return$ @\color@var$true$;
	}
\end{Verbatim}

This takes $O(\log (\ensuremath{\mathtt{{\color{var}n}}}+1)) = O(\log \ensuremath{\mathtt{{\color{var}n}}})$ expected time.

The \ensuremath{\mathtt{remove()}} operation is similarly easy.  The node we want to remove
is the root, so we just merge its two children and make the result the root:


 \ \vspace{.5ex}\begin{Verbatim}[gobble=1,tabsize=2,frame=single,framerule=1.2pt,commandchars=\\@\$,label=\texttt{MeldableHeap},labelposition=topline]
	@\color@keyword$T$ remove() {
		@\color@keyword$T$ @\color@var$x$ = @\color@var$r$.@\color@var$x$;
		@\color@var$r$ = merge(@\color@var$r$.@\color@var$left$, @\color@var$r$.@\color@var$right$);
		@\color@keyword$if$ (@\color@var$r$ != @\color@var$nil$) @\color@var$r$.@\color@var$parent$ = @\color@var$nil$;
		@\color@var$n$--;
		@\color@var$return$ @\color@var$x$;
	}
\end{Verbatim}

Again, this takes $O(\log \ensuremath{\mathtt{{\color{var}n}}})$ expected time.

Additionally, a \texttt{MeldableHeap} can implement many other operations in
$O(\log \ensuremath{\mathtt{{\color{var}n}}})$ expected time, including:
\begin{itemize}
\item \ensuremath{\mathtt{remove({\color{var}u})}}: remove the node \ensuremath{\mathtt{{\color{var}u}}} (and its key \ensuremath{\mathtt{{\color{var}u}.{\color{var}x}}}) from the heap.
\item \ensuremath{\mathtt{absorb({\color{var}h})}}: add all the elements of the \texttt{MeldableHeap} \ensuremath{\mathtt{{\color{var}h}}} to this heap, emptying \ensuremath{\mathtt{{\color{var}h}}} in the process.
\end{itemize}
Each of these operations can be implemented using a constant number of
\ensuremath{\mathtt{merge({\color{var}h1},{\color{var}h2})}} operations that each take $O(\log \ensuremath{\mathtt{{\color{var}n}}})$ expected time.

\subsection{Analysis of \ensuremath{\mathtt{merge({\color{var}h1},{\color{var}h2})}}}

The analysis of \ensuremath{\mathtt{merge({\color{var}h1},{\color{var}h2})}} is based on the analysis of a random walk
in a binary tree.  A \emph{random walk} in a binary tree starts at the
root of the tree.  At each step in the random walk, a coin is tossed and,
depending on the result of this coin toss, the walk proceeds to the left
or to the right child of the current node.  The walk ends when it falls
off the tree (the current node becomes \ensuremath{\mathtt{{\color{var}nil}}}).

The following lemma is somewhat remarkable because it does not depend
at all on the shape of the binary tree:

\begin{lem}\lemlabel{tree-random-walk}
The expected length of a random walk in a binary tree with \ensuremath{\mathtt{{\color{var}n}}} nodes is at most \ensuremath{\mathtt{\log ({\color{var}n}+1)}}.
\end{lem}

\begin{proof}
The proof is by induction on \ensuremath{\mathtt{{\color{var}n}}}. In the base case, $\ensuremath{\mathtt{{\color{var}n}}}=0$ and the walk
has length $0=\log (\ensuremath{\mathtt{{\color{var}n}}}+1)$.  Suppose now that the result is true for
all non-negative integers $\ensuremath{\mathtt{{\color{var}n}}}'< \ensuremath{\mathtt{{\color{var}n}}}$.

Let $\ensuremath{\mathtt{{\color{var}n}}}_1$ denote the size of the root's left subtree, so that
$\ensuremath{\mathtt{{\color{var}n}}}_2=\ensuremath{\mathtt{{\color{var}n}}}-\ensuremath{\mathtt{{\color{var}n}}}_1-1$ is the size of the root's right subtree.  Starting at
the root, the walk takes one step and then continues in a subtree of
size $\ensuremath{\mathtt{{\color{var}n}}}_1$ or $\ensuremath{\mathtt{{\color{var}n}}}_2$.  By our inductive hypothesis, the expected
length of the walk is then
\[
    \E[W] = 1 + \frac{1}{2}\log (\ensuremath{\mathtt{{\color{var}n}}}_1+1) + \frac{1}{2}\log (\ensuremath{\mathtt{{\color{var}n}}}_2+1)  \enspace , 
\] 
since each of $\ensuremath{\mathtt{{\color{var}n}}}_1$ and $\ensuremath{\mathtt{{\color{var}n}}}_2$ are less than $\ensuremath{\mathtt{{\color{var}n}}}$.  Since $\log$
is a concave function, $\E[W]$ is maximized when $\ensuremath{\mathtt{{\color{var}n}}}_1=\ensuremath{\mathtt{{\color{var}n}}}_2=(\ensuremath{\mathtt{{\color{var}n}}}-1)/2$.
Therefore,
 the expected number of steps taken by the random walk is 
\begin{align*}
    \E[W] 
    & = 1 + \frac{1}{2}\log (\ensuremath{\mathtt{{\color{var}n}}}_1+1) + \frac{1}{2}\log (\ensuremath{\mathtt{{\color{var}n}}}_2+1) \\
   & \le  1 + \log ((\ensuremath{\mathtt{{\color{var}n}}}-1)/2+1) \\
   & =  1 + \log ((\ensuremath{\mathtt{{\color{var}n}}}+1)/2) \\
   & =  \log (\ensuremath{\mathtt{{\color{var}n}}}+1)  \enspace . \qedhere 
\end{align*}
\end{proof}

We make a quick digression to note that, for readers who know a little
about information theory, the proof of \lemref{tree-random-walk} can
be stated in terms of entropy.  
\begin{proof}[Information Theoretic Proof of \lemref{tree-random-walk}]
Let $d_i$ denote the depth of the $i$th external node and recall that a
binary tree with \ensuremath{\mathtt{{\color{var}n}}} nodes has \ensuremath{\mathtt{{\color{var}n}+1}} external nodes.  The probability
of the random walk reaching the $i$th external node is exactly
$p_i=1/2^{d_i}$, so the expected length of the random walk is given by
\[
   H=\sum_{i=0}^{\ensuremath{\mathtt{{\color{var}n}}}} p_id_i
    =\sum_{i=0}^{\ensuremath{\mathtt{{\color{var}n}}}} p_i\log\left(2^{d_i}\right)
    = \sum_{i=0}^{\ensuremath{\mathtt{{\color{var}n}}}}p_i\log({1/p_i})
\]
The right hand side of this equation is easily recognizable as the
entropy of a probability distribution over $\ensuremath{\mathtt{{\color{var}n}}}+1$ elements.  A basic
fact about the entropy of a distribution over $\ensuremath{\mathtt{{\color{var}n}}}+1$ elements is that
it does not exceed $\log(\ensuremath{\mathtt{{\color{var}n}}}+1)$, which proves the lemma.
\end{proof}

With this result on random walks, we can now easily prove that the
running time of the \ensuremath{\mathtt{merge({\color{var}h1},{\color{var}h2})}} operation is $O(\log \ensuremath{\mathtt{{\color{var}n}}})$.

\begin{lem}
  If \ensuremath{\mathtt{{\color{var}h1}}} and \ensuremath{\mathtt{{\color{var}h2}}} are the roots of two heaps containing $\ensuremath{\mathtt{{\color{var}n}}}_1$
  and $\ensuremath{\mathtt{{\color{var}n}}}_2$ nodes, respectively, then the expected running time of
  \ensuremath{\mathtt{merge({\color{var}h1},{\color{var}h2})}} is at most $O(\log \ensuremath{\mathtt{{\color{var}n}}})$, where $\ensuremath{\mathtt{{\color{var}n}}}=\ensuremath{\mathtt{{\color{var}n}}}_1+\ensuremath{\mathtt{{\color{var}n}}}_2$.
\end{lem}

\begin{proof}
  Each step of the merge algorithm takes one step of a random walk,
  either in the heap rooted at \ensuremath{\mathtt{{\color{var}h1}}} or the heap rooted at \ensuremath{\mathtt{{\color{var}h2}}}.
  The algorithm terminates when either of these two random walks fall
  out of its corresponding tree (when $\ensuremath{\mathtt{{\color{var}h1}}}=\ensuremath{\mathtt{{\color{keyword}null}}}$ or $\ensuremath{\mathtt{{\color{var}h2}}}=\ensuremath{\mathtt{{\color{keyword}null}}}$).
  Therefore, the expected number of steps performed by the merge algorithm
  is at most
  \[
     \log (\ensuremath{\mathtt{{\color{var}n}}}_1+1) + \log (\ensuremath{\mathtt{{\color{var}n}}}_2+1) \le 2\log \ensuremath{\mathtt{{\color{var}n}}} \enspace . \qedhere
  \]
\end{proof}

\subsection{Summary}

The following theorem summarizes the performance of a \texttt{MeldableHeap}:

\begin{thm}\thmlabel{meldableheap}
  A \texttt{MeldableHeap} implements the (priority) \texttt{Queue} interface.
  A \texttt{MeldableHeap} supports the operations \ensuremath{\mathtt{add({\color{var}x})}} and \ensuremath{\mathtt{remove()}}
  in $O(\log \ensuremath{\mathtt{{\color{var}n}}})$ expected time per operation.
\end{thm}

\section{Discussion and Exercises}

The implicit representation of a complete binary tree as an array,
or list, seems to have been first proposed by Eytzinger \cite{e1590}.
He used this representation in books containing pedigree family trees
\index{pedigree family tree}%
of noble families.  The \texttt{BinaryHeap} data structure described here was
first introduced by Williams \cite{w64}.

The randomized \texttt{MeldableHeap} data structure described here appears
to have first been proposed by Gambin and Malinowski \cite{gm98}.
Other meldable heap implementations exist, including 
leftist heaps \cite[Section~5.3.2]{c72,k97v3},
\index{leftist heap}%
\index{heap!leftist}%
binomial heaps \cite{v78},
\index{binomial heap}%
\index{heap!binomial}%
Fibonacci heaps \cite{ft87}, 
\index{Fibonacci heap}%
\index{heap!Fibonacci}%
pairing heaps \cite{fsst86},\
\index{pairing heap}%
\index{heap!pairing}%
 and skew heaps \cite{st83}, 
\index{skew heap}%
\index{heap!skew}%
although none of these are as simple as the \texttt{MeldableHeap}
structure.

Some of the above structures also support a \ensuremath{\mathtt{decreaseKey({\color{var}u},{\color{var}y})}} operation
\index{decreaseKey@\ensuremath{\mathtt{decreaseKey({\color{var}u},{\color{var}y})}}}%
in which the value stored at node \ensuremath{\mathtt{{\color{var}u}}} is decreased to \ensuremath{\mathtt{{\color{var}y}}}.  (It is a
pre-condition that $\ensuremath{\mathtt{{\color{var}y}}}\le\ensuremath{\mathtt{{\color{var}u}.{\color{var}x}}}$.)  In most of the preceding structures,
this operation can be supported in $O(\log \ensuremath{\mathtt{{\color{var}n}}})$ time by removing node
\ensuremath{\mathtt{{\color{var}u}}} and adding  \ensuremath{\mathtt{{\color{var}y}}}.  However, some of these structures can implement
\ensuremath{\mathtt{decreaseKey({\color{var}u},{\color{var}y})}} more efficiently.  In particular, \ensuremath{\mathtt{decreaseKey({\color{var}u},{\color{var}y})}}
takes $O(1)$ amortized time in Fibonacci heaps and $O(\log\log \ensuremath{\mathtt{{\color{var}n}}})$
amortized time in a special version of pairing heaps \cite{e09}.
This more efficient \ensuremath{\mathtt{decreaseKey({\color{var}u},{\color{var}y})}} operation has applications in
speeding up several graph algorithms, including Dijkstra's shortest path
algorithm \cite{ft87}.

\begin{exc}
  Illustrate the addition of the values 7 and then 3 to the \texttt{BinaryHeap}
  shown at the end of \figref{heap-insert}.
\end{exc}

\begin{exc}
  Illustrate the removal of the next two values (6 and 8) on the
  \texttt{BinaryHeap} shown at the end of \figref{heap-remove}.
\end{exc}

\begin{exc}
  Implement the \ensuremath{\mathtt{remove({\color{var}i})}} method, that removes the value stored in
  \ensuremath{\mathtt{{\color{var}a}[{\color{var}i}]}} in a \texttt{BinaryHeap}.  This method should run in $O(\log \ensuremath{\mathtt{{\color{var}n}}})$ time.
  Next, explain why this method is not likely to be useful.
\end{exc}

\begin{exc}\exclabel{general-eytzinger}
  \index{tree!$d$-ary}%
  A $d$-ary tree is a generalization of a binary tree in which each
  internal node has $d$ children.  Using Eytzinger's method it is also
  possible to represent complete $d$-ary trees using arrays.  Work out
  the equations that, given an index \ensuremath{\mathtt{{\color{var}i}}}, determine the index of \ensuremath{\mathtt{{\color{var}i}}}'s
  parent and each of \ensuremath{\mathtt{{\color{var}i}}}'s $d$ children in this representation.
\end{exc}

\begin{exc}
  \index{DaryHeap@\texttt{DaryHeap}}%
  Using what you learned in \excref{general-eytzinger}, design and
  implement a \emph{\texttt{DaryHeap}}, the $d$-ary generalization of a
  \texttt{BinaryHeap}. Analyze the running times of operations on a \texttt{DaryHeap}
  and test the performance of your \texttt{DaryHeap} implementation against
  that of the \texttt{BinaryHeap} implementation given here.
\end{exc}



\begin{exc}
  Illustrate the addition of the values 17 and then 82 in the
  \texttt{MeldableHeap} \ensuremath{\mathtt{{\color{var}h1}}} shown in \figref{meldable-merge}.  Use a coin to
  simulate a random bit when needed.
\end{exc}

\begin{exc}
  Illustrate the removal of the next two values (4 and 8) in the
  \texttt{MeldableHeap} \ensuremath{\mathtt{{\color{var}h1}}} shown in \figref{meldable-merge}.  Use a coin to
  simulate a random bit when needed.
\end{exc}

\begin{exc}
  Implement the \ensuremath{\mathtt{remove({\color{var}u})}} method, that removes the node \ensuremath{\mathtt{{\color{var}u}}} from
  a \texttt{MeldableHeap}.  This method should run in $O(\log \ensuremath{\mathtt{{\color{var}n}}})$ expected time.
\end{exc}

\begin{exc}
  Show how to find the second smallest value in a \texttt{BinaryHeap} or
  \texttt{MeldableHeap} in constant time.
\end{exc}

\begin{exc}
  Show how to find the $k$th smallest value in a \texttt{BinaryHeap} or
  \texttt{MeldableHeap} in $O(k\log k)$ time.  (Hint: Using another heap
  might help.)
\end{exc}

\begin{exc}
  Suppose you are given \ensuremath{\mathtt{{\color{var}k}}} sorted lists, of total length \ensuremath{\mathtt{{\color{var}n}}}.  Using
  a heap, show how to merge these into a single sorted list in $O(n\log
  k)$ time.  (Hint: Starting with the case $k=2$ can be instructive.)
\end{exc}









\clearpage

\clearpage
\chapter{Sorting Algorithms}
\chaplabel{sorting}

This chapter discusses algorithms for sorting a set of \ensuremath{\mathtt{{\color{var}n}}} items.
This might seem like a strange topic for a book on data structures, but
there are several good reasons for including it here.  The most obvious
reason is that two of these sorting algorithms (quicksort and heap-sort)
are intimately related to two of the data structures we have already
studied (random binary search trees and heaps, respectively).

The first part of this chapter discusses algorithms that sort using only
comparisons and presents three algorithms that run in $O(\ensuremath{\mathtt{{\color{var}n}}}\log \ensuremath{\mathtt{{\color{var}n}}})$
time.  As it turns out, all three algorithms are asymptotically optimal;
no algorithm that uses only comparisons can avoid doing roughly $\ensuremath{\mathtt{{\color{var}n}}}\log
\ensuremath{\mathtt{{\color{var}n}}}$ comparisons in the worst case and even the average case.

Before continuing, we should note that any of the \texttt{SSet} or priority
\texttt{Queue} implementations presented in previous chapters can also
be used to obtain an $O(\ensuremath{\mathtt{{\color{var}n}}}\log \ensuremath{\mathtt{{\color{var}n}}})$ time sorting algorithm.
For example, we can sort \ensuremath{\mathtt{{\color{var}n}}} items by performing \ensuremath{\mathtt{{\color{var}n}}} \ensuremath{\mathtt{add({\color{var}x})}}
operations followed by \ensuremath{\mathtt{{\color{var}n}}} \ensuremath{\mathtt{remove()}} operations on a \texttt{BinaryHeap}
or \texttt{MeldableHeap}. Alternatively, we can use \ensuremath{\mathtt{{\color{var}n}}} \ensuremath{\mathtt{add({\color{var}x})}} operations
on any of the binary search tree data structures and then perform an
in-order traversal (\excref{tree-traversal}) to extract the elements in
sorted order.  However, in both cases we go through a lot of overhead to
build a structure that is never fully used.  Sorting is such an important
problem that it is worthwhile developing direct methods that are as fast,
simple, and space-efficient as possible.

The second part of this chapter shows that, if we allow other
operations besides comparisons, then all bets are off.  Indeed, by using
array-indexing, it is possible to sort a set of \ensuremath{\mathtt{{\color{var}n}}} integers in the range
$\{0,\ldots,\ensuremath{\mathtt{{\color{var}n}}}^c-1\}$ in $O(c\ensuremath{\mathtt{{\color{var}n}}})$ time.



\section{Comparison-Based Sorting}

\index{comparison-based sorting}%
\index{sorting algorithm!comparison-based}%
In this section, we present three sorting algorithms: merge-sort,
quicksort, and heap-sort.  Each of these algorithms takes an input array \ensuremath{\mathtt{{\color{var}a}}}
and sorts the elements of \ensuremath{\mathtt{{\color{var}a}}} into non-decreasing order in $O(\ensuremath{\mathtt{{\color{var}n}}}\log \ensuremath{\mathtt{{\color{var}n}}})$
(expected) time.  These algorithms are all \emph{comparison-based}.
\javaonly{Their second argument, \ensuremath{\mathtt{{\color{var}c}}}, is a \texttt{Comparator}
\index{Comparator@\texttt{Comparator}}%
that implements
the \ensuremath{\mathtt{compare({\color{var}a},{\color{var}b})}}
\index{compare@\ensuremath{\mathtt{compare({\color{var}a},{\color{var}b})}}}%
method.}  These algorithms don't care what type
of data is being sorted; the only operation they do on the data is
comparisons using the \ensuremath{\mathtt{compare({\color{var}a},{\color{var}b})}} method. Recall, from \secref{sset},
that \ensuremath{\mathtt{compare({\color{var}a},{\color{var}b})}} returns a negative value if $\ensuremath{\mathtt{{\color{var}a}}}<\ensuremath{\mathtt{{\color{var}b}}}$, a positive
value if $\ensuremath{\mathtt{{\color{var}a}}}>\ensuremath{\mathtt{{\color{var}b}}}$, and zero if $\ensuremath{\mathtt{{\color{var}a}}}=\ensuremath{\mathtt{{\color{var}b}}}$.

\subsection{Merge-Sort}
\seclabel{merge-sort}

\index{merge-sort}%
The \emph{merge-sort} algorithm is a classic example of recursive divide
and conquer: 
\index{divide-and-conquer}%
If the length of \ensuremath{\mathtt{{\color{var}a}}} is at most 1, then \ensuremath{\mathtt{{\color{var}a}}} is already
sorted, so we do nothing.  Otherwise, we split \ensuremath{\mathtt{{\color{var}a}}} into two halves,
$\ensuremath{\mathtt{{\color{var}a0}}}=\ensuremath{\mathtt{{\color{var}a}[0]}},\ldots,\ensuremath{\mathtt{{\color{var}a}[{\color{var}n}/2-1]}}$ and $\ensuremath{\mathtt{{\color{var}a1}}}=\ensuremath{\mathtt{{\color{var}a}[{\color{var}n}/2]}},\ldots,\ensuremath{\mathtt{{\color{var}a}[{\color{var}n}-1]}}$.
We recursively sort \ensuremath{\mathtt{{\color{var}a0}}} and \ensuremath{\mathtt{{\color{var}a1}}}, and then we merge (the now sorted)
\ensuremath{\mathtt{{\color{var}a0}}} and \ensuremath{\mathtt{{\color{var}a1}}} to get our fully sorted array \ensuremath{\mathtt{{\color{var}a}}}:


 \ \vspace{.5ex}\begin{Verbatim}[gobble=1,tabsize=2,frame=single,framerule=1.2pt,commandchars=\\@\$,label=\texttt{Algorithms},labelposition=topline]
	<@\color@keyword$T$> @\color@keyword$void$ mergeSort(@\color@keyword$T$[] @\color@var$a$, Comparator<@\color@keyword$T$> @\color@var$c$) {
		@\color@keyword$if$ (@\color@var$a$.@\color@var$length$ <= 1) @\color@var$return$;
		@\color@keyword$T$[] @\color@var$a0$ = Arrays.copyOfRange(@\color@var$a$, 0, @\color@var$a$.@\color@var$length$/2);
		@\color@keyword$T$[] @\color@var$a1$ = Arrays.copyOfRange(@\color@var$a$, @\color@var$a$.@\color@var$length$/2, @\color@var$a$.@\color@var$length$);
		mergeSort(@\color@var$a0$, @\color@var$c$);
		mergeSort(@\color@var$a1$, @\color@var$c$);
		merge(@\color@var$a0$, @\color@var$a1$, @\color@var$a$, @\color@var$c$);
	}
\end{Verbatim}

\pcodeimport{ods/Algorithms.mergeSort(a)}
An example is shown in \figref{merge-sort}.
\begin{figure}
  \begin{center}
    \includegraphics[width=\ScaleIfNeeded]{figs/mergesort}
  \end{center}
  \caption[Merge sort]{The execution of \ensuremath{\mathtt{mergeSort({\color{var}a},{\color{var}c})}}}
  \figlabel{merge-sort}
\end{figure}

Compared to sorting, merging the two sorted arrays \ensuremath{\mathtt{{\color{var}a0}}} and \ensuremath{\mathtt{{\color{var}a1}}} is
fairly easy.  We add elements to \ensuremath{\mathtt{{\color{var}a}}} one at a time.  If \ensuremath{\mathtt{{\color{var}a0}}} or \ensuremath{\mathtt{{\color{var}a1}}}
is empty, then we add the next elements from the other (non-empty)
array. Otherwise, we take the minimum of the next element in \ensuremath{\mathtt{{\color{var}a0}}} and
the next element in \ensuremath{\mathtt{{\color{var}a1}}} and add it to \ensuremath{\mathtt{{\color{var}a}}}:


 \ \vspace{.5ex}\begin{Verbatim}[gobble=1,tabsize=2,frame=single,framerule=1.2pt,commandchars=\\@\$,label=\texttt{Algorithms},labelposition=topline]
	<@\color@keyword$T$> @\color@keyword$void$ merge(@\color@keyword$T$[] @\color@var$a0$, @\color@keyword$T$[] @\color@var$a1$, @\color@keyword$T$[] @\color@var$a$, Comparator<@\color@keyword$T$> @\color@var$c$) {
		@\color@keyword$int$ @\color@var$i0$ = 0, @\color@var$i1$ = 0;
		@\color@keyword$for$ (@\color@keyword$int$ @\color@var$i$ = 0; @\color@var$i$ < @\color@var$a$.@\color@var$length$; @\color@var$i$++) {
			@\color@keyword$if$ (@\color@var$i0$ == @\color@var$a0$.@\color@var$length$)
				@\color@var$a$[@\color@var$i$] = @\color@var$a1$[@\color@var$i1$++];
			@\color@keyword$else$ @\color@keyword$if$ (@\color@var$i1$ == @\color@var$a1$.@\color@var$length$)
				@\color@var$a$[@\color@var$i$] = @\color@var$a0$[@\color@var$i0$++];
			@\color@keyword$else$ @\color@keyword$if$ (compare(@\color@var$a0$[@\color@var$i0$], @\color@var$a1$[@\color@var$i1$]) < 0)
				@\color@var$a$[@\color@var$i$] = @\color@var$a0$[@\color@var$i0$++];
			@\color@keyword$else$ 
				@\color@var$a$[@\color@var$i$] = @\color@var$a1$[@\color@var$i1$++];
		}
	}
\end{Verbatim}

\pcodeimport{ods/Algorithms.merge(a0,a1,a)}
Notice that the \ensuremath{\mathtt{merge({\color{var}a0},{\color{var}a1},{\color{var}a},{\color{var}c})}} algorithm performs at most $\ensuremath{\mathtt{{\color{var}n}}}-1$
comparisons before running out of elements in one of \ensuremath{\mathtt{{\color{var}a0}}} or \ensuremath{\mathtt{{\color{var}a1}}}.

To understand the running-time of merge-sort, it is easiest to think
of it in terms of its recursion tree.  Suppose for now that \ensuremath{\mathtt{{\color{var}n}}} is a
power of two, so that $\ensuremath{\mathtt{{\color{var}n}}}=2^{\log \ensuremath{\mathtt{{\color{var}n}}}}$, and $\log \ensuremath{\mathtt{{\color{var}n}}}$ is an integer.
Refer to \figref{mergesort-recursion}. Merge-sort turns the problem of
sorting \ensuremath{\mathtt{{\color{var}n}}} elements into two problems, each of sorting $\ensuremath{\mathtt{{\color{var}n}}}/2$ elements.
These two subproblem are then turned into two problems each, for a total
of four subproblems, each of size $\ensuremath{\mathtt{{\color{var}n}}}/4$. These four subproblems become eight
subproblems, each of size $\ensuremath{\mathtt{{\color{var}n}}}/8$, and so on.  At the bottom of this process,
$\ensuremath{\mathtt{{\color{var}n}}}/2$ subproblems, each of size two, are converted into \ensuremath{\mathtt{{\color{var}n}}} problems,
each of size one.  For each subproblem of size $\ensuremath{\mathtt{{\color{var}n}}}/2^{i}$, the time
spent merging and copying data is $O(\ensuremath{\mathtt{{\color{var}n}}}/2^i)$.  Since there are $2^i$
subproblems of size $\ensuremath{\mathtt{{\color{var}n}}}/2^i$, the total time spent working on problems
of size $2^i$, not counting recursive calls, is
\[
       2^i\times O(\ensuremath{\mathtt{{\color{var}n}}}/2^i) = O(\ensuremath{\mathtt{{\color{var}n}}}) \enspace .
\]
Therefore, the total amount of time taken by merge-sort is
\[
   \sum_{i=0}^{\log \ensuremath{\mathtt{{\color{var}n}}}} O(\ensuremath{\mathtt{{\color{var}n}}}) = O(\ensuremath{\mathtt{{\color{var}n}}}\log \ensuremath{\mathtt{{\color{var}n}}}) \enspace .
\]

\begin{figure}
  \begin{center}
    \includegraphics[width=\ScaleIfNeeded]{figs/mergesort-recursion}
    \caption{The merge-sort recursion tree.}
    \figlabel{mergesort-recursion}
  \end{center}
\end{figure}

The proof of the following theorem is based on preceding analysis,
but has to be a little more careful to deal with the cases where \ensuremath{\mathtt{{\color{var}n}}}
is not a power of 2.
\begin{thm}
  The \javaonly{\ensuremath{\mathtt{mergeSort({\color{var}a},{\color{var}c})}}}\cpponly{mergeSort(a)}\pcodeonly{merge\_sort(a)} algorithm runs in $O(\ensuremath{\mathtt{{\color{var}n}}}\log \ensuremath{\mathtt{{\color{var}n}}})$ time and
  performs at most $\ensuremath{\mathtt{{\color{var}n}}}\log \ensuremath{\mathtt{{\color{var}n}}}$ comparisons.
\end{thm}

\begin{proof}
The proof is by induction on $\ensuremath{\mathtt{{\color{var}n}}}$.  The base case, in which $\ensuremath{\mathtt{{\color{var}n}}}=1$,
is trivial; when presented with an array of length 0 or 1 the algorithm
simply returns without performing any comparisons.

Merging two sorted lists of total length $\ensuremath{\mathtt{{\color{var}n}}}$ requires at most $\ensuremath{\mathtt{{\color{var}n}}}-1$
comparisons. Let $C(\ensuremath{\mathtt{{\color{var}n}}})$ denote the maximum number of comparisons performed by
\ensuremath{\mathtt{mergeSort({\color{var}a},{\color{var}c})}} on an array \ensuremath{\mathtt{{\color{var}a}}} of length \ensuremath{\mathtt{{\color{var}n}}}.  If $\ensuremath{\mathtt{{\color{var}n}}}$ is even, then we apply the inductive hypothesis to
the two subproblems and obtain
\begin{align*}
  C(\ensuremath{\mathtt{{\color{var}n}}}) 
  &\le \ensuremath{\mathtt{{\color{var}n}}}-1 + 2C(\ensuremath{\mathtt{{\color{var}n}}}/2) \\
  &\le \ensuremath{\mathtt{{\color{var}n}}}-1 + 2((\ensuremath{\mathtt{{\color{var}n}}}/2)\log(\ensuremath{\mathtt{{\color{var}n}}}/2)) \\
  &= \ensuremath{\mathtt{{\color{var}n}}}-1 + \ensuremath{\mathtt{{\color{var}n}}}\log(\ensuremath{\mathtt{{\color{var}n}}}/2) \\
  &= \ensuremath{\mathtt{{\color{var}n}}}-1 + \ensuremath{\mathtt{{\color{var}n}}}\log \ensuremath{\mathtt{{\color{var}n}}}-\ensuremath{\mathtt{{\color{var}n}}} \\
  &< \ensuremath{\mathtt{{\color{var}n}}}\log \ensuremath{\mathtt{{\color{var}n}}} \enspace .
\end{align*}
The case where $\ensuremath{\mathtt{{\color{var}n}}}$ is odd is slightly more complicated.  For this case,
we use two inequalities that are easy to verify:
\begin{equation}
  \log(x+1) \le \log(x) + 1 \enspace , \eqlabel{log-ineq-a}
\end{equation}
for all $x\ge 1$ and
\begin{equation}
  \log(x+1/2) + \log(x-1/2) \le 2\log(x) \enspace , \eqlabel{log-ineq-b}
\end{equation}
for all $x\ge 1/2$.  Inequality~\myeqref{log-ineq-a} comes from the fact that $\log(x)+1 = \log(2x)$ while \myeqref{log-ineq-b} follows from the fact that $\log$ is a concave function.  With these tools in hand we have, for odd \ensuremath{\mathtt{{\color{var}n}}},
\begin{align*}
  C(\ensuremath{\mathtt{{\color{var}n}}}) 
  &\le \ensuremath{\mathtt{{\color{var}n}}}-1 + C(\lceil \ensuremath{\mathtt{{\color{var}n}}}/2 \rceil) + C(\lfloor \ensuremath{\mathtt{{\color{var}n}}}/2 \rfloor) \\
  &\le \ensuremath{\mathtt{{\color{var}n}}}-1 + \lceil \ensuremath{\mathtt{{\color{var}n}}}/2 \rceil\log \lceil \ensuremath{\mathtt{{\color{var}n}}}/2 \rceil 
           + \lfloor \ensuremath{\mathtt{{\color{var}n}}}/2 \rfloor\log \lfloor \ensuremath{\mathtt{{\color{var}n}}}/2 \rfloor \\
  &= \ensuremath{\mathtt{{\color{var}n}}}-1 + (\ensuremath{\mathtt{{\color{var}n}}}/2 + 1/2)\log (\ensuremath{\mathtt{{\color{var}n}}}/2+1/2) 
           + (\ensuremath{\mathtt{{\color{var}n}}}/2 - 1/2) \log (\ensuremath{\mathtt{{\color{var}n}}}/2-1/2) \\
  &\le \ensuremath{\mathtt{{\color{var}n}}}-1 + \ensuremath{\mathtt{{\color{var}n}}}\log(\ensuremath{\mathtt{{\color{var}n}}}/2) + (1/2)(\log (\ensuremath{\mathtt{{\color{var}n}}}/2+1/2) 
           - \log (\ensuremath{\mathtt{{\color{var}n}}}/2-1/2)) \\
  &\le \ensuremath{\mathtt{{\color{var}n}}}-1 + \ensuremath{\mathtt{{\color{var}n}}}\log(\ensuremath{\mathtt{{\color{var}n}}}/2) + 1/2 \\
  &< \ensuremath{\mathtt{{\color{var}n}}} + \ensuremath{\mathtt{{\color{var}n}}}\log(\ensuremath{\mathtt{{\color{var}n}}}/2) \\
  &= \ensuremath{\mathtt{{\color{var}n}}} + \ensuremath{\mathtt{{\color{var}n}}}(\log\ensuremath{\mathtt{{\color{var}n}}}-1) \\
  &= \ensuremath{\mathtt{{\color{var}n}}}\log\ensuremath{\mathtt{{\color{var}n}}} \enspace . \qedhere
\end{align*} 
\end{proof}

\subsection{Quicksort}

\index{quicksort}%
The \emph{quicksort} algorithm is another classic divide and conquer
algorithm.  Unlike merge-sort, which does merging after solving the two
subproblems, quicksort does all of its work upfront.

Quicksort is simple to describe:  Pick a random \emph{pivot} element,
\index{pivot element}%
\ensuremath{\mathtt{{\color{var}x}}}, from \ensuremath{\mathtt{{\color{var}a}}}; partition \ensuremath{\mathtt{{\color{var}a}}} into the set of elements less than \ensuremath{\mathtt{{\color{var}x}}}, the
set of elements equal to \ensuremath{\mathtt{{\color{var}x}}}, and the set of elements greater than \ensuremath{\mathtt{{\color{var}x}}};
and, finally, recursively sort the first and third sets in this partition.
An example is shown in \figref{quicksort}.


 \ \vspace{.5ex}\begin{Verbatim}[gobble=1,tabsize=2,frame=single,framerule=1.2pt,commandchars=\\@\$,label=\texttt{Algorithms},labelposition=topline]
	<@\color@keyword$T$> @\color@keyword$void$ quickSort(@\color@keyword$T$[] @\color@var$a$, Comparator<@\color@keyword$T$> @\color@var$c$) {
		quickSort(@\color@var$a$, 0, @\color@var$a$.@\color@var$length$, @\color@var$c$);
	}
	<@\color@keyword$T$> @\color@keyword$void$ quickSort(@\color@keyword$T$[] @\color@var$a$, @\color@keyword$int$ @\color@var$i$, @\color@keyword$int$ @\color@var$n$, Comparator<@\color@keyword$T$> @\color@var$c$) {
		@\color@keyword$if$ (@\color@var$n$ <= 1) @\color@var$return$;
		@\color@keyword$T$ @\color@var$x$ = @\color@var$a$[@\color@var$i$ + @\color@var$rand$.nextInt(@\color@var$n$)];
		@\color@keyword$int$ @\color@var$p$ = @\color@var$i$-1, @\color@var$j$ = @\color@var$i$, @\color@var$q$ = @\color@var$i$+@\color@var$n$;
		@\color@comment$// a[i..p]<x,  a[p+1..q-1]??x, a[q..i+n-1]>x $
		@\color@keyword$while$ (@\color@var$j$ < @\color@var$q$) {
			@\color@keyword$int$ @\color@var$comp$ = compare(@\color@var$a$[@\color@var$j$], @\color@var$x$);
			@\color@keyword$if$ (@\color@var$comp$ < 0) {       @\color@comment$// move to beginning of array$
				swap(@\color@var$a$, @\color@var$j$++, ++@\color@var$p$);
			} @\color@keyword$else$ @\color@keyword$if$ (@\color@var$comp$ > 0) {
				swap(@\color@var$a$, @\color@var$j$, --@\color@var$q$);  @\color@comment$// move to end of array$
			} @\color@keyword$else$ {
				@\color@var$j$++;              @\color@comment$// keep in the middle$
			}
		}
		@\color@comment$// a[i..p]<x,  a[p+1..q-1]=x, a[q..i+n-1]>x $
		quickSort(@\color@var$a$, @\color@var$i$, @\color@var$p$-@\color@var$i$+1, @\color@var$c$);
		quickSort(@\color@var$a$, @\color@var$q$, @\color@var$n$-(@\color@var$q$-@\color@var$i$), @\color@var$c$);
	}
\end{Verbatim}

\pcodeimport{ods/Algorithms.quickSort(a).quickSort(a,i,n)}
\begin{figure}
  \begin{center}
    \includegraphics[scale=0.90909]{figs/quicksort}
    \caption[Quicksort]{An example execution of \javaonly{\ensuremath{\mathtt{quickSort({\color{var}a},0,14,{\color{var}c})}}}\cpponly{\ensuremath{\mathtt{quickSort({\color{var}a},0,14)}}}\pcodeonly{\ensuremath{\mathtt{quickSort({\color{var}a},0,14)}}}}
    \figlabel{quicksort}
  \end{center}
\end{figure}
All of this is done in place, so that instead of making copies of
subarrays being sorted, the \ensuremath{\mathtt{quickSort({\color{var}a},{\color{var}i},{\color{var}n},{\color{var}c})}} method only sorts the
subarray $\ensuremath{\mathtt{{\color{var}a}[{\color{var}i}]}},\ldots,\ensuremath{\mathtt{{\color{var}a}[{\color{var}i}+{\color{var}n}-1]}}$.  Initially, this method is invoked
with the arguments
\ensuremath{\mathtt{quickSort({\color{var}a},0,{\color{var}a}.{\color{var}length},{\color{var}c})}}.

At the heart of the quicksort algorithm is the in-place partitioning
algorithm.  This algorithm, without using any extra space, swaps elements
in \ensuremath{\mathtt{{\color{var}a}}} and computes indices \ensuremath{\mathtt{{\color{var}p}}} and \ensuremath{\mathtt{{\color{var}q}}} so that
\[
   \ensuremath{\mathtt{{\color{var}a}[{\color{var}i}]}} \begin{cases} 
         {}< \ensuremath{\mathtt{{\color{var}x}}} & \text{if $0\le \ensuremath{\mathtt{{\color{var}i}}}\le \ensuremath{\mathtt{{\color{var}p}}}$} \\
         {}= \ensuremath{\mathtt{{\color{var}x}}} & \text{if $\ensuremath{\mathtt{{\color{var}p}}}< \ensuremath{\mathtt{{\color{var}i}}} < \ensuremath{\mathtt{{\color{var}q}}}$} \\
         {}> \ensuremath{\mathtt{{\color{var}x}}} & \text{if $\ensuremath{\mathtt{{\color{var}q}}}\le \ensuremath{\mathtt{{\color{var}i}}} \le \ensuremath{\mathtt{{\color{var}n}}}-1$}
     \end{cases}
\]
This partitioning, which is done by the \ensuremath{\mathtt{{\color{keyword}while}}} loop in the code, works
by iteratively increasing \ensuremath{\mathtt{{\color{var}p}}} and decreasing \ensuremath{\mathtt{{\color{var}q}}} while maintaining the
first and last of these conditions.  At each step, the element at position
\ensuremath{\mathtt{{\color{var}j}}} is either moved to the front, left where it is, or moved to the back.
In the first two cases, \ensuremath{\mathtt{{\color{var}j}}} is incremented, while in the last case, \ensuremath{\mathtt{{\color{var}j}}}
is not incremented since the new element at position \ensuremath{\mathtt{{\color{var}j}}} has not yet been
processed.

Quicksort is very closely related to the random binary search trees
studied in \secref{rbst}.  In fact, if the input to quicksort consists
of \ensuremath{\mathtt{{\color{var}n}}} distinct elements, then the quicksort recursion tree is a random
binary search tree.  To see this, recall that when constructing a random
binary search tree the first thing we do is pick a random element \ensuremath{\mathtt{{\color{var}x}}} and
make it the root of the tree.  After this, every element will eventually
be compared to \ensuremath{\mathtt{{\color{var}x}}}, with smaller elements going into the left subtree
and larger elements into the right.

In quicksort, we select a random element \ensuremath{\mathtt{{\color{var}x}}} and immediately compare
everything to \ensuremath{\mathtt{{\color{var}x}}}, putting the smaller elements at the beginning of
the array and larger elements at the end of the array.  Quicksort then
recursively sorts the beginning of the array and the end of the array,
while the random binary search tree recursively inserts smaller elements
in the left subtree of the root and larger elements in the right subtree
of the root.

The above correspondence between random binary search trees and quicksort
means that we can translate \lemref{rbs} to a statement about quicksort:

\begin{lem}\lemlabel{quicksort}
  When quicksort is called to sort an array containing the integers
  $0,\ldots,\ensuremath{\mathtt{{\color{var}n}}}-1$, the expected number of times element \ensuremath{\mathtt{{\color{var}i}}} is compared
  to a pivot element is at most $H_{\ensuremath{\mathtt{{\color{var}i}}}+1} + H_{\ensuremath{\mathtt{{\color{var}n}}}-\ensuremath{\mathtt{{\color{var}i}}}}$.
\end{lem}

A little summing up of harmonic numbers gives us the following theorem
about the running time of quicksort:

\begin{thm}\thmlabel{quicksort-i}
  When quicksort is called to sort an array containing \ensuremath{\mathtt{{\color{var}n}}} distinct
  elements, the expected number of comparisons performed is at most
  $2\ensuremath{\mathtt{{\color{var}n}}}\ln \ensuremath{\mathtt{{\color{var}n}}} + O(\ensuremath{\mathtt{{\color{var}n}}})$.
\end{thm}

\begin{proof}
Let $T$ be the number of comparisons performed by quicksort when sorting
\ensuremath{\mathtt{{\color{var}n}}} distinct elements.  Using \lemref{quicksort} and linearity of
expectation, we have:
\begin{align*}
  \E[T] &= \sum_{i=0}^{\ensuremath{\mathtt{{\color{var}n}}}-1}(H_{\ensuremath{\mathtt{{\color{var}i}}}+1}+H_{\ensuremath{\mathtt{{\color{var}n}}}-\ensuremath{\mathtt{{\color{var}i}}}}) \\ 
        &= 2\sum_{i=1}^{\ensuremath{\mathtt{{\color{var}n}}}}H_i \\ 
        &\le 2\sum_{i=1}^{\ensuremath{\mathtt{{\color{var}n}}}}H_{\ensuremath{\mathtt{{\color{var}n}}}} \\ 
        &\le 2\ensuremath{\mathtt{{\color{var}n}}}\ln\ensuremath{\mathtt{{\color{var}n}}} + 2\ensuremath{\mathtt{{\color{var}n}}} = 2\ensuremath{\mathtt{{\color{var}n}}}\ln \ensuremath{\mathtt{{\color{var}n}}} + O(\ensuremath{\mathtt{{\color{var}n}}}) \qedhere
\end{align*}
\end{proof}

\thmref{quicksort} describes the case where the elements being sorted are
all distinct.  When the input array, \ensuremath{\mathtt{{\color{var}a}}}, contains duplicate elements,
the expected running time of quicksort is no worse, and can be even
better; any time a duplicate element \ensuremath{\mathtt{{\color{var}x}}} is chosen as a pivot, all
occurrences of \ensuremath{\mathtt{{\color{var}x}}} get grouped together and do not take part in either
of the two subproblems.

\begin{thm}\thmlabel{quicksort}
  The \javaonly{\ensuremath{\mathtt{quickSort({\color{var}a},{\color{var}c})}}} \cpponly{\ensuremath{\mathtt{quickSort({\color{var}a})}}}
  \pcodeonly{\ensuremath{\mathtt{quickSort({\color{var}a},{\color{var}c})}}} method runs in $O(\ensuremath{\mathtt{{\color{var}n}}}\log \ensuremath{\mathtt{{\color{var}n}}})$ expected
  time and the expected number of comparisons it performs is at most
  $2\ensuremath{\mathtt{{\color{var}n}}}\ln \ensuremath{\mathtt{{\color{var}n}}} +O(\ensuremath{\mathtt{{\color{var}n}}})$.
\end{thm}

\subsection{Heap-sort}
\seclabel{heapsort}

\index{heap-sort}%
The \emph{heap-sort} algorithm is another in-place sorting algorithm.
Heap-sort uses the binary heaps discussed in \secref{binaryheap}.
Recall that the \texttt{BinaryHeap} data structure represents a heap using
a single array.  The heap-sort algorithm converts the input array \ensuremath{\mathtt{{\color{var}a}}}
into a heap and then repeatedly extracts the minimum value.

More specifically, a heap stores \ensuremath{\mathtt{{\color{var}n}}} elements in an array, \ensuremath{\mathtt{{\color{var}a}}}, at array locations
$\ensuremath{\mathtt{{\color{var}a}[0]}},\ldots,\ensuremath{\mathtt{{\color{var}a}[{\color{var}n}-1]}}$ with the smallest value stored at the root,
\ensuremath{\mathtt{{\color{var}a}[0]}}.  After transforming \ensuremath{\mathtt{{\color{var}a}}} into a \texttt{BinaryHeap}, the heap-sort
algorithm repeatedly swaps \ensuremath{\mathtt{{\color{var}a}[0]}} and \ensuremath{\mathtt{{\color{var}a}[{\color{var}n}-1]}}, decrements \ensuremath{\mathtt{{\color{var}n}}}, and
calls \ensuremath{\mathtt{trickleDown(0)}} so that $\ensuremath{\mathtt{{\color{var}a}[0]}},\ldots,\ensuremath{\mathtt{{\color{var}a}[{\color{var}n}-2]}}$ once again are
a valid heap representation. When this process ends (because $\ensuremath{\mathtt{{\color{var}n}}}=0$)
the elements of \ensuremath{\mathtt{{\color{var}a}}} are stored in decreasing order, so \ensuremath{\mathtt{{\color{var}a}}} is reversed
to obtain the final sorted order.\footnote{The algorithm
could alternatively redefine the \ensuremath{\mathtt{compare({\color{var}x},{\color{var}y})}} function so that the
heap sort algorithm stores the elements directly in ascending order.}
\figref{heapsort} shows an example of the execution of \ensuremath{\mathtt{heapSort({\color{var}a},{\color{var}c})}}.

\begin{figure}
  \begin{center}
    \includegraphics[scale=0.90909]{figs/heapsort}
  \end{center}
  \caption[Heap sort]{A snapshot of the execution of \ensuremath{\mathtt{heapSort({\color{var}a},{\color{var}c})}}.
      The shaded part of the
      array is already sorted.  The unshaded part is a \texttt{BinaryHeap}.
      During the next iteration, element $5$ will be placed into array
      location $8$.}
  \figlabel{heapsort}
\end{figure}



 \ \vspace{.5ex}\begin{Verbatim}[gobble=1,tabsize=2,frame=single,framerule=1.2pt,commandchars=\\@\$,label=\texttt{BinaryHeap},labelposition=topline]
	<@\color@keyword$T$> @\color@keyword$void$ sort(@\color@keyword$T$[] @\color@var$a$, Comparator<@\color@keyword$T$> @\color@var$c$) {
		BinaryHeap<@\color@keyword$T$> @\color@var$h$ = @\color@keyword$new$ BinaryHeap<@\color@keyword$T$>(@\color@var$a$, @\color@var$c$);
		@\color@keyword$while$ (@\color@var$h$.@\color@var$n$ > 1) {
			@\color@var$h$.swap(--@\color@var$h$.@\color@var$n$, 0);
			@\color@var$h$.trickleDown(0);
		}
		Collections.reverse(Arrays.asList(@\color@var$a$));
	}
\end{Verbatim}

\pcodeimport{ods/Algorithms.heapSort(a)}

A key subroutine in heap sort is the constructor for turning
an unsorted array \ensuremath{\mathtt{{\color{var}a}}} into a heap.  It would be easy to do this
in $O(\ensuremath{\mathtt{{\color{var}n}}}\log\ensuremath{\mathtt{{\color{var}n}}})$ time by repeatedly calling the \texttt{BinaryHeap}
\ensuremath{\mathtt{add({\color{var}x})}} method, but we can do better by using a bottom-up algorithm.
Recall that, in a binary heap, the children of \ensuremath{\mathtt{{\color{var}a}[{\color{var}i}]}} are stored at
positions \ensuremath{\mathtt{{\color{var}a}[2i+1]}} and \ensuremath{\mathtt{{\color{var}a}[2i+2]}}.  This implies that the elements
$\ensuremath{\mathtt{{\color{var}a}}}[\lfloor\ensuremath{\mathtt{{\color{var}n}}}/2\rfloor],\ldots,\ensuremath{\mathtt{{\color{var}a}[{\color{var}n}-1]}}$ have no children. In other
words, each of $\ensuremath{\mathtt{{\color{var}a}}}[\lfloor\ensuremath{\mathtt{{\color{var}n}}}/2\rfloor],\ldots,\ensuremath{\mathtt{{\color{var}a}[{\color{var}n}-1]}}$ is a sub-heap
of size 1.  Now, working backwards, we can call \ensuremath{\mathtt{trickleDown({\color{var}i})}} for
each $\ensuremath{\mathtt{{\color{var}i}}}\in\{\lfloor \ensuremath{\mathtt{{\color{var}n}}}/2\rfloor-1,\ldots,0\}$. This works, because by
the time we call \ensuremath{\mathtt{trickleDown({\color{var}i})}}, each of the two children of \ensuremath{\mathtt{{\color{var}a}[{\color{var}i}]}}
are the root of a sub-heap, so calling \ensuremath{\mathtt{trickleDown({\color{var}i})}} makes \ensuremath{\mathtt{{\color{var}a}[{\color{var}i}]}}
into the root of its own subheap.


 \ \vspace{.5ex}\begin{Verbatim}[gobble=1,tabsize=2,frame=single,framerule=1.2pt,commandchars=\\@\$,label=\texttt{BinaryHeap},labelposition=topline]
	BinaryHeap(@\color@keyword$T$[] @\color@var$a$, Comparator<@\color@keyword$T$> @\color@var$c$) {
		@\color@var$this$.@\color@var$c$ = @\color@var$c$;
		@\color@var$this$.@\color@var$a$ = @\color@var$a$;
		@\color@var$n$ = @\color@var$a$.@\color@var$length$;
		@\color@keyword$for$ (@\color@keyword$int$ @\color@var$i$ = @\color@var$n$/2-1; @\color@var$i$ >= 0; @\color@var$i$--) {
			trickleDown(@\color@var$i$);
		}
	}
\end{Verbatim}


The interesting thing about this bottom-up strategy is that it is more
efficient than calling \ensuremath{\mathtt{add({\color{var}x})}} \ensuremath{\mathtt{{\color{var}n}}} times.  To see this, notice that,
for $\ensuremath{\mathtt{{\color{var}n}}}/2$ elements, we do no work at all, for $\ensuremath{\mathtt{{\color{var}n}}}/4$ elements, we call
\ensuremath{\mathtt{trickleDown({\color{var}i})}} on a subheap rooted at \ensuremath{\mathtt{{\color{var}a}[{\color{var}i}]}} and whose height is one, for
$\ensuremath{\mathtt{{\color{var}n}}}/8$ elements, we call \ensuremath{\mathtt{trickleDown({\color{var}i})}} on a subheap whose height is two,
and so on.  Since the work done by \ensuremath{\mathtt{trickleDown({\color{var}i})}} is proportional to
the height of the sub-heap rooted at \ensuremath{\mathtt{{\color{var}a}[{\color{var}i}]}}, this means that the total
work done is at most
\[
    \sum_{i=1}^{\log\ensuremath{\mathtt{{\color{var}n}}}} O((i-1)\ensuremath{\mathtt{{\color{var}n}}}/2^{i})
    \le \sum_{i=1}^{\infty} O(i\ensuremath{\mathtt{{\color{var}n}}}/2^{i})
    = O(\ensuremath{\mathtt{{\color{var}n}}})\sum_{i=1}^{\infty} i/2^{i}
    =  O(2\ensuremath{\mathtt{{\color{var}n}}}) = O(\ensuremath{\mathtt{{\color{var}n}}}) \enspace .
\]
The second-last equality follows by recognizing that the sum
$\sum_{i=1}^{\infty} i/2^{i}$ is equal, by definition of expected value,
to the expected number of times we toss a coin up to and including the
first time the coin comes up as heads and applying \lemref{coin-tosses}.

The following theorem describes the performance of \ensuremath{\mathtt{heapSort({\color{var}a},{\color{var}c})}}.
\begin{thm}
  The \ensuremath{\mathtt{heapSort({\color{var}a},{\color{var}c})}} method runs in $O(\ensuremath{\mathtt{{\color{var}n}}}\log \ensuremath{\mathtt{{\color{var}n}}})$ time and performs at
  most $2\ensuremath{\mathtt{{\color{var}n}}}\log \ensuremath{\mathtt{{\color{var}n}}} + O(\ensuremath{\mathtt{{\color{var}n}}})$ comparisons.
\end{thm}

\begin{proof}
The algorithm runs in three steps:  (1)~transforming \ensuremath{\mathtt{{\color{var}a}}} into a heap,
(2)~repeatedly extracting the minimum element from \ensuremath{\mathtt{{\color{var}a}}}, and (3)~reversing
the elements in \ensuremath{\mathtt{{\color{var}a}}}.  We have just argued that step~1 takes $O(\ensuremath{\mathtt{{\color{var}n}}})$
time and performs $O(\ensuremath{\mathtt{{\color{var}n}}})$ comparisons.  Step~3 takes $O(\ensuremath{\mathtt{{\color{var}n}}})$ time and
performs no comparisons.  Step~2 performs \ensuremath{\mathtt{{\color{var}n}}} calls to \ensuremath{\mathtt{trickleDown(0)}}.
The $i$th such call operates on a heap of size $\ensuremath{\mathtt{{\color{var}n}}}-i$ and performs
at most $2\log(\ensuremath{\mathtt{{\color{var}n}}}-i)$ comparisons.  Summing this over $i$ gives
\[
   \sum_{i=0}^{\ensuremath{\mathtt{{\color{var}n}}}-i} 2\log(\ensuremath{\mathtt{{\color{var}n}}}-i) 
   \le \sum_{i=0}^{\ensuremath{\mathtt{{\color{var}n}}}-i} 2\log \ensuremath{\mathtt{{\color{var}n}}}
   =  2\ensuremath{\mathtt{{\color{var}n}}}\log \ensuremath{\mathtt{{\color{var}n}}}
\]
Adding the number of comparisons performed in each of the three steps
completes the proof.
\end{proof}

\subsection{A Lower-Bound for Comparison-Based Sorting}

\index{lower-bound}%
\index{sorting lower-bound}%
We have now seen three comparison-based sorting algorithms that each run
in $O(\ensuremath{\mathtt{{\color{var}n}}}\log \ensuremath{\mathtt{{\color{var}n}}})$ time.  By now, we should be wondering if faster
algorithms exist.  The short answer to this question is no.  If the
only operations allowed on the elements of \ensuremath{\mathtt{{\color{var}a}}} are comparisons, then no
algorithm can avoid doing roughly $\ensuremath{\mathtt{{\color{var}n}}}\log \ensuremath{\mathtt{{\color{var}n}}}$ comparisons.  This is
not difficult to prove, but requires a little imagination.  Ultimately,
it follows from the fact that
\[
   \log(\ensuremath{\mathtt{{\color{var}n}}}!) 
     = \log \ensuremath{\mathtt{{\color{var}n}}} + \log (\ensuremath{\mathtt{{\color{var}n}}}-1) + \dots + \log(1) 
     = \ensuremath{\mathtt{{\color{var}n}}}\log \ensuremath{\mathtt{{\color{var}n}}} - O(\ensuremath{\mathtt{{\color{var}n}}})
    \enspace .
\]
(Proving this fact is left as \excref{log-factorial}.)

We will start by focusing our attention on deterministic algorithms like
merge-sort and heap-sort and on a particular fixed value of \ensuremath{\mathtt{{\color{var}n}}}.  Imagine
such an algorithm is being used to sort \ensuremath{\mathtt{{\color{var}n}}} distinct elements.  The key
to proving the lower-bound is to observe that, for a deterministic
algorithm with a fixed value of \ensuremath{\mathtt{{\color{var}n}}}, the first pair of elements that are
compared is always the same.  For example, in \ensuremath{\mathtt{heapSort({\color{var}a},{\color{var}c})}}, when \ensuremath{\mathtt{{\color{var}n}}}
is even, the first call to \ensuremath{\mathtt{trickleDown({\color{var}i})}} is with \ensuremath{\mathtt{{\color{var}i}={\color{var}n}/2-1}} and the
first comparison is between elements \ensuremath{\mathtt{{\color{var}a}[{\color{var}n}/2-1]}} and \ensuremath{\mathtt{{\color{var}a}[{\color{var}n}-1]}}.

Since all input elements are distinct, this first comparison has only
two possible outcomes.  The second comparison done by the algorithm may
depend on the outcome of the first comparison.  The third comparison
may depend on the results of the first two, and so on.  In this way,
any deterministic comparison-based sorting algorithm can be viewed
as a rooted binary \emph{comparison tree}.
\index{comparison tree}%
Each internal node, \ensuremath{\mathtt{{\color{var}u}}},
of this tree is labelled with a pair of indices \ensuremath{\mathtt{{\color{var}u}.{\color{var}i}}} and \ensuremath{\mathtt{{\color{var}u}.{\color{var}j}}}.
If $\ensuremath{\mathtt{{\color{var}a}[{\color{var}u}.{\color{var}i}]}}<\ensuremath{\mathtt{{\color{var}a}[{\color{var}u}.{\color{var}j}]}}$ the algorithm proceeds to the left subtree,
otherwise it proceeds to the right subtree.  Each leaf \ensuremath{\mathtt{{\color{var}w}}} of this
tree is labelled with a permutation $\ensuremath{\mathtt{{\color{var}w}.{\color{var}p}[0]}},\ldots,\ensuremath{\mathtt{{\color{var}w}.{\color{var}p}[{\color{var}n}-1]}}$ of
$0,\ldots,\ensuremath{\mathtt{{\color{var}n}}}-1$.  This permutation represents the one that is
required to sort \ensuremath{\mathtt{{\color{var}a}}} if the comparison tree reaches this leaf.  That is,
\[
   \ensuremath{\mathtt{{\color{var}a}[{\color{var}w}.{\color{var}p}[0]]}}<\ensuremath{\mathtt{{\color{var}a}[{\color{var}w}.{\color{var}p}[1]]}}<\cdots<\ensuremath{\mathtt{{\color{var}a}[{\color{var}w}.{\color{var}p}[{\color{var}n}-1]]}} \enspace .
\]
An example of a comparison tree for an array of size \ensuremath{\mathtt{{\color{var}n}=3}} is shown in
\figref{comparison-tree}.
\begin{figure}
  \begin{center}
    \includegraphics[width=\ScaleIfNeeded]{figs/comparison-tree}
  \end{center}
  \caption[A comparison tree]{A comparison tree for sorting an array $\ensuremath{\mathtt{{\color{var}a}[0]}},\ensuremath{\mathtt{{\color{var}a}[1]}},\ensuremath{\mathtt{{\color{var}a}[2]}}$ of length \ensuremath{\mathtt{{\color{var}n}=3}}.}
  \figlabel{comparison-tree}
\end{figure}

The comparison tree for a sorting algorithm tells us everything about
the algorithm.  It tells us exactly the sequence of comparisons that
will be performed for any input array, \ensuremath{\mathtt{{\color{var}a}}}, having \ensuremath{\mathtt{{\color{var}n}}} distinct elements
and it tells us how the algorithm will reorder \ensuremath{\mathtt{{\color{var}a}}} in order to sort it.
Consequently, the comparison tree must have at least $\ensuremath{\mathtt{{\color{var}n}}}!$ leaves;
if not, then there are two distinct permutations that lead to the same
leaf; therefore, the algorithm does not correctly sort at least one of
these permutations.

For example, the comparison tree in \figref{comparison-tree-2} has only
$4< 3!=6$ leaves. Inspecting this tree, we see that the two input arrays
$3,1,2$ and $3,2,1$ both lead to the rightmost leaf.  On the input $3,1,2$
this leaf correctly outputs $\ensuremath{\mathtt{{\color{var}a}[1]}}=1,\ensuremath{\mathtt{{\color{var}a}[2]}}=2,\ensuremath{\mathtt{{\color{var}a}[0]}}=3$.  However, on the
input $3,2,1$, this node incorrectly outputs $\ensuremath{\mathtt{{\color{var}a}[1]}}=2,\ensuremath{\mathtt{{\color{var}a}[2]}}=1,\ensuremath{\mathtt{{\color{var}a}[0]}}=3$.
This discussion leads to the primary lower-bound for comparison-based
algorithms.

\begin{figure}
  \begin{center}
    \includegraphics[width=\ScaleIfNeeded]{figs/comparison-tree-b}
  \end{center}
  \caption{A comparison tree that does not correctly sort every input
  permutation.}
  \figlabel{comparison-tree-2}
\end{figure}

\begin{thm}\thmlabel{deterministic-sorting-lower-bound}
  For any deterministic comparison-based sorting algorithm $\mathcal{A}$
  and any integer $\ensuremath{\mathtt{{\color{var}n}}}\ge 1$, there exists an input array \ensuremath{\mathtt{{\color{var}a}}} of
  length \ensuremath{\mathtt{{\color{var}n}}} such that $\mathcal{A}$ performs at least $\log(\ensuremath{\mathtt{{\color{var}n}}}!) =
  \ensuremath{\mathtt{{\color{var}n}}}\log\ensuremath{\mathtt{{\color{var}n}}}-O(\ensuremath{\mathtt{{\color{var}n}}})$ comparisons when sorting \ensuremath{\mathtt{{\color{var}a}}}.
\end{thm}

\begin{proof}
  By the preceding discussion, the comparison tree defined by $\mathcal{A}$
  must have at least $\ensuremath{\mathtt{{\color{var}n}}}!$ leaves.  An easy inductive proof shows that
  any binary tree with $k$ leaves has a height of at least $\log k$.
  Therefore, the comparison tree for $\mathcal{A}$ has a leaf, \ensuremath{\mathtt{{\color{var}w}}},
  with a depth of at least $\log(\ensuremath{\mathtt{{\color{var}n}}}!)$ and there is an input array \ensuremath{\mathtt{{\color{var}a}}}
  that leads to this leaf.  The input array \ensuremath{\mathtt{{\color{var}a}}} is an input for which
  $\mathcal{A}$ does at least $\log(\ensuremath{\mathtt{{\color{var}n}}}!)$ comparisons.
\end{proof}

\thmref{deterministic-sorting-lower-bound} deals with deterministic
algorithms like merge-sort and heap-sort, but doesn't tell us anything
about randomized algorithms like quicksort.  Could a randomized algorithm
beat the $\log(\ensuremath{\mathtt{{\color{var}n}}}!)$ lower bound on the number of comparisons?
The answer, again, is no.  Again, the way to prove it is to think
differently about what a randomized algorithm is.

In the following discussion, we will assume that our decision
trees have been ``cleaned up'' in the following way: Any node that can not
be reached by some input array \ensuremath{\mathtt{{\color{var}a}}} is removed.  This cleaning up implies
that the tree has exactly $\ensuremath{\mathtt{{\color{var}n}}}!$ leaves.  It has at least $\ensuremath{\mathtt{{\color{var}n}}}!$ leaves
because, otherwise, it could not sort correctly.  It has at most $\ensuremath{\mathtt{{\color{var}n}}}!$
leaves since each of the possible $\ensuremath{\mathtt{{\color{var}n}}}!$ permutation of \ensuremath{\mathtt{{\color{var}n}}} distinct
elements follows exactly one root to leaf path in the decision tree.

We can think of a randomized sorting algorithm, $\mathcal{R}$, as a
deterministic algorithm that takes two inputs: The input array \ensuremath{\mathtt{{\color{var}a}}}
that should be sorted and a long sequence $b=b_1,b_2,b_3,\ldots,b_m$
of random real numbers in the range $[0,1]$.  The random numbers provide
the randomization for the algorithm.  When the algorithm wants to toss a
coin or make a random choice, it does so by using some element from $b$.
For example, to compute the index of the first pivot in quicksort,
the algorithm could use the formula $\lfloor n b_1\rfloor$.

Now, notice that if we fix $b$ to some particular sequence $\hat{b}$
then $\mathcal{R}$ becomes a deterministic sorting algorithm,
$\mathcal{R}(\hat{b})$, that has an associated comparison tree,
$\mathcal{T}(\hat{b})$.  Next, notice that if we select \ensuremath{\mathtt{{\color{var}a}}} to be a random
permutation of $\{1,\ldots,\ensuremath{\mathtt{{\color{var}n}}}\}$, then this is equivalent to selecting
a random leaf, \ensuremath{\mathtt{{\color{var}w}}}, from the $\ensuremath{\mathtt{{\color{var}n}}}!$ leaves of $\mathcal{T}(\hat{b})$.

\excref{randomized-lower-bound} asks you to prove that, if we select
a random leaf from any binary tree with $k$ leaves, then the expected
depth of that leaf is at least $\log k$.  Therefore, the expected
number of comparisons performed by the (deterministic) algorithm
$\mathcal{R}(\hat{b})$ when given an input array containing a random
permutation of $\{1,\ldots,n\}$ is at least $\log(\ensuremath{\mathtt{{\color{var}n}}}!)$.  Finally,
notice that this is true for every choice of $\hat{b}$, therefore it holds even for $\mathcal{R}$.  This completes the proof of the lower-bound for randomized algorithms.

\begin{thm}\thmlabel{randomized-sorting-lower-bound}
  For any integer $n\ge 1$ and any  (deterministic or randomized)
  comparison-based sorting algorithm $\mathcal{A}$, the expected number
  of comparisons done by $\mathcal{A}$ when sorting a random permutation
  of $\{1,\ldots,n\}$ is at least $\log(\ensuremath{\mathtt{{\color{var}n}}}!) = \ensuremath{\mathtt{{\color{var}n}}}\log\ensuremath{\mathtt{{\color{var}n}}}-O(\ensuremath{\mathtt{{\color{var}n}}})$.
\end{thm}



\section{Counting Sort and Radix Sort}

In this section we study two sorting algorithms that are not
comparison-based. Specialized for sorting small integers, these algorithms
elude the lower-bounds of \thmref{deterministic-sorting-lower-bound}
by using (parts of) the elements in \ensuremath{\mathtt{{\color{var}a}}} as indices into an array.
Consider a statement of the form
\[
  \ensuremath{\mathtt{{\color{var}c}[{\color{var}a}[{\color{var}i}]]}} = 1 \enspace .
\]
This statement executes in constant time, but has \ensuremath{\mathtt{{\color{var}c}.{\color{var}length}}} possible
different outcomes, depending on the value of \ensuremath{\mathtt{{\color{var}a}[{\color{var}i}]}}.  This means that the
execution of an algorithm that makes such a statement cannot be modelled
as a binary tree.  Ultimately, this is the reason that the algorithms
in this section are able to sort faster than comparison-based algorithms.

\subsection{Counting Sort}

Suppose we have an input array \ensuremath{\mathtt{{\color{var}a}}} consisting of \ensuremath{\mathtt{{\color{var}n}}} integers, each in
the range $0,\ldots,\ensuremath{\mathtt{{\color{var}k}}}-1$.  The \emph{counting-sort}
\index{counting-sort}%
algorithm sorts \ensuremath{\mathtt{{\color{var}a}}}
using an auxiliary array \ensuremath{\mathtt{{\color{var}c}}} of counters.  It outputs a sorted version
of \ensuremath{\mathtt{{\color{var}a}}} as an auxiliary array \ensuremath{\mathtt{{\color{var}b}}}.

The idea behind counting-sort is simple:  For each
$\ensuremath{\mathtt{{\color{var}i}}}\in\{0,\ldots,\ensuremath{\mathtt{{\color{var}k}}}-1\}$, count the number of occurrences of \ensuremath{\mathtt{{\color{var}i}}} in \ensuremath{\mathtt{{\color{var}a}}}
and store this in \ensuremath{\mathtt{{\color{var}c}[{\color{var}i}]}}.  Now, after sorting, the output will look like
\ensuremath{\mathtt{{\color{var}c}[0]}} occurrences of 0, followed by \ensuremath{\mathtt{{\color{var}c}[1]}} occurrences of 1, followed by
\ensuremath{\mathtt{{\color{var}c}[2]}} occurrences of 2,\ldots, followed by \ensuremath{\mathtt{{\color{var}c}[{\color{var}k}-1]}} occurrences of \ensuremath{\mathtt{{\color{var}k}-1}}.
The code that does this is very slick, and its execution is illustrated in
\figref{countingsort}:


 \ \vspace{.5ex}\begin{Verbatim}[gobble=1,tabsize=2,frame=single,framerule=1.2pt,commandchars=\\@\$,label=\texttt{Algorithms},labelposition=topline]
	@\color@keyword$int$[] countingSort(@\color@keyword$int$[] @\color@var$a$, @\color@keyword$int$ @\color@var$k$) {
		@\color@keyword$int$ @\color@var$c$[] = @\color@keyword$new$ @\color@keyword$int$[@\color@var$k$];
		@\color@keyword$for$ (@\color@keyword$int$ @\color@var$i$ = 0; @\color@var$i$ < @\color@var$a$.@\color@var$length$; @\color@var$i$++)
			@\color@var$c$[@\color@var$a$[@\color@var$i$]]++;
		@\color@keyword$for$ (@\color@keyword$int$ @\color@var$i$ = 1; @\color@var$i$ < @\color@var$k$; @\color@var$i$++)
			@\color@var$c$[@\color@var$i$] += @\color@var$c$[@\color@var$i$-1];
		@\color@keyword$int$ @\color@var$b$[] = @\color@keyword$new$ @\color@keyword$int$[@\color@var$a$.@\color@var$length$];
		@\color@keyword$for$ (@\color@keyword$int$ @\color@var$i$ = @\color@var$a$.@\color@var$length$-1; @\color@var$i$ >= 0; @\color@var$i$--)
			@\color@var$b$[--@\color@var$c$[@\color@var$a$[@\color@var$i$]]] = @\color@var$a$[@\color@var$i$];
		@\color@var$return$ @\color@var$b$;
	}
\end{Verbatim}


\begin{figure}
  \begin{center}
    \includegraphics[width=\ScaleIfNeeded]{figs/countingsort}
  \end{center}
  \caption{The operation of counting sort on an array of length $\ensuremath{\mathtt{{\color{var}n}}}=20$ that stores integers $0,\ldots,\ensuremath{\mathtt{{\color{var}k}}}-1=9$.}
  \figlabel{countingsort}
\end{figure}

The first \ensuremath{\mathtt{{\color{keyword}for}}} loop in this code sets each counter \ensuremath{\mathtt{{\color{var}c}[{\color{var}i}]}} so that it
counts the number of occurrences of \ensuremath{\mathtt{{\color{var}i}}} in \ensuremath{\mathtt{{\color{var}a}}}.  By using the values
of \ensuremath{\mathtt{{\color{var}a}}} as indices, these counters can all be computed in $O(\ensuremath{\mathtt{{\color{var}n}}})$ time
with a single for loop.  At this point, we could use \ensuremath{\mathtt{{\color{var}c}}} to
fill in the output array \ensuremath{\mathtt{{\color{var}b}}} directly. However, this would not work if
the elements of \ensuremath{\mathtt{{\color{var}a}}} have associated data.  Therefore we spend a little
extra effort to copy the elements of \ensuremath{\mathtt{{\color{var}a}}} into \ensuremath{\mathtt{{\color{var}b}}}.

The next \ensuremath{\mathtt{{\color{keyword}for}}} loop, which takes $O(\ensuremath{\mathtt{{\color{var}k}}})$ time, computes a running-sum
of the counters so that \ensuremath{\mathtt{{\color{var}c}[{\color{var}i}]}} becomes the number of elements in
\ensuremath{\mathtt{{\color{var}a}}} that are less than or equal to \ensuremath{\mathtt{{\color{var}i}}}.  In particular, for every
$\ensuremath{\mathtt{{\color{var}i}}}\in\{0,\ldots,\ensuremath{\mathtt{{\color{var}k}}}-1\}$, the output array, \ensuremath{\mathtt{{\color{var}b}}}, will have
\[
   \ensuremath{\mathtt{{\color{var}b}[{\color{var}c}[{\color{var}i}-1]]}}=\ensuremath{\mathtt{{\color{var}b}[{\color{var}c}[{\color{var}i}-1]+1]=}}\cdots=\ensuremath{\mathtt{{\color{var}b}[{\color{var}c}[{\color{var}i}]-1]}}=\ensuremath{\mathtt{{\color{var}i}}} \enspace .
\]
Finally, the algorithm scans \ensuremath{\mathtt{{\color{var}a}}} backwards to place its elements, in order,
into an output array \ensuremath{\mathtt{{\color{var}b}}}.  When scanning, the element \ensuremath{\mathtt{{\color{var}a}[{\color{var}i}]={\color{var}j}}} is placed
at location \ensuremath{\mathtt{{\color{var}b}[{\color{var}c}[{\color{var}j}]-1]}} and the value \ensuremath{\mathtt{{\color{var}c}[{\color{var}j}]}} is decremented.

\begin{thm}
  The \ensuremath{\mathtt{countingSort({\color{var}a},{\color{var}k})}} method can sort an array \ensuremath{\mathtt{{\color{var}a}}} containing \ensuremath{\mathtt{{\color{var}n}}}
  integers in the set $\{0,\ldots,\ensuremath{\mathtt{{\color{var}k}}}-1\}$ in $O(\ensuremath{\mathtt{{\color{var}n}}}+\ensuremath{\mathtt{{\color{var}k}}})$ time.
\end{thm}

The counting-sort algorithm has the nice property of being \emph{stable};
\index{stable sorting algorithm}%
it preserves the relative order of equal elements.  If two elements
\ensuremath{\mathtt{{\color{var}a}[{\color{var}i}]}} and \ensuremath{\mathtt{{\color{var}a}[{\color{var}j}]}} have the same value, and $\ensuremath{\mathtt{{\color{var}i}}}<\ensuremath{\mathtt{{\color{var}j}}}$ then \ensuremath{\mathtt{{\color{var}a}[{\color{var}i}]}} will
appear before \ensuremath{\mathtt{{\color{var}a}[{\color{var}j}]}} in \ensuremath{\mathtt{{\color{var}b}}}.  This will be useful in the next section.

\subsection{Radix-Sort}

Counting-sort is very efficient for sorting an array of integers when the
length, \ensuremath{\mathtt{{\color{var}n}}}, of the array is not much smaller than the maximum value,
$\ensuremath{\mathtt{{\color{var}k}}}-1$, that appears in the array.  The \emph{radix-sort}
\index{radix-sort}%
algorithm,
which we now describe, uses several passes of counting-sort to allow
for a much greater range of maximum values.

Radix-sort sorts \ensuremath{\mathtt{{\color{var}w}}}-bit integers by using $\ensuremath{\mathtt{{\color{var}w}}}/\ensuremath{\mathtt{{\color{var}d}}}$ passes of counting-sort
to sort these integers \ensuremath{\mathtt{{\color{var}d}}} bits at a time.\footnote{We assume that
\ensuremath{\mathtt{{\color{var}d}}} divides \ensuremath{\mathtt{{\color{var}w}}}, otherwise we can always increase \ensuremath{\mathtt{{\color{var}w}}} to $\ensuremath{\mathtt{{\color{var}d}}}\lceil
\ensuremath{\mathtt{{\color{var}w}}}/\ensuremath{\mathtt{{\color{var}d}}}\rceil$.}  More precisely, radix sort first sorts the integers by
their least significant \ensuremath{\mathtt{{\color{var}d}}} bits, then their next significant \ensuremath{\mathtt{{\color{var}d}}} bits,
and so on until, in the last pass, the integers are sorted by their most
significant \ensuremath{\mathtt{{\color{var}d}}} bits.


 \ \vspace{.5ex}\begin{Verbatim}[gobble=1,tabsize=2,frame=single,framerule=1.2pt,commandchars=\\@\$,label=\texttt{Algorithms},labelposition=topline]
	@\color@keyword$int$[] radixSort(@\color@keyword$int$[] @\color@var$a$) {
		@\color@keyword$int$[] @\color@var$b$ = @\color@keyword$null$;
		@\color@keyword$for$ (@\color@keyword$int$ @\color@var$p$ = 0; @\color@var$p$ < @\color@var$w$/@\color@var$d$; @\color@var$p$++) {
			@\color@keyword$int$ @\color@var$c$[] = @\color@keyword$new$ @\color@keyword$int$[1<<@\color@var$d$];
			@\color@comment$// the next three for loops implement counting-sort$
			@\color@var$b$ = @\color@keyword$new$ @\color@keyword$int$[@\color@var$a$.@\color@var$length$];
			@\color@keyword$for$ (@\color@keyword$int$ @\color@var$i$ = 0; @\color@var$i$ < @\color@var$a$.@\color@var$length$; @\color@var$i$++)
				@\color@var$c$[(@\color@var$a$[@\color@var$i$] >> @\color@var$d$*@\color@var$p$)&((1<<@\color@var$d$)-1)]++;
			@\color@keyword$for$ (@\color@keyword$int$ @\color@var$i$ = 1; @\color@var$i$ < 1<<@\color@var$d$; @\color@var$i$++)
				@\color@var$c$[@\color@var$i$] += @\color@var$c$[@\color@var$i$-1];
			@\color@keyword$for$ (@\color@keyword$int$ @\color@var$i$ = @\color@var$a$.@\color@var$length$-1; @\color@var$i$ >= 0; @\color@var$i$--)
				@\color@var$b$[--@\color@var$c$[(@\color@var$a$[@\color@var$i$] >> @\color@var$d$*@\color@var$p$)&((1<<@\color@var$d$)-1)]] = @\color@var$a$[@\color@var$i$];
			@\color@var$a$ = @\color@var$b$;
		}
		@\color@var$return$ @\color@var$b$;
	}
\end{Verbatim}

(In this code, the expression \ensuremath{\mathtt{({\color{var}a}[{\color{var}i}]\text{\ttfamily >>}{\color{var}d}*{\color{var}p})\text{\ttfamily \&}((1\text{\ttfamily <<}{\color{var}d})-1)}} extracts the integer
whose binary representation is given by bits
$(\ensuremath{\mathtt{{\color{var}p}}}+1)\ensuremath{\mathtt{{\color{var}d}}}-1,\ldots,\ensuremath{\mathtt{{\color{var}p}}}\ensuremath{\mathtt{{\color{var}d}}}$ of \ensuremath{\mathtt{{\color{var}a}[{\color{var}i}]}}.)
An example of the steps of this algorithm is shown in \figref{radixsort}.

\begin{figure}
  \begin{center}
    \includegraphics[width=\ScaleIfNeeded]{figs/radixsort}
  \end{center}
  \caption{Using radixsort to sort $\ensuremath{\mathtt{{\color{var}w}}}=8$-bit integers by using 4 passes
   of counting sort on $\ensuremath{\mathtt{{\color{var}d}}}=2$-bit integers.}
  \figlabel{radixsort}
\end{figure}

This remarkable algorithm sorts correctly because counting-sort is
a stable sorting algorithm.  If $\ensuremath{\mathtt{{\color{var}x}}} < \ensuremath{\mathtt{{\color{var}y}}}$ are two elements of \ensuremath{\mathtt{{\color{var}a}}},
and the most significant bit at which \ensuremath{\mathtt{{\color{var}x}}} differs from \ensuremath{\mathtt{{\color{var}y}}} has index $r$,
then \ensuremath{\mathtt{{\color{var}x}}} will be placed before \ensuremath{\mathtt{{\color{var}y}}} during pass $\lfloor r/\ensuremath{\mathtt{{\color{var}d}}}\rfloor$
and subsequent passes will not change the relative order of \ensuremath{\mathtt{{\color{var}x}}} and \ensuremath{\mathtt{{\color{var}y}}}.

Radix-sort performs \ensuremath{\mathtt{{\color{var}w}/{\color{var}d}}} passes of counting-sort.  Each pass requires
$O(\ensuremath{\mathtt{{\color{var}n}}}+2^{\ensuremath{\mathtt{{\color{var}d}}}})$ time. Therefore, the performance of radix-sort is given
by the following theorem.
\begin{thm}\thmlabel{radix-sort}
  For any integer $\ensuremath{\mathtt{{\color{var}d}}}>0$, the \ensuremath{\mathtt{radixSort({\color{var}a},{\color{var}k})}} method can sort an array
  \ensuremath{\mathtt{{\color{var}a}}} containing \ensuremath{\mathtt{{\color{var}n}}} \ensuremath{\mathtt{{\color{var}w}}}-bit integers in $O((\ensuremath{\mathtt{{\color{var}w}}}/\ensuremath{\mathtt{{\color{var}d}}})(\ensuremath{\mathtt{{\color{var}n}}}+2^{\ensuremath{\mathtt{{\color{var}d}}}}))$ time.
\end{thm}

If we think, instead, of the elements of the array being in the range
$\{0,\ldots,\ensuremath{\mathtt{{\color{var}n}}}^c-1\}$, and take $\ensuremath{\mathtt{{\color{var}d}}}=\lceil\log\ensuremath{\mathtt{{\color{var}n}}}\rceil$ we obtain
the following version of \thmref{radix-sort}.
\begin{cor}\corlabel{radix-sort-poly}
  The \ensuremath{\mathtt{radixSort({\color{var}a},{\color{var}k})}} method can sort an array \ensuremath{\mathtt{{\color{var}a}}} containing \ensuremath{\mathtt{{\color{var}n}}}
  integer values in the range $\{0,\ldots,\ensuremath{\mathtt{{\color{var}n}}}^c-1\}$ in $O(c\ensuremath{\mathtt{{\color{var}n}}})$ time.
\end{cor}

\section{Discussion and Exercises}

Sorting is \emph{the} fundamental algorithmic problem in computer science,
and it has a long history.  Knuth \cite{k97v3} attributes the merge-sort
algorithm to von~Neumann (1945).  Quicksort is due to Hoare \cite{h61}.
The original heap-sort algorithm is due to Williams \cite{w64}, but the
version presented here (in which the heap is constructed bottom-up
in $O(\ensuremath{\mathtt{{\color{var}n}}})$ time) is due to Floyd \cite{f64}.  Lower-bounds for
comparison-based sorting appear to be folklore.  The following table
summarizes the performance of these comparison-based algorithms:

\begin{center}
  \begin{tabular}{|l|r@{}l@{ }l|l|} \hline
      & \multicolumn{3}{c|}{comparisons} & in-place  \\ \hline
    Merge-sort & $\ensuremath{\mathtt{{\color{var}n}}}\log \ensuremath{\mathtt{{\color{var}n}}}$ & &  worst-case & No  \\
    Quicksort & $1.38\ensuremath{\mathtt{{\color{var}n}}}\log \ensuremath{\mathtt{{\color{var}n}}}$ & ${}+ O(\ensuremath{\mathtt{{\color{var}n}}})$ & expected & Yes \\
    Heap-sort & $2\ensuremath{\mathtt{{\color{var}n}}}\log \ensuremath{\mathtt{{\color{var}n}}}$ & ${}+ O(\ensuremath{\mathtt{{\color{var}n}}})$ & worst-case & Yes \\ \hline
  \end{tabular}
\end{center}

Each of these comparison-based algorithms has its advantages and
disadvantages.  Merge-sort does the fewest comparisons and does not rely
on randomization.  Unfortunately, it uses an auxilliary array during its
merge phase.  Allocating this array can be expensive and is a potential
point of failure if memory is limited.  Quicksort is an \emph{in-place}
\index{in-place algorithm}%
algorithm and is a close second in terms of the number of comparisons,
but is randomized, so this running time is not always guaranteed.
Heap-sort does the most comparisons, but it is in-place and deterministic.

There is one setting in which merge-sort is a clear-winner;  this
occurs when sorting a linked-list.  In this case, the auxiliary
array is not needed;  two sorted linked lists are very easily merged
into a single sorted linked-list by pointer manipulations (see
\excref{list-merge-sort}).

The counting-sort and radix-sort algorithms described here are due
to Seward \cite[Section~2.4.6]{s54}.  However, variants of radix-sort
have been used since the 1920s to sort punch cards using punched card
sorting machines.  These machines can sort a stack of cards into two
piles based on the existence (or not) of a hole in a specific location
on the card.  Repeating this process for different hole locations gives
an implementation of radix-sort.

Finally, we note that counting sort and radix-sort can be used to sort
other types of numbers besides non-negative integers.  Straightforward
modifications of counting sort can sort integers, in any interval
$\{a,\ldots,b\}$, in $O(\ensuremath{\mathtt{{\color{var}n}}}+b-a)$ time.  Similarly, radix sort can sort
integers in the same interval in $O(\ensuremath{\mathtt{{\color{var}n}}}(\log_{\ensuremath{\mathtt{{\color{var}n}}}}(b-a))$ time.  Finally, both of these
algorithms can also be used to sort floating point numbers in the IEEE 754
floating point format.  This is because the IEEE format is designed to
allow the comparison of two floating point numbers by comparing their values
as if they were integers in a signed-magnitude binary representation
\cite{ieee754}.

\begin{exc}
  Illustrate the execution of merge-sort and heap-sort on an input array
  containing $1,7,4,6,2,8,3,5$.  Give a sample illustration of one possible
  execution of quicksort on the same array.
\end{exc}

\begin{exc}\exclabel{list-merge-sort}
  Implement a version of the merge-sort algorithm that sorts a \texttt{DLList}
  without using an auxiliary array. (See \excref{dllist-sort}.) 
\end{exc}

\begin{exc}
  Some implementations of \ensuremath{\mathtt{quickSort({\color{var}a},{\color{var}i},{\color{var}n},{\color{var}c})}} always use \ensuremath{\mathtt{{\color{var}a}[{\color{var}i}]}}
  as a pivot.  Give an example of an input array of length \ensuremath{\mathtt{{\color{var}n}}} in which
  such an implementation would perform $\binom{\ensuremath{\mathtt{{\color{var}n}}}}{2}$ comparisons.
\end{exc}

\begin{exc}
  Some implementations of \ensuremath{\mathtt{quickSort({\color{var}a},{\color{var}i},{\color{var}n},{\color{var}c})}} always use \ensuremath{\mathtt{{\color{var}a}[{\color{var}i}+{\color{var}n}/2]}}
  as a pivot.  Given an example of an input array of length \ensuremath{\mathtt{{\color{var}n}}} in which
  such an implementation would perform $\binom{\ensuremath{\mathtt{{\color{var}n}}}}{2}$ comparisons.
\end{exc}

\begin{exc}
  Show that, for any implementation of \ensuremath{\mathtt{quickSort({\color{var}a},{\color{var}i},{\color{var}n},{\color{var}c})}}
  that chooses a pivot deterministically, without first looking at
  any values in $\ensuremath{\mathtt{{\color{var}a}[{\color{var}i}]}},\ldots,\ensuremath{\mathtt{{\color{var}a}[{\color{var}i}+{\color{var}n}-1]}}$, there exists an input array of length \ensuremath{\mathtt{{\color{var}n}}}
  that causes this implementation to perform $\binom{\ensuremath{\mathtt{{\color{var}n}}}}{2}$ comparisons.
\end{exc}

\begin{exc}
  Design a \texttt{Comparator}, \ensuremath{\mathtt{{\color{var}c}}}, that you could pass as an argument
  to \ensuremath{\mathtt{quickSort({\color{var}a},{\color{var}i},{\color{var}n},{\color{var}c})}} and that would cause quicksort to perform
  $\binom{\ensuremath{\mathtt{{\color{var}n}}}}{2}$ comparisons.  (Hint: Your comparator does not actually
  need to look at the values being compared.)
\end{exc}

\begin{exc}
  Analyze the expected number of comparisons done by Quicksort a little
  more carefully than the proof of \thmref{quicksort}.  In particular, show
  that the expected number of comparisons is $2\ensuremath{\mathtt{{\color{var}n}}}H_\ensuremath{\mathtt{{\color{var}n}}} -\ensuremath{\mathtt{{\color{var}n}}} + H_\ensuremath{\mathtt{{\color{var}n}}}$.
\end{exc}

\begin{exc}
  Describe an input array that causes heap sort to perform at least
  $2\ensuremath{\mathtt{{\color{var}n}}}\log \ensuremath{\mathtt{{\color{var}n}}}-O(\ensuremath{\mathtt{{\color{var}n}}})$ comparisons.  Justify your answer.
\end{exc}

\javaonly{
\begin{exc}
  The heap sort implementation described here sorts the elements into
  reverse sorted order and then reverses the array.  This last step
  could be avoided by defining a new \texttt{Comparator} that negates the
  results of the input \texttt{Comparator}, \ensuremath{\mathtt{{\color{var}c}}}.  Explain why this would not
  be a good optimization.  (Hint: Consider how many negations would need
  to be done in relation to how long it takes to reverse the array.)
\end{exc}
}

\begin{exc}
  Find another pair of permutations of $1,2,3$ that are not correctly
  sorted by the comparison tree in \figref{comparison-tree-2}.
\end{exc}

\begin{exc}\exclabel{log-factorial}
  Prove that $\log \ensuremath{\mathtt{{\color{var}n}}}! = \ensuremath{\mathtt{{\color{var}n}}}\log \ensuremath{\mathtt{{\color{var}n}}}-O(\ensuremath{\mathtt{{\color{var}n}}})$.
\end{exc}

\begin{exc}
  Prove that a binary tree with $k$ leaves has height at least $\log k$.
\end{exc}

\begin{exc}\exclabel{randomized-lower-bound}
  Prove that, if we pick a random leaf from a binary tree with $k$
  leaves, then the expected height of this leaf is at least $\log k$.
\end{exc}


\begin{exc}
  The implementation of \ensuremath{\mathtt{radixSort({\color{var}a},{\color{var}k})}} given here works when the input
  array, \ensuremath{\mathtt{{\color{var}a}}} contains only \cpponly{unsigned}\javaonly{non-negative}
  integers.  \javaonly{Extend this implementation so that it also
  works correctly when \ensuremath{\mathtt{{\color{var}a}}} contains both negative and non-negative
  integers.}\cpponly{Write a version that works correctly for signed
  integers.}
\end{exc}


\clearpage

\clearpage
\chapter{Graphs}
\chaplabel{graphs}


In this chapter, we study two representations of graphs and basic
algorithms that use these representations.  

Mathematically, a \emph{(directed) graph}
\index{graph}%
\index{directed graph}%
is a pair $G=(V,E)$ where
$V$ is a set of \emph{vertices}
\index{vertex}%
and $E$ is a set of ordered pairs
of vertices called \emph{edges}.
\index{edge}%
An edge \ensuremath{\mathtt{({\color{var}i},{\color{var}j})}} is \emph{directed}
\index{directed edge}%
from \ensuremath{\mathtt{{\color{var}i}}} to \ensuremath{\mathtt{{\color{var}j}}};  \ensuremath{\mathtt{{\color{var}i}}} is called the \emph{source}
\index{source} of the edge and \ensuremath{\mathtt{{\color{var}j}}}
is called the \emph{target}.
\index{target}  A \emph{path}%
\index{path} in $G$ is a sequence of
vertices $v_0,\ldots,v_k$ such that, for every $i\in\{1,\ldots,k\}$,
the edge $(v_{i-1},v_{i})$ is in $E$.  A path $v_0,\ldots,v_k$ is a
\emph{cycle}
\index{cycle}%
if, additionally, the edge $(v_k,v_0)$ is in $E$.  A path (or
cycle) is \emph{simple}
\index{simple path/cycle}%
if all of its vertices are unique.  If there
is a path from some vertex $v_i$ to some vertex $v_j$ then we say that
$v_j$ is \emph{reachable}
\index{reachable vertex} from $v_i$.  An example of a graph is shown
in \figref{graph}.

\begin{figure}
  \begin{center}
    \includegraphics[scale=0.90909]{figs/graph}
  \end{center}
  \caption{A graph with twelve vertices.  Vertices are drawn as numbered
    circles and edges are drawn as pointed curves pointing from source
    to target.}
  \figlabel{graph}
\end{figure}

Due to their ability to model so many phenomena, graphs have an enormous
number of applications. There are many obvious examples. Computer
networks can be modelled as graphs, with vertices corresponding to
computers and edges corresponding to (directed) communication links
between those computers.  City streets can be modelled as graphs,
with vertices representing intersections and edges representing streets
joining consecutive intersections.

Less obvious examples occur as soon as we realize that graphs can model
any pairwise relationships within a set. For example, in a university
setting we might have a timetable \emph{conflict graph}
\index{conflict graph}%
whose vertices represent
courses offered in the university and in which the edge \ensuremath{\mathtt{({\color{var}i},{\color{var}j})}} is present
if and only if there is at least one student that is taking both class
\ensuremath{\mathtt{{\color{var}i}}} and class \ensuremath{\mathtt{{\color{var}j}}}.  Thus, an edge indicates that the exam for class \ensuremath{\mathtt{{\color{var}i}}}
should not be scheduled at the same time as the exam for class \ensuremath{\mathtt{{\color{var}j}}}.

Throughout this section, we will use \ensuremath{\mathtt{{\color{var}n}}} to denote the number of vertices
of $G$ and \ensuremath{\mathtt{{\color{var}m}}} to denote the number of edges of $G$.  That is, $\ensuremath{\mathtt{{\color{var}n}}}=|V|$
and $\ensuremath{\mathtt{{\color{var}m}}}=|E|$. Furthermore, we will assume that $V=\{0,\ldots,\ensuremath{\mathtt{{\color{var}n}}}-1\}$.
Any other data that we would like to associate with the elements of $V$
can be stored in an array of length $\ensuremath{\mathtt{{\color{var}n}}}$.

Some typical operations performed on graphs are:
\begin{itemize}
  \item \ensuremath{\mathtt{addEdge({\color{var}i},{\color{var}j})}}: Add the edge $(\ensuremath{\mathtt{{\color{var}i}}},\ensuremath{\mathtt{{\color{var}j}}})$ to $E$.
  \item \ensuremath{\mathtt{removeEdge({\color{var}i},{\color{var}j})}}: Remove the edge $(\ensuremath{\mathtt{{\color{var}i}}},\ensuremath{\mathtt{{\color{var}j}}})$ from $E$.
  \item \ensuremath{\mathtt{hasEdge({\color{var}i},{\color{var}j})}}: Check if the edge $(\ensuremath{\mathtt{{\color{var}i}}},\ensuremath{\mathtt{{\color{var}j}}})\in E$ 
  \item \ensuremath{\mathtt{outEdges({\color{var}i})}}: Return a \texttt{List} of all integers $\ensuremath{\mathtt{{\color{var}j}}}$ such that
  $(\ensuremath{\mathtt{{\color{var}i}}},\ensuremath{\mathtt{{\color{var}j}}})\in E$
  \item \ensuremath{\mathtt{inEdges({\color{var}i})}}: Return a \texttt{List} of all integers $\ensuremath{\mathtt{{\color{var}j}}}$ such that
  $(\ensuremath{\mathtt{{\color{var}j}}},\ensuremath{\mathtt{{\color{var}i}}})\in E$
\end{itemize}

Note that these operations are not terribly difficult to implement
efficiently. For example, the first three operations can be implemented
directly by using a \texttt{USet}, so they can be implemented in constant
expected time using the hash tables discussed in \chapref{hashing}.
The last two operations can be implemented in constant time by storing,
for each vertex, a list of its adjacent vertices.

However, different applications of graphs have different performance
requirements for these operations and, ideally, we can use the simplest
implementation that satisfies all the application's requirements.
For this reason, we discuss two broad categories of graph representations.

\section{\texttt{AdjacencyMatrix}: Representing a Graph by a Matrix}
\seclabel{adjacency-matrix}

\index{adjacency matrix}%
An \emph{adjacency matrix} is a way of representing an \ensuremath{\mathtt{{\color{var}n}}} vertex graph
$G=(V,E)$ by an $\ensuremath{\mathtt{{\color{var}n}}}\times\ensuremath{\mathtt{{\color{var}n}}}$ matrix, \ensuremath{\mathtt{{\color{var}a}}}, whose entries are boolean
values.


 \ \vspace{.5ex}\begin{Verbatim}[gobble=1,tabsize=2,frame=single,framerule=1.2pt,commandchars=\\@\$,label=\texttt{AdjacencyMatrix},labelposition=topline]
	@\color@keyword$int$ @\color@var$n$;
	@\color@var$boolean$[][] @\color@var$a$;
	AdjacencyMatrix(@\color@keyword$int$ @\color@var$n0$) {
		@\color@var$n$ = @\color@var$n0$;
		@\color@var$a$ = @\color@keyword$new$ @\color@var$boolean$[@\color@var$n$][@\color@var$n$];
	}
\end{Verbatim}


The matrix entry \ensuremath{\mathtt{{\color{var}a}[{\color{var}i}][{\color{var}j}]}} is defined as
\[  \ensuremath{\mathtt{{\color{var}a}[{\color{var}i}][{\color{var}j}]}}= 
    \begin{cases}
      \ensuremath{\mathtt{{\color{var}true}}} & \text{if $\ensuremath{\mathtt{({\color{var}i},{\color{var}j})}}\in E$} \\
      \ensuremath{\mathtt{{\color{var}false}}} & \text{otherwise}
    \end{cases}
\]
The adjacency matrix for the graph in \figref{graph} is
shown in \figref{graph-adj}.

In this representation, the operations \ensuremath{\mathtt{addEdge({\color{var}i},{\color{var}j})}},
\ensuremath{\mathtt{removeEdge({\color{var}i},{\color{var}j})}}, and \ensuremath{\mathtt{hasEdge({\color{var}i},{\color{var}j})}} just
involve setting or reading the matrix entry \ensuremath{\mathtt{{\color{var}a}[{\color{var}i}][{\color{var}j}]}}:


 \ \vspace{.5ex}\begin{Verbatim}[gobble=1,tabsize=2,frame=single,framerule=1.2pt,commandchars=\\@\$,label=\texttt{AdjacencyMatrix},labelposition=topline]
	@\color@keyword$void$ addEdge(@\color@keyword$int$ @\color@var$i$, @\color@keyword$int$ @\color@var$j$) {
		@\color@var$a$[@\color@var$i$][@\color@var$j$] = @\color@var$true$;
	}
	@\color@keyword$void$ removeEdge(@\color@keyword$int$ @\color@var$i$, @\color@keyword$int$ @\color@var$j$) {
		@\color@var$a$[@\color@var$i$][@\color@var$j$] = @\color@var$false$;
	}
	@\color@var$boolean$ hasEdge(@\color@keyword$int$ @\color@var$i$, @\color@keyword$int$ @\color@var$j$) {
		@\color@var$return$ @\color@var$a$[@\color@var$i$][@\color@var$j$];
	}
\end{Verbatim}

These operations clearly take constant time per operation.

\begin{figure}
  \begin{center}
    \includegraphics[scale=0.90909]{figs/graph} \\[3ex]
    \begin{tabular}{c|cccccccccccc}
        &0&1&2&3&4&5&6&7&8&9&10&11 \\\hline
       0&0&1&0&0&1&0&0&0&0&0&0 &0\\
       1&1&0&1&0&0&1&1&0&0&0&0 &0\\
       2&1&0&0&1&0&0&1&0&0&0&0 &0\\
       3&0&0&1&0&0&0&0&1&0&0&0 &0\\
       4&1&0&0&0&0&1&0&0&1&0&0 &0\\
       5&0&1&1&0&1&0&1&0&0&1&0 &0\\
       6&0&0&1&0&0&1&0&1&0&0&1 &0\\
       7&0&0&0&1&0&0&1&0&0&0&0 &1\\
       8&0&0&0&0&1&0&0&0&0&1&0 &0\\
       9&0&0&0&0&0&1&0&0&1&0&1 &0\\
      10&0&0&0&0&0&0&1&0&0&1&0 &1\\
      11&0&0&0&0&0&0&0&1&0&0&1 &0\\
    \end{tabular} 
  \end{center}
  \caption{A graph and its adjacency matrix.}
  \figlabel{graph-adj}
\end{figure}

Where the adjacency matrix performs poorly is with the \ensuremath{\mathtt{outEdges({\color{var}i})}} and
\ensuremath{\mathtt{inEdges({\color{var}i})}} operations.  To implement these, we must scan all \ensuremath{\mathtt{{\color{var}n}}}
entries in the corresponding row or column of \ensuremath{\mathtt{{\color{var}a}}} and gather up all the
indices, \ensuremath{\mathtt{{\color{var}j}}}, where \ensuremath{\mathtt{{\color{var}a}[{\color{var}i}][{\color{var}j}]}}, respectively \ensuremath{\mathtt{{\color{var}a}[{\color{var}j}][{\color{var}i}]}}, is true.


 \ \vspace{.5ex}\begin{Verbatim}[gobble=1,tabsize=2,frame=single,framerule=1.2pt,commandchars=\\@\$,label=\texttt{AdjacencyMatrix},labelposition=topline]
	List<Integer> outEdges(@\color@keyword$int$ @\color@var$i$) {
		List<Integer> @\color@var$edges$ = @\color@keyword$new$ ArrayList<Integer>();
		@\color@keyword$for$ (@\color@keyword$int$ @\color@var$j$ = 0; @\color@var$j$ < @\color@var$n$; @\color@var$j$++) 
			@\color@keyword$if$ (@\color@var$a$[@\color@var$i$][@\color@var$j$]) @\color@var$edges$.add(@\color@var$j$);
		@\color@var$return$ @\color@var$edges$;
	}
	List<Integer> inEdges(@\color@keyword$int$ @\color@var$i$) {
		List<Integer> @\color@var$edges$ = @\color@keyword$new$ ArrayList<Integer>();
		@\color@keyword$for$ (@\color@keyword$int$ @\color@var$j$ = 0; @\color@var$j$ < @\color@var$n$; @\color@var$j$++)
			@\color@keyword$if$ (@\color@var$a$[@\color@var$j$][@\color@var$i$]) @\color@var$edges$.add(@\color@var$j$);
		@\color@var$return$ @\color@var$edges$;
	}
\end{Verbatim}

These operations clearly take $O(\ensuremath{\mathtt{{\color{var}n}}})$ time per operation.  

Another drawback of the adjacency matrix representation is that it
is large.  It stores an $\ensuremath{\mathtt{{\color{var}n}}}\times \ensuremath{\mathtt{{\color{var}n}}}$ boolean matrix, so it requires at
least $\ensuremath{\mathtt{{\color{var}n}}}^2$ bits of memory.  The implementation here uses a matrix
of \javaonly{\ensuremath{\mathtt{{\color{var}boolean}}}}\cpponly{\ensuremath{\mathtt{{\color{var}bool}}}} values so it actually uses on the
order of $\ensuremath{\mathtt{{\color{var}n}}}^2$ bytes of memory.  A more careful implementation, which
packs \ensuremath{\mathtt{{\color{var}w}}} boolean values into each word of memory, could reduce this
space usage to $O(\ensuremath{\mathtt{{\color{var}n}}}^2/\ensuremath{\mathtt{{\color{var}w}}})$ words of memory.

\begin{thm}
The \texttt{AdjacencyMatrix} data structure implements the \texttt{Graph} interface.
An \texttt{AdjacencyMatrix} supports the operations
\begin{itemize}
  \item \ensuremath{\mathtt{addEdge({\color{var}i},{\color{var}j})}}, \ensuremath{\mathtt{removeEdge({\color{var}i},{\color{var}j})}}, and \ensuremath{\mathtt{hasEdge({\color{var}i},{\color{var}j})}} in constant
  time per operation; and
  \item \ensuremath{\mathtt{inEdges({\color{var}i})}}, and \ensuremath{\mathtt{outEdges({\color{var}i})}} in $O(\ensuremath{\mathtt{{\color{var}n}}})$ time per operation.
\end{itemize}
The space used by an \texttt{AdjacencyMatrix} is  $O(\ensuremath{\mathtt{{\color{var}n}}}^2)$.
\end{thm}

Despite its high memory requirements and poor performance of the \ensuremath{\mathtt{inEdges({\color{var}i})}}
and \ensuremath{\mathtt{outEdges({\color{var}i})}} operations, an \texttt{AdjacencyMatrix} can still be useful for
some applications.  In particular, when the graph $G$ is \emph{dense},
i.e., it has close to $\ensuremath{\mathtt{{\color{var}n}}}^2$ edges, then a memory usage of $\ensuremath{\mathtt{{\color{var}n}}}^2$
may be acceptable.

The \texttt{AdjacencyMatrix} data structure is also commonly used because
algebraic operations on the matrix \ensuremath{\mathtt{{\color{var}a}}} can be used to efficiently compute
properties of the graph $G$.  This is a topic for a course on algorithms,
but we point out one such property here:  If we treat the entries of \ensuremath{\mathtt{{\color{var}a}}}
as integers (1 for \ensuremath{\mathtt{{\color{var}true}}} and 0 for \ensuremath{\mathtt{{\color{var}false}}}) and multiply \ensuremath{\mathtt{{\color{var}a}}} by itself
using matrix multiplication then we get the matrix $\ensuremath{\mathtt{{\color{var}a}}}^2$.  Recall,
from the definition of matrix multiplication, that
\[
    \ensuremath{\mathtt{{\color{var}a}^2[{\color{var}i}][{\color{var}j}]}} = \sum_{k=0}^{\ensuremath{\mathtt{{\color{var}n}}}-1} \ensuremath{\mathtt{{\color{var}a}[{\color{var}i}][{\color{var}k}]}}\cdot \ensuremath{\mathtt{{\color{var}a}[{\color{var}k}][{\color{var}j}]}} \enspace .
\]
Interpreting this sum in terms of the graph $G$, this formula counts the
number of vertices, $\ensuremath{\mathtt{{\color{var}k}}}$, such that $G$ contains both edges \ensuremath{\mathtt{({\color{var}i},{\color{var}k})}}
and \ensuremath{\mathtt{({\color{var}k},{\color{var}j})}}.  That is, it counts the number of paths from $\ensuremath{\mathtt{{\color{var}i}}}$ to $\ensuremath{\mathtt{{\color{var}j}}}$
(through intermediate vertices, $\ensuremath{\mathtt{{\color{var}k}}}$) whose length is exactly two.
This observation is the foundation of an algorithm that computes the
shortest paths between all pairs of vertices in $G$ using only $O(\log
\ensuremath{\mathtt{{\color{var}n}}})$ matrix multiplications.

\section{\texttt{AdjacencyLists}: A Graph as a Collection of Lists}
\seclabel{adjacency-list}

\index{adjacency list}%
\emph{Adjacency list} representations of graphs take a more vertex-centric
approach.  There are many possible implementations of adjacency lists.
In this section, we present a simple one.  At the end of the section,
we discuss different possibilities.  In an adjacency list representation,
the graph $G=(V,E)$ is represented as an array, \ensuremath{\mathtt{{\color{var}adj}}}, of lists.  The list
\ensuremath{\mathtt{{\color{var}adj}[{\color{var}i}]}} contains a list of all the vertices adjacent to vertex \ensuremath{\mathtt{{\color{var}i}}}.
That is, it contains every index \ensuremath{\mathtt{{\color{var}j}}} such that $\ensuremath{\mathtt{({\color{var}i},{\color{var}j})}}\in E$.


 \ \vspace{.5ex}\begin{Verbatim}[gobble=1,tabsize=2,frame=single,framerule=1.2pt,commandchars=\\@\$,label=\texttt{AdjacencyLists},labelposition=topline]
	@\color@keyword$int$ @\color@var$n$;
	List<Integer>[] @\color@var$adj$;
	AdjacencyLists(@\color@keyword$int$ @\color@var$n0$) {
		@\color@var$n$ = @\color@var$n0$;
		@\color@var$adj$ = (List<Integer>[])@\color@keyword$new$ List[@\color@var$n$];
		@\color@keyword$for$ (@\color@keyword$int$ @\color@var$i$ = 0; @\color@var$i$ < @\color@var$n$; @\color@var$i$++) 
			@\color@var$adj$[@\color@var$i$] = @\color@keyword$new$ ArrayStack<Integer>();
	}
\end{Verbatim}

(An example is shown in \figref{graph-adjlist}.)  In this particular
implementation, we represent each list in \ensuremath{\mathtt{{\color{var}adj}}} as \javaonly{an}\cpponly{a
subclass of} \texttt{ArrayStack}, because we would like constant time access
by position. Other options are also possible.  Specifically, we could
have implemented \ensuremath{\mathtt{{\color{var}adj}}} as a \texttt{DLList}.


\begin{figure}
  \begin{center}
    \includegraphics[scale=0.90909]{figs/graph} \\[3ex]
    \begin{tabular}{|c|c|c|c|c|c|c|c|c|c|c|c|c|}\hline
        0&1&2&3&4&5&6 &7 &8&9 &10&11 \\\hline
        1&0&1&2&0&1&5 &6 &4&8 &9 &10 \\
        4&2&3&7&5&2&2 &3 &9&5 &6 &7 \\
         &6&6& &8&6&7 &11& &10&11& \\
         &5& & & &9&10&  & &  &  & \\
         & & & & &4&  &  & &  &  & \\
    \end{tabular} 
  \end{center}
  \caption{A graph and its adjacency lists}
  \figlabel{graph-adjlist}
\end{figure}



The \ensuremath{\mathtt{addEdge({\color{var}i},{\color{var}j})}} operation just appends the value \ensuremath{\mathtt{{\color{var}j}}} to the list \ensuremath{\mathtt{{\color{var}adj}[{\color{var}i}]}}:


 \ \vspace{.5ex}\begin{Verbatim}[gobble=1,tabsize=2,frame=single,framerule=1.2pt,commandchars=\\@\$,label=\texttt{AdjacencyLists},labelposition=topline]
	@\color@keyword$void$ addEdge(@\color@keyword$int$ @\color@var$i$, @\color@keyword$int$ @\color@var$j$) {
		@\color@var$adj$[@\color@var$i$].add(@\color@var$j$);
	}
\end{Verbatim}

This takes constant time.

The \ensuremath{\mathtt{removeEdge({\color{var}i},{\color{var}j})}} operation searches through the list \ensuremath{\mathtt{{\color{var}adj}[{\color{var}i}]}}
until it finds \ensuremath{\mathtt{{\color{var}j}}} and then removes it:


 \ \vspace{.5ex}\begin{Verbatim}[gobble=1,tabsize=2,frame=single,framerule=1.2pt,commandchars=\\@\$,label=\texttt{AdjacencyLists},labelposition=topline]
	@\color@keyword$void$ removeEdge(@\color@keyword$int$ @\color@var$i$, @\color@keyword$int$ @\color@var$j$) {
		Iterator<Integer> @\color@var$it$ = @\color@var$adj$[@\color@var$i$].iterator();
		@\color@keyword$while$ (@\color@var$it$.hasNext()) {
			@\color@keyword$if$ (@\color@var$it$.next() == @\color@var$j$) {
				@\color@var$it$.remove();
				@\color@var$return$;
			}
		}	
	}
\end{Verbatim}

This takes $O(\deg(\ensuremath{\mathtt{{\color{var}i}}}))$ time, where $\deg(\ensuremath{\mathtt{{\color{var}i}}})$ (the \emph{degree}
\index{degree}%
of
$\ensuremath{\mathtt{{\color{var}i}}}$) counts the number of edges in $E$ that have $\ensuremath{\mathtt{{\color{var}i}}}$ as their source.

The \ensuremath{\mathtt{hasEdge({\color{var}i},{\color{var}j})}} operation is similar;  it searches through the list
\ensuremath{\mathtt{{\color{var}adj}[{\color{var}i}]}} until it finds \ensuremath{\mathtt{{\color{var}j}}} (and returns true), or reaches the end of
the list (and returns false):


 \ \vspace{.5ex}\begin{Verbatim}[gobble=1,tabsize=2,frame=single,framerule=1.2pt,commandchars=\\@\$,label=\texttt{AdjacencyLists},labelposition=topline]
	@\color@var$boolean$ hasEdge(@\color@keyword$int$ @\color@var$i$, @\color@keyword$int$ @\color@var$j$) {
		@\color@var$return$ @\color@var$adj$[@\color@var$i$].contains(@\color@var$j$);
	}
\end{Verbatim}

This also takes $O(\deg(\ensuremath{\mathtt{{\color{var}i}}}))$ time.

The \ensuremath{\mathtt{outEdges({\color{var}i})}} operation is very simple;
\pcodeonly{it returns the list \ensuremath{\mathtt{{\color{var}adj}[{\color{var}i}]}}}
\javaonly{it returns the list \ensuremath{\mathtt{{\color{var}adj}[{\color{var}i}]}}}
\cpponly{it copies the values in \ensuremath{\mathtt{{\color{var}adj}[{\color{var}i}]}} into the output list}:
\pcodeimport{ods/AdjacencyLists.outEdges(i)}


 \ \vspace{.5ex}\begin{Verbatim}[gobble=1,tabsize=2,frame=single,framerule=1.2pt,commandchars=\\@\$,label=\texttt{AdjacencyLists},labelposition=topline]
	List<Integer> outEdges(@\color@keyword$int$ @\color@var$i$) {
		@\color@var$return$ @\color@var$adj$[@\color@var$i$];
	}
\end{Verbatim}

\javaonly{This clearly takes constant time.}\cpponly{This clearly takes $O(\deg(\ensuremath{\mathtt{{\color{var}i}}}))$ time.}

The \ensuremath{\mathtt{inEdges({\color{var}i})}} operation is much more work.  It scans over every
vertex $j$ checking if the edge \ensuremath{\mathtt{({\color{var}i},{\color{var}j})}} exists and, if so, adding \ensuremath{\mathtt{{\color{var}j}}}
to the output list:
\pcodeimport{ods/AdjacencyLists.inEdges(i)}


 \ \vspace{.5ex}\begin{Verbatim}[gobble=1,tabsize=2,frame=single,framerule=1.2pt,commandchars=\\@\$,label=\texttt{AdjacencyLists},labelposition=topline]
	List<Integer> inEdges(@\color@keyword$int$ @\color@var$i$) {
		List<Integer> @\color@var$edges$ = @\color@keyword$new$ ArrayStack<Integer>();
		@\color@keyword$for$ (@\color@keyword$int$ @\color@var$j$ = 0; @\color@var$j$ < @\color@var$n$; @\color@var$j$++)
			@\color@keyword$if$ (@\color@var$adj$[@\color@var$j$].contains(@\color@var$i$))	@\color@var$edges$.add(@\color@var$j$);
		@\color@var$return$ @\color@var$edges$;
	}
\end{Verbatim}

This operation is very slow. It scans the adjacency list of every vertex,
so it takes $O(\ensuremath{\mathtt{{\color{var}n}}} + \ensuremath{\mathtt{{\color{var}m}}})$ time.

The following theorem summarizes the performance of the above data structure:

\begin{thm}
The \texttt{AdjacencyLists} data structure implements the \texttt{Graph} interface.
An \texttt{AdjacencyLists} supports the operations
\begin{itemize}
  \item \ensuremath{\mathtt{addEdge({\color{var}i},{\color{var}j})}} in constant time per operation;
  \item \ensuremath{\mathtt{removeEdge({\color{var}i},{\color{var}j})}} and \ensuremath{\mathtt{hasEdge({\color{var}i},{\color{var}j})}} in $O(\deg(\ensuremath{\mathtt{{\color{var}i}}}))$ time
    per operation;
  \javaonly{\item \ensuremath{\mathtt{outEdges({\color{var}i})}} in constant time per operation; and}
  \cpponly{\item \ensuremath{\mathtt{outEdges({\color{var}i})}} in $O(\deg(\ensuremath{\mathtt{{\color{var}i}}}))$ time per operation; and}
  \item \ensuremath{\mathtt{inEdges({\color{var}i})}} in $O(\ensuremath{\mathtt{{\color{var}n}}}+\ensuremath{\mathtt{{\color{var}m}}})$ time per operation.
\end{itemize}
The space used by a \texttt{AdjacencyLists} is  $O(\ensuremath{\mathtt{{\color{var}n}}}+\ensuremath{\mathtt{{\color{var}m}}})$.
\end{thm}

As alluded to earlier, there are many different choices to be made when
implementing a graph as an adjacency list.  Some questions that come
up include:
\begin{itemize}
  \item What type of collection should be used to store each element
  of \ensuremath{\mathtt{{\color{var}adj}}}?  One could use an array-based list,  a linked-list, or even
  a hashtable.
  \item Should there be a second adjacency list, \ensuremath{\mathtt{{\color{var}inadj}}}, that stores,
  for each \ensuremath{\mathtt{{\color{var}i}}}, the list of vertices, \ensuremath{\mathtt{{\color{var}j}}}, such that $\ensuremath{\mathtt{({\color{var}j},{\color{var}i})}}\in E$?
  This can greatly reduce the running-time of the \ensuremath{\mathtt{inEdges({\color{var}i})}}
  operation, but requires slightly more work when adding or removing
  edges.
  \item Should the entry for the edge \ensuremath{\mathtt{({\color{var}i},{\color{var}j})}} in \ensuremath{\mathtt{{\color{var}adj}[{\color{var}i}]}} be linked by
  a reference to the corresponding entry in \ensuremath{\mathtt{{\color{var}inadj}[{\color{var}j}]}}?
  \item Should edges be first-class objects with their own associated data?
  In this way, \ensuremath{\mathtt{{\color{var}adj}}} would contain lists of edges rather than lists of vertices (integers).
\end{itemize}
Most of these questions come down to a tradeoff between complexity (and
space) of implementation and performance features of the implementation.

\section{Graph Traversal}

In this section we present two algorithms for exploring a graph,
starting at one of its vertices, \ensuremath{\mathtt{{\color{var}i}}}, and finding all vertices that
are reachable from \ensuremath{\mathtt{{\color{var}i}}}.  Both of these algorithms are best suited to
graphs represented using an adjacency list representation.  Therefore,
when analyzing these algorithms we will assume that the underlying
representation is an \texttt{AdjacencyLists}.

\subsection{Breadth-First Search}

\index{breadth-first-search}%
The \emph{breadth-first-search} algorithm starts at a vertex \ensuremath{\mathtt{{\color{var}i}}} and visits,
first the neighbours of \ensuremath{\mathtt{{\color{var}i}}}, then the neighbours of the neighbours of \ensuremath{\mathtt{{\color{var}i}}},
then the neighbours of the neighbours of the neighbours of \ensuremath{\mathtt{{\color{var}i}}}, and so on.

This algorithm is a generalization of the breadth-first traversal
algorithm for binary trees (\secref{bintree:traversal}), and is
very similar; it uses a queue, \ensuremath{\mathtt{{\color{var}q}}}, that initially contains only \ensuremath{\mathtt{{\color{var}i}}}.
It then repeatedly extracts an element from \ensuremath{\mathtt{{\color{var}q}}} and adds its neighbours
to \ensuremath{\mathtt{{\color{var}q}}}, provided that these neighbours have never been in \ensuremath{\mathtt{{\color{var}q}}} before.
The only major difference between the breadth-first-search algorithm
for graphs and the one for trees is that the algorithm for graphs has
to ensure that it does not add the same vertex to \ensuremath{\mathtt{{\color{var}q}}} more than once.
It does this by using an auxiliary boolean array, \ensuremath{\mathtt{{\color{var}seen}}}, that tracks
which vertices have already been discovered.


 \ \vspace{.5ex}\begin{Verbatim}[gobble=1,tabsize=2,frame=single,framerule=1.2pt,commandchars=\\@\$,label=\texttt{Algorithms},labelposition=topline]
	@\color@keyword$void$ bfs(Graph @\color@var$g$, @\color@keyword$int$ @\color@var$r$) {
		@\color@var$boolean$[] @\color@var$seen$ = @\color@keyword$new$ @\color@var$boolean$[@\color@var$g$.nVertices()];
		Queue<Integer> @\color@var$q$ = @\color@keyword$new$ SLList<Integer>();
		@\color@var$q$.add(@\color@var$r$);
		@\color@var$seen$[@\color@var$r$] = @\color@var$true$;
		@\color@keyword$while$ (!@\color@var$q$.isEmpty()) {
			@\color@keyword$int$ @\color@var$i$ = @\color@var$q$.remove();
			@\color@keyword$for$ (Integer @\color@var$j$ : @\color@var$g$.outEdges(@\color@var$i$)) {
				@\color@keyword$if$ (!@\color@var$seen$[@\color@var$j$]) {
					@\color@var$q$.add(@\color@var$j$);
					@\color@var$seen$[@\color@var$j$] = @\color@var$true$;
				}
			}
		}
	}
\end{Verbatim}

An example of running \ensuremath{\mathtt{bfs({\color{var}g},0)}} on the graph from \figref{graph}
is shown in \figref{graph-bfs}.  Different executions are possible,
depending on the ordering of the adjacency lists; \figref{graph-bfs}
uses the adjacency lists in \figref{graph-adjlist}.

\begin{figure}
  \begin{center}
    \includegraphics[scale=0.90909]{figs/graph-bfs}
  \end{center}
  \caption[Breadth-first-search]{An example of breadth-first-search starting at node 0. Nodes are
  labelled with the order in which they are added to \ensuremath{\mathtt{{\color{var}q}}}.  Edges that
  result in nodes being added to \ensuremath{\mathtt{{\color{var}q}}} are drawn in black, other edges
  are drawn in grey.}
  \figlabel{graph-bfs}
\end{figure}

Analyzing the running-time of the \ensuremath{\mathtt{bfs({\color{var}g},{\color{var}i})}} routine is fairly
straightforward.  The use of the \ensuremath{\mathtt{{\color{var}seen}}} array ensures that no vertex is
added to \ensuremath{\mathtt{{\color{var}q}}} more than once.  Adding (and later removing) each vertex
from \ensuremath{\mathtt{{\color{var}q}}} takes constant time per vertex for a total of $O(\ensuremath{\mathtt{{\color{var}n}}})$ time.
Since each vertex is processed by the inner loop at most once, each
adjacency list is processed at most once, so each edge of $G$ is processed
at most once.  This processing, which is done in the inner loop takes
constant time per iteration, for a total of $O(\ensuremath{\mathtt{{\color{var}m}}})$ time.  Therefore,
the entire algorithm runs in $O(\ensuremath{\mathtt{{\color{var}n}}}+\ensuremath{\mathtt{{\color{var}m}}})$ time.

The following theorem summarizes the performance of the \ensuremath{\mathtt{bfs({\color{var}g},{\color{var}r})}} algorithm.
\begin{thm}\thmlabel{bfs-graph}
  When given as input a \texttt{Graph}, \ensuremath{\mathtt{{\color{var}g}}}, that is implemented using the
  \texttt{AdjacencyLists} data structure, the \ensuremath{\mathtt{bfs({\color{var}g},{\color{var}r})}} algorithm runs in $O(\ensuremath{\mathtt{{\color{var}n}}}+\ensuremath{\mathtt{{\color{var}m}}})$
  time.
\end{thm}

A breadth-first traversal has some very special properties.  Calling
\ensuremath{\mathtt{bfs({\color{var}g},{\color{var}r})}} will eventually enqueue (and eventually dequeue) every vertex
\ensuremath{\mathtt{{\color{var}j}}} such that there is a directed path from \ensuremath{\mathtt{{\color{var}r}}} to \ensuremath{\mathtt{{\color{var}j}}}.  Moreover,
the vertices at distance 0 from \ensuremath{\mathtt{{\color{var}r}}} (\ensuremath{\mathtt{{\color{var}r}}} itself) will enter \ensuremath{\mathtt{{\color{var}q}}} before
the vertices at distance 1, which will enter \ensuremath{\mathtt{{\color{var}q}}} before the vertices at
distance 2, and so on.  Thus, the \ensuremath{\mathtt{bfs({\color{var}g},{\color{var}r})}} method visits vertices
in increasing order of distance from \ensuremath{\mathtt{{\color{var}r}}} and vertices that cannot be
reached from \ensuremath{\mathtt{{\color{var}r}}} are never visited at all.

A particularly useful application of the breadth-first-search algorithm
is, therefore, in computing shortest paths.  To compute the shortest
path from \ensuremath{\mathtt{{\color{var}r}}} to every other vertex, we use a variant of \ensuremath{\mathtt{bfs({\color{var}g},{\color{var}r})}}
that uses an auxilliary array, \ensuremath{\mathtt{{\color{var}p}}}, of length \ensuremath{\mathtt{{\color{var}n}}}.  When a new vertex
\ensuremath{\mathtt{{\color{var}j}}} is added to \ensuremath{\mathtt{{\color{var}q}}}, we set \ensuremath{\mathtt{{\color{var}p}[{\color{var}j}]={\color{var}i}}}.  In this way, \ensuremath{\mathtt{{\color{var}p}[{\color{var}j}]}} becomes the
second last node on a shortest path from \ensuremath{\mathtt{{\color{var}r}}} to \ensuremath{\mathtt{{\color{var}j}}}.  Repeating this,
by taking \ensuremath{\mathtt{{\color{var}p}[{\color{var}p}[{\color{var}j}]}}, \ensuremath{\mathtt{{\color{var}p}[{\color{var}p}[{\color{var}p}[{\color{var}j}]]]}}, and so on we can reconstruct the
(reversal of) a shortest path from \ensuremath{\mathtt{{\color{var}r}}} to \ensuremath{\mathtt{{\color{var}j}}}.



\subsection{Depth-First Search}

The \emph{depth-first-search}
\index{depth-first-search}%
algorithm is similar to the standard
algorithm for traversing binary trees;  it first fully explores one
subtree before returning to the current node and then exploring the
other subtree.  Another way to think of depth-first-search is by saying
that it is similar to breadth-first search except that it uses a stack
instead of a queue.

During the execution of the depth-first-search algorithm, each vertex,
\ensuremath{\mathtt{{\color{var}i}}}, is assigned a colour, \ensuremath{\mathtt{{\color{var}c}[{\color{var}i}]}}: \ensuremath{\mathtt{{\color{var}white}}} if we have never seen
the vertex before, \ensuremath{\mathtt{{\color{var}grey}}} if we are currently visiting that vertex,
and \ensuremath{\mathtt{{\color{var}black}}} if we are done visiting that vertex.  The easiest way to
think of depth-first-search is as a recursive algorithm.  It starts by
visiting \ensuremath{\mathtt{{\color{var}r}}}.  When visiting a vertex \ensuremath{\mathtt{{\color{var}i}}}, we first mark \ensuremath{\mathtt{{\color{var}i}}} as \ensuremath{\mathtt{{\color{var}grey}}}.
Next, we scan \ensuremath{\mathtt{{\color{var}i}}}'s adjacency list and recursively visit any white vertex
we find in this list.  Finally, we are done processing \ensuremath{\mathtt{{\color{var}i}}}, so we colour
\ensuremath{\mathtt{{\color{var}i}}} black and return.


 \ \vspace{.5ex}\begin{Verbatim}[gobble=1,tabsize=2,frame=single,framerule=1.2pt,commandchars=\\@\$,label=\texttt{Algorithms},labelposition=topline]
	@\color@keyword$void$ dfs(Graph @\color@var$g$, @\color@keyword$int$ @\color@var$r$) {
		@\color@keyword$byte$[] @\color@var$c$ = @\color@keyword$new$ @\color@keyword$byte$[@\color@var$g$.nVertices()];
		dfs(@\color@var$g$, @\color@var$r$, @\color@var$c$);
	}
	@\color@keyword$void$ dfs(Graph @\color@var$g$, @\color@keyword$int$ @\color@var$i$, @\color@keyword$byte$[] @\color@var$c$) {
		@\color@var$c$[@\color@var$i$] = @\color@var$grey$;  @\color@comment$// currently visiting i$
		@\color@keyword$for$ (Integer @\color@var$j$ : @\color@var$g$.outEdges(@\color@var$i$)) {
			@\color@keyword$if$ (@\color@var$c$[@\color@var$j$] == @\color@var$white$) {
				@\color@var$c$[@\color@var$j$] = @\color@var$grey$;
				dfs(@\color@var$g$, @\color@var$j$, @\color@var$c$);
			} 
		}
		@\color@var$c$[@\color@var$i$] = @\color@var$black$; @\color@comment$// done visiting i$
	}
\end{Verbatim}

An example of the execution of this algorithm is shown in \figref{graph-dfs}.

\begin{figure}
  \begin{center}
    \includegraphics[scale=0.90909]{figs/graph-dfs}
  \end{center}
  \caption[Depth-first-search]{An example of depth-first-search starting at node 0. Nodes are
  labelled with the order in which they are processed.  Edges that
  result in a recursive call are drawn in black, other edges
  are drawn in \ensuremath{\mathtt{{\color{var}grey}}}.}
  \figlabel{graph-dfs}
\end{figure}

Although depth-first-search may best be thought of as a recursive
algorithm, recursion is not the best way to implement it. Indeed, the code
given above will fail for many large graphs by causing a stack overflow.
An alternative implementation is to replace the recursion stack with an
explicit stack, \ensuremath{\mathtt{{\color{var}s}}}.  The following implementation does just that:


 \ \vspace{.5ex}\begin{Verbatim}[gobble=1,tabsize=2,frame=single,framerule=1.2pt,commandchars=\\@\$,label=\texttt{Algorithms},labelposition=topline]
	@\color@keyword$void$ dfs2(Graph @\color@var$g$, @\color@keyword$int$ @\color@var$r$) {
		@\color@keyword$byte$[] @\color@var$c$ = @\color@keyword$new$ @\color@keyword$byte$[@\color@var$g$.nVertices()];
		Stack<Integer> @\color@var$s$ = @\color@keyword$new$ Stack<Integer>();
		@\color@var$s$.push(@\color@var$r$);
		@\color@keyword$while$ (!@\color@var$s$.isEmpty()) {
			@\color@keyword$int$ @\color@var$i$ = @\color@var$s$.pop();
			@\color@keyword$if$ (@\color@var$c$[@\color@var$i$] == @\color@var$white$) {
				@\color@var$c$[@\color@var$i$] = @\color@var$grey$;
				@\color@keyword$for$ (@\color@keyword$int$ @\color@var$j$ : @\color@var$g$.outEdges(@\color@var$i$))
					@\color@var$s$.push(@\color@var$j$);
			}
		}
	}
\end{Verbatim}

In the preceding code, when the next vertex, \ensuremath{\mathtt{{\color{var}i}}}, is processed, \ensuremath{\mathtt{{\color{var}i}}} is coloured
\ensuremath{\mathtt{{\color{var}grey}}} and then replaced, on the stack, with its adjacent vertices.
During the next iteration, one of these vertices will be visited.

Not surprisingly, the running times of \ensuremath{\mathtt{dfs({\color{var}g},{\color{var}r})}} and \ensuremath{\mathtt{dfs2({\color{var}g},{\color{var}r})}} are the
same as that of \ensuremath{\mathtt{bfs({\color{var}g},{\color{var}r})}}:
\begin{thm}\thmlabel{dfs-graph}
  When given as input a \texttt{Graph}, \ensuremath{\mathtt{{\color{var}g}}}, that is implemented using the
  \texttt{AdjacencyLists} data structure, the \ensuremath{\mathtt{dfs({\color{var}g},{\color{var}r})}} and \ensuremath{\mathtt{dfs2({\color{var}g},{\color{var}r})}} algorithms
  each run in $O(\ensuremath{\mathtt{{\color{var}n}}}+\ensuremath{\mathtt{{\color{var}m}}})$ time.
\end{thm}

As with the breadth-first-search algorithm, there is an underlying
tree associated with each execution of depth-first-search.  When a node
$\ensuremath{\mathtt{{\color{var}i}}}\neq \ensuremath{\mathtt{{\color{var}r}}}$ goes from \ensuremath{\mathtt{{\color{var}white}}} to \ensuremath{\mathtt{{\color{var}grey}}}, this is because \ensuremath{\mathtt{dfs({\color{var}g},{\color{var}i},{\color{var}c})}}
was called recursively while processing some node \ensuremath{\mathtt{{\color{var}i}'}}.  (In the case
of \ensuremath{\mathtt{dfs2({\color{var}g},{\color{var}r})}} algorithm, \ensuremath{\mathtt{{\color{var}i}}} is one of the nodes that replaced \ensuremath{\mathtt{{\color{var}i}'}}
on the stack.)  If we think of \ensuremath{\mathtt{{\color{var}i}'}} as the parent of \ensuremath{\mathtt{{\color{var}i}}}, then we obtain
a tree rooted at \ensuremath{\mathtt{{\color{var}r}}}.  In \figref{graph-dfs}, this tree is a path from
vertex 0 to vertex 11.

An important property of the depth-first-search algorithm is the
following: Suppose that when node \ensuremath{\mathtt{{\color{var}i}}} is coloured \ensuremath{\mathtt{{\color{var}grey}}}, there exists a path
from \ensuremath{\mathtt{{\color{var}i}}} to some other node \ensuremath{\mathtt{{\color{var}j}}} that uses only white vertices.  Then \ensuremath{\mathtt{{\color{var}j}}}
will be coloured first \ensuremath{\mathtt{{\color{var}grey}}} then \ensuremath{\mathtt{{\color{var}black}}} before \ensuremath{\mathtt{{\color{var}i}}} is coloured \ensuremath{\mathtt{{\color{var}black}}}.
(This can be proven by contradiction, by considering any path $P$ from \ensuremath{\mathtt{{\color{var}i}}}
to \ensuremath{\mathtt{{\color{var}j}}}.)

One application of this property is the detection of cycles.
\index{cycle detection}%
Refer
to \figref{dfs-cycle}.  Consider some cycle, $C$, that can be reached
from \ensuremath{\mathtt{{\color{var}r}}}.  Let \ensuremath{\mathtt{{\color{var}i}}} be the first node of $C$ that is coloured \ensuremath{\mathtt{{\color{var}grey}}},
and let \ensuremath{\mathtt{{\color{var}j}}} be the node that precedes \ensuremath{\mathtt{{\color{var}i}}} on the cycle $C$.  Then,
by the above property, \ensuremath{\mathtt{{\color{var}j}}} will be coloured \ensuremath{\mathtt{{\color{var}grey}}} and the edge \ensuremath{\mathtt{({\color{var}j},{\color{var}i})}}
will be considered by the algorithm while \ensuremath{\mathtt{{\color{var}i}}} is still \ensuremath{\mathtt{{\color{var}grey}}}.  Thus,
the algorithm can conclude that there is a path, $P$, from \ensuremath{\mathtt{{\color{var}i}}} to \ensuremath{\mathtt{{\color{var}j}}}
in the depth-first-search tree and the edge \ensuremath{\mathtt{({\color{var}j},{\color{var}i})}} exists.  Therefore,
$P$ is also a cycle.

\begin{figure}
  \begin{center}
    \includegraphics[scale=0.90909]{figs/dfs-cycle}
  \end{center}
  \caption[Cycle detection]{The depth-first-search algorithm can be used to detect cycles
  in $G$. The node \ensuremath{\mathtt{{\color{var}j}}} is coloured \ensuremath{\mathtt{{\color{var}grey}}} while \ensuremath{\mathtt{{\color{var}i}}} is still \ensuremath{\mathtt{{\color{var}grey}}}.  This
  implies that there is a path, $P$, from \ensuremath{\mathtt{{\color{var}i}}} to \ensuremath{\mathtt{{\color{var}j}}} in the depth-first-search
  tree, and the edge \ensuremath{\mathtt{({\color{var}j},{\color{var}i})}} implies that $P$ is also a cycle.}
  \figlabel{dfs-cycle}
\end{figure}

\section{Discussion and Exercises}

The running times of the depth-first-search and breadth-first-search
algorithms are somewhat overstated by the Theorems~\ref{thm:bfs-graph} and
\ref{thm:dfs-graph}.  Define $\ensuremath{\mathtt{{\color{var}n}}}_{\ensuremath{\mathtt{{\color{var}r}}}}$ as the number of vertices, \ensuremath{\mathtt{{\color{var}i}}},
of $G$, for which there exists a path from \ensuremath{\mathtt{{\color{var}r}}} to \ensuremath{\mathtt{{\color{var}i}}}.  Define $\ensuremath{\mathtt{{\color{var}m}}}_\ensuremath{\mathtt{{\color{var}r}}}$
as the number of edges that have these vertices as their sources.
Then the following theorem is a more precise statement of the running
times of the breadth-first-search and depth-first-search algorithms.
(This more refined statement of the running time is useful in some of
the applications of these algorithms outlined in the exercises.)
\begin{thm}\thmlabel{graph-traversal}
  When given as input a \texttt{Graph}, \ensuremath{\mathtt{{\color{var}g}}}, that is implemented using the
  \texttt{AdjacencyLists} data structure, the \ensuremath{\mathtt{bfs({\color{var}g},{\color{var}r})}}, \ensuremath{\mathtt{dfs({\color{var}g},{\color{var}r})}} and \ensuremath{\mathtt{dfs2({\color{var}g},{\color{var}r})}}
  algorithms each run in $O(\ensuremath{\mathtt{{\color{var}n}}}_{\ensuremath{\mathtt{{\color{var}r}}}}+\ensuremath{\mathtt{{\color{var}m}}}_{\ensuremath{\mathtt{{\color{var}r}}}})$ time.
\end{thm}

Breadth-first search seems to have been discovered independently by
Moore \cite{m59} and Lee \cite{l61} in the contexts of maze exploration
and circuit routing, respectively.

Adjacency-list representations of graphs were presented by
Hopcroft and Tarjan \cite{ht73} as an alternative to the (then more
common) adjacency-matrix representation.  This representation, as well as
depth-first-search, played a major part in the celebrated Hopcroft-Tarjan
planarity testing algorithm 
\index{planarity testing}%
that can determine, in $O(\ensuremath{\mathtt{{\color{var}n}}})$ time, if
a graph can be drawn, in the plane, and in such a way that no pair of
edges cross each other \cite{ht74}.

In the following exercises, an undirected graph is one in which, for
every \ensuremath{\mathtt{{\color{var}i}}} and \ensuremath{\mathtt{{\color{var}j}}}, the edge $(\ensuremath{\mathtt{{\color{var}i}}},\ensuremath{\mathtt{{\color{var}j}}})$ is present if and only if the
edge $(\ensuremath{\mathtt{{\color{var}j}}},\ensuremath{\mathtt{{\color{var}i}}})$ is present.
\index{undirected graph}%
\index{graph!undirected}%

\begin{exc}
  Draw an adjacency list representation and an adjacency matrix
  representation of the graph in \figref{graph-example2}.
\end{exc}

\begin{figure}
  \centering{\includegraphics[scale=0.90909]{figs/graph-example2}}
  \caption{An example graph.}
  \figlabel{graph-example2}
\end{figure}

\begin{exc}
  \index{incidence matrix}%
  The \emph{incidence matrix} representation of a graph,
  $G$, is an $\ensuremath{\mathtt{{\color{var}n}}}\times\ensuremath{\mathtt{{\color{var}m}}}$ matrix, $A$, where
  \[
     A_{i,j} = \begin{cases}
        -1 & \text{if vertex $i$ the source of edge $j$} \\
        +1 & \text{if vertex $i$ the target of edge $j$} \\
        0 & \text{otherwise.}
     \end{cases}
  \]
  \begin{enumerate}
    \item Draw the incident matrix representation of the graph in
      \figref{graph-example2}.
    \item Design, analyze and implement an incidence matrix representation
      of a graph.  Be sure to analyze the space, the cost of
      \ensuremath{\mathtt{addEdge({\color{var}i},{\color{var}j})}}, \ensuremath{\mathtt{removeEdge({\color{var}i},{\color{var}j})}}, \ensuremath{\mathtt{hasEdge({\color{var}i},{\color{var}j})}}, \ensuremath{\mathtt{inEdges({\color{var}i})}},
      and \ensuremath{\mathtt{outEdges({\color{var}i})}}.
  \end{enumerate}
\end{exc}

\begin{exc}
  Illustrate an execution of the \ensuremath{\mathtt{bfs(G,0)}} and \ensuremath{\mathtt{dfs(G,0)}} on the graph, $\ensuremath{\mathtt{G}}$,
  in \figref{graph-example2}.
\end{exc}

\begin{exc}
  \index{connected graph}%
  \index{graph!connected}%
  Let $G$ be an undirected graph.  We say $G$ is \emph{connected} if,
  for every pair of vertices \ensuremath{\mathtt{{\color{var}i}}} and \ensuremath{\mathtt{{\color{var}j}}} in $G$, there is a path from
  $\ensuremath{\mathtt{{\color{var}i}}}$ to $\ensuremath{\mathtt{{\color{var}j}}}$ (since $G$ is undirected, there is also a path from \ensuremath{\mathtt{{\color{var}j}}}
  to \ensuremath{\mathtt{{\color{var}i}}}). Show how to test if $G$ is connected in $O(\ensuremath{\mathtt{{\color{var}n}}}+\ensuremath{\mathtt{{\color{var}m}}})$ time.
\end{exc}

\begin{exc}
  \index{connected components}%
  Let $G$ be an undirected graph.  A \emph{connected-component labelling}
  of $G$ partitions the vertices of $G$ into maximal sets, each of which
  forms a connected subgraph.  Show how to compute a connected component
  labelling of $G$ in $O(\ensuremath{\mathtt{{\color{var}n}}}+\ensuremath{\mathtt{{\color{var}m}}})$ time.
\end{exc}

\begin{exc}
  \index{spanning forest}%
  Let $G$ be an undirected graph.  A \emph{spanning forest} of $G$ is a
  collection of trees, one per component, whose edges are edges of $G$
  and whose vertices contain all vertices of $G$.  Show how to compute
  a spanning forest of of $G$ in $O(\ensuremath{\mathtt{{\color{var}n}}}+\ensuremath{\mathtt{{\color{var}m}}})$ time.
\end{exc}

\begin{exc}
  \index{strongly-connected graph}%
  \index{graph!strongly-connected}%
  We say that a graph $G$ is \emph{strongly-connected} if, for every
  pair of vertices \ensuremath{\mathtt{{\color{var}i}}} and \ensuremath{\mathtt{{\color{var}j}}} in $G$, there is a path from $\ensuremath{\mathtt{{\color{var}i}}}$ to
  $\ensuremath{\mathtt{{\color{var}j}}}$. Show how to test if $G$ is strongly-connected in $O(\ensuremath{\mathtt{{\color{var}n}}}+\ensuremath{\mathtt{{\color{var}m}}})$
  time.
\end{exc}

\begin{exc}
  Given a graph $G=(V,E)$ and some special vertex $\ensuremath{\mathtt{{\color{var}r}}}\in V$, show how
  to compute the length of the shortest path from $\ensuremath{\mathtt{{\color{var}r}}}$ to \ensuremath{\mathtt{{\color{var}i}}} for every
  vertex $\ensuremath{\mathtt{{\color{var}i}}}\in V$.
\end{exc}

\begin{exc}
  Give a (simple) example where the \ensuremath{\mathtt{dfs({\color{var}g},{\color{var}r})}} code visits the nodes of a
  graph in an order that is different from that of the \ensuremath{\mathtt{dfs2({\color{var}g},{\color{var}r})}} code.
  Write a version of \ensuremath{\mathtt{dfs2({\color{var}g},{\color{var}r})}} that always visits nodes in exactly
  the same order as \ensuremath{\mathtt{dfs({\color{var}g},{\color{var}r})}}.  (Hint: Just start tracing the execution
  of each algorithm on some graph where \ensuremath{\mathtt{{\color{var}r}}} is the source of more than
  1 edge.)
\end{exc}

\begin{exc}
  \index{universal sink}%
  \index{celebrity|see{universal sink}}%
  A \emph{universal sink} in a graph $G$ is a vertex that is the target
  of $\ensuremath{\mathtt{{\color{var}n}}}-1$ edges and the source of no edges.\footnote{A universal sink,
  \ensuremath{\mathtt{{\color{var}v}}}, is also sometimes called a \emph{celebrity}: Everyone in the room
  recognizes \ensuremath{\mathtt{{\color{var}v}}}, but \ensuremath{\mathtt{{\color{var}v}}} doesn't recognize anyone else in the room.}
  Design and implement an algorithm that tests if a graph $G$, represented
  as an \texttt{AdjacencyMatrix}, has a universal sink.  Your algorithm should
  run in $O(\ensuremath{\mathtt{{\color{var}n}}})$ time.
\end{exc}




\clearpage

\clearpage
\chapter{Data Structures for Integers}

In this chapter, we return to the problem of implementing an \texttt{SSet}.
The difference now is that we assume the elements stored in the \texttt{SSet} are
\ensuremath{\mathtt{{\color{var}w}}}-bit integers.  That is, we want to implement \ensuremath{\mathtt{add({\color{var}x})}}, \ensuremath{\mathtt{remove({\color{var}x})}},
and \ensuremath{\mathtt{find({\color{var}x})}} where $\ensuremath{\mathtt{{\color{var}x}}}\in\{0,\ldots,2^{\ensuremath{\mathtt{{\color{var}w}}}}-1\}$.  It is not too hard
to think of plenty of applications where the data---or at least the key
that we use for sorting the data---is an integer.

We will discuss three data structures, each building on the ideas of
the previous.  The first structure, the \texttt{BinaryTrie} performs all three
\texttt{SSet} operations in $O(\ensuremath{\mathtt{{\color{var}w}}})$ time. This is not very impressive, since
any subset of $\{0,\ldots,2^{\ensuremath{\mathtt{{\color{var}w}}}}-1\}$ has size $\ensuremath{\mathtt{{\color{var}n}}}\le 2^{\ensuremath{\mathtt{{\color{var}w}}}}$, so that
$\log \ensuremath{\mathtt{{\color{var}n}}} \le \ensuremath{\mathtt{{\color{var}w}}}$.  All the other \texttt{SSet} implementations discussed in
this book perform all operations in $O(\log \ensuremath{\mathtt{{\color{var}n}}})$ time so they are all
at least as fast as a \texttt{BinaryTrie}.

The second structure, the \texttt{XFastTrie}, speeds up the search in a
\texttt{BinaryTrie} by using hashing.  With this speedup, the \ensuremath{\mathtt{find({\color{var}x})}}
operation runs in $O(\log \ensuremath{\mathtt{{\color{var}w}}})$ time.  However, \ensuremath{\mathtt{add({\color{var}x})}} and \ensuremath{\mathtt{remove({\color{var}x})}}
operations in an \texttt{XFastTrie} still take $O(\ensuremath{\mathtt{{\color{var}w}}})$ time and the space used
by an \texttt{XFastTrie} is $O(\ensuremath{\mathtt{{\color{var}n}}}\cdot\ensuremath{\mathtt{{\color{var}w}}})$.

The third data structure, the \texttt{YFastTrie}, uses an \texttt{XFastTrie} to store
only a sample of roughly one out of every $\ensuremath{\mathtt{{\color{var}w}}}$ elements and stores the
remaining elements in a standard \texttt{SSet} structure.  This trick reduces the
running time of \ensuremath{\mathtt{add({\color{var}x})}} and \ensuremath{\mathtt{remove({\color{var}x})}} to $O(\log \ensuremath{\mathtt{{\color{var}w}}})$ and decreases
the space to $O(\ensuremath{\mathtt{{\color{var}n}}})$.

The implementations used as examples in this chapter can store any type of
data, as long as an integer can be associated with it.  In the code samples,
the variable \ensuremath{\mathtt{{\color{var}ix}}} is always the integer value associated with \ensuremath{\mathtt{{\color{var}x}}}, and the method \javaonly{\ensuremath{\mathtt{{\color{var}in}.}}}\ensuremath{\mathtt{intValue({\color{var}x})}} converts \ensuremath{\mathtt{{\color{var}x}}} to its associated integer. In
the text, however, we will simply treat \ensuremath{\mathtt{{\color{var}x}}} as if it is an integer.

\section{\texttt{BinaryTrie}: A digital search tree}
\seclabel{binarytrie}

\index{BinaryTrie@\texttt{BinaryTrie}}%
A \texttt{BinaryTrie} encodes a set of \ensuremath{\mathtt{{\color{var}w}}} bit integers in a binary tree.
All leaves in the tree have depth \ensuremath{\mathtt{{\color{var}w}}} and each integer is encoded as a
root-to-leaf path.  The path for the integer \ensuremath{\mathtt{{\color{var}x}}} turns left at level \ensuremath{\mathtt{{\color{var}i}}}
if the \ensuremath{\mathtt{{\color{var}i}}}th most significant bit of \ensuremath{\mathtt{{\color{var}x}}} is a 0 and turns right if it
is a 1.  \figref{binarytrie-ex} shows an example for the case $\ensuremath{\mathtt{{\color{var}w}}}=4$,
in which the trie stores the integers 3(0011), 9(1001), 12(1100),
and 13(1101).
\begin{figure}
  \begin{center}
    \includegraphics[width=\ScaleIfNeeded]{figs/binarytrie-ex-1}
  \end{center}
  \caption{The integers stored in a binary trie are encoded as
    root-to-leaf paths.}
  \figlabel{binarytrie-ex}
\end{figure}

Because the search path
\index{search path!in a \texttt{BinaryTrie}}%
for a value \ensuremath{\mathtt{{\color{var}x}}} depends on the bits of \ensuremath{\mathtt{{\color{var}x}}}, it will
be helpful to name the children of a node, \ensuremath{\mathtt{{\color{var}u}}}, \ensuremath{\mathtt{{\color{var}u}.{\color{var}child}[0]}} (\ensuremath{\mathtt{{\color{var}left}}})
and \ensuremath{\mathtt{{\color{var}u}.{\color{var}child}[1]}} (\ensuremath{\mathtt{{\color{var}right}}}).  These child pointers will actually serve
double-duty.  Since the leaves in a binary trie have no children, the
pointers are used to string the leaves together into a doubly-linked list.
For a leaf in the binary trie \ensuremath{\mathtt{{\color{var}u}.{\color{var}child}[0]}} (\ensuremath{\mathtt{{\color{var}prev}}}) is the node that
comes before \ensuremath{\mathtt{{\color{var}u}}} in the list and \ensuremath{\mathtt{{\color{var}u}.{\color{var}child}[1]}} (\ensuremath{\mathtt{{\color{var}next}}}) is the node that
follows \ensuremath{\mathtt{{\color{var}u}}} in the list.  A special node, \ensuremath{\mathtt{{\color{var}dummy}}}, is used both before
the first node and after the last node in the list (see \secref{dllist}).
\cpponly{In the code samples, \ensuremath{\mathtt{{\color{var}u}.{\color{var}child}[0]}}, \ensuremath{\mathtt{{\color{var}u}.{\color{var}left}}}, and \ensuremath{\mathtt{{\color{var}u}.{\color{var}prev}}} refer to the same field in the node \ensuremath{\mathtt{{\color{var}u}}}, as do \ensuremath{\mathtt{{\color{var}u}.{\color{var}child}[1]}}, \ensuremath{\mathtt{{\color{var}u}.{\color{var}right}}}, and \ensuremath{\mathtt{{\color{var}u}.{\color{var}next}}}.}

Each node, \ensuremath{\mathtt{{\color{var}u}}}, also contains an additional pointer \ensuremath{\mathtt{{\color{var}u}.{\color{var}jump}}}.  If \ensuremath{\mathtt{{\color{var}u}}}'s
left child is missing, then \ensuremath{\mathtt{{\color{var}u}.{\color{var}jump}}} points to the smallest leaf in
\ensuremath{\mathtt{{\color{var}u}}}'s subtree.  If \ensuremath{\mathtt{{\color{var}u}}}'s right child is missing, then \ensuremath{\mathtt{{\color{var}u}.{\color{var}jump}}} points
to the largest leaf in \ensuremath{\mathtt{{\color{var}u}}}'s subtree.  An example of a \texttt{BinaryTrie},
showing \ensuremath{\mathtt{{\color{var}jump}}} pointers and the doubly-linked list at the leaves, is
shown in \figref{binarytrie-ex2}.

\begin{figure}
  \begin{center}
    \includegraphics[width=\ScaleIfNeeded]{figs/binarytrie-ex-2}
  \end{center}
  \caption[A BinaryTrie]{A \texttt{BinaryTrie} with \ensuremath{\mathtt{{\color{var}jump}}} pointers shown as curved dashed
  edges.}
  \figlabel{binarytrie-ex2}
\end{figure}



The \ensuremath{\mathtt{find({\color{var}x})}} operation in a \texttt{BinaryTrie} is fairly straightforward.
We try to follow the search path for \ensuremath{\mathtt{{\color{var}x}}} in the trie.  If we reach a leaf,
then we have found \ensuremath{\mathtt{{\color{var}x}}}.  If we reach a node \ensuremath{\mathtt{{\color{var}u}}} where we cannot proceed
(because \ensuremath{\mathtt{{\color{var}u}}} is missing a child), then we follow \ensuremath{\mathtt{{\color{var}u}.{\color{var}jump}}}, which takes us
either to the smallest leaf larger than \ensuremath{\mathtt{{\color{var}x}}} or the largest leaf smaller than
\ensuremath{\mathtt{{\color{var}x}}}. Which of these two cases occurs depends on whether \ensuremath{\mathtt{{\color{var}u}}} is missing
its left or right child, respectively.  In the former case (\ensuremath{\mathtt{{\color{var}u}}} is missing
its left child), we have found the node we want.  In the latter case (\ensuremath{\mathtt{{\color{var}u}}}
is missing its right child), we can use the linked list to reach the node
we want. Each of these cases is illustrated in \figref{binarytrie-find}.


 \ \vspace{.5ex}\begin{Verbatim}[gobble=1,tabsize=2,frame=single,framerule=1.2pt,commandchars=\\@\$,label=\texttt{BinaryTrie},labelposition=topline]
	@\color@keyword$T$ find(@\color@keyword$T$ @\color@var$x$) {
		@\color@keyword$int$ @\color@var$i$, @\color@var$c$ = 0, @\color@var$ix$ = @\color@var$it$.intValue(@\color@var$x$);
		Node @\color@var$u$ = @\color@var$r$;
		@\color@keyword$for$ (@\color@var$i$ = 0; @\color@var$i$ < @\color@var$w$; @\color@var$i$++) {
			@\color@var$c$ = (@\color@var$ix$ >>> @\color@var$w$-@\color@var$i$-1) & 1;
			@\color@keyword$if$ (@\color@var$u$.@\color@var$child$[@\color@var$c$] == @\color@keyword$null$) @\color@var$break$;
			@\color@var$u$ = @\color@var$u$.@\color@var$child$[@\color@var$c$];
		}
		@\color@keyword$if$ (@\color@var$i$ == @\color@var$w$) @\color@var$return$ @\color@var$u$.@\color@var$x$;  @\color@comment$// found it$
		@\color@var$u$ = (@\color@var$c$ == 0) ? @\color@var$u$.@\color@var$jump$ : @\color@var$u$.@\color@var$jump$.@\color@var$child$[@\color@var$next$]; 
		@\color@var$return$ @\color@var$u$ == @\color@var$dummy$ ? @\color@keyword$null$ : @\color@var$u$.@\color@var$x$;
	}
\end{Verbatim}

\begin{figure}
  \begin{center}
    \includegraphics[width=\ScaleIfNeeded]{figs/binarytrie-ex-3}
  \end{center}
  \caption[Search paths in a BinaryTrie]{The paths followed by \ensuremath{\mathtt{find(5)}} and \ensuremath{\mathtt{find(8)}}.}
  \figlabel{binarytrie-find}
\end{figure}
The running-time of the \ensuremath{\mathtt{find({\color{var}x})}} method is dominated by the time it
takes to follow a root-to-leaf path, so it runs in $O(\ensuremath{\mathtt{{\color{var}w}}})$ time.

The \ensuremath{\mathtt{add({\color{var}x})}} operation in a \texttt{BinaryTrie} is also fairly straightforward,
but has a lot of work to do:
\begin{enumerate}
  \item It follows the search path for \ensuremath{\mathtt{{\color{var}x}}} until reaching a node \ensuremath{\mathtt{{\color{var}u}}}
    where it can no longer proceed.
  \item It creates the remainder of the search path from \ensuremath{\mathtt{{\color{var}u}}} to a leaf
    that contains \ensuremath{\mathtt{{\color{var}x}}}.
  \item It adds the node, \ensuremath{\mathtt{{\color{var}u}'}}, containing \ensuremath{\mathtt{{\color{var}x}}} to the linked list
    of leaves (it has access to the predecessor, \ensuremath{\mathtt{{\color{var}pred}}}, of \ensuremath{\mathtt{{\color{var}u}'}} in
    the linked list from the \ensuremath{\mathtt{{\color{var}jump}}} pointer of the last node, \ensuremath{\mathtt{{\color{var}u}}},
    encountered during step~1.)
  \item It walks back up the search path for \ensuremath{\mathtt{{\color{var}x}}} adjusting \ensuremath{\mathtt{{\color{var}jump}}}
    pointers at the nodes whose \ensuremath{\mathtt{{\color{var}jump}}} pointer should now point to \ensuremath{\mathtt{{\color{var}x}}}.
\end{enumerate}
An addition is illustrated in \figref{binarytrie-add}.
\begin{figure}
  \begin{center}
    \includegraphics[width=\ScaleIfNeeded]{figs/binarytrie-add}
  \end{center}
  \caption[Adding to a BinaryTrie]{Adding the values 2 and 15 to the \texttt{BinaryTrie} in
  \figref{binarytrie-ex2}.}
  \figlabel{binarytrie-add}
\end{figure}


 \ \vspace{.5ex}\begin{Verbatim}[gobble=1,tabsize=2,frame=single,framerule=1.2pt,commandchars=\\@\$,label=\texttt{BinaryTrie},labelposition=topline]
	@\color@var$boolean$ add(@\color@keyword$T$ @\color@var$x$) {
		@\color@keyword$int$ @\color@var$i$, @\color@var$c$ = 0, @\color@var$ix$ = @\color@var$it$.intValue(@\color@var$x$);
		Node @\color@var$u$ = @\color@var$r$;
		@\color@comment$// 1 - search for ix until falling out of the trie$
		@\color@keyword$for$ (@\color@var$i$ = 0; @\color@var$i$ < @\color@var$w$; @\color@var$i$++) {
			@\color@var$c$ = (@\color@var$ix$ >>> @\color@var$w$-@\color@var$i$-1) & 1;
			@\color@keyword$if$ (@\color@var$u$.@\color@var$child$[@\color@var$c$] == @\color@keyword$null$) @\color@var$break$;
			@\color@var$u$ = @\color@var$u$.@\color@var$child$[@\color@var$c$];
		}		
		@\color@keyword$if$ (@\color@var$i$ == @\color@var$w$) @\color@var$return$ @\color@var$false$; @\color@comment$// already contains x - abort$
		Node @\color@var$pred$ = (@\color@var$c$ == @\color@var$right$) ? @\color@var$u$.@\color@var$jump$ : @\color@var$u$.@\color@var$jump$.@\color@var$child$[0];
		@\color@var$u$.@\color@var$jump$ = @\color@keyword$null$;  @\color@comment$// u will have two children shortly$
		@\color@comment$// 2 - add path to ix$
		@\color@keyword$for$ (; @\color@var$i$ < @\color@var$w$; @\color@var$i$++) {
			@\color@var$c$ = (@\color@var$ix$ >>> @\color@var$w$-@\color@var$i$-1) & 1;
			@\color@var$u$.@\color@var$child$[@\color@var$c$] = newNode();
			@\color@var$u$.@\color@var$child$[@\color@var$c$].@\color@var$parent$ = @\color@var$u$;
			@\color@var$u$ = @\color@var$u$.@\color@var$child$[@\color@var$c$];
		}
		@\color@var$u$.@\color@var$x$ = @\color@var$x$;
		@\color@comment$// 3 - add u to linked list$
		@\color@var$u$.@\color@var$child$[@\color@var$prev$] = @\color@var$pred$;
		@\color@var$u$.@\color@var$child$[@\color@var$next$] = @\color@var$pred$.@\color@var$child$[@\color@var$next$];
		@\color@var$u$.@\color@var$child$[@\color@var$prev$].@\color@var$child$[@\color@var$next$] = @\color@var$u$;
		@\color@var$u$.@\color@var$child$[@\color@var$next$].@\color@var$child$[@\color@var$prev$] = @\color@var$u$;
		@\color@comment$// 4 - walk back up, updating jump pointers$
		Node @\color@var$v$ = @\color@var$u$.@\color@var$parent$;
		@\color@keyword$while$ (@\color@var$v$ != @\color@keyword$null$) {
			@\color@keyword$if$ ((@\color@var$v$.@\color@var$child$[@\color@var$left$] == @\color@keyword$null$ 
					&& (@\color@var$v$.@\color@var$jump$ == @\color@keyword$null$ || @\color@var$it$.intValue(@\color@var$v$.@\color@var$jump$.@\color@var$x$) > @\color@var$ix$))
			|| (@\color@var$v$.@\color@var$child$[@\color@var$right$] == @\color@keyword$null$ 
					&& (@\color@var$v$.@\color@var$jump$ == @\color@keyword$null$ || @\color@var$it$.intValue(@\color@var$v$.@\color@var$jump$.@\color@var$x$) < @\color@var$ix$)))
				@\color@var$v$.@\color@var$jump$ = @\color@var$u$;
			@\color@var$v$ = @\color@var$v$.@\color@var$parent$;
		}
		@\color@var$n$++;
		@\color@var$return$ @\color@var$true$;
	}
\end{Verbatim}

This method performs one walk down the search path for \ensuremath{\mathtt{{\color{var}x}}} and one walk
back up.  Each step of these walks takes constant time, so the \ensuremath{\mathtt{add({\color{var}x})}}
method runs in $O(\ensuremath{\mathtt{{\color{var}w}}})$ time.


The \ensuremath{\mathtt{remove({\color{var}x})}} operation undoes the work of \ensuremath{\mathtt{add({\color{var}x})}}.  Like \ensuremath{\mathtt{add({\color{var}x})}},
it has a lot of work to do:
\begin{enumerate}
  \item It follows the search path for \ensuremath{\mathtt{{\color{var}x}}} until reaching the leaf, \ensuremath{\mathtt{{\color{var}u}}},
  containing \ensuremath{\mathtt{{\color{var}x}}}.
  \item It removes \ensuremath{\mathtt{{\color{var}u}}} from the doubly-linked list.
  \item It deletes \ensuremath{\mathtt{{\color{var}u}}} and then walks back up the search path for \ensuremath{\mathtt{{\color{var}x}}}
  deleting nodes until reaching a node \ensuremath{\mathtt{{\color{var}v}}} that has a child that is not
  on the search path for \ensuremath{\mathtt{{\color{var}x}}}.
  \item It walks upwards from \ensuremath{\mathtt{{\color{var}v}}} to the root updating any \ensuremath{\mathtt{{\color{var}jump}}} pointers
  that point to \ensuremath{\mathtt{{\color{var}u}}}.
\end{enumerate}
A removal is illustrated in \figref{binarytrie-remove}.
\begin{figure}
  \begin{center}
    \includegraphics[scale=0.90909]{figs/binarytrie-remove}
  \end{center}
  \caption[Removing from a BinaryTrie]{Removing the value 9 from the \texttt{BinaryTrie} in
  \figref{binarytrie-ex2}.}
  \figlabel{binarytrie-remove}
\end{figure}


 \ \vspace{.5ex}\begin{Verbatim}[gobble=1,tabsize=2,frame=single,framerule=1.2pt,commandchars=\\@\$,label=\texttt{BinaryTrie},labelposition=topline]
	@\color@var$boolean$ remove(@\color@keyword$T$ @\color@var$x$) {
		@\color@comment$// 1 - find leaf, u, containing x$
		@\color@keyword$int$ @\color@var$i$, @\color@var$c$, @\color@var$ix$ = @\color@var$it$.intValue(@\color@var$x$);
		Node @\color@var$u$ = @\color@var$r$;
		@\color@keyword$for$ (@\color@var$i$ = 0; @\color@var$i$ < @\color@var$w$; @\color@var$i$++) {
			@\color@var$c$ = (@\color@var$ix$ >>> @\color@var$w$-@\color@var$i$-1) & 1;
			@\color@keyword$if$ (@\color@var$u$.@\color@var$child$[@\color@var$c$] == @\color@keyword$null$) @\color@var$return$ @\color@var$false$;
			@\color@var$u$ = @\color@var$u$.@\color@var$child$[@\color@var$c$];
		}
		@\color@comment$// 2 - remove u from linked list$
		@\color@var$u$.@\color@var$child$[@\color@var$prev$].@\color@var$child$[@\color@var$next$] = @\color@var$u$.@\color@var$child$[@\color@var$next$];
		@\color@var$u$.@\color@var$child$[@\color@var$next$].@\color@var$child$[@\color@var$prev$] = @\color@var$u$.@\color@var$child$[@\color@var$prev$];
		Node @\color@var$v$ = @\color@var$u$;
		@\color@comment$// 3 - delete nodes on path to u$
		@\color@keyword$for$ (@\color@var$i$ = @\color@var$w$-1; @\color@var$i$ >= 0; @\color@var$i$--) {
			@\color@var$c$ = (@\color@var$ix$ >>> @\color@var$w$-@\color@var$i$-1) & 1;
			@\color@var$v$ = @\color@var$v$.@\color@var$parent$;
			@\color@var$v$.@\color@var$child$[@\color@var$c$] = @\color@keyword$null$;
			@\color@keyword$if$ (@\color@var$v$.@\color@var$child$[1-@\color@var$c$] != @\color@keyword$null$) @\color@var$break$;
		}
		@\color@comment$// 4 - update jump pointers$
		@\color@var$c$ = (@\color@var$ix$ >>> @\color@var$w$-@\color@var$i$-1) & 1;
		@\color@var$v$.@\color@var$jump$ = @\color@var$u$.@\color@var$child$[1-@\color@var$c$];
		@\color@var$v$ = @\color@var$v$.@\color@var$parent$;
		@\color@var$i$--;
		@\color@keyword$for$ (; @\color@var$i$ >= 0; @\color@var$i$--) {
			@\color@var$c$ = (@\color@var$ix$ >>> @\color@var$w$-@\color@var$i$-1) & 1;
			@\color@keyword$if$ (@\color@var$v$.@\color@var$jump$ == @\color@var$u$) 
				@\color@var$v$.@\color@var$jump$ = @\color@var$u$.@\color@var$child$[1-@\color@var$c$];
			@\color@var$v$ = @\color@var$v$.@\color@var$parent$;
		}
		@\color@var$n$--;
		@\color@var$return$ @\color@var$true$;
	}
\end{Verbatim}


\begin{thm}
A \texttt{BinaryTrie} implements the \texttt{SSet} interface for \ensuremath{\mathtt{{\color{var}w}}}-bit integers. A
\texttt{BinaryTrie} supports the operations \ensuremath{\mathtt{add({\color{var}x})}}, \ensuremath{\mathtt{remove({\color{var}x})}}, and \ensuremath{\mathtt{find({\color{var}x})}}
in $O(\ensuremath{\mathtt{{\color{var}w}}})$ time per operation.  The space used by a \texttt{BinaryTrie} that
stores \ensuremath{\mathtt{{\color{var}n}}} values is $O(\ensuremath{\mathtt{{\color{var}n}}}\cdot\ensuremath{\mathtt{{\color{var}w}}})$.
\end{thm}

\section{\texttt{XFastTrie}: Searching in Doubly-Logarithmic Time}
\seclabel{xfast}

\index{XFastTrie@\texttt{XFastTrie}}%
The performance of the \texttt{BinaryTrie} structure is not very impressive.
The number of elements, \ensuremath{\mathtt{{\color{var}n}}}, stored in the structure is at most $2^{\ensuremath{\mathtt{{\color{var}w}}}}$,
so $\log \ensuremath{\mathtt{{\color{var}n}}}\le \ensuremath{\mathtt{{\color{var}w}}}$.  In other words, any of the comparison-based \texttt{SSet}
structures described in other parts of this book are at least as efficient
as a \texttt{BinaryTrie}, and are not restricted to only storing integers.

Next we describe the \texttt{XFastTrie}, which is just a \texttt{BinaryTrie} with
\ensuremath{\mathtt{{\color{var}w}+1}} hash tables---one for each level of the trie. These hash tables
are used to speed up the \ensuremath{\mathtt{find({\color{var}x})}} operation to $O(\log \ensuremath{\mathtt{{\color{var}w}}})$ time.
Recall that the \ensuremath{\mathtt{find({\color{var}x})}} operation in a \texttt{BinaryTrie} is almost complete
once we reach a node, \ensuremath{\mathtt{{\color{var}u}}}, where the search path for \ensuremath{\mathtt{{\color{var}x}}} would like to
proceed to \ensuremath{\mathtt{{\color{var}u}.{\color{var}right}}} (or \ensuremath{\mathtt{{\color{var}u}.{\color{var}left}}}) but \ensuremath{\mathtt{{\color{var}u}}} has no right (respectively,
left) child.  At this point, the search uses \ensuremath{\mathtt{{\color{var}u}.{\color{var}jump}}} to jump to a leaf,
\ensuremath{\mathtt{{\color{var}v}}}, of the \texttt{BinaryTrie} and either return \ensuremath{\mathtt{{\color{var}v}}} or its successor in the
linked list of leaves.  An \texttt{XFastTrie} speeds up the search process by
using binary search
\index{binary search}%
on the levels of the trie to locate the node \ensuremath{\mathtt{{\color{var}u}}}.

To use binary search, we need a way to determine if the node \ensuremath{\mathtt{{\color{var}u}}} we are
looking for is above a particular level, \ensuremath{\mathtt{{\color{var}i}}}, of if \ensuremath{\mathtt{{\color{var}u}}} is at or below
level \ensuremath{\mathtt{{\color{var}i}}}.  This information is given by the highest-order \ensuremath{\mathtt{{\color{var}i}}} bits
in the binary representation of \ensuremath{\mathtt{{\color{var}x}}}; these bits determine the search
path that \ensuremath{\mathtt{{\color{var}x}}} takes from the root to level \ensuremath{\mathtt{{\color{var}i}}}.   For an example,
refer to \figref{xfast-path}; in this figure the last node, \ensuremath{\mathtt{{\color{var}u}}}, on
search path for 14 (whose binary representation is 1110) is the node
labelled $11{\star\star}$ at level 2 because there is no node labelled
$111{\star}$ at level 3.  Thus, we can label each node at level \ensuremath{\mathtt{{\color{var}i}}}
with an \ensuremath{\mathtt{{\color{var}i}}}-bit integer.  Then, the node \ensuremath{\mathtt{{\color{var}u}}} we are searching for would
be at or below level \ensuremath{\mathtt{{\color{var}i}}} if and only if there is a node at level \ensuremath{\mathtt{{\color{var}i}}}
whose label matches the highest-order \ensuremath{\mathtt{{\color{var}i}}} bits of \ensuremath{\mathtt{{\color{var}x}}}.

\begin{figure}
  \begin{center}
    \includegraphics[scale=0.90909]{figs/xfast-path}
  \end{center}
  \caption{Since there is no node labelled $111\star$, the search path
    for 14 (1110) ends at the node labelled $11{\star\star}$ .}
  \figlabel{xfast-path}
\end{figure}

In an \texttt{XFastTrie}, we store, for each $\ensuremath{\mathtt{{\color{var}i}}}\in\{0,\ldots,\ensuremath{\mathtt{{\color{var}w}}}\}$, all
the nodes at level \ensuremath{\mathtt{{\color{var}i}}} in a \texttt{USet}, \ensuremath{\mathtt{{\color{var}t}[{\color{var}i}]}}, that is implemented as a
hash table (\chapref{hashing}).  Using this \texttt{USet} allows us to check
in constant expected time if there is a node at level \ensuremath{\mathtt{{\color{var}i}}} whose label
matches the highest-order \ensuremath{\mathtt{{\color{var}i}}} bits of \ensuremath{\mathtt{{\color{var}x}}}.  In fact, we can even find
this node using
\javaonly{\ensuremath{\mathtt{{\color{var}t}[{\color{var}i}].find({\color{var}x}\text{\ttfamily >>>}({\color{var}w}-{\color{var}i}))}}}%
\cpponly{\ensuremath{\mathtt{{\color{var}t}[{\color{var}i}].find({\color{var}x}\text{\ttfamily >>}({\color{var}w}-{\color{var}i}))}}}%
\pcodeonly{\ensuremath{\mathtt{{\color{var}t}[{\color{var}i}].find({\color{var}x}\text{\ttfamily >>}({\color{var}w}-{\color{var}i}))}}}%

The hash tables $\ensuremath{\mathtt{{\color{var}t}[0]}},\ldots,\ensuremath{\mathtt{{\color{var}t}[{\color{var}w}]}}$ allow us to use binary search
to find \ensuremath{\mathtt{{\color{var}u}}}.  Initially, we know that \ensuremath{\mathtt{{\color{var}u}}} is at some level \ensuremath{\mathtt{{\color{var}i}}} with
$0\le \ensuremath{\mathtt{{\color{var}i}}}< \ensuremath{\mathtt{{\color{var}w}}}+1$. We therefore initialize $\ensuremath{\mathtt{{\color{var}l}}}=0$ and $\ensuremath{\mathtt{{\color{var}h}}}=\ensuremath{\mathtt{{\color{var}w}}}+1$
and repeatedly look at the hash table \ensuremath{\mathtt{{\color{var}t}[{\color{var}i}]}}, where $\ensuremath{\mathtt{{\color{var}i}}}=\lfloor
(\ensuremath{\mathtt{{\color{var}l}+{\color{var}h}}})/2\rfloor$.  If $\ensuremath{\mathtt{{\color{var}t}[{\color{var}i}]}}$ contains a node whose label matches
\ensuremath{\mathtt{{\color{var}x}}}'s highest-order \ensuremath{\mathtt{{\color{var}i}}} bits then we set \ensuremath{\mathtt{{\color{var}l}={\color{var}i}}} (\ensuremath{\mathtt{{\color{var}u}}} is at or below level
\ensuremath{\mathtt{{\color{var}i}}}); otherwise we set \ensuremath{\mathtt{{\color{var}h}={\color{var}i}}} (\ensuremath{\mathtt{{\color{var}u}}} is above level \ensuremath{\mathtt{{\color{var}i}}}).  This process
terminates when $\ensuremath{\mathtt{{\color{var}h}-{\color{var}l}}}\le 1$, in which case we determine that \ensuremath{\mathtt{{\color{var}u}}} is
at level \ensuremath{\mathtt{{\color{var}l}}}.  We then complete the \ensuremath{\mathtt{find({\color{var}x})}} operation using \ensuremath{\mathtt{{\color{var}u}.{\color{var}jump}}}
and the doubly-linked list of leaves.


 \ \vspace{.5ex}\begin{Verbatim}[gobble=1,tabsize=2,frame=single,framerule=1.2pt,commandchars=\\@\$,label=\texttt{XFastTrie},labelposition=topline]
	@\color@keyword$T$ find(@\color@keyword$T$ @\color@var$x$) {
		@\color@keyword$int$ @\color@var$l$ = 0, @\color@var$h$ = @\color@var$w$+1, @\color@var$ix$ = @\color@var$it$.intValue(@\color@var$x$);
		Node @\color@var$v$, @\color@var$u$ = @\color@var$r$, @\color@var$q$ = newNode();
		@\color@keyword$while$ (@\color@var$h$-@\color@var$l$ > 1) {
			@\color@keyword$int$ @\color@var$i$ = (@\color@var$l$+@\color@var$h$)/2;
			@\color@var$q$.@\color@var$prefix$ = @\color@var$ix$ >>> @\color@var$w$-@\color@var$i$;
			@\color@keyword$if$ ((@\color@var$v$ = @\color@var$t$[@\color@var$i$].find(@\color@var$q$)) == @\color@keyword$null$) {
				@\color@var$h$ = @\color@var$i$;
			} @\color@keyword$else$ {
				@\color@var$u$ = @\color@var$v$;
				@\color@var$l$ = @\color@var$i$;
			}
		}
		@\color@keyword$if$ (@\color@var$l$ == @\color@var$w$) @\color@var$return$ @\color@var$u$.@\color@var$x$;
		Node @\color@var$pred$ = (((@\color@var$ix$ >>> @\color@var$w$-@\color@var$l$-1) & 1) == 1) 
			     ? @\color@var$u$.@\color@var$jump$ : @\color@var$u$.@\color@var$jump$.@\color@var$child$[0];
		return (@\color@var$pred$.@\color@var$child$[@\color@var$next$] == @\color@var$dummy$) 
		             ? @\color@keyword$null$ : @\color@var$pred$.@\color@var$child$[@\color@var$next$].@\color@var$x$;
	}
\end{Verbatim}

Each iteration of the \ensuremath{\mathtt{{\color{keyword}while}}} loop in the above method decreases \ensuremath{\mathtt{{\color{var}h}-{\color{var}l}}}
by roughly a factor of two, so this loop finds \ensuremath{\mathtt{{\color{var}u}}} after $O(\log \ensuremath{\mathtt{{\color{var}w}}})$
iterations.  Each iteration performs a constant amount of work and one
\ensuremath{\mathtt{find({\color{var}x})}} operation in a \texttt{USet}, which takes a constant expected amount
of time.  The remaining work takes only constant time, so the \ensuremath{\mathtt{find({\color{var}x})}}
method in an \texttt{XFastTrie} takes only $O(\log\ensuremath{\mathtt{{\color{var}w}}})$ expected time.

The \ensuremath{\mathtt{add({\color{var}x})}} and \ensuremath{\mathtt{remove({\color{var}x})}} methods for an \texttt{XFastTrie} are almost
identical to the same methods in a \texttt{BinaryTrie}.  The only modifications
are for managing the hash tables \ensuremath{\mathtt{{\color{var}t}[0]}},\ldots,\ensuremath{\mathtt{{\color{var}t}[{\color{var}w}]}}.  During the
\ensuremath{\mathtt{add({\color{var}x})}} operation, when a new node is created at level \ensuremath{\mathtt{{\color{var}i}}}, this node
is added to \ensuremath{\mathtt{{\color{var}t}[{\color{var}i}]}}.  During a \ensuremath{\mathtt{remove({\color{var}x})}} operation, when a node is
removed form level \ensuremath{\mathtt{{\color{var}i}}}, this node is removed from \ensuremath{\mathtt{{\color{var}t}[{\color{var}i}]}}.  Since adding
and removing from a hash table take constant expected time, this does
not increase the running times of \ensuremath{\mathtt{add({\color{var}x})}} and \ensuremath{\mathtt{remove({\color{var}x})}} by more than
a constant factor. We omit a code listing for \ensuremath{\mathtt{add({\color{var}x})}} and \ensuremath{\mathtt{remove({\color{var}x})}}
since the code is almost identical to the (long) code listing already
provided for the same methods in a \texttt{BinaryTrie}.

The following theorem summarizes the performance of an \texttt{XFastTrie}:

\begin{thm}
An \texttt{XFastTrie} implements the \texttt{SSet} interface for \ensuremath{\mathtt{{\color{var}w}}}-bit integers. An
\texttt{XFastTrie} supports the operations
\begin{itemize}
\item \ensuremath{\mathtt{add({\color{var}x})}} and \ensuremath{\mathtt{remove({\color{var}x})}} in $O(\ensuremath{\mathtt{{\color{var}w}}})$ expected time per operation and 
\item \ensuremath{\mathtt{find({\color{var}x})}} in $O(\log \ensuremath{\mathtt{{\color{var}w}}})$ expected time per operation.
\end{itemize}
The space used by an \texttt{XFastTrie} that
stores \ensuremath{\mathtt{{\color{var}n}}} values is $O(\ensuremath{\mathtt{{\color{var}n}}}\cdot\ensuremath{\mathtt{{\color{var}w}}})$.
\end{thm}

\section{\texttt{YFastTrie}: A Doubly-Logarithmic Time \texttt{SSet}}
\seclabel{yfast}

The \texttt{XFastTrie} is a vast---even exponential---improvement over the
\texttt{BinaryTrie} in terms of query time, but the \ensuremath{\mathtt{add({\color{var}x})}} and \ensuremath{\mathtt{remove({\color{var}x})}}
operations are still not terribly fast.  Furthermore, the space usage,
$O(\ensuremath{\mathtt{{\color{var}n}}}\cdot\ensuremath{\mathtt{{\color{var}w}}})$, is higher than the other \texttt{SSet} implementations
described in this book, which all use $O(\ensuremath{\mathtt{{\color{var}n}}})$ space.  These two
problems are related; if \ensuremath{\mathtt{{\color{var}n}}} \ensuremath{\mathtt{add({\color{var}x})}} operations build a structure of
size $\ensuremath{\mathtt{{\color{var}n}}}\cdot\ensuremath{\mathtt{{\color{var}w}}}$, then the \ensuremath{\mathtt{add({\color{var}x})}} operation requires at least on the
order of \ensuremath{\mathtt{{\color{var}w}}} time (and space) per operation.

\index{YFastTrie@\texttt{YFastTrie}}%
The \texttt{YFastTrie}, discussed next, simultaneously improves the space and
speed of \texttt{XFastTrie}s.  A \texttt{YFastTrie} uses an \texttt{XFastTrie}, \ensuremath{\mathtt{{\color{var}xft}}}, but only
stores $O(\ensuremath{\mathtt{{\color{var}n}}}/\ensuremath{\mathtt{{\color{var}w}}})$ values in \ensuremath{\mathtt{{\color{var}xft}}}.  In this way, the total space used by
\ensuremath{\mathtt{{\color{var}xft}}} is only $O(\ensuremath{\mathtt{{\color{var}n}}})$.  Furthermore, only one out of every \ensuremath{\mathtt{{\color{var}w}}} \ensuremath{\mathtt{add({\color{var}x})}}
or \ensuremath{\mathtt{remove({\color{var}x})}} operations in the \texttt{YFastTrie} results in an \ensuremath{\mathtt{add({\color{var}x})}} or
\ensuremath{\mathtt{remove({\color{var}x})}} operation in \ensuremath{\mathtt{{\color{var}xft}}}.  By doing this, the average cost incurred
by calls to \ensuremath{\mathtt{{\color{var}xft}}}'s \ensuremath{\mathtt{add({\color{var}x})}} and \ensuremath{\mathtt{remove({\color{var}x})}} operations is only constant.

The obvious question becomes:  If \ensuremath{\mathtt{{\color{var}xft}}} only stores \ensuremath{\mathtt{{\color{var}n}}}/\ensuremath{\mathtt{{\color{var}w}}} elements,
where do the remaining $\ensuremath{\mathtt{{\color{var}n}}}(1-1/\ensuremath{\mathtt{{\color{var}w}}})$ elements go?  These elements move 
into \emph{secondary structures},
\index{secondary structure}%
in this case an extended version of
treaps (\secref{treap}).  There are roughly \ensuremath{\mathtt{{\color{var}n}}}/\ensuremath{\mathtt{{\color{var}w}}} of these secondary
structures so, on average, each of them stores $O(\ensuremath{\mathtt{{\color{var}w}}})$ items.  Treaps
support logarithmic time \texttt{SSet} operations, so the operations on these
treaps will run in $O(\log \ensuremath{\mathtt{{\color{var}w}}})$ time, as required.

More concretely, a \texttt{YFastTrie} contains an \texttt{XFastTrie}, \ensuremath{\mathtt{{\color{var}xft}}},
that contains a random sample of the data, where each element
appears in the sample independently with probability $1/\ensuremath{\mathtt{{\color{var}w}}}$.
For convenience, the value $2^{\ensuremath{\mathtt{{\color{var}w}}}}-1$, is always contained in \ensuremath{\mathtt{{\color{var}xft}}}.
Let $\ensuremath{\mathtt{{\color{var}x}}}_0<\ensuremath{\mathtt{{\color{var}x}}}_1<\cdots<\ensuremath{\mathtt{{\color{var}x}}}_{k-1}$ denote the elements stored in \ensuremath{\mathtt{{\color{var}xft}}}.
Associated with each element, $\ensuremath{\mathtt{{\color{var}x}}}_i$, is a treap, $\ensuremath{\mathtt{{\color{var}t}}}_i$, that stores
all values in the range $\ensuremath{\mathtt{{\color{var}x}}}_{i-1}+1,\ldots,\ensuremath{\mathtt{{\color{var}x}}}_i$.  This is illustrated
in \figref{yfast}.

\begin{figure}
  \begin{center}
    \includegraphics[scale=0.90909]{figs/yfast}
  \end{center}
  \caption[A YFastTrie]{A \texttt{YFastTrie} containing the values 0, 1, 3, 4,
  6, 8, 9, 10, 11, and 13.}
  \figlabel{yfast}
\end{figure}

The \ensuremath{\mathtt{find({\color{var}x})}} operation in a \texttt{YFastTrie} is fairly easy.  We search
for \ensuremath{\mathtt{{\color{var}x}}} in \ensuremath{\mathtt{{\color{var}xft}}} and find some value $\ensuremath{\mathtt{{\color{var}x}}}_i$ associated with the treap
$\ensuremath{\mathtt{{\color{var}t}}}_i$.  We then use the treap \ensuremath{\mathtt{find({\color{var}x})}} method on $\ensuremath{\mathtt{{\color{var}t}}}_i$ to answer
the query.  The entire method is a one-liner:


 \ \vspace{.5ex}\begin{Verbatim}[gobble=1,tabsize=2,frame=single,framerule=1.2pt,commandchars=\\@\$,label=\texttt{YFastTrie},labelposition=topline]
	@\color@keyword$T$ find(@\color@keyword$T$ @\color@var$x$) {
		@\color@var$return$ @\color@var$xft$.find(@\color@keyword$new$ Pair<@\color@keyword$T$>(@\color@var$it$.intValue(@\color@var$x$))).@\color@var$t$.find(@\color@var$x$);
	}
\end{Verbatim}

The first \ensuremath{\mathtt{find({\color{var}x})}} operation (on \ensuremath{\mathtt{{\color{var}xft}}}) takes $O(\log\ensuremath{\mathtt{{\color{var}w}}})$ time.
The second \ensuremath{\mathtt{find({\color{var}x})}} operation (on a treap) takes $O(\log r)$ time, where
$r$ is the size of the treap.  Later in this section, we will show that
the expected size of the treap is $O(\ensuremath{\mathtt{{\color{var}w}}})$ so that this operation takes
$O(\log \ensuremath{\mathtt{{\color{var}w}}})$ time.\footnote{This is an application of \emph{Jensen's Inequality}: If $\E[r]=\ensuremath{\mathtt{{\color{var}w}}}$, then $\E[\log r]
\le \log w$.}

Adding an element to a \texttt{YFastTrie} is also fairly simple---most of
the time.  The \ensuremath{\mathtt{add({\color{var}x})}} method calls \ensuremath{\mathtt{{\color{var}xft}.find({\color{var}x})}} to locate the treap,
\ensuremath{\mathtt{{\color{var}t}}}, into which \ensuremath{\mathtt{{\color{var}x}}} should be inserted.  It then calls \ensuremath{\mathtt{{\color{var}t}.add({\color{var}x})}} to
add \ensuremath{\mathtt{{\color{var}x}}} to \ensuremath{\mathtt{{\color{var}t}}}.  At this point, it tosses a biased coin that comes up as
heads with probability $1/\ensuremath{\mathtt{{\color{var}w}}}$ and as tails with probability $1-1/\ensuremath{\mathtt{{\color{var}w}}}$.
If this coin comes up heads, then \ensuremath{\mathtt{{\color{var}x}}} will be added to \ensuremath{\mathtt{{\color{var}xft}}}.

This is where things get a little more complicated.  When \ensuremath{\mathtt{{\color{var}x}}} is added to
\ensuremath{\mathtt{{\color{var}xft}}}, the treap \ensuremath{\mathtt{{\color{var}t}}} needs to be split into two treaps, \ensuremath{\mathtt{{\color{var}t1}}} and \ensuremath{\mathtt{{\color{var}t}'}}.
The treap \ensuremath{\mathtt{{\color{var}t1}}} contains all the values less than or equal to \ensuremath{\mathtt{{\color{var}x}}};
\ensuremath{\mathtt{{\color{var}t}'}} is the original treap, \ensuremath{\mathtt{{\color{var}t}}}, with the elements of \ensuremath{\mathtt{{\color{var}t1}}} removed.
Once this is done, we add the pair \ensuremath{\mathtt{({\color{var}x},{\color{var}t1})}} to \ensuremath{\mathtt{{\color{var}xft}}}.  \figref{yfast-add}
shows an example.


 \ \vspace{.5ex}\begin{Verbatim}[gobble=1,tabsize=2,frame=single,framerule=1.2pt,commandchars=\\@\$,label=\texttt{YFastTrie},labelposition=topline]
	@\color@var$boolean$ add(@\color@keyword$T$ @\color@var$x$) {
		@\color@keyword$int$ @\color@var$ix$ = @\color@var$it$.intValue(@\color@var$x$);
		STreap<@\color@keyword$T$> @\color@var$t$ = @\color@var$xft$.find(@\color@keyword$new$ Pair<@\color@keyword$T$>(@\color@var$ix$)).@\color@var$t$;
		@\color@keyword$if$ (@\color@var$t$.add(@\color@var$x$)) {
			@\color@var$n$++;
			@\color@keyword$if$ (@\color@var$rand$.nextInt(@\color@var$w$) == 0) {
				STreap<@\color@keyword$T$> @\color@var$t1$ = @\color@var$t$.split(@\color@var$x$);
				@\color@var$xft$.add(@\color@keyword$new$ Pair<@\color@keyword$T$>(@\color@var$ix$, @\color@var$t1$));
			}
			@\color@var$return$ @\color@var$true$;
		} 
		@\color@var$return$ @\color@var$false$;
	}
\end{Verbatim}

\begin{figure}
  \begin{center}
    \includegraphics[scale=0.90909]{figs/yfast-add}
  \end{center}
  \caption[Adding to a YFastTrie]{Adding the values 2 and 6 to a \texttt{YFastTrie}. The coin toss
    for 6 came up heads, so 6 was added to \ensuremath{\mathtt{{\color{var}xft}}} and the treap containing
    $4,5,6,8,9$ was split.}
  \figlabel{yfast-add}
\end{figure}
Adding \ensuremath{\mathtt{{\color{var}x}}} to \ensuremath{\mathtt{{\color{var}t}}} takes $O(\log \ensuremath{\mathtt{{\color{var}w}}})$ time.  \excref{treap-split} shows
that splitting \ensuremath{\mathtt{{\color{var}t}}} into \ensuremath{\mathtt{{\color{var}t1}}} and \ensuremath{\mathtt{{\color{var}t}'}} can also be done in $O(\log \ensuremath{\mathtt{{\color{var}w}}})$
expected time.  Adding the pair (\ensuremath{\mathtt{{\color{var}x}}},\ensuremath{\mathtt{{\color{var}t1}}}) to \ensuremath{\mathtt{{\color{var}xft}}} takes $O(\ensuremath{\mathtt{{\color{var}w}}})$ time,
but only happens with probability $1/\ensuremath{\mathtt{{\color{var}w}}}$.  Therefore, the expected
running time of the \ensuremath{\mathtt{add({\color{var}x})}} operation is
\[
    O(\log\ensuremath{\mathtt{{\color{var}w}}}) + \frac{1}{\ensuremath{\mathtt{{\color{var}w}}}}O(\ensuremath{\mathtt{{\color{var}w}}}) = O(\log \ensuremath{\mathtt{{\color{var}w}}}) \enspace .
\]

The \ensuremath{\mathtt{remove({\color{var}x})}} method undoes the work performed by \ensuremath{\mathtt{add({\color{var}x})}}.
We use \ensuremath{\mathtt{{\color{var}xft}}} to find the leaf, \ensuremath{\mathtt{{\color{var}u}}}, in \ensuremath{\mathtt{{\color{var}xft}}} that contains the answer
to \ensuremath{\mathtt{{\color{var}xft}.find({\color{var}x})}}.  From \ensuremath{\mathtt{{\color{var}u}}}, we get the treap, \ensuremath{\mathtt{{\color{var}t}}}, containing \ensuremath{\mathtt{{\color{var}x}}}
and remove \ensuremath{\mathtt{{\color{var}x}}} from \ensuremath{\mathtt{{\color{var}t}}}.  If \ensuremath{\mathtt{{\color{var}x}}} was also stored in \ensuremath{\mathtt{{\color{var}xft}}} (and \ensuremath{\mathtt{{\color{var}x}}}
is not equal to $2^{\ensuremath{\mathtt{{\color{var}w}}}}-1$) then we remove \ensuremath{\mathtt{{\color{var}x}}} from \ensuremath{\mathtt{{\color{var}xft}}} and add the
elements from \ensuremath{\mathtt{{\color{var}x}}}'s treap to the treap, \ensuremath{\mathtt{{\color{var}t2}}}, that is stored by \ensuremath{\mathtt{{\color{var}u}}}'s
successor in the linked list.   This is illustrated in
\figref{yfast-remove}.


 \ \vspace{.5ex}\begin{Verbatim}[gobble=1,tabsize=2,frame=single,framerule=1.2pt,commandchars=\\@\$,label=\texttt{YFastTrie},labelposition=topline]
	@\color@var$boolean$ remove(@\color@keyword$T$ @\color@var$x$) {
		@\color@keyword$int$ @\color@var$ix$ = @\color@var$it$.intValue(@\color@var$x$);
		Node<@\color@keyword$T$> @\color@var$u$ = @\color@var$xft$.findNode(@\color@var$ix$);
		@\color@var$boolean$ @\color@var$ret$ = @\color@var$u$.@\color@var$x$.@\color@var$t$.remove(@\color@var$x$);
		@\color@keyword$if$ (@\color@var$ret$) @\color@var$n$--;
		@\color@keyword$if$ (@\color@var$u$.@\color@var$x$.@\color@var$x$ == @\color@var$ix$ && @\color@var$ix$ != 0xffffffff) {
			STreap<@\color@keyword$T$> @\color@var$t2$ = @\color@var$u$.@\color@var$child$[1].@\color@var$x$.@\color@var$t$;
			@\color@var$t2$.absorb(@\color@var$u$.@\color@var$x$.@\color@var$t$);
			@\color@var$xft$.remove(@\color@var$u$.@\color@var$x$);
		}
		@\color@var$return$ @\color@var$ret$;
	}
\end{Verbatim}

\begin{figure}
  \begin{center}
    \includegraphics[scale=0.90909]{figs/yfast-remove}
  \end{center}
  \caption[Removing from a YFastTrie]{Removing the values 1 and 9 from a \texttt{YFastTrie} in \figref{yfast-add}.}
  \figlabel{yfast-remove}
\end{figure}
Finding the node \ensuremath{\mathtt{{\color{var}u}}} in \ensuremath{\mathtt{{\color{var}xft}}} takes $O(\log\ensuremath{\mathtt{{\color{var}w}}})$ expected time.
Removing \ensuremath{\mathtt{{\color{var}x}}} from \ensuremath{\mathtt{{\color{var}t}}} takes $O(\log\ensuremath{\mathtt{{\color{var}w}}})$ expected time.  Again,
\excref{treap-split} shows that merging all the elements of \ensuremath{\mathtt{{\color{var}t}}} into
\ensuremath{\mathtt{{\color{var}t2}}} can be done in $O(\log\ensuremath{\mathtt{{\color{var}w}}})$ time.  If necessary, removing \ensuremath{\mathtt{{\color{var}x}}}
from \ensuremath{\mathtt{{\color{var}xft}}} takes $O(\ensuremath{\mathtt{{\color{var}w}}})$ time, but \ensuremath{\mathtt{{\color{var}x}}} is only contained in \ensuremath{\mathtt{{\color{var}xft}}} with
probability $1/\ensuremath{\mathtt{{\color{var}w}}}$.  Therefore, the expected time to remove an element
from a \texttt{YFastTrie} is $O(\log \ensuremath{\mathtt{{\color{var}w}}})$.

Earlier in the discussion, we delayed arguing about the sizes of treaps
in this structure until later.  Before finishing this chapter, we prove
the result we need.

\begin{lem}\lemlabel{yfast-subtreesize}
Let \ensuremath{\mathtt{{\color{var}x}}} be an integer stored in a \texttt{YFastTrie} and let $\ensuremath{\mathtt{{\color{var}n}}}_\ensuremath{\mathtt{{\color{var}x}}}$
denote the number of elements in the treap, \ensuremath{\mathtt{{\color{var}t}}}, that contains \ensuremath{\mathtt{{\color{var}x}}}.
Then $\E[\ensuremath{\mathtt{{\color{var}n}}}_\ensuremath{\mathtt{{\color{var}x}}}] \le 2\ensuremath{\mathtt{{\color{var}w}}}-1$.
\end{lem}

\begin{proof}
Refer to \figref{yfast-sample}. Let
$\ensuremath{\mathtt{{\color{var}x}}}_1<\ensuremath{\mathtt{{\color{var}x}}}_2<\cdots<\ensuremath{\mathtt{{\color{var}x}}}_i=\ensuremath{\mathtt{{\color{var}x}}}<\ensuremath{\mathtt{{\color{var}x}}}_{i+1}<\cdots<\ensuremath{\mathtt{{\color{var}x}}}_\ensuremath{\mathtt{{\color{var}n}}}$ denote
the elements stored
in the \texttt{YFastTrie}.  The treap \ensuremath{\mathtt{{\color{var}t}}} contains some elements greater than
or equal to \ensuremath{\mathtt{{\color{var}x}}}.  These are $\ensuremath{\mathtt{{\color{var}x}}}_i,\ensuremath{\mathtt{{\color{var}x}}}_{i+1},\ldots,\ensuremath{\mathtt{{\color{var}x}}}_{i+j-1}$,
where $\ensuremath{\mathtt{{\color{var}x}}}_{i+j-1}$ is the only one of these elements in which the
biased coin toss performed in the \ensuremath{\mathtt{add({\color{var}x})}} method turned up as heads.
In other words, $\E[j]$ is equal to the expected number of biased coin
tosses required to obtain the first heads.\footnote{This analysis ignores
the fact that $j$ never exceeds $\ensuremath{\mathtt{{\color{var}n}}}-i+1$.  However, this only decreases
$\E[j]$, so the upper bound still holds.}  Each coin toss is independent
and turns up as heads with probability $1/\ensuremath{\mathtt{{\color{var}w}}}$, so $\E[j]\le\ensuremath{\mathtt{{\color{var}w}}}$.
(See \lemref{coin-tosses} for an analysis of this for the case $\ensuremath{\mathtt{{\color{var}w}}}=2$.)

Similarly, the elements of \ensuremath{\mathtt{{\color{var}t}}} smaller than \ensuremath{\mathtt{{\color{var}x}}} are
$\ensuremath{\mathtt{{\color{var}x}}}_{i-1},\ldots,\ensuremath{\mathtt{{\color{var}x}}}_{i-k}$ where all these $k$ coin tosses turn up as
tails and the coin toss for $\ensuremath{\mathtt{{\color{var}x}}}_{i-k-1}$ turns up as heads.  Therefore,
$\E[k]\le\ensuremath{\mathtt{{\color{var}w}}}-1$, since this is the same coin tossing experiment considered
in the preceding paragraph, but one in which the last toss is not counted.
In summary, $\ensuremath{\mathtt{{\color{var}n}}}_\ensuremath{\mathtt{{\color{var}x}}}=j+k$, so
\[  \E[\ensuremath{\mathtt{{\color{var}n}}}_\ensuremath{\mathtt{{\color{var}x}}}] = \E[j+k] = \E[j] + \E[k] \le 2\ensuremath{\mathtt{{\color{var}w}}}-1 \enspace .  \qedhere \]
\end{proof}
\begin{figure}
  \begin{center}
    \includegraphics[width=\ScaleIfNeeded]{figs/yfast-sample}
  \end{center}
  \caption[The query time in a YFastTrie]{The number of elements in
  the treap \ensuremath{\mathtt{{\color{var}t}}} containing \ensuremath{\mathtt{{\color{var}x}}} is determined by two coin tossing
  experiments.}
  \figlabel{yfast-sample}
\end{figure}

\lemref{yfast-subtreesize} was the last piece in the proof of the
following theorem, which summarizes the performance of the \texttt{YFastTrie}:

\begin{thm}
A \texttt{YFastTrie} implements the \texttt{SSet} interface for \ensuremath{\mathtt{{\color{var}w}}}-bit integers. A
\texttt{YFastTrie} supports the operations \ensuremath{\mathtt{add({\color{var}x})}}, \ensuremath{\mathtt{remove({\color{var}x})}}, and \ensuremath{\mathtt{find({\color{var}x})}}
in $O(\log \ensuremath{\mathtt{{\color{var}w}}})$ expected time per operation.  The space used by a
\texttt{YFastTrie} that stores \ensuremath{\mathtt{{\color{var}n}}} values is $O(\ensuremath{\mathtt{{\color{var}n}}}+\ensuremath{\mathtt{{\color{var}w}}})$.
\end{thm}

The \ensuremath{\mathtt{{\color{var}w}}} term in the space requirement comes from the fact that \ensuremath{\mathtt{{\color{var}xft}}} always
stores the value $2^\ensuremath{\mathtt{{\color{var}w}}}-1$.  The implementation could be modified (at the
expense of adding some extra cases to the code) so that it is unnecessary
to store this value.  In this case, the space requirement in the theorem
becomes $O(\ensuremath{\mathtt{{\color{var}n}}})$.

\section{Discussion and Exercises}

The first data structure to provide $O(\log\ensuremath{\mathtt{{\color{var}w}}})$ time \ensuremath{\mathtt{add({\color{var}x})}},
\ensuremath{\mathtt{remove({\color{var}x})}}, and \ensuremath{\mathtt{find({\color{var}x})}} operations was proposed by van~Emde~Boas and
has since become known as the \emph{van~Emde~Boas}
\index{van Emde Boas tree}%
(or \emph{stratified})
\index{stratified tree}%
\emph{tree} \cite{e77}.  The original van~Emde~Boas structure had size
$2^{\ensuremath{\mathtt{{\color{var}w}}}}$, making it impractical for large integers.

The \texttt{XFastTrie} and \texttt{YFastTrie} data structures were discovered by
Willard \cite{w83}.  The \texttt{XFastTrie} structure is closely related
to van~Emde~Boas trees;  for instance, the hash tables in an \texttt{XFastTrie}
replace arrays in a van~Emde~Boas tree.  That is, instead of storing
the hash table \ensuremath{\mathtt{{\color{var}t}[{\color{var}i}]}}, a van~Emde~Boas tree stores an array of length
$2^{\ensuremath{\mathtt{{\color{var}i}}}}$.

Another structure for storing integers is Fredman and Willard's fusion
trees \cite{fw93}.
\index{fusion tree}%
This structure can store \ensuremath{\mathtt{{\color{var}n}}} \ensuremath{\mathtt{{\color{var}w}}}-bit integers in
$O(\ensuremath{\mathtt{{\color{var}n}}})$ space so that the \ensuremath{\mathtt{find({\color{var}x})}} operation runs in $O((\log \ensuremath{\mathtt{{\color{var}n}}})/(\log
\ensuremath{\mathtt{{\color{var}w}}}))$ time.  By using a fusion tree when $\log \ensuremath{\mathtt{{\color{var}w}}} > \sqrt{\log \ensuremath{\mathtt{{\color{var}n}}}}$ and
a \texttt{YFastTrie} when $\log \ensuremath{\mathtt{{\color{var}w}}} \le \sqrt{\log \ensuremath{\mathtt{{\color{var}n}}}}$, one obtains an $O(\ensuremath{\mathtt{{\color{var}n}}})$
space data structure that can implement the \ensuremath{\mathtt{find({\color{var}x})}} operation in
$O(\sqrt{\log \ensuremath{\mathtt{{\color{var}n}}}})$ time.  Recent lower-bound results of P\v{a}tra\c{s}cu
and Thorup \cite{pt07} show that these results are more or less optimal,
at least for structures that use only $O(\ensuremath{\mathtt{{\color{var}n}}})$ space.

\begin{exc}
  Design and implement a simplified version of a \texttt{BinaryTrie} that
  does not have a linked list or jump pointers, but for which \ensuremath{\mathtt{find({\color{var}x})}}

  still runs in $O(\ensuremath{\mathtt{{\color{var}w}}})$ time.
\end{exc}

\begin{exc}
  Design and implement a simplified implementation of an \texttt{XFastTrie}
  that doesn't use a binary trie at all. Instead, your implementation
  should store everything in a doubly-linked list and $\ensuremath{\mathtt{{\color{var}w}}}+1$
  hash tables.
\end{exc}

\begin{exc}
  We can think of a \texttt{BinaryTrie} as a structure that stores bit strings
  of length \ensuremath{\mathtt{{\color{var}w}}} in such a way that each bitstring is represented as a
  root to leaf path.  Extend this idea into an \texttt{SSet} implementation that
  stores variable-length strings and implements \ensuremath{\mathtt{add({\color{var}s})}}, \ensuremath{\mathtt{remove({\color{var}s})}},
  and \ensuremath{\mathtt{find({\color{var}s})}} in time proporitional to the length of \ensuremath{\mathtt{{\color{var}s}}}.

  \noindent Hint: Each node in your data structure should store a hash
  table that is indexed by character values.
\end{exc}

\begin{exc}
  For an integer $\ensuremath{\mathtt{{\color{var}x}}}\in\{0,\ldots2^{\ensuremath{\mathtt{{\color{var}w}}}}-1\}$, let $d(\ensuremath{\mathtt{{\color{var}x}}})$ denote
  the difference between \ensuremath{\mathtt{{\color{var}x}}} and the value returned by \ensuremath{\mathtt{find({\color{var}x})}}
  [if \ensuremath{\mathtt{find({\color{var}x})}} returns \ensuremath{\mathtt{{\color{keyword}null}}}, then define $d(\ensuremath{\mathtt{{\color{var}x}}})$ as $2^\ensuremath{\mathtt{{\color{var}w}}}$].
  For example, if \ensuremath{\mathtt{find(23)}} returns 43, then $d(23)=20$.
  \begin{enumerate}
    \item Design and implement a modified version of the \ensuremath{\mathtt{find({\color{var}x})}}
      operation in an \texttt{XFastTrie} that runs in $O(1+\log d(\ensuremath{\mathtt{{\color{var}x}}}))$
      expected time. Hint: The hash table $t[\ensuremath{\mathtt{{\color{var}w}}}]$ contains all the
      values, \ensuremath{\mathtt{{\color{var}x}}}, such that $d(\ensuremath{\mathtt{{\color{var}x}}})=0$, so that would be a good place
      to start.
    \item Design and implement a modified version of the \ensuremath{\mathtt{find({\color{var}x})}}
      operation in an  \texttt{XFastTrie} that runs in $O(1+\log\log d(\ensuremath{\mathtt{{\color{var}x}}}))$
      expected time.
  \end{enumerate}
\end{exc}



\clearpage

\clearpage
\chapter{External Memory Searching}
\chaplabel{btree}

Throughout this book, we have been using the \ensuremath{\mathtt{{\color{var}w}}}-bit word-RAM model
of computation defined in \secref{model}.   An implicit assumption of
this model is that our computer has a large enough random access memory
to store all of the data in the data structure.  In some situations,
this assumption is not valid.  There exist collections of data so large
that no computer has enough memory to store them.  In such cases, the
application must resort to storing the data on some external storage
medium such as a hard disk, a solid state disk, or even a network file
server (which has its own external storage).

\index{external storage}%
\index{external memory}%
\index{hard disk}%
\index{solid-state drive}%
Accessing an item from external storage is extremely slow.  The hard
disk attached to the computer on which this book was written has an
average access time of 19ms and the solid state drive attached to the
computer has an average access time of 0.3ms.  In contrast, the random
access memory in the computer has an average access time of less than
0.000113ms. Accessing RAM is more than 2\,500 times faster than accessing
the solid state drive and more than 160\,000 times faster than accessing
the hard drive.


These speeds are fairly typical;  accessing a random byte from RAM is
thousands of times faster than accessing a random byte from a hard disk
or solid-state drive.  Access time, however, does not tell the whole
story.  When we access a byte from a hard disk or solid state disk, an
entire \emph{block}
\index{block}%
of the disk is read.  Each of the drives attached
to the computer has a block size of 4\,096; each time we read one byte,
the drive gives us a block containing 4\,096 bytes.  If we organize our
data structure carefully, this means that each disk access could yield
4\,096 bytes that are helpful in completing whatever operation we are doing.


This is the idea behind the \emph{external memory model}
\index{external memory model}%
of computation,
illustrated schematically in \figref{em}.  In this model, the computer
has access to a large external memory in which all of the data resides.
This memory is divided into memory \emph{blocks}
\index{block}%
each containing $B$
words.  The computer also has limited internal memory on which it can
perform computations.  Transferring a block between internal memory and
external memory takes constant time.  Computations performed within the
internal memory are \emph{free}; they take no time at all.  The fact
that internal memory computations are free may seem a bit strange, but
it simply emphasizes the fact that external memory is so much slower
than RAM.

\begin{figure}
  \centering{\includegraphics[width=\ScaleIfNeeded]{figs/em}}
  \caption[The external memory model]{In the external memory model,
  accessing an individual item, \ensuremath{\mathtt{{\color{var}x}}}, in the external memory requires
  reading the entire block containing \ensuremath{\mathtt{{\color{var}x}}} into RAM.}
  \figlabel{em}
\end{figure}

In the full-blown external memory model, the size of the internal
memory is also a parameter.  However, for the data structures described
in this chapter, it is sufficient to have an internal memory of size
$O(B+\log_B \ensuremath{\mathtt{{\color{var}n}}})$.  That is, the memory needs to be capable of storing
a constant number of blocks and a recursion stack of height $O(\log_B
\ensuremath{\mathtt{{\color{var}n}}})$.  In most cases, the $O(B)$ term dominates the memory requirement.
For example, even with the relatively small value $B=32$, $B\ge \log_B
\ensuremath{\mathtt{{\color{var}n}}}$ for all $\ensuremath{\mathtt{{\color{var}n}}}\le 2^{160}$.  In decimal, $B\ge \log_B \ensuremath{\mathtt{{\color{var}n}}}$ for any
\[
\ensuremath{\mathtt{{\color{var}n}}} \le 1\,461\,501\,637\,330\,902\,918\,203\,684\,832\,716\,283\,019\,655\,932\,542\,976 \enspace 
. \]

\section{The Block Store}

\index{block store}%
\index{BlockStore@\texttt{BlockStore}}%
The notion of external memory includes a large number of possible
different devices, each of which has its own block size and is
accessed with its own collection of system calls.  To simplify the
exposition of this chapter so that we can focus on the common ideas, we
encapsulate external memory devices with an object called a \texttt{BlockStore}.
A \texttt{BlockStore} stores a collection of memory blocks, each of size $B$.
Each block is uniquely identified by its integer index.  A \texttt{BlockStore}
supports these operations:

\begin{enumerate}
  \item \ensuremath{\mathtt{readBlock({\color{var}i})}}: Return the contents of the block whose index is \ensuremath{\mathtt{{\color{var}i}}}.

  \item \ensuremath{\mathtt{writeBlock({\color{var}i},{\color{var}b})}}: Write contents of \ensuremath{\mathtt{{\color{var}b}}} to the block whose
    index is \ensuremath{\mathtt{{\color{var}i}}}.

  \item \ensuremath{\mathtt{placeBlock({\color{var}b})}}: Return a new index and store the contents of \ensuremath{\mathtt{{\color{var}b}}}
    at this index.

  \item \ensuremath{\mathtt{freeBlock({\color{var}i})}}: Free the block whose index is \ensuremath{\mathtt{{\color{var}i}}}.  This indicates
    that the contents of this block are no longer used so the external
    memory allocated by this block may be reused.
\end{enumerate}

The easiest way to imagine a \texttt{BlockStore} is to imagine it as storing a
file on disk that is partitioned into blocks, each containing $B$ bytes.
In this way, \ensuremath{\mathtt{readBlock({\color{var}i})}} and \ensuremath{\mathtt{writeBlock({\color{var}i},{\color{var}b})}} simply read and write
bytes $\ensuremath{\mathtt{{\color{var}i}}}B,\ldots,(\ensuremath{\mathtt{{\color{var}i}}}+1)B-1$ of this file.  In addition, a simple
\texttt{BlockStore} could keep a \emph{free list} of blocks that are available
for use. Blocks freed with \ensuremath{\mathtt{freeBlock({\color{var}i})}} are added to the free list.
In this way, \ensuremath{\mathtt{placeBlock({\color{var}b})}} can use a block from the free list or,
if none is available, append a new block to the end of the file.


\section{B-Trees}
\seclabel{btree}

In this section, we discuss a generalization of binary trees,
called $B$-trees, which is efficient in the external memory model.
Alternatively, $B$-trees can be viewed as the natural generalization of
2-4 trees described in \secref{twofour}. (A 2-4 tree is a special case
of a $B$-tree that we get by setting $B=2$.)

\index{B-tree@$B$-tree}%
For any integer $B\ge 2$, a \emph{$B$-tree} is a tree in which all of
the leaves have the same depth and every non-root internal node, \ensuremath{\mathtt{{\color{var}u}}},
has at least $B$ children and at most $2B$ children.  The children of \ensuremath{\mathtt{{\color{var}u}}}
are stored in an array, \ensuremath{\mathtt{{\color{var}u}.{\color{var}children}}}.  The required number of children is
relaxed at the root, which can have anywhere between 2 and $2B$ children.

If the height of a $B$-tree is $h$, then it follows that the number,
$\ell$, of leaves in the $B$-tree satisfies
\[
    2B^{h-1} \le \ell \le (2B)^{h} \enspace .
\]
Taking the logarithm of the first inequality and rearranging terms yields:
\begin{align*}
    h & \le \frac{\log \ell-1}{\log B} + 1  \\
      & \le \frac{\log \ell}{\log B} + 1 \\
      & = \log_B \ell + 1 \enspace .
\end{align*}
That is, the height of a $B$-tree is proportional to the base-$B$
logarithm of the number of leaves.

Each node, \ensuremath{\mathtt{{\color{var}u}}}, in $B$-tree stores an array of keys
$\ensuremath{\mathtt{{\color{var}u}.{\color{var}keys}}}[0],\ldots,\ensuremath{\mathtt{{\color{var}u}.{\color{var}keys}}}[2B-1]$.  If \ensuremath{\mathtt{{\color{var}u}}} is an internal node with $k$
children, then the number of keys stored at \ensuremath{\mathtt{{\color{var}u}}} is exactly $k-1$ and these
are stored in $\ensuremath{\mathtt{{\color{var}u}.{\color{var}keys}}}[0],\ldots,\ensuremath{\mathtt{{\color{var}u}.{\color{var}keys}}}[k-2]$.  The remaining $2B-k+1$
array entries in \ensuremath{\mathtt{{\color{var}u}.{\color{var}keys}}} are set to \ensuremath{\mathtt{{\color{keyword}null}}}.  If \ensuremath{\mathtt{{\color{var}u}}} is a non-root leaf
node, then \ensuremath{\mathtt{{\color{var}u}}} contains between $B-1$ and $2B-1$ keys. The keys in a
$B$-tree respect an order similar to the keys in a binary search tree.
For any node, \ensuremath{\mathtt{{\color{var}u}}}, that stores $k-1$ keys,
\[
   \ensuremath{\mathtt{{\color{var}u}.{\color{var}keys}[0]}} < \ensuremath{\mathtt{{\color{var}u}.{\color{var}keys}[1]}} < \cdots < \ensuremath{\mathtt{{\color{var}u}.{\color{var}keys}}}[k-2] \enspace .
\]
If \ensuremath{\mathtt{{\color{var}u}}} is an internal node, then for every $\ensuremath{\mathtt{{\color{var}i}}}\in\{0,\ldots,k-2\}$,
$\ensuremath{\mathtt{{\color{var}u}.{\color{var}keys}[{\color{var}i}]}}$ is larger than every key stored in the subtree rooted at
\ensuremath{\mathtt{{\color{var}u}.{\color{var}children}[{\color{var}i}]}} but smaller than every key stored in the subtree rooted
at $\ensuremath{\mathtt{{\color{var}u}.{\color{var}children}[{\color{var}i}+1]}}$.  Informally,
\[
   \ensuremath{\mathtt{{\color{var}u}.{\color{var}children}[{\color{var}i}]}} \prec \ensuremath{\mathtt{{\color{var}u}.{\color{var}keys}[{\color{var}i}]}} \prec \ensuremath{\mathtt{{\color{var}u}.{\color{var}children}[{\color{var}i}+1]}} \enspace .
\]
An example of a $B$-tree with $B=2$ is shown in \figref{btree}.

\begin{figure}
  \centering{\includegraphics[width=\ScaleIfNeeded]{figs/btree-1}}
  \caption{A $B$-tree with $B=2$.}
  \figlabel{btree}
\end{figure}

Note that the data stored in a $B$-tree node has size $O(B)$.  Therefore,
in an external memory setting, the value of $B$ in a $B$-tree is chosen
so that  a node fits into a single external memory block.  In this way,
the time it takes to perform a $B$-tree operation in the external memory
model is proportional to the number of nodes that are accessed (read or
written) by the operation.

For example, if the keys are 4 byte integers and the node indices are
also 4 bytes, then setting $B=256$ means that each node stores
\[
(4+4)\times 2B
 = 8\times512=4096
\]
bytes of data.  This would be a perfect value of $B$ for the hard disk
or solid state drive discussed in the introduction to this chaper,
which have a block size of $4096$ bytes.

The \texttt{BTree} class, which implements a $B$-tree, stores a \texttt{BlockStore},
\ensuremath{\mathtt{{\color{var}bs}}}, that stores \texttt{BTree} nodes as well as the index, \ensuremath{\mathtt{{\color{var}ri}}}, of the
root node.  As usual, an integer, \ensuremath{\mathtt{{\color{var}n}}}, is used to keep track of the number
of items in the data structure:


 \ \vspace{.5ex}\begin{Verbatim}[gobble=1,tabsize=2,frame=single,framerule=1.2pt,commandchars=\\@\$,label=\texttt{BTree},labelposition=topline]
	@\color@keyword$int$ @\color@var$n$;
	BlockStore<Node> @\color@var$bs$;
	@\color@keyword$int$ @\color@var$ri$;
\end{Verbatim}

\pcodeimport{ods/BTree.initialize(b)}

\subsection{Searching}

The implementation of the \ensuremath{\mathtt{find({\color{var}x})}} operation, which is illustrated in
\figref{btree-find}, generalizes the \ensuremath{\mathtt{find({\color{var}x})}} operation in a binary
search tree.  The search for \ensuremath{\mathtt{{\color{var}x}}} starts at the root and uses the keys
stored at a node, \ensuremath{\mathtt{{\color{var}u}}}, to determine in which of \ensuremath{\mathtt{{\color{var}u}}}'s children the search
should continue.

\begin{figure}
  \centering{\includegraphics[width=\ScaleIfNeeded]{figs/btree-2}}
  \caption[Searching in a $B$-tree]{A successful search (for the value 4)
    and an unsuccessful search (for the value 16.5) in a $B$-tree. Shaded nodes show where the value of \ensuremath{\mathtt{{\color{var}z}}} is updated during the searches.}
  \figlabel{btree-find}
\end{figure}
More specifically, at a node \ensuremath{\mathtt{{\color{var}u}}}, the search checks if \ensuremath{\mathtt{{\color{var}x}}} is stored
in \ensuremath{\mathtt{{\color{var}u}.{\color{var}keys}}}.  If so, \ensuremath{\mathtt{{\color{var}x}}} has been found and the search is complete.
Otherwise, the search finds the smallest integer, \ensuremath{\mathtt{{\color{var}i}}}, such that
$\ensuremath{\mathtt{{\color{var}u}.{\color{var}keys}[{\color{var}i}]}} > \ensuremath{\mathtt{{\color{var}x}}}$ and continues the search in the subtree rooted at
\ensuremath{\mathtt{{\color{var}u}.{\color{var}children}[{\color{var}i}]}}.  If no key in \ensuremath{\mathtt{{\color{var}u}.{\color{var}keys}}} is greater than \ensuremath{\mathtt{{\color{var}x}}}, then the
search continues in \ensuremath{\mathtt{{\color{var}u}}}'s rightmost child.  Just like binary search
trees, the algorithm keeps track of the most recently seen key, \ensuremath{\mathtt{{\color{var}z}}},
that is larger than \ensuremath{\mathtt{{\color{var}x}}}.  In case \ensuremath{\mathtt{{\color{var}x}}} is not found, \ensuremath{\mathtt{{\color{var}z}}} is returned as
the smallest value that is greater or equal to \ensuremath{\mathtt{{\color{var}x}}}.


 \ \vspace{.5ex}\begin{Verbatim}[gobble=1,tabsize=2,frame=single,framerule=1.2pt,commandchars=\\@\$,label=\texttt{BTree},labelposition=topline]
	@\color@keyword$T$ find(@\color@keyword$T$ @\color@var$x$) {
		@\color@keyword$T$ @\color@var$z$ = @\color@keyword$null$;
		@\color@keyword$int$ @\color@var$ui$ = @\color@var$ri$;
		@\color@keyword$while$ (@\color@var$ui$ >= 0) {
			Node @\color@var$u$ = @\color@var$bs$.readBlock(@\color@var$ui$);
			@\color@keyword$int$ @\color@var$i$ = findIt(@\color@var$u$.@\color@var$keys$, @\color@var$x$);
			@\color@keyword$if$ (@\color@var$i$ < 0) @\color@var$return$ @\color@var$u$.@\color@var$keys$[-(@\color@var$i$+1)]; @\color@comment$// found it$
			@\color@keyword$if$ (@\color@var$u$.@\color@var$keys$[@\color@var$i$] != @\color@keyword$null$)
				@\color@var$z$ = @\color@var$u$.@\color@var$keys$[@\color@var$i$];
			@\color@var$ui$ = @\color@var$u$.@\color@var$children$[@\color@var$i$];
		}
		@\color@var$return$ @\color@var$z$;
	}
\end{Verbatim}

Central to the \ensuremath{\mathtt{find({\color{var}x})}} method is the \ensuremath{\mathtt{findIt({\color{var}a},{\color{var}x})}} method that
searches in a \ensuremath{\mathtt{{\color{keyword}null}}}-padded sorted array, \ensuremath{\mathtt{{\color{var}a}}}, for the value \ensuremath{\mathtt{{\color{var}x}}}.
This method, illustrated in \figref{findit}, works for any array,
\ensuremath{\mathtt{{\color{var}a}}}, where $\ensuremath{\mathtt{{\color{var}a}}}[0],\ldots,\ensuremath{\mathtt{{\color{var}a}}}[k-1]$ is a sequence of keys in sorted
order and $\ensuremath{\mathtt{{\color{var}a}}}[k],\ldots,\ensuremath{\mathtt{{\color{var}a}}}[\ensuremath{\mathtt{{\color{var}a}.{\color{var}length}}}-1]$ are all set to \ensuremath{\mathtt{{\color{keyword}null}}}.
If \ensuremath{\mathtt{{\color{var}x}}} is in the array at position \ensuremath{\mathtt{{\color{var}i}}}, then \ensuremath{\mathtt{findIt({\color{var}a},{\color{var}x})}} returns
$-\ensuremath{\mathtt{{\color{var}i}}}-1$. Otherwise, it returns the smallest index, \ensuremath{\mathtt{{\color{var}i}}}, such that
$\ensuremath{\mathtt{{\color{var}a}[{\color{var}i}]}}>\ensuremath{\mathtt{{\color{var}x}}}$ or $\ensuremath{\mathtt{{\color{var}a}[{\color{var}i}]}}=\ensuremath{\mathtt{{\color{keyword}null}}}$.
\begin{figure}
  \centering{\includegraphics[scale=0.90909]{figs/findit}}
  \caption[The findIt(a,x) method]{The execution of \ensuremath{\mathtt{findIt({\color{var}a},27)}}.}
  \figlabel{findit}
\end{figure}


 \ \vspace{.5ex}\begin{Verbatim}[gobble=1,tabsize=2,frame=single,framerule=1.2pt,commandchars=\\@\$,label=\texttt{BTree},labelposition=topline]
	@\color@keyword$int$ findIt(@\color@keyword$T$[] @\color@var$a$, @\color@keyword$T$ @\color@var$x$) {
		@\color@keyword$int$ @\color@var$lo$ = 0, @\color@var$hi$ = @\color@var$a$.@\color@var$length$;
		@\color@keyword$while$ (@\color@var$hi$ != @\color@var$lo$) {
			@\color@keyword$int$ @\color@var$m$ = (@\color@var$hi$+@\color@var$lo$)/2;
			@\color@keyword$int$ @\color@var$cmp$ = @\color@var$a$[@\color@var$m$] == @\color@keyword$null$ ? -1 : compare(@\color@var$x$, @\color@var$a$[@\color@var$m$]);
			@\color@keyword$if$ (@\color@var$cmp$ < 0)
				@\color@var$hi$ = @\color@var$m$;      @\color@comment$// look in first half$
			@\color@keyword$else$ @\color@keyword$if$ (@\color@var$cmp$ > 0)
				@\color@var$lo$ = @\color@var$m$+1;    @\color@comment$// look in second half$
			@\color@keyword$else$
				@\color@var$return$ -@\color@var$m$-1; @\color@comment$// found it$
		}
		@\color@var$return$ @\color@var$lo$;
	}
\end{Verbatim}

The \ensuremath{\mathtt{findIt({\color{var}a},{\color{var}x})}} method uses a binary search 
\index{binary search}%
that halves the search
space at each step, so it runs in $O(\log(\ensuremath{\mathtt{{\color{var}a}.{\color{var}length}}}))$ time.  In our setting, $\ensuremath{\mathtt{{\color{var}a}.{\color{var}length}}}=2B$, so \ensuremath{\mathtt{findIt({\color{var}a},{\color{var}x})}} runs in $O(\log B)$ time.

We can analyze the running time of a $B$-tree \ensuremath{\mathtt{find({\color{var}x})}} operation both
in the usual word-RAM model (where every instruction counts) and in the
external memory model (where we only count the number of nodes accessed).
Since each leaf in a $B$-tree stores at least one key and the height
of a $B$-Tree with $\ell$ leaves is $O(\log_B\ell)$, the height of a
$B$-tree that stores \ensuremath{\mathtt{{\color{var}n}}} keys is $O(\log_B \ensuremath{\mathtt{{\color{var}n}}})$.  Therefore, in the
external memory model, the time taken by the \ensuremath{\mathtt{find({\color{var}x})}} operation is
$O(\log_B \ensuremath{\mathtt{{\color{var}n}}})$.  To determine the running time in the word-RAM model,
we have to account for the cost of calling \ensuremath{\mathtt{findIt({\color{var}a},{\color{var}x})}} for each node
we access, so the running time of \ensuremath{\mathtt{find({\color{var}x})}} in the word-RAM model is
\[
   O(\log_B \ensuremath{\mathtt{{\color{var}n}}})\times O(\log B) = O(\log \ensuremath{\mathtt{{\color{var}n}}}) \enspace .
\]

\subsection{Addition}

One important difference between $B$-trees and the \texttt{BinarySearchTree}
data structure from \secref{binarysearchtree} is that the nodes of a
$B$-tree do not store pointers to their parents.  The reason for this
will be explained shortly.  The lack of parent pointers means that
the \ensuremath{\mathtt{add({\color{var}x})}} and \ensuremath{\mathtt{remove({\color{var}x})}} operations on $B$-trees are most easily
implemented using recursion.

Like all balanced search trees, some form of rebalancing is required
during an \ensuremath{\mathtt{add({\color{var}x})}} operation.  In a $B$-tree, this is done by
\emph{splitting} nodes.
\index{split}%
Refer to \figref{btree-split} for what follows.
Although splitting takes place across two levels of recursion, it is
best understood as an operation that takes a node \ensuremath{\mathtt{{\color{var}u}}} containing $2B$
keys and having $2B+1$ children.  It creates a new node, \ensuremath{\mathtt{{\color{var}w}}}, that
adopts $\ensuremath{\mathtt{{\color{var}u}.{\color{var}children}}}[B],\ldots,\ensuremath{\mathtt{{\color{var}u}.{\color{var}children}}}[2B]$.  The new node \ensuremath{\mathtt{{\color{var}w}}}
also takes \ensuremath{\mathtt{{\color{var}u}}}'s $B$ largest keys, $\ensuremath{\mathtt{{\color{var}u}.{\color{var}keys}}}[B],\ldots,\ensuremath{\mathtt{{\color{var}u}.{\color{var}keys}}}[2B-1]$.
At this point, \ensuremath{\mathtt{{\color{var}u}}} has $B$ children and $B$ keys.  The extra key,
$\ensuremath{\mathtt{{\color{var}u}.{\color{var}keys}}}[B-1]$, is passed up to the parent of \ensuremath{\mathtt{{\color{var}u}}}, which also adopts \ensuremath{\mathtt{{\color{var}w}}}.

Notice that the splitting operation modifies three nodes: \ensuremath{\mathtt{{\color{var}u}}}, \ensuremath{\mathtt{{\color{var}u}}}'s
parent, and the new node, \ensuremath{\mathtt{{\color{var}w}}}.   This is why it is important that the
nodes of a $B$-tree do not maintain parent pointers.  If they did, then
the $B+1$ children adopted by \ensuremath{\mathtt{{\color{var}w}}} would all need to have their parent
pointers modified. This would increase the number of external memory
accesses from 3 to $B+4$ and would make $B$-trees much less efficient for
large values of $B$.

\begin{figure}
   \centering{\begin{tabular}{@{}l@{}}
     \includegraphics[width=\ScaleIfNeeded]{figs/btree-split-1} \\[2ex]
     \multicolumn{1}{c}{\ensuremath{\mathtt{{\color{var}u}.split()}}} \\ 
     \multicolumn{1}{c}{$\Downarrow$} \\[2ex]
     \includegraphics[width=\ScaleIfNeeded]{figs/btree-split-2} \\
   \end{tabular}}
   \caption[Splitting a $B$-tree node]{Splitting the node \ensuremath{\mathtt{{\color{var}u}}} in a
     $B$-tree ($B=3$). Notice that the key $\ensuremath{\mathtt{{\color{var}u}.{\color{var}keys}}}[2]=\mathrm{m}$
     passes from \ensuremath{\mathtt{{\color{var}u}}} to its parent.}
   \figlabel{btree-split}
\end{figure}

The \ensuremath{\mathtt{add({\color{var}x})}} method in a $B$-tree is illustrated in \figref{btree-add}.
At a high level, this method finds a leaf, \ensuremath{\mathtt{{\color{var}u}}}, at which to add the
value \ensuremath{\mathtt{{\color{var}x}}}.  If this causes \ensuremath{\mathtt{{\color{var}u}}} to become overfull (because it already
contained $B-1$ keys), then \ensuremath{\mathtt{{\color{var}u}}} is split.  If this causes \ensuremath{\mathtt{{\color{var}u}}}'s parent to
become overfull, then \ensuremath{\mathtt{{\color{var}u}}}'s parent is also split, which may cause \ensuremath{\mathtt{{\color{var}u}}}'s
grandparent to become overfull, and so on. This process continues,
moving up the tree one level at a time until reaching a node that
is not overfull or until the root is split. In the former case, the
process stops.  In the latter case, a new root is created whose two
children become the nodes obtained when the original root was split.

\begin{figure}
   \centering{\begin{tabular}{@{}l@{}}
     \includegraphics[width=\ScaleIfNeeded]{figs/btree-add-1} \\[2ex]
     \multicolumn{1}{c}{$\Downarrow$} \\[2ex]
     \includegraphics[width=\ScaleIfNeeded]{figs/btree-add-2} \\[2ex]
     \multicolumn{1}{c}{$\Downarrow$} \\[2ex]
     \includegraphics[width=\ScaleIfNeeded]{figs/btree-add-3} 
   \end{tabular}}
   \caption[Adding to a $B$-tree]{The \ensuremath{\mathtt{add({\color{var}x})}} operation in a
      \texttt{BTree}. Adding the value 21 results in two nodes being split.}
   \figlabel{btree-add}
\end{figure}

The executive summary of the \ensuremath{\mathtt{add({\color{var}x})}} method is that it walks
from the root to a leaf searching for \ensuremath{\mathtt{{\color{var}x}}}, adds \ensuremath{\mathtt{{\color{var}x}}} to this leaf, and
then walks back up to the root, splitting any overfull nodes it encounters
along the way.  With this high level view in mind, we can now delve into
the details of how this method can be implemented recursively.

The real work of \ensuremath{\mathtt{add({\color{var}x})}} is done by the \ensuremath{\mathtt{addRecursive({\color{var}x},{\color{var}ui})}} method,
which adds the value \ensuremath{\mathtt{{\color{var}x}}} to the subtree whose root, \ensuremath{\mathtt{{\color{var}u}}}, has the
identifier \ensuremath{\mathtt{{\color{var}ui}}}.  If \ensuremath{\mathtt{{\color{var}u}}} is a leaf, then \ensuremath{\mathtt{{\color{var}x}}} is simply inserted into
\ensuremath{\mathtt{{\color{var}u}.{\color{var}keys}}}.  Otherwise, \ensuremath{\mathtt{{\color{var}x}}} is added recursively into the appropriate
child, $\ensuremath{\mathtt{{\color{var}u}}}'$, of \ensuremath{\mathtt{{\color{var}u}}}.  The result of this recursive call is normally
\ensuremath{\mathtt{{\color{keyword}null}}} but may also be a reference to a newly-created node, \ensuremath{\mathtt{{\color{var}w}}}, that
was created because $\ensuremath{\mathtt{{\color{var}u}}}'$ was split.  In this case, \ensuremath{\mathtt{{\color{var}u}}} adopts \ensuremath{\mathtt{{\color{var}w}}}
and takes its first key, completing the splitting operation on $\ensuremath{\mathtt{{\color{var}u}}}'$.

After the value \ensuremath{\mathtt{{\color{var}x}}} has been added (either to \ensuremath{\mathtt{{\color{var}u}}} or to a descendant of \ensuremath{\mathtt{{\color{var}u}}}),
the \ensuremath{\mathtt{addRecursive({\color{var}x},{\color{var}ui})}} method checks to see if \ensuremath{\mathtt{{\color{var}u}}} is storing too many
(more than $2B-1$) keys.  If so, then \ensuremath{\mathtt{{\color{var}u}}} needs to be \emph{split}
with a call to the \ensuremath{\mathtt{{\color{var}u}.split()}} method.  The result of calling \ensuremath{\mathtt{{\color{var}u}.split()}}
is a new node that is used as the return value for \ensuremath{\mathtt{addRecursive({\color{var}x},{\color{var}ui})}}.


 \ \vspace{.5ex}\begin{Verbatim}[gobble=1,tabsize=2,frame=single,framerule=1.2pt,commandchars=\\@\$,label=\texttt{BTree},labelposition=topline]
	Node addRecursive(@\color@keyword$T$ @\color@var$x$, @\color@keyword$int$ @\color@var$ui$) {
		Node @\color@var$u$ = @\color@var$bs$.readBlock(@\color@var$ui$);
		@\color@keyword$int$ @\color@var$i$ = findIt(@\color@var$u$.@\color@var$keys$, @\color@var$x$);
		@\color@keyword$if$ (@\color@var$i$ < 0) @\color@keyword$throw$ @\color@keyword$new$ DuplicateValueException();
		@\color@keyword$if$ (@\color@var$u$.@\color@var$children$[@\color@var$i$] < 0) { @\color@comment$// leaf node, just add it$
			@\color@var$u$.add(@\color@var$x$, -1);
			@\color@var$bs$.writeBlock(@\color@var$u$.@\color@var$id$, @\color@var$u$);
		} @\color@keyword$else$ {
			Node @\color@var$w$ = addRecursive(@\color@var$x$, @\color@var$u$.@\color@var$children$[@\color@var$i$]);
			@\color@keyword$if$ (@\color@var$w$ != @\color@keyword$null$) {  @\color@comment$// child was split, w is new child $
				@\color@var$x$ = @\color@var$w$.remove(0);
				@\color@var$bs$.writeBlock(@\color@var$w$.@\color@var$id$, @\color@var$w$);
				@\color@var$u$.add(@\color@var$x$, @\color@var$w$.@\color@var$id$);
				@\color@var$bs$.writeBlock(@\color@var$u$.@\color@var$id$, @\color@var$u$);
			}
		}
		@\color@var$return$ @\color@var$u$.isFull() ? @\color@var$u$.split() : @\color@keyword$null$;
	}
\end{Verbatim}


The \ensuremath{\mathtt{addRecursive({\color{var}x},{\color{var}ui})}} method is a helper for the \ensuremath{\mathtt{add({\color{var}x})}} method, which
calls \ensuremath{\mathtt{addRecursive({\color{var}x},{\color{var}ri})}} to insert \ensuremath{\mathtt{{\color{var}x}}} into the root of the $B$-tree.
If \ensuremath{\mathtt{addRecursive({\color{var}x},{\color{var}ri})}} causes the root to split, then a new root is
created that takes as its children both the old root and the new node
created by the splitting of the old root.


 \ \vspace{.5ex}\begin{Verbatim}[gobble=1,tabsize=2,frame=single,framerule=1.2pt,commandchars=\\@\$,label=\texttt{BTree},labelposition=topline]
	@\color@var$boolean$ add(@\color@keyword$T$ @\color@var$x$) {
		Node @\color@var$w$;
		@\color@var$try$ {
			@\color@var$w$ = addRecursive(@\color@var$x$, @\color@var$ri$);
		} catch (DuplicateValueException @\color@var$e$) {
			@\color@var$return$ @\color@var$false$;
		}
		@\color@keyword$if$ (@\color@var$w$ != @\color@keyword$null$) {   @\color@comment$// root was split, make new root$
			Node @\color@var$newroot$ = @\color@keyword$new$ Node();
			@\color@var$x$ = @\color@var$w$.remove(0);
			@\color@var$bs$.writeBlock(@\color@var$w$.@\color@var$id$, @\color@var$w$);
			@\color@var$newroot$.@\color@var$children$[0] = @\color@var$ri$;
			@\color@var$newroot$.@\color@var$keys$[0] = @\color@var$x$;
			@\color@var$newroot$.@\color@var$children$[1] = @\color@var$w$.@\color@var$id$;
			@\color@var$ri$ = @\color@var$newroot$.@\color@var$id$;
			@\color@var$bs$.writeBlock(@\color@var$ri$, @\color@var$newroot$);
		}
		@\color@var$n$++;
		@\color@var$return$ @\color@var$true$;
	}
\end{Verbatim}


The \ensuremath{\mathtt{add({\color{var}x})}} method and its helper, \ensuremath{\mathtt{addRecursive({\color{var}x},{\color{var}ui})}}, can be analyzed
in two phases:

\begin{description}
  \item[Downward phase:]
    During the downward phase of the recursion, before \ensuremath{\mathtt{{\color{var}x}}} has been added,
    they access a sequence of \texttt{BTree} nodes and call \ensuremath{\mathtt{findIt({\color{var}a},{\color{var}x})}} on each node.
    As with the \ensuremath{\mathtt{find({\color{var}x})}} method, this takes $O(\log_B \ensuremath{\mathtt{{\color{var}n}}})$ time in the
    external memory model and $O(\log \ensuremath{\mathtt{{\color{var}n}}})$ time in the word-RAM model.
  
  \item[Upward phase:]
    During the upward phase of the recursion, after \ensuremath{\mathtt{{\color{var}x}}} has been added,
    these methods perform a sequence of at most $O(\log_B \ensuremath{\mathtt{{\color{var}n}}})$ splits.
    Each split involves only three nodes, so this phase takes $O(\log_B
    \ensuremath{\mathtt{{\color{var}n}}})$ time in the external memory model.  However, each split
    involves moving $B$ keys and children from one node to another, so
    in the word-RAM model, this takes $O(B\log \ensuremath{\mathtt{{\color{var}n}}})$ time.
\end{description}

Recall that the value of $B$ can be quite large, much larger
than even $\log \ensuremath{\mathtt{{\color{var}n}}}$.  Therefore, in the word-RAM model, adding a value
to a $B$-tree can be much slower than adding into a balanced binary
search tree.  Later, in \secref{btree-amortized}, we will show that the
situation is not quite so bad; the amortized number of split operations
done during an \ensuremath{\mathtt{add({\color{var}x})}} operation is constant.  This shows that the
(amortized) running time of the \ensuremath{\mathtt{add({\color{var}x})}} operation in the word-RAM model
is $O(B+\log \ensuremath{\mathtt{{\color{var}n}}})$.


\subsection{Removal}

The \ensuremath{\mathtt{remove({\color{var}x})}} operation in a \texttt{BTree} is, again, most easily implemented
as a recursive method.  Although the recursive implementation of
\ensuremath{\mathtt{remove({\color{var}x})}} spreads the complexity across several methods, the overall
process, which is illustrated in \figref{btree-remove-full}, is fairly
straightforward.  By shuffling keys around, removal is reduced to the
problem of removing a value, $\ensuremath{\mathtt{{\color{var}x}}}'$, from some leaf, \ensuremath{\mathtt{{\color{var}u}}}.  Removing $\ensuremath{\mathtt{{\color{var}x}}}'$
may leave \ensuremath{\mathtt{{\color{var}u}}} with less than $B-1$  keys;  this situation is called
an \emph{underflow}.
\index{underflow}%

\begin{figure}
   \centering{\begin{tabular}{@{}l@{}}
     \includegraphics[width=\ScaleIfNeeded]{figs/btree-remove-full-1} \\[2ex]
     \multicolumn{1}{c}{$\Downarrow$} \\[2ex]
     \includegraphics[width=\ScaleIfNeeded]{figs/btree-remove-full-2} \\[2ex]
     \multicolumn{1}{c}{\ensuremath{\mathtt{merge({\color{var}v},{\color{var}w})}}} \\
     \multicolumn{1}{c}{$\Downarrow$} \\[2ex]
     \includegraphics[width=\ScaleIfNeeded]{figs/btree-remove-full-3} \\[2ex]
     \multicolumn{1}{c}{\ensuremath{\mathtt{shiftLR({\color{var}w},{\color{var}v})}}} \\
     \multicolumn{1}{c}{$\Downarrow$} \\[2ex]
     \includegraphics[width=\ScaleIfNeeded]{figs/btree-remove-full-4} \\[2ex]
   \end{tabular}}
   \caption[Removing from a $B$-tree]{Removing the value 4 from a $B$-tree
     results in one merge and one borrowing operation.}
   \figlabel{btree-remove-full}
\end{figure}

When an underflow occurs, \ensuremath{\mathtt{{\color{var}u}}} either borrows keys from, or is merged with,
one of its siblings.  If \ensuremath{\mathtt{{\color{var}u}}} is merged with a sibling, then \ensuremath{\mathtt{{\color{var}u}}}'s parent
will now have one less child and one less key, which can cause \ensuremath{\mathtt{{\color{var}u}}}'s
parent to underflow; this is again corrected by borrowing or merging,
but merging may cause \ensuremath{\mathtt{{\color{var}u}}}'s grandparent to underflow.  This process
works its way back up to the root until there is no more underflow or
until the root has its last two children merged into a single child.
When the latter case occurs, the root is removed and its lone child
becomes the new root.

Next we delve into the details of how each of these steps is implemented.
The first job of the \ensuremath{\mathtt{remove({\color{var}x})}} method is to find the element \ensuremath{\mathtt{{\color{var}x}}} that
should be removed.  If \ensuremath{\mathtt{{\color{var}x}}} is found in a leaf, then \ensuremath{\mathtt{{\color{var}x}}} is removed from
this leaf.  Otherwise, if \ensuremath{\mathtt{{\color{var}x}}} is found at \ensuremath{\mathtt{{\color{var}u}.{\color{var}keys}[{\color{var}i}]}} for some internal
node, \ensuremath{\mathtt{{\color{var}u}}}, then the algorithm removes the smallest value, \ensuremath{\mathtt{{\color{var}x}'}}, in the
subtree rooted at \ensuremath{\mathtt{{\color{var}u}.{\color{var}children}[{\color{var}i}+1]}}.  The value \ensuremath{\mathtt{{\color{var}x}'}} is the smallest
value stored in the \texttt{BTree} that is greater than \ensuremath{\mathtt{{\color{var}x}}}.  The value of \ensuremath{\mathtt{{\color{var}x}'}}
is then used to replace \ensuremath{\mathtt{{\color{var}x}}} in \ensuremath{\mathtt{{\color{var}u}.{\color{var}keys}[{\color{var}i}]}}.  This process is illustrated
in \figref{btree-remove}.

\begin{figure}
   \centering{\begin{tabular}{@{}l@{}}
     \includegraphics[width=\ScaleIfNeeded]{figs/btree-remove-1} \\[2ex]
     \multicolumn{1}{c}{$\Downarrow$} \\[2ex]
     \includegraphics[width=\ScaleIfNeeded]{figs/btree-remove-2} 
   \end{tabular}}
   \caption[The remove operation in a $B$-tree] {The \ensuremath{\mathtt{remove({\color{var}x})}} operation
      in a \texttt{BTree}. To remove the value $\ensuremath{\mathtt{{\color{var}x}}}=10$ we replace it with the
      the value $\ensuremath{\mathtt{{\color{var}x}'}}=11$ and remove 11 from the leaf that contains it.}
   \figlabel{btree-remove}
\end{figure}

The \ensuremath{\mathtt{removeRecursive({\color{var}x},{\color{var}ui})}} method is a recursive implementation of the
preceding algorithm:


 \ \vspace{.5ex}\begin{Verbatim}[gobble=1,tabsize=2,frame=single,framerule=1.2pt,commandchars=\\@\$,label=\texttt{BTree},labelposition=topline]
	@\color@var$boolean$ removeRecursive(@\color@keyword$T$ @\color@var$x$, @\color@keyword$int$ @\color@var$ui$) {
		@\color@keyword$if$ (@\color@var$ui$ < 0) @\color@var$return$ @\color@var$false$;  @\color@comment$// didn't find it$
		Node @\color@var$u$ = @\color@var$bs$.readBlock(@\color@var$ui$);
		@\color@keyword$int$ @\color@var$i$ = findIt(@\color@var$u$.@\color@var$keys$, @\color@var$x$);
		@\color@keyword$if$ (@\color@var$i$ < 0) { @\color@comment$// found it$
			@\color@var$i$ = -(@\color@var$i$+1);
			@\color@keyword$if$ (@\color@var$u$.isLeaf()) {
				@\color@var$u$.remove(@\color@var$i$);
			} @\color@keyword$else$ {
				@\color@var$u$.@\color@var$keys$[@\color@var$i$] = removeSmallest(@\color@var$u$.@\color@var$children$[@\color@var$i$+1]);
				checkUnderflow(@\color@var$u$, @\color@var$i$+1);
			}
			@\color@var$return$ @\color@var$true$;  
		} @\color@keyword$else$ @\color@keyword$if$ (removeRecursive(@\color@var$x$, @\color@var$u$.@\color@var$children$[@\color@var$i$])) {
			checkUnderflow(@\color@var$u$, @\color@var$i$);
			@\color@var$return$ @\color@var$true$;
		}
		@\color@var$return$ @\color@var$false$;
	}
	@\color@keyword$T$ removeSmallest(@\color@keyword$int$ @\color@var$ui$) {
		Node @\color@var$u$ = @\color@var$bs$.readBlock(@\color@var$ui$);
		@\color@keyword$if$ (@\color@var$u$.isLeaf()) 
			@\color@var$return$ @\color@var$u$.remove(0);
		@\color@keyword$T$ @\color@var$y$ = removeSmallest(@\color@var$u$.@\color@var$children$[0]);  
		checkUnderflow(@\color@var$u$, 0);
		@\color@var$return$ @\color@var$y$;
	}
\end{Verbatim}


Note that, after recursively removing the value \ensuremath{\mathtt{{\color{var}x}}} from the \ensuremath{\mathtt{{\color{var}i}}}th child of \ensuremath{\mathtt{{\color{var}u}}},
\ensuremath{\mathtt{removeRecursive({\color{var}x},{\color{var}ui})}} needs to ensure that this child still has at
least $B-1$ keys.  In the preceding code, this is done using a
method called \ensuremath{\mathtt{checkUnderflow({\color{var}x},{\color{var}i})}}, which checks for and corrects an
underflow in the \ensuremath{\mathtt{{\color{var}i}}}th child of \ensuremath{\mathtt{{\color{var}u}}}.  Let \ensuremath{\mathtt{{\color{var}w}}} be the \ensuremath{\mathtt{{\color{var}i}}}th child of \ensuremath{\mathtt{{\color{var}u}}}.
If \ensuremath{\mathtt{{\color{var}w}}} has only $B-2$ keys, then this needs to be fixed.  The fix
requires using a sibling of \ensuremath{\mathtt{{\color{var}w}}}.  This can be either child $\ensuremath{\mathtt{{\color{var}i}}}+1$ of
\ensuremath{\mathtt{{\color{var}u}}} or child $\ensuremath{\mathtt{{\color{var}i}}}-1$ of \ensuremath{\mathtt{{\color{var}u}}}.  We will usually use child $\ensuremath{\mathtt{{\color{var}i}}}-1$ of \ensuremath{\mathtt{{\color{var}u}}},
which is the sibling, \ensuremath{\mathtt{{\color{var}v}}}, of \ensuremath{\mathtt{{\color{var}w}}} directly to its left.  The only time
this doesn't work is when $\ensuremath{\mathtt{{\color{var}i}}}=0$, in which case we use the sibling
directly to \ensuremath{\mathtt{{\color{var}w}}}'s right.


 \ \vspace{.5ex}\begin{Verbatim}[gobble=1,tabsize=2,frame=single,framerule=1.2pt,commandchars=\\@\$,label=\texttt{BTree},labelposition=topline]
	@\color@keyword$void$ checkUnderflow(Node @\color@var$u$, @\color@keyword$int$ @\color@var$i$) {
		@\color@keyword$if$ (@\color@var$u$.@\color@var$children$[@\color@var$i$] < 0) @\color@var$return$;
		@\color@keyword$if$ (@\color@var$i$ == 0) 
			checkUnderflowZero(@\color@var$u$, @\color@var$i$); @\color@comment$// use u's right sibling$
		@\color@keyword$else$
			checkUnderflowNonZero(@\color@var$u$,@\color@var$i$);
	}
\end{Verbatim}

In the following, we focus on the case when $\ensuremath{\mathtt{{\color{var}i}}}\neq 0$ so that any
underflow at the \ensuremath{\mathtt{{\color{var}i}}}th child of \ensuremath{\mathtt{{\color{var}u}}} will be corrected with the help
of the $(\ensuremath{\mathtt{{\color{var}i}}}-1)$st child of \ensuremath{\mathtt{{\color{var}u}}}.  The case $\ensuremath{\mathtt{{\color{var}i}}}=0$ is similar and the
details can be found in the accompanying source code.

To fix an underflow at node \ensuremath{\mathtt{{\color{var}w}}}, we need to find more keys (and possibly
also children), for \ensuremath{\mathtt{{\color{var}w}}}.  There are two ways to do this:

\begin{description}
  \item[Borrowing:]
  \index{borrow}%
  If \ensuremath{\mathtt{{\color{var}w}}} has a sibling, \ensuremath{\mathtt{{\color{var}v}}}, with more than $B-1$ keys,
  then \ensuremath{\mathtt{{\color{var}w}}} can borrow some keys (and possibly also children) from \ensuremath{\mathtt{{\color{var}v}}}.
  More specifically, if \ensuremath{\mathtt{{\color{var}v}}} stores \ensuremath{\mathtt{size({\color{var}v})}} keys, then between them,
  \ensuremath{\mathtt{{\color{var}v}}} and \ensuremath{\mathtt{{\color{var}w}}} have a total of
  \[
     B-2 + \ensuremath{\mathtt{size({\color{var}w})}} \ge 2B-2
  \]
  keys.  We can therefore shift keys from \ensuremath{\mathtt{{\color{var}v}}} to \ensuremath{\mathtt{{\color{var}w}}} so that each of
  \ensuremath{\mathtt{{\color{var}v}}} and \ensuremath{\mathtt{{\color{var}w}}} has at least $B-1$ keys.  This process is illustrated in
  \figref{btree-borrow}.

  \begin{figure}
      \centering{\begin{tabular}{@{}l@{}}
       \includegraphics[width=\ScaleIfNeeded]{figs/btree-borrow-1} \\[2ex]
       \multicolumn{1}{c}{\ensuremath{\mathtt{shiftRL({\color{var}v},{\color{var}w})}}} \\ 
       \multicolumn{1}{c}{$\Downarrow$} \\[2ex]
       \includegraphics[width=\ScaleIfNeeded]{figs/btree-borrow-2} \\
     \end{tabular}}
    \caption[Borrowing in a $B$-tree]{If \ensuremath{\mathtt{{\color{var}v}}} has more than $B-1$ keys,
       then \ensuremath{\mathtt{{\color{var}w}}} can borrow keys from \ensuremath{\mathtt{{\color{var}v}}}.}
    \figlabel{btree-borrow}
  \end{figure}
  
  \item[Merging:]
  \index{merge}%
  If \ensuremath{\mathtt{{\color{var}v}}} has only $B-1$ keys, we must do something more
  drastic, since \ensuremath{\mathtt{{\color{var}v}}} cannot afford to give any keys to \ensuremath{\mathtt{{\color{var}w}}}.  Therefore,
  we \emph{merge} \ensuremath{\mathtt{{\color{var}v}}} and \ensuremath{\mathtt{{\color{var}w}}} as shown in \figref{btree-merge}.  The merge
  operation is the opposite of the split operation.  It takes two nodes
  that contain a total of $2B-3$ keys and merges them into a single
  node that contains $2B-2$ keys.  (The additional key comes from the
  fact that, when we merge \ensuremath{\mathtt{{\color{var}v}}} and \ensuremath{\mathtt{{\color{var}w}}}, their common parent, \ensuremath{\mathtt{{\color{var}u}}}, now
  has one less child and therefore needs to give up one of its keys.)
  
  \begin{figure}
     \centering{\begin{tabular}{@{}l@{}}
       \includegraphics[width=\ScaleIfNeeded]{figs/btree-merge-1} \\[2ex]
       \multicolumn{1}{c}{\ensuremath{\mathtt{merge({\color{var}v},{\color{var}w})}}} \\ 
       \multicolumn{1}{c}{$\Downarrow$} \\[2ex]
       \includegraphics[width=\ScaleIfNeeded]{figs/btree-merge-2} \\
     \end{tabular}}
     \caption[Merging in a $B$-tree]{Merging two siblings \ensuremath{\mathtt{{\color{var}v}}} and \ensuremath{\mathtt{{\color{var}w}}}
     in a $B$-tree ($B=3$).}
     \figlabel{btree-merge}
  \end{figure}
\end{description}



 \ \vspace{.5ex}\begin{Verbatim}[gobble=1,tabsize=2,frame=single,framerule=1.2pt,commandchars=\\@\$,label=\texttt{BTree},labelposition=topline]
	@\color@keyword$void$ checkUnderflowNonZero(Node @\color@var$u$, @\color@keyword$int$ @\color@var$i$) {
		Node @\color@var$w$ = @\color@var$bs$.readBlock(@\color@var$u$.@\color@var$children$[@\color@var$i$]);  @\color@comment$// w is child of u$
		@\color@keyword$if$ (@\color@var$w$.size() < B-1) {  @\color@comment$// underflow at w$
			Node @\color@var$v$ = @\color@var$bs$.readBlock(@\color@var$u$.@\color@var$children$[@\color@var$i$-1]); @\color@comment$// v left of w$
			@\color@keyword$if$ (@\color@var$v$.size() > B) {  @\color@comment$// w can borrow from v$
				shiftLR(@\color@var$u$, @\color@var$i$-1, @\color@var$v$, @\color@var$w$);
			} @\color@keyword$else$ { @\color@comment$// v will absorb w$
				merge(@\color@var$u$, @\color@var$i$-1, @\color@var$v$, @\color@var$w$);
			}
		}
	}
	@\color@keyword$void$ checkUnderflowZero(Node @\color@var$u$, @\color@keyword$int$ @\color@var$i$) {
		Node @\color@var$w$ = @\color@var$bs$.readBlock(@\color@var$u$.@\color@var$children$[@\color@var$i$]); @\color@comment$// w is child of u$
		@\color@keyword$if$ (@\color@var$w$.size() < B-1) {  @\color@comment$// underflow at w$
			Node @\color@var$v$ = @\color@var$bs$.readBlock(@\color@var$u$.@\color@var$children$[@\color@var$i$+1]); @\color@comment$// v right of w$
			@\color@keyword$if$ (@\color@var$v$.size() > B) { @\color@comment$// w can borrow from v$
				shiftRL(@\color@var$u$, @\color@var$i$, @\color@var$v$, @\color@var$w$);
			} @\color@keyword$else$ { @\color@comment$// w will absorb w$
				merge(@\color@var$u$, @\color@var$i$, @\color@var$w$, @\color@var$v$);
				@\color@var$u$.@\color@var$children$[@\color@var$i$] = @\color@var$w$.@\color@var$id$;
			}
		}
	}
\end{Verbatim}


To summarize, the \ensuremath{\mathtt{remove({\color{var}x})}} method in a $B$-tree follows a root to
leaf path, removes a key \ensuremath{\mathtt{{\color{var}x}'}} from a leaf, \ensuremath{\mathtt{{\color{var}u}}}, and then performs zero
or more merge operations involving \ensuremath{\mathtt{{\color{var}u}}} and its ancestors, and performs
at most one borrowing operation.  Since each merge and borrow operation
involves modifying only three nodes, and only $O(\log_B \ensuremath{\mathtt{{\color{var}n}}})$ of these
operations occur, the entire process takes $O(\log_B \ensuremath{\mathtt{{\color{var}n}}})$ time in the
external memory model.  Again, however, each merge and borrow operation
takes $O(B)$ time in the word-RAM model, so (for now) the most we can
say about the running time required by \ensuremath{\mathtt{remove({\color{var}x})}} in the word-RAM model
is that it is $O(B\log_B \ensuremath{\mathtt{{\color{var}n}}})$.

\subsection{Amortized Analysis of $B$-Trees}
\seclabel{btree-amortized}

Thus far, we have shown that
\begin{enumerate}
  \item In the external memory model, the running time of \ensuremath{\mathtt{find({\color{var}x})}},
    \ensuremath{\mathtt{add({\color{var}x})}}, and \ensuremath{\mathtt{remove({\color{var}x})}} in a $B$-tree is $O(\log_B \ensuremath{\mathtt{{\color{var}n}}})$.
  \item In the word-RAM model, the running time of \ensuremath{\mathtt{find({\color{var}x})}} is $O(\log \ensuremath{\mathtt{{\color{var}n}}})$
    and the running time of \ensuremath{\mathtt{add({\color{var}x})}} and \ensuremath{\mathtt{remove({\color{var}x})}} is $O(B\log \ensuremath{\mathtt{{\color{var}n}}})$.
\end{enumerate}

The following lemma shows that, so far, we have overestimated the number of merge and split operations performed by $B$-trees.

\begin{lem}\lemlabel{btree-split}
  Starting with an empty $B$-tree and performing any sequence of
  $m$ \ensuremath{\mathtt{add({\color{var}x})}} and \ensuremath{\mathtt{remove({\color{var}x})}} operations results in at most $3m/2$ splits,
  merges, and borrows being performed.
\end{lem}

\begin{proof}
   The proof of this has already been sketched in
   \secref{redblack-summary} for the special case in which $B=2$.
   The lemma can be proven using a credit scheme,
   \index{credit scheme}%
   in which
  \begin{enumerate}
    \item each split, merge, or borrow operation is paid for with two
      credits, i.e., a credit is removed each time one of these operations
      occurs; and
    \item at most three credits are created during any \ensuremath{\mathtt{add({\color{var}x})}} or
      \ensuremath{\mathtt{remove({\color{var}x})}} operation.
  \end{enumerate}
  Since at most $3m$ credits are ever created and each split,
  merge, and borrow is paid for with with two credits, it follows
  that at most $3m/2$ splits, merges, and borrows are performed.
  These credits are illustrated using the \cent\ symbol in
  Figures~\ref{fig:btree-split}, \ref{fig:btree-borrow}, and
  \ref{fig:btree-merge}.

  To keep track of these credits the proof maintains the following
  \emph{credit invariant}:
  \index{credit invariant}%
  Any non-root node with $B-1$ keys stores one
  credit and any node with $2B-1$ keys stores three credits.  A node
  that stores at least $B$ keys and most $2B-2$ keys need not store
  any credits.  What remains is to show that we can maintain the credit
  invariant and satisfy properties 1 and 2, above, during each \ensuremath{\mathtt{add({\color{var}x})}}
  and \ensuremath{\mathtt{remove({\color{var}x})}} operation.

  \paragraph{Adding:}
  The \ensuremath{\mathtt{add({\color{var}x})}} method does not perform any merges or borrows, so we
  need only consider split operations that occur as a result of calls
  to \ensuremath{\mathtt{add({\color{var}x})}}.

  Each split operation occurs because a key is added to a node, \ensuremath{\mathtt{{\color{var}u}}}, that
  already contains $2B-1$ keys.  When this happens, \ensuremath{\mathtt{{\color{var}u}}} is split into two
  nodes, \ensuremath{\mathtt{{\color{var}u}'}} and \ensuremath{\mathtt{{\color{var}u}''}} having $B-1$ and $B$ keys, respectively.  Prior to
  this operation, \ensuremath{\mathtt{{\color{var}u}}} was storing $2B-1$ keys, and hence three credits.
  Two of these credits can be used to pay for the split and the other
  credit can be given to \ensuremath{\mathtt{{\color{var}u}'}} (which has $B-1$ keys) to maintain the
  credit invariant.  Therefore, we can pay for the split and maintain
  the credit invariant during any split.

  The only other modification to nodes that occur during an \ensuremath{\mathtt{add({\color{var}x})}}
  operation happens after all splits, if any, are complete.  This
  modification involves adding a new key to some node \ensuremath{\mathtt{{\color{var}u}'}}.  If, prior
  to this, \ensuremath{\mathtt{{\color{var}u}'}} had $2B-2$ children, then it now has $2B-1$ children and
  must therefore receive three credits.  These are the only credits given
  out by the \ensuremath{\mathtt{add({\color{var}x})}} method.

  \paragraph{Removing:}
  During a call to \ensuremath{\mathtt{remove({\color{var}x})}}, zero or more merges occur and are possibly
  followed by a single borrow.  Each merge occurs because two nodes,
  \ensuremath{\mathtt{{\color{var}v}}} and \ensuremath{\mathtt{{\color{var}w}}}, each of which had exactly $B-1$ keys prior to calling
  \ensuremath{\mathtt{remove({\color{var}x})}} were merged into a single node with exactly $2B-2$ keys.
  Each such merge therefore frees up two credits that can be used to
  pay for the merge.

  After any merges are performed, at most one borrow operation occurs,
  after which no further merges or borrows occur.  This borrow operation
  only occurs if we remove a key from a leaf, \ensuremath{\mathtt{{\color{var}v}}}, that has $B-1$ keys.
  The node \ensuremath{\mathtt{{\color{var}v}}} therefore has one credit, and this credit goes towards
  the cost of the borrow.  This single credit is not enough to pay for
  the borrow, so we create one credit to complete the payment.

  At this point, we have created one credit and we still need to show
  that the credit invariant can be maintained.  In the worst case,
  \ensuremath{\mathtt{{\color{var}v}}}'s sibling, \ensuremath{\mathtt{{\color{var}w}}}, has exactly $B$ keys before the borrow so that,
  afterwards, both \ensuremath{\mathtt{{\color{var}v}}} and \ensuremath{\mathtt{{\color{var}w}}} have $B-1$ keys.  This means that \ensuremath{\mathtt{{\color{var}v}}} and
  \ensuremath{\mathtt{{\color{var}w}}} each should be storing a credit when the operation is complete.
  Therefore, in this case, we create an additional two credits to give to
  \ensuremath{\mathtt{{\color{var}v}}} and \ensuremath{\mathtt{{\color{var}w}}}.  Since a borrow happens at most once during a \ensuremath{\mathtt{remove({\color{var}x})}}
  operation, this means that we create at most three credits, as required.

  If the \ensuremath{\mathtt{remove({\color{var}x})}} operation does not include a borrow operation, this
  is because it finishes by removing a key from some node that, prior
  to the operation, had $B$ or more keys.  In the worst case, this node
  had exactly $B$ keys, so that it now has $B-1$ keys and must be given
  one credit, which we create.

  In either case---whether the removal finishes with a borrow
  operation or not---at most three credits need to be created during a
  call to \ensuremath{\mathtt{remove({\color{var}x})}} to maintain the credit invariant and pay for all
  borrows and merges that occur. This completes the proof of the lemma.
\end{proof}

The purpose of \lemref{btree-split} is to show that, in the word-RAM
model the cost of splits, merges and joins during a sequence of $m$
\ensuremath{\mathtt{add({\color{var}x})}} and \ensuremath{\mathtt{remove({\color{var}x})}} operations is only $O(Bm)$.  That is, the
amortized cost per operation is only $O(B)$, so the amortized cost
of \ensuremath{\mathtt{add({\color{var}x})}} and \ensuremath{\mathtt{remove({\color{var}x})}} in the word-RAM model is $O(B+\log \ensuremath{\mathtt{{\color{var}n}}})$.
This is summarized by the following pair of theorems:

\begin{thm}[External Memory $B$-Trees]
  A \texttt{BTree} implements the \texttt{SSet} interface. In the external memory model,
  a \texttt{BTree} supports the operations \ensuremath{\mathtt{add({\color{var}x})}}, \ensuremath{\mathtt{remove({\color{var}x})}}, and \ensuremath{\mathtt{find({\color{var}x})}}
  in $O(\log_B \ensuremath{\mathtt{{\color{var}n}}})$ time per operation.
\end{thm}

\begin{thm}[Word-RAM $B$-Trees]
  A \texttt{BTree} implements the \texttt{SSet} interface. In the word-RAM model, and
  ignoring the cost of splits, merges, and borrows, a \texttt{BTree} supports
  the operations \ensuremath{\mathtt{add({\color{var}x})}}, \ensuremath{\mathtt{remove({\color{var}x})}}, and \ensuremath{\mathtt{find({\color{var}x})}} in $O(\log \ensuremath{\mathtt{{\color{var}n}}})$
  time per operation.
  Furthermore, beginning with an empty \texttt{BTree}, any sequence of $m$
  \ensuremath{\mathtt{add({\color{var}x})}} and \ensuremath{\mathtt{remove({\color{var}x})}} operations results in a total of $O(Bm)$
  time spent performing splits, merges, and borrows.
\end{thm}

\section{Discussion and Exercises}

The external memory model of computation was introduced by Aggarwal and
Vitter \cite{av88}.  It is sometimes also called the \emph{I/O model}
\index{I/O model}%
or the \emph{disk access model}. 
\index{disk access model}%

$B$-Trees are to external memory searching what binary search trees
are to internal memory searching.  $B$-trees were introduced by Bayer
and McCreight \cite{bm70} in 1970 and, less than ten years later, the
title of Comer's ACM Computing Surveys article referred to
them as ubiquitous \cite{c79}.

Like binary search trees, there are many variants of $B$-Trees, including
$B^+$-trees,
\index{B+-tree@$B^+$-tree}%
$B^*$-trees,
\index{B*-tree@$B^*$-tree}%
and counted $B$-trees.
\index{conted $B$-tree}%
$B$-trees are indeed
ubiquitous and are the primary data structure in many file systems,
including Apple's HFS+,
\index{HFS+}%
Microsoft's NTFS, 
\index{NTFS}%
and Linux's Ext4;
\index{Ext4}%
every
major database system; and key-value stores used in cloud computing.
Graefe's recent survey \cite{g10} provides a 200+ page overview of the
many modern applications, variants, and optimizations of $B$-trees.

$B$-trees implement the \texttt{SSet} interface.  If only the \texttt{USet} interface
is needed, then external memory hashing
\index{external memory hashing}%
could be used as an alternative
to $B$-trees.  External memory hashing schemes do exist; see, for example,
Jensen and Pagh~\cite{jp08}.  These schemes implement the \texttt{USet} operations
in $O(1)$ expected time in the external memory model. However, for a
variety of reasons, many applications still use $B$-trees even though
they only require \texttt{USet} operations.

One reason $B$-trees are such a popular choice is that they often perform
better than their $O(\log_B \ensuremath{\mathtt{{\color{var}n}}})$ running time bounds suggest. The
reason for this is that, in external memory settings, the value of $B$
is typically quite large---in the hundreds or even thousands.  This means
that 99\% or even 99.9\% of the data in a $B$-tree is stored in the
leaves.  In a database system with a large memory, it may be possible
to cache all the internal nodes of a $B$-tree in RAM, since they only
represent 1\% or 0.1\% of the total data set.  When this happens,
this means that a search in a $B$-tree involves a very fast search in
RAM, through the internal nodes, followed by a single external memory
access to retrieve a leaf.

\begin{exc}
  Show what happens when the keys 1.5 and then 7.5 are added to the
  $B$-tree in \figref{btree}.
\end{exc}

\begin{exc}
  Show what happens when the keys 3 and then 4 are removed from the
  $B$-tree in \figref{btree}.
\end{exc}

\begin{exc}
  What is the maximum number of internal nodes in a $B$-tree that stores
  \ensuremath{\mathtt{{\color{var}n}}} keys (as a function of \ensuremath{\mathtt{{\color{var}n}}} and $B$)?
\end{exc}

\begin{exc}
  The introduction to this chapter claims that $B$-trees only need an
  internal memory of size $O(B+\log_B\ensuremath{\mathtt{{\color{var}n}}})$.  However, the implementation
  given here actually requires more memory.
  \begin{enumerate}
    \item Show that the implementation of the \ensuremath{\mathtt{add({\color{var}x})}} and \ensuremath{\mathtt{remove({\color{var}x})}}
      methods given in this chapter use an internal memory
      proportional to $B\log_B \ensuremath{\mathtt{{\color{var}n}}}$.
    \item Describe how these methods could be modified in order to reduce their memory
      consumption to $O(B + \log_B \ensuremath{\mathtt{{\color{var}n}}})$.
  \end{enumerate}
\end{exc}

\begin{exc}
  Draw the credits used in the proof of \lemref{btree-split} on the trees
  in Figures~\ref{fig:btree-add} and \ref{fig:btree-remove-full}.  Verify
  that (with three additional credits) it is possible to pay for the splits, 
  merges, and borrows and maintain the credit invariant.
\end{exc}

\begin{exc}
  Design a modified version of a $B$-tree in which nodes can have anywhere
  from $B$ up to $3B$ children (and hence $B-1$ up to $3B-1$ keys).
  Show that this new version of $B$-trees performs only $O(m/B)$ splits,
  merges, and borrows during a sequence of $m$ operations.  (Hint:
  For this to work, you will have to be more agressive with merging,
  sometimes merging two nodes before it is strictly necessary.)
\end{exc}

\begin{exc}
  In this exercise, you will design a modified method of splitting and
  merging in $B$-trees that asymptotically reduces the number of splits,
  borrows and merges by considering up to three nodes at a time.
  \begin{enumerate}
    \item Let \ensuremath{\mathtt{{\color{var}u}}} be an overfull node and let \ensuremath{\mathtt{{\color{var}v}}} be a sibling immediately
    to the right of \ensuremath{\mathtt{{\color{var}u}}}.  There are two ways to fix the overflow at \ensuremath{\mathtt{{\color{var}u}}}:
    \begin{enumerate}
       \item \ensuremath{\mathtt{{\color{var}u}}} can give some of its keys to \ensuremath{\mathtt{{\color{var}v}}}; or
       \item \ensuremath{\mathtt{{\color{var}u}}} can be split and the keys of \ensuremath{\mathtt{{\color{var}u}}} and \ensuremath{\mathtt{{\color{var}v}}} can be evenly
        distributed among \ensuremath{\mathtt{{\color{var}u}}}, \ensuremath{\mathtt{{\color{var}v}}}, and the newly created node, \ensuremath{\mathtt{{\color{var}w}}}.
    \end{enumerate}
    Show that this can always be done in such a way that, after the
    operation, each of the (at most 3) affected nodes has at least
    $B+\alpha B$ keys and at most $2B-\alpha B$ keys, for some constant
    $\alpha > 0$.
    \item Let \ensuremath{\mathtt{{\color{var}u}}} be an underfull node and let \ensuremath{\mathtt{{\color{var}v}}} and \ensuremath{\mathtt{{\color{var}w}}} be siblings of \ensuremath{\mathtt{{\color{var}u}}}
    There are two ways to fix the underflow at \ensuremath{\mathtt{{\color{var}u}}}:
    \begin{enumerate}
       \item keys can be redistributed among \ensuremath{\mathtt{{\color{var}u}}}, \ensuremath{\mathtt{{\color{var}v}}}, and \ensuremath{\mathtt{{\color{var}w}}}; or
       \item \ensuremath{\mathtt{{\color{var}u}}}, \ensuremath{\mathtt{{\color{var}v}}}, and \ensuremath{\mathtt{{\color{var}w}}} can be merged into two nodes and the keys
        of \ensuremath{\mathtt{{\color{var}u}}}, \ensuremath{\mathtt{{\color{var}v}}}, and \ensuremath{\mathtt{{\color{var}w}}} can be redistributed amongst these nodes.
    \end{enumerate}
    Show that this can always be done in such a way that, after the
    operation, each of the (at most 3) affected nodes has at least
    $B+\alpha B$ keys and at most $2B-\alpha B$ keys, for some constant
    $\alpha > 0$.
    \item Show that, with these modifications, the number of 
      merges, borrows, and splits that occur during $m$ operations is $O(m/B)$.
  \end{enumerate}
\end{exc}


\begin{exc}
  A $B^+$-tree, illustrated in \figref{bplustree} stores every key in a
  leaf and keeps its leaves stored as a doubly-linked list.  As usual,
  each leaf stores between $B-1$ and $2B-1$ keys.  Above this list is
  a standard $B$-tree that stores the largest value from each leaf but
  the last.
  \begin{enumerate}
    \item Describe fast implementations of \ensuremath{\mathtt{add({\color{var}x})}}, \ensuremath{\mathtt{remove({\color{var}x})}},
      and \ensuremath{\mathtt{find({\color{var}x})}} in a $B^+$-tree.
    \item Explain how to efficiently implement the \ensuremath{\mathtt{findRange({\color{var}x},{\color{var}y})}} method,
      that reports all values
      greater than \ensuremath{\mathtt{{\color{var}x}}} and less than or equal to \ensuremath{\mathtt{{\color{var}y}}}, in
      a $B^+$-tree.
    \item Implement a class, \texttt{BPlusTree}, that implements \ensuremath{\mathtt{find({\color{var}x})}},
      \ensuremath{\mathtt{add({\color{var}x})}}, \ensuremath{\mathtt{remove({\color{var}x})}}, and \ensuremath{\mathtt{findRange({\color{var}x},{\color{var}y})}}.
      \index{BPlusTree@\texttt{BPlusTree}}%
    \item $B^+$-trees duplicate some of the keys because they are stored
      both in the $B$-tree and in the list.  Explain why this duplication
      does not add up to much for large values of $B$.
  \end{enumerate}
\end{exc}

\begin{figure}
  \centering{\includegraphics[width=\ScaleIfNeeded]{figs/bplustree}} 
  \caption{A $B^+$-tree is a $B$-tree on top of a doubly-linked list of blocks.}
  \figlabel{bplustree}
\end{figure}



\clearpage


\fancyhead[RO]{} % section number on the outside
\fancyhead[LE]{} % section number on the outside
\renewcommand{\chaptermark}[1]{\markboth{#1}{}} 
\renewcommand{\sectionmark}[1]{\markright{#1}} 
\fancyhead[CO]{\small\nouppercase\rightmark}

\cleardoublepage
\addcontentsline{toc}{chapter}{Bibliography}
\bibliographystyle{abbrvurl}
\begin{thebibliography}{10}

\bibitem{gutenberg}
Free e{B}ooks by {P}roject {G}utenberg.
\newblock URL: \url{http://www.gutenberg.org/} [cited 2011-10-12].

\bibitem{ieee754}
{IEEE Standard for Floating-Point Arithmetic}.
\newblock Technical report, Microprocessor Standards Committee of the IEEE
  Computer Society, 3 Park Avenue, New York, NY 10016-5997, USA, August 2008.
\newblock \href {https://doi.org/10.1109/IEEESTD.2008.4610935}
  {\path{doi:10.1109/IEEESTD.2008.4610935}}.

\bibitem{avl62}
G.~Adelson-Velskii and E.~Landis.
\newblock An algorithm for the organization of information.
\newblock {\em Soviet Mathematics Doklady}, 3(1259-1262):4, 1962.

\bibitem{av88}
A.~Aggarwal and J.~S. Vitter.
\newblock The input/output complexity of sorting and related problems.
\newblock {\em Communications of the ACM}, 31(9):1116--1127, 1988.

\bibitem{a89}
A.~Andersson.
\newblock Improving partial rebuilding by using simple balance criteria.
\newblock In F.~K. H.~A. Dehne, J.-R. Sack, and N.~Santoro, editors, {\em
  Algorithms and Data Structures, Workshop WADS '89, Ottawa, Canada, August
  17--19, 1989, Proceedings}, volume 382 of {\em Lecture Notes in Computer
  Science}, pages 393--402. Springer, 1989.

\bibitem{a93}
A.~Andersson.
\newblock Balanced search trees made simple.
\newblock In F.~K. H.~A. Dehne, J.-R. Sack, N.~Santoro, and S.~Whitesides,
  editors, {\em Algorithms and Data Structures, Third Workshop, WADS '93,
  Montr{\'e}al, Canada, August 11--13, 1993, Proceedings}, volume 709 of {\em
  Lecture Notes in Computer Science}, pages 60--71. Springer, 1993.

\bibitem{a99}
A.~Andersson.
\newblock General balanced trees.
\newblock {\em Journal of Algorithms}, 30(1):1--18, 1999.

\bibitem{bbg02}
A.~Bagchi, A.~L. Buchsbaum, and M.~T. Goodrich.
\newblock Biased skip lists.
\newblock In P.~Bose and P.~Morin, editors, {\em Algorithms and Computation,
  13th International Symposium, ISAAC 2002 Vancouver, BC, Canada, November
  21--23, 2002, Proceedings}, volume 2518 of {\em Lecture Notes in Computer
  Science}, pages 1--13. Springer, 2002.

\bibitem{bm70}
R.~Bayer and E.~M. McCreight.
\newblock Organization and maintenance of large ordered indexes.
\newblock In {\em SIGFIDET Workshop}, pages 107--141. ACM, 1970.

\bibitem{hashing}
Bibliography on hashing.
\newblock URL: \url{http://liinwww.ira.uka.de/bibliography/Theory/hash.html}
  [cited 2011-07-20].

\bibitem{bhkkr99}
J.~Black, S.~Halevi, H.~Krawczyk, T.~Krovetz, and P.~Rogaway.
\newblock {UMAC}: Fast and secure message authentication.
\newblock In M.~J. Wiener, editor, {\em Advances in Cryptology - CRYPTO '99,
  19th Annual International Cryptology Conference, Santa Barbara, California,
  USA, August 15--19, 1999, Proceedings}, volume 1666 of {\em Lecture Notes in
  Computer Science}, pages 79--79. Springer, 1999.

\bibitem{bdl08}
P.~Bose, K.~Dou\"{\i}eb, and S.~Langerman.
\newblock Dynamic optimality for skip lists and b-trees.
\newblock In S.-H. Teng, editor, {\em Proceedings of the Nineteenth Annual
  ACM-SIAM Symposium on Discrete Algorithms, SODA 2008, San Francisco,
  California, USA, January 20--22, 2008}, pages 1106--1114. SIAM, 2008.

\bibitem{bcdms99}
A.~Brodnik, S.~Carlsson, E.~D. Demaine, J.~I. Munro, and R.~Sedgewick.
\newblock Resizable arrays in optimal time and space.
\newblock In Dehne et~al. \cite{wads99proc}, pages 37--48.

\bibitem{cw79}
J.~Carter and M.~Wegman.
\newblock Universal classes of hash functions.
\newblock {\em Journal of computer and system sciences}, 18(2):143--154, 1979.

\bibitem{c79}
D.~Comer.
\newblock The ubiquitous {B}-tree.
\newblock {\em ACM Computing Surveys}, 11(2):121--137, 1979.

\bibitem{c72}
C.~Crane.
\newblock Linear lists and priority queues as balanced binary trees.
\newblock Technical Report STAN-CS-72-259, Computer Science Department,
  Stanford University, 1972.

\bibitem{cw03}
S.~Crosby and D.~Wallach.
\newblock Denial of service via algorithmic complexity attacks.
\newblock In {\em Proceedings of the 12th USENIX Security Symposium}, pages
  29--44, 2003.

\bibitem{wads99proc}
F.~K. H.~A. Dehne, A.~Gupta, J.-R. Sack, and R.~Tamassia, editors.
\newblock {\em Algorithms and Data Structures, 6th International Workshop, WADS
  '99, Vancouver, British Columbia, Canada, August 11--14, 1999, Proceedings},
  volume 1663 of {\em Lecture Notes in Computer Science}. Springer, 1999.

\bibitem{d88}
L.~Devroye.
\newblock Applications of the theory of records in the study of random trees.
\newblock {\em Acta Informatica}, 26(1):123--130, 1988.

\bibitem{d90}
P.~Dietz and J.~Zhang.
\newblock Lower bounds for monotonic list labeling.
\newblock In J.~R. Gilbert and R.~G. Karlsson, editors, {\em SWAT 90, 2nd
  Scandinavian Workshop on Algorithm Theory, Bergen, Norway, July 11--14, 1990,
  Proceedings}, volume 447 of {\em Lecture Notes in Computer Science}, pages
  173--180. Springer, 1990.

\bibitem{d96}
M.~Dietzfelbinger.
\newblock Universal hashing and {$k$}-wise independent random variables via
  integer arithmetic without primes.
\newblock In C.~Puech and R.~Reischuk, editors, {\em STACS 96, 13th Annual
  Symposium on Theoretical Aspects of Computer Science, Grenoble, France,
  February 22--24, 1996, Proceedings}, volume 1046 of {\em Lecture Notes in
  Computer Science}, pages 567--580. Springer, 1996.

\bibitem{dgmp92}
M.~Dietzfelbinger, J.~Gil, Y.~Matias, and N.~Pippenger.
\newblock Polynomial hash functions are reliable.
\newblock In W.~Kuich, editor, {\em Automata, Languages and Programming, 19th
  International Colloquium, ICALP92, Vienna, Austria, July 13--17, 1992,
  Proceedings}, volume 623 of {\em Lecture Notes in Computer Science}, pages
  235--246. Springer, 1992.

\bibitem{dhkp97}
M.~Dietzfelbinger, T.~Hagerup, J.~Katajainen, and M.~Penttonen.
\newblock A reliable randomized algorithm for the closest-pair problem.
\newblock {\em Journal of Algorithms}, 25(1):19--51, 1997.

\bibitem{dkkmrt94}
M.~Dietzfelbinger, A.~R. Karlin, K.~Mehlhorn, F.~M. auf~der Heide, H.~Rohnert,
  and R.~E. Tarjan.
\newblock Dynamic perfect hashing: Upper and lower bounds.
\newblock {\em SIAM J. Comput.}, 23(4):738--761, 1994.

\bibitem{e09}
A.~Elmasry.
\newblock Pairing heaps with {$O(\log\log n)$} decrease cost.
\newblock In {\em Proceedings of the twentieth Annual ACM-SIAM Symposium on
  Discrete Algorithms}, pages 471--476. Society for Industrial and Applied
  Mathematics, 2009.

\bibitem{esss01}
F.~Ergun, S.~C. Sahinalp, J.~Sharp, and R.~Sinha.
\newblock Biased dictionaries with fast insert/deletes.
\newblock In {\em Proceedings of the thirty-third annual ACM symposium on
  Theory of computing}, pages 483--491, New York, NY, USA, 2001. ACM.

\bibitem{e1590}
M.~Eytzinger.
\newblock {\em Thesaurus principum hac aetate in Europa viventium (Cologne)}.
\newblock 1590.
\newblock In commentaries, `Eytzinger' may appear in variant forms, including:
  Aitsingeri, Aitsingero, Aitsingerum, Eyzingern.

\bibitem{f64}
R.~W. Floyd.
\newblock Algorithm 245: Treesort 3.
\newblock {\em Communications of the ACM}, 7(12):701, 1964.

\bibitem{fsst86}
M.~Fredman, R.~Sedgewick, D.~Sleator, and R.~Tarjan.
\newblock The pairing heap: A new form of self-adjusting heap.
\newblock {\em Algorithmica}, 1(1):111--129, 1986.

\bibitem{ft87}
M.~Fredman and R.~Tarjan.
\newblock Fibonacci heaps and their uses in improved network optimization
  algorithms.
\newblock {\em Journal of the ACM}, 34(3):596--615, 1987.

\bibitem{fks84}
M.~L. Fredman, J.~Koml{\'o}s, and E.~Szemer{\'e}di.
\newblock Storing a sparse table with 0 (1) worst case access time.
\newblock {\em Journal of the ACM}, 31(3):538--544, 1984.

\bibitem{fw93}
M.~L. Fredman and D.~E. Willard.
\newblock Surpassing the information theoretic bound with fusion trees.
\newblock {\em Journal of computer and system sciences}, 47(3):424--436, 1993.

\bibitem{gr93}
I.~Galperin and R.~Rivest.
\newblock Scapegoat trees.
\newblock In {\em Proceedings of the fourth annual ACM-SIAM Symposium on
  Discrete algorithms}, pages 165--174. Society for Industrial and Applied
  Mathematics, 1993.

\bibitem{gm98}
A.~Gambin and A.~Malinowski.
\newblock Randomized meldable priority queues.
\newblock In {\em SOFSEM’98: Theory and Practice of Informatics}, pages
  344--349. Springer, 1998.

\bibitem{gk99}
M.~T. Goodrich and J.~G. Kloss.
\newblock Tiered vectors: Efficient dynamic arrays for rank-based sequences.
\newblock In Dehne et~al. \cite{wads99proc}, pages 205--216.

\bibitem{g10}
G.~Graefe.
\newblock Modern b-tree techniques.
\newblock {\em Foundations and Trends in Databases}, 3(4):203--402, 2010.

\bibitem{gkp94}
R.~L. Graham, D.~E. Knuth, and O.~Patashnik.
\newblock {\em Concrete Mathematics}.
\newblock Addison-Wesley, 2nd edition, 1994.

\bibitem{gs78}
L.~Guibas and R.~Sedgewick.
\newblock A dichromatic framework for balanced trees.
\newblock In {\em 19th Annual Symposium on Foundations of Computer Science, Ann
  Arbor, Michigan, 16--18 October 1978, Proceedings}, pages 8--21. IEEE
  Computer Society, 1978.

\bibitem{h61}
C.~A.~R. Hoare.
\newblock Algorithm 64: Quicksort.
\newblock {\em Communications of the ACM}, 4(7):321, 1961.

\bibitem{ht73}
J.~E. Hopcroft and R.~E. Tarjan.
\newblock Algorithm 447: Efficient algorithms for graph manipulation.
\newblock {\em Communications of the ACM}, 16(6):372--378, 1973.

\bibitem{ht74}
J.~E. Hopcroft and R.~E. Tarjan.
\newblock Efficient planarity testing.
\newblock {\em Journal of the ACM}, 21(4):549--568, 1974.

\bibitem{hpux}
{HP-UX} process management white paper, version 1.3, 1997.
\newblock URL:
  \url{http://h21007.www2.hp.com/portal/download/files/prot/files/STK/pdfs/proc_mgt.pdf}
  [cited 2011-07-20].

\bibitem{jp08}
M.~S. Jensen and R.~Pagh.
\newblock Optimality in external memory hashing.
\newblock {\em Algorithmica}, 52(3):403--411, 2008.

\bibitem{kmp95}
P.~Kirschenhofer, C.~Martinez, and H.~Prodinger.
\newblock Analysis of an optimized search algorithm for skip lists.
\newblock {\em Theoretical Computer Science}, 144:199--220, 1995.

\bibitem{kp94}
P.~Kirschenhofer and H.~Prodinger.
\newblock The path length of random skip lists.
\newblock {\em Acta Informatica}, 31:775--792, 1994.

\bibitem{k97v1}
D.~Knuth.
\newblock {\em Fundamental Algorithms}, volume~1 of {\em The Art of Computer
  Programming}.
\newblock Addison-Wesley, third edition, 1997.

\bibitem{k97v2}
D.~Knuth.
\newblock {\em Seminumerical Algorithms}, volume~2 of {\em The Art of Computer
  Programming}.
\newblock Addison-Wesley, third edition, 1997.

\bibitem{k97v3}
D.~Knuth.
\newblock {\em Sorting and Searching}, volume~3 of {\em The Art of Computer
  Programming}.
\newblock Addison-Wesley, second edition, 1997.

\bibitem{l61}
C.~Y. Lee.
\newblock An algorithm for path connection and its applications.
\newblock {\em IRE Transaction on Electronic Computers}, EC-10(3):346--365,
  1961.

\bibitem{llm11}
E.~Lehman, F.~T. Leighton, and A.~R. Meyer.
\newblock {\em Mathematics for Computer Science}.
\newblock 2011.
\newblock URL: \url{http://people.csail.mit.edu/meyer/mcs.pdf} [cited
  2012-09-06].

\bibitem{mr98}
C.~Mart\'{\i}nez and S.~Roura.
\newblock Randomized binary search trees.
\newblock {\em Journal of the ACM}, 45(2):288--323, 1998.

\bibitem{m59}
E.~F. Moore.
\newblock The shortest path through a maze.
\newblock In {\em Proceedings of the International Symposium on the Theory of
  Switching}, pages 285--292, 1959.

\bibitem{mps92}
J.~I. Munro, T.~Papadakis, and R.~Sedgewick.
\newblock Deterministic skip lists.
\newblock In {\em Proceedings of the third annual ACM-SIAM symposium on
  Discrete algorithms (SODA'92)}, pages 367--375, Philadelphia, PA, USA, 1992.
  Society for Industrial and Applied Mathematics.

\bibitem{oracle_collections}
Oracle.
\newblock {\em The Collections Framework}.
\newblock URL:
  \url{http://download.oracle.com/javase/1.5.0/docs/guide/collections/} [cited
  2011-07-19].

\bibitem{oracle_jdk6}
Oracle.
\newblock {\em Java Platform Standard Ed. 6}.
\newblock URL: \url{http://download.oracle.com/javase/6/docs/api/} [cited
  2011-07-19].

\bibitem{oracle_tutorials}
Oracle.
\newblock {\em The Java Tutorials}.
\newblock URL: \url{http://download.oracle.com/javase/tutorial/} [cited
  2011-07-19].

\bibitem{pr04}
R.~Pagh and F.~Rodler.
\newblock Cuckoo hashing.
\newblock {\em Journal of Algorithms}, 51(2):122--144, 2004.

\bibitem{pmp92}
T.~Papadakis, J.~I. Munro, and P.~V. Poblete.
\newblock Average search and update costs in skip lists.
\newblock {\em BIT}, 32:316--332, 1992.

\bibitem{pt07}
M.~P{\v a}tra{\c s}cu and M.~Thorup.
\newblock Randomization does not help searching predecessors.
\newblock In N.~Bansal, K.~Pruhs, and C.~Stein, editors, {\em Proceedings of
  the Eighteenth Annual ACM-SIAM Symposium on Discrete Algorithms, SODA 2007,
  New Orleans, Louisiana, USA, January 7--9, 2007}, pages 555--564. SIAM, 2007.

\bibitem{pt12}
M.~P{\v a}tra{\c s}cu and M.~Thorup.
\newblock The power of simple tabulation hashing.
\newblock {\em Journal of the ACM}, 59(3):14, 2012.

\bibitem{p89}
W.~Pugh.
\newblock A skip list cookbook.
\newblock Technical report, Institute for Advanced Computer Studies, Department
  of Computer Science, University of Maryland, College Park, 1989.
\newblock URL: \url{ftp://ftp.cs.umd.edu/pub/skipLists/cookbook.pdf} [cited
  2011-07-20].

\bibitem{p91}
W.~Pugh.
\newblock Skip lists: A probabilistic alternative to balanced trees.
\newblock {\em Communications of the ACM}, 33(6):668--676, 1990.

\bibitem{redis}
Redis.
\newblock URL: \url{http://redis.io/} [cited 2011-07-20].

\bibitem{r03}
B.~Reed.
\newblock The height of a random binary search tree.
\newblock {\em Journal of the ACM}, 50(3):306--332, 2003.

\bibitem{r01}
S.~M. Ross.
\newblock {\em Probability Models for Computer Science}.
\newblock Academic Press, Inc., Orlando, FL, USA, 2001.

\bibitem{s08}
R.~Sedgewick.
\newblock Left-leaning red-black trees, September 2008.
\newblock URL: \url{http://www.cs.princeton.edu/~rs/talks/LLRB/LLRB.pdf} [cited
  2011-07-21].

\bibitem{as96}
R.~Seidel and C.~Aragon.
\newblock Randomized search trees.
\newblock {\em Algorithmica}, 16(4):464--497, 1996.

\bibitem{s54}
H.~H. Seward.
\newblock Information sorting in the application of electronic digital
  computers to business operations.
\newblock Master's thesis, Massachusetts Institute of Technology, Digital
  Computer Laboratory, 1954.

\bibitem{sra94}
Z.~Shao, J.~H. Reppy, and A.~W. Appel.
\newblock Unrolling lists.
\newblock In {\em Proceedings of the 1994 ACM conference LISP and Functional
  Programming (LFP'94)}, pages 185--195, New York, 1994. ACM.

\bibitem{s04}
P.~Sinha.
\newblock A memory-efficient doubly linked list.
\newblock {\em Linux Journal}, 129, 2005.
\newblock URL: \url{http://www.linuxjournal.com/article/6828} [cited
  2013-06-05].

\bibitem{skipdb}
Skip{DB}.
\newblock URL: \url{http://dekorte.com/projects/opensource/SkipDB/} [cited
  2011-07-20].

\bibitem{st83}
D.~Sleator and R.~Tarjan.
\newblock Self-adjusting binary trees.
\newblock In {\em Proceedings of the 15th Annual ACM Symposium on Theory of
  Computing, 25--27 April, 1983, Boston, Massachusetts, USA}, pages 235--245.
  ACM, ACM, 1983.

\bibitem{t14}
S.~P. Thompson.
\newblock {\em Calculus Made Easy}.
\newblock MacMillan, Toronto, 1914.
\newblock Project Gutenberg EBook 33283.
\newblock URL: \url{http://www.gutenberg.org/ebooks/33283} [cited 2012-06-14].

\bibitem{e77}
P.~van Emde~Boas.
\newblock Preserving order in a forest in less than logarithmic time and linear
  space.
\newblock {\em Inf. Process. Lett.}, 6(3):80--82, 1977.

\bibitem{v78}
J.~Vuillemin.
\newblock A data structure for manipulating priority queues.
\newblock {\em Communications of the ACM}, 21(4):309--315, 1978.

\bibitem{v80}
J.~Vuillemin.
\newblock A unifying look at data structures.
\newblock {\em Communications of the ACM}, 23(4):229--239, 1980.

\bibitem{w83}
D.~E. Willard.
\newblock Log-logarithmic worst-case range queries are possible in space
  {$\Theta(N)$}.
\newblock {\em Inf. Process. Lett.}, 17(2):81--84, 1983.

\bibitem{w64}
J.~Williams.
\newblock Algorithm 232: Heapsort.
\newblock {\em Communications of the ACM}, 7(6):347--348, 1964.

\end{thebibliography}


\cleardoublepage
\addcontentsline{toc}{chapter}{Index}
\begin{theindex}

  \item 9-1-1, \hyperpage{2}

  \indexspace

  \item abstract data type, \hyperxindexformat{\see{interface}}{4}
  \item adjacency list, \hyperpage{252}
  \item adjacency matrix, \hyperpage{249}
  \item algorithmic complexity attack, \hyperpage{132}
  \item amortized cost, \hyperpage{21}
  \item amortized running time, \hyperpage{20}
  \item ancestor, \hyperpage{133}
  \item array
    \subitem circular, \hyperpage{38}
  \item \texttt{ArrayDeque}, \hyperpage{40}
  \item \texttt{ArrayQueue}, \hyperpage{36}
  \item arrays, \hyperpage{29}
  \item \texttt{ArrayStack}, \hyperpage{30}
  \item asymptotic notation, \hyperpage{12}
  \item AVL tree, \hyperpage{206}

  \indexspace

  \item $B^*$-tree, \hyperpage{304}
  \item $B^+$-tree, \hyperpage{304}
  \item $B$-tree, \hyperpage{286}
  \item backing array, \hyperpage{29}
  \item \texttt{Bag}, \hyperpage{28}
  \item \texttt{BDeque}, \hyperpage{71}
  \item Bibliography on Hashing, \hyperpage{128}
  \item big-Oh notation, \hyperpage{12}
  \item binary heap, \hyperpage{211}
  \item binary logarithm, \hyperpage{10}
  \item binary search, \hyperpage{272}, \hyperpage{289}
  \item binary search tree, \hyperpage{140}
    \subitem height balanced, \hyperpage{206}
    \subitem partial rebuilding, \hyperpage{173}
    \subitem random, \hyperpage{154}
    \subitem randomized, \hyperpage{169}
    \subitem red-black, \hyperpage{185}
    \subitem size-balanced, \hyperpage{148}
    \subitem versus skiplist, \hyperpage{105}
  \item binary search tree property, \hyperpage{140}
  \item binary tree, \hyperpage{133}
    \subitem complete, \hyperpage{215}
    \subitem heap-ordered, \hyperpage{212}
    \subitem search, \hyperpage{140}
  \item binary-tree traversal, \hyperpage{136}
  \item \texttt{BinaryHeap}, \hyperpage{211}
  \item \texttt{BinarySearchTree}, \hyperpage{140}
  \item \texttt{BinaryTree}, \hyperpage{135}
  \item \texttt{BinaryTrie}, \hyperpage{266}
  \item binomial coefficients, \hyperpage{12}
  \item binomial heap, \hyperpage{222}
  \item black node, \hyperpage{190}
  \item black-height property, \hyperpage{190}
  \item block, \hyperpage{283, 284}
  \item block store, \hyperpage{285}
  \item \texttt{BlockStore}, \hyperpage{285}
  \item borrow, \hyperpage{298}
  \item bounded deque, \hyperpage{71}
  \item \texttt{BPlusTree}, \hyperpage{307}
  \item breadth-first traversal, \hyperpage{139}
  \item breadth-first-search, \hyperpage{256}

  \indexspace

  \item celebrity, \hyperxindexformat{\see{universal sink}}{263}
  \item \texttt{ChainedHashTable}, \hyperpage{107}
  \item chaining, \hyperpage{107}
  \item child, \hyperpage{133}
    \subitem left, \hyperpage{133}
    \subitem right, \hyperpage{133}
  \item circular array, \hyperpage{38}
  \item coin toss, \hyperpage{17}, \hyperpage{99}
  \item collision resolution, \hyperpage{128}
  \item colour, \hyperpage{190}
  \item \texttt  {Comparator}, \hyperpage{226}
  \item \ensuremath  {\mathtt  {compare({\color  {var}a},{\color  {var}b})}}, 
		\hyperpage{226}
  \item \ensuremath{\mathtt{compare({\color{var}x},{\color{var}y})}}, 
		\hyperpage{9}
  \item comparison tree, \hyperpage{236}
  \item comparison-based sorting, \hyperpage{226}
  \item complete binary tree, \hyperpage{215}
  \item complexity
    \subitem space, \hyperpage{20}
    \subitem time, \hyperpage{20}
  \item conflict graph, \hyperpage{247}
  \item connected components, \hyperpage{263}
  \item connected graph, \hyperpage{263}
  \item contact list, \hyperpage{1}
  \item conted $B$-tree, \hyperpage{304}
  \item correctness, \hyperpage{20}
  \item \texttt{CountdownTree}, \hyperpage{183}
  \item counting-sort, \hyperpage{239}
  \item credit invariant, \hyperpage{302}
  \item credit scheme, \hyperpage{179}, \hyperpage{302}
  \item \texttt{CubishArrayStack}, \hyperpage{61}
  \item cuckoo hashing, \hyperpage{129}
  \item cycle, \hyperpage{247}
  \item cycle detection, \hyperpage{260}

  \indexspace

  \item \texttt{DaryHeap}, \hyperpage{223}
  \item \ensuremath{\mathtt{decreaseKey({\color{var}u},{\color{var}y})}}, 
		\hyperpage{222}
  \item degree, \hyperpage{254}
  \item dependencies, \hyperpage{22}
  \item depth, \hyperpage{133}
  \item depth-first-search, \hyperpage{258}
  \item deque, \hyperpage{6}
    \subitem bounded, \hyperpage{71}
  \item descendant, \hyperpage{133}
  \item dictionary, \hyperpage{8}
  \item directed edge, \hyperpage{247}
  \item directed graph, \hyperpage{247}
  \item disk access model, \hyperpage{304}
  \item divide-and-conquer, \hyperpage{226}
  \item \texttt{DLList}, \hyperpage{67}
  \item doubly-linked list, \hyperpage{67}
  \item \texttt{DualArrayDeque}, \hyperpage{43}
  \item dummy node, \hyperpage{67}
  \item Dyck word, \hyperpage{28}
  \item \texttt{DynamiteTree}, \hyperpage{183}

  \indexspace

  \item $e$ (Euler's constant), \hyperpage{10}
  \item edge, \hyperpage{247}
  \item emergency services, \hyperpage{2}
  \item Euler's constant, \hyperpage{10}
  \item expected cost, \hyperpage{21}
  \item expected running time, \hyperpage{17}, \hyperpage{20}
  \item expected value, \hyperpage{17}
  \item exponential, \hyperpage{10}
  \item Ext4, \hyperpage{304}
  \item external memory, \hyperpage{283}
  \item external memory hashing, \hyperpage{305}
  \item external memory model, \hyperpage{284}
  \item external storage, \hyperpage{283}
  \item Eytzinger's method, \hyperpage{211}

  \indexspace

  \item factorial, \hyperpage{11}
  \item family tree, \hyperpage{147}
  \item \texttt{FastArrayStack}, \hyperpage{35}
  \item Fibonacci heap, \hyperpage{222}
  \item FIFO queue, \hyperpage{5}
  \item file system, \hyperpage{1}
  \item finger, \hyperpage{104}, \hyperpage{171}
  \item finger search
    \subitem in a skiplist, \hyperpage{104}
    \subitem in a treap, \hyperpage{171}
  \item fusion tree, \hyperpage{281}

  \indexspace

  \item general balanced tree, \hyperpage{181}
  \item \texttt{git}, \hyperpage{xii}
  \item Google, \hyperpage{3}
  \item graph, \hyperpage{247}
    \subitem connected, \hyperpage{263}
    \subitem strongly-connected, \hyperpage{263}
    \subitem undirected, \hyperpage{262}

  \indexspace

  \item $H_k$ (harmonic number), \hyperpage{154}
  \item hard disk, \hyperpage{283}
  \item harmonic number, \hyperpage{154}
  \item hash code, \hyperpage{107}, \hyperpage{122}
    \subitem for arrays, \hyperpage{125}
    \subitem for compound objects, \hyperpage{123}
    \subitem for primitive data, \hyperpage{123}
    \subitem for strings, \hyperpage{125}
  \item hash function
    \subitem perfect, \hyperpage{129}
  \item hash table, \hyperpage{107}
    \subitem cuckoo, \hyperpage{129}
    \subitem two-level, \hyperpage{129}
  \item hash value, \hyperpage{107}
  \item \ensuremath{\mathtt{hash({\color{var}x})}}, \hyperpage{107}
  \item hashing
    \subitem multiplicative, \hyperpage{110}, \hyperpage{129}
    \subitem multiply-add, \hyperpage{129}
    \subitem tabulation, \hyperpage{169}
    \subitem universal, \hyperpage{129}
  \item hashing with chaining, \hyperpage{107}, \hyperpage{128}
  \item heap, \hyperpage{211}
    \subitem binary, \hyperpage{211}
    \subitem binomial, \hyperpage{222}
    \subitem Fibonacci, \hyperpage{222}
    \subitem leftist, \hyperpage{222}
    \subitem pairing, \hyperpage{222}
    \subitem skew, \hyperpage{222}
  \item heap order, \hyperpage{212}
  \item heap property, \hyperpage{159}
  \item heap-ordered binary tree, \hyperpage{212}
  \item heap-sort, \hyperpage{233}
  \item height
    \subitem in a tree, \hyperpage{133}
    \subitem of a skiplist, \hyperpage{87}
    \subitem of a tree, \hyperpage{133}
  \item height-balanced, \hyperpage{206}
  \item HFS+, \hyperpage{304}

  \indexspace

  \item I/O model, \hyperpage{304}
  \item in-order number, \hyperpage{148}
  \item in-order traversal, \hyperpage{148}
  \item in-place algorithm, \hyperpage{243}
  \item incidence matrix, \hyperpage{262}
  \item indicator random variable, \hyperpage{17}
  \item interface, \hyperpage{4}

  \indexspace

  \item Java Collections Framework, \hyperpage{26}
  \item Java Runtime Environment, \hyperpage{60}

  \indexspace

  \item leaf, \hyperpage{133}
  \item left child, \hyperpage{133}
  \item left rotation, \hyperpage{161}
  \item left-leaning property, \hyperpage{194}
  \item left-leaning red-black tree, \hyperpage{194}
  \item leftist heap, \hyperpage{222}
  \item LIFO queue, \hyperpage{5}, 
		\hyperxindexformat{\seealso{stack}}{5}
  \item linear probing, \hyperpage{114}
  \item \texttt{LinearHashTable}, \hyperpage{114}
  \item linearity of expectation, \hyperpage{17}
  \item linked list, \hyperpage{63}
    \subitem doubly-, \hyperpage{67}
    \subitem singly-, \hyperpage{63}
    \subitem space-efficient, \hyperpage{71}
    \subitem unrolled, 
		\hyperxindexformat{\seealso{\texttt{SEList}}}{82}
  \item \texttt{List}, \hyperpage{6}
  \item logarithm, \hyperpage{10}
    \subitem binary, \hyperpage{10}
    \subitem natural, \hyperpage{10}
  \item lower-bound, \hyperpage{235}

  \indexspace

  \item map, \hyperpage{8}
  \item matched string, \hyperpage{28}
  \item \texttt{MeldableHeap}, \hyperpage{217}
  \item \ensuremath{\mathtt{memcpy({\color{var}d},{\color{var}s},{\color{var}n})}}, 
		\hyperpage{36}
  \item memory manager, \hyperpage{60}
  \item merge, \hyperpage{187}, \hyperpage{299}
  \item merge-sort, \hyperpage{84}, \hyperpage{226}
  \item min-wise independence, \hyperpage{169}
  \item \texttt{MinDeque}, \hyperpage{85}
  \item \texttt{MinQueue}, \hyperpage{85}
  \item \texttt{MinStack}, \hyperpage{85}
  \item modular arithmetic, \hyperpage{37}
  \item multiplicative hashing, \hyperpage{110}, \hyperpage{129}
  \item multiply-add hashing, \hyperpage{129}

  \indexspace

  \item \ensuremath  {\mathtt  {{\color  {var}n}}}, \hyperpage{22}
  \item natural logarithm, \hyperpage{10}
  \item no-red-edge property, \hyperpage{190}
  \item NTFS, \hyperpage{304}
  \item number
    \subitem in-order, \hyperpage{148}
    \subitem post-order, \hyperpage{148}
    \subitem pre-order, \hyperpage{148}

  \indexspace

  \item $O$ notation, \hyperpage{12}
  \item open addressing, \hyperpage{114}, \hyperpage{128}
  \item Open Source, \hyperpage{xi}
  \item ordered tree, \hyperpage{133}

  \indexspace

  \item pair, \hyperpage{8}
  \item pairing heap, \hyperpage{222}
  \item palindrome, \hyperpage{83}
  \item parent, \hyperpage{133}
  \item partial rebuilding, \hyperpage{173}
  \item path, \hyperpage{247}
  \item pedigree family tree, \hyperpage{147}, \hyperpage{222}
  \item perfect hash function, \hyperpage{129}
  \item perfect hashing, \hyperpage{129}
  \item permutation, \hyperpage{11}
    \subitem random, \hyperpage{154}
  \item pivot element, \hyperpage{230}
  \item planarity testing, \hyperpage{262}
  \item post-order number, \hyperpage{148}
  \item post-order traversal, \hyperpage{148}
  \item potential, \hyperpage{48}
  \item potential method, \hyperpage{48}, \hyperpage{80}, 
		\hyperpage{205}
  \item pre-order number, \hyperpage{148}
  \item pre-order traversal, \hyperpage{148}
  \item prime field, \hyperpage{126}
  \item priority queue, \hyperpage{5}, 
		\hyperxindexformat{\seealso{heap}}{5}
  \item probability, \hyperpage{15}

  \indexspace

  \item queue
    \subitem FIFO, \hyperpage{5}
    \subitem LIFO, \hyperpage{5}
    \subitem priority, \hyperpage{5}
  \item quicksort, \hyperpage{230}

  \indexspace

  \item radix-sort, \hyperpage{241}
  \item RAM, \hyperpage{18}
  \item random binary search tree, \hyperpage{154}
  \item random permutation, \hyperpage{154}
  \item randomization, \hyperpage{15}
  \item randomized algorithm, \hyperpage{15}
  \item randomized binary search tree, \hyperpage{169}
  \item randomized data structure, \hyperpage{15}
  \item \texttt{RandomQueue}, \hyperpage{60}
  \item reachable vertex, \hyperpage{247}
  \item recursive algorithm, \hyperpage{136}
  \item red node, \hyperpage{190}
  \item red-black tree, \hyperpage{185}, \hyperpage{194}
  \item \texttt{RedBlackTree}, \hyperpage{194}
  \item remix, \hyperpage{xi}
  \item right child, \hyperpage{133}
  \item right rotation, \hyperpage{161}
  \item rooted tree, \hyperpage{133}
  \item \texttt{RootishArrayStack}, \hyperpage{49}
  \item rotation, \hyperpage{161}
  \item run, \hyperpage{118}
  \item running time, \hyperpage{20}
    \subitem amortized, \hyperpage{20}
    \subitem expected, \hyperpage{17}, \hyperpage{20}
    \subitem worst-case, \hyperpage{20}

  \indexspace

  \item scapegoat, \hyperpage{173}
  \item \texttt{ScapegoatTree}, \hyperpage{174}
  \item search path
    \subitem in a \texttt{BinaryTrie}, \hyperpage{266}
    \subitem in a binary search tree, \hyperpage{140}
    \subitem in a skiplist, \hyperpage{88}
  \item secondary structure, \hyperpage{275}
  \item \texttt{SEList}, \hyperpage{71}
  \item sentinel node, \hyperpage{88}
  \item \texttt{Sequence}, \hyperpage{184}
  \item share, \hyperpage{xi}
  \item simple path/cycle, \hyperpage{247}
  \item singly-linked list, \hyperpage{63}
  \item size-balanced, \hyperpage{148}
  \item skew heap, \hyperpage{222}
  \item skiplist, \hyperpage{87}
    \subitem versus binary search tree, \hyperpage{105}
  \item \texttt{SkiplistList}, \hyperpage{93}
  \item \texttt{SkiplistSSet}, \hyperpage{90}
  \item \texttt{SLList}, \hyperpage{63}
  \item social network, \hyperpage{1}
  \item solid-state drive, \hyperpage{283}
  \item sorting algorithm
    \subitem comparison-based, \hyperpage{226}
  \item sorting lower-bound, \hyperpage{235}
  \item source, \hyperpage{247}
  \item space complexity, \hyperpage{20}
  \item spanning forest, \hyperpage{263}
  \item speciation event, \hyperpage{147}
  \item species tree, \hyperpage{147}
  \item split, \hyperpage{187}, \hyperpage{290}
  \item square roots, \hyperpage{56}
  \item \texttt{SSet}, \hyperpage{9}
  \item stable sorting algorithm, \hyperpage{241}
  \item stack, \hyperpage{5}
  \item \ensuremath{\mathtt{{\color{var}std}::copy({\color{var}a0},{\color{var}a1},{\color{var}b})}}, 
		\hyperpage{36}
  \item Stirling's Approximation, \hyperpage{11}
  \item stratified tree, \hyperpage{280}
  \item string
    \subitem matched, \hyperpage{28}
  \item strongly-connected graph, \hyperpage{263}
  \item successor search, \hyperpage{9}
  \item \ensuremath{\mathtt{System.arraycopy({\color{var}s},{\color{var}i},{\color{var}d},{\color{var}j},{\color{var}n})}}, 
		\hyperpage{36}

  \indexspace

  \item tabulation hashing, \hyperpage{121}, \hyperpage{169}
  \item target, \hyperpage{247}
  \item tiered-vector, \hyperpage{59}
  \item time complexity, \hyperpage{20}
  \item traversal
    \subitem breadth-first, \hyperpage{139}
    \subitem in-order, \hyperpage{148}
    \subitem of a binary tree, \hyperpage{136}
    \subitem post-order, \hyperpage{148}
    \subitem pre-order, \hyperpage{148}
  \item \texttt{Treap}, \hyperpage{159}
  \item \texttt{TreapList}, \hyperpage{172}
  \item tree, \hyperpage{133}
    \subitem $d$-ary, \hyperpage{222}
    \subitem binary, \hyperpage{133}
    \subitem ordered, \hyperpage{133}
    \subitem rooted, \hyperpage{133}
  \item tree traversal, \hyperpage{136}
  \item \texttt{Treque}, \hyperpage{60}
  \item two-level hash table, \hyperpage{129}

  \indexspace

  \item underflow, \hyperpage{295}
  \item undirected graph, \hyperpage{262}
  \item universal hashing, \hyperpage{129}
  \item universal sink, \hyperpage{263}
  \item unrolled linked list, 
		\hyperxindexformat{\seealso{\texttt{SEList}}}{82}
  \item \texttt{USet}, \hyperpage{8}

  \indexspace

  \item van Emde Boas tree, \hyperpage{280}
  \item vertex, \hyperpage{247}

  \indexspace

  \item wasted space, \hyperpage{54}
  \item web search, \hyperpage{1}
  \item \texttt{WeightBalancedTree}, \hyperpage{183}
  \item word, \hyperpage{19}
  \item word-RAM, \hyperpage{18}
  \item worst-case running time, \hyperpage{20}

  \indexspace

  \item \texttt{XFastTrie}, \hyperpage{272}
  \item XOR-list, \hyperpage{82}

  \indexspace

  \item \texttt{YFastTrie}, \hyperpage{275}

\end{theindex}


\end{document}

